%============================================================
%  Bd. 28 - Halbgruppen
%============================================================

\documentclass[main.tex]{subfiles}

\ifSubfilesClassLoaded{
  \BandSetupStandalone{B28}
}{
}

\title{Bd. 28 - Halbgruppen}
\author{Martin Kunze}
\date{}
\setcounter{file}{28}

\hypersetup{
  pdftitle={Bd. 28 - Halbgruppen},
  pdfauthor={Martin Kunze}
}

\begin{document}

\ProofWideReasonsNearFormula

\title{Bd. 28 - Halbgruppen}
\author{Martin Kunze}
\date{}
\hypersetup{pageanchor=false}
\maketitle
\hypersetup{pageanchor=true}
\setcounter{tocdepth}{3}
\tableofcontents
\setcounter{file}{28}
\renewcommand{\FormulaBandID}{28}
% Gemeinsame Deklarationen: unveränderte Aussagen, Kontexte und IDs.
% Der Hauptsatz wird nur im Fachband, die sechs Hilfsresultate nur im Beweisband aufgerufen.

\newcommand{\BracketingWordBlockBase}{%
\FormulaThmDeltaK[Induktionsanfang des Blockgesetzes]{%
  \SemigroupAlg{A}{\star}\dsep u\in\NonemptyWords{A}\dsep a\in A
  \vdash
  \WordFold{\star}{\WordConcat{u}{\LetterWord{a}}}
  =\WordFold{\star}{u}\star\WordFold{\star}{\LetterWord{a}}%
}{SemigroupWordBlockBase}{%
  \DeltaRow{Mengen}{A\dsep a\dsep u}
  \DeltaRow{Funktionen}{\star}
}
}

\newcommand{\BracketingWordBlockStep}{%
\FormulaThmDeltaK[Induktionsschritt des Blockgesetzes]{%
  \begin{aligned}[t]
    &\SemigroupAlg{A}{\star}\dsep
      u,v\in\NonemptyWords{A}\dsep a\in A\dsep
      \WordFold{\star}{\WordConcat{u}{v}}
        =\WordFold{\star}{u}\star\WordFold{\star}{v}\vdash{}\\[-2pt]
    &\WordFold{\star}{\WordConcat{u}{\WordConcat{v}{\LetterWord{a}}}}
      =\WordFold{\star}{u}\star
        \WordFold{\star}{\WordConcat{v}{\LetterWord{a}}}
  \end{aligned}%
}{SemigroupWordBlockStep}{%
  \DeltaRow{Mengen}{A\dsep a\dsep u\dsep v}
  \DeltaRow{Funktionen}{\star}
}
}

\newcommand{\BracketingWordBlockLaw}{%
\FormulaThmDeltaK[Blockgesetz für Halbgruppenprodukte]{%
  \SemigroupAlg{A}{\star}\dsep u,v\in\NonemptyWords{A}
  \vdash
  \WordFold{\star}{\WordConcat{u}{v}}
    =\WordFold{\star}{u}\star\WordFold{\star}{v}%
}{SemigroupWordBlockLaw}{%
  \DeltaRow{Mengen}{A\dsep u\dsep v}
  \DeltaRow{Funktionen}{\star}
}
}

\newcommand{\BracketingTreeNormalFormLeaf}{%
\FormulaThmDeltaK[Blattfall der Baum-Normalform]{%
  \SemigroupAlg{A}{\star}\dsep a\in A\vdash
  \TreeEvaluation{\star}{\LeafTree{A}{a}}
  =\WordFold{\star}{\TreeLeafWord{A}{\LeafTree{A}{a}}}%
}{SemigroupTreeNormalFormLeaf}{%
  \DeltaRow{Mengen}{A\dsep a}
  \DeltaRow{Funktionen}{\star}
}
}

\newcommand{\BracketingTreeNormalFormNode}{%
\FormulaThmDeltaK[Knotenschritt der Baum-Normalform]{%
  \begin{aligned}[t]
    &\SemigroupAlg{A}{\star}\dsep
      S,T\in\BracketTreeSet{A}\dsep
      \TreeEvaluation{\star}{S}=\WordFold{\star}{\TreeLeafWord{A}{S}}\dsep{}\\[-2pt]
    &\TreeEvaluation{\star}{T}=\WordFold{\star}{\TreeLeafWord{A}{T}}
      \vdash{}\\[-2pt]
    &\TreeEvaluation{\star}{\NodeTree{A}{S}{T}}
      =\WordFold{\star}{\TreeLeafWord{A}{\NodeTree{A}{S}{T}}}
  \end{aligned}%
}{SemigroupTreeNormalFormNode}{%
  \DeltaRow{Mengen}{A\dsep S\dsep T}
  \DeltaRow{Funktionen}{\star}
}
}

\newcommand{\BracketingTreeNormalForm}{%
\FormulaThmDeltaK[Baum-Normalform für Halbgruppenprodukte]{%
  \SemigroupAlg{A}{\star}\dsep T\in\BracketTreeSet{A}\vdash
  \TreeEvaluation{\star}{T}=\WordFold{\star}{\TreeLeafWord{A}{T}}%
}{SemigroupTreeNormalForm}{%
  \DeltaRow{Mengen}{A\dsep T}
  \DeltaRow{Funktionen}{\star}
}
}

\newcommand{\BracketingIndependence}{%
\FormulaThmDeltaK[Unabhängigkeit von der Klammerung]{%
  \begin{aligned}[t]
    &\SemigroupAlg{A}{\star}\dsep w\in\NonemptyWords{A}\dsep
      S,T\in\WordBracketings{A}{w}\vdash{}\\[-2pt]
    &\TreeEvaluation{\star}{S}=\TreeEvaluation{\star}{T}
  \end{aligned}%
}{SemigroupBracketingIndependence}{%
  \DeltaRow{Mengen}{A\dsep S\dsep T\dsep w}
  \DeltaRow{Funktionen}{\star}
}
}


\chapter{Einführung}

Dieser Band entwickelt Halbgruppen aus einer assoziativen binären Operation.
Behandelt werden Beispiele, Teilbarkeit, Homomorphismen, Isomorphismen und
die funktorielle Konstruktion der Potenzhalbgruppe. Zusätzliche
Strukturaxiome werden in den folgenden Bänden getrennt untersucht.

\chapter{Halbgruppen}

\section{Definition der Halbgruppe}

\FormulaDefDeltaK[Begriff der Halbgruppe]{%
  \SemigroupAlg{A}{\star}%
}{SemigroupDef}{%
  \DeltaRow{Mengen}{A\dsep x\dsep y\dsep z}
  \DeltaRow{\textbf{Axiome}}{}
  \DeltaRow{Abgeschlossenheit}{%
    x,y\in A\vdash x\star y\in A}
  \DeltaRow{Assoziativität}{%
    x,y,z\in A\vdash (x\star y)\star z=x\star(y\star z)}
  \DeltaRow{\textbf{Neues Symbol}}{}
  \DeltaPrem{Halbgruppe}{\SemigroupAlg{A}{\star}}
}

\FormulaAxiomDeltaK[Zugriff auf die Abgeschlossenheit einer Halbgruppe]{%
  \SemigroupAlg{A}{\star}\dsep x,y\in A
  \vdash x\star y\in A%
}{SemigroupClosureAxiom}{%
  \DeltaRow{Mengen}{A\dsep x\dsep y}
}

\FormulaAxiomDeltaK[Zugriff auf die Assoziativität einer Halbgruppe]{%
  \SemigroupAlg{A}{\star}\dsep x,y,z\in A
  \vdash (x\star y)\star z=x\star(y\star z)%
}{SemigroupAssociativityAxiom}{%
  \DeltaRow{Mengen}{A\dsep x\dsep y\dsep z}
}

\FormulaAxiomDeltaK[Zugriff auf die binäre Operation einer Halbgruppe]{%
  \SemigroupAlg{A}{\star}\vdash
  \mathsf{BinOp}\!\left(A,\star\right)%
}{SemigroupBinaryOperationAxiom}{%
  \DeltaRow{Mengen}{A}
}

\FormulaThmDeltaK[Fünfgliedrige Umklammerung in Halbgruppen]{%
  \begin{aligned}[t]
    &\SemigroupAlg{A}{\star}\dsep a,b,c,d,e\in A\vdash{}\\[-2pt]
    &(a\star b)\star\bigl((c\star d)\star e\bigr)
      =\bigl((a\star(b\star c))\star d\bigr)\star e
  \end{aligned}%
}{SemigroupFiveFactorReassociation}{%
  \DeltaRow{Mengen}{A\dsep a\dsep b\dsep c\dsep d\dsep e}
  \DeltaRow{Funktionen}{\star}
}
\begin{tabproofwide}
  \proofstepwidestar[1]{\SemigroupAlg{A}{\star}}{\rA}
  \proofstepwidestar[2]{a\in A}{\rA}
  \proofstepwidestar[3]{b\in A}{\rA}
  \proofstepwidestar[4]{c\in A}{\rA}
  \proofstepwidestar[5]{d\in A}{\rA}
  \proofstepwidestar[6]{e\in A}{\rA}
  \proofstepwidestar[1,2,3]{a\star b\in A}{%
    \FormulaRefAuto{SemigroupClosureAxiom}[ax]{1,2,3}}
  \proofstepwidestar[1,4,5]{c\star d\in A}{%
    \FormulaRefAuto{SemigroupClosureAxiom}[ax]{1,4,5}}
  \proofstepwidestar[1,2,3,4,5]{%
    ((a\star b)\star c)\star d=(a\star b)\star(c\star d)}{%
    \FormulaRefAuto{SemigroupAssociativityAxiom}[ax]{1,7,4,5}}
  \proofstepwidestar[1,2,3,4,5]{%
    (a\star b)\star(c\star d)=((a\star b)\star c)\star d}{%
    \FormulaRefAuto{a=b\vdash b=a}[thm]{9}}
  \proofstepwidestar[1,2,3,4,5,6]{%
    ((a\star b)\star(c\star d))\star e
      =(a\star b)\star((c\star d)\star e)}{%
    \FormulaRefAuto{SemigroupAssociativityAxiom}[ax]{1,7,8,6}}
  \proofstepwidestar[1,2,3,4,5,6]{%
    (a\star b)\star((c\star d)\star e)
      =((a\star b)\star(c\star d))\star e}{%
    \FormulaRefAuto{a=b\vdash b=a}[thm]{11}}
  \proofstepwidestar[]{%
    ((a\star b)\star(c\star d))\star e
      =((a\star b)\star(c\star d))\star e}{\rII}
  \proofstepwidestar[1,2,3,4,5]{%
    ((a\star b)\star(c\star d))\star e
      =(((a\star b)\star c)\star d)\star e}{\rIE{10,13}}
  \proofstepwidestar[1,2,3,4]{%
    (a\star b)\star c=a\star(b\star c)}{%
    \FormulaRefAuto{SemigroupAssociativityAxiom}[ax]{1,2,3,4}}
  \proofstepwidestar[]{%
    (((a\star b)\star c)\star d)\star e
      =(((a\star b)\star c)\star d)\star e}{\rII}
  \proofstepwidestar[1,2,3,4]{%
    (((a\star b)\star c)\star d)\star e
      =((a\star(b\star c))\star d)\star e}{\rIE{15,16}}
  \proofstepwidestar[1,2,3,4,5,6]{%
    (a\star b)\star((c\star d)\star e)
      =(((a\star b)\star c)\star d)\star e}{\rIE{14,12}}
  \proofstepwidestar[1,2,3,4,5,6]{%
    (a\star b)\star((c\star d)\star e)
      =((a\star(b\star c))\star d)\star e}{\rIE{17,18}}
\end{tabproofwide}

% Der Hauptsatz gehört zum Fachband, die sechs Hilfsresultate zum Beweisband.
% Abschnittsnummern und anonyme Anker der späteren Ergebnisse bleiben stabil.
\section{Produkte endlicher Wörter}

Für eine Halbgruppe \(\SemigroupAlg{A}{\star}\) bezeichnet
\(\WordFold{\star}{u}\) die in Band~27 rekursiv definierte Linksfaltung
des nichtleeren Wortes \(u\in\NonemptyWords{A}\). Sie wertet die Faktoren
von links nach rechts aus. Das
\BracketingResultLink{SemigroupWordBlockLaw}{Blockgesetz in den Beweistabellen}
zeigt, dass man auch zunächst zusammenhängende Blöcke auswerten und danach
deren Werte verknüpfen darf.

% Die drei bisherigen Satzanker reservieren, ohne Aussagen zu deklarieren.
\phantomsection\phantomsection\phantomsection

\section{Klammerungsunabhängigkeit}

Alle vollständigen Klammerungen desselben nichtleeren Wortes liefern in
einer Halbgruppe denselben Wert. Der folgende Hauptsatz hält diese Aussage
in der Sprache der Klammerungsbäume fest.

\Needspace{16\baselineskip}
% Drei ausgelagerte Hilfssatzanker; der Hauptsatz behält seine Originalnummer.
\phantomsection\phantomsection\phantomsection
\setcounter{formulaThm}{3}
\hypertarget{bracketing.main.independence}{}
\BracketingIndependence

\noindent\emph{Beweis:} Die \BracketingReadingLink{} erklärt die
Beweiskette mit Beispielen, Diagrammen und den beiden Induktionsargumenten.
Die \BracketingProofsLink{} enthalten die sechs Hilfsresultate und alle
Ableitungen einschließlich des \BracketingIndependenceProofLink{}.

\par\medskip
\hypertarget{bracketing.notation}{}

\noindent\textbf{Anmerkung zur Produktschreibweise.}
In einer Halbgruppe haben alle vollständigen Klammerungen eines endlichen,
nichtleeren Produkts denselben Wert, sofern die Reihenfolge der Faktoren
unverändert bleibt. Für jede natürliche Zahl \(n\geq 1\) und Faktoren
\(a_1,\ldots,a_n\in A\) dürfen wir deshalb einfach
\[
  a_1\star a_2\star\cdots\star a_n
\]
schreiben. Gemeint ist der gemeinsame Wert aller dieser Klammerungen,
beispielsweise der von links ausgewertete Wert. Bei einem einzigen Faktor
ist das Produkt gleich \(a_1\).

Die Klammern können somit weggelassen oder nach Bedarf gesetzt werden.
Ein beliebiges Vertauschen von Faktoren ist damit nicht erlaubt; dazu wäre zusätzlich
Kommutativität nötig. Ein leeres Produkt wird hier nicht eingeführt, da
eine Halbgruppe kein neutrales Element besitzen muss.

Die im Beweis verwendete
\BracketingResultLink{SemigroupTreeNormalForm}{Baum-Normalform}, die jede
Baumauswertung auf die Linksfaltung ihres Blattwortes zurückführt,
steht zusammen mit ihrem Beweis in den Beweistabellen.


\section{Beispiele von Halbgruppen}

\subsection{Die Potenzhalbgruppe}

Aus jeder Halbgruppe entsteht eine weitere Halbgruppe, deren Elemente die
nichtleeren Teilmengen des ursprünglichen Trägers sind. Die Menge dieser
Teilmengen schreiben wir als \(\PowNE{A}\).

Für \(X,Y\in\PowNE{A}\) wird die ursprüngliche Operation mengenweise
fortgesetzt. Der Index \(\mathcal P\) unterscheidet diese Operation auf
Teilmengen von der Operation \(\star\) auf einzelnen Elementen.

\FormulaDefDeltaK[Mengenweise Fortsetzung einer Halbgruppenoperation]{%
  \begin{aligned}[t]
    &\SemigroupAlg{A}{\star}\dsep X,Y\in\PowNE{A}\vdash{}\\[-2pt]
    &X\star_{\mathcal P}Y\coloneqq
      \bigl\{z\in A\bigm|
        \exists x\in X\,\exists y\in Y\,(z=x\star y)\bigr\}
  \end{aligned}
}{PowerSemigroupProductDef}{%
  \DeltaRow{Mengen}{A\dsep X\dsep Y\dsep x\dsep y\dsep z}
  \DeltaRow{Funktionen}{\star\dsep\star_{\mathcal P}}
  \DeltaRow{\textbf{Neues Symbol}}{}
  \DeltaPrem{Mengenweise Operation}{\star_{\mathcal P}}
}

\Needspace{12\baselineskip}
\FormulaThmDeltaK[Elementkriterium für das Mengenprodukt]{%
  \begin{aligned}[t]
    &\SemigroupAlg{A}{\star}\dsep X,Y\in\PowNE{A}\vdash{}\\[-2pt]
    &z\in X\star_{\mathcal P}Y
      \leftrightarrow
      \Bigl(z\in A\land
        \exists x\in X\,\exists y\in Y\,(z=x\star y)\Bigr)
  \end{aligned}
}{PowerSemigroupProductMembership}{%
  \DeltaRow{Mengen}{A\dsep X\dsep Y\dsep x\dsep y\dsep z}
  \DeltaRow{Funktionen}{\star\dsep\star_{\mathcal P}}
}
\begin{tabproofwide}
  \proofstepwidestar[1]{\SemigroupAlg{A}{\star}}{\rA}
  \proofstepwidestar[2]{X\in\PowNE{A}}{\rA}
  \proofstepwidestar[3]{Y\in\PowNE{A}}{\rA}
  \proofstepwidestar[1,2,3]{%
    X\star_{\mathcal P}Y=
    \bigl\{z\in A\bigm|
      \exists x\in X\,\exists y\in Y\,(z=x\star y)\bigr\}}{%
    \FormulaRefAuto{PowerSemigroupProductDef}[def]{1,2,3}}
  \proofstepwidestar[1,2,3]{\begin{aligned}[t]
 &z\in X\star_{\mathcal P}Y\leftrightarrow{}\\[-2pt]
 &\quad\Bigl(z\in A\land\exists x\in X\,\exists y\in Y\,(z=x\star y)\Bigr)
 \end{aligned}}{%
    \FormulaRefAuto{x\in\{u\in A\mid P(u)\}\eqvdash
      x\in A\land P(x)}[thm]{4}}
\end{tabproofwide}

Im folgenden Beweis fassen \(\exists I^{*}\) und \(\exists E^{*}\) endlich
viele geschachtelte Anwendungen von \(\exists I\) beziehungsweise
\(\exists E\) zusammen. Bei beschränkten Quantoren schließen diese
Kurzschreibweisen die erforderlichen \(\land\)-Schritte ein.

\FormulaThmDeltaK[Die nichtleeren Teilmengen bilden eine Halbgruppe]{%
  \SemigroupAlg{A}{\star}\vdash
  \SemigroupAlg{\PowNE{A}}{\star_{\mathcal P}}
}{NonemptyPowerSemigroup}{%
  \DeltaRow{Mengen}{%
    A\dsep X\dsep Y\dsep Z\dsep u\dsep v\dsep w\dsep x\dsep y\dsep z}
  \DeltaRow{Funktionen}{\star\dsep\star_{\mathcal P}}
}
\begin{tabproofsplitwide}
  \proofpartwideindR[Trägerelement einer nichtleeren Teilmenge]{%
    X\in\PowNE{A}\dsep x\in X\vdash x\in A
  }{%
    X\in\PowNE{A}\dsep x\in X\vdash x\in A
  }[PowerSemigroupCarrierElement]
    \proofstepwidestar[1]{X\in\PowNE{A}}{\rA}
    \proofstepwidestar[2]{x\in X}{\rA}
    \proofstepwidestar[1]{X\subseteq A\land X\neq\varnothing}{%
      \FormulaRefAuto{NonemptyPowersetMembership}[thm]{1}}
    \proofstepwidestar[1]{X\subseteq A}{%
      \rAEa{3}}
    \proofstepwidestar[1,2]{x\in A}{%
      \FormulaRefAuto{A\subseteq B,\,x\in A\vdash x\in B}[thm]{4,2}}
  \closeproofpartwideindR

  \proofpartwideindR[Elementare Produkte liegen im Mengenprodukt]{%
    \begin{aligned}[t]
      &\SemigroupAlg{A}{\star}\dsep X,Y\in\PowNE{A}\dsep
        x\in X\dsep y\in Y\vdash{}\\[-2pt]
      &x\star y\in X\star_{\mathcal P}Y
    \end{aligned}
  }{%
    \begin{aligned}[t]
      &\SemigroupAlg{A}{\star}\dsep X,Y\in\PowNE{A}\dsep
        x\in X\dsep y\in Y\vdash{}\\[-2pt]
      &x\star y\in X\star_{\mathcal P}Y
    \end{aligned}
  }[PowerSemigroupProductContains]
    \proofstepwidestar[1]{\SemigroupAlg{A}{\star}}{\rA}
    \proofstepwidestar[2]{X\in\PowNE{A}}{\rA}
    \proofstepwidestar[3]{Y\in\PowNE{A}}{\rA}
    \proofstepwidestar[4]{x\in X}{\rA}
    \proofstepwidestar[5]{y\in Y}{\rA}
    \proofstepwidestar[2,4]{x\in A}{%
      \FormulaRefAuto{PowerSemigroupCarrierElement}[thm]{2,4}}
    \proofstepwidestar[3,5]{y\in A}{%
      \FormulaRefAuto{PowerSemigroupCarrierElement}[thm]{3,5}}
    \proofstepwidestar[1,2,3,4,5]{x\star y\in A}{%
      \FormulaRefAuto{SemigroupClosureAxiom}[ax]{1,6,7}}
    \proofstepwidestar[4,5]{%
      \exists u\in X\,\exists v\in Y\,(x\star y=u\star v)}{%
      \rEI{4,\rEI{5,\rII}}}
    \proofstepwidestar[1,2,3,4,5]{%
      x\star y\in A\land
      \exists u\in X\,\exists v\in Y\,(x\star y=u\star v)}{%
      \rAI{8,9}}
    \proofstepwidestar[1,2,3,4,5]{x\star y\in X\star_{\mathcal P}Y}{%
      \FormulaRefAuto{PowerSemigroupProductMembership}[thm]{1,2,3,10}}
  \closeproofpartwideindR

  \proofpartwideindR[Das Mengenprodukt bleibt im ursprünglichen Träger]{%
    \SemigroupAlg{A}{\star}\dsep X,Y\in\PowNE{A}
    \vdash X\star_{\mathcal P}Y\subseteq A
  }{%
    \SemigroupAlg{A}{\star}\dsep X,Y\in\PowNE{A}
    \vdash X\star_{\mathcal P}Y\subseteq A
  }[PowerSemigroupProductSubset]
    \proofstepwidestar[1]{\SemigroupAlg{A}{\star}}{\rA}
    \proofstepwidestar[2]{X\in\PowNE{A}}{\rA}
    \proofstepwidestar[3]{Y\in\PowNE{A}}{\rA}
    \proofstepwidestar[4]{w\in X\star_{\mathcal P}Y}{\rA}
    \proofstepwidestar[1,2,3,4]{%
      w\in A\land
      \exists x\in X\,\exists y\in Y\,(w=x\star y)}{%
      \FormulaRefAuto{PowerSemigroupProductMembership}[thm]{1,2,3,4}}
    \proofstepwidestar[1,2,3,4]{w\in A}{\rAEa{5}}
    \proofstepwidestar[1,2,3]{X\star_{\mathcal P}Y\subseteq A}{%
      \FormulaRefAuto{SubsetIntroductionRule}[thm]{[4]\vdots 6}}
  \closeproofpartwideindR

  \proofpartwideindR[Das Mengenprodukt ist nichtleer]{%
    \SemigroupAlg{A}{\star}\dsep X,Y\in\PowNE{A}
    \vdash X\star_{\mathcal P}Y\neq\varnothing
  }{%
    \SemigroupAlg{A}{\star}\dsep X,Y\in\PowNE{A}
    \vdash X\star_{\mathcal P}Y\neq\varnothing
  }[PowerSemigroupProductNonempty]
    \proofstepwidestar[1]{\SemigroupAlg{A}{\star}}{\rA}
    \proofstepwidestar[2]{X\in\PowNE{A}}{\rA}
    \proofstepwidestar[3]{Y\in\PowNE{A}}{\rA}
    \proofstepwidestar[2]{X\subseteq A\land X\neq\varnothing}{%
      \FormulaRefAuto{NonemptyPowersetMembership}[thm]{2}}
    \proofstepwidestar[2]{X\neq\varnothing}{%
      \rAEb{4}}
    \proofstepwidestar[3]{Y\subseteq A\land Y\neq\varnothing}{%
      \FormulaRefAuto{NonemptyPowersetMembership}[thm]{3}}
    \proofstepwidestar[3]{Y\neq\varnothing}{%
      \rAEb{6}}
    \proofstepwidestar[2]{\exists x\,(x\in X)}{%
      \FormulaRefAuto{A\neq\varnothing\eqvdash
        \exists x\,(x\in A)}[thm]{5}}
    \proofstepwidestar[3]{\exists y\,(y\in Y)}{%
      \FormulaRefAuto{A\neq\varnothing\eqvdash
        \exists x\,(x\in A)}[thm]{7}}
    \proofstepwidestar[8]{x\in X}{\rA}
    \proofstepwidestar[9]{y\in Y}{\rA}
    \proofstepwidestar[1,2,3,10,11]{x\star y\in X\star_{\mathcal P}Y}{%
      \FormulaRefAuto{PowerSemigroupProductContains}[thm]{1,2,3,10,11}}
    \proofstepwidestar[1,2,3,10,11]{X\star_{\mathcal P}Y\neq\varnothing}{%
      \FormulaRefAuto{a\in A\vdash A\neq\varnothing}[thm]{12}}
    \proofstepwidestar[1,2,3,10]{X\star_{\mathcal P}Y\neq\varnothing}{%
      \rEE{9,11,13}}
    \proofstepwidestar[1,2,3]{X\star_{\mathcal P}Y\neq\varnothing}{%
      \rEE{8,10,14}}
  \closeproofpartwideindR

  \proofpartwideindR[Abgeschlossenheit der mengenweisen Operation]{%
    \SemigroupAlg{A}{\star}\dsep X,Y\in\PowNE{A}
    \vdash X\star_{\mathcal P}Y\in\PowNE{A}
  }{%
    \SemigroupAlg{A}{\star}\dsep X,Y\in\PowNE{A}
    \vdash X\star_{\mathcal P}Y\in\PowNE{A}
  }[PowerSemigroupProductClosure]
    \proofstepwidestar[1]{\SemigroupAlg{A}{\star}}{\rA}
    \proofstepwidestar[2]{X\in\PowNE{A}}{\rA}
    \proofstepwidestar[3]{Y\in\PowNE{A}}{\rA}
    \proofstepwidestar[1,2,3]{X\star_{\mathcal P}Y\subseteq A}{%
      \FormulaRefAuto{PowerSemigroupProductSubset}[thm]{1,2,3}}
    \proofstepwidestar[1,2,3]{X\star_{\mathcal P}Y\neq\varnothing}{%
      \FormulaRefAuto{PowerSemigroupProductNonempty}[thm]{1,2,3}}
    \proofstepwidestar[1,2,3]{%
      X\star_{\mathcal P}Y\subseteq A\land
      X\star_{\mathcal P}Y\neq\varnothing}{\rAI{4,5}}
    \proofstepwidestar[1,2,3]{X\star_{\mathcal P}Y\in\PowNE{A}}{%
      \FormulaRefAuto{NonemptyPowersetMembership}[thm]{6}}
  \closeproofpartwideindR

  \proofpartwideindR[Zeugen eines Mengenprodukts]{%
    \begin{aligned}[t]
      &\SemigroupAlg{A}{\star}\dsep X,Y\in\PowNE{A}\dsep
        w\in X\star_{\mathcal P}Y\vdash{}\\[-2pt]
      &\exists x\in X\,\exists y\in Y\,(w=x\star y)
    \end{aligned}
  }{%
    \begin{aligned}[t]
      &\SemigroupAlg{A}{\star}\dsep X,Y\in\PowNE{A}\dsep
        w\in X\star_{\mathcal P}Y\vdash{}\\[-2pt]
      &\exists x\in X\,\exists y\in Y\,(w=x\star y)
    \end{aligned}
  }[PowerSemigroupProductWitnesses]
    \proofstepwidestar[1]{\SemigroupAlg{A}{\star}}{\rA}
    \proofstepwidestar[2]{X\in\PowNE{A}}{\rA}
    \proofstepwidestar[3]{Y\in\PowNE{A}}{\rA}
    \proofstepwidestar[4]{w\in X\star_{\mathcal P}Y}{\rA}
    \proofstepwidestar[1,2,3,4]{%
      w\in A\land
      \exists x\in X\,\exists y\in Y\,(w=x\star y)}{%
      \FormulaRefAuto{PowerSemigroupProductMembership}[thm]{1,2,3,4}}
    \proofstepwidestar[1,2,3,4]{%
      \exists x\in X\,\exists y\in Y\,(w=x\star y)}{\rAEb{5}}
  \closeproofpartwideindR

  \proofpartwideindR[Assoziativität auf Teilmengenelementen]{%
    \begin{aligned}[t]
      &\SemigroupAlg{A}{\star}\dsep X,Y,Z\in\PowNE{A}\dsep
        x\in X\dsep y\in Y\dsep z\in Z\vdash{}\\[-2pt]
      &(x\star y)\star z=x\star(y\star z)
    \end{aligned}
  }{%
    \begin{aligned}[t]
      &\SemigroupAlg{A}{\star}\dsep X,Y,Z\in\PowNE{A}\dsep
        x\in X\dsep y\in Y\dsep z\in Z\vdash{}\\[-2pt]
      &(x\star y)\star z=x\star(y\star z)
    \end{aligned}
  }[PowerSemigroupElementAssociativity]
    \proofstepwidestar[1]{\SemigroupAlg{A}{\star}}{\rA}
    \proofstepwidestar[2]{X\in\PowNE{A}}{\rA}
    \proofstepwidestar[3]{Y\in\PowNE{A}}{\rA}
    \proofstepwidestar[4]{Z\in\PowNE{A}}{\rA}
    \proofstepwidestar[5]{x\in X}{\rA}
    \proofstepwidestar[6]{y\in Y}{\rA}
    \proofstepwidestar[7]{z\in Z}{\rA}
    \proofstepwidestar[2,5]{x\in A}{%
      \FormulaRefAuto{PowerSemigroupCarrierElement}[thm]{2,5}}
    \proofstepwidestar[3,6]{y\in A}{%
      \FormulaRefAuto{PowerSemigroupCarrierElement}[thm]{3,6}}
    \proofstepwidestar[4,7]{z\in A}{%
      \FormulaRefAuto{PowerSemigroupCarrierElement}[thm]{4,7}}
    \proofstepwidestar[1,2,3,4,5,6,7]{%
      (x\star y)\star z=x\star(y\star z)}{%
      \FormulaRefAuto{SemigroupAssociativityAxiom}[ax]{1,8,9,10}}
  \closeproofpartwideindR

  \proofpartwideindR[Dreifachzeugen der links geklammerten Seite]{%
    \begin{aligned}[t]
      &\SemigroupAlg{A}{\star}\dsep X,Y,Z\in\PowNE{A}\dsep
        w\in (X\star_{\mathcal P}Y)\star_{\mathcal P}Z\vdash{}\\[-2pt]
      &\exists x\in X\,\exists y\in Y\,\exists z\in Z\,
        \bigl(w=(x\star y)\star z\bigr)
    \end{aligned}
  }{%
    \begin{aligned}[t]
      &\SemigroupAlg{A}{\star}\dsep X,Y,Z\in\PowNE{A}\dsep
        w\in (X\star_{\mathcal P}Y)\star_{\mathcal P}Z\vdash{}\\[-2pt]
      &\exists x\in X\,\exists y\in Y\,\exists z\in Z\,
        \bigl(w=(x\star y)\star z\bigr)
    \end{aligned}
  }[PowerSemigroupLeftTripleMembershipForward]
    \proofstepwidestar[1]{\SemigroupAlg{A}{\star}}{\rA}
    \proofstepwidestar[2]{X\in\PowNE{A}}{\rA}
    \proofstepwidestar[3]{Y\in\PowNE{A}}{\rA}
    \proofstepwidestar[4]{Z\in\PowNE{A}}{\rA}
    \proofstepwidestar[5]{%
      w\in (X\star_{\mathcal P}Y)\star_{\mathcal P}Z}{\rA}
    \proofstepwidestar[1,2,3]{X\star_{\mathcal P}Y\in\PowNE{A}}{%
      \FormulaRefAuto{PowerSemigroupProductClosure}[thm]{1,2,3}}
    \proofstepwidestar[1,2,3,4,5]{%
      \exists u\in X\star_{\mathcal P}Y\,
      \exists z\in Z\,(w=u\star z)}{%
      \FormulaRefAuto{PowerSemigroupProductWitnesses}[thm]{1,6,4,5}}
    \proofstepwidestar[7]{u\in X\star_{\mathcal P}Y}{\rA}
    \proofstepwidestar[7]{z\in Z}{\rA}
    \proofstepwidestar[7]{w=u\star z}{\rA}
    \proofstepwidestar[1,2,3,8]{%
      \exists x\in X\,\exists y\in Y\,(u=x\star y)}{%
      \FormulaRefAuto{PowerSemigroupProductWitnesses}[thm]{1,2,3,8}}
    \proofstepwidestar[11]{x\in X}{\rA}
    \proofstepwidestar[11]{y\in Y}{\rA}
    \proofstepwidestar[11]{u=x\star y}{\rA}
    \proofstepwidestar[7,11]{w=(x\star y)\star z}{\rIE{14,10}}
    \proofstepwidestar[7,11]{%
      \exists x\in X\,\exists y\in Y\,\exists z\in Z\,
      \bigl(w=(x\star y)\star z\bigr)}{%
      \rEIStar{12,13,9,15}}
    \proofstepwidestar[1,2,3,4,5]{%
      \exists x\in X\,\exists y\in Y\,\exists z\in Z\,
      \bigl(w=(x\star y)\star z\bigr)}{%
      \rEEStar{7\text{--}16}}
  \closeproofpartwideindR

  \proofpartwideindR[Aufbau der links geklammerten Seite]{%
    \begin{aligned}[t]
      &\SemigroupAlg{A}{\star}\dsep X,Y,Z\in\PowNE{A}\dsep{}\\[-2pt]
      &\exists x\in X\,\exists y\in Y\,\exists z\in Z\,
        \bigl(w=(x\star y)\star z\bigr)\vdash{}\\[-2pt]
      &w\in (X\star_{\mathcal P}Y)\star_{\mathcal P}Z
    \end{aligned}
  }{%
    \begin{aligned}[t]
      &\SemigroupAlg{A}{\star}\dsep X,Y,Z\in\PowNE{A}\dsep{}\\[-2pt]
      &\exists x\in X\,\exists y\in Y\,\exists z\in Z\,
        \bigl(w=(x\star y)\star z\bigr)\vdash{}\\[-2pt]
      &w\in (X\star_{\mathcal P}Y)\star_{\mathcal P}Z
    \end{aligned}
  }[PowerSemigroupLeftTripleMembershipBackward]
    \proofstepwidestar[1]{\SemigroupAlg{A}{\star}}{\rA}
    \proofstepwidestar[2]{X\in\PowNE{A}}{\rA}
    \proofstepwidestar[3]{Y\in\PowNE{A}}{\rA}
    \proofstepwidestar[4]{Z\in\PowNE{A}}{\rA}
    \proofstepwidestar[5]{%
      \exists x\in X\,\exists y\in Y\,\exists z\in Z\,
      \bigl(w=(x\star y)\star z\bigr)}{\rA}
    \proofstepwidestar[5]{x\in X}{\rA}
    \proofstepwidestar[5]{y\in Y}{\rA}
    \proofstepwidestar[5]{z\in Z}{\rA}
    \proofstepwidestar[5]{w=(x\star y)\star z}{\rA}
    \proofstepwidestar[1,2,3,5]{%
      x\star y\in X\star_{\mathcal P}Y}{%
      \FormulaRefAuto{PowerSemigroupProductContains}[thm]{1,2,3,6,7}}
    \proofstepwidestar[1,2,3]{X\star_{\mathcal P}Y\in\PowNE{A}}{%
      \FormulaRefAuto{PowerSemigroupProductClosure}[thm]{1,2,3}}
    \proofstepwidestarbreak[1,2,3,4,5]{%
      (x\star y)\star z
      \in (X\star_{\mathcal P}Y)\star_{\mathcal P}Z}{%
      \FormulaRefAuto{PowerSemigroupProductContains}[thm]{1,11,4,10,8}}
    \proofstepwidestar[5]{(x\star y)\star z=w}{%
      \FormulaRefAuto{a=b\vdash b=a}[thm]{9}}
    \proofstepwidestar[1,2,3,4,5]{%
      w\in (X\star_{\mathcal P}Y)\star_{\mathcal P}Z}{\rIE{13,12}}
    \proofstepwidestar[1,2,3,4,5]{%
      w\in (X\star_{\mathcal P}Y)\star_{\mathcal P}Z}{%
      \rEEStar{5\text{--}14}}
  \closeproofpartwideindR

  \proofpartwideindR[Aufbau der rechts geklammerten Seite]{%
    \begin{aligned}[t]
      &\SemigroupAlg{A}{\star}\dsep X,Y,Z\in\PowNE{A}\dsep{}\\[-2pt]
      &\exists x\in X\,\exists y\in Y\,\exists z\in Z\,
        \bigl(w=x\star(y\star z)\bigr)\vdash{}\\[-2pt]
      &w\in X\star_{\mathcal P}(Y\star_{\mathcal P}Z)
    \end{aligned}
  }{%
    \begin{aligned}[t]
      &\SemigroupAlg{A}{\star}\dsep X,Y,Z\in\PowNE{A}\dsep{}\\[-2pt]
      &\exists x\in X\,\exists y\in Y\,\exists z\in Z\,
        \bigl(w=x\star(y\star z)\bigr)\vdash{}\\[-2pt]
      &w\in X\star_{\mathcal P}(Y\star_{\mathcal P}Z)
    \end{aligned}
  }[PowerSemigroupRightTripleMembershipForward]
    \proofstepwidestar[1]{\SemigroupAlg{A}{\star}}{\rA}
    \proofstepwidestar[2]{X\in\PowNE{A}}{\rA}
    \proofstepwidestar[3]{Y\in\PowNE{A}}{\rA}
    \proofstepwidestar[4]{Z\in\PowNE{A}}{\rA}
    \proofstepwidestar[5]{%
      \exists x\in X\,\exists y\in Y\,\exists z\in Z\,
      \bigl(w=x\star(y\star z)\bigr)}{\rA}
    \proofstepwidestar[5]{x\in X}{\rA}
    \proofstepwidestar[5]{y\in Y}{\rA}
    \proofstepwidestar[5]{z\in Z}{\rA}
    \proofstepwidestar[5]{w=x\star(y\star z)}{\rA}
    \proofstepwidestar[1,3,4,5]{%
      y\star z\in Y\star_{\mathcal P}Z}{%
      \FormulaRefAuto{PowerSemigroupProductContains}[thm]{1,3,4,7,8}}
    \proofstepwidestar[1,3,4]{Y\star_{\mathcal P}Z\in\PowNE{A}}{%
      \FormulaRefAuto{PowerSemigroupProductClosure}[thm]{1,3,4}}
    \proofstepwidestarbreak[1,2,3,4,5]{%
      x\star(y\star z)
      \in X\star_{\mathcal P}(Y\star_{\mathcal P}Z)}{%
      \FormulaRefAuto{PowerSemigroupProductContains}[thm]{1,2,11,6,10}}
    \proofstepwidestar[5]{x\star(y\star z)=w}{%
      \FormulaRefAuto{a=b\vdash b=a}[thm]{9}}
    \proofstepwidestar[1,2,3,4,5]{%
      w\in X\star_{\mathcal P}(Y\star_{\mathcal P}Z)}{\rIE{13,12}}
    \proofstepwidestar[1,2,3,4,5]{%
      w\in X\star_{\mathcal P}(Y\star_{\mathcal P}Z)}{%
      \rEEStar{5\text{--}14}}
  \closeproofpartwideindR

  \proofpartwideindR[Dreifachzeugen der rechts geklammerten Seite]{%
    \begin{aligned}[t]
      &\SemigroupAlg{A}{\star}\dsep X,Y,Z\in\PowNE{A}\dsep
        w\in X\star_{\mathcal P}(Y\star_{\mathcal P}Z)\vdash{}\\[-2pt]
      &\exists x\in X\,\exists y\in Y\,\exists z\in Z\,
        \bigl(w=x\star(y\star z)\bigr)
    \end{aligned}
  }{%
    \begin{aligned}[t]
      &\SemigroupAlg{A}{\star}\dsep X,Y,Z\in\PowNE{A}\dsep
        w\in X\star_{\mathcal P}(Y\star_{\mathcal P}Z)\vdash{}\\[-2pt]
      &\exists x\in X\,\exists y\in Y\,\exists z\in Z\,
        \bigl(w=x\star(y\star z)\bigr)
    \end{aligned}
  }[PowerSemigroupRightTripleMembershipBackward]
    \proofstepwidestar[1]{\SemigroupAlg{A}{\star}}{\rA}
    \proofstepwidestar[2]{X\in\PowNE{A}}{\rA}
    \proofstepwidestar[3]{Y\in\PowNE{A}}{\rA}
    \proofstepwidestar[4]{Z\in\PowNE{A}}{\rA}
    \proofstepwidestar[5]{%
      w\in X\star_{\mathcal P}(Y\star_{\mathcal P}Z)}{\rA}
    \proofstepwidestar[1,3,4]{Y\star_{\mathcal P}Z\in\PowNE{A}}{%
      \FormulaRefAuto{PowerSemigroupProductClosure}[thm]{1,3,4}}
    \proofstepwidestar[1,2,3,4,5]{%
      \exists x\in X\,
      \exists v\in Y\star_{\mathcal P}Z\,(w=x\star v)}{%
      \FormulaRefAuto{PowerSemigroupProductWitnesses}[thm]{1,2,6,5}}
    \proofstepwidestar[7]{x\in X}{\rA}
    \proofstepwidestar[7]{v\in Y\star_{\mathcal P}Z}{\rA}
    \proofstepwidestar[7]{w=x\star v}{\rA}
    \proofstepwidestar[1,3,4,9]{%
      \exists y\in Y\,\exists z\in Z\,(v=y\star z)}{%
      \FormulaRefAuto{PowerSemigroupProductWitnesses}[thm]{1,3,4,9}}
    \proofstepwidestar[11]{y\in Y}{\rA}
    \proofstepwidestar[11]{z\in Z}{\rA}
    \proofstepwidestar[11]{v=y\star z}{\rA}
    \proofstepwidestar[7,11]{w=x\star(y\star z)}{\rIE{14,10}}
    \proofstepwidestar[7,11]{%
      \exists x\in X\,\exists y\in Y\,\exists z\in Z\,
      \bigl(w=x\star(y\star z)\bigr)}{%
      \rEIStar{8,12,13,15}}
    \proofstepwidestar[1,2,3,4,5]{%
      \exists x\in X\,\exists y\in Y\,\exists z\in Z\,
      \bigl(w=x\star(y\star z)\bigr)}{%
      \rEEStar{7\text{--}16}}
  \closeproofpartwideindR

  \par\Needspace{12\baselineskip}
  \proofpartwideindR[Umklammern der Dreifachzeugen nach rechts]{%
    \begin{aligned}[t]
      &\SemigroupAlg{A}{\star}\dsep X,Y,Z\in\PowNE{A}\dsep{}\\[-2pt]
      &\exists x\in X\,\exists y\in Y\,\exists z\in Z\,
        \bigl(w=(x\star y)\star z\bigr)\vdash{}\\[-2pt]
      &\exists x\in X\,\exists y\in Y\,\exists z\in Z\,
        \bigl(w=x\star(y\star z)\bigr)
    \end{aligned}
  }{%
    \begin{aligned}[t]
      &\SemigroupAlg{A}{\star}\dsep X,Y,Z\in\PowNE{A}\dsep{}\\[-2pt]
      &\exists x\in X\,\exists y\in Y\,\exists z\in Z\,
        \bigl(w=(x\star y)\star z\bigr)\vdash{}\\[-2pt]
      &\exists x\in X\,\exists y\in Y\,\exists z\in Z\,
        \bigl(w=x\star(y\star z)\bigr)
    \end{aligned}
  }[PowerSemigroupTripleWitnessAssociativityForward]
    \proofstepwidestar[1]{\SemigroupAlg{A}{\star}}{\rA}
    \proofstepwidestar[2]{X\in\PowNE{A}}{\rA}
    \proofstepwidestar[3]{Y\in\PowNE{A}}{\rA}
    \proofstepwidestar[4]{Z\in\PowNE{A}}{\rA}

    \proofstepwidestar[5]{\exists x\in X\,\exists y\in Y\,
      \exists z\in Z\;(w=(x\star y)\star z)}{\rA}

    \proofstepwidestar[6]{x\in X}{\rA}
    \proofstepwidestar[7]{y\in Y}{\rA}
    \proofstepwidestar[8]{z\in Z}{\rA}
    \proofstepwidestar[1,2,3,4,6,7,8]{
      (x\star y)\star z=x\star(y\star z)}{%
      \FormulaRefAuto{PowerSemigroupElementAssociativity}[thm]{
        1,2,3,4,6,7,8}}
    \proofstepwidestar[1,2,3,4]{\begin{aligned}[t]
      &\forall x\in X\,\forall y\in Y\,\forall z\in Z\;\\[-2pt]
      &(x\star y)\star z=x\star(y\star z)
    \end{aligned}}{%
      \rUI{\rRI{6,\rUI{\rRI{7,\rUI{\rRI{8,9}}}}}}}
    \proofstepwidestar[1,2,3,4]{\begin{aligned}[t]
      &(\exists x\in X\,\exists y\in Y\,\exists z\in Z\;
        w=(x\star y)\star z)\\[-2pt]
      &\quad\leftrightarrow
        (\exists x\in X\,\exists y\in Y\,\exists z\in Z\;
        w=x\star(y\star z))
    \end{aligned}}{%
      \FormulaRefAuto{TripleTermWitnessPointwiseEquality}[thm]{10}}

    \proofstepwidestar[1,2,3,4,5]{\exists x\in X\,\exists y\in Y\,
      \exists z\in Z\;(w=x\star(y\star z))}{\rRE{\rLREa{11},5}}
  \closeproofpartwideindR

  \proofpartwideindR[Umklammern der Dreifachzeugen nach links]{%
    \begin{aligned}[t]
      &\SemigroupAlg{A}{\star}\dsep X,Y,Z\in\PowNE{A}\dsep{}\\[-2pt]
      &\exists x\in X\,\exists y\in Y\,\exists z\in Z\,
        \bigl(w=x\star(y\star z)\bigr)\vdash{}\\[-2pt]
      &\exists x\in X\,\exists y\in Y\,\exists z\in Z\,
        \bigl(w=(x\star y)\star z\bigr)
    \end{aligned}
  }{%
    \begin{aligned}[t]
      &\SemigroupAlg{A}{\star}\dsep X,Y,Z\in\PowNE{A}\dsep{}\\[-2pt]
      &\exists x\in X\,\exists y\in Y\,\exists z\in Z\,
        \bigl(w=x\star(y\star z)\bigr)\vdash{}\\[-2pt]
      &\exists x\in X\,\exists y\in Y\,\exists z\in Z\,
        \bigl(w=(x\star y)\star z\bigr)
    \end{aligned}
  }[PowerSemigroupTripleWitnessAssociativityBackward]
    \proofstepwidestar[1]{\SemigroupAlg{A}{\star}}{\rA}
    \proofstepwidestar[2]{X\in\PowNE{A}}{\rA}
    \proofstepwidestar[3]{Y\in\PowNE{A}}{\rA}
    \proofstepwidestar[4]{Z\in\PowNE{A}}{\rA}

    \proofstepwidestar[5]{\exists x\in X\,\exists y\in Y\,
      \exists z\in Z\;(w=x\star(y\star z))}{\rA}

    \proofstepwidestar[6]{x\in X}{\rA}
    \proofstepwidestar[7]{y\in Y}{\rA}
    \proofstepwidestar[8]{z\in Z}{\rA}
    \proofstepwidestar[1,2,3,4,6,7,8]{
      (x\star y)\star z=x\star(y\star z)}{%
      \FormulaRefAuto{PowerSemigroupElementAssociativity}[thm]{
        1,2,3,4,6,7,8}}
    \proofstepwidestar[1,2,3,4]{\begin{aligned}[t]
      &\forall x\in X\,\forall y\in Y\,\forall z\in Z\;\\[-2pt]
      &(x\star y)\star z=x\star(y\star z)
    \end{aligned}}{%
      \rUI{\rRI{6,\rUI{\rRI{7,\rUI{\rRI{8,9}}}}}}}
    \proofstepwidestar[1,2,3,4]{\begin{aligned}[t]
      &(\exists x\in X\,\exists y\in Y\,\exists z\in Z\;
        w=(x\star y)\star z)\\[-2pt]
      &\quad\leftrightarrow
        (\exists x\in X\,\exists y\in Y\,\exists z\in Z\;
        w=x\star(y\star z))
    \end{aligned}}{%
      \FormulaRefAuto{TripleTermWitnessPointwiseEquality}[thm]{10}}

    \proofstepwidestar[1,2,3,4,5]{\exists x\in X\,\exists y\in Y\,
      \exists z\in Z\;(w=(x\star y)\star z)}{\rRE{\rLREb{11},5}}
  \closeproofpartwideindR

  \proofpartwideindR[Zeugenform der links geklammerten Seite]{%
    \begin{aligned}[t]
      &\SemigroupAlg{A}{\star}\dsep X,Y,Z\in\PowNE{A}\vdash{}\\[-2pt]
      &w\in (X\star_{\mathcal P}Y)\star_{\mathcal P}Z
        \leftrightarrow
        \exists x\in X\,\exists y\in Y\,\exists z\in Z\,
        \bigl(w=(x\star y)\star z\bigr)
    \end{aligned}
  }{%
    \begin{aligned}[t]
      &\SemigroupAlg{A}{\star}\dsep X,Y,Z\in\PowNE{A}\vdash{}\\[-2pt]
      &w\in (X\star_{\mathcal P}Y)\star_{\mathcal P}Z
        \leftrightarrow
        \exists x\in X\,\exists y\in Y\,\exists z\in Z\,
        \bigl(w=(x\star y)\star z\bigr)
    \end{aligned}
  }[PowerSemigroupLeftTripleMembershipEquivalence]
  \proofstepwidestar[1]{\SemigroupAlg{A}{\star}}{\rA}
  \proofstepwidestar[2]{X\in\PowNE{A}}{\rA}
  \proofstepwidestar[3]{Y\in\PowNE{A}}{\rA}
  \proofstepwidestar[4]{Z\in\PowNE{A}}{\rA}
  \proofstepwidestar[5]{w\in(X\star_{\mathcal P}Y)\star_{\mathcal P}Z}{\rA}
  \proofstepwidestar[1,2,3,4,5]{\exists x\in X\,\exists y\in Y\,\exists z\in Z\;(w=(x\star y)\star z)}{\FormulaRefAuto{PowerSemigroupLeftTripleMembershipForward}[thm]{1,2,3,4,5}}
  \proofstepwidestar[1,2,3,4]{\begin{aligned}[t]&w\in(X\star_{\mathcal P}Y)\star_{\mathcal P}Z\rightarrow\\[-2pt]&\exists x\in X\,\exists y\in Y\,\exists z\in Z\;(w=(x\star y)\star z)\end{aligned}}{\rRI{5,6}}
  \proofstepwidestar[8]{\exists x\in X\,\exists y\in Y\,\exists z\in Z\;(w=(x\star y)\star z)}{\rA}
  \proofstepwidestar[1,2,3,4,8]{w\in(X\star_{\mathcal P}Y)\star_{\mathcal P}Z}{\FormulaRefAuto{PowerSemigroupLeftTripleMembershipBackward}[thm]{1,2,3,4,8}}
  \proofstepwidestar[1,2,3,4]{\begin{aligned}[t]&\exists x\in X\,\exists y\in Y\,\exists z\in Z\;(w=(x\star y)\star z)\rightarrow\\[-2pt]&w\in(X\star_{\mathcal P}Y)\star_{\mathcal P}Z\end{aligned}}{\rRI{8,9}}
  \proofstepwidestar[1,2,3,4]{\begin{aligned}[t]&w\in(X\star_{\mathcal P}Y)\star_{\mathcal P}Z\leftrightarrow\\[-2pt]&\exists x\in X\,\exists y\in Y\,\exists z\in Z\;(w=(x\star y)\star z)\end{aligned}}{\rLRI{7,10}}
\closeproofpartwideindR

  \proofpartwideindR[Umklammern der Zeugen]{%
    \begin{aligned}[t]
      &\SemigroupAlg{A}{\star}\dsep X,Y,Z\in\PowNE{A}\vdash{}\\[-2pt]
      &\exists x\in X\,\exists y\in Y\,\exists z\in Z\,
        \bigl(w=(x\star y)\star z\bigr)
        \leftrightarrow
        \exists x\in X\,\exists y\in Y\,\exists z\in Z\,
        \bigl(w=x\star(y\star z)\bigr)
    \end{aligned}
  }{%
    \begin{aligned}[t]
      &\SemigroupAlg{A}{\star}\dsep X,Y,Z\in\PowNE{A}\vdash{}\\[-2pt]
      &\exists x\in X\,\exists y\in Y\,\exists z\in Z\,
        \bigl(w=(x\star y)\star z\bigr)
        \leftrightarrow
        \exists x\in X\,\exists y\in Y\,\exists z\in Z\,
        \bigl(w=x\star(y\star z)\bigr)
    \end{aligned}
  }[PowerSemigroupTripleWitnessAssociativityEquivalence]
    \proofstepwidestar[1]{\SemigroupAlg{A}{\star}}{\rA}
    \proofstepwidestar[2]{X\in\PowNE{A}}{\rA}
    \proofstepwidestar[3]{Y\in\PowNE{A}}{\rA}
    \proofstepwidestar[4]{Z\in\PowNE{A}}{\rA}

    \proofstepwidestar[5]{x\in X}{\rA}
    \proofstepwidestar[6]{y\in Y}{\rA}
    \proofstepwidestar[7]{z\in Z}{\rA}
    \proofstepwidestar[1,2,3,4,5,6,7]{
      (x\star y)\star z=x\star(y\star z)}{%
      \FormulaRefAuto{PowerSemigroupElementAssociativity}[thm]{
        1,2,3,4,5,6,7}}
    \proofstepwidestar[1,2,3,4]{\begin{aligned}[t]
      &\forall x\in X\,\forall y\in Y\,\forall z\in Z\;\\[-2pt]
      &(x\star y)\star z=x\star(y\star z)
    \end{aligned}}{%
      \rUI{\rRI{5,\rUI{\rRI{6,\rUI{\rRI{7,8}}}}}}}
    \proofstepwidestar[1,2,3,4]{\begin{aligned}[t]
      &(\exists x\in X\,\exists y\in Y\,\exists z\in Z\;
        w=(x\star y)\star z)\\[-2pt]
      &\quad\leftrightarrow
        (\exists x\in X\,\exists y\in Y\,\exists z\in Z\;
        w=x\star(y\star z))
    \end{aligned}}{%
      \FormulaRefAuto{TripleTermWitnessPointwiseEquality}[thm]{9}}
  \closeproofpartwideindR

  \par\Needspace{12\baselineskip}
  \proofpartwideindR[Zeugenform der rechts geklammerten Seite]{%
    \begin{aligned}[t]
      &\SemigroupAlg{A}{\star}\dsep X,Y,Z\in\PowNE{A}\vdash{}\\[-2pt]
      &\exists x\in X\,\exists y\in Y\,\exists z\in Z\,
        \bigl(w=x\star(y\star z)\bigr)
        \leftrightarrow
        w\in X\star_{\mathcal P}(Y\star_{\mathcal P}Z)
    \end{aligned}
  }{%
    \begin{aligned}[t]
      &\SemigroupAlg{A}{\star}\dsep X,Y,Z\in\PowNE{A}\vdash{}\\[-2pt]
      &\exists x\in X\,\exists y\in Y\,\exists z\in Z\,
        \bigl(w=x\star(y\star z)\bigr)
        \leftrightarrow
        w\in X\star_{\mathcal P}(Y\star_{\mathcal P}Z)
    \end{aligned}
  }[PowerSemigroupRightTripleMembershipEquivalence]
  \proofstepwidestar[1]{\SemigroupAlg{A}{\star}}{\rA}
  \proofstepwidestar[2]{X\in\PowNE{A}}{\rA}
  \proofstepwidestar[3]{Y\in\PowNE{A}}{\rA}
  \proofstepwidestar[4]{Z\in\PowNE{A}}{\rA}
  \proofstepwidestar[5]{\exists x\in X\,\exists y\in Y\,\exists z\in Z\;(w=x\star(y\star z))}{\rA}
  \proofstepwidestar[1,2,3,4,5]{w\in X\star_{\mathcal P}(Y\star_{\mathcal P}Z)}{\FormulaRefAuto{PowerSemigroupRightTripleMembershipForward}[thm]{1,2,3,4,5}}
  \proofstepwidestar[1,2,3,4]{\begin{aligned}[t]&\exists x\in X\,\exists y\in Y\,\exists z\in Z\;(w=x\star(y\star z))\rightarrow\\[-2pt]&w\in X\star_{\mathcal P}(Y\star_{\mathcal P}Z)\end{aligned}}{\rRI{5,6}}
  \proofstepwidestar[8]{w\in X\star_{\mathcal P}(Y\star_{\mathcal P}Z)}{\rA}
  \proofstepwidestar[1,2,3,4,8]{\exists x\in X\,\exists y\in Y\,\exists z\in Z\;(w=x\star(y\star z))}{\FormulaRefAuto{PowerSemigroupRightTripleMembershipBackward}[thm]{1,2,3,4,8}}
  \proofstepwidestar[1,2,3,4]{\begin{aligned}[t]&w\in X\star_{\mathcal P}(Y\star_{\mathcal P}Z)\rightarrow\\[-2pt]&\exists x\in X\,\exists y\in Y\,\exists z\in Z\;(w=x\star(y\star z))\end{aligned}}{\rRI{8,9}}
  \proofstepwidestar[1,2,3,4]{\begin{aligned}[t]&\exists x\in X\,\exists y\in Y\,\exists z\in Z\;(w=x\star(y\star z))\leftrightarrow\\[-2pt]&w\in X\star_{\mathcal P}(Y\star_{\mathcal P}Z)\end{aligned}}{\rLRI{7,10}}
\closeproofpartwideindR

  \proofpartwideindR[Assoziativität der mengenweisen Operation]{%
    \begin{aligned}[t]
      &\SemigroupAlg{A}{\star}\dsep X,Y,Z\in\PowNE{A}\vdash{}\\[-2pt]
      &(X\star_{\mathcal P}Y)\star_{\mathcal P}Z
        =X\star_{\mathcal P}(Y\star_{\mathcal P}Z)
    \end{aligned}
  }{%
    \begin{aligned}[t]
      &\SemigroupAlg{A}{\star}\dsep X,Y,Z\in\PowNE{A}\vdash{}\\[-2pt]
      &(X\star_{\mathcal P}Y)\star_{\mathcal P}Z
        =X\star_{\mathcal P}(Y\star_{\mathcal P}Z)
    \end{aligned}
  }[PowerSemigroupAssociativity]
    \proofstepwidestar[1]{\SemigroupAlg{A}{\star}}{\rA}
    \proofstepwidestar[2]{X\in\PowNE{A}}{\rA}
    \proofstepwidestar[3]{Y\in\PowNE{A}}{\rA}
    \proofstepwidestar[4]{Z\in\PowNE{A}}{\rA}
    \proofstepwide[1,2,3,4]{%
      w\in (X\star_{\mathcal P}Y)\star_{\mathcal P}Z}{\leftrightarrow}{%
      \begin{aligned}[t]
        &\exists x\in X\,\exists y\in Y\\[-2pt]
        &\quad\exists z\in Z\,\bigl(w=(x\star y)\star z\bigr)
      \end{aligned}}{%
      \FormulaRefAuto{PowerSemigroupLeftTripleMembershipEquivalence}[thm]}
    \proofstepwide[1,2,3,4]{}{\leftrightarrow}{%
      \begin{aligned}[t]
        &\exists x\in X\,\exists y\in Y\\[-2pt]
        &\quad\exists z\in Z\,\bigl(w=x\star(y\star z)\bigr)
      \end{aligned}}{%
      \FormulaRefAuto{PowerSemigroupTripleWitnessAssociativityEquivalence}[thm]}
    \proofstepwide[1,2,3,4]{}{\leftrightarrow}{%
      w\in X\star_{\mathcal P}(Y\star_{\mathcal P}Z)}{%
      \FormulaRefAuto{PowerSemigroupRightTripleMembershipEquivalence}[thm]}
    \proofstepwide[1,2,3,4]{%
      w\in (X\star_{\mathcal P}Y)\star_{\mathcal P}Z}{\leftrightarrow}{%
      w\in X\star_{\mathcal P}(Y\star_{\mathcal P}Z)}{\rChain{5,7}}
    \proofstepwidestar[1,2,3,4]{%
      \forall w\,\Bigl(
        w\in (X\star_{\mathcal P}Y)\star_{\mathcal P}Z
        \leftrightarrow
        w\in X\star_{\mathcal P}(Y\star_{\mathcal P}Z)
      \Bigr)}{\rUI{8}}
    \proofstepwidestar[1,2,3,4]{%
      (X\star_{\mathcal P}Y)\star_{\mathcal P}Z
      =X\star_{\mathcal P}(Y\star_{\mathcal P}Z)}{%
      \FormulaRefAuto{\forall x\,(x\in A\leftrightarrow x\in B)
        \vdash A=B}[ax]{9}}
  \closeproofpartwideindR

  \par\Needspace{12\baselineskip}
  \proofpartwide{Schluss}
    \proofstepwidestar[1]{\SemigroupAlg{A}{\star}}{\rA}
    \proofstepwidestar[2]{X\in\PowNE{A}}{\rA}
    \proofstepwidestar[3]{Y\in\PowNE{A}}{\rA}
    \proofstepwidestar[1,2,3]{X\star_{\mathcal P}Y\in\PowNE{A}}{%
      \FormulaRefAuto{PowerSemigroupProductClosure}[thm]{1,2,3}}
    \proofstepwidestar[5]{Z\in\PowNE{A}}{\rA}
    \proofstepwidestar[1,2,3,5]{%
      (X\star_{\mathcal P}Y)\star_{\mathcal P}Z
      =X\star_{\mathcal P}(Y\star_{\mathcal P}Z)}{%
      \FormulaRefAuto{PowerSemigroupAssociativity}[thm]{1,2,3,5}}
    \proofstepwidestar[1]{\SemigroupAlg{\PowNE{A}}{\star_{\mathcal P}}}{%
      \FormulaRefAuto{SemigroupDef}[def]{4,6}}
  \closeproofpartwide
\end{tabproofsplitwide}

Die Halbgruppe \((\PowNE{A},\star_{\mathcal P})\) heißt die
\emph{Potenzhalbgruppe} von \((A,\star)\). Ihre Konstruktion setzt keine
Endlichkeit von \(A\) voraus.

\FormulaDefDeltaK[Einermengenkopie einer Halbgruppe]{%
  \begin{aligned}[t]
    &\SemigroupAlg{A}{\star}\vdash{}\\[-2pt]
    &\operatorname{Sing}(A)\coloneqq
      \bigl\{\{a\}\bigm|a\in A\bigr\},\\[-2pt]
    &\eta_A\colon A\to\operatorname{Sing}(A),
      \qquad \eta_A(a)\coloneqq\{a\}
  \end{aligned}%
}{SemigroupSingletonCopyDef}{%
  \DeltaRow{Mengen}{A\dsep a}
  \DeltaRow{Funktionen}{\star\dsep\eta_A}
  \DeltaRow{\textbf{Neue Symbole}}{}
  \DeltaPrem{Einermengenschicht}{\operatorname{Sing}(A)}
  \DeltaPrem{Einermengenabbildung}{\eta_A}
}

\section{Weitere Beispiele}

Auf jeder total geordneten Menge liefert Band~15 die binären Operationen
\(\min_{A,\leq}\) und \(\max_{A,\leq}\); siehe
\FormulaRefAuto{TotalOrderBinaryMinimum}[def] und
\FormulaRefAuto{TotalOrderBinaryMaximum}[def]. Ihre dort bewiesene
Funktionstypisierung und Assoziativität geben unmittelbar zwei Halbgruppen.

\FormulaThmDeltaK[Die Minimumsoperation bildet eine Halbgruppe]{%
  \TotOrd{A,\leq}\vdash\SemigroupAlg{A}{\min_{A,\leq}}%
}{TotalOrderMinimumSemigroup}{%
  \DeltaRow{Mengen}{A\dsep x\dsep y\dsep z}
  \DeltaRow{Relationsoperatoren}{\leq}
  \DeltaRow{Funktionen}{\min_{A,\leq}}
}
\begin{tabproofwide}
  \proofstepwidestar[1]{\TotOrd{A,\leq}}{\rA}
  \proofstepwidestar[1]{\min_{A,\leq}\colon A\times A\to A}{%
    \FormulaRefAuto{TotalOrderBinaryMinimumFunctionType}[thm]{1}}
  \proofstepwidestar[3]{x\in A}{\rA}
  \proofstepwidestar[4]{y\in A}{\rA}
  \proofstepwidestar[3,4]{(x,y)\in A\times A}{%
    \FormulaRefAuto{a\in A\dsep b\in B
      \vdash(a,b)\in A\times B}[thm]{3,4}}
  \proofstepwidestar[1,3,4]{\min_{A,\leq}(x,y)\in A}{%
    \FormulaRefAuto{F\colon A\to B\dsep x\in A
      \vdash F(x)\in B}[thm]{2,5}}
  \proofstepwidestar[7]{z\in A}{\rA}
  \proofstepwidestar[1,3,4,7]{%
    \begin{aligned}[t]
      &\min_{A,\leq}(\min_{A,\leq}(x,y),z)\\[-2pt]
      &\qquad=\min_{A,\leq}(x,\min_{A,\leq}(y,z))
    \end{aligned}}{%
    \FormulaRefAuto{TotalOrderBinaryMinimumAssociative}[thm]{1,3,4,7}}
  \proofstepwidestar[1]{\SemigroupAlg{A}{\min_{A,\leq}}}{%
    \FormulaRefAuto{SemigroupDef}[def]{6,8}}
\end{tabproofwide}

\newpage

\FormulaThmDeltaK[Die Maximumsoperation bildet eine Halbgruppe]{%
  \TotOrd{A,\leq}\vdash\SemigroupAlg{A}{\max_{A,\leq}}%
}{TotalOrderMaximumSemigroup}{%
  \DeltaRow{Mengen}{A\dsep x\dsep y\dsep z}
  \DeltaRow{Relationsoperatoren}{\leq}
  \DeltaRow{Funktionen}{\max_{A,\leq}}
}
\begin{tabproofwide}
  \proofstepwidestar[1]{\TotOrd{A,\leq}}{\rA}
  \proofstepwidestar[1]{\max_{A,\leq}\colon A\times A\to A}{%
    \FormulaRefAuto{TotalOrderBinaryMaximumFunctionType}[thm]{1}}
  \proofstepwidestar[3]{x\in A}{\rA}
  \proofstepwidestar[4]{y\in A}{\rA}
  \proofstepwidestar[3,4]{(x,y)\in A\times A}{%
    \FormulaRefAuto{a\in A\dsep b\in B
      \vdash(a,b)\in A\times B}[thm]{3,4}}
  \proofstepwidestar[1,3,4]{\max_{A,\leq}(x,y)\in A}{%
    \FormulaRefAuto{F\colon A\to B\dsep x\in A
      \vdash F(x)\in B}[thm]{2,5}}
  \proofstepwidestar[7]{z\in A}{\rA}
  \proofstepwidestar[1,3,4,7]{%
    \begin{aligned}[t]
      &\max_{A,\leq}(\max_{A,\leq}(x,y),z)\\[-2pt]
      &\qquad=\max_{A,\leq}(x,\max_{A,\leq}(y,z))
    \end{aligned}}{%
    \FormulaRefAuto{TotalOrderBinaryMaximumAssociative}[thm]{1,3,4,7}}
  \proofstepwidestar[1]{\SemigroupAlg{A}{\max_{A,\leq}}}{%
    \FormulaRefAuto{SemigroupDef}[def]{6,8}}
\end{tabproofwide}

\FormulaThmDeltaK[Die Addition bildet eine Halbgruppe]{%
  \SemigroupAlg{\mathbb{N}}{+}%
}{NaturalAdditionSemigroup}{%
  \DeltaRow{Mengen}{\mathbb{N}\dsep a\dsep b\dsep c}
}
\begin{tabproof}
  \proofstep{1}{a\in\mathbb{N}}{\rA}
  \proofstep{2}{b\in\mathbb{N}}{\rA}
  \proofstep{1,2}{a+b\in\mathbb{N}}{%
    \FormulaRefAuto{PeanoAddClosure}[thm]{1,2}}
  \proofstep{4}{c\in\mathbb{N}}{\rA}
  \proofstep{1,2,4}{(a+b)+c=a+(b+c)}{%
    \FormulaRefAuto{PeanoAddAssociative}[thm]{1,2,4}}
  \proofstep{}{\SemigroupAlg{\mathbb{N}}{+}}{%
    \FormulaRefAuto{SemigroupDef}[def]{3,5}}
\end{tabproof}

\newpage
\FormulaThmDeltaK[Die Multiplikation bildet eine Halbgruppe]{%
  \SemigroupAlg{\mathbb{N}}{\cdot}%
}{NaturalMultiplicationSemigroup}{%
  \DeltaRow{Mengen}{\mathbb{N}\dsep a\dsep b\dsep c}
}
\begin{tabproof}
  \proofstep{1}{a\in\mathbb{N}}{\rA}
  \proofstep{2}{b\in\mathbb{N}}{\rA}
  \proofstep{1,2}{a\cdot b\in\mathbb{N}}{%
    \FormulaRefAuto{PeanoMulClosure}[thm]{1,2}}
  \proofstep{4}{c\in\mathbb{N}}{\rA}
  \proofstep{1,2,4}{(a\cdot b)\cdot c=a\cdot(b\cdot c)}{%
    \FormulaRefAuto{PeanoMulAssociative}[thm]{1,2,4}}
  \proofstep{}{\SemigroupAlg{\mathbb{N}}{\cdot}}{%
    \FormulaRefAuto{SemigroupDef}[def]{3,5}}
\end{tabproof}

\section{Potenzen in Halbgruppen}

In einer Halbgruppe gibt es im Allgemeinen kein neutrales Element. Daher
werden zun\"achst nur \emph{positive} Potenzen eingef\"uhrt. Die Konstruktion
wird auf die Rekursionsabbildung aus dem Band \emph{Nat\"urliche Zahlen}
zur\"uckgef\"uhrt. Auf diese Weise ist die Potenzfolge eine wirkliche
Abbildung und nicht lediglich eine informelle Rekursionsvorschrift.

\subsection{Rechtstranslation und Potenzfolge}

\FormulaDefDeltaK[Rechtstranslation]{%
  \SemigroupAlg{A}{\star}\dsep a\in A\vdash
  \begin{aligned}[t]
    &\rho_{A,\star,a}\colon A\to A,\\[-2pt]
    &\forall x\in A\,
      \bigl(\rho_{A,\star,a}(x)\coloneqq x\star a\bigr)
  \end{aligned}%
}{SemigroupRightTranslationDef}{%
  \DeltaRow{Mengen}{A\dsep a\dsep x}
  \DeltaRow{Funktionen}{\star\dsep\rho_{A,\star,a}}
  \DeltaRow{\textbf{Neues Symbol}}{}
  \DeltaPrem{Rechtstranslation}{\rho_{A,\star,a}}
}
\begin{tabproof}
  \proofstep{1}{\SemigroupAlg{A}{\star}}{\rA}
  \proofstep{2}{a\in A}{\rA}
  \proofstep{3}{x\in A}{\rA}
  \proofstep{1,2,3}{x\star a\in A}{%
    \FormulaRefAuto{SemigroupClosureAxiom}[ax]{1,3,2}}
  \proofstep{1,2}{\forall x\in A\,x\star a\in A}{%
    \rUI{\rRI{3,4}}}
\end{tabproof}

\FormulaThmDeltaK[Die Rechtstranslation ist eine Abbildung]{%
  \SemigroupAlg{A}{\star}\dsep a\in A\vdash
  \rho_{A,\star,a}\colon A\to A%
}{SemigroupRightTranslationFunction}{%
  \DeltaRow{Mengen}{A\dsep a}
  \DeltaRow{Funktionen}{\star\dsep\rho_{A,\star,a}}
}
\begin{tabproofwide}
  \proofstepwidestar[1]{\SemigroupAlg{A}{\star}}{\rA}
  \proofstepwidestar[2]{a\in A}{\rA}
  \proofstepwidestar[1,2]{\rho_{A,\star,a}\colon A\to A}{%
    \FormulaRefAuto{SemigroupRightTranslationDef}[def]{1,2}}
\end{tabproofwide}

\FormulaThmDeltaK[Auswertung der Rechtstranslation]{%
  \SemigroupAlg{A}{\star}\dsep a,x\in A\vdash
  \rho_{A,\star,a}(x)=x\star a%
}{SemigroupRightTranslationEvaluation}{%
  \DeltaRow{Mengen}{A\dsep a\dsep x}
  \DeltaRow{Funktionen}{\star\dsep\rho_{A,\star,a}}
}
\begin{tabproofwide}
  \proofstepwidestar[1]{\SemigroupAlg{A}{\star}}{\rA}
  \proofstepwidestar[2]{a\in A}{\rA}
  \proofstepwidestar[3]{x\in A}{\rA}
  \proofstepwidestar[1,2]{%
    \forall u\in A\,\rho_{A,\star,a}(u)=u\star a}{%
    \FormulaRefAuto{SemigroupRightTranslationDef}[def]{1,2}}
  \proofstepwidestar[1,2,3]{\rho_{A,\star,a}(x)=x\star a}{%
    \rRE{\rUE{4},3}}
\end{tabproofwide}

\FormulaDefDeltaK[Potenzfolge eines Halbgruppenelements]{%
  \SemigroupAlg{A}{\star}\dsep a\in A\vdash
  \operatorname{pow}_{A,\star}(a)\coloneqq
  \RecFun{a}{\rho_{A,\star,a}}%
}{SemigroupPowerSequenceDef}{%
  \DeltaRow{Mengen}{A\dsep a}
  \DeltaRow{Funktionen}{\star\dsep\rho_{A,\star,a}\dsep
    \operatorname{pow}_{A,\star}}
  \DeltaRow{\textbf{Neues Symbol}}{}
  \DeltaPrem{Potenzfolge}{\operatorname{pow}_{A,\star}(a)}
}

\FormulaThmDeltaK[Die Potenzfolge ist eine Abbildung]{%
  \SemigroupAlg{A}{\star}\dsep a\in A\vdash
  \operatorname{pow}_{A,\star}(a)\colon\mathbb N\to A%
}{SemigroupPowerSequenceFunction}{%
  \DeltaRow{Mengen}{A\dsep a}
  \DeltaRow{Funktionen}{\star\dsep\operatorname{pow}_{A,\star}}
}
\begin{tabproofwide}
  \proofstepwidestar[1]{\SemigroupAlg{A}{\star}}{\rA}
  \proofstepwidestar[2]{a\in A}{\rA}
  \proofstepwidestar[1,2]{\rho_{A,\star,a}\colon A\to A}{%
    \FormulaRefAuto{SemigroupRightTranslationFunction}[thm]{1,2}}
  \proofstepwidestar[1,2]{%
    \RecFun{a}{\rho_{A,\star,a}}\colon\mathbb N\to A}{%
    \FormulaRefAuto{m_0\in M\dsep f\colon M\to M\vdash
      \RecFun{m_0}{f}\colon\mathbb{N}\to M}[thm]{2,3}}
  \proofstepwidestarbreakkeep[1,2]{%
    \operatorname{pow}_{A,\star}(a)=
    \RecFun{a}{\rho_{A,\star,a}}}{%
    \FormulaRefAuto{SemigroupPowerSequenceDef}[def]{1,2}}
  \proofstepwidestar[1,2]{%
    \operatorname{pow}_{A,\star}(a)\colon\mathbb N\to A}{\rIE{5,4}}
\end{tabproofwide}

\subsection{Positive Potenzen}

Ist \(n\in\mathbb N\) und \(n\neq0\), so ist
\((n-1)+1=n\). Der Index \(n-1\) stimmt daher die bei \(0\)
beginnende Potenzfolge auf die bei \(1\) beginnenden Exponenten ab.

\FormulaDefDeltaK[Positive Potenz eines Halbgruppenelements]{%
  \begin{aligned}[t]
    &\SemigroupAlg{A}{\star}\dsep a\in A\dsep
      n\in\mathbb N\dsep n\neq0\vdash{}\\[-2pt]
    &a^n\coloneqq
      \operatorname{pow}_{A,\star}(a)(n-1)
  \end{aligned}
}{SemigroupPositivePowerDef}{%
  \DeltaRow{Mengen}{A\dsep a\dsep n}
  \DeltaRow{Funktionen}{\star\dsep\operatorname{pow}_{A,\star}}
  \DeltaRow{\textbf{Neues Symbol}}{}
  \DeltaPrem{Positive Potenz}{a^n}
}

\FormulaThmDeltaK[Nachfolgerindex der Potenzfolge]{%
  \SemigroupAlg{A}{\star}\dsep a\in A\dsep n\in\mathbb N\vdash
  a^{n+1}=\operatorname{pow}_{A,\star}(a)(n)%
}{SemigroupPowerSuccessorIndex}{%
  \DeltaRow{Mengen}{A\dsep a\dsep n}
  \DeltaRow{Funktionen}{\star\dsep\operatorname{pow}_{A,\star}}
}
\begin{tabproofwide}
  \proofstepwidestar[1]{\SemigroupAlg{A}{\star}}{\rA}
  \proofstepwidestar[2]{a\in A}{\rA}
  \proofstepwidestar[3]{n\in\mathbb N}{\rA}
  \proofstepwidestar[3]{n+1\in\mathbb N}{%
    \FormulaRefAuto{PeanoAddOneClosure}[thm]{3}}
  \proofstepwidestar[3]{n+1\neq0}{%
    \FormulaRefAuto{PeanoAddOneNonzero}[thm]{3}}
  \proofstepwidestar[1,2,3]{%
    a^{n+1}=\operatorname{pow}_{A,\star}(a)((n+1)-1)}{%
    \FormulaRefAuto{SemigroupPositivePowerDef}[def]{1,2,4,5}}
  \proofstepwidestar[3]{(n+1)-1=n}{%
    \FormulaRefAuto{PeanoAddOneMinusOne}[thm]{3}}
  \proofstepwidestar[1,2,3]{%
    a^{n+1}=\operatorname{pow}_{A,\star}(a)(n)}{\rIE{7,6}}
\end{tabproofwide}

\FormulaThmDeltaK[Erste positive Potenz]{%
  \SemigroupAlg{A}{\star}\dsep a\in A\vdash a^1=a%
}{SemigroupPowerOne}{%
  \DeltaRow{Mengen}{A\dsep a}
  \DeltaRow{Funktionen}{\star\dsep\operatorname{pow}_{A,\star}}
}
\begin{tabproofwide}
  \proofstepwidestar[1]{\SemigroupAlg{A}{\star}}{\rA}
  \proofstepwidestar[2]{a\in A}{\rA}
  \proofstepwidestar[]{0\in\mathbb N}{\FormulaRefAuto{PeanoZero}[ax]}
  \proofstepwidestar[1,2]{%
    a^{0+1}=\operatorname{pow}_{A,\star}(a)(0)}{%
    \FormulaRefAuto{SemigroupPowerSuccessorIndex}[thm]{1,2,3}}
  \proofstepwidestar[1,2]{%
    \operatorname{pow}_{A,\star}(a)(0)=
    \RecFun{a}{\rho_{A,\star,a}}(0)}{%
    \FormulaRefAuto{SemigroupPowerSequenceDef}[def]{1,2}}
  \proofstepwidestar[1,2]{\rho_{A,\star,a}\colon A\to A}{%
    \FormulaRefAuto{SemigroupRightTranslationFunction}[thm]{1,2}}
  \proofstepwidestar[1,2]{\RecFun{a}{\rho_{A,\star,a}}(0)=a}{%
    \FormulaRefAuto{m_0\in M\dsep f\colon M\to M\vdash
      \RecFun{m_0}{f}(0)=m_0}[thm]{2,6}}
  \proofstepwidestar[1,2]{\operatorname{pow}_{A,\star}(a)(0)=a}{%
      \FormulaRefAuto{a = b,\, b = c \vdash a = c}[thm]{5,7}}
  \proofstepwidestar[1,2]{a^{0+1}=a}{%
    \FormulaRefAuto{a = b,\, b = c \vdash a = c}[thm]{%
      4,8}}
  \proofstepwidestar[]{1\in\mathbb N}{%
      \FormulaRefAuto{PeanoOneInN}[thm]}
  \proofstepwidestar[]{0+1=1}{%
    \FormulaRefAuto{PeanoAddLeftZero}[thm]{%
      10}}
  \proofstepwidestar[1,2]{a^1=a}{\rIE{11,9}}
\end{tabproofwide}

\FormulaThmDeltaK[Nachfolgerregel positiver Potenzen]{%
  \begin{aligned}[t]
    &\SemigroupAlg{A}{\star}\dsep a\in A\dsep
      n\in\mathbb N\dsep n\neq0\vdash{}\\[-2pt]
    &a^{n+1}=a^n\star a
  \end{aligned}
}{SemigroupPowerSuccessor}{%
  \DeltaRow{Mengen}{A\dsep a\dsep n}
  \DeltaRow{Funktionen}{\star\dsep\operatorname{pow}_{A,\star}}
}
\begin{tabproofwide}
  \proofstepwidestar[1]{\SemigroupAlg{A}{\star}}{\rA}
  \proofstepwidestar[2]{a\in A}{\rA}
  \proofstepwidestar[3]{n\in\mathbb N}{\rA}
  \proofstepwidestar[4]{n\neq0}{\rA}
  \proofstepwidestar[1,2,3]{%
    a^{n+1}=\operatorname{pow}_{A,\star}(a)(n)}{%
    \FormulaRefAuto{SemigroupPowerSuccessorIndex}[thm]{1,2,3}}
  \proofstepwidestar[1,2]{\rho_{A,\star,a}\colon A\to A}{%
    \FormulaRefAuto{SemigroupRightTranslationFunction}[thm]{1,2}}
  \proofstepwidestarbreak[1,2,3,4]{%
    \operatorname{pow}_{A,\star}(a)(n)=
    \RecFun{a}{\rho_{A,\star,a}}(n)}{%
    \FormulaRefAuto{SemigroupPowerSequenceDef}[def]{1,2}}
  \proofstepwidestarbreak[1,2,3,4]{%
    \RecFun{a}{\rho_{A,\star,a}}(n)=
    \rho_{A,\star,a}
      (\RecFun{a}{\rho_{A,\star,a}}(n-1))}{%
    \FormulaRefAuto{PeanoRecursionMinusOne}[thm]{2,6,3,4}}
  \proofstepwidestar[1,2,3,4]{n-1\in\mathbb N}{%
      \FormulaRefAuto{PeanoMinusOneNatural}[thm]{3,4}}
  \proofstepwidestarbreak[1,2,3,4]{%
    \RecFun{a}{\rho_{A,\star,a}}(n-1)\in A}{%
    \FormulaRefAuto{m_0\in M\dsep f\colon M\to M\dsep k\in\mathbb{N}
      \vdash\RecFun{m_0}{f}(k)\in M}[thm]{%
      2,6,9}}
  \proofstepwidestarbreak[1,2,3,4]{%
    \rho_{A,\star,a}
      (\RecFun{a}{\rho_{A,\star,a}}(n-1))
    =\RecFun{a}{\rho_{A,\star,a}}(n-1)\star a}{%
    \FormulaRefAuto{SemigroupRightTranslationEvaluation}[thm]{1,2,10}}
  \proofstepwidestarbreak[1,2,3,4]{%
    a^n=\operatorname{pow}_{A,\star}(a)(n-1)}{%
    \FormulaRefAuto{SemigroupPositivePowerDef}[def]{1,2,3,4}}
  \proofstepwidestar[1,2,3,4]{\operatorname{pow}_{A,\star}(a)=\RecFun{a}{\rho_{A,\star,a}}}{%
      \FormulaRefAuto{SemigroupPowerSequenceDef}[def]{1,2}}
  \proofstepwidestarbreak[1,2,3,4]{%
    a^n=\RecFun{a}{\rho_{A,\star,a}}(n-1)}{%
    \rIE{13,12}}
  \proofstepwidestar[1,2,3,4]{\RecFun{a}{\rho_{A,\star,a}}(n-1)=a^n}{%
      \FormulaRefAuto{a=b\vdash b=a}[thm]{14}}
  \proofstepwidestarbreak[1,2,3,4]{%
    \RecFun{a}{\rho_{A,\star,a}}(n-1)\star a=a^n\star a}{%
    \rIE{15}}
  \proofstepwidestar[1,2,3,4]{\rho_{A,\star,a} (\RecFun{a}{\rho_{A,\star,a}}(n-1))=a^n\star a}{%
      \FormulaRefAuto{a = b,\, b = c \vdash a = c}[thm]{11,16}}
  \proofstepwidestar[1,2,3,4]{\RecFun{a}{\rho_{A,\star,a}}(n)=a^n\star a}{%
      \FormulaRefAuto{a = b,\, b = c \vdash a = c}[thm]{%
          8,17}}
  \proofstepwidestar[1,2,3,4]{\operatorname{pow}_{A,\star}(a)(n)=a^n\star a}{%
      \FormulaRefAuto{a = b,\, b = c \vdash a = c}[thm]{%
        7,18}}
  \proofstepwidestarbreak[1,2,3,4]{a^{n+1}=a^n\star a}{%
    \FormulaRefAuto{a = b,\, b = c \vdash a = c}[thm]{%
      5,19}}
\end{tabproofwide}

\FormulaThmDeltaK[Abgeschlossenheit positiver Potenzen]{%
  \SemigroupAlg{A}{\star}\dsep a\in A\dsep
  n\in\mathbb N\dsep n\neq0\vdash a^n\in A%
}{SemigroupPowerClosure}{%
  \DeltaRow{Mengen}{A\dsep a\dsep n}
  \DeltaRow{Funktionen}{\star\dsep\operatorname{pow}_{A,\star}}
}
\begin{tabproofwide}
  \proofstepwidestar[1]{\SemigroupAlg{A}{\star}}{\rA}
  \proofstepwidestar[2]{a\in A}{\rA}
  \proofstepwidestar[3]{n\in\mathbb N}{\rA}
  \proofstepwidestar[4]{n\neq0}{\rA}
  \proofstepwidestar[1,2]{%
    \operatorname{pow}_{A,\star}(a)\colon\mathbb N\to A}{%
    \FormulaRefAuto{SemigroupPowerSequenceFunction}[thm]{1,2}}
  \proofstepwidestar[3,4]{(n-1)\in\mathbb N}{%
    \FormulaRefAuto{PeanoMinusOneNatural}[thm]{3,4}}
  \proofstepwidestar[1,2,3,4]{%
    \operatorname{pow}_{A,\star}(a)(n-1)\in A}{%
    \FormulaRefAuto{F\colon A\to B\dsep x\in A\vdash F(x)\in B}[thm]{5,6}}
  \proofstepwidestar[1,2,3,4]{%
    a^n=\operatorname{pow}_{A,\star}(a)(n-1)}{%
    \FormulaRefAuto{SemigroupPositivePowerDef}[def]{1,2,3,4}}
  \proofstepwidestar[1,2,3,4]{a^n\in A}{\rIE{8,7}}
\end{tabproofwide}

\subsection{Potenzgesetze}

Die folgenden Gesetze zeigen, dass die Rekursionsklammerung wegen der
Assoziativit\"at keine zus\"atzliche Information tr\"agt.

\FormulaThmDeltaK[Additionsgesetz positiver Potenzen]{%
  \begin{aligned}[t]
    &\SemigroupAlg{A}{\star}\dsep a\in A\dsep
      m,n\in\mathbb N\dsep m\neq0\dsep n\neq0\vdash{}\\[-2pt]
    &a^{m+n}=a^m\star a^n
  \end{aligned}
}{SemigroupPowerAdditionLaw}{%
  \DeltaRow{Mengen}{A\dsep a\dsep m\dsep n\dsep k}
  \DeltaRow{Funktionen}{\star}
}
\begin{tabproofsplitwide}
  \proofpartwideindR[Induktionsanfang bei 1]{%
\SemigroupAlg{A}{\star}\dsep a\in A\dsep m\in\mathbb N\dsep m\neq0\vdash a^{m+1}=a^m\star a^1
}{%
\SemigroupAlg{A}{\star}\dsep a\in A\dsep m\in\mathbb N\dsep m\neq0\vdash a^{m+1}=a^m\star a^1
}[SemigroupPowerAdditionBase]
  \proofstepwidestar[1]{\SemigroupAlg{A}{\star}}{\rA}
  \proofstepwidestar[2]{a\in A}{\rA}
  \proofstepwidestar[3]{m\in\mathbb N}{\rA}
  \proofstepwidestar[4]{m\neq0}{\rA}
  \proofstepwidestar[1,2,3,4]{a^{m+1}=a^m\star a}{\FormulaRefAuto{SemigroupPowerSuccessor}[thm]{1,2,3,4}}
  \proofstepwidestar[1,2]{a^1=a}{\FormulaRefAuto{SemigroupPowerOne}[thm]{1,2}}
  \proofstepwidestar[1,2,3,4]{a^{m+1}=a^m\star a^1}{\rIE{6,5}}
  \closeproofpartwideindR

  \newpage
  \proofpartwideindR[Induktionsschritt von \(k\) auf \(k+1\)]{%
\begin{aligned}[t]&\SemigroupAlg{A}{\star}\dsep a\in A\dsep m,k\in\mathbb N\dsep m\neq0\dsep k\neq0\dsep{}\\[-2pt]&a^{m+k}=a^m\star a^k\vdash a^{m+(k+1)}=a^m\star a^{k+1}\end{aligned}
}{%
\begin{aligned}[t]&\SemigroupAlg{A}{\star}\dsep a\in A\dsep m,k\in\mathbb N\dsep m\neq0\dsep k\neq0\dsep{}\\[-2pt]&a^{m+k}=a^m\star a^k\vdash a^{m+(k+1)}=a^m\star a^{k+1}\end{aligned}
}[SemigroupPowerAdditionStep]
  \proofstepwidestar[1]{\SemigroupAlg{A}{\star}}{\rA}
  \proofstepwidestar[2]{a\in A}{\rA}
  \proofstepwidestar[3]{m\in\mathbb N}{\rA}
  \proofstepwidestar[4]{k\in\mathbb N}{\rA}
  \proofstepwidestar[5]{m\neq0}{\rA}
  \proofstepwidestar[6]{k\neq0}{\rA}
  \proofstepwidestar[7]{a^{m+k}=a^m\star a^k}{\rA}
  \proofstepwidestar[3,4]{m+k\in\mathbb N}{\FormulaRefAuto{PeanoAddClosure}[thm]{3,4}}
  \proofstepwidestar[3,4,5]{m+k\neq0}{\FormulaRefAuto{PeanoAddNonzeroLeft}[thm]{3,4,5}}
  \proofstepwidestar[3,4]{m+(k+1)=(m+k)+1}{\FormulaRefAuto{PeanoAddAddOneRight}[thm]{3,4}}
  \proofstepwidestar[1,2,3,4,5]{a^{(m+k)+1}=a^{m+k}\star a}{\FormulaRefAuto{SemigroupPowerSuccessor}[thm]{1,2,8,9}}
  \proofstepwidestar[1,2,3,5]{a^m\in A}{\FormulaRefAuto{SemigroupPowerClosure}[thm]{1,2,3,5}}
  \proofstepwidestar[1,2,4,6]{a^k\in A}{\FormulaRefAuto{SemigroupPowerClosure}[thm]{1,2,4,6}}
  \proofstepwidestar[1,2,3,4,5,6]{(a^m\star a^k)\star a=a^m\star(a^k\star a)}{\FormulaRefAuto{SemigroupAssociativityAxiom}[ax]{1,12,13,2}}
  \proofstepwidestar[1,2,4,6]{a^{k+1}=a^k\star a}{\FormulaRefAuto{SemigroupPowerSuccessor}[thm]{1,2,4,6}}
  \proofstepwidestar[1,2,3,4,5,6,7]{a^{m+(k+1)}=(a^m\star a^k)\star a}{\rIE{10,\rIE{7,11}}}
  \proofstepwidestar[1,2,3,4,5,6]{(a^m\star a^k)\star a=a^m\star a^{k+1}}{\rIE{15,14}}
  \proofstepwidestar[1,2,3,4,5,6,7]{a^{m+(k+1)}=a^m\star a^{k+1}}{\FormulaRefAuto{a = b,\, b = c \vdash a = c}[thm]{16,17}}
  \closeproofpartwideindR

  \proofpartwide{Induktionsschluss}
    \proofstepwidestarbreak[1]{\SemigroupAlg{A}{\star}}{%
      \rA}
    \proofstepwidestarbreak[2]{a\in A}{%
      \rA}
    \proofstepwidestarbreak[3]{m\in\mathbb N}{%
      \rA}
    \proofstepwidestarbreak[4]{n\in\mathbb N}{%
      \rA}
    \proofstepwidestarbreak[5]{m\neq0}{%
      \rA}
    \proofstepwidestarbreak[6]{n\neq0}{%
      \rA}
    \proofstepwidestarbreak[1,2,3,5]{a^{m+1}=a^m\star a^1}{%
      \FormulaRefAuto{SemigroupPowerAdditionBase}[thm]{1,2,3,5}}
    \proofstepwidestarbreak[8]{k\in\mathbb N}{%
      \rA}
    \proofstepwidestarbreak[9]{k\neq0}{%
      \rA}
    \proofstepwidestarbreak[10]{a^{m+k}=a^m\star a^k}{%
      \rA}
    \proofstepwidestarbreak[1,2,3,5,8,9,10]{a^{m+(k+1)}=a^m\star a^{k+1}}{%
      \FormulaRefAuto{SemigroupPowerAdditionStep}[thm]{1,2,3,8,5,9,10}}
    \proofstepwidestarbreak[1,2,3,5]{\begin{aligned}[t]&\forall k\in\mathbb N\,\bigl(k\neq0\rightarrow{}\\[-2pt]&\quad(a^{m+k}=a^m\star a^k\rightarrow a^{m+(k+1)}=a^m\star a^{k+1})\bigr)\end{aligned}}{%
      \rUI{\rRI{8,\rRI{9,\rRI{10,11}}}}}
    \proofstepwidestarbreak[1,2,3,5]{\forall k\in\mathbb N\,(k\neq0\rightarrow a^{m+k}=a^m\star a^k)}{%
      \FormulaRefAuto{PeanoInductionFromOne}[thm]{7,12}}
    \proofstepwidestarbreak[1,2,3,4,5,6]{a^{m+n}=a^m\star a^n}{%
      \rRE{6,\rRE{4,\rUE{13}}}}
  \closeproofpartwide
\end{tabproofsplitwide}

\FormulaThmDeltaK[Potenzen desselben Elements kommutieren]{%
  \begin{aligned}[t]
    &\SemigroupAlg{A}{\star}\dsep a\in A\dsep
      m,n\in\mathbb N\dsep m\neq0\dsep n\neq0\vdash{}\\[-2pt]
    &a^m\star a^n=a^n\star a^m
  \end{aligned}
}{SemigroupPowersCommute}{%
  \DeltaRow{Mengen}{A\dsep a\dsep m\dsep n}
  \DeltaRow{Funktionen}{\star}
}
\begin{tabproofwide}
  \proofstepwidestar[1]{\SemigroupAlg{A}{\star}}{\rA}
  \proofstepwidestar[2]{a\in A}{\rA}
  \proofstepwidestar[3]{m\in\mathbb N}{\rA}
  \proofstepwidestar[4]{n\in\mathbb N}{\rA}
  \proofstepwidestar[5]{m\neq0}{\rA}
  \proofstepwidestar[6]{n\neq0}{\rA}
  \proofstepwidestar[1,2,3,4,5,6]{a^{m+n}=a^m\star a^n}{%
      \FormulaRefAuto{SemigroupPowerAdditionLaw}[thm]{1,2,3,4,5,6}}
  \proofstepwide[1,2,3,4,5,6]{a^m\star a^n}{=}{a^{m+n}}{%
    \FormulaRefAuto{a=b\vdash b=a}[thm]{%
      7}}
  \proofstepwidestar[3,4]{m+n=n+m}{%
      \FormulaRefAuto{PeanoAddCommutative}[thm]{3,4}}
  \proofstepwide[3,4]{a^{m+n}}{=}{a^{n+m}}{%
    \rIE{9}}
  \proofstepwide[1,2,3,4,5,6]{}{=}{a^n\star a^m}{%
    \FormulaRefAuto{SemigroupPowerAdditionLaw}[thm]{1,2,4,3,6,5}}
  \proofstepwide[1,2,3,4,5,6]{a^m\star a^n}{=}{a^n\star a^m}{\rChain{8,10,11}}
\end{tabproofwide}

\section{Idempotente Elemente}

Ein Element ist idempotent, wenn seine zweite Potenz bereits wieder das
Element selbst ist. F\"ur eine feste Halbgruppe sammeln wir alle solchen
Elemente in einer ausgezeichneten Teilmenge des Tr\"agers.

\FormulaDefDeltaK[Idempotentenmenge einer Halbgruppe]{%
  \SemigroupAlg{A}{\star}\vdash
  E_{A,\star}\coloneqq
  \bigl\{e\in A\bigm|e\star e=e\bigr\}%
}{SemigroupIdempotentSetDef}{%
  \DeltaRow{Mengen}{A\dsep e}
  \DeltaRow{Funktionen}{\star}
  \DeltaRow{\textbf{Neues Symbol}}{}
  \DeltaPrem{Idempotentenmenge}{E_{A,\star}}
}

\FormulaThmDeltaK[Elementkriterium f\"ur Idempotente]{%
  \SemigroupAlg{A}{\star}\vdash
  e\in E_{A,\star}\leftrightarrow
  \bigl(e\in A\land e\star e=e\bigr)%
}{SemigroupIdempotentMembership}{%
  \DeltaRow{Mengen}{A\dsep e}
  \DeltaRow{Funktionen}{\star}
}
\begin{tabproofwide}
  \proofstepwidestar[1]{\SemigroupAlg{A}{\star}}{\rA}
  \proofstepwidestar[1]{%
    E_{A,\star}=\{u\in A\mid u\star u=u\}}{%
    \FormulaRefAuto{SemigroupIdempotentSetDef}[def]{1}}
  \proofstepwidestar[1]{e\in\{u\in A\mid u\star u=u\}\leftrightarrow(e\in A\land e\star e=e)}{%
      \FormulaRefAuto{x\in\{u\in A\mid P(u)\}\eqvdash
      x\in A\land P(x)}[thm]}
  \proofstepwidestar[1]{%
    e\in E_{A,\star}\leftrightarrow
    \bigl(e\in A\land e\star e=e\bigr)}{%
    \rIE{2,3}}
\end{tabproofwide}

\FormulaThmDeltaK[Ein idempotentes Element besitzt nur eine positive Potenz]{%
  \begin{aligned}[t]
    &\SemigroupAlg{A}{\star}\dsep e\in E_{A,\star}\dsep
      n\in\mathbb N\dsep n\neq0\vdash{}\\[-2pt]
    &e^n=e
  \end{aligned}
}{SemigroupIdempotentPower}{%
  \DeltaRow{Mengen}{A\dsep e\dsep n\dsep k}
  \DeltaRow{Funktionen}{\star}
}
\begin{tabproofsplitwide}
  \proofpartwideindR[Induktionsanfang bei 1]{%
\SemigroupAlg{A}{\star}\dsep e\in E_{A,\star}\vdash e^1=e
}{%
\SemigroupAlg{A}{\star}\dsep e\in E_{A,\star}\vdash e^1=e
}[SemigroupIdempotentPowerBase]
  \proofstepwidestar[1]{\SemigroupAlg{A}{\star}}{\rA}
  \proofstepwidestar[2]{e\in E_{A,\star}}{\rA}
  \proofstepwidestar[1,2]{e\in A\land e\star e=e}{%
      \FormulaRefAuto{SemigroupIdempotentMembership}[thm]{1,2}}
  \proofstepwidestar[1,2]{e\in A}{\rAEa{3}}
  \proofstepwidestar[1,2]{e^1=e}{\FormulaRefAuto{SemigroupPowerOne}[thm]{1,4}}
  \closeproofpartwideindR

  \proofpartwideindR[Induktionsschritt von \(k\) auf \(k+1\)]{%
\SemigroupAlg{A}{\star}\dsep e\in E_{A,\star}\dsep k\in\mathbb N\dsep k\neq0\dsep e^k=e\vdash e^{k+1}=e
}{%
\SemigroupAlg{A}{\star}\dsep e\in E_{A,\star}\dsep k\in\mathbb N\dsep k\neq0\dsep e^k=e\vdash e^{k+1}=e
}[SemigroupIdempotentPowerStep]
  \proofstepwidestar[1]{\SemigroupAlg{A}{\star}}{\rA}
  \proofstepwidestar[2]{e\in E_{A,\star}}{\rA}
  \proofstepwidestar[3]{k\in\mathbb N}{\rA}
  \proofstepwidestar[4]{k\neq0}{\rA}
  \proofstepwidestar[5]{e^k=e}{\rA}
  \proofstepwidestar[1,2]{e\in A\land e\star e=e}{\FormulaRefAuto{SemigroupIdempotentMembership}[thm]{1,2}}
  \proofstepwidestar[1,2,3,4]{e^{k+1}=e^k\star e}{\FormulaRefAuto{SemigroupPowerSuccessor}[thm]{1,\rAEa{6},3,4}}
  \proofstepwidestar[1,2,3,4,5]{e^{k+1}=e\star e}{\rIE{5,7}}
  \proofstepwidestar[1,2,3,4,5]{e^{k+1}=e}{\FormulaRefAuto{a = b,\, b = c \vdash a = c}[thm]{8,\rAEb{6}}}
  \closeproofpartwideindR

  \proofpartwide{Induktionsschluss}
    \proofstepwidestarbreak[1]{\SemigroupAlg{A}{\star}}{%
      \rA}
    \proofstepwidestarbreak[2]{e\in E_{A,\star}}{%
      \rA}
    \proofstepwidestarbreak[3]{n\in\mathbb N}{%
      \rA}
    \proofstepwidestarbreak[4]{n\neq0}{%
      \rA}
    \proofstepwidestarbreak[1,2]{e^1=e}{%
      \FormulaRefAuto{SemigroupIdempotentPowerBase}[thm]{1,2}}
    \proofstepwidestarbreak[6]{k\in\mathbb N}{%
      \rA}
    \proofstepwidestarbreak[7]{k\neq0}{%
      \rA}
    \proofstepwidestarbreak[8]{e^k=e}{%
      \rA}
    \proofstepwidestarbreak[1,2,6,7,8]{e^{k+1}=e}{%
      \FormulaRefAuto{SemigroupIdempotentPowerStep}[thm]{1,2,6,7,8}}
    \proofstepwidestarbreak[1,2]{\forall k\in\mathbb N\,(k\neq0\rightarrow(e^k=e\rightarrow e^{k+1}=e))}{%
      \rUI{\rRI{6,\rRI{7,\rRI{8,9}}}}}
    \proofstepwidestarbreak[1,2]{\forall k\in\mathbb N\,(k\neq0\rightarrow e^k=e)}{%
      \FormulaRefAuto{PeanoInductionFromOne}[thm]{5,10}}
    \proofstepwidestarbreak[1,2,3,4]{e^n=e}{%
      \rRE{4,\rRE{3,\rUE{11}}}}
  \closeproofpartwide
\end{tabproofsplitwide}

\section{Unterhalbgruppen}

Eine Unterhalbgruppe erbt die Assoziativit\"at von der umgebenden
Halbgruppe. Deshalb gen\"ugen Nichtleerheit und Abgeschlossenheit unter der
eingeschr\"ankten Operation.

Eine \emph{Unterhalbgruppe} ist dabei noch keine Gruppe: Gemeint ist nur
eine nichtleere, produktabgeschlossene Teilmenge. Echte Untergruppen mit
neutralem Element und Inversen werden erst nach Einführung des
Gruppenbegriffs behandelt.

Für ihren Träger verwenden wir die in der Standardliteratur übliche
Bezeichnung (U); das Prädikat schreiben wir als
\(\mathsf{UHGr}(U;A,\star)\).

\par\noindent\begin{minipage}{\linewidth}
\FormulaDefDeltaK[Unterhalbgruppe]{%
\mathsf{UHGr}(U;A,\star)%
}{SemigroupSubsemigroupDef}{%
  \DeltaRow{Mengen}{A\dsep U\dsep x\dsep y}
  \DeltaRow{Funktionen}{\star}
  \DeltaRow{\textbf{Voraussetzung}}{}
  \DeltaPrem{Umgebende Halbgruppe}{\SemigroupAlg{A}{\star}}
  \DeltaRow{\textbf{Axiome}}{}
  \DeltaRow{Träger}{U\subseteq A}
  \DeltaRow{Nichtleerheit}{U\neq\varnothing}
  \DeltaRow{Produktabschluss}{x\in U\dsep y\in U\vdash x\star y\in U}
  \DeltaRow{\textbf{Neues Symbol}}{}
  \DeltaPrem{Unterhalbgruppe}{\mathsf{UHGr}(U;A,\star)}
}
\end{minipage}\par

\par\noindent\begin{minipage}{\linewidth}
\FormulaAxiomDeltaK[Träger einer Unterhalbgruppe]{%
\SemigroupAlg{A}{\star}\dsep \mathsf{UHGr}(U;A,\star)\vdash U\subseteq A%
}{SemigroupSubsemigroupSubsetAxiom}{%
  \DeltaRow{Mengen}{A\dsep U\dsep x\dsep y}
  \DeltaRow{Funktionen}{\star}
}
\end{minipage}\par

\par\noindent\begin{minipage}{\linewidth}
\FormulaAxiomDeltaK[Nichtleerheit einer Unterhalbgruppe]{%
\SemigroupAlg{A}{\star}\dsep \mathsf{UHGr}(U;A,\star)\vdash U\neq\varnothing%
}{SemigroupSubsemigroupNonemptyAxiom}{%
  \DeltaRow{Mengen}{A\dsep U\dsep x\dsep y}
  \DeltaRow{Funktionen}{\star}
}
\end{minipage}\par

\par\noindent\begin{minipage}{\linewidth}
\FormulaAxiomDeltaK[Produktabschluss einer Unterhalbgruppe]{%
\SemigroupAlg{A}{\star}\dsep \mathsf{UHGr}(U;A,\star)\dsep x\in U\dsep y\in U\vdash x\star y\in U%
}{SemigroupSubsemigroupClosureAxiom}{%
  \DeltaRow{Mengen}{A\dsep U\dsep x\dsep y}
  \DeltaRow{Funktionen}{\star}
}
\end{minipage}\par

\par\noindent\begin{minipage}{\linewidth}
\FormulaThmDeltaK[Kriterium für eine Unterhalbgruppe]{%
\begin{aligned}[t]&\SemigroupAlg{A}{\star}\vdash \mathsf{UHGr}(U;A,\star)\leftrightarrow{}\\[-2pt]&\quad \bigl(U\subseteq A\land U\neq\varnothing\land \forall x\in U\,\forall y\in U\,(x\star y\in U)\bigr)\end{aligned}%
}{SemigroupSubsemigroupCriterion}{%
  \DeltaRow{Mengen}{A\dsep U\dsep x\dsep y}
  \DeltaRow{Funktionen}{\star}
}
\end{minipage}\par
\begin{tabproofwide}
  \proofstepwidestarbreak[1]{\SemigroupAlg{A}{\star}}{\rA}
  \proofstepwidestarbreak[2]{\mathsf{UHGr}(U;A,\star)}{\rA}
  \proofstepwidestarbreak[1,2]{U\subseteq A}{\FormulaRefAuto{SemigroupSubsemigroupSubsetAxiom}[ax]{1,2}}
  \proofstepwidestarbreak[1,2]{U\neq\varnothing}{\FormulaRefAuto{SemigroupSubsemigroupNonemptyAxiom}[ax]{1,2}}
  \proofstepwidestarbreak[5]{x\in U}{\rA}
  \proofstepwidestarbreak[6]{y\in U}{\rA}
  \proofstepwidestarbreak[1,2,5,6]{x\star y\in U}{\FormulaRefAuto{SemigroupSubsemigroupClosureAxiom}[ax]{1,2,5,6}}
  \proofstepwidestarbreak[1,2]{\forall x\in U\,\forall y\in U\,(x\star y\in U)}{\rUI{\rRI{5,\rUI{\rRI{6,7}}}}}
  \proofstepwidestarbreak[1,2]{U\subseteq A\land U\neq\varnothing\land \forall x\in U\,\forall y\in U\,(x\star y\in U)}{\rAI{3,\rAI{4,8}}}
  \proofstepwidestarbreak[1]{\begin{aligned}[t]&\mathsf{UHGr}(U;A,\star)\rightarrow{}\\[-2pt]&\quad (U\subseteq A\land U\neq\varnothing\land \forall x\in U\,\forall y\in U\,(x\star y\in U))\end{aligned}}{\rRI{2,9}}
  \proofstepwidestarbreak[11]{U\subseteq A\land U\neq\varnothing\land \forall x\in U\,\forall y\in U\,(x\star y\in U)}{\rA}
  \proofstepwidestarbreak[11]{U\subseteq A}{\rAEa{11}}
  \proofstepwidestarbreak[11]{U\neq\varnothing}{\rAEa{\rAEb{11}}}
  \proofstepwidestarbreak[11]{\forall x\in U\,\forall y\in U\,(x\star y\in U)}{\rAEb{\rAEb{11}}}
  \proofstepwidestarbreak[15]{x\in U}{\rA}
  \proofstepwidestarbreak[16]{y\in U}{\rA}
  \proofstepwidestarbreak[11,15,16]{x\star y\in U}{\rRE{16,\rUE{\rRE{15,\rUE{14}}}}}
  \proofstepwidestarbreak[1,11]{\mathsf{UHGr}(U;A,\star)}{\FormulaRefAuto{SemigroupSubsemigroupDef}[def]{1,12,13,17}}
  \proofstepwidestarbreak[1]{\begin{aligned}[t]
 &(U\subseteq A\land U\neq\varnothing\land{}\\[-2pt]
 &\quad\forall x\in U\,\forall y\in U\,(x\star y\in U))\rightarrow{}\\[-2pt]
 &\quad\mathsf{UHGr}(U;A,\star)
 \end{aligned}}{\rRI{11,18}}
  \proofstepwidestarbreak[1]{\begin{aligned}[t]&\mathsf{UHGr}(U;A,\star)\leftrightarrow{}\\[-2pt]&\quad (U\subseteq A\land U\neq\varnothing\land \forall x\in U\,\forall y\in U\,(x\star y\in U))\end{aligned}}{\rLRI{10,19}}
\end{tabproofwide}







\FormulaThmDeltaK[Eine Unterhalbgruppe ist eine Halbgruppe]{%
  \SemigroupAlg{A}{\star}\dsep\mathsf{UHGr}(U;A,\star)
  \vdash\SemigroupAlg{U}{\star}%
}{SubsemigroupFormsSemigroup}{%
  \DeltaRow{Mengen}{A\dsep U\dsep x\dsep y\dsep z}
  \DeltaRow{Funktionen}{\star}
}
\begin{tabproofsplitwide}
  \proofpartwideindR[Abgeschlossenheit]{%
    \SemigroupAlg{A}{\star}\dsep\mathsf{UHGr}(U;A,\star)\dsep
    x,y\in U\vdash x\star y\in U
  }{%
    \SemigroupAlg{A}{\star}\dsep\mathsf{UHGr}(U;A,\star)\dsep
    x,y\in U\vdash x\star y\in U
  }[SubsemigroupClosurePart]
    \proofstepwidestar[1]{\SemigroupAlg{A}{\star}}{\rA}
    \proofstepwidestar[2]{\mathsf{UHGr}(U;A,\star)}{\rA}
    \proofstepwidestar[3]{x\in U}{\rA}
    \proofstepwidestar[4]{y\in U}{\rA}
    \proofstepwidestar[1,2,3,4]{x\star y\in U}{%
      \FormulaRefAuto{SemigroupSubsemigroupClosureAxiom}[ax]{1,2,3,4}}
  \closeproofpartwideindR

  \proofpartwideindR[Vererbte Assoziativit\"at]{%
    \begin{aligned}[t]
      &\SemigroupAlg{A}{\star}\dsep\mathsf{UHGr}(U;A,\star)\dsep
        x,y,z\in U\vdash{}\\[-2pt]
      &(x\star y)\star z=x\star(y\star z)
    \end{aligned}
  }{%
    \begin{aligned}[t]
      &\SemigroupAlg{A}{\star}\dsep\mathsf{UHGr}(U;A,\star)\dsep
        x,y,z\in U\vdash{}\\[-2pt]
      &(x\star y)\star z=x\star(y\star z)
    \end{aligned}
  }[SubsemigroupAssociativityPart]
    \proofstepwidestar[1]{\SemigroupAlg{A}{\star}}{\rA}
    \proofstepwidestar[2]{\mathsf{UHGr}(U;A,\star)}{\rA}
    \proofstepwidestar[3]{x\in U}{\rA}
    \proofstepwidestar[4]{y\in U}{\rA}
    \proofstepwidestar[5]{z\in U}{\rA}
    \proofstepwidestar[1,2]{U\subseteq A}{%
      \FormulaRefAuto{SemigroupSubsemigroupSubsetAxiom}[ax]{1,2}}
    \proofstepwidestar[1,2,3]{x\in A}{%
      \FormulaRefAuto{A\subseteq B,\,x\in A\vdash x\in B}[thm]{6,3}}
    \proofstepwidestar[1,2,4]{y\in A}{%
      \FormulaRefAuto{A\subseteq B,\,x\in A\vdash x\in B}[thm]{6,4}}
    \proofstepwidestar[1,2,5]{z\in A}{%
      \FormulaRefAuto{A\subseteq B,\,x\in A\vdash x\in B}[thm]{6,5}}
    \proofstepwidestar[1,2,3,4,5]{%
      (x\star y)\star z=x\star(y\star z)}{%
      \FormulaRefAuto{SemigroupAssociativityAxiom}[ax]{1,7,8,9}}
  \closeproofpartwideindR

  \proofpartwide{Schluss}
    \proofstepwidestar[1]{\SemigroupAlg{A}{\star}}{\rA}
    \proofstepwidestar[2]{\mathsf{UHGr}(U;A,\star)}{\rA}
    \proofstepwidestar[1,2]{\forall x\in U\,\forall y\in U\,(x\star y\in U)}{%
      \FormulaRefAuto{SubsemigroupClosurePart}[thm]{1,2}}
    \proofstepwidestar[1,2]{\forall x,y,z\in U\;((x\star y)\star z=x\star(y\star z))}{%
      \FormulaRefAuto{SubsemigroupAssociativityPart}[thm]{1,2}}
    \proofstepwidestar[1,2]{\SemigroupAlg{U}{\star}}{%
      \FormulaRefAuto{SemigroupDef}[def]{%
        3,
        4}}
  \closeproofpartwide
\end{tabproofsplitwide}

\FormulaThmDeltaK[Der Tr\"ager ist eine Unterhalbgruppe]{%
  \SemigroupAlg{A}{\star}\dsep A\neq\varnothing
  \vdash\mathsf{UHGr}(A;A,\star)%
}{SemigroupCarrierIsSubsemigroup}{%
  \DeltaRow{Mengen}{A\dsep x\dsep y}
  \DeltaRow{Funktionen}{\star}
}
\begin{tabproofwide}
  \proofstepwidestar[1]{\SemigroupAlg{A}{\star}}{\rA}
  \proofstepwidestar[2]{A\neq\varnothing}{\rA}
  \proofstepwidestar[]{A\subseteq A}{%
    \FormulaRefAuto{A\subseteq A}[thm]}
  \proofstepwidestar[4]{x\in A}{\rA}
  \proofstepwidestar[5]{y\in A}{\rA}
  \proofstepwidestar[1,4,5]{x\star y\in A}{%
    \FormulaRefAuto{SemigroupClosureAxiom}[ax]{1,4,5}}
  \proofstepwidestar[1]{%
    \forall x\in A\,\forall y\in A\,(x\star y\in A)}{%
    \rUI{\rRI{4,\rUI{\rRI{5,6}}}}}
  \proofstepwidestar[1,2]{%
    A\subseteq A\land A\neq\varnothing\land
    \forall x\in A\,\forall y\in A\,(x\star y\in A)}{%
    \rAI{3,\rAI{2,7}}}
  \proofstepwidestar[1,2]{\mathsf{UHGr}(A;A,\star)}{%
    \FormulaRefAuto{SemigroupSubsemigroupCriterion}[thm]{1,8}}
\end{tabproofwide}

\subsection{Erzeugte Unterhalbgruppen}

\FormulaDefDeltaK[Von einer nichtleeren Menge erzeugte Unterhalbgruppe]{%
  \begin{aligned}[t]
    &\SemigroupAlg{A}{\star}\dsep X\subseteq A\dsep
      X\neq\varnothing\vdash{}\\[-2pt]
    &\langle X\rangle_{A,\star}\coloneqq
      \bigcap_{\mathsf{UHGr}(U;A,\star)\land X\subseteq U}U
  \end{aligned}
}{SemigroupGeneratedSubsemigroupDef}{%
  \DeltaRow{Mengen}{A\dsep U\dsep X}
  \DeltaRow{Funktionen}{\star}
  \DeltaRow{\textbf{Neues Symbol}}{}
  \DeltaPrem{Erzeugte Unterhalbgruppe}{\langle X\rangle_{A,\star}}
}

\FormulaThmDeltaK[Der Tr\"ager enth\"alt die Erzeuger]{%
  \begin{aligned}[t]
    &\SemigroupAlg{A}{\star}\dsep X\subseteq A\dsep
      X\neq\varnothing\vdash{}\\[-2pt]
    &X\subseteq\langle X\rangle_{A,\star}
  \end{aligned}
}{GeneratedSubsemigroupContainsGenerators}{%
  \DeltaRow{Mengen}{A\dsep U\dsep X\dsep x\dsep u}
  \DeltaRow{Funktionen}{\star}
}
\begin{tabproofwide}
  \proofstepwidestar[1]{\SemigroupAlg{A}{\star}}{\rA}
  \proofstepwidestar[2]{X\subseteq A}{\rA}
  \proofstepwidestar[3]{X\neq\varnothing}{\rA}
  \proofstepwidestar[3]{\exists u\,(u\in X)}{%
    \FormulaRefAuto{A\neq\varnothing\eqvdash\exists x\,(x\in A)}[thm]{3}}
  \proofstepwidestar[5]{u\in X}{\rA}
  \proofstepwidestar[2,5]{u\in A}{%
    \FormulaRefAuto{A\subseteq B,\,x\in A\vdash x\in B}[thm]{2,5}}
  \proofstepwidestar[2,5]{A\neq\varnothing}{%
    \FormulaRefAuto{a\in A\vdash A\neq\varnothing}[thm]{6}}
  \proofstepwidestar[2,3]{A\neq\varnothing}{\rEE{4,5,7}}
  \proofstepwidestar[1,2,3]{\mathsf{UHGr}(A;A,\star)}{%
    \FormulaRefAuto{SemigroupCarrierIsSubsemigroup}[thm]{1,8}}
  \proofstepwidestar[1,2,3]{%
    \mathsf{UHGr}(A;A,\star)\land X\subseteq A}{\rAI{9,2}}
  \proofstepwidestar[11]{x\in X}{\rA}
  \proofstepwidestar[12]{%
    \mathsf{UHGr}(U;A,\star)\land X\subseteq U}{\rA}
  \proofstepwidestar[12]{X\subseteq U}{\rAEb{12}}
  \proofstepwidestar[11,12]{x\in U}{%
    \FormulaRefAuto{A\subseteq B,\,x\in A\vdash x\in B}[thm]{13,11}}
  \proofstepwidestar[11]{%
    \forall U\,\Bigl(
      \mathsf{UHGr}(U;A,\star)\land X\subseteq U
      \rightarrow x\in U\Bigr)}{\rUI{\rRI{12,14}}}
  \proofstepwidestar[1,2,3,11]{%
    x\in\bigcap_{\mathsf{UHGr}(U;A,\star)\land X\subseteq U}U}{%
    \FormulaRefAuto{P(A)\vdash x\in\bigcap_{P(B)}B
      \leftrightarrow\forall C\,(P(C)\rightarrow x\in C)}[thm]{10,15}}
  \proofstepwidestar[1,2,3]{%
    \langle X\rangle_{A,\star}=
    \bigcap_{\mathsf{UHGr}(U;A,\star)\land X\subseteq U}U}{%
    \FormulaRefAuto{SemigroupGeneratedSubsemigroupDef}[def]{1,2,3}}
  \proofstepwidestar[1,2,3,11]{\bigcap_{\mathsf{UHGr}(U;A,\star)\land X\subseteq U}U=\langle X\rangle_{A,\star}}{%
      \FormulaRefAuto{a=b\vdash b=a}[thm]{17}}
  \proofstepwidestar[1,2,3,11]{x\in\langle X\rangle_{A,\star}}{%
    \rIE{18,16}}
  \proofstepwidestar[1,2,3]{X\subseteq\langle X\rangle_{A,\star}}{%
    \FormulaRefAuto{SubsetIntroductionRule}[thm]{[11]\vdots 19}}
\end{tabproofwide}

\FormulaThmDeltaK[Minimalit\"at der erzeugten Unterhalbgruppe]{%
  \begin{aligned}[t]
    &\SemigroupAlg{A}{\star}\dsep X\subseteq A\dsep X\neq\varnothing\dsep
      \mathsf{UHGr}(U;A,\star)\dsep X\subseteq U\vdash{}\\[-2pt]
    &\langle X\rangle_{A,\star}\subseteq U
  \end{aligned}
}{GeneratedSubsemigroupMinimal}{%
  \DeltaRow{Mengen}{A\dsep U\dsep V\dsep X}
  \DeltaRow{Funktionen}{\star}
}
\begin{tabproofwide}
  \proofstepwidestar[1]{\SemigroupAlg{A}{\star}}{\rA}
  \proofstepwidestar[2]{X\subseteq A}{\rA}
  \proofstepwidestar[3]{X\neq\varnothing}{\rA}
  \proofstepwidestar[4]{\mathsf{UHGr}(U;A,\star)}{\rA}
  \proofstepwidestar[5]{X\subseteq U}{\rA}
  \proofstepwidestar[4,5]{%
    \mathsf{UHGr}(U;A,\star)\land X\subseteq U}{\rAI{4,5}}
  \proofstepwidestar[4,5]{%
    \bigcap_{\mathsf{UHGr}(V;A,\star)\land X\subseteq V}V
    \subseteq U}{%
    \FormulaRefAuto{P(C)\vdash\bigcap_{P(A)}A\subseteq C}[thm]{6}}
  \proofstepwidestar[1,2,3]{%
    \langle X\rangle_{A,\star}=
    \bigcap_{\mathsf{UHGr}(V;A,\star)\land X\subseteq V}V}{%
    \FormulaRefAuto{SemigroupGeneratedSubsemigroupDef}[def]{1,2,3}}
  \proofstepwidestar[1,2,3,4,5]{\langle X\rangle_{A,\star}\subseteq U}{%
    \rIE{8,7}}
\end{tabproofwide}

\FormulaThmDeltaK[Produktabschluss der erzeugten Unterhalbgruppe]{%
  \begin{aligned}[t]
    &\SemigroupAlg{A}{\star}\dsep X\subseteq A\dsep X\neq\varnothing\dsep
      x,y\in\langle X\rangle_{A,\star}\vdash{}\\[-2pt]
    &x\star y\in\langle X\rangle_{A,\star}
  \end{aligned}
}{GeneratedSubsemigroupClosure}{%
  \DeltaRow{Mengen}{A\dsep U\dsep X\dsep x\dsep y\dsep u}
  \DeltaRow{Funktionen}{\star}
}
\begin{tabproofsplitwide}
\proofpartwide{Trägerzeugen und Schnittdarstellung}
  \proofstepwidestar[1]{\SemigroupAlg{A}{\star}}{\rA}
  \proofstepwidestar[2]{X\subseteq A}{\rA}
  \proofstepwidestar[3]{X\neq\varnothing}{\rA}
  \proofstepwidestar[4]{x\in\langle X\rangle_{A,\star}}{\rA}
  \proofstepwidestar[5]{y\in\langle X\rangle_{A,\star}}{\rA}
  \proofstepwidestar[3]{\exists u\,(u\in X)}{%
    \FormulaRefAuto{A\neq\varnothing\eqvdash\exists u\,(u\in A)}[thm]{3}}
  \proofstepwidestar[7]{u\in X}{\rA}
  \proofstepwidestar[2,7]{u\in A}{%
    \FormulaRefAuto{A\subseteq B,\,x\in A\vdash x\in B}[thm]{2,7}}
  \proofstepwidestar[2,7]{A\neq\varnothing}{%
    \FormulaRefAuto{a\in A\vdash A\neq\varnothing}[thm]{8}}
  \proofstepwidestar[2,3]{A\neq\varnothing}{\rEE{6,7,9}}
  \proofstepwidestar[1,2,3]{\mathsf{UHGr}(A;A,\star)}{%
    \FormulaRefAuto{SemigroupCarrierIsSubsemigroup}[thm]{1,10}}
  \proofstepwidestar[1,2,3]{%
    \mathsf{UHGr}(A;A,\star)\land X\subseteq A}{\rAI{11,2}}
  \proofstepwidestar[1,2,3]{%
    \langle X\rangle_{A,\star}=
    \bigcap_{\mathsf{UHGr}(U;A,\star)\land X\subseteq U}U}{%
    \FormulaRefAuto{SemigroupGeneratedSubsemigroupDef}[def]{1,2,3}}
  \proofstepwidestar[1,2,3,4]{%
    x\in\bigcap_{\mathsf{UHGr}(U;A,\star)\land X\subseteq U}U}{\rIE{13,4}}
  \proofstepwidestar[1,2,3,5]{%
    y\in\bigcap_{\mathsf{UHGr}(U;A,\star)\land X\subseteq U}U}{\rIE{13,5}}
  \proofstepwidestar[1,2,3,4]{%
    \forall U\,\Bigl(
      \mathsf{UHGr}(U;A,\star)\land X\subseteq U
      \rightarrow x\in U\Bigr)}{%
    \FormulaRefAuto{P(A)\vdash x\in\bigcap_{P(B)}B
      \leftrightarrow\forall C\,(P(C)\rightarrow x\in C)}[thm]{12,14}}
  \proofstepwidestar[1,2,3,5]{%
    \forall U\,\Bigl(
      \mathsf{UHGr}(U;A,\star)\land X\subseteq U
      \rightarrow y\in U\Bigr)}{%
    \FormulaRefAuto{P(A)\vdash x\in\bigcap_{P(B)}B
      \leftrightarrow\forall C\,(P(C)\rightarrow x\in C)}[thm]{12,15}}

\closeproofpartwide
\proofpartwide{Abschluss in jedem Oberträger und Schnittschluss}
\setcounter{proofstepnr}{17}
  \proofstepwidestar[18]{%
    \mathsf{UHGr}(U;A,\star)\land X\subseteq U}{\rA}
  \proofstepwidestar[18]{\mathsf{UHGr}(U;A,\star)}{\rAEa{18}}
  \proofstepwidestar[1,2,3,4,18]{x\in U}{\rRE{18,\rUE{16}}}
  \proofstepwidestar[1,2,3,5,18]{y\in U}{\rRE{18,\rUE{17}}}
  \proofstepwidestar[1,2,3,4,5,18]{x\star y\in U}{%
    \FormulaRefAuto{SemigroupSubsemigroupClosureAxiom}[ax]{1,19,20,21}}
  \proofstepwidestar[1,2,3,4,5]{%
    \forall U\,\Bigl(
      \mathsf{UHGr}(U;A,\star)\land X\subseteq U
      \rightarrow x\star y\in U\Bigr)}{\rUI{\rRI{18,22}}}
  \proofstepwidestar[1,2,3,4,5]{%
    x\star y\in
    \bigcap_{\mathsf{UHGr}(U;A,\star)\land X\subseteq U}U}{%
    \FormulaRefAuto{P(A)\vdash x\in\bigcap_{P(B)}B
      \leftrightarrow\forall C\,(P(C)\rightarrow x\in C)}[thm]{12,23}}
  \proofstepwidestar[1,2,3,4,5]{\bigcap_{\mathsf{UHGr}(U;A,\star)\land X\subseteq U}U=\langle X\rangle_{A,\star}}{%
      \FormulaRefAuto{a=b\vdash b=a}[thm]{13}}
  \proofstepwidestar[1,2,3,4,5]{%
    x\star y\in\langle X\rangle_{A,\star}}{%
    \rIE{25,24}}
\closeproofpartwide
\end{tabproofsplitwide}

\FormulaThmDeltaK[Das Erzeugnis ist eine Unterhalbgruppe]{%
  \begin{aligned}[t]
    &\SemigroupAlg{A}{\star}\dsep X\subseteq A\dsep
      X\neq\varnothing\vdash{}\\[-2pt]
    &\mathsf{UHGr}(\langle X\rangle_{A,\star};A,\star)
  \end{aligned}
}{GeneratedSubsemigroupIsSubsemigroup}{%
  \DeltaRow{Mengen}{A\dsep X\dsep x\dsep y\dsep u}
  \DeltaRow{Funktionen}{\star}
}
\begin{tabproofsplitwide}
  \proofpartwideindR[Träger und Nichtleerheit des Erzeugnisses]{%
    \begin{aligned}[t]
      &\SemigroupAlg{A}{\star}\dsep X\subseteq A\dsep
        X\neq\varnothing\vdash{}\\[-2pt]
      &\langle X\rangle_{A,\star}\subseteq A\land
        \langle X\rangle_{A,\star}\neq\varnothing
    \end{aligned}
  }{%
    \begin{aligned}[t]
      &\SemigroupAlg{A}{\star}\dsep X\subseteq A\dsep
        X\neq\varnothing\vdash{}\\[-2pt]
      &\langle X\rangle_{A,\star}\subseteq A\land
        \langle X\rangle_{A,\star}\neq\varnothing
    \end{aligned}%
  }[GeneratedSubsemigroupCarrierAndNonempty]
    \proofstepwidestar[1]{\SemigroupAlg{A}{\star}}{\rA}
    \proofstepwidestar[2]{X\subseteq A}{\rA}
    \proofstepwidestar[3]{X\neq\varnothing}{\rA}
    \proofstepwidestar[3]{\exists u\,(u\in X)}{%
      \FormulaRefAuto{A\neq\varnothing\eqvdash
        \exists x\,(x\in A)}[thm]{3}}
    \proofstepwidestar[5]{u\in X}{\rA}
    \proofstepwidestar[2,5]{u\in A}{%
      \FormulaRefAuto{A\subseteq B,\,x\in A\vdash x\in B}[thm]{2,5}}
    \proofstepwidestar[2,5]{A\neq\varnothing}{%
      \FormulaRefAuto{a\in A\vdash A\neq\varnothing}[thm]{6}}
    \proofstepwidestar[2,3]{A\neq\varnothing}{\rEE{4,5,7}}
    \proofstepwidestar[1,2,3]{\mathsf{UHGr}(A;A,\star)}{%
      \FormulaRefAuto{SemigroupCarrierIsSubsemigroup}[thm]{1,8}}
    \proofstepwidestarbreak[1,2,3]{%
      \langle X\rangle_{A,\star}\subseteq A}{%
      \FormulaRefAuto{GeneratedSubsemigroupMinimal}[thm]{1,2,3,9,2}}
    \proofstepwidestar[1,2,3]{%
      X\subseteq\langle X\rangle_{A,\star}}{%
      \FormulaRefAuto{GeneratedSubsemigroupContainsGenerators}[thm]{1,2,3}}
    \proofstepwidestarbreak[2,3]{%
      \exists u\,(u\in\langle X\rangle_{A,\star})}{%
      \FormulaRefAuto{A\subseteq B\dsep A\neq\varnothing
        \vdash\exists a\in B\,a\in A}[thm]{11,3}}
    \proofstepwidestar[1,2,3]{%
      \langle X\rangle_{A,\star}\neq\varnothing}{%
      \FormulaRefAuto{A\neq\varnothing\eqvdash
        \exists x\,(x\in A)}[thm]{12}}
    \proofstepwidestarbreak[1,2,3]{%
      \langle X\rangle_{A,\star}\subseteq A\land
      \langle X\rangle_{A,\star}\neq\varnothing}{\rAI{10,13}}
  \closeproofpartwideindR

  \proofpartwideindR[Produktabschluss des Erzeugnisses]{%
    \begin{aligned}[t]
      &\SemigroupAlg{A}{\star}\dsep X\subseteq A\dsep
        X\neq\varnothing\vdash{}\\[-2pt]
      &\forall x\in\langle X\rangle_{A,\star}\,
        \forall y\in\langle X\rangle_{A,\star}\,
        (x\star y\in\langle X\rangle_{A,\star})
    \end{aligned}
  }{%
    \begin{aligned}[t]
      &\SemigroupAlg{A}{\star}\dsep X\subseteq A\dsep
        X\neq\varnothing\vdash{}\\[-2pt]
      &\forall x\in\langle X\rangle_{A,\star}\,
        \forall y\in\langle X\rangle_{A,\star}\,
        (x\star y\in\langle X\rangle_{A,\star})
    \end{aligned}%
  }[GeneratedSubsemigroupProductClosurePart]
    \proofstepwidestar[1]{\SemigroupAlg{A}{\star}}{\rA}
    \proofstepwidestar[2]{X\subseteq A}{\rA}
    \proofstepwidestar[3]{X\neq\varnothing}{\rA}
    \proofstepwidestar[4]{x\in\langle X\rangle_{A,\star}}{\rA}
    \proofstepwidestar[5]{y\in\langle X\rangle_{A,\star}}{\rA}
    \proofstepwidestarbreak[1,2,3,4,5]{%
      x\star y\in\langle X\rangle_{A,\star}}{%
      \FormulaRefAuto{GeneratedSubsemigroupClosure}[thm]{1,2,3,4,5}}
    \proofstepwidestarbreak[1,2,3]{%
      \forall x\in\langle X\rangle_{A,\star}\,
      \forall y\in\langle X\rangle_{A,\star}\,
      (x\star y\in\langle X\rangle_{A,\star})}{%
      \rUI{\rRI{4,\rUI{\rRI{5,6}}}}}
  \closeproofpartwideindR

  \proofpartwideindR[Schluss]{%
    \begin{aligned}[t]
      &\SemigroupAlg{A}{\star}\dsep X\subseteq A\dsep
        X\neq\varnothing\vdash{}\\[-2pt]
      &\mathsf{UHGr}(\langle X\rangle_{A,\star};A,\star)
    \end{aligned}
  }{%
    \begin{aligned}[t]
      &\SemigroupAlg{A}{\star}\dsep X\subseteq A\dsep
        X\neq\varnothing\vdash{}\\[-2pt]
      &\mathsf{UHGr}(\langle X\rangle_{A,\star};A,\star)
    \end{aligned}%
  }[GeneratedSubsemigroupConclusion]
    \proofstepwidestar[1]{\SemigroupAlg{A}{\star}}{\rA}
    \proofstepwidestar[2]{X\subseteq A}{\rA}
    \proofstepwidestar[3]{X\neq\varnothing}{\rA}
    \proofstepwidestarbreak[1,2,3]{%
      \begin{aligned}[t]
        &\langle X\rangle_{A,\star}\subseteq A\land{}\\[-2pt]
        &\langle X\rangle_{A,\star}\neq\varnothing
      \end{aligned}}{%
      \FormulaRefAuto{GeneratedSubsemigroupCarrierAndNonempty}[thm]{1,2,3}}
    \proofstepwidestar[1,2,3]{%
      \langle X\rangle_{A,\star}\subseteq A}{\rAEa{4}}
    \proofstepwidestar[1,2,3]{%
      \langle X\rangle_{A,\star}\neq\varnothing}{\rAEb{4}}
    \proofstepwidestarbreak[1,2,3]{%
      \begin{aligned}[t]
        &\forall x\in\langle X\rangle_{A,\star}\,\\[-2pt]
        &\quad\forall y\in\langle X\rangle_{A,\star}\,
          (x\star y\in\langle X\rangle_{A,\star})
      \end{aligned}}{%
      \FormulaRefAuto{GeneratedSubsemigroupProductClosurePart}[thm]{1,2,3}}
    \proofstepwidestarbreak[1,2,3]{%
      \begin{aligned}[t]
        &\langle X\rangle_{A,\star}\subseteq A\land
          \langle X\rangle_{A,\star}\neq\varnothing\land{}\\[-2pt]
        &\forall x\in\langle X\rangle_{A,\star}\,\\[-2pt]
        &\quad\forall y\in\langle X\rangle_{A,\star}\,
          (x\star y\in\langle X\rangle_{A,\star})
      \end{aligned}}{%
      \rAI{5,\rAI{6,7}}}
    \proofstepwidestarbreak[1,2,3]{%
      \mathsf{UHGr}(\langle X\rangle_{A,\star};A,\star)}{%
      \FormulaRefAuto{SemigroupSubsemigroupCriterion}[thm]{1,8}}
  \closeproofpartwideindR
\end{tabproofsplitwide}

\subsection{Monogene Unterhalbgruppen}

\FormulaDefDeltaK[Von einem Element erzeugte Unterhalbgruppe]{%
  \SemigroupAlg{A}{\star}\dsep a\in A\vdash
  \langle a\rangle_{A,\star}\coloneqq
  \langle\{a\}\rangle_{A,\star}%
}{SemigroupMonogenicSubsemigroupDef}{%
  \DeltaRow{Mengen}{A\dsep a}
  \DeltaRow{Funktionen}{\star}
  \DeltaRow{\textbf{Neues Symbol}}{}
  \DeltaPrem{Monogene Unterhalbgruppe}{\langle a\rangle_{A,\star}}
}

\FormulaThmDeltaK[Das Monogen ist eine Unterhalbgruppe]{%
  \SemigroupAlg{A}{\star}\dsep a\in A\vdash
  \mathsf{UHGr}(\langle a\rangle_{A,\star};A,\star)%
}{MonogenicSubsemigroupIsSubsemigroup}{%
  \DeltaRow{Mengen}{A\dsep a}
  \DeltaRow{Funktionen}{\star}
}
\begin{tabproofwide}
  \proofstepwidestar[1]{\SemigroupAlg{A}{\star}}{\rA}
  \proofstepwidestar[2]{a\in A}{\rA}
  \proofstepwidestar[2]{\{a\}\subseteq A}{%
    \FormulaRefAuto{a\in A\vdash\{a\}\subseteq A}[thm]{2}}
  \proofstepwidestar[]{a\in\{a\}}{\FormulaRefAuto{a\in\{a\}}[thm]}
  \proofstepwidestar[]{\{a\}\neq\varnothing}{%
    \FormulaRefAuto{a\in A\vdash A\neq\varnothing}[thm]{4}}
  \proofstepwidestar[1,2]{%
    \mathsf{UHGr}(\langle\{a\}\rangle_{A,\star};A,\star)}{%
    \FormulaRefAuto{GeneratedSubsemigroupIsSubsemigroup}[thm]{1,3,5}}
  \proofstepwidestar[1,2]{%
    \langle a\rangle_{A,\star}=\langle\{a\}\rangle_{A,\star}}{%
    \FormulaRefAuto{SemigroupMonogenicSubsemigroupDef}[def]{1,2}}
  \proofstepwidestar[1,2]{%
    \mathsf{UHGr}(\langle a\rangle_{A,\star};A,\star)}{\rIE{7,6}}
\end{tabproofwide}

\FormulaDefDeltaK[Menge der positiven Potenzen]{%
  \SemigroupAlg{A}{\star}\dsep a\in A\vdash
  P_{A,\star}(a)\coloneqq
  \bigl\{x\in A\bigm|
    \exists n\in\mathbb N\,(n\neq0\land x=a^n)\bigr\}%
}{SemigroupPositivePowerSetDef}{%
  \DeltaRow{Mengen}{A\dsep a\dsep x\dsep n}
  \DeltaRow{Funktionen}{\star}
  \DeltaRow{\textbf{Neues Symbol}}{}
  \DeltaPrem{Positive Potenzmenge}{P_{A,\star}(a)}
}

\FormulaThmDeltaK[Elementkriterium der positiven Potenzmenge]{%
  \begin{aligned}[t]
    &\SemigroupAlg{A}{\star}\dsep a\in A\vdash{}\\[-2pt]
    &x\in P_{A,\star}(a)\leftrightarrow
      \Bigl(x\in A\land
        \exists n\in\mathbb N\,(n\neq0\land x=a^n)\Bigr)
  \end{aligned}
}{SemigroupPositivePowerSetMembership}{%
  \DeltaRow{Mengen}{A\dsep a\dsep x\dsep n}
  \DeltaRow{Funktionen}{\star}
}
\begin{tabproofwide}
  \proofstepwidestar[1]{\SemigroupAlg{A}{\star}}{\rA}
  \proofstepwidestar[2]{a\in A}{\rA}
  \proofstepwidestar[1,2]{%
    P_{A,\star}(a)=\{u\in A\mid
      \exists n\in\mathbb N\,(n\neq0\land u=a^n)\}}{%
    \FormulaRefAuto{SemigroupPositivePowerSetDef}[def]{1,2}}
  \proofstepwidestar[1,2]{x\in\{u\in A\mid\exists n\in\mathbb N\;(n\neq0\land u=a^n)\}\leftrightarrow(x\in A\land\exists n\in\mathbb N\;(n\neq0\land x=a^n))}{%
      \FormulaRefAuto{x\in\{u\in A\mid P(u)\}\eqvdash
      x\in A\land P(x)}[thm]}
  \proofstepwidestar[1,2]{%
    x\in P_{A,\star}(a)\leftrightarrow
    \Bigl(x\in A\land
      \exists n\in\mathbb N\,(n\neq0\land x=a^n)\Bigr)}{%
    \rIE{3,4}}
\end{tabproofwide}

\FormulaThmDeltaK[Die positiven Potenzen bilden eine Unterhalbgruppe]{%
  \SemigroupAlg{A}{\star}\dsep a\in A\vdash
  \mathsf{UHGr}(P_{A,\star}(a);A,\star)%
}{SemigroupPositivePowersFormSubsemigroup}{%
  \DeltaRow{Mengen}{A\dsep a\dsep x\dsep y\dsep m\dsep n}
  \DeltaRow{Funktionen}{\star}
}
\begin{tabproofsplitwide}
  \proofpartwideindR[Teilmenge und Nichtleerheit]{%
    \begin{aligned}[t]
      &\SemigroupAlg{A}{\star}\dsep a\in A\vdash{}\\[-2pt]
      &P_{A,\star}(a)\subseteq A\land
        P_{A,\star}(a)\neq\varnothing
    \end{aligned}
  }{%
    \begin{aligned}[t]
      &\SemigroupAlg{A}{\star}\dsep a\in A\vdash{}\\[-2pt]
      &P_{A,\star}(a)\subseteq A\land
        P_{A,\star}(a)\neq\varnothing
    \end{aligned}%
  }[SemigroupPositivePowerSetCarrierAndNonempty]
    \proofstepwidestar[1]{\SemigroupAlg{A}{\star}}{\rA}
    \proofstepwidestar[2]{a\in A}{\rA}
    \proofstepwidestar[]{P_{A,\star}(a)\subseteq A}{%
      \FormulaRefAuto{\{x\in A\mid P(x)\}\subseteq A}[thm]}
    \proofstepwidestar[]{1\in\mathbb N}{%
      \FormulaRefAuto{PeanoOneInN}[thm]}
    \proofstepwidestar[]{1\neq0}{%
      \FormulaRefAuto{PeanoOneNeZero}[thm]}
    \proofstepwidestar[1,2]{a^1=a}{%
      \FormulaRefAuto{SemigroupPowerOne}[thm]{1,2}}
    \proofstepwidestar[1,2]{a=a^1}{%
      \FormulaRefAuto{a=b\vdash b=a}[thm]{6}}
    \proofstepwidestarbreak[1,2]{%
      a\in A\land
      \exists n\in\mathbb N\,(n\neq0\land a=a^n)}{%
      \rAI{2,\rEI{4,
        \rAI{5,7}}}}
    \proofstepwidestar[1,2]{a\in P_{A,\star}(a)}{%
      \FormulaRefAuto{SemigroupPositivePowerSetMembership}[thm]{1,2,8}}
    \proofstepwidestar[1,2]{P_{A,\star}(a)\neq\varnothing}{%
      \FormulaRefAuto{a\in A\vdash A\neq\varnothing}[thm]{9}}
    \proofstepwidestarbreak[1,2]{%
      P_{A,\star}(a)\subseteq A\land
      P_{A,\star}(a)\neq\varnothing}{\rAI{3,10}}
  \closeproofpartwideindR

  \proofpartwideindR[Produktabschluss der positiven Potenzen]{%
    \begin{aligned}[t]
      &\SemigroupAlg{A}{\star}\dsep a\in A\vdash{}\\[-2pt]
      &\forall x\in P_{A,\star}(a)\,
        \forall y\in P_{A,\star}(a)\,
        (x\star y\in P_{A,\star}(a))
    \end{aligned}
  }{%
    \begin{aligned}[t]
      &\SemigroupAlg{A}{\star}\dsep a\in A\vdash{}\\[-2pt]
      &\forall x\in P_{A,\star}(a)\,
        \forall y\in P_{A,\star}(a)\,
        (x\star y\in P_{A,\star}(a))
    \end{aligned}%
  }[SemigroupPositivePowerSetProductClosure]
    \proofstepwidestar[1]{\SemigroupAlg{A}{\star}}{\rA}
    \proofstepwidestar[2]{a\in A}{\rA}
    \proofstepwidestar[3]{x\in P_{A,\star}(a)}{\rA}
    \proofstepwidestar[4]{y\in P_{A,\star}(a)}{\rA}
    \proofstepwidestar[1,2,3]{x\in A\land\exists m\in\mathbb N\;(m\neq0\land x=a^m)}{%
      \FormulaRefAuto{SemigroupPositivePowerSetMembership}[thm]{%
        1,2,3}}
    \proofstepwidestarbreak[1,2,3]{%
      \exists m\in\mathbb N\,(m\neq0\land x=a^m)}{%
      \rAEb{5}}
    \proofstepwidestar[1,2,4]{y\in A\land\exists n\in\mathbb N\;(n\neq0\land y=a^n)}{%
      \FormulaRefAuto{SemigroupPositivePowerSetMembership}[thm]{%
        1,2,4}}
    \proofstepwidestarbreak[1,2,4]{%
      \exists n\in\mathbb N\,(n\neq0\land y=a^n)}{%
      \rAEb{7}}
    \proofstepwidestar[9]{%
      m\in\mathbb N\land(m\neq0\land x=a^m)}{\rA}
    \proofstepwidestar[10]{%
      n\in\mathbb N\land(n\neq0\land y=a^n)}{\rA}
    \proofstepwidestar[9]{m\in\mathbb N}{\rAEa{9}}
    \proofstepwidestar[9]{m\neq0}{\rAEa{\rAEb{9}}}
    \proofstepwidestar[9]{x=a^m}{\rAEb{\rAEb{9}}}
    \proofstepwidestar[10]{n\in\mathbb N}{\rAEa{10}}
    \proofstepwidestar[10]{n\neq0}{\rAEa{\rAEb{10}}}
    \proofstepwidestar[10]{y=a^n}{\rAEb{\rAEb{10}}}
    \proofstepwidestar[9,10]{m+n\in\mathbb N}{%
      \FormulaRefAuto{PeanoAddClosure}[thm]{11,14}}
    \proofstepwidestar[9,10]{m+n\neq0}{%
      \FormulaRefAuto{PeanoAddNonzeroLeft}[thm]{11,14,12}}
    \proofstepwidestarbreak[1,2,9,10]{a^{m+n}=a^m\star a^n}{%
      \FormulaRefAuto{SemigroupPowerAdditionLaw}[thm]{1,2,11,14,12,15}}
    \proofstepwidestar[1,2,9,10]{x\star y=a^{m+n}}{\rIE{13,16,19}}
    \proofstepwidestar[1,2,9,10]{a^{m+n}\in A}{%
      \FormulaRefAuto{SemigroupPowerClosure}[thm]{1,2,17,18}}
    \proofstepwidestarbreak[1,2,9,10]{x\star y\in A}{%
      \rIE{20,21}}
    \proofstepwidestarbreak[1,2,9,10]{%
      x\star y\in A\land
      \exists k\in\mathbb N\,(k\neq0\land x\star y=a^k)}{%
      \rAI{22,\rEI{17,\rAI{18,20}}}}
    \proofstepwidestarbreak[1,2,9,10]{x\star y\in P_{A,\star}(a)}{%
      \FormulaRefAuto{SemigroupPositivePowerSetMembership}[thm]{1,2,23}}
    \proofstepwidestar[1,2,3,4]{x\star y\in P_{A,\star}(a)}{%
      \rEE{6,9,\rEE{8,10,24}}}
    \proofstepwidestarbreak[1,2]{%
      \forall x\in P_{A,\star}(a)\,
      \forall y\in P_{A,\star}(a)\,
      (x\star y\in P_{A,\star}(a))}{%
      \rUI{\rRI{3,\rUI{\rRI{4,25}}}}}
  \closeproofpartwideindR

  \proofpartwideindR[Schluss]{%
    \begin{aligned}[t]
      &\SemigroupAlg{A}{\star}\dsep a\in A\vdash{}\\[-2pt]
      &\mathsf{UHGr}(P_{A,\star}(a);A,\star)
    \end{aligned}
  }{%
    \begin{aligned}[t]
      &\SemigroupAlg{A}{\star}\dsep a\in A\vdash{}\\[-2pt]
      &\mathsf{UHGr}(P_{A,\star}(a);A,\star)
    \end{aligned}%
  }[SemigroupPositivePowerSetSubsemigroupConclusion]
    \proofstepwidestar[1]{\SemigroupAlg{A}{\star}}{\rA}
    \proofstepwidestar[2]{a\in A}{\rA}
    \proofstepwidestarbreak[1,2]{%
      \begin{aligned}[t]
        &P_{A,\star}(a)\subseteq A\land{}\\[-2pt]
        &P_{A,\star}(a)\neq\varnothing
      \end{aligned}}{%
      \FormulaRefAuto{SemigroupPositivePowerSetCarrierAndNonempty}[thm]{1,2}}
    \proofstepwidestar[1,2]{P_{A,\star}(a)\subseteq A}{\rAEa{3}}
    \proofstepwidestar[1,2]{P_{A,\star}(a)\neq\varnothing}{\rAEb{3}}
    \proofstepwidestarbreak[1,2]{%
      \begin{aligned}[t]
        &\forall x\in P_{A,\star}(a)\,\\[-2pt]
        &\quad\forall y\in P_{A,\star}(a)\,
          (x\star y\in P_{A,\star}(a))
      \end{aligned}}{%
      \FormulaRefAuto{SemigroupPositivePowerSetProductClosure}[thm]{1,2}}
    \proofstepwidestarbreak[1,2]{%
      \begin{aligned}[t]
        &P_{A,\star}(a)\subseteq A\land
          P_{A,\star}(a)\neq\varnothing\land{}\\[-2pt]
        &\forall x\in P_{A,\star}(a)\,\\[-2pt]
        &\quad\forall y\in P_{A,\star}(a)\,
          (x\star y\in P_{A,\star}(a))
      \end{aligned}}{%
      \rAI{4,\rAI{5,6}}}
    \proofstepwidestarbreak[1,2]{%
      \mathsf{UHGr}(P_{A,\star}(a);A,\star)}{%
      \FormulaRefAuto{SemigroupSubsemigroupCriterion}[thm]{1,7}}
  \closeproofpartwideindR
\end{tabproofsplitwide}

\FormulaThmDeltaK[Beschreibung der monogenen Unterhalbgruppe durch Potenzen]{%
  \SemigroupAlg{A}{\star}\dsep a\in A\vdash
  \langle a\rangle_{A,\star}=P_{A,\star}(a)%
}{SemigroupMonogenicPowerDescription}{%
  \DeltaRow{Mengen}{A\dsep a\dsep x\dsep n\dsep k\dsep j}
  \DeltaRow{Funktionen}{\star}
}
\begin{tabproofsplitwide}
  \proofpartwideindR[Positive Potenzen liegen im Monogen]{%
\SemigroupAlg{A}{\star}\dsep a\in A\dsep k\in\mathbb N\dsep k\neq0\vdash a^k\in\langle a\rangle_{A,\star}
}{%
\SemigroupAlg{A}{\star}\dsep a\in A\dsep k\in\mathbb N\dsep k\neq0\vdash a^k\in\langle a\rangle_{A,\star}
}[SemigroupPowerInMonogenicSubsemigroup]
  \proofstepwidestar[1]{\SemigroupAlg{A}{\star}}{\rA}
  \proofstepwidestar[2]{a\in A}{\rA}
  \proofstepwidestar[3]{k\in\mathbb N}{\rA}
  \proofstepwidestar[4]{k\neq0}{\rA}
  \proofstepwidestar[2]{\{a\}\subseteq A}{\FormulaRefAuto{a\in A\vdash\{a\}\subseteq A}[thm]{2}}
  \proofstepwidestar[]{a\in\{a\}}{%
      \FormulaRefAuto{a\in\{a\}}[thm]}
  \proofstepwidestar[]{\{a\}\neq\varnothing}{\FormulaRefAuto{a\in A\vdash A\neq\varnothing}[thm]{6}}
  \proofstepwidestar[1,2]{a\in\{a\}}{%
      \FormulaRefAuto{a\in\{a\}}[thm]}
  \proofstepwidestar[1,2]{a\in\langle\{a\}\rangle_{A,\star}}{\FormulaRefAuto{GeneratedSubsemigroupContainsGenerators}[thm]{1,5,7,8}}
  \proofstepwidestar[1,2]{\langle a\rangle_{A,\star}=\langle\{a\}\rangle_{A,\star}}{%
      \FormulaRefAuto{SemigroupMonogenicSubsemigroupDef}[def]{1,2}}
  \proofstepwidestar[1,2]{a\in\langle a\rangle_{A,\star}}{\rIE{10,9}}
  \proofstepwidestar[1,2]{a^1=a}{%
      \FormulaRefAuto{SemigroupPowerOne}[thm]{1,2}}
  \proofstepwidestar[1,2]{a^1\in\langle a\rangle_{A,\star}}{\rIE{12,11}}
  \proofstepwidestar[1,2]{\mathsf{UHGr}(\langle a\rangle_{A,\star};A,\star)}{\FormulaRefAuto{MonogenicSubsemigroupIsSubsemigroup}[thm]{1,2}}
  \proofstepwidestar[15]{j\in\mathbb N}{\rA}
  \proofstepwidestar[16]{j\neq0}{\rA}
  \proofstepwidestar[17]{a^j\in\langle a\rangle_{A,\star}}{\rA}
  \proofstepwidestar[1,2,15,16]{a^{j+1}=a^j\star a}{\FormulaRefAuto{SemigroupPowerSuccessor}[thm]{1,2,15,16}}
  \proofstepwidestar[1,2,17]{a^j\star a\in\langle a\rangle_{A,\star}}{\FormulaRefAuto{SemigroupSubsemigroupClosureAxiom}[ax]{1,14,17,11}}
  \proofstepwidestar[1,2,15,16,17]{a^{j+1}\in\langle a\rangle_{A,\star}}{\rIE{18,19}}
    \proofstepwidestarbreak[1,2]{\begin{aligned}[t]&\forall j\in\mathbb N\,\bigl(j\neq0\rightarrow{}\\[-2pt]&\quad(a^j\in\langle a\rangle_{A,\star}\rightarrow a^{j+1}\in\langle a\rangle_{A,\star})\bigr)\end{aligned}}{%
      \rUI{\rRI{15,\rRI{16,\rRI{17,20}}}}}
    \proofstepwidestarbreak[1,2]{\forall j\in\mathbb N\,(j\neq0\rightarrow a^j\in\langle a\rangle_{A,\star})}{%
      \FormulaRefAuto{PeanoInductionFromOne}[thm]{13,21}}
    \proofstepwidestarbreak[1,2,3,4]{a^k\in\langle a\rangle_{A,\star}}{%
      \rRE{4,\rRE{3,\rUE{22}}}}
  \closeproofpartwideindR

  \proofpartwideindR[Monogen liegt in der Menge positiver Potenzen]{%
    \begin{aligned}[t]
      &\SemigroupAlg{A}{\star}\dsep a\in A\vdash{}\\[-2pt]
      &\langle a\rangle_{A,\star}\subseteq P_{A,\star}(a)
    \end{aligned}
  }{%
    \SemigroupAlg{A}{\star}\dsep a\in A\vdash
    \langle a\rangle_{A,\star}\subseteq P_{A,\star}(a)%
  }[SemigroupMonogenicSubsemigroupInPositivePowers]
    \proofstepwidestar[1]{\SemigroupAlg{A}{\star}}{\rA}
    \proofstepwidestar[2]{a\in A}{\rA}
    \proofstepwidestarbreak[1,2]{\mathsf{UHGr}(P_{A,\star}(a);A,\star)}{%
      \FormulaRefAuto{SemigroupPositivePowersFormSubsemigroup}[thm]{1,2}}
    \proofstepwidestar[]{1\in\mathbb N}{\FormulaRefAuto{PeanoOneInN}[thm]}
    \proofstepwidestar[]{1\neq0}{%
      \FormulaRefAuto{PeanoOneNeZero}[thm]}
    \proofstepwidestar[1,2]{a^1=a}{%
      \FormulaRefAuto{SemigroupPowerOne}[thm]{1,2}}
    \proofstepwidestar[1,2]{a=a^1}{%
      \FormulaRefAuto{a=b\vdash b=a}[thm]{6}}
    \proofstepwidestarbreak[1,2]{%
      \exists n\in\mathbb N\,(n\neq0\land a=a^n)}{%
      \rEI{4,\rAI{5,7}}}
    \proofstepwidestar[1,2]{%
      a\in A\land\exists n\in\mathbb N\,(n\neq0\land a=a^n)}{%
      \rAI{2,8}}
    \proofstepwidestarbreak[1,2]{a\in P_{A,\star}(a)}{%
      \FormulaRefAuto{SemigroupPositivePowerSetMembership}[thm]{1,2,9}}
    \proofstepwidestarbreak[1,2]{\{a\}\subseteq P_{A,\star}(a)}{%
      \FormulaRefAuto{a\in A\vdash\{a\}\subseteq A}[thm]{10}}
    \proofstepwidestar[2]{\{a\}\subseteq A}{%
      \FormulaRefAuto{a\in A\vdash\{a\}\subseteq A}[thm]{2}}
    \proofstepwidestar[]{a\in\{a\}}{%
      \FormulaRefAuto{a\in\{a\}}[thm]}
    \proofstepwidestar[]{\{a\}\neq\varnothing}{%
      \FormulaRefAuto{a\in A\vdash A\neq\varnothing}[thm]{%
        13}}
    \proofstepwidestarbreak[1,2]{%
      \langle\{a\}\rangle_{A,\star}\subseteq P_{A,\star}(a)}{%
      \FormulaRefAuto{GeneratedSubsemigroupMinimal}[thm]{1,12,14,3,11}}
    \proofstepwidestar[1,2]{\langle a\rangle_{A,\star}=\langle\{a\}\rangle_{A,\star}}{%
      \FormulaRefAuto{SemigroupMonogenicSubsemigroupDef}[def]{1,2}}
    \proofstepwidestarbreak[1,2]{%
      \langle a\rangle_{A,\star}\subseteq P_{A,\star}(a)}{%
      \rIE{16,15}}
  \closeproofpartwideindR

  \proofpartwideindR[Menge positiver Potenzen liegt im Monogen]{%
    \begin{aligned}[t]
      &\SemigroupAlg{A}{\star}\dsep a\in A\vdash{}\\[-2pt]
      &P_{A,\star}(a)\subseteq\langle a\rangle_{A,\star}
    \end{aligned}
  }{%
    \SemigroupAlg{A}{\star}\dsep a\in A\vdash
    P_{A,\star}(a)\subseteq\langle a\rangle_{A,\star}%
  }[SemigroupPositivePowersInMonogenicSubsemigroup]
    \proofstepwidestar[1]{\SemigroupAlg{A}{\star}}{\rA}
    \proofstepwidestar[2]{a\in A}{\rA}
    \proofstepwidestar[3]{x\in P_{A,\star}(a)}{\rA}
    \proofstepwidestar[1,2,3]{x\in A\land\exists n\in\mathbb N\;(n\neq0\land x=a^n)}{%
      \FormulaRefAuto{SemigroupPositivePowerSetMembership}[thm]{1,2,3}}
    \proofstepwidestarbreak[1,2,3]{%
      \exists n\in\mathbb N\,(n\neq0\land x=a^n)}{%
      \rAEb{4}}
    \proofstepwidestar[6]{n\in\mathbb N\land(n\neq0\land x=a^n)}{\rA}
    \proofstepwidestar[6]{n\in\mathbb N}{\rAEa{6}}
    \proofstepwidestar[6]{n\neq0}{\rAEa{\rAEb{6}}}
    \proofstepwidestar[6]{x=a^n}{\rAEb{\rAEb{6}}}
  \proofstepwidestar[1,2,6]{a^n\in\langle a\rangle_{A,\star}}{\FormulaRefAuto{SemigroupPowerInMonogenicSubsemigroup}[thm]{1,2,7,8}}
  \proofstepwidestar[1,2,6]{x\in\langle a\rangle_{A,\star}}{\rIE{9,10}}
  \proofstepwidestar[1,2,3]{x\in\langle a\rangle_{A,\star}}{\rEE{5,6,11}}
  \proofstepwidestar[1,2]{P_{A,\star}(a)\subseteq\langle a\rangle_{A,\star}}{\FormulaRefAuto{SubsetIntroductionRule}[thm]{[3]\vdots 12}}
  \closeproofpartwideindR

  \proofpartwideindR[Gleichheit aus beiden Inklusionen]{%
    \begin{aligned}[t]
      &\SemigroupAlg{A}{\star}\dsep a\in A\vdash{}\\[-2pt]
      &\langle a\rangle_{A,\star}=P_{A,\star}(a)
    \end{aligned}
  }{%
    \SemigroupAlg{A}{\star}\dsep a\in A\vdash
    \langle a\rangle_{A,\star}=P_{A,\star}(a)%
  }[SemigroupMonogenicPowerDescriptionConclusion]
    \proofstepwidestar[1]{\SemigroupAlg{A}{\star}}{\rA}
    \proofstepwidestar[2]{a\in A}{\rA}
    \proofstepwidestarbreak[1,2]{%
      \langle a\rangle_{A,\star}\subseteq P_{A,\star}(a)}{%
      \FormulaRefAuto{SemigroupMonogenicSubsemigroupInPositivePowers}[thm]{1,2}}
    \proofstepwidestarbreak[1,2]{%
      P_{A,\star}(a)\subseteq\langle a\rangle_{A,\star}}{%
      \FormulaRefAuto{SemigroupPositivePowersInMonogenicSubsemigroup}[thm]{1,2}}
    \proofstepwidestarbreak[1,2]{%
      \langle a\rangle_{A,\star}=P_{A,\star}(a)}{%
      \FormulaRefAuto{A\subseteq B,\,B\subseteq A\vdash A=B}[thm]{3,4}}
  \closeproofpartwideindR
\end{tabproofsplitwide}

\section{Ideale}

Im Unterschied zu Unterhalbgruppen werden Ideale durch Multiplikation mit
beliebigen Elementen der umgebenden Halbgruppe absorbiert. Wir verlangen
ausdr\"ucklich Nichtleerheit; andernfalls w\"are die leere Menge stets das
kleinste Ideal und die sp\"atere Minimalidealtheorie w\"urde entarten.

\par\noindent\begin{minipage}{\linewidth}
\FormulaDefDeltaK[Linksideal]{%
\mathsf{LId}(I;A,\star)%
}{SemigroupLeftIdealDef}{%
  \DeltaRow{Mengen}{A\dsep I\dsep a\dsep x}
  \DeltaRow{Funktionen}{\star}
  \DeltaRow{\textbf{Voraussetzung}}{}
  \DeltaPrem{Umgebende Halbgruppe}{\SemigroupAlg{A}{\star}}
  \DeltaRow{\textbf{Axiome}}{}
  \DeltaRow{Träger}{I\subseteq A}
  \DeltaRow{Nichtleerheit}{I\neq\varnothing}
  \DeltaRow{Linksabsorption}{a\in A\dsep x\in I\vdash a\star x\in I}
  \DeltaRow{\textbf{Neues Symbol}}{}
  \DeltaPrem{Linksideal}{\mathsf{LId}(I;A,\star)}
}
\end{minipage}\par

\par\noindent\begin{minipage}{\linewidth}
\FormulaAxiomDeltaK[Träger eines Linksideals]{%
\SemigroupAlg{A}{\star}\dsep \mathsf{LId}(I;A,\star)\vdash I\subseteq A%
}{SemigroupLeftIdealSubsetAxiom}{%
  \DeltaRow{Mengen}{A\dsep I\dsep a\dsep x}
  \DeltaRow{Funktionen}{\star}
}
\end{minipage}\par

\par\noindent\begin{minipage}{\linewidth}
\FormulaAxiomDeltaK[Nichtleerheit eines Linksideals]{%
\SemigroupAlg{A}{\star}\dsep \mathsf{LId}(I;A,\star)\vdash I\neq\varnothing%
}{SemigroupLeftIdealNonemptyAxiom}{%
  \DeltaRow{Mengen}{A\dsep I\dsep a\dsep x}
  \DeltaRow{Funktionen}{\star}
}
\end{minipage}\par

\par\noindent\begin{minipage}{\linewidth}
\FormulaAxiomDeltaK[Linksabsorption eines Linksideals]{%
\SemigroupAlg{A}{\star}\dsep \mathsf{LId}(I;A,\star)\dsep a\in A\dsep x\in I\vdash a\star x\in I%
}{SemigroupLeftIdealClosureAxiom}{%
  \DeltaRow{Mengen}{A\dsep I\dsep a\dsep x}
  \DeltaRow{Funktionen}{\star}
}
\end{minipage}\par

\par\noindent\begin{minipage}{\linewidth}
\FormulaThmDeltaK[Kriterium für ein Linksideal]{%
\begin{aligned}[t]&\SemigroupAlg{A}{\star}\vdash \mathsf{LId}(I;A,\star)\leftrightarrow{}\\[-2pt]&\quad \bigl(I\subseteq A\land I\neq\varnothing\land \forall a\in A\,\forall x\in I\,(a\star x\in I)\bigr)\end{aligned}%
}{SemigroupLeftIdealCriterion}{%
  \DeltaRow{Mengen}{A\dsep I\dsep a\dsep x}
  \DeltaRow{Funktionen}{\star}
}
\end{minipage}\par
\begin{tabproofwide}
  \proofstepwidestarbreak[1]{\SemigroupAlg{A}{\star}}{\rA}
  \proofstepwidestarbreak[2]{\mathsf{LId}(I;A,\star)}{\rA}
  \proofstepwidestarbreak[1,2]{I\subseteq A}{\FormulaRefAuto{SemigroupLeftIdealSubsetAxiom}[ax]{1,2}}
  \proofstepwidestarbreak[1,2]{I\neq\varnothing}{\FormulaRefAuto{SemigroupLeftIdealNonemptyAxiom}[ax]{1,2}}
  \proofstepwidestarbreak[5]{a\in A}{\rA}
  \proofstepwidestarbreak[6]{x\in I}{\rA}
  \proofstepwidestarbreak[1,2,5,6]{a\star x\in I}{\FormulaRefAuto{SemigroupLeftIdealClosureAxiom}[ax]{1,2,5,6}}
  \proofstepwidestarbreak[1,2]{\forall a\in A\,\forall x\in I\,(a\star x\in I)}{\rUI{\rRI{5,\rUI{\rRI{6,7}}}}}
  \proofstepwidestarbreak[1,2]{I\subseteq A\land I\neq\varnothing\land \forall a\in A\,\forall x\in I\,(a\star x\in I)}{\rAI{3,\rAI{4,8}}}
  \proofstepwidestarbreak[1]{\begin{aligned}[t]&\mathsf{LId}(I;A,\star)\rightarrow{}\\[-2pt]&\quad (I\subseteq A\land I\neq\varnothing\land \forall a\in A\,\forall x\in I\,(a\star x\in I))\end{aligned}}{\rRI{2,9}}
  \proofstepwidestarbreak[11]{I\subseteq A\land I\neq\varnothing\land \forall a\in A\,\forall x\in I\,(a\star x\in I)}{\rA}
  \proofstepwidestarbreak[11]{I\subseteq A}{\rAEa{11}}
  \proofstepwidestarbreak[11]{I\neq\varnothing}{\rAEa{\rAEb{11}}}
  \proofstepwidestarbreak[11]{\forall a\in A\,\forall x\in I\,(a\star x\in I)}{\rAEb{\rAEb{11}}}
  \proofstepwidestarbreak[15]{a\in A}{\rA}
  \proofstepwidestarbreak[16]{x\in I}{\rA}
  \proofstepwidestarbreak[11,15,16]{a\star x\in I}{\rRE{16,\rUE{\rRE{15,\rUE{14}}}}}
  \proofstepwidestarbreak[1,11]{\mathsf{LId}(I;A,\star)}{\FormulaRefAuto{SemigroupLeftIdealDef}[def]{1,12,13,17}}
  \proofstepwidestarbreak[1]{\begin{aligned}[t]&(I\subseteq A\land I\neq\varnothing\land \forall a\in A\,\forall x\in I\,(a\star x\in I))\rightarrow{}\\[-2pt]&\quad \mathsf{LId}(I;A,\star)\end{aligned}}{\rRI{11,18}}
  \proofstepwidestarbreak[1]{\begin{aligned}[t]&\mathsf{LId}(I;A,\star)\leftrightarrow{}\\[-2pt]&\quad (I\subseteq A\land I\neq\varnothing\land \forall a\in A\,\forall x\in I\,(a\star x\in I))\end{aligned}}{\rLRI{10,19}}
\end{tabproofwide}

\par\noindent\begin{minipage}{\linewidth}
\FormulaDefDeltaK[Rechtsideal]{%
\mathsf{RId}(I;A,\star)%
}{SemigroupRightIdealDef}{%
  \DeltaRow{Mengen}{A\dsep I\dsep x\dsep a}
  \DeltaRow{Funktionen}{\star}
  \DeltaRow{\textbf{Voraussetzung}}{}
  \DeltaPrem{Umgebende Halbgruppe}{\SemigroupAlg{A}{\star}}
  \DeltaRow{\textbf{Axiome}}{}
  \DeltaRow{Träger}{I\subseteq A}
  \DeltaRow{Nichtleerheit}{I\neq\varnothing}
  \DeltaRow{Rechtsabsorption}{x\in I\dsep a\in A\vdash x\star a\in I}
  \DeltaRow{\textbf{Neues Symbol}}{}
  \DeltaPrem{Rechtsideal}{\mathsf{RId}(I;A,\star)}
}
\end{minipage}\par

\par\noindent\begin{minipage}{\linewidth}
\FormulaAxiomDeltaK[Träger eines Rechtsideals]{%
\SemigroupAlg{A}{\star}\dsep \mathsf{RId}(I;A,\star)\vdash I\subseteq A%
}{SemigroupRightIdealSubsetAxiom}{%
  \DeltaRow{Mengen}{A\dsep I\dsep x\dsep a}
  \DeltaRow{Funktionen}{\star}
}
\end{minipage}\par

\par\noindent\begin{minipage}{\linewidth}
\FormulaAxiomDeltaK[Nichtleerheit eines Rechtsideals]{%
\SemigroupAlg{A}{\star}\dsep \mathsf{RId}(I;A,\star)\vdash I\neq\varnothing%
}{SemigroupRightIdealNonemptyAxiom}{%
  \DeltaRow{Mengen}{A\dsep I\dsep x\dsep a}
  \DeltaRow{Funktionen}{\star}
}
\end{minipage}\par

\par\noindent\begin{minipage}{\linewidth}
\FormulaAxiomDeltaK[Rechtsabsorption eines Rechtsideals]{%
\SemigroupAlg{A}{\star}\dsep \mathsf{RId}(I;A,\star)\dsep x\in I\dsep a\in A\vdash x\star a\in I%
}{SemigroupRightIdealClosureAxiom}{%
  \DeltaRow{Mengen}{A\dsep I\dsep x\dsep a}
  \DeltaRow{Funktionen}{\star}
}
\end{minipage}\par

\par\noindent\begin{minipage}{\linewidth}
\FormulaThmDeltaK[Kriterium für ein Rechtsideal]{%
\begin{aligned}[t]&\SemigroupAlg{A}{\star}\vdash \mathsf{RId}(I;A,\star)\leftrightarrow{}\\[-2pt]&\quad \bigl(I\subseteq A\land I\neq\varnothing\land \forall x\in I\,\forall a\in A\,(x\star a\in I)\bigr)\end{aligned}%
}{SemigroupRightIdealCriterion}{%
  \DeltaRow{Mengen}{A\dsep I\dsep x\dsep a}
  \DeltaRow{Funktionen}{\star}
}
\end{minipage}\par
\begin{tabproofwide}
  \proofstepwidestarbreak[1]{\SemigroupAlg{A}{\star}}{\rA}
  \proofstepwidestarbreak[2]{\mathsf{RId}(I;A,\star)}{\rA}
  \proofstepwidestarbreak[1,2]{I\subseteq A}{\FormulaRefAuto{SemigroupRightIdealSubsetAxiom}[ax]{1,2}}
  \proofstepwidestarbreak[1,2]{I\neq\varnothing}{\FormulaRefAuto{SemigroupRightIdealNonemptyAxiom}[ax]{1,2}}
  \proofstepwidestarbreak[5]{x\in I}{\rA}
  \proofstepwidestarbreak[6]{a\in A}{\rA}
  \proofstepwidestarbreak[1,2,5,6]{x\star a\in I}{\FormulaRefAuto{SemigroupRightIdealClosureAxiom}[ax]{1,2,5,6}}
  \proofstepwidestarbreak[1,2]{\forall x\in I\,\forall a\in A\,(x\star a\in I)}{\rUI{\rRI{5,\rUI{\rRI{6,7}}}}}
  \proofstepwidestarbreak[1,2]{I\subseteq A\land I\neq\varnothing\land \forall x\in I\,\forall a\in A\,(x\star a\in I)}{\rAI{3,\rAI{4,8}}}
  \proofstepwidestarbreak[1]{\begin{aligned}[t]&\mathsf{RId}(I;A,\star)\rightarrow{}\\[-2pt]&\quad (I\subseteq A\land I\neq\varnothing\land \forall x\in I\,\forall a\in A\,(x\star a\in I))\end{aligned}}{\rRI{2,9}}
  \proofstepwidestarbreak[11]{I\subseteq A\land I\neq\varnothing\land \forall x\in I\,\forall a\in A\,(x\star a\in I)}{\rA}
  \proofstepwidestarbreak[11]{I\subseteq A}{\rAEa{11}}
  \proofstepwidestarbreak[11]{I\neq\varnothing}{\rAEa{\rAEb{11}}}
  \proofstepwidestarbreak[11]{\forall x\in I\,\forall a\in A\,(x\star a\in I)}{\rAEb{\rAEb{11}}}
  \proofstepwidestarbreak[15]{x\in I}{\rA}
  \proofstepwidestarbreak[16]{a\in A}{\rA}
  \proofstepwidestarbreak[11,15,16]{x\star a\in I}{\rRE{16,\rUE{\rRE{15,\rUE{14}}}}}
  \proofstepwidestarbreak[1,11]{\mathsf{RId}(I;A,\star)}{\FormulaRefAuto{SemigroupRightIdealDef}[def]{1,12,13,17}}
  \proofstepwidestarbreak[1]{\begin{aligned}[t]&(I\subseteq A\land I\neq\varnothing\land \forall x\in I\,\forall a\in A\,(x\star a\in I))\rightarrow{}\\[-2pt]&\quad \mathsf{RId}(I;A,\star)\end{aligned}}{\rRI{11,18}}
  \proofstepwidestarbreak[1]{\begin{aligned}[t]&\mathsf{RId}(I;A,\star)\leftrightarrow{}\\[-2pt]&\quad (I\subseteq A\land I\neq\varnothing\land \forall x\in I\,\forall a\in A\,(x\star a\in I))\end{aligned}}{\rLRI{10,19}}
\end{tabproofwide}

\par\noindent\begin{minipage}{\linewidth}
\FormulaDefDeltaK[Zweiseitiges Ideal]{%
\mathsf{Id}(I;A,\star)%
}{SemigroupIdealDef}{%
  \DeltaRow{Mengen}{A\dsep I}
  \DeltaRow{Funktionen}{\star}
  \DeltaRow{\textbf{Voraussetzung}}{}
  \DeltaPrem{Umgebende Halbgruppe}{\SemigroupAlg{A}{\star}}
  \DeltaRow{\textbf{Axiome}}{}
  \DeltaRow{Linksideal}{\mathsf{LId}(I;A,\star)}
  \DeltaRow{Rechtsideal}{\mathsf{RId}(I;A,\star)}
  \DeltaRow{\textbf{Neues Symbol}}{}
  \DeltaPrem{Zweiseitiges Ideal}{\mathsf{Id}(I;A,\star)}
}
\end{minipage}\par

\par\noindent\begin{minipage}{\linewidth}
\FormulaAxiomDeltaK[Zugriff auf das Linksideal eines Ideals]{%
\SemigroupAlg{A}{\star}\dsep \mathsf{Id}(I;A,\star)\vdash \mathsf{LId}(I;A,\star)%
}{SemigroupIdealLeftAxiom}{%
  \DeltaRow{Mengen}{A\dsep I}
  \DeltaRow{Funktionen}{\star}
}
\end{minipage}\par

\par\noindent\begin{minipage}{\linewidth}
\FormulaAxiomDeltaK[Zugriff auf das Rechtsideal eines Ideals]{%
\SemigroupAlg{A}{\star}\dsep \mathsf{Id}(I;A,\star)\vdash \mathsf{RId}(I;A,\star)%
}{SemigroupIdealRightAxiom}{%
  \DeltaRow{Mengen}{A\dsep I}
  \DeltaRow{Funktionen}{\star}
}
\end{minipage}\par

\par\noindent\begin{minipage}{\linewidth}
\FormulaThmDeltaK[Kriterium für ein zweiseitiges Ideal]{%
\SemigroupAlg{A}{\star}\vdash \mathsf{Id}(I;A,\star)\leftrightarrow(\mathsf{LId}(I;A,\star)\land \mathsf{RId}(I;A,\star))%
}{SemigroupIdealCriterion}{%
  \DeltaRow{Mengen}{A\dsep I}
  \DeltaRow{Funktionen}{\star}
}
\end{minipage}\par
\begin{tabproofwide}
  \proofstepwidestarbreak[1]{\SemigroupAlg{A}{\star}}{\rA}
  \proofstepwidestarbreak[2]{\mathsf{Id}(I;A,\star)}{\rA}
  \proofstepwidestarbreak[1,2]{\mathsf{LId}(I;A,\star)}{\FormulaRefAuto{SemigroupIdealLeftAxiom}[ax]{1,2}}
  \proofstepwidestarbreak[1,2]{\mathsf{RId}(I;A,\star)}{\FormulaRefAuto{SemigroupIdealRightAxiom}[ax]{1,2}}
  \proofstepwidestarbreak[1,2]{\mathsf{LId}(I;A,\star)\land \mathsf{RId}(I;A,\star)}{\rAI{3,4}}
  \proofstepwidestarbreak[1]{\mathsf{Id}(I;A,\star)\rightarrow(\mathsf{LId}(I;A,\star)\land \mathsf{RId}(I;A,\star))}{\rRI{2,5}}
  \proofstepwidestarbreak[7]{\mathsf{LId}(I;A,\star)\land \mathsf{RId}(I;A,\star)}{\rA}
  \proofstepwidestarbreak[1,7]{\mathsf{Id}(I;A,\star)}{\FormulaRefAuto{SemigroupIdealDef}[def]{1,\rAEa{7},\rAEb{7}}}
  \proofstepwidestarbreak[1]{(\mathsf{LId}(I;A,\star)\land \mathsf{RId}(I;A,\star))\rightarrow \mathsf{Id}(I;A,\star)}{\rRI{7,8}}
  \proofstepwidestarbreak[1]{\mathsf{Id}(I;A,\star)\leftrightarrow(\mathsf{LId}(I;A,\star)\land \mathsf{RId}(I;A,\star))}{\rLRI{6,9}}
\end{tabproofwide}

\FormulaThmDeltaK[Tr\"agermenge eines Ideals]{%
  \SemigroupAlg{A}{\star}\dsep\mathsf{Id}(I;A,\star)
  \vdash I\subseteq A%
}{SemigroupIdealSubset}{%
  \DeltaRow{Mengen}{A\dsep I}
  \DeltaRow{Funktionen}{\star}
}
\begin{tabproofwide}
  \proofstepwidestarbreak[1]{\SemigroupAlg{A}{\star}}{\rA}
  \proofstepwidestarbreak[2]{\mathsf{Id}(I;A,\star)}{\rA}
  \proofstepwidestarbreak[1,2]{\mathsf{LId}(I;A,\star)}{\FormulaRefAuto{SemigroupIdealLeftAxiom}[ax]{1,2}}
  \proofstepwidestarbreak[1,2]{I\subseteq A}{\FormulaRefAuto{SemigroupLeftIdealSubsetAxiom}[ax]{1,3}}
\end{tabproofwide}

\FormulaThmDeltaK[Nichtleerheit eines Ideals]{%
  \SemigroupAlg{A}{\star}\dsep\mathsf{Id}(I;A,\star)
  \vdash I\neq\varnothing%
}{SemigroupIdealNonempty}{%
  \DeltaRow{Mengen}{A\dsep I}
  \DeltaRow{Funktionen}{\star}
}
\begin{tabproofwide}
  \proofstepwidestarbreak[1]{\SemigroupAlg{A}{\star}}{\rA}
  \proofstepwidestarbreak[2]{\mathsf{Id}(I;A,\star)}{\rA}
  \proofstepwidestarbreak[1,2]{\mathsf{LId}(I;A,\star)}{\FormulaRefAuto{SemigroupIdealLeftAxiom}[ax]{1,2}}
  \proofstepwidestarbreak[1,2]{I\neq\varnothing}{\FormulaRefAuto{SemigroupLeftIdealNonemptyAxiom}[ax]{1,3}}
\end{tabproofwide}

\FormulaThmDeltaK[Linksabsorption eines Ideals]{%
  \SemigroupAlg{A}{\star}\dsep\mathsf{Id}(I;A,\star)\dsep
  a\in A\dsep x\in I\vdash a\star x\in I%
}{SemigroupIdealLeftAbsorption}{%
  \DeltaRow{Mengen}{A\dsep I\dsep a\dsep x}
  \DeltaRow{Funktionen}{\star}
}
\begin{tabproofwide}
  \proofstepwidestarbreak[1]{\SemigroupAlg{A}{\star}}{\rA}
  \proofstepwidestarbreak[2]{\mathsf{Id}(I;A,\star)}{\rA}
  \proofstepwidestarbreak[3]{a\in A}{\rA}
  \proofstepwidestarbreak[4]{x\in I}{\rA}
  \proofstepwidestarbreak[1,2]{\mathsf{LId}(I;A,\star)}{\FormulaRefAuto{SemigroupIdealLeftAxiom}[ax]{1,2}}
  \proofstepwidestarbreak[1,2,3,4]{a\star x\in I}{\FormulaRefAuto{SemigroupLeftIdealClosureAxiom}[ax]{1,5,3,4}}
\end{tabproofwide}

\FormulaThmDeltaK[Rechtsabsorption eines Ideals]{%
  \SemigroupAlg{A}{\star}\dsep\mathsf{Id}(I;A,\star)\dsep
  x\in I\dsep a\in A\vdash x\star a\in I%
}{SemigroupIdealRightAbsorption}{%
  \DeltaRow{Mengen}{A\dsep I\dsep a\dsep x}
  \DeltaRow{Funktionen}{\star}
}
\begin{tabproofwide}
  \proofstepwidestarbreak[1]{\SemigroupAlg{A}{\star}}{\rA}
  \proofstepwidestarbreak[2]{\mathsf{Id}(I;A,\star)}{\rA}
  \proofstepwidestarbreak[3]{x\in I}{\rA}
  \proofstepwidestarbreak[4]{a\in A}{\rA}
  \proofstepwidestarbreak[1,2]{\mathsf{RId}(I;A,\star)}{\FormulaRefAuto{SemigroupIdealRightAxiom}[ax]{1,2}}
  \proofstepwidestarbreak[1,2,3,4]{x\star a\in I}{\FormulaRefAuto{SemigroupRightIdealClosureAxiom}[ax]{1,5,3,4}}
\end{tabproofwide}

\FormulaThmDeltaK[Der nichtleere Tr\"ager ist ein Ideal]{%
  \SemigroupAlg{A}{\star}\dsep A\neq\varnothing
  \vdash\mathsf{Id}(A;A,\star)%
}{SemigroupCarrierIsIdeal}{%
  \DeltaRow{Mengen}{A\dsep a\dsep x}
  \DeltaRow{Funktionen}{\star}
}
\begin{tabproofwide}
  \proofstepwidestar[1]{\SemigroupAlg{A}{\star}}{\rA}
  \proofstepwidestar[2]{A\neq\varnothing}{\rA}
  \proofstepwidestar[]{A\subseteq A}{\FormulaRefAuto{A\subseteq A}[thm]}
  \proofstepwidestar[4]{a\in A}{\rA}
  \proofstepwidestar[5]{x\in A}{\rA}
  \proofstepwidestar[1,4,5]{a\star x\in A}{%
    \FormulaRefAuto{SemigroupClosureAxiom}[ax]{1,4,5}}
  \proofstepwidestar[1]{\forall a\in A\,\forall x\in A\,(a\star x\in A)}{%
    \rUI{\rRI{4,\rUI{\rRI{5,6}}}}}
  \proofstepwidestar[1,2]{\mathsf{LId}(A;A,\star)}{%
    \FormulaRefAuto{SemigroupLeftIdealCriterion}[thm]{1,\rAI{3,\rAI{2,7}}}}
  \proofstepwidestar[1]{\forall x\in A\,\forall a\in A\,(x\star a\in A)}{%
    \rUI{\rRI{5,\rUI{\rRI{4,6}}}}}
  \proofstepwidestar[1,2]{\mathsf{RId}(A;A,\star)}{%
    \FormulaRefAuto{SemigroupRightIdealCriterion}[thm]{1,\rAI{3,\rAI{2,9}}}}
  \proofstepwidestar[1,2]{\mathsf{Id}(A;A,\star)}{%
    \FormulaRefAuto{SemigroupIdealCriterion}[thm]{1,\rAI{8,10}}}
\end{tabproofwide}

\FormulaThmDeltaK[Schnitt zweier Ideale]{%
  \begin{aligned}[t]
    &\SemigroupAlg{A}{\star}\dsep\mathsf{Id}(I;A,\star)\dsep
      \mathsf{Id}(J;A,\star)\vdash{}\\[-2pt]
    &\mathsf{Id}(I\cap J;A,\star)
  \end{aligned}
}{SemigroupIdealIntersection}{%
  \DeltaRow{Mengen}{A\dsep I\dsep J\dsep a\dsep x\dsep i\dsep j}
  \DeltaRow{Funktionen}{\star}
}
Der Schnitt ist nichtleer, weil das Produkt eines Elements von (I) mit
einem Element von (J) zugleich in beiden Idealen liegt.  Die linke und
die rechte Absorption werden danach komponentenweise vererbt.

\begin{tabproofsplitwide}
  \proofpartwideindR[Nichtleerheit]{%
    \SemigroupAlg{A}{\star}\dsep\mathsf{Id}(I;A,\star)\dsep
    \mathsf{Id}(J;A,\star)\vdash I\cap J\neq\varnothing
  }{%
    \SemigroupAlg{A}{\star}\dsep\mathsf{Id}(I;A,\star)\dsep
    \mathsf{Id}(J;A,\star)\vdash I\cap J\neq\varnothing
  }[SemigroupIdealIntersectionNonempty]
    \proofstepwidestar[1]{\SemigroupAlg{A}{\star}}{\rA}
    \proofstepwidestar[2]{\mathsf{Id}(I;A,\star)}{\rA}
    \proofstepwidestar[3]{\mathsf{Id}(J;A,\star)}{\rA}
    \proofstepwidestar[1,2]{I\subseteq A}{%
      \FormulaRefAuto{SemigroupIdealSubset}[thm]{1,2}}
    \proofstepwidestar[1,3]{J\subseteq A}{%
      \FormulaRefAuto{SemigroupIdealSubset}[thm]{1,3}}
    \proofstepwidestar[1,2]{I\neq\varnothing}{%
      \FormulaRefAuto{SemigroupIdealNonempty}[thm]{1,2}}
    \proofstepwidestar[1,3]{J\neq\varnothing}{%
      \FormulaRefAuto{SemigroupIdealNonempty}[thm]{1,3}}
    \proofstepwidestar[1,2]{\exists i\,(i\in I)}{%
      \FormulaRefAuto{A\neq\varnothing\eqvdash\exists x\,(x\in A)}[thm]{6}}
    \proofstepwidestar[1,3]{\exists j\,(j\in J)}{%
      \FormulaRefAuto{A\neq\varnothing\eqvdash\exists x\,(x\in A)}[thm]{7}}
    \proofstepwidestar[10]{i\in I}{\rA}
    \proofstepwidestar[11]{j\in J}{\rA}
    \proofstepwidestar[1,2,10]{i\in A}{%
      \FormulaRefAuto{A\subseteq B,\,x\in A\vdash x\in B}[thm]{4,10}}
    \proofstepwidestar[1,3,11]{j\in A}{%
      \FormulaRefAuto{A\subseteq B,\,x\in A\vdash x\in B}[thm]{5,11}}
    \proofstepwidestar[1,3,10,11]{i\star j\in J}{%
      \FormulaRefAuto{SemigroupIdealLeftAbsorption}[thm]{1,3,12,11}}
    \proofstepwidestar[1,2,10,11]{i\star j\in I}{%
      \FormulaRefAuto{SemigroupIdealRightAbsorption}[thm]{1,2,10,13}}
    \proofstepwidestar[1,2,3,10,11]{i\star j\in I\cap J}{%
      \FormulaRefAuto{x\in A\cap B\eqvdash x\in A\land x\in B}[thm]{%
        \rAI{15,14}}}
    \proofstepwidestar[1,2,3,10,11]{I\cap J\neq\varnothing}{%
      \FormulaRefAuto{a\in A\vdash A\neq\varnothing}[thm]{16}}
    \proofstepwidestar[1,2,3,10]{I\cap J\neq\varnothing}{%
      \rEE{9,11,17}}
    \proofstepwidestar[1,2,3]{I\cap J\neq\varnothing}{%
      \rEE{8,10,18}}
  \closeproofpartwideindR

  \proofpartwideindR[Linksidealeigenschaft]{%
    \SemigroupAlg{A}{\star}\dsep\mathsf{Id}(I;A,\star)\dsep
    \mathsf{Id}(J;A,\star)\vdash\mathsf{LId}(I\cap J;A,\star)
  }{%
    \SemigroupAlg{A}{\star}\dsep\mathsf{Id}(I;A,\star)\dsep
    \mathsf{Id}(J;A,\star)\vdash\mathsf{LId}(I\cap J;A,\star)
  }[SemigroupIdealIntersectionLeftIdeal]
    \proofstepwidestar[1]{\SemigroupAlg{A}{\star}}{\rA}
    \proofstepwidestar[2]{\mathsf{Id}(I;A,\star)}{\rA}
    \proofstepwidestar[3]{\mathsf{Id}(J;A,\star)}{\rA}
    \proofstepwidestar[1,2]{I\subseteq A}{%
      \FormulaRefAuto{SemigroupIdealSubset}[thm]{1,2}}
    \proofstepwidestar[]{I\cap J\subseteq I}{%
      \FormulaRefAuto{A\cap B\subseteq A}[thm]}
    \proofstepwidestar[1,2]{I\cap J\subseteq A}{%
      \FormulaRefAuto{A\subseteq B\dsep B\subseteq C\vdash
        A\subseteq C}[thm]{5,4}}
    \proofstepwidestar[1,2,3]{I\cap J\neq\varnothing}{%
      \FormulaRefAuto{SemigroupIdealIntersectionNonempty}[thm]{1,2,3}}
    \proofstepwidestar[8]{a\in A}{\rA}
    \proofstepwidestar[9]{x\in I\cap J}{\rA}
    \proofstepwidestar[9]{x\in I}{%
      \FormulaRefAuto{x\in A\cap B\vdash x\in A}[thm]{9}}
    \proofstepwidestar[9]{x\in J}{%
      \FormulaRefAuto{x\in A\cap B\vdash x\in B}[thm]{9}}
    \proofstepwidestar[1,2,8,9]{a\star x\in I}{%
      \FormulaRefAuto{SemigroupIdealLeftAbsorption}[thm]{1,2,8,10}}
    \proofstepwidestar[1,3,8,9]{a\star x\in J}{%
      \FormulaRefAuto{SemigroupIdealLeftAbsorption}[thm]{1,3,8,11}}
    \proofstepwidestar[1,2,3,8,9]{a\star x\in I\cap J}{%
      \FormulaRefAuto{x\in A\cap B\eqvdash x\in A\land x\in B}[thm]{%
        \rAI{12,13}}}
    \proofstepwidestar[1,2,3]{%
      \forall a\in A\,\forall x\in I\cap J\,(a\star x\in I\cap J)}{%
      \rUI{\rRI{8,\rUI{\rRI{9,14}}}}}
    \proofstepwidestar[1,2,3]{\mathsf{LId}(I\cap J;A,\star)}{%
      \FormulaRefAuto{SemigroupLeftIdealCriterion}[thm]{1,\rAI{6,\rAI{7,15}}}}
  \closeproofpartwideindR

  \proofpartwideindR[Rechtsidealeigenschaft]{%
    \SemigroupAlg{A}{\star}\dsep\mathsf{Id}(I;A,\star)\dsep
    \mathsf{Id}(J;A,\star)\vdash\mathsf{RId}(I\cap J;A,\star)
  }{%
    \SemigroupAlg{A}{\star}\dsep\mathsf{Id}(I;A,\star)\dsep
    \mathsf{Id}(J;A,\star)\vdash\mathsf{RId}(I\cap J;A,\star)
  }[SemigroupIdealIntersectionRightIdeal]
    \proofstepwidestar[1]{\SemigroupAlg{A}{\star}}{\rA}
    \proofstepwidestar[2]{\mathsf{Id}(I;A,\star)}{\rA}
    \proofstepwidestar[3]{\mathsf{Id}(J;A,\star)}{\rA}
    \proofstepwidestar[1,2]{I\subseteq A}{%
      \FormulaRefAuto{SemigroupIdealSubset}[thm]{1,2}}
    \proofstepwidestar[]{I\cap J\subseteq I}{%
      \FormulaRefAuto{A\cap B\subseteq A}[thm]}
    \proofstepwidestar[1,2]{I\cap J\subseteq A}{%
      \FormulaRefAuto{A\subseteq B\dsep B\subseteq C\vdash
        A\subseteq C}[thm]{5,4}}
    \proofstepwidestar[1,2,3]{I\cap J\neq\varnothing}{%
      \FormulaRefAuto{SemigroupIdealIntersectionNonempty}[thm]{1,2,3}}
    \proofstepwidestar[8]{x\in I\cap J}{\rA}
    \proofstepwidestar[9]{a\in A}{\rA}
    \proofstepwidestar[8]{x\in I}{%
      \FormulaRefAuto{x\in A\cap B\vdash x\in A}[thm]{8}}
    \proofstepwidestar[8]{x\in J}{%
      \FormulaRefAuto{x\in A\cap B\vdash x\in B}[thm]{8}}
    \proofstepwidestar[1,2,8,9]{x\star a\in I}{%
      \FormulaRefAuto{SemigroupIdealRightAbsorption}[thm]{1,2,10,9}}
    \proofstepwidestar[1,3,8,9]{x\star a\in J}{%
      \FormulaRefAuto{SemigroupIdealRightAbsorption}[thm]{1,3,11,9}}
    \proofstepwidestar[1,2,3,8,9]{x\star a\in I\cap J}{%
      \FormulaRefAuto{x\in A\cap B\eqvdash x\in A\land x\in B}[thm]{%
        \rAI{12,13}}}
    \proofstepwidestar[1,2,3]{%
      \forall x\in I\cap J\,\forall a\in A\,(x\star a\in I\cap J)}{%
      \rUI{\rRI{8,\rUI{\rRI{9,14}}}}}
    \proofstepwidestar[1,2,3]{\mathsf{RId}(I\cap J;A,\star)}{%
      \FormulaRefAuto{SemigroupRightIdealCriterion}[thm]{1,\rAI{6,\rAI{7,15}}}}
  \closeproofpartwideindR

  \proofpartwide{Schluss}
    \proofstepwidestar[1]{\SemigroupAlg{A}{\star}}{\rA}
    \proofstepwidestar[2]{\mathsf{Id}(I;A,\star)}{\rA}
    \proofstepwidestar[3]{\mathsf{Id}(J;A,\star)}{\rA}
    \proofstepwidestar[1,2,3]{\mathsf{LId}(I\cap J;A,\star)}{%
      \FormulaRefAuto{SemigroupIdealIntersectionLeftIdeal}[thm]{1,2,3}}
    \proofstepwidestar[1,2,3]{\mathsf{RId}(I\cap J;A,\star)}{%
      \FormulaRefAuto{SemigroupIdealIntersectionRightIdeal}[thm]{1,2,3}}
    \proofstepwidestar[1,2,3]{\mathsf{Id}(I\cap J;A,\star)}{%
      \FormulaRefAuto{SemigroupIdealCriterion}[thm]{1,\rAI{4,5}}}
  \closeproofpartwide
\end{tabproofsplitwide}

\subsection{Hauptideale}

Statt die symbolische Einheitsadjunktion \(A^1\) informell zu benutzen,
definieren wir Hauptideale unmittelbar durch ihre Elemente. Damit sind auch
die F\"alle ohne linken oder rechten Faktor ausdr\"ucklich erfasst.

\FormulaDefDeltaK[Linkshauptideal]{%
  \SemigroupAlg{A}{\star}\dsep a\in A\vdash
  L_{A,\star}(a)\coloneqq
  \bigl\{b\in A\bigm|b=a\lor
    \exists x\in A\,(b=x\star a)\bigr\}%
}{SemigroupPrincipalLeftIdealDef}{%
  \DeltaRow{Mengen}{A\dsep a\dsep b\dsep x}
  \DeltaRow{Funktionen}{\star}
  \DeltaRow{\textbf{Neues Symbol}}{}
  \DeltaPrem{Linkshauptideal}{L_{A,\star}(a)}
}

\FormulaDefDeltaK[Rechtshauptideal]{%
  \SemigroupAlg{A}{\star}\dsep a\in A\vdash
  R_{A,\star}(a)\coloneqq
  \bigl\{b\in A\bigm|b=a\lor
    \exists y\in A\,(b=a\star y)\bigr\}%
}{SemigroupPrincipalRightIdealDef}{%
  \DeltaRow{Mengen}{A\dsep a\dsep b\dsep y}
  \DeltaRow{Funktionen}{\star}
  \DeltaRow{\textbf{Neues Symbol}}{}
  \DeltaPrem{Rechtshauptideal}{R_{A,\star}(a)}
}

\FormulaDefDeltaK[Zweiseitiges Hauptideal]{%
  \begin{aligned}[t]
    &\SemigroupAlg{A}{\star}\dsep a\in A\vdash{}\\[-2pt]
    &J_{A,\star}(a)\coloneqq
      \Bigl\{b\in A\Bigm|
        b=a\lor\exists x\in A\,(b=x\star a)\\[-2pt]
    &\hspace{93pt}{}\lor\exists y\in A\,(b=a\star y)
        \lor\exists x,y\in A\,(b=(x\star a)\star y)
      \Bigr\}
  \end{aligned}
}{SemigroupPrincipalTwoSidedIdealDef}{%
  \DeltaRow{Mengen}{A\dsep a\dsep b\dsep x\dsep y}
  \DeltaRow{Funktionen}{\star}
  \DeltaRow{\textbf{Neues Symbol}}{}
  \DeltaPrem{Zweiseitiges Hauptideal}{J_{A,\star}(a)}
}

\FormulaThmDeltaK[Elementkriterium des Linkshauptideals]{%
  \begin{aligned}[t]
    &\SemigroupAlg{A}{\star}\dsep a\in A\vdash{}\\[-2pt]
    &b\in L_{A,\star}(a)\leftrightarrow
      \bigl(b\in A\land(b=a\lor\exists x\in A\,(b=x\star a))\bigr)
  \end{aligned}
}{SemigroupPrincipalLeftIdealMembership}{%
  \DeltaRow{Mengen}{A\dsep a\dsep b\dsep x}
  \DeltaRow{Funktionen}{\star}
}
\begin{tabproofwide}
  \proofstepwidestar[1]{\SemigroupAlg{A}{\star}}{\rA}
  \proofstepwidestar[2]{a\in A}{\rA}
  \proofstepwidestar[1,2]{%
    L_{A,\star}(a)=\{u\in A\mid
      u=a\lor\exists x\in A\,(u=x\star a)\}}{%
    \FormulaRefAuto{SemigroupPrincipalLeftIdealDef}[def]{1,2}}
  \proofstepwidestar[1,2]{b\in\{u\in A\mid u=a\lor\exists x\in A\;(u=x\star a)\}\leftrightarrow(b\in A\land(b=a\lor\exists x\in A\;(b=x\star a)))}{%
      \FormulaRefAuto{x\in\{u\in A\mid P(u)\}\eqvdash
      x\in A\land P(x)}[thm]}
  \proofstepwidestar[1,2]{%
    b\in L_{A,\star}(a)\leftrightarrow
    \bigl(b\in A\land(b=a\lor\exists x\in A\,(b=x\star a))\bigr)}{%
    \rIE{3,4}}
\end{tabproofwide}

\FormulaThmDeltaK[Elementkriterium des Rechtshauptideals]{%
  \begin{aligned}[t]
    &\SemigroupAlg{A}{\star}\dsep a\in A\vdash{}\\[-2pt]
    &b\in R_{A,\star}(a)\leftrightarrow
      \bigl(b\in A\land(b=a\lor\exists y\in A\,(b=a\star y))\bigr)
  \end{aligned}
}{SemigroupPrincipalRightIdealMembership}{%
  \DeltaRow{Mengen}{A\dsep a\dsep b\dsep y}
  \DeltaRow{Funktionen}{\star}
}
\begin{tabproofwide}
  \proofstepwidestar[1]{\SemigroupAlg{A}{\star}}{\rA}
  \proofstepwidestar[2]{a\in A}{\rA}
  \proofstepwidestar[1,2]{%
    R_{A,\star}(a)=\{u\in A\mid
      u=a\lor\exists y\in A\,(u=a\star y)\}}{%
    \FormulaRefAuto{SemigroupPrincipalRightIdealDef}[def]{1,2}}
  \proofstepwidestar[1,2]{b\in\{u\in A\mid u=a\lor\exists y\in A\;(u=a\star y)\}\leftrightarrow(b\in A\land(b=a\lor\exists y\in A\;(b=a\star y)))}{%
      \FormulaRefAuto{x\in\{u\in A\mid P(u)\}\eqvdash
      x\in A\land P(x)}[thm]}
  \proofstepwidestar[1,2]{%
    b\in R_{A,\star}(a)\leftrightarrow
    \bigl(b\in A\land(b=a\lor\exists y\in A\,(b=a\star y))\bigr)}{%
    \rIE{3,4}}
\end{tabproofwide}

\FormulaThmDeltaK[Elementkriterium des zweiseitigen Hauptideals]{%
  \begin{aligned}[t]
    &\SemigroupAlg{A}{\star}\dsep a\in A\vdash{}\\[-2pt]
    &b\in J_{A,\star}(a)\leftrightarrow
      \Bigl(b\in A\land\bigl(
        b=a\lor\exists x\in A\,(b=x\star a)\\[-2pt]
    &\hspace{112pt}{}\lor\exists y\in A\,(b=a\star y)
        \lor\exists x,y\in A\,(b=(x\star a)\star y)
      \bigr)\Bigr)
  \end{aligned}
}{SemigroupPrincipalTwoSidedIdealMembership}{%
  \DeltaRow{Mengen}{A\dsep a\dsep b\dsep x\dsep y}
  \DeltaRow{Funktionen}{\star}
}
\begin{tabproofwide}
  \proofstepwidestar[1]{\SemigroupAlg{A}{\star}}{\rA}
  \proofstepwidestar[2]{a\in A}{\rA}
  \proofstepwidestarbreak[1,2]{\begin{aligned}[t]
 &J_{A,\star}(a)=\Bigl\{u\in A\Bigm|u=a\\[-2pt]
 &\quad\lor\exists x\in A\,(u=x\star a)\\[-2pt]
 &\quad\lor\exists y\in A\,(u=a\star y)\\[-2pt]
 &\quad\lor\exists x,y\in A\,(u=(x\star a)\star y)\Bigr\}
 \end{aligned}}{%
    \FormulaRefAuto{SemigroupPrincipalTwoSidedIdealDef}[def]{1,2}}
  \proofstepwidestar[1,2]{\begin{aligned}[t]
 &b\in\{u\in A\mid u=a\\[-2pt]
 &\quad\lor\exists x\in A\;(u=x\star a)\\[-2pt]
 &\quad\lor\exists y\in A\;(u=a\star y)\\[-2pt]
 &\quad\lor\exists x,y\in A\;(u=(x\star a)\star y)\}\\[-2pt]
 &\leftrightarrow\bigl(b\in A\land(b=a\\[-2pt]
 &\quad\lor\exists x\in A\;(b=x\star a)\\[-2pt]
 &\quad\lor\exists y\in A\;(b=a\star y)\\[-2pt]
 &\quad\lor\exists x,y\in A\;(b=(x\star a)\star y))\bigr)
 \end{aligned}}{%
      \FormulaRefAuto{x\in\{u\in A\mid P(u)\}\eqvdash
      x\in A\land P(x)}[thm]}
  \proofstepwidestarbreak[1,2]{\begin{aligned}[t]
 &b\in J_{A,\star}(a)\leftrightarrow\Bigl(b\in A\land\bigl(b=a\\[-2pt]
 &\quad\lor\exists x\in A\,(b=x\star a)\\[-2pt]
 &\quad\lor\exists y\in A\,(b=a\star y)\\[-2pt]
 &\quad\lor\exists x,y\in A\,(b=(x\star a)\star y)\bigr)\Bigr)
 \end{aligned}}{%
    \rIE{3,4}}
\end{tabproofwide}

\FormulaThmDeltaK[Der Erzeuger liegt in seinen Hauptidealen]{%
  \begin{aligned}[t]
    &\SemigroupAlg{A}{\star}\dsep a\in A\vdash{}\\[-2pt]
    &a\in L_{A,\star}(a)\land
      a\in R_{A,\star}(a)\land
      a\in J_{A,\star}(a)
  \end{aligned}
}{SemigroupPrincipalIdealsContainGenerator}{%
  \DeltaRow{Mengen}{A\dsep a}
  \DeltaRow{Funktionen}{\star}
}
\begin{tabproofwide}
  \proofstepwidestar[1]{\SemigroupAlg{A}{\star}}{\rA}
  \proofstepwidestar[2]{a\in A}{\rA}
  \proofstepwidestar[1,2]{a\in L_{A,\star}(a)}{%
    \FormulaRefAuto{SemigroupPrincipalLeftIdealMembership}[thm]{%
      1,2,\rAI{2,\rOIa{\rII}}}}
  \proofstepwidestar[1,2]{a\in R_{A,\star}(a)}{%
    \FormulaRefAuto{SemigroupPrincipalRightIdealMembership}[thm]{%
      1,2,\rAI{2,\rOIa{\rII}}}}
  \proofstepwidestar[1,2]{a\in J_{A,\star}(a)}{%
    \FormulaRefAuto{SemigroupPrincipalTwoSidedIdealMembership}[thm]{%
      1,2,\rAI{2,\rOIa{\rII}}}}
  \proofstepwidestar[1,2]{%
    a\in L_{A,\star}(a)\land a\in R_{A,\star}(a)\land
    a\in J_{A,\star}(a)}{\rAI{3,\rAI{4,5}}}
\end{tabproofwide}

\par\noindent\begin{minipage}{\linewidth}
\FormulaThmDeltaK[Das Linkshauptideal ist ein Linksideal]{%
  \SemigroupAlg{A}{\star}\dsep a\in A\vdash
  \mathsf{LId}(L_{A,\star}(a);A,\star)%
}{SemigroupPrincipalLeftIdealIsLeftIdeal}{%
  \DeltaRow{Mengen}{A\dsep a\dsep b\dsep s\dsep x}
  \DeltaRow{Funktionen}{\star}
}
\end{minipage}\par
\begin{tabproofsplitwide}
\proofpartwide{Teilträger und Nichtleerheit}
  \proofstepwidestarbreak[1]{\SemigroupAlg{A}{\star}}{\rA}
  \proofstepwidestarbreak[2]{a\in A}{\rA}
  \proofstepwidestarbreak[]{L_{A,\star}(a)\subseteq A}{%
    \FormulaRefAuto{\{x\in A\mid P(x)\}\subseteq A}[thm]}
  \proofstepwidestar[1,2]{a\in L_{A,\star}(a)\land a\in R_{A,\star}(a)\land a\in J_{A,\star}(a)}{%
      \FormulaRefAuto{SemigroupPrincipalIdealsContainGenerator}[thm]{1,2}}
  \proofstepwidestarbreak[1,2]{a\in L_{A,\star}(a)}{%
    \rAEa{4}}
  \proofstepwidestarbreak[1,2]{L_{A,\star}(a)\neq\varnothing}{%
    \FormulaRefAuto{a\in A\vdash A\neq\varnothing}[thm]{5}}

\closeproofpartwide
\proofpartwide{Linksabsorption}
\setcounter{proofstepnr}{6}
  \proofstepwidestarbreak[7]{s\in A}{\rA}
  \proofstepwidestarbreak[8]{b\in L_{A,\star}(a)}{\rA}
  \proofstepwidestar[1,2,8]{b\in A\land\bigl(b=a\lor \exists u\in A\;(b=u\star a)\bigr)}{%
      \FormulaRefAuto{SemigroupPrincipalLeftIdealMembership}[thm]{1,2,8}}
  \proofstepwidestarbreak[1,2,8]{b\in A}{%
    \rAEa{9}}
  \proofstepwidestarbreak[1,2,7,8]{s\star b\in A}{%
    \FormulaRefAuto{SemigroupClosureAxiom}[ax]{1,7,10}}
  \proofstepwidestar[1,2,8]{b\in A\land\bigl(b=a\lor \exists u\in A\;(b=u\star a)\bigr)}{%
      \FormulaRefAuto{SemigroupPrincipalLeftIdealMembership}[thm]{1,2,8}}
  \proofstepwidestarbreak[1,2,8]{b=a\lor\exists x\in A\,(b=x\star a)}{%
    \rAEb{12}}
  \proofstepwidestarbreak[14]{b=a}{\rA}
  \proofstepwidestarbreak[7,14]{s\star b=s\star a}{\rIE{14}}
  \proofstepwidestarbreak[7,14]{\exists u\in A\,(s\star b=u\star a)}{%
    \rEI{7,15}}
  \proofstepwidestarbreak[1,2,7,8,14]{s\star b\in L_{A,\star}(a)}{%
    \FormulaRefAuto{SemigroupPrincipalLeftIdealMembership}[thm]{%
      1,2,\rAI{11,\rOIb{16}}}}
  \proofstepwidestarbreak[18]{\exists x\in A\,(b=x\star a)}{\rA}
  \proofstepwidestarbreak[19]{x\in A\land b=x\star a}{\rA}
  \proofstepwidestarbreak[19]{x\in A}{\rAEa{19}}
  \proofstepwidestarbreak[19]{b=x\star a}{\rAEb{19}}
  \proofstepwidestarbreak[1,7,19]{s\star x\in A}{%
    \FormulaRefAuto{SemigroupClosureAxiom}[ax]{1,7,20}}
  \proofstepwidestar[1,2,7,19]{(s\star x)\star a=s\star(x\star a)}{%
      \FormulaRefAuto{SemigroupAssociativityAxiom}[ax]{1,7,20,2}}
  \proofstepwidestarbreak[1,2,7,19]{s\star b=(s\star x)\star a}{%
    \rIE{21,23}}
  \proofstepwidestarbreak[1,2,7,19]{%
    \exists u\in A\,(s\star b=u\star a)}{\rEI{22,24}}
  \proofstepwidestarbreak[1,2,7,8,19]{s\star b\in L_{A,\star}(a)}{%
    \FormulaRefAuto{SemigroupPrincipalLeftIdealMembership}[thm]{%
      1,2,\rAI{11,\rOIb{25}}}}
  \proofstepwidestarbreak[1,2,7,8,18]{s\star b\in L_{A,\star}(a)}{%
    \rEE{18,19,26}}
  \proofstepwidestarbreak[1,2,7,8]{s\star b\in L_{A,\star}(a)}{%
    \rOE{13,14,17,18,27}}

\closeproofpartwide
\proofpartwide{Allgemeingültigkeit und Idealabschluss}
\setcounter{proofstepnr}{28}
  \proofstepwidestarbreak[1,2]{%
    \forall s\in A\,\forall b\in L_{A,\star}(a)\,
      (s\star b\in L_{A,\star}(a))}{%
    \rUI{\rRI{7,\rUI{\rRI{8,28}}}}}
  \proofstepwidestarbreak[1,2]{\mathsf{LId}(L_{A,\star}(a);A,\star)}{%
    \FormulaRefAuto{SemigroupLeftIdealCriterion}[thm]{1,\rAI{3,\rAI{6,29}}}}
\closeproofpartwide
\end{tabproofsplitwide}

\FormulaThmDeltaK[Das Rechtshauptideal ist ein Rechtsideal]{%
  \SemigroupAlg{A}{\star}\dsep a\in A\vdash
  \mathsf{RId}(R_{A,\star}(a);A,\star)%
}{SemigroupPrincipalRightIdealIsRightIdeal}{%
  \DeltaRow{Mengen}{A\dsep a\dsep b\dsep s\dsep y}
  \DeltaRow{Funktionen}{\star}
}
\begin{tabproofsplitwide}
\proofpartwide{Teilträger und Nichtleerheit}
  \proofstepwidestarbreak[1]{\SemigroupAlg{A}{\star}}{\rA}
  \proofstepwidestarbreak[2]{a\in A}{\rA}
  \proofstepwidestarbreak[]{R_{A,\star}(a)\subseteq A}{%
    \FormulaRefAuto{\{x\in A\mid P(x)\}\subseteq A}[thm]}
  \proofstepwidestar[1,2]{a\in L_{A,\star}(a)\land a\in R_{A,\star}(a)\land a\in J_{A,\star}(a)}{%
      \FormulaRefAuto{SemigroupPrincipalIdealsContainGenerator}[thm]{1,2}}
  \proofstepwidestarbreak[1,2]{a\in R_{A,\star}(a)}{%
    \rAEa{\rAEb{4}}}
  \proofstepwidestarbreak[1,2]{R_{A,\star}(a)\neq\varnothing}{%
    \FormulaRefAuto{a\in A\vdash A\neq\varnothing}[thm]{5}}

\closeproofpartwide
\proofpartwide{Rechtsabsorption}
\setcounter{proofstepnr}{6}
  \proofstepwidestarbreak[7]{b\in R_{A,\star}(a)}{\rA}
  \proofstepwidestarbreak[8]{s\in A}{\rA}
  \proofstepwidestar[1,2,7]{b\in A\land\bigl(b=a\lor \exists v\in A\;(b=a\star v)\bigr)}{%
      \FormulaRefAuto{SemigroupPrincipalRightIdealMembership}[thm]{1,2,7}}
  \proofstepwidestarbreak[1,2,7]{b\in A}{%
    \rAEa{9}}
  \proofstepwidestarbreak[1,2,7,8]{b\star s\in A}{%
    \FormulaRefAuto{SemigroupClosureAxiom}[ax]{1,10,8}}
  \proofstepwidestar[1,2,7]{b\in A\land\bigl(b=a\lor \exists v\in A\;(b=a\star v)\bigr)}{%
      \FormulaRefAuto{SemigroupPrincipalRightIdealMembership}[thm]{1,2,7}}
  \proofstepwidestarbreak[1,2,7]{b=a\lor\exists y\in A\,(b=a\star y)}{%
    \rAEb{12}}
  \proofstepwidestarbreak[14]{b=a}{\rA}
  \proofstepwidestarbreak[8,14]{b\star s=a\star s}{\rIE{14}}
  \proofstepwidestarbreak[8,14]{\exists u\in A\,(b\star s=a\star u)}{\rEI{8,15}}
  \proofstepwidestarbreak[1,2,7,8,14]{b\star s\in R_{A,\star}(a)}{%
    \FormulaRefAuto{SemigroupPrincipalRightIdealMembership}[thm]{%
      1,2,\rAI{11,\rOIb{16}}}}
  \proofstepwidestarbreak[18]{\exists y\in A\,(b=a\star y)}{\rA}
  \proofstepwidestarbreak[19]{y\in A\land b=a\star y}{\rA}
  \proofstepwidestarbreak[19]{y\in A}{\rAEa{19}}
  \proofstepwidestarbreak[19]{b=a\star y}{\rAEb{19}}
  \proofstepwidestarbreak[1,8,19]{y\star s\in A}{%
    \FormulaRefAuto{SemigroupClosureAxiom}[ax]{1,20,8}}
  \proofstepwidestar[1,2,8,19]{(a\star y)\star s=a\star(y\star s)}{%
      \FormulaRefAuto{SemigroupAssociativityAxiom}[ax]{1,2,20,8}}
  \proofstepwidestarbreak[1,2,8,19]{b\star s=a\star(y\star s)}{%
    \rIE{21,23}}
  \proofstepwidestarbreak[1,2,8,19]{\exists u\in A\,(b\star s=a\star u)}{\rEI{22,24}}
  \proofstepwidestarbreak[1,2,7,8,19]{b\star s\in R_{A,\star}(a)}{%
    \FormulaRefAuto{SemigroupPrincipalRightIdealMembership}[thm]{%
      1,2,\rAI{11,\rOIb{25}}}}
  \proofstepwidestarbreak[1,2,7,8,18]{b\star s\in R_{A,\star}(a)}{\rEE{18,19,26}}
  \proofstepwidestarbreak[1,2,7,8]{b\star s\in R_{A,\star}(a)}{%
    \rOE{13,14,17,18,27}}

\closeproofpartwide
\proofpartwide{Allgemeingültigkeit und Idealabschluss}
\setcounter{proofstepnr}{28}
  \proofstepwidestarbreak[1,2]{%
    \forall b\in R_{A,\star}(a)\,\forall s\in A\,
      (b\star s\in R_{A,\star}(a))}{%
    \rUI{\rRI{7,\rUI{\rRI{8,28}}}}}
  \proofstepwidestarbreak[1,2]{\mathsf{RId}(R_{A,\star}(a);A,\star)}{%
    \FormulaRefAuto{SemigroupRightIdealCriterion}[thm]{1,\rAI{3,\rAI{6,29}}}}
\closeproofpartwide
\end{tabproofsplitwide}

\FormulaThmDeltaK[Minimalität des Linkshauptideals]{%
  \begin{aligned}[t]
    &\SemigroupAlg{A}{\star}\dsep a\in A\dsep
      \mathsf{LId}(I;A,\star)\dsep a\in I\vdash{}\\[-2pt]
    &L_{A,\star}(a)\subseteq I
  \end{aligned}
}{SemigroupPrincipalLeftIdealMinimal}{%
  \DeltaRow{Mengen}{A\dsep I\dsep a\dsep b\dsep x}
  \DeltaRow{Funktionen}{\star}
}
\begin{tabproofwide}
  \proofstepwidestar[1]{\SemigroupAlg{A}{\star}}{\rA}
  \proofstepwidestar[2]{a\in A}{\rA}
  \proofstepwidestar[3]{\mathsf{LId}(I;A,\star)}{\rA}
  \proofstepwidestar[4]{a\in I}{\rA}
  \proofstepwidestar[5]{b\in L_{A,\star}(a)}{\rA}
  \proofstepwidestar[1,2,5]{b\in A\land\bigl(b=a\lor \exists u\in A\;(b=u\star a)\bigr)}{%
      \FormulaRefAuto{SemigroupPrincipalLeftIdealMembership}[thm]{1,2,5}}
  \proofstepwidestar[1,2,5]{b=a\lor\exists x\in A\,(b=x\star a)}{%
    \rAEb{6}}
  \proofstepwidestar[8]{b=a}{\rA}
  \proofstepwidestar[4,8]{a=b}{%
      \FormulaRefAuto{a=b\vdash b=a}[thm]{8}}
  \proofstepwidestar[4,8]{b\in I}{\rIE{9,4}}
  \proofstepwidestar[11]{\exists x\in A\,(b=x\star a)}{\rA}
  \proofstepwidestar[12]{x\in A\land b=x\star a}{\rA}
  \proofstepwidestar[12]{x\in A}{\rAEa{12}}
  \proofstepwidestar[12]{b=x\star a}{\rAEb{12}}
  \proofstepwidestar[1,3]{I\subseteq A\land I\neq\varnothing\land\forall u\in A\,\forall v\in I\,(u\star v\in I)}{%
      \FormulaRefAuto{SemigroupLeftIdealCriterion}[thm]{1,3}}
  \proofstepwidestar[1,3]{%
    \forall u\in A\,\forall v\in I\,(u\star v\in I)}{%
    \rAEb{\rAEb{15}}}
  \proofstepwidestar[1,3,4,12]{x\star a\in I}{\rRE{4,\rUE{\rRE{13,\rUE{16}}}}}
  \proofstepwidestar[1,3,4,12]{b\in I}{\rIE{14,17}}
  \proofstepwidestar[1,3,4,11]{b\in I}{\rEE{11,12,18}}
  \proofstepwidestar[1,2,3,4,5]{b\in I}{\rOE{7,8,10,11,19}}
  \proofstepwidestar[1,2,3,4]{L_{A,\star}(a)\subseteq I}{%
    \FormulaRefAuto{SubsetIntroductionRule}[thm]{[5]\vdots 20}}
\end{tabproofwide}

\FormulaThmDeltaK[Minimalität des Rechtshauptideals]{%
  \begin{aligned}[t]
    &\SemigroupAlg{A}{\star}\dsep a\in A\dsep
      \mathsf{RId}(I;A,\star)\dsep a\in I\vdash{}\\[-2pt]
    &R_{A,\star}(a)\subseteq I
  \end{aligned}
}{SemigroupPrincipalRightIdealMinimal}{%
  \DeltaRow{Mengen}{A\dsep I\dsep a\dsep b\dsep y}
  \DeltaRow{Funktionen}{\star}
}
\begin{tabproofwide}
  \proofstepwidestar[1]{\SemigroupAlg{A}{\star}}{\rA}
  \proofstepwidestar[2]{a\in A}{\rA}
  \proofstepwidestar[3]{\mathsf{RId}(I;A,\star)}{\rA}
  \proofstepwidestar[4]{a\in I}{\rA}
  \proofstepwidestar[5]{b\in R_{A,\star}(a)}{\rA}
  \proofstepwidestar[1,2,5]{b\in A\land\bigl(b=a\lor \exists v\in A\;(b=a\star v)\bigr)}{%
      \FormulaRefAuto{SemigroupPrincipalRightIdealMembership}[thm]{1,2,5}}
  \proofstepwidestar[1,2,5]{b=a\lor\exists y\in A\,(b=a\star y)}{%
    \rAEb{6}}
  \proofstepwidestar[8]{b=a}{\rA}
  \proofstepwidestar[4,8]{a=b}{%
      \FormulaRefAuto{a=b\vdash b=a}[thm]{8}}
  \proofstepwidestar[4,8]{b\in I}{\rIE{9,4}}
  \proofstepwidestar[11]{\exists y\in A\,(b=a\star y)}{\rA}
  \proofstepwidestar[12]{y\in A\land b=a\star y}{\rA}
  \proofstepwidestar[12]{y\in A}{\rAEa{12}}
  \proofstepwidestar[12]{b=a\star y}{\rAEb{12}}
  \proofstepwidestar[1,3]{I\subseteq A\land I\neq\varnothing\land\forall u\in A\,\forall v\in I\,(v\star u\in I)}{%
      \FormulaRefAuto{SemigroupRightIdealCriterion}[thm]{1,3}}
  \proofstepwidestar[1,3]{%
    \forall u\in I\,\forall v\in A\,(u\star v\in I)}{%
    \rAEb{\rAEb{15}}}
  \proofstepwidestar[1,3,4,12]{a\star y\in I}{\rRE{13,\rUE{\rRE{4,\rUE{16}}}}}
  \proofstepwidestar[1,3,4,12]{b\in I}{\rIE{14,17}}
  \proofstepwidestar[1,3,4,11]{b\in I}{\rEE{11,12,18}}
  \proofstepwidestar[1,2,3,4,5]{b\in I}{\rOE{7,8,10,11,19}}
  \proofstepwidestar[1,2,3,4]{R_{A,\star}(a)\subseteq I}{%
    \FormulaRefAuto{SubsetIntroductionRule}[thm]{[5]\vdots 20}}
\end{tabproofwide}

\FormulaThmDeltaK[Das zweiseitige Hauptideal ist ein Ideal]{%
  \SemigroupAlg{A}{\star}\dsep a\in A\vdash
  \mathsf{Id}(J_{A,\star}(a);A,\star)%
}{SemigroupPrincipalTwoSidedIdealIsIdeal}{%
  \DeltaRow{Mengen}{A\dsep a\dsep b\dsep s\dsep x\dsep y}
  \DeltaRow{Funktionen}{\star}
}
Nach dem Elementkriterium besitzt jedes Element des Hauptideals eine von
vier Formen: Es ist der Erzeuger selbst, hat nur einen linken oder rechten
Faktor oder hat beide Faktoren.  Multiplikation von links beziehungsweise
rechts erh\"alt jede dieser Formen; die beiden Absorptionseigenschaften
werden deshalb getrennt bewiesen und erst im Schluss zusammengef\"uhrt.

\begin{tabproofsplitwide}
  \proofpartwideindR[Linksidealeigenschaft]{%
    \SemigroupAlg{A}{\star}\dsep a\in A\vdash
    \mathsf{LId}(J_{A,\star}(a);A,\star)
  }{%
    \SemigroupAlg{A}{\star}\dsep a\in A\vdash
    \mathsf{LId}(J_{A,\star}(a);A,\star)
  }[SemigroupPrincipalTwoSidedIdealLeftPart]
    \proofstepwidestarbreak[1]{\SemigroupAlg{A}{\star}}{%
      \rA}
    \proofstepwidestarbreak[2]{a\in A}{%
      \rA}
    \proofstepwidestarbreak[]{J_{A,\star}(a)\subseteq A}{%
      \FormulaRefAuto{\{x\in A\mid P(x)\}\subseteq A}[thm]}
    \proofstepwidestar[1,2]{a\in L_{A,\star}(a)\land a\in R_{A,\star}(a)\land a\in J_{A,\star}(a)}{%
      \FormulaRefAuto{SemigroupPrincipalIdealsContainGenerator}[thm]{1,2}}
    \proofstepwidestarbreak[1,2]{a\in J_{A,\star}(a)}{%
      \rAEb{\rAEb{4}}}
    \proofstepwidestarbreak[1,2]{J_{A,\star}(a)\neq\varnothing}{%
      \FormulaRefAuto{a\in A\vdash A\neq\varnothing}[thm]{5}}
    \proofstepwidestarbreak[7]{s\in A}{%
      \rA}
    \proofstepwidestarbreak[8]{b\in J_{A,\star}(a)}{%
      \rA}
    \proofstepwidestar[1,2,8]{b\in A\land\bigl(b=a\lor \exists u\in A\;(b=u\star a)\lor \exists v\in A\;(b=a\star v)\lor \exists u,v\in A\;(b=(u\star a)\star v)\bigr)}{%
      \FormulaRefAuto{SemigroupPrincipalTwoSidedIdealMembership}[thm]{1,2,8}}
    \proofstepwidestarbreak[1,2,8]{b\in A}{%
      \rAEa{9}}
    \proofstepwidestarbreak[1,2,7,8]{s\star b\in A}{%
      \FormulaRefAuto{SemigroupClosureAxiom}[ax]{1,7,10}}
    \proofstepwidestar[1,2,8]{b\in A\land\bigl(b=a\lor \exists u\in A\;(b=u\star a)\lor \exists v\in A\;(b=a\star v)\lor \exists u,v\in A\;(b=(u\star a)\star v)\bigr)}{%
      \FormulaRefAuto{SemigroupPrincipalTwoSidedIdealMembership}[thm]{1,2,8}}
    \proofstepwidestarbreak[1,2,8]{%
      \begin{aligned}[t]
        b=a\lor{}&\exists x\in A\,(b=x\star a)\\[-2pt]
        &{}\lor\exists y\in A\,(b=a\star y)\\[-2pt]
        &{}\lor\exists x,y\in A\,(b=(x\star a)\star y)
      \end{aligned}}{%
      \rAEb{12}}
    \proofstepwidestarbreak[14]{b=a}{%
      \rA}
    \proofstepwidestarbreak[7,14]{s\star b=s\star a}{%
      \FormulaRefAuto{BinaryFunctionRightArgumentCongruence}[thm]{14}}
    \proofstepwidestarbreak[7,14]{\exists u\in A\,(s\star b=u\star a)}{%
      \rEI{7,15}}
    \proofstepwidestarbreak[1,2,7,8,14]{s\star b\in J_{A,\star}(a)}{%
      \FormulaRefAuto{SemigroupPrincipalTwoSidedIdealMembership}[thm]{%
        1,2,\rAI{11,\rOIb{\rOIa{16}}}}}
    \proofstepwidestarbreak[18]{\exists x\in A\,(b=x\star a)}{%
      \rA}
    \proofstepwidestarbreak[19]{x\in A\land b=x\star a}{%
      \rA}
    \proofstepwidestarbreak[19]{x\in A}{%
      \rAEa{19}}
    \proofstepwidestarbreak[19]{b=x\star a}{%
      \rAEb{19}}
    \proofstepwidestarbreak[1,7,19]{s\star x\in A}{%
      \FormulaRefAuto{SemigroupClosureAxiom}[ax]{1,7,20}}
    \proofstepwidestarbreak[7,19]{s\star b=s\star(x\star a)}{%
      \FormulaRefAuto{BinaryFunctionRightArgumentCongruence}[thm]{21}}
    \proofstepwidestarbreak[1,2,7,19]{s\star(x\star a)=(s\star x)\star a}{%
      \FormulaRefAuto{SemigroupAssociativityAxiom}[ax]{1,7,20,2}}
    \proofstepwidestarbreak[1,2,7,19]{s\star b=(s\star x)\star a}{%
      \rIE{24,23}}
    \proofstepwidestarbreak[1,2,7,19]{\exists u\in A\,(s\star b=u\star a)}{%
      \rEI{22,25}}
    \proofstepwidestarbreak[1,2,7,8,19]{s\star b\in J_{A,\star}(a)}{%
      \FormulaRefAuto{SemigroupPrincipalTwoSidedIdealMembership}[thm]{%
        1,2,\rAI{11,\rOIb{\rOIa{26}}}}}
    \proofstepwidestarbreak[28]{\exists y\in A\,(b=a\star y)}{%
      \rA}
    \proofstepwidestarbreak[29]{y\in A\land b=a\star y}{%
      \rA}
    \proofstepwidestarbreak[29]{y\in A}{%
      \rAEa{29}}
    \proofstepwidestarbreak[29]{b=a\star y}{%
      \rAEb{29}}
    \proofstepwidestarbreak[7,29]{s\star b=s\star(a\star y)}{%
      \FormulaRefAuto{BinaryFunctionRightArgumentCongruence}[thm]{31}}
    \proofstepwidestarbreak[1,2,7,29]{s\star(a\star y)=(s\star a)\star y}{%
      \FormulaRefAuto{SemigroupAssociativityAxiom}[ax]{1,7,2,30}}
    \proofstepwidestarbreak[1,2,7,29]{s\star b=(s\star a)\star y}{%
      \rIE{33,32}}
    \proofstepwidestarbreak[1,2,7,29]{%
      \exists u,v\in A\,(s\star b=(u\star a)\star v)}{%
      \rEI{7,\rEI{30,34}}}
    \proofstepwidestarbreak[1,2,7,8,29]{s\star b\in J_{A,\star}(a)}{%
      \FormulaRefAuto{SemigroupPrincipalTwoSidedIdealMembership}[thm]{%
        1,2,\rAI{11,\rOIb{\rOIb{\rOIb{35}}}}}}
    \proofstepwidestarbreak[37]{\exists x,y\in A\,(b=(x\star a)\star y)}{%
      \rA}
    \proofstepwidestarbreak[38]{x\in A\land(y\in A\land b=(x\star a)\star y)}{%
      \rA}
    \proofstepwidestarbreak[38]{x\in A}{%
      \rAEa{38}}
    \proofstepwidestarbreak[38]{y\in A}{%
      \rAEa{\rAEb{38}}}
    \proofstepwidestarbreak[38]{b=(x\star a)\star y}{%
      \rAEb{\rAEb{38}}}
    \proofstepwidestarbreak[1,7,38]{s\star x\in A}{%
      \FormulaRefAuto{SemigroupClosureAxiom}[ax]{1,7,39}}
    \proofstepwidestarbreak[1,2,38]{x\star a\in A}{%
      \FormulaRefAuto{SemigroupClosureAxiom}[ax]{1,39,2}}
    \proofstepwidestarbreak[7,38]{s\star b=s\star((x\star a)\star y)}{%
      \FormulaRefAuto{BinaryFunctionRightArgumentCongruence}[thm]{41}}
    \proofstepwidestarbreak[1,2,7,38]{s\star((x\star a)\star y)=(s\star(x\star a))\star y}{%
      \FormulaRefAuto{SemigroupAssociativityAxiom}[ax]{1,7,43,40}}
    \proofstepwidestarbreak[1,2,7,38]{s\star(x\star a)=(s\star x)\star a}{%
      \FormulaRefAuto{SemigroupAssociativityAxiom}[ax]{1,7,39,2}}
    \proofstepwidestarbreak[1,2,7,38]{%
      s\star b=((s\star x)\star a)\star y}{%
      \rIE{45,46,44}}
    \proofstepwidestarbreak[1,2,7,38]{%
      \exists u,v\in A\,(s\star b=(u\star a)\star v)}{%
      \rEI{42,\rEI{40,47}}}
    \proofstepwidestarbreak[1,2,7,8,38]{s\star b\in J_{A,\star}(a)}{%
      \FormulaRefAuto{SemigroupPrincipalTwoSidedIdealMembership}[thm]{%
        1,2,\rAI{11,\rOIb{\rOIb{\rOIb{48}}}}}}
    \proofstepwidestarbreak[1,2,7,8,37]{s\star b\in J_{A,\star}(a)}{%
      \rEE{37,38,49}}
    \proofstepwidestarbreak[1,2,7,8,28]{s\star b\in J_{A,\star}(a)}{%
      \rEE{28,29,36}}
    \proofstepwidestarbreak[1,2,7,8,18]{s\star b\in J_{A,\star}(a)}{%
      \rEE{18,19,27}}
    \proofstepwidestarbreak[53]{\begin{aligned}[t]
        &\exists x\in A\,(b=x\star a)\\[-2pt]
        &{}\lor\exists y\in A\,(b=a\star y)\\[-2pt]
        &{}\lor\exists x,y\in A\,(b=(x\star a)\star y)
      \end{aligned}}{%
      \rA}
    \proofstepwidestarbreak[54]{\begin{aligned}[t]
        &\exists y\in A\,(b=a\star y)\\[-2pt]
        &{}\lor\exists x,y\in A\,(b=(x\star a)\star y)
      \end{aligned}}{%
      \rA}
    \proofstepwidestarbreak[1,2,7,8,54]{s\star b\in J_{A,\star}(a)}{%
      \rOE{54,28,51,37,50}}
    \proofstepwidestarbreak[1,2,7,8,53]{s\star b\in J_{A,\star}(a)}{%
      \rOE{53,18,52,54,55}}
    \proofstepwidestarbreak[1,2,7,8]{s\star b\in J_{A,\star}(a)}{%
      \rOE{13,14,17,53,56}}
    \proofstepwidestarbreak[1,2]{%
      \forall s\in A\,\forall b\in J_{A,\star}(a)\,
      (s\star b\in J_{A,\star}(a))}{%
      \rUI{\rRI{7,\rUI{\rRI{8,57}}}}}
    \proofstepwidestarbreak[1,2]{\mathsf{LId}(J_{A,\star}(a);A,\star)}{%
      \FormulaRefAuto{SemigroupLeftIdealCriterion}[thm]{1,\rAI{3,\rAI{6,58}}}}
  \closeproofpartwideindR

  \proofpartwideindR[Rechtsidealeigenschaft]{%
    \SemigroupAlg{A}{\star}\dsep a\in A\vdash
    \mathsf{RId}(J_{A,\star}(a);A,\star)
  }{%
    \SemigroupAlg{A}{\star}\dsep a\in A\vdash
    \mathsf{RId}(J_{A,\star}(a);A,\star)
  }[SemigroupPrincipalTwoSidedIdealRightPart]
    \proofstepwidestarbreak[1]{\SemigroupAlg{A}{\star}}{%
      \rA}
    \proofstepwidestarbreak[2]{a\in A}{%
      \rA}
    \proofstepwidestarbreak[]{J_{A,\star}(a)\subseteq A}{%
      \FormulaRefAuto{\{x\in A\mid P(x)\}\subseteq A}[thm]}
    \proofstepwidestar[1,2]{a\in L_{A,\star}(a)\land a\in R_{A,\star}(a)\land a\in J_{A,\star}(a)}{%
      \FormulaRefAuto{SemigroupPrincipalIdealsContainGenerator}[thm]{1,2}}
    \proofstepwidestarbreak[1,2]{a\in J_{A,\star}(a)}{%
      \rAEb{\rAEb{4}}}
    \proofstepwidestarbreak[1,2]{J_{A,\star}(a)\neq\varnothing}{%
      \FormulaRefAuto{a\in A\vdash A\neq\varnothing}[thm]{5}}
    \proofstepwidestarbreak[7]{b\in J_{A,\star}(a)}{%
      \rA}
    \proofstepwidestarbreak[8]{s\in A}{%
      \rA}
    \proofstepwidestar[1,2,7]{b\in A\land\bigl(b=a\lor \exists u\in A\;(b=u\star a)\lor \exists v\in A\;(b=a\star v)\lor \exists u,v\in A\;(b=(u\star a)\star v)\bigr)}{%
      \FormulaRefAuto{SemigroupPrincipalTwoSidedIdealMembership}[thm]{1,2,7}}
    \proofstepwidestarbreak[1,2,7]{b\in A}{%
      \rAEa{9}}
    \proofstepwidestarbreak[1,2,7,8]{b\star s\in A}{%
      \FormulaRefAuto{SemigroupClosureAxiom}[ax]{1,10,8}}
    \proofstepwidestar[1,2,7]{b\in A\land\bigl(b=a\lor \exists u\in A\;(b=u\star a)\lor \exists v\in A\;(b=a\star v)\lor \exists u,v\in A\;(b=(u\star a)\star v)\bigr)}{%
      \FormulaRefAuto{SemigroupPrincipalTwoSidedIdealMembership}[thm]{1,2,7}}
    \proofstepwidestarbreak[1,2,7]{%
      \begin{aligned}[t]
        b=a\lor{}&\exists x\in A\,(b=x\star a)\\[-2pt]
        &{}\lor\exists y\in A\,(b=a\star y)\\[-2pt]
        &{}\lor\exists x,y\in A\,(b=(x\star a)\star y)
      \end{aligned}}{%
      \rAEb{12}}
    \proofstepwidestarbreak[14]{b=a}{%
      \rA}
    \proofstepwidestarbreak[8,14]{b\star s=a\star s}{%
      \FormulaRefAuto{BinaryFunctionLeftArgumentCongruence}[thm]{14}}
    \proofstepwidestarbreak[8,14]{\exists v\in A\,(b\star s=a\star v)}{%
      \rEI{8,15}}
    \proofstepwidestarbreak[1,2,7,8,14]{b\star s\in J_{A,\star}(a)}{%
      \FormulaRefAuto{SemigroupPrincipalTwoSidedIdealMembership}[thm]{%
        1,2,\rAI{11,\rOIb{\rOIb{\rOIa{16}}}}}}
    \proofstepwidestarbreak[18]{\exists x\in A\,(b=x\star a)}{%
      \rA}
    \proofstepwidestarbreak[19]{x\in A\land b=x\star a}{%
      \rA}
    \proofstepwidestarbreak[19]{x\in A}{%
      \rAEa{19}}
    \proofstepwidestarbreak[19]{b=x\star a}{%
      \rAEb{19}}
    \proofstepwidestarbreak[1,2,8,19]{b\star s=(x\star a)\star s}{%
      \FormulaRefAuto{BinaryFunctionLeftArgumentCongruence}[thm]{21}}
    \proofstepwidestarbreak[1,2,8,19]{%
      \exists u,v\in A\,(b\star s=(u\star a)\star v)}{%
      \rEI{20,\rEI{8,22}}}
    \proofstepwidestarbreak[1,2,7,8,19]{b\star s\in J_{A,\star}(a)}{%
      \FormulaRefAuto{SemigroupPrincipalTwoSidedIdealMembership}[thm]{%
        1,2,\rAI{11,\rOIb{\rOIb{\rOIb{23}}}}}}
    \proofstepwidestarbreak[25]{\exists y\in A\,(b=a\star y)}{%
      \rA}
    \proofstepwidestarbreak[26]{y\in A\land b=a\star y}{%
      \rA}
    \proofstepwidestarbreak[26]{y\in A}{%
      \rAEa{26}}
    \proofstepwidestarbreak[26]{b=a\star y}{%
      \rAEb{26}}
    \proofstepwidestarbreak[1,8,26]{y\star s\in A}{%
      \FormulaRefAuto{SemigroupClosureAxiom}[ax]{1,27,8}}
    \proofstepwidestarbreak[8,26]{b\star s=(a\star y)\star s}{%
      \FormulaRefAuto{BinaryFunctionLeftArgumentCongruence}[thm]{28}}
    \proofstepwidestar[1,2,8,26]{(a\star y)\star s=a\star(y\star s)}{%
      \FormulaRefAuto{SemigroupAssociativityAxiom}[ax]{1,2,27,8}}
    \proofstepwidestarbreak[1,2,8,26]{b\star s=a\star(y\star s)}{%
      \rIE{31,30}}
    \proofstepwidestarbreak[1,2,8,26]{\exists v\in A\,(b\star s=a\star v)}{%
      \rEI{29,32}}
    \proofstepwidestarbreak[1,2,7,8,26]{b\star s\in J_{A,\star}(a)}{%
      \FormulaRefAuto{SemigroupPrincipalTwoSidedIdealMembership}[thm]{%
        1,2,\rAI{11,\rOIb{\rOIb{\rOIa{33}}}}}}
    \proofstepwidestarbreak[35]{\exists x,y\in A\,(b=(x\star a)\star y)}{%
      \rA}
    \proofstepwidestarbreak[36]{x\in A\land(y\in A\land b=(x\star a)\star y)}{%
      \rA}
    \proofstepwidestarbreak[36]{x\in A}{%
      \rAEa{36}}
    \proofstepwidestarbreak[36]{y\in A}{%
      \rAEa{\rAEb{36}}}
    \proofstepwidestarbreak[36]{b=(x\star a)\star y}{%
      \rAEb{\rAEb{36}}}
    \proofstepwidestarbreak[1,8,36]{y\star s\in A}{%
      \FormulaRefAuto{SemigroupClosureAxiom}[ax]{1,38,8}}
    \proofstepwidestarbreak[1,2,36]{x\star a\in A}{%
      \FormulaRefAuto{SemigroupClosureAxiom}[ax]{1,37,2}}
    \proofstepwidestarbreak[8,36]{b\star s=((x\star a)\star y)\star s}{%
      \FormulaRefAuto{BinaryFunctionLeftArgumentCongruence}[thm]{39}}
    \proofstepwidestarbreak[1,2,8,36]{((x\star a)\star y)\star s=(x\star a)\star(y\star s)}{%
      \FormulaRefAuto{SemigroupAssociativityAxiom}[ax]{1,41,38,8}}
    \proofstepwidestarbreak[1,2,8,36]{b\star s=(x\star a)\star(y\star s)}{%
      \rIE{43,42}}
    \proofstepwidestarbreak[1,2,8,36]{%
      \exists u,v\in A\,(b\star s=(u\star a)\star v)}{%
      \rEI{37,\rEI{40,44}}}
    \proofstepwidestarbreak[1,2,7,8,36]{b\star s\in J_{A,\star}(a)}{%
      \FormulaRefAuto{SemigroupPrincipalTwoSidedIdealMembership}[thm]{%
        1,2,\rAI{11,\rOIb{\rOIb{\rOIb{45}}}}}}
    \proofstepwidestarbreak[1,2,7,8,35]{b\star s\in J_{A,\star}(a)}{%
      \rEE{35,36,46}}
    \proofstepwidestarbreak[1,2,7,8,25]{b\star s\in J_{A,\star}(a)}{%
      \rEE{25,26,34}}
    \proofstepwidestarbreak[1,2,7,8,18]{b\star s\in J_{A,\star}(a)}{%
      \rEE{18,19,24}}
    \proofstepwidestarbreak[50]{\begin{aligned}[t]
        &\exists x\in A\,(b=x\star a)\\[-2pt]
        &{}\lor\exists y\in A\,(b=a\star y)\\[-2pt]
        &{}\lor\exists x,y\in A\,(b=(x\star a)\star y)
      \end{aligned}}{%
      \rA}
    \proofstepwidestarbreak[51]{\begin{aligned}[t]
        &\exists y\in A\,(b=a\star y)\\[-2pt]
        &{}\lor\exists x,y\in A\,(b=(x\star a)\star y)
      \end{aligned}}{%
      \rA}
    \proofstepwidestarbreak[1,2,7,8,51]{b\star s\in J_{A,\star}(a)}{%
      \rOE{51,25,48,35,47}}
    \proofstepwidestarbreak[1,2,7,8,50]{b\star s\in J_{A,\star}(a)}{%
      \rOE{50,18,49,51,52}}
    \proofstepwidestarbreak[1,2,7,8]{b\star s\in J_{A,\star}(a)}{%
      \rOE{13,14,17,50,53}}
    \proofstepwidestarbreak[1,2]{%
      \forall b\in J_{A,\star}(a)\,\forall s\in A\,
      (b\star s\in J_{A,\star}(a))}{%
      \rUI{\rRI{7,\rUI{\rRI{8,54}}}}}
    \proofstepwidestarbreak[1,2]{\mathsf{RId}(J_{A,\star}(a);A,\star)}{%
      \FormulaRefAuto{SemigroupRightIdealCriterion}[thm]{1,\rAI{3,\rAI{6,55}}}}
  \closeproofpartwideindR

  \proofpartwide{Schluss}
    \proofstepwidestarbreak[1]{\SemigroupAlg{A}{\star}}{\rA}
    \proofstepwidestarbreak[2]{a\in A}{\rA}
    \proofstepwidestarbreak[1,2]{\mathsf{LId}(J_{A,\star}(a);A,\star)}{%
      \FormulaRefAuto{SemigroupPrincipalTwoSidedIdealLeftPart}[thm]{1,2}}
    \proofstepwidestarbreak[1,2]{\mathsf{RId}(J_{A,\star}(a);A,\star)}{%
      \FormulaRefAuto{SemigroupPrincipalTwoSidedIdealRightPart}[thm]{1,2}}
    \proofstepwidestarbreak[1,2]{\mathsf{Id}(J_{A,\star}(a);A,\star)}{%
      \FormulaRefAuto{SemigroupIdealCriterion}[thm]{1,\rAI{3,4}}}
  \closeproofpartwide
\end{tabproofsplitwide}

\FormulaThmDeltaK[Minimalit\"at des zweiseitigen Hauptideals]{%
  \begin{aligned}[t]
    &\SemigroupAlg{A}{\star}\dsep a\in A\dsep
      \mathsf{Id}(I;A,\star)\dsep a\in I\vdash{}\\[-2pt]
    &J_{A,\star}(a)\subseteq I
  \end{aligned}
}{SemigroupPrincipalTwoSidedIdealMinimal}{%
  \DeltaRow{Mengen}{A\dsep I\dsep a\dsep b\dsep x\dsep y}
  \DeltaRow{Funktionen}{\star}
}
Das Elementkriterium liefert vier und nur vier Formen.  In jedem Fall folgt
die Zugeh\"origkeit zu (I) unmittelbar aus (a\in I) und der linken
beziehungsweise rechten Idealeigenschaft.

\begin{tabproofsplitwide}
  \proofpartwideindR[Erzeugerfall]{%
    \SemigroupAlg{A}{\star}\dsep a\in A\dsep\mathsf{Id}(I;A,\star)\dsep
    a\in I\dsep b=a\vdash b\in I
  }{%
    \SemigroupAlg{A}{\star}\dsep a\in A\dsep\mathsf{Id}(I;A,\star)\dsep
    a\in I\dsep b=a\vdash b\in I
  }[SemigroupPrincipalTwoSidedIdealMinimalGeneratorCase]
    \proofstepwidestar[1]{\SemigroupAlg{A}{\star}}{\rA}
    \proofstepwidestar[2]{a\in A}{\rA}
    \proofstepwidestar[3]{\mathsf{Id}(I;A,\star)}{\rA}
    \proofstepwidestar[4]{a\in I}{\rA}
    \proofstepwidestar[5]{b=a}{\rA}
    \proofstepwidestar[4,5]{a=b}{%
      \FormulaRefAuto{a=b\vdash b=a}[thm]{5}}
    \proofstepwidestar[4,5]{b\in I}{%
      \rIE{6,4}}
  \closeproofpartwideindR

  \proofpartwideindR[Fall eines linken Faktors]{%
    \begin{aligned}[t]
      &\SemigroupAlg{A}{\star}\dsep a\in A\dsep
        \mathsf{Id}(I;A,\star)\dsep a\in I\dsep{}\\[-2pt]
      &\exists x\in A\,(b=x\star a)\vdash b\in I
    \end{aligned}
  }{%
    \SemigroupAlg{A}{\star}\dsep a\in A\dsep
    \mathsf{Id}(I;A,\star)\dsep a\in I\dsep
    \exists x\in A\,(b=x\star a)\vdash b\in I
  }[SemigroupPrincipalTwoSidedIdealMinimalLeftFactorCase]
    \proofstepwidestar[1]{\SemigroupAlg{A}{\star}}{\rA}
    \proofstepwidestar[2]{a\in A}{\rA}
    \proofstepwidestar[3]{\mathsf{Id}(I;A,\star)}{\rA}
    \proofstepwidestar[4]{a\in I}{\rA}
    \proofstepwidestar[5]{\exists x\in A\,(b=x\star a)}{\rA}
    \proofstepwidestar[6]{x\in A\land b=x\star a}{\rA}
    \proofstepwidestar[6]{x\in A}{\rAEa{6}}
    \proofstepwidestar[6]{b=x\star a}{\rAEb{6}}
    \proofstepwidestar[1,3,4,6]{x\star a\in I}{%
      \FormulaRefAuto{SemigroupIdealLeftAbsorption}[thm]{1,3,7,4}}
    \proofstepwidestar[1,3,4,6]{b\in I}{\rIE{8,9}}
    \proofstepwidestar[1,3,4,5]{b\in I}{\rEE{5,6,10}}
  \closeproofpartwideindR

  \proofpartwideindR[Fall eines rechten Faktors]{%
    \begin{aligned}[t]
      &\SemigroupAlg{A}{\star}\dsep a\in A\dsep
        \mathsf{Id}(I;A,\star)\dsep a\in I\dsep{}\\[-2pt]
      &\exists y\in A\,(b=a\star y)\vdash b\in I
    \end{aligned}
  }{%
    \SemigroupAlg{A}{\star}\dsep a\in A\dsep
    \mathsf{Id}(I;A,\star)\dsep a\in I\dsep
    \exists y\in A\,(b=a\star y)\vdash b\in I
  }[SemigroupPrincipalTwoSidedIdealMinimalRightFactorCase]
    \proofstepwidestar[1]{\SemigroupAlg{A}{\star}}{\rA}
    \proofstepwidestar[2]{a\in A}{\rA}
    \proofstepwidestar[3]{\mathsf{Id}(I;A,\star)}{\rA}
    \proofstepwidestar[4]{a\in I}{\rA}
    \proofstepwidestar[5]{\exists y\in A\,(b=a\star y)}{\rA}
    \proofstepwidestar[6]{y\in A\land b=a\star y}{\rA}
    \proofstepwidestar[6]{y\in A}{\rAEa{6}}
    \proofstepwidestar[6]{b=a\star y}{\rAEb{6}}
    \proofstepwidestar[1,3,4,6]{a\star y\in I}{%
      \FormulaRefAuto{SemigroupIdealRightAbsorption}[thm]{1,3,4,7}}
    \proofstepwidestar[1,3,4,6]{b\in I}{\rIE{8,9}}
    \proofstepwidestar[1,3,4,5]{b\in I}{\rEE{5,6,10}}
  \closeproofpartwideindR

  \proofpartwideindR[Fall zweier Faktoren]{%
    \begin{aligned}[t]
      &\SemigroupAlg{A}{\star}\dsep a\in A\dsep
        \mathsf{Id}(I;A,\star)\dsep a\in I\dsep{}\\[-2pt]
      &\exists x,y\in A\,(b=(x\star a)\star y)\vdash b\in I
    \end{aligned}
  }{%
    \begin{aligned}[t]
      &\SemigroupAlg{A}{\star}\dsep a\in A\dsep
        \mathsf{Id}(I;A,\star)\dsep a\in I\dsep{}\\[-2pt]
      &\exists x,y\in A\,(b=(x\star a)\star y)\vdash b\in I
    \end{aligned}
  }[SemigroupPrincipalTwoSidedIdealMinimalTwoFactorCase]
    \proofstepwidestar[1]{\SemigroupAlg{A}{\star}}{\rA}
    \proofstepwidestar[2]{a\in A}{\rA}
    \proofstepwidestar[3]{\mathsf{Id}(I;A,\star)}{\rA}
    \proofstepwidestar[4]{a\in I}{\rA}
    \proofstepwidestar[5]{\exists x,y\in A\,(b=(x\star a)\star y)}{\rA}
    \proofstepwidestar[6]{%
      x\in A\land\bigl(y\in A\land b=(x\star a)\star y\bigr)}{\rA}
    \proofstepwidestar[6]{x\in A}{\rAEa{6}}
    \proofstepwidestar[6]{y\in A}{\rAEa{\rAEb{6}}}
    \proofstepwidestar[6]{b=(x\star a)\star y}{\rAEb{\rAEb{6}}}
    \proofstepwidestar[1,3,4,6]{x\star a\in I}{%
      \FormulaRefAuto{SemigroupIdealLeftAbsorption}[thm]{1,3,7,4}}
    \proofstepwidestar[1,3,4,6]{(x\star a)\star y\in I}{%
      \FormulaRefAuto{SemigroupIdealRightAbsorption}[thm]{1,3,10,8}}
    \proofstepwidestar[1,3,4,6]{b\in I}{\rIE{9,11}}
    \proofstepwidestar[1,3,4,5]{b\in I}{\rEE{5,6,12}}
  \closeproofpartwideindR

  \proofpartwide{Fallunterscheidung und Schluss}
    \proofstepwidestar[1]{\SemigroupAlg{A}{\star}}{\rA}
    \proofstepwidestar[2]{a\in A}{\rA}
    \proofstepwidestar[3]{\mathsf{Id}(I;A,\star)}{\rA}
    \proofstepwidestar[4]{a\in I}{\rA}
    \proofstepwidestar[5]{b\in J_{A,\star}(a)}{\rA}
    \proofstepwidestar[1,2,5]{b\in A\land\bigl(b=a\lor \exists u\in A\;(b=u\star a)\lor \exists v\in A\;(b=a\star v)\lor \exists u,v\in A\;(b=(u\star a)\star v)\bigr)}{%
      \FormulaRefAuto{SemigroupPrincipalTwoSidedIdealMembership}[thm]{1,2,5}}
    \proofstepwidestarbreak[1,2,5]{%
      \begin{aligned}[t]
        b=a\lor{}&\exists x\in A\,(b=x\star a)\\[-2pt]
        &{}\lor\exists y\in A\,(b=a\star y)\\[-2pt]
        &{}\lor\exists x,y\in A\,(b=(x\star a)\star y)
      \end{aligned}}{%
      \rAEb{6}}
    \proofstepwidestar[8]{b=a}{\rA}
    \proofstepwidestar[1,2,3,4,8]{b\in I}{%
      \FormulaRefAuto{SemigroupPrincipalTwoSidedIdealMinimalGeneratorCase}[thm]{%
        1,2,3,4,8}}
    \proofstepwidestar[10]{\exists x\in A\,(b=x\star a)}{\rA}
    \proofstepwidestar[1,2,3,4,10]{b\in I}{%
      \FormulaRefAuto{SemigroupPrincipalTwoSidedIdealMinimalLeftFactorCase}[thm]{%
        1,2,3,4,10}}
    \proofstepwidestar[12]{\exists y\in A\,(b=a\star y)}{\rA}
    \proofstepwidestar[1,2,3,4,12]{b\in I}{%
      \FormulaRefAuto{SemigroupPrincipalTwoSidedIdealMinimalRightFactorCase}[thm]{%
        1,2,3,4,12}}
    \proofstepwidestar[14]{\exists x,y\in A\,(b=(x\star a)\star y)}{\rA}
    \proofstepwidestar[1,2,3,4,14]{b\in I}{%
      \FormulaRefAuto{SemigroupPrincipalTwoSidedIdealMinimalTwoFactorCase}[thm]{%
        1,2,3,4,14}}
    \proofstepwidestar[1,2,3,4,5]{b\in I}{%
      \rOE{7,8,9,10,11,12,13,14,15}}
    \proofstepwidestar[1,2,3,4]{J_{A,\star}(a)\subseteq I}{%
      \FormulaRefAuto{SubsetIntroductionRule}[thm]{[5]\vdots 16}}
  \closeproofpartwide
\end{tabproofsplitwide}

\section{Teilbarkeit in Halbgruppen}

Ohne neutrales Element wäre die Bedingung \(b=a\star c\) nicht notwendig
reflexiv. Deshalb nehmen die folgenden Definitionen den Fall \(a=b\)
ausdrücklich auf. Die zuvor definierten Hauptideale erfassen genau diese
fehlenden Faktoren; eine unbestimmte Schreibweise wie \(A^1\) wird daher
nicht benötigt.

\FormulaDefDeltaK[Linksteilbarkeit in einer Halbgruppe]{%
  \SemigroupLeftDivides{A,\star}{a}{b}%
}{SemigroupLeftDividesDef}{%
  \DeltaRow{Mengen}{A\dsep a\dsep b\dsep c}
  \DeltaRow{Funktionen}{\star}
  \DeltaPrem{Halbgruppe}{\SemigroupAlg{A}{\star}}
  \DeltaRow{Elemente}{a,b\in A}
  \DeltaRow{Linksteilbarkeit}{%
    b=a\lor\exists c\in A\,(b=a\star c)}
  \DeltaRow{\textbf{Neues Symbol}}{}
  \DeltaPrem{Linksteilbarkeit}{%
    \SemigroupLeftDivides{A,\star}{a}{b}}
}

\FormulaDefDeltaK[Rechtsteilbarkeit in einer Halbgruppe]{%
  \SemigroupRightDivides{A,\star}{a}{b}%
}{SemigroupRightDividesDef}{%
  \DeltaRow{Mengen}{A\dsep a\dsep b\dsep c}
  \DeltaRow{Funktionen}{\star}
  \DeltaPrem{Halbgruppe}{\SemigroupAlg{A}{\star}}
  \DeltaRow{Elemente}{a,b\in A}
  \DeltaRow{Rechtsteilbarkeit}{%
    b=a\lor\exists c\in A\,(b=c\star a)}
  \DeltaRow{\textbf{Neues Symbol}}{}
  \DeltaPrem{Rechtsteilbarkeit}{%
    \SemigroupRightDivides{A,\star}{a}{b}}
}

\FormulaDefDeltaK[Zweiseitige Teilbarkeit in einer Halbgruppe]{%
  \SemigroupDivides{A,\star}{a}{b}%
}{SemigroupTwoSidedDividesDef}{%
  \DeltaRow{Mengen}{A\dsep a\dsep b\dsep x\dsep y}
  \DeltaRow{Funktionen}{\star}
  \DeltaPrem{Halbgruppe}{\SemigroupAlg{A}{\star}}
  \DeltaRow{Elemente}{a,b\in A}
  \DeltaRow{Zweiseitige Teilbarkeit}{%
    \begin{aligned}[t]
      &b=a\lor\exists x\in A\,(b=x\star a)\\[-2pt]
      &{}\lor\exists y\in A\,(b=a\star y)\\[-2pt]
      &{}\lor\exists x,y\in A\,
        \bigl(b=(x\star a)\star y\bigr)
    \end{aligned}}
  \DeltaRow{\textbf{Neues Symbol}}{}
  \DeltaPrem{Zweiseitige Teilbarkeit}{%
    \SemigroupDivides{A,\star}{a}{b}}
}

\FormulaThmDeltaK[Linksteilbarkeit als Mitgliedschaft im Rechtshauptideal]{%
  \begin{aligned}[t]
    &\SemigroupAlg{A}{\star}\dsep a,b\in A\vdash{}\\[-2pt]
    &\SemigroupLeftDivides{A,\star}{a}{b}
      \leftrightarrow b\in R_{A,\star}(a)
  \end{aligned}
}{SemigroupLeftDividesPrincipalRightIdeal}{%
  \DeltaRow{Mengen}{A\dsep a\dsep b\dsep c}
  \DeltaRow{Funktionen}{\star}
}
\begin{tabproofwide}
  \proofstepwidestar[1]{\SemigroupAlg{A}{\star}}{\rA}
  \proofstepwidestar[2]{a\in A}{\rA}
  \proofstepwidestar[3]{b\in A}{\rA}
  \proofstepwidestar[4]{\SemigroupLeftDivides{A,\star}{a}{b}}{\rA}
  \proofstepwidestar[1,2,3,4]{b=a\lor\exists c\in A\,(b=a\star c)}{%
    \FormulaRefAuto{SemigroupLeftDividesDef}[def]{1,2,3,4}}
  \proofstepwidestar[1,2,3,4]{b\in R_{A,\star}(a)}{%
    \FormulaRefAuto{SemigroupPrincipalRightIdealMembership}[thm]{%
      1,2,\rAI{3,5}}}
  \proofstepwidestar[1,2,3]{%
    \SemigroupLeftDivides{A,\star}{a}{b}\rightarrow
    b\in R_{A,\star}(a)}{\rRI{4,6}}
  \proofstepwidestar[8]{b\in R_{A,\star}(a)}{\rA}
  \proofstepwidestar[1,2,8]{b\in A\land\bigl(b=a\lor \exists v\in A\;(b=a\star v)\bigr)}{%
      \FormulaRefAuto{SemigroupPrincipalRightIdealMembership}[thm]{1,2,8}}
  \proofstepwidestar[1,2,8]{b=a\lor\exists c\in A\,(b=a\star c)}{%
    \rAEb{9}}
  \proofstepwidestar[1,2,3,8]{\SemigroupLeftDivides{A,\star}{a}{b}}{%
    \FormulaRefAuto{SemigroupLeftDividesDef}[def]{1,2,3,10}}
  \proofstepwidestar[1,2,3]{%
    b\in R_{A,\star}(a)\rightarrow
    \SemigroupLeftDivides{A,\star}{a}{b}}{\rRI{8,11}}
  \proofstepwidestar[1,2,3]{%
    \SemigroupLeftDivides{A,\star}{a}{b}
    \leftrightarrow b\in R_{A,\star}(a)}{\rLRI{7,12}}
\end{tabproofwide}

\FormulaThmDeltaK[Rechtsteilbarkeit als Mitgliedschaft im Linkshauptideal]{%
  \begin{aligned}[t]
    &\SemigroupAlg{A}{\star}\dsep a,b\in A\vdash{}\\[-2pt]
    &\SemigroupRightDivides{A,\star}{a}{b}
      \leftrightarrow b\in L_{A,\star}(a)
  \end{aligned}
}{SemigroupRightDividesPrincipalLeftIdeal}{%
  \DeltaRow{Mengen}{A\dsep a\dsep b\dsep c}
  \DeltaRow{Funktionen}{\star}
}
\begin{tabproofwide}
  \proofstepwidestar[1]{\SemigroupAlg{A}{\star}}{\rA}
  \proofstepwidestar[2]{a\in A}{\rA}
  \proofstepwidestar[3]{b\in A}{\rA}
  \proofstepwidestar[4]{\SemigroupRightDivides{A,\star}{a}{b}}{\rA}
  \proofstepwidestar[1,2,3,4]{b=a\lor\exists c\in A\,(b=c\star a)}{%
    \FormulaRefAuto{SemigroupRightDividesDef}[def]{1,2,3,4}}
  \proofstepwidestar[1,2,3,4]{b\in L_{A,\star}(a)}{%
    \FormulaRefAuto{SemigroupPrincipalLeftIdealMembership}[thm]{%
      1,2,\rAI{3,5}}}
  \proofstepwidestar[1,2,3]{%
    \SemigroupRightDivides{A,\star}{a}{b}\rightarrow
    b\in L_{A,\star}(a)}{\rRI{4,6}}
  \proofstepwidestar[8]{b\in L_{A,\star}(a)}{\rA}
  \proofstepwidestar[1,2,8]{b\in A\land\bigl(b=a\lor \exists u\in A\;(b=u\star a)\bigr)}{%
      \FormulaRefAuto{SemigroupPrincipalLeftIdealMembership}[thm]{1,2,8}}
  \proofstepwidestar[1,2,8]{b=a\lor\exists c\in A\,(b=c\star a)}{%
    \rAEb{9}}
  \proofstepwidestar[1,2,3,8]{\SemigroupRightDivides{A,\star}{a}{b}}{%
    \FormulaRefAuto{SemigroupRightDividesDef}[def]{1,2,3,10}}
  \proofstepwidestar[1,2,3]{%
    b\in L_{A,\star}(a)\rightarrow
    \SemigroupRightDivides{A,\star}{a}{b}}{\rRI{8,11}}
  \proofstepwidestar[1,2,3]{%
    \SemigroupRightDivides{A,\star}{a}{b}
    \leftrightarrow b\in L_{A,\star}(a)}{\rLRI{7,12}}
\end{tabproofwide}

\FormulaThmDeltaK[Zweiseitige Teilbarkeit als Hauptidealmitgliedschaft]{%
  \begin{aligned}[t]
    &\SemigroupAlg{A}{\star}\dsep a,b\in A\vdash{}\\[-2pt]
    &\SemigroupDivides{A,\star}{a}{b}
      \leftrightarrow b\in J_{A,\star}(a)
  \end{aligned}
}{SemigroupTwoSidedDividesPrincipalIdeal}{%
  \DeltaRow{Mengen}{A\dsep a\dsep b\dsep x\dsep y}
  \DeltaRow{Funktionen}{\star}
}
\begin{tabproofwide}
  \proofstepwidestar[1]{\SemigroupAlg{A}{\star}}{\rA}
  \proofstepwidestar[2]{a\in A}{\rA}
  \proofstepwidestar[3]{b\in A}{\rA}
  \proofstepwidestar[4]{\SemigroupDivides{A,\star}{a}{b}}{\rA}
  \proofstepwidestarbreak[1,2,3,4]{%
    \begin{aligned}[t]
      b=a\lor{}&\exists x\in A\,(b=x\star a)\\[-2pt]
      &{}\lor\exists y\in A\,(b=a\star y)\\[-2pt]
      &{}\lor\exists x,y\in A\,(b=(x\star a)\star y)
    \end{aligned}}{%
    \FormulaRefAuto{SemigroupTwoSidedDividesDef}[def]{1,2,3,4}}
  \proofstepwidestar[1,2,3,4]{b\in J_{A,\star}(a)}{%
    \FormulaRefAuto{SemigroupPrincipalTwoSidedIdealMembership}[thm]{%
      1,2,\rAI{3,5}}}
  \proofstepwidestar[1,2,3]{%
    \SemigroupDivides{A,\star}{a}{b}\rightarrow
    b\in J_{A,\star}(a)}{\rRI{4,6}}
  \proofstepwidestar[8]{b\in J_{A,\star}(a)}{\rA}
  \proofstepwidestar[1,2,8]{b\in A\land\bigl(b=a\lor \exists u\in A\;(b=u\star a)\lor \exists v\in A\;(b=a\star v)\lor \exists u,v\in A\;(b=(u\star a)\star v)\bigr)}{%
      \FormulaRefAuto{SemigroupPrincipalTwoSidedIdealMembership}[thm]{1,2,8}}
  \proofstepwidestarbreak[1,2,8]{%
    \begin{aligned}[t]
      b=a\lor{}&\exists x\in A\,(b=x\star a)\\[-2pt]
      &{}\lor\exists y\in A\,(b=a\star y)\\[-2pt]
      &{}\lor\exists x,y\in A\,(b=(x\star a)\star y)
    \end{aligned}}{%
    \rAEb{9}}
  \proofstepwidestar[1,2,3,8]{\SemigroupDivides{A,\star}{a}{b}}{%
    \FormulaRefAuto{SemigroupTwoSidedDividesDef}[def]{1,2,3,10}}
  \proofstepwidestar[1,2,3]{%
    b\in J_{A,\star}(a)\rightarrow
    \SemigroupDivides{A,\star}{a}{b}}{\rRI{8,11}}
  \proofstepwidestar[1,2,3]{%
    \SemigroupDivides{A,\star}{a}{b}
    \leftrightarrow b\in J_{A,\star}(a)}{\rLRI{7,12}}
\end{tabproofwide}

\FormulaThmDeltaK[Reflexivität der drei Teilbarkeiten]{%
  \begin{aligned}[t]
    &\SemigroupAlg{A}{\star}\dsep a\in A\vdash{}\\[-2pt]
    &\SemigroupLeftDivides{A,\star}{a}{a}\land
      \SemigroupRightDivides{A,\star}{a}{a}\land
      \SemigroupDivides{A,\star}{a}{a}
  \end{aligned}
}{SemigroupDivisibilityReflexive}{%
  \DeltaRow{Mengen}{A\dsep a}
  \DeltaRow{Funktionen}{\star}
}
\begin{tabproofwide}
  \proofstepwidestar[1]{\SemigroupAlg{A}{\star}}{\rA}
  \proofstepwidestar[2]{a\in A}{\rA}
  \proofstepwidestar[1,2]{\SemigroupLeftDivides{A,\star}{a}{a}}{%
    \FormulaRefAuto{SemigroupLeftDividesDef}[def]{1,2,2,\rOIa{\rII}}}
  \proofstepwidestar[1,2]{\SemigroupRightDivides{A,\star}{a}{a}}{%
    \FormulaRefAuto{SemigroupRightDividesDef}[def]{1,2,2,\rOIa{\rII}}}
  \proofstepwidestar[1,2]{\SemigroupDivides{A,\star}{a}{a}}{%
    \FormulaRefAuto{SemigroupTwoSidedDividesDef}[def]{1,2,2,\rOIa{\rII}}}
  \proofstepwidestar[1,2]{%
    \SemigroupLeftDivides{A,\star}{a}{a}\land
    \SemigroupRightDivides{A,\star}{a}{a}\land
    \SemigroupDivides{A,\star}{a}{a}}{\rAI{3,\rAI{4,5}}}
\end{tabproofwide}

\FormulaThmDeltaK[Transitivität der Linksteilbarkeit]{%
  \begin{aligned}[t]
    &\SemigroupAlg{A}{\star}\dsep a,b,c\in A\dsep
      \SemigroupLeftDivides{A,\star}{a}{b}\dsep
      \SemigroupLeftDivides{A,\star}{b}{c}\vdash{}\\[-2pt]
    &\SemigroupLeftDivides{A,\star}{a}{c}
  \end{aligned}
}{SemigroupLeftDividesTransitive}{%
  \DeltaRow{Mengen}{A\dsep a\dsep b\dsep c}
  \DeltaRow{Funktionen}{\star}
}
\begin{tabproofwide}
  \proofstepwidestar[1]{\SemigroupAlg{A}{\star}}{\rA}
  \proofstepwidestar[2]{a\in A}{\rA}
  \proofstepwidestar[3]{b\in A}{\rA}
  \proofstepwidestar[4]{c\in A}{\rA}
  \proofstepwidestar[5]{\SemigroupLeftDivides{A,\star}{a}{b}}{\rA}
  \proofstepwidestar[6]{\SemigroupLeftDivides{A,\star}{b}{c}}{\rA}
  \proofstepwidestar[1,2,3,5]{b\in R_{A,\star}(a)}{%
    \FormulaRefAuto{SemigroupLeftDividesPrincipalRightIdeal}[thm]{1,2,3,5}}
  \proofstepwidestar[1,3,4,6]{c\in R_{A,\star}(b)}{%
    \FormulaRefAuto{SemigroupLeftDividesPrincipalRightIdeal}[thm]{1,3,4,6}}
  \proofstepwidestar[1,2]{\mathsf{RId}(R_{A,\star}(a);A,\star)}{%
    \FormulaRefAuto{SemigroupPrincipalRightIdealIsRightIdeal}[thm]{1,2}}
  \proofstepwidestar[1,2,3,5]{R_{A,\star}(b)\subseteq R_{A,\star}(a)}{%
    \FormulaRefAuto{SemigroupPrincipalRightIdealMinimal}[thm]{1,3,9,7}}
  \proofstepwidestar[1,2,3,4,5,6]{c\in R_{A,\star}(a)}{%
    \FormulaRefAuto{A\subseteq B,\,x\in A\vdash x\in B}[thm]{10,8}}
  \proofstepwidestar[1,2,3,4,5,6]{\SemigroupLeftDivides{A,\star}{a}{c}}{%
    \FormulaRefAuto{SemigroupLeftDividesPrincipalRightIdeal}[thm]{1,2,4,11}}
\end{tabproofwide}

\FormulaThmDeltaK[Transitivität der Rechtsteilbarkeit]{%
  \begin{aligned}[t]
    &\SemigroupAlg{A}{\star}\dsep a,b,c\in A\dsep
      \SemigroupRightDivides{A,\star}{a}{b}\dsep
      \SemigroupRightDivides{A,\star}{b}{c}\vdash{}\\[-2pt]
    &\SemigroupRightDivides{A,\star}{a}{c}
  \end{aligned}
}{SemigroupRightDividesTransitive}{%
  \DeltaRow{Mengen}{A\dsep a\dsep b\dsep c}
  \DeltaRow{Funktionen}{\star}
}
\begin{tabproofwide}
  \proofstepwidestar[1]{\SemigroupAlg{A}{\star}}{\rA}
  \proofstepwidestar[2]{a\in A}{\rA}
  \proofstepwidestar[3]{b\in A}{\rA}
  \proofstepwidestar[4]{c\in A}{\rA}
  \proofstepwidestar[5]{\SemigroupRightDivides{A,\star}{a}{b}}{\rA}
  \proofstepwidestar[6]{\SemigroupRightDivides{A,\star}{b}{c}}{\rA}
  \proofstepwidestar[1,2,3,5]{b\in L_{A,\star}(a)}{%
    \FormulaRefAuto{SemigroupRightDividesPrincipalLeftIdeal}[thm]{1,2,3,5}}
  \proofstepwidestar[1,3,4,6]{c\in L_{A,\star}(b)}{%
    \FormulaRefAuto{SemigroupRightDividesPrincipalLeftIdeal}[thm]{1,3,4,6}}
  \proofstepwidestar[1,2]{\mathsf{LId}(L_{A,\star}(a);A,\star)}{%
    \FormulaRefAuto{SemigroupPrincipalLeftIdealIsLeftIdeal}[thm]{1,2}}
  \proofstepwidestar[1,2,3,5]{L_{A,\star}(b)\subseteq L_{A,\star}(a)}{%
    \FormulaRefAuto{SemigroupPrincipalLeftIdealMinimal}[thm]{1,3,9,7}}
  \proofstepwidestar[1,2,3,4,5,6]{c\in L_{A,\star}(a)}{%
    \FormulaRefAuto{A\subseteq B,\,x\in A\vdash x\in B}[thm]{10,8}}
  \proofstepwidestar[1,2,3,4,5,6]{\SemigroupRightDivides{A,\star}{a}{c}}{%
    \FormulaRefAuto{SemigroupRightDividesPrincipalLeftIdeal}[thm]{1,2,4,11}}
\end{tabproofwide}

\FormulaThmDeltaK[Transitivität der zweiseitigen Teilbarkeit]{%
  \begin{aligned}[t]
    &\SemigroupAlg{A}{\star}\dsep a,b,c\in A\dsep
      \SemigroupDivides{A,\star}{a}{b}\dsep
      \SemigroupDivides{A,\star}{b}{c}\vdash{}\\[-2pt]
    &\SemigroupDivides{A,\star}{a}{c}
  \end{aligned}
}{SemigroupTwoSidedDividesTransitive}{%
  \DeltaRow{Mengen}{A\dsep a\dsep b\dsep c}
  \DeltaRow{Funktionen}{\star}
}
\begin{tabproofwide}
  \proofstepwidestar[1]{\SemigroupAlg{A}{\star}}{\rA}
  \proofstepwidestar[2]{a\in A}{\rA}
  \proofstepwidestar[3]{b\in A}{\rA}
  \proofstepwidestar[4]{c\in A}{\rA}
  \proofstepwidestar[5]{\SemigroupDivides{A,\star}{a}{b}}{\rA}
  \proofstepwidestar[6]{\SemigroupDivides{A,\star}{b}{c}}{\rA}
  \proofstepwidestar[1,2,3,5]{b\in J_{A,\star}(a)}{%
    \FormulaRefAuto{SemigroupTwoSidedDividesPrincipalIdeal}[thm]{1,2,3,5}}
  \proofstepwidestar[1,3,4,6]{c\in J_{A,\star}(b)}{%
    \FormulaRefAuto{SemigroupTwoSidedDividesPrincipalIdeal}[thm]{1,3,4,6}}
  \proofstepwidestar[1,2]{\mathsf{Id}(J_{A,\star}(a);A,\star)}{%
    \FormulaRefAuto{SemigroupPrincipalTwoSidedIdealIsIdeal}[thm]{1,2}}
  \proofstepwidestar[1,2,3,5]{J_{A,\star}(b)\subseteq J_{A,\star}(a)}{%
    \FormulaRefAuto{SemigroupPrincipalTwoSidedIdealMinimal}[thm]{1,3,9,7}}
  \proofstepwidestar[1,2,3,4,5,6]{c\in J_{A,\star}(a)}{%
    \FormulaRefAuto{A\subseteq B,\,x\in A\vdash x\in B}[thm]{10,8}}
  \proofstepwidestar[1,2,3,4,5,6]{\SemigroupDivides{A,\star}{a}{c}}{%
    \FormulaRefAuto{SemigroupTwoSidedDividesPrincipalIdeal}[thm]{1,2,4,11}}
\end{tabproofwide}

Linksteilbarkeit, Rechtsteilbarkeit und zweiseitige Teilbarkeit sind
reflexiv und transitiv.  In einer kommutativen Halbgruppe stimmen die drei
Relationen überein.  Für Monoide ist die ausdrückliche Alternative \(b=a\)
überflüssig, weil das neutrale Element als fehlender Faktor dient.

\section{Green-Relationen}

Die Green-Relationen vergleichen Elemente danach, welche Hauptideale sie
erzeugen. Die Indizes (A,\star) bleiben sichtbar, damit die Relation stets
eindeutig zu ihrer umgebenden Halbgruppe gehört.

\FormulaDefDeltaK[Greensche \(\mathcal L\)-Relation]{%
  \SemigroupAlg{A}{\star}\dsep a,b\in A\vdash
  a\mathrel{\mathcal L_{A,\star}}b\coloneqq
  L_{A,\star}(a)=L_{A,\star}(b)%
}{SemigroupGreenLDef}{%
  \DeltaRow{Mengen}{A\dsep a\dsep b}
  \DeltaRow{Funktionen}{\star}
  \DeltaRow{Relationsoperatoren}{\mathcal L_{A,\star}}
}

\FormulaDefDeltaK[Greensche \(\mathcal R\)-Relation]{%
  \SemigroupAlg{A}{\star}\dsep a,b\in A\vdash
  a\mathrel{\mathcal R_{A,\star}}b\coloneqq
  R_{A,\star}(a)=R_{A,\star}(b)%
}{SemigroupGreenRDef}{%
  \DeltaRow{Mengen}{A\dsep a\dsep b}
  \DeltaRow{Funktionen}{\star}
  \DeltaRow{Relationsoperatoren}{\mathcal R_{A,\star}}
}

\FormulaDefDeltaK[Greensche \(\mathcal J\)-Relation]{%
  \SemigroupAlg{A}{\star}\dsep a,b\in A\vdash
  a\mathrel{\mathcal J_{A,\star}}b\coloneqq
  J_{A,\star}(a)=J_{A,\star}(b)%
}{SemigroupGreenJDef}{%
  \DeltaRow{Mengen}{A\dsep a\dsep b}
  \DeltaRow{Funktionen}{\star}
  \DeltaRow{Relationsoperatoren}{\mathcal J_{A,\star}}
}

\FormulaDefDeltaK[Greensche \(\mathcal H\)-Relation]{%
  \begin{aligned}[t]
    &\SemigroupAlg{A}{\star}\dsep a,b\in A\vdash{}\\[-2pt]
    &a\mathrel{\mathcal H_{A,\star}}b\coloneqq
      \bigl(a\mathrel{\mathcal L_{A,\star}}b\land
            a\mathrel{\mathcal R_{A,\star}}b\bigr)
  \end{aligned}
}{SemigroupGreenHDef}{%
  \DeltaRow{Mengen}{A\dsep a\dsep b}
  \DeltaRow{Funktionen}{\star}
  \DeltaRow{Relationsoperatoren}{%
    \mathcal L_{A,\star}\dsep\mathcal R_{A,\star}\dsep
    \mathcal H_{A,\star}}
}

\FormulaThmDeltaK[\(\mathcal L\) als gegenseitige Rechtsteilbarkeit]{%
  \begin{aligned}[t]
    &\SemigroupAlg{A}{\star}\dsep a,b\in A\vdash{}\\[-2pt]
    &a\mathrel{\mathcal L_{A,\star}}b\leftrightarrow
      \Bigl(\SemigroupRightDivides{A,\star}{a}{b}\land
            \SemigroupRightDivides{A,\star}{b}{a}\Bigr)
  \end{aligned}
}{SemigroupGreenLMutualRightDivisibility}{%
  \DeltaRow{Mengen}{A\dsep a\dsep b}
  \DeltaRow{Funktionen}{\star}
  \DeltaRow{Relationsoperatoren}{\mathcal L_{A,\star}}
}
\begin{tabproofsplitwide}
\proofpartwide{Von gleichen Hauptidealen zur Teilbarkeit}
  \proofstepwidestar[1]{\SemigroupAlg{A}{\star}}{\rA}
  \proofstepwidestar[2]{a\in A}{\rA}
  \proofstepwidestar[3]{b\in A}{\rA}
  \proofstepwidestar[4]{a\mathrel{\mathcal L_{A,\star}}b}{\rA}
  \proofstepwidestar[1,2,3,4]{L_{A,\star}(a)=L_{A,\star}(b)}{%
    \FormulaRefAuto{SemigroupGreenLDef}[def]{1,2,3,4}}
  \proofstepwidestar[1,3]{b\in L_{A,\star}(b)\land b\in R_{A,\star}(b)\land b\in J_{A,\star}(b)}{%
      \FormulaRefAuto{SemigroupPrincipalIdealsContainGenerator}[thm]{1,3}}
  \proofstepwidestar[1,3]{b\in L_{A,\star}(b)}{%
    \rAEa{6}}
  \proofstepwidestar[1,2,3,4]{L_{A,\star}(b)=L_{A,\star}(a)}{%
      \FormulaRefAuto{a=b\vdash b=a}[thm]{5}}
  \proofstepwidestar[1,2,3,4]{b\in L_{A,\star}(a)}{%
    \rIE{8,7}}
  \proofstepwidestar[1,2,3,4]{\SemigroupRightDivides{A,\star}{a}{b}}{%
    \FormulaRefAuto{SemigroupRightDividesPrincipalLeftIdeal}[thm]{1,2,3,9}}
  \proofstepwidestar[1,2]{a\in L_{A,\star}(a)\land a\in R_{A,\star}(a)\land a\in J_{A,\star}(a)}{%
      \FormulaRefAuto{SemigroupPrincipalIdealsContainGenerator}[thm]{1,2}}
  \proofstepwidestar[1,2]{a\in L_{A,\star}(a)}{%
    \rAEa{11}}
  \proofstepwidestar[1,2,3,4]{a\in L_{A,\star}(b)}{\rIE{5,12}}
  \proofstepwidestar[1,2,3,4]{\SemigroupRightDivides{A,\star}{b}{a}}{%
    \FormulaRefAuto{SemigroupRightDividesPrincipalLeftIdeal}[thm]{1,3,2,13}}
  \proofstepwidestar[1,2,3,4]{%
    \SemigroupRightDivides{A,\star}{a}{b}\land
    \SemigroupRightDivides{A,\star}{b}{a}}{\rAI{10,14}}
  \proofstepwidestar[1,2,3]{%
    a\mathrel{\mathcal L_{A,\star}}b\rightarrow
    \bigl(\SemigroupRightDivides{A,\star}{a}{b}\land
          \SemigroupRightDivides{A,\star}{b}{a}\bigr)}{\rRI{4,15}}

\closeproofpartwide
\proofpartwide{Von gegenseitiger Teilbarkeit zu gleichen Hauptidealen}
\setcounter{proofstepnr}{16}
  \proofstepwidestar[17]{%
    \SemigroupRightDivides{A,\star}{a}{b}\land
    \SemigroupRightDivides{A,\star}{b}{a}}{\rA}
  \proofstepwidestar[17]{\SemigroupRightDivides{A,\star}{a}{b}}{\rAEa{17}}
  \proofstepwidestar[17]{\SemigroupRightDivides{A,\star}{b}{a}}{\rAEb{17}}
  \proofstepwidestar[1,2,3,17]{b\in L_{A,\star}(a)}{%
    \FormulaRefAuto{SemigroupRightDividesPrincipalLeftIdeal}[thm]{1,2,3,18}}
  \proofstepwidestar[1,2,3,17]{a\in L_{A,\star}(b)}{%
    \FormulaRefAuto{SemigroupRightDividesPrincipalLeftIdeal}[thm]{1,3,2,19}}
  \proofstepwidestar[1,2]{\mathsf{LId}(L_{A,\star}(a);A,\star)}{%
    \FormulaRefAuto{SemigroupPrincipalLeftIdealIsLeftIdeal}[thm]{1,2}}
  \proofstepwidestar[1,3]{\mathsf{LId}(L_{A,\star}(b);A,\star)}{%
    \FormulaRefAuto{SemigroupPrincipalLeftIdealIsLeftIdeal}[thm]{1,3}}
  \proofstepwidestar[1,2,3,17]{L_{A,\star}(a)\subseteq L_{A,\star}(b)}{%
    \FormulaRefAuto{SemigroupPrincipalLeftIdealMinimal}[thm]{1,2,23,21}}
  \proofstepwidestar[1,2,3,17]{L_{A,\star}(b)\subseteq L_{A,\star}(a)}{%
    \FormulaRefAuto{SemigroupPrincipalLeftIdealMinimal}[thm]{1,3,22,20}}
  \proofstepwidestar[1,2,3,17]{L_{A,\star}(a)=L_{A,\star}(b)}{%
    \FormulaRefAuto{A\subseteq B,\,B\subseteq A\vdash A=B}[thm]{24,25}}
  \proofstepwidestar[1,2,3,17]{a\mathrel{\mathcal L_{A,\star}}b}{%
    \FormulaRefAuto{SemigroupGreenLDef}[def]{1,2,3,26}}
  \proofstepwidestar[1,2,3]{%
    \bigl(\SemigroupRightDivides{A,\star}{a}{b}\land
          \SemigroupRightDivides{A,\star}{b}{a}\bigr)
    \rightarrow a\mathrel{\mathcal L_{A,\star}}b}{\rRI{17,27}}
  \proofstepwidestar[1,2,3]{%
    a\mathrel{\mathcal L_{A,\star}}b\leftrightarrow
    \bigl(\SemigroupRightDivides{A,\star}{a}{b}\land
          \SemigroupRightDivides{A,\star}{b}{a}\bigr)}{\rLRI{16,28}}
\closeproofpartwide
\end{tabproofsplitwide}

\FormulaThmDeltaK[\(\mathcal R\) als gegenseitige Linksteilbarkeit]{%
  \begin{aligned}[t]
    &\SemigroupAlg{A}{\star}\dsep a,b\in A\vdash{}\\[-2pt]
    &a\mathrel{\mathcal R_{A,\star}}b\leftrightarrow
      \Bigl(\SemigroupLeftDivides{A,\star}{a}{b}\land
            \SemigroupLeftDivides{A,\star}{b}{a}\Bigr)
  \end{aligned}
}{SemigroupGreenRMutualLeftDivisibility}{%
  \DeltaRow{Mengen}{A\dsep a\dsep b}
  \DeltaRow{Funktionen}{\star}
  \DeltaRow{Relationsoperatoren}{\mathcal R_{A,\star}}
}
\begin{tabproofsplitwide}
\proofpartwide{Von gleichen Hauptidealen zur Teilbarkeit}
  \proofstepwidestar[1]{\SemigroupAlg{A}{\star}}{\rA}
  \proofstepwidestar[2]{a\in A}{\rA}
  \proofstepwidestar[3]{b\in A}{\rA}
  \proofstepwidestar[4]{a\mathrel{\mathcal R_{A,\star}}b}{\rA}
  \proofstepwidestar[1,2,3,4]{R_{A,\star}(a)=R_{A,\star}(b)}{%
    \FormulaRefAuto{SemigroupGreenRDef}[def]{1,2,3,4}}
  \proofstepwidestar[1,3]{b\in L_{A,\star}(b)\land b\in R_{A,\star}(b)\land b\in J_{A,\star}(b)}{%
      \FormulaRefAuto{SemigroupPrincipalIdealsContainGenerator}[thm]{1,3}}
  \proofstepwidestar[1,3]{b\in R_{A,\star}(b)}{%
    \rAEa{\rAEb{6}}}
  \proofstepwidestar[1,2,3,4]{R_{A,\star}(b)=R_{A,\star}(a)}{%
      \FormulaRefAuto{a=b\vdash b=a}[thm]{5}}
  \proofstepwidestar[1,2,3,4]{b\in R_{A,\star}(a)}{%
    \rIE{8,7}}
  \proofstepwidestar[1,2,3,4]{\SemigroupLeftDivides{A,\star}{a}{b}}{%
    \FormulaRefAuto{SemigroupLeftDividesPrincipalRightIdeal}[thm]{1,2,3,9}}
  \proofstepwidestar[1,2]{a\in L_{A,\star}(a)\land a\in R_{A,\star}(a)\land a\in J_{A,\star}(a)}{%
      \FormulaRefAuto{SemigroupPrincipalIdealsContainGenerator}[thm]{1,2}}
  \proofstepwidestar[1,2]{a\in R_{A,\star}(a)}{%
    \rAEa{\rAEb{11}}}
  \proofstepwidestar[1,2,3,4]{a\in R_{A,\star}(b)}{\rIE{5,12}}
  \proofstepwidestar[1,2,3,4]{\SemigroupLeftDivides{A,\star}{b}{a}}{%
    \FormulaRefAuto{SemigroupLeftDividesPrincipalRightIdeal}[thm]{1,3,2,13}}
  \proofstepwidestar[1,2,3,4]{%
    \SemigroupLeftDivides{A,\star}{a}{b}\land
    \SemigroupLeftDivides{A,\star}{b}{a}}{\rAI{10,14}}
  \proofstepwidestar[1,2,3]{%
    a\mathrel{\mathcal R_{A,\star}}b\rightarrow
    \bigl(\SemigroupLeftDivides{A,\star}{a}{b}\land
          \SemigroupLeftDivides{A,\star}{b}{a}\bigr)}{\rRI{4,15}}

\closeproofpartwide
\proofpartwide{Von gegenseitiger Teilbarkeit zu gleichen Hauptidealen}
\setcounter{proofstepnr}{16}
  \proofstepwidestar[17]{%
    \SemigroupLeftDivides{A,\star}{a}{b}\land
    \SemigroupLeftDivides{A,\star}{b}{a}}{\rA}
  \proofstepwidestar[17]{\SemigroupLeftDivides{A,\star}{a}{b}}{\rAEa{17}}
  \proofstepwidestar[17]{\SemigroupLeftDivides{A,\star}{b}{a}}{\rAEb{17}}
  \proofstepwidestar[1,2,3,17]{b\in R_{A,\star}(a)}{%
    \FormulaRefAuto{SemigroupLeftDividesPrincipalRightIdeal}[thm]{1,2,3,18}}
  \proofstepwidestar[1,2,3,17]{a\in R_{A,\star}(b)}{%
    \FormulaRefAuto{SemigroupLeftDividesPrincipalRightIdeal}[thm]{1,3,2,19}}
  \proofstepwidestar[1,2]{\mathsf{RId}(R_{A,\star}(a);A,\star)}{%
    \FormulaRefAuto{SemigroupPrincipalRightIdealIsRightIdeal}[thm]{1,2}}
  \proofstepwidestar[1,3]{\mathsf{RId}(R_{A,\star}(b);A,\star)}{%
    \FormulaRefAuto{SemigroupPrincipalRightIdealIsRightIdeal}[thm]{1,3}}
  \proofstepwidestar[1,2,3,17]{R_{A,\star}(a)\subseteq R_{A,\star}(b)}{%
    \FormulaRefAuto{SemigroupPrincipalRightIdealMinimal}[thm]{1,2,23,21}}
  \proofstepwidestar[1,2,3,17]{R_{A,\star}(b)\subseteq R_{A,\star}(a)}{%
    \FormulaRefAuto{SemigroupPrincipalRightIdealMinimal}[thm]{1,3,22,20}}
  \proofstepwidestar[1,2,3,17]{R_{A,\star}(a)=R_{A,\star}(b)}{%
    \FormulaRefAuto{A\subseteq B,\,B\subseteq A\vdash A=B}[thm]{24,25}}
  \proofstepwidestar[1,2,3,17]{a\mathrel{\mathcal R_{A,\star}}b}{%
    \FormulaRefAuto{SemigroupGreenRDef}[def]{1,2,3,26}}
  \proofstepwidestar[1,2,3]{%
    \bigl(\SemigroupLeftDivides{A,\star}{a}{b}\land
          \SemigroupLeftDivides{A,\star}{b}{a}\bigr)
    \rightarrow a\mathrel{\mathcal R_{A,\star}}b}{\rRI{17,27}}
  \proofstepwidestar[1,2,3]{%
    a\mathrel{\mathcal R_{A,\star}}b\leftrightarrow
    \bigl(\SemigroupLeftDivides{A,\star}{a}{b}\land
          \SemigroupLeftDivides{A,\star}{b}{a}\bigr)}{\rLRI{16,28}}
\closeproofpartwide
\end{tabproofsplitwide}

\FormulaThmDeltaK[\(\mathcal J\) als gegenseitige zweiseitige Teilbarkeit]{%
  \begin{aligned}[t]
    &\SemigroupAlg{A}{\star}\dsep a,b\in A\vdash{}\\[-2pt]
    &a\mathrel{\mathcal J_{A,\star}}b\leftrightarrow
      \Bigl(\SemigroupDivides{A,\star}{a}{b}\land
            \SemigroupDivides{A,\star}{b}{a}\Bigr)
  \end{aligned}
}{SemigroupGreenJMutualTwoSidedDivisibility}{%
  \DeltaRow{Mengen}{A\dsep a\dsep b}
  \DeltaRow{Funktionen}{\star}
  \DeltaRow{Relationsoperatoren}{\mathcal J_{A,\star}}
}
\begin{tabproofsplitwide}
\proofpartwide{Von gleichen Hauptidealen zur Teilbarkeit}
  \proofstepwidestar[1]{\SemigroupAlg{A}{\star}}{\rA}
  \proofstepwidestar[2]{a\in A}{\rA}
  \proofstepwidestar[3]{b\in A}{\rA}
  \proofstepwidestar[4]{a\mathrel{\mathcal J_{A,\star}}b}{\rA}
  \proofstepwidestar[1,2,3,4]{J_{A,\star}(a)=J_{A,\star}(b)}{%
    \FormulaRefAuto{SemigroupGreenJDef}[def]{1,2,3,4}}
  \proofstepwidestar[1,3]{b\in L_{A,\star}(b)\land b\in R_{A,\star}(b)\land b\in J_{A,\star}(b)}{%
      \FormulaRefAuto{SemigroupPrincipalIdealsContainGenerator}[thm]{1,3}}
  \proofstepwidestar[1,3]{b\in J_{A,\star}(b)}{%
    \rAEb{\rAEb{6}}}
  \proofstepwidestar[1,2,3,4]{J_{A,\star}(b)=J_{A,\star}(a)}{%
      \FormulaRefAuto{a=b\vdash b=a}[thm]{5}}
  \proofstepwidestar[1,2,3,4]{b\in J_{A,\star}(a)}{%
    \rIE{8,7}}
  \proofstepwidestar[1,2,3,4]{\SemigroupDivides{A,\star}{a}{b}}{%
    \FormulaRefAuto{SemigroupTwoSidedDividesPrincipalIdeal}[thm]{1,2,3,9}}
  \proofstepwidestar[1,2]{a\in L_{A,\star}(a)\land a\in R_{A,\star}(a)\land a\in J_{A,\star}(a)}{%
      \FormulaRefAuto{SemigroupPrincipalIdealsContainGenerator}[thm]{1,2}}
  \proofstepwidestar[1,2]{a\in J_{A,\star}(a)}{%
    \rAEb{\rAEb{11}}}
  \proofstepwidestar[1,2,3,4]{a\in J_{A,\star}(b)}{\rIE{5,12}}
  \proofstepwidestar[1,2,3,4]{\SemigroupDivides{A,\star}{b}{a}}{%
    \FormulaRefAuto{SemigroupTwoSidedDividesPrincipalIdeal}[thm]{1,3,2,13}}
  \proofstepwidestar[1,2,3,4]{%
    \SemigroupDivides{A,\star}{a}{b}\land
    \SemigroupDivides{A,\star}{b}{a}}{\rAI{10,14}}
  \proofstepwidestar[1,2,3]{%
    a\mathrel{\mathcal J_{A,\star}}b\rightarrow
    \bigl(\SemigroupDivides{A,\star}{a}{b}\land
          \SemigroupDivides{A,\star}{b}{a}\bigr)}{\rRI{4,15}}

\closeproofpartwide
\proofpartwide{Von gegenseitiger Teilbarkeit zu gleichen Hauptidealen}
\setcounter{proofstepnr}{16}
  \proofstepwidestar[17]{%
    \SemigroupDivides{A,\star}{a}{b}\land
    \SemigroupDivides{A,\star}{b}{a}}{\rA}
  \proofstepwidestar[17]{\SemigroupDivides{A,\star}{a}{b}}{\rAEa{17}}
  \proofstepwidestar[17]{\SemigroupDivides{A,\star}{b}{a}}{\rAEb{17}}
  \proofstepwidestar[1,2,3,17]{b\in J_{A,\star}(a)}{%
    \FormulaRefAuto{SemigroupTwoSidedDividesPrincipalIdeal}[thm]{1,2,3,18}}
  \proofstepwidestar[1,2,3,17]{a\in J_{A,\star}(b)}{%
    \FormulaRefAuto{SemigroupTwoSidedDividesPrincipalIdeal}[thm]{1,3,2,19}}
  \proofstepwidestar[1,2]{\mathsf{Id}(J_{A,\star}(a);A,\star)}{%
    \FormulaRefAuto{SemigroupPrincipalTwoSidedIdealIsIdeal}[thm]{1,2}}
  \proofstepwidestar[1,3]{\mathsf{Id}(J_{A,\star}(b);A,\star)}{%
    \FormulaRefAuto{SemigroupPrincipalTwoSidedIdealIsIdeal}[thm]{1,3}}
  \proofstepwidestar[1,2,3,17]{J_{A,\star}(a)\subseteq J_{A,\star}(b)}{%
    \FormulaRefAuto{SemigroupPrincipalTwoSidedIdealMinimal}[thm]{1,2,23,21}}
  \proofstepwidestar[1,2,3,17]{J_{A,\star}(b)\subseteq J_{A,\star}(a)}{%
    \FormulaRefAuto{SemigroupPrincipalTwoSidedIdealMinimal}[thm]{1,3,22,20}}
  \proofstepwidestar[1,2,3,17]{J_{A,\star}(a)=J_{A,\star}(b)}{%
    \FormulaRefAuto{A\subseteq B,\,B\subseteq A\vdash A=B}[thm]{24,25}}
  \proofstepwidestar[1,2,3,17]{a\mathrel{\mathcal J_{A,\star}}b}{%
    \FormulaRefAuto{SemigroupGreenJDef}[def]{1,2,3,26}}
  \proofstepwidestar[1,2,3]{%
    \bigl(\SemigroupDivides{A,\star}{a}{b}\land
          \SemigroupDivides{A,\star}{b}{a}\bigr)
    \rightarrow a\mathrel{\mathcal J_{A,\star}}b}{\rRI{17,27}}
  \proofstepwidestar[1,2,3]{%
    a\mathrel{\mathcal J_{A,\star}}b\leftrightarrow
    \bigl(\SemigroupDivides{A,\star}{a}{b}\land
          \SemigroupDivides{A,\star}{b}{a}\bigr)}{\rLRI{16,28}}
\closeproofpartwide
\end{tabproofsplitwide}

\FormulaThmDeltaK[\(\mathcal H\) als gegenseitige einseitige Teilbarkeit]{%
  \begin{aligned}[t]
    &\SemigroupAlg{A}{\star}\dsep a,b\in A\vdash{}\\[-2pt]
    &a\mathrel{\mathcal H_{A,\star}}b\leftrightarrow
      \Bigl(
        \SemigroupRightDivides{A,\star}{a}{b}\land
        \SemigroupRightDivides{A,\star}{b}{a}\land
        \SemigroupLeftDivides{A,\star}{a}{b}\land
        \SemigroupLeftDivides{A,\star}{b}{a}\Bigr)
  \end{aligned}
}{SemigroupGreenHMutualOneSidedDivisibility}{%
  \DeltaRow{Mengen}{A\dsep a\dsep b}
  \DeltaRow{Funktionen}{\star}
  \DeltaRow{Relationsoperatoren}{\mathcal H_{A,\star}}
}
Die Definition von \(\mathcal H\) zerf\"allt in eine \(\mathcal L\)- und
eine \(\mathcal R\)-Bedingung.  Die beiden Richtungen bestehen daher nur
darin, diese Bedingungen in die zugeh\"origen Teilbarkeitsaussagen zu
\"ubersetzen beziehungsweise wieder zusammenzusetzen.

\begin{tabproofsplitwide}
  \proofpartwideindR[Hinrichtung]{%
    \begin{aligned}[t]
      &\SemigroupAlg{A}{\star}\dsep a,b\in A\dsep
        a\mathrel{\mathcal H_{A,\star}}b\vdash{}\\[-2pt]
      &\SemigroupRightDivides{A,\star}{a}{b}\land
        \SemigroupRightDivides{A,\star}{b}{a}\land
        \SemigroupLeftDivides{A,\star}{a}{b}\land
        \SemigroupLeftDivides{A,\star}{b}{a}
    \end{aligned}
  }{%
    \begin{aligned}[t]
      &\SemigroupAlg{A}{\star}\dsep a,b\in A\dsep
        a\mathrel{\mathcal H_{A,\star}}b\vdash{}\\[-2pt]
      &\SemigroupRightDivides{A,\star}{a}{b}\land
        \SemigroupRightDivides{A,\star}{b}{a}\land
        \SemigroupLeftDivides{A,\star}{a}{b}\land
        \SemigroupLeftDivides{A,\star}{b}{a}
    \end{aligned}
  }[SemigroupGreenHMutualOneSidedDivisibilityForward]
    \proofstepwidestar[1]{\SemigroupAlg{A}{\star}}{\rA}
    \proofstepwidestar[2]{a\in A}{\rA}
    \proofstepwidestar[3]{b\in A}{\rA}
    \proofstepwidestar[4]{a\mathrel{\mathcal H_{A,\star}}b}{\rA}
    \proofstepwidestar[1,2,3,4]{%
      a\mathrel{\mathcal L_{A,\star}}b\land
      a\mathrel{\mathcal R_{A,\star}}b}{%
      \FormulaRefAuto{SemigroupGreenHDef}[def]{1,2,3,4}}
    \proofstepwidestar[1,2,3,4]{%
      \SemigroupRightDivides{A,\star}{a}{b}\land
      \SemigroupRightDivides{A,\star}{b}{a}}{%
      \FormulaRefAuto{SemigroupGreenLMutualRightDivisibility}[thm]{%
        1,2,3,\rAEa{5}}}
    \proofstepwidestar[1,2,3,4]{%
      \SemigroupLeftDivides{A,\star}{a}{b}\land
      \SemigroupLeftDivides{A,\star}{b}{a}}{%
      \FormulaRefAuto{SemigroupGreenRMutualLeftDivisibility}[thm]{%
        1,2,3,\rAEb{5}}}
    \proofstepwidestarbreak[1,2,3,4]{%
      \begin{aligned}[t]
        &\SemigroupRightDivides{A,\star}{a}{b}\land
          \SemigroupRightDivides{A,\star}{b}{a}\land{}\\[-2pt]
        &\SemigroupLeftDivides{A,\star}{a}{b}\land
          \SemigroupLeftDivides{A,\star}{b}{a}
      \end{aligned}}{%
      \rAI{\rAEa{6},\rAI{\rAEb{6},\rAI{\rAEa{7},\rAEb{7}}}}}
  \closeproofpartwideindR

  \proofpartwideindR[R\"uckrichtung]{%
    \begin{aligned}[t]
      &\SemigroupAlg{A}{\star}\dsep a,b\in A\dsep{}\\[-2pt]
      &\SemigroupRightDivides{A,\star}{a}{b}\land
        \SemigroupRightDivides{A,\star}{b}{a}\land
        \SemigroupLeftDivides{A,\star}{a}{b}\land
        \SemigroupLeftDivides{A,\star}{b}{a}
        \vdash a\mathrel{\mathcal H_{A,\star}}b
    \end{aligned}
  }{%
    \begin{aligned}[t]
      &\SemigroupAlg{A}{\star}\dsep a,b\in A\dsep{}\\[-2pt]
      &\SemigroupRightDivides{A,\star}{a}{b}\land
        \SemigroupRightDivides{A,\star}{b}{a}\land
        \SemigroupLeftDivides{A,\star}{a}{b}\land
        \SemigroupLeftDivides{A,\star}{b}{a}
        \vdash a\mathrel{\mathcal H_{A,\star}}b
    \end{aligned}
  }[SemigroupGreenHMutualOneSidedDivisibilityBackward]
    \proofstepwidestar[1]{\SemigroupAlg{A}{\star}}{\rA}
    \proofstepwidestar[2]{a\in A}{\rA}
    \proofstepwidestar[3]{b\in A}{\rA}
    \proofstepwidestar[4]{%
      \begin{aligned}[t]
        &\SemigroupRightDivides{A,\star}{a}{b}\land
          \SemigroupRightDivides{A,\star}{b}{a}\land{}\\[-2pt]
        &\SemigroupLeftDivides{A,\star}{a}{b}\land
          \SemigroupLeftDivides{A,\star}{b}{a}
      \end{aligned}}{\rA}
    \proofstepwidestar[4]{%
      \SemigroupRightDivides{A,\star}{a}{b}\land
      \SemigroupRightDivides{A,\star}{b}{a}}{%
      \rAI{\rAEa{4},\rAEa{\rAEb{4}}}}
    \proofstepwidestar[4]{%
      \SemigroupLeftDivides{A,\star}{a}{b}\land
      \SemigroupLeftDivides{A,\star}{b}{a}}{%
      \rAI{\rAEa{\rAEb{\rAEb{4}}},\rAEb{\rAEb{\rAEb{4}}}}}
    \proofstepwidestar[1,2,3,4]{a\mathrel{\mathcal L_{A,\star}}b}{%
      \FormulaRefAuto{SemigroupGreenLMutualRightDivisibility}[thm]{1,2,3,5}}
    \proofstepwidestar[1,2,3,4]{a\mathrel{\mathcal R_{A,\star}}b}{%
      \FormulaRefAuto{SemigroupGreenRMutualLeftDivisibility}[thm]{1,2,3,6}}
    \proofstepwidestar[1,2,3,4]{a\mathrel{\mathcal H_{A,\star}}b}{%
      \FormulaRefAuto{SemigroupGreenHDef}[def]{1,2,3,\rAI{7,8}}}
  \closeproofpartwideindR

  \proofpartwide{Schluss}
    \proofstepwidestar[1]{\SemigroupAlg{A}{\star}}{\rA}
    \proofstepwidestar[2]{a\in A}{\rA}
    \proofstepwidestar[3]{b\in A}{\rA}
    \proofstepwidestar[4]{a\mathrel{\mathcal H_{A,\star}}b}{\rA}
    \proofstepwidestarbreak[1,2,3,4]{%
      \begin{aligned}[t]
        &\SemigroupRightDivides{A,\star}{a}{b}\land
          \SemigroupRightDivides{A,\star}{b}{a}\land{}\\[-2pt]
        &\SemigroupLeftDivides{A,\star}{a}{b}\land
          \SemigroupLeftDivides{A,\star}{b}{a}
      \end{aligned}}{%
      \FormulaRefAuto{SemigroupGreenHMutualOneSidedDivisibilityForward}[thm]{%
        1,2,3,4}}
    \proofstepwidestar[1,2,3]{%
      \begin{aligned}[t]
        a\mathrel{\mathcal H_{A,\star}}b\rightarrow\bigl(&
          \SemigroupRightDivides{A,\star}{a}{b}\land
          \SemigroupRightDivides{A,\star}{b}{a}\land{}\\[-2pt]
        &\SemigroupLeftDivides{A,\star}{a}{b}\land
          \SemigroupLeftDivides{A,\star}{b}{a}\bigr)
      \end{aligned}}{\rRI{4,5}}
    \proofstepwidestar[7]{%
      \begin{aligned}[t]
        &\SemigroupRightDivides{A,\star}{a}{b}\land
          \SemigroupRightDivides{A,\star}{b}{a}\land{}\\[-2pt]
        &\SemigroupLeftDivides{A,\star}{a}{b}\land
          \SemigroupLeftDivides{A,\star}{b}{a}
      \end{aligned}}{\rA}
    \proofstepwidestar[1,2,3,7]{a\mathrel{\mathcal H_{A,\star}}b}{%
      \FormulaRefAuto{SemigroupGreenHMutualOneSidedDivisibilityBackward}[thm]{%
        1,2,3,7}}
    \proofstepwidestar[1,2,3]{%
      \begin{aligned}[t]
        \bigl(&\SemigroupRightDivides{A,\star}{a}{b}\land
          \SemigroupRightDivides{A,\star}{b}{a}\land{}\\[-2pt]
        &\SemigroupLeftDivides{A,\star}{a}{b}\land
          \SemigroupLeftDivides{A,\star}{b}{a}\bigr)
          \rightarrow a\mathrel{\mathcal H_{A,\star}}b
      \end{aligned}}{\rRI{7,8}}
    \proofstepwidestarbreak[1,2,3]{%
      \begin{aligned}[t]
        a\mathrel{\mathcal H_{A,\star}}b\leftrightarrow\bigl(&
          \SemigroupRightDivides{A,\star}{a}{b}\land
          \SemigroupRightDivides{A,\star}{b}{a}\land{}\\[-2pt]
        &\SemigroupLeftDivides{A,\star}{a}{b}\land
          \SemigroupLeftDivides{A,\star}{b}{a}\bigr)
      \end{aligned}}{\rLRI{6,9}}
  \closeproofpartwide
\end{tabproofsplitwide}

\subsection{Äquivalenzeigenschaften}

\FormulaThmDeltaK[Reflexivität von \(\mathcal L\)]{%
  \SemigroupAlg{A}{\star}\dsep a\in A\vdash
  a\mathrel{\mathcal L_{A,\star}}a%
}{SemigroupGreenLReflexive}{%
  \DeltaRow{Mengen}{A\dsep a}
  \DeltaRow{Funktionen}{\star}
  \DeltaRow{Relationsoperatoren}{\mathcal L_{A,\star}}
}
\begin{tabproofwide}
  \proofstepwidestar[1]{\SemigroupAlg{A}{\star}}{\rA}
  \proofstepwidestar[2]{a\in A}{\rA}
  \proofstepwidestar[]{L_{A,\star}(a)=L_{A,\star}(a)}{\rII}
  \proofstepwidestar[1,2]{a\mathrel{\mathcal L_{A,\star}}a}{%
    \FormulaRefAuto{SemigroupGreenLDef}[def]{1,2,2,3}}
\end{tabproofwide}

\FormulaThmDeltaK[Symmetrie von \(\mathcal L\)]{%
  \begin{aligned}[t]
    &\SemigroupAlg{A}{\star}\dsep a,b\in A\dsep
      a\mathrel{\mathcal L_{A,\star}}b\vdash{}\\[-2pt]
    &b\mathrel{\mathcal L_{A,\star}}a
  \end{aligned}
}{SemigroupGreenLSymmetric}{%
  \DeltaRow{Mengen}{A\dsep a\dsep b}
  \DeltaRow{Funktionen}{\star}
  \DeltaRow{Relationsoperatoren}{\mathcal L_{A,\star}}
}
\begin{tabproofwide}
  \proofstepwidestar[1]{\SemigroupAlg{A}{\star}}{\rA}
  \proofstepwidestar[2]{a\in A}{\rA}
  \proofstepwidestar[3]{b\in A}{\rA}
  \proofstepwidestar[4]{a\mathrel{\mathcal L_{A,\star}}b}{\rA}
  \proofstepwidestar[1,2,3,4]{L_{A,\star}(a)=L_{A,\star}(b)}{%
    \FormulaRefAuto{SemigroupGreenLDef}[def]{1,2,3,4}}
  \proofstepwidestar[1,2,3,4]{L_{A,\star}(b)=L_{A,\star}(a)}{%
    \FormulaRefAuto{a=b\vdash b=a}[thm]{5}}
  \proofstepwidestar[1,2,3,4]{b\mathrel{\mathcal L_{A,\star}}a}{%
    \FormulaRefAuto{SemigroupGreenLDef}[def]{1,3,2,6}}
\end{tabproofwide}

\FormulaThmDeltaK[Transitivität von \(\mathcal L\)]{%
  \begin{aligned}[t]
    &\SemigroupAlg{A}{\star}\dsep a,b,c\in A\dsep
      a\mathrel{\mathcal L_{A,\star}}b\dsep
      b\mathrel{\mathcal L_{A,\star}}c\vdash{}\\[-2pt]
    &a\mathrel{\mathcal L_{A,\star}}c
  \end{aligned}
}{SemigroupGreenLTransitive}{%
  \DeltaRow{Mengen}{A\dsep a\dsep b\dsep c}
  \DeltaRow{Funktionen}{\star}
  \DeltaRow{Relationsoperatoren}{\mathcal L_{A,\star}}
}
\begin{tabproofwide}
  \proofstepwidestar[1]{\SemigroupAlg{A}{\star}}{\rA}
  \proofstepwidestar[2]{a\in A}{\rA}
  \proofstepwidestar[3]{b\in A}{\rA}
  \proofstepwidestar[4]{c\in A}{\rA}
  \proofstepwidestar[5]{a\mathrel{\mathcal L_{A,\star}}b}{\rA}
  \proofstepwidestar[6]{b\mathrel{\mathcal L_{A,\star}}c}{\rA}
  \proofstepwidestar[1,2,3,5]{L_{A,\star}(a)=L_{A,\star}(b)}{%
    \FormulaRefAuto{SemigroupGreenLDef}[def]{1,2,3,5}}
  \proofstepwidestar[1,3,4,6]{L_{A,\star}(b)=L_{A,\star}(c)}{%
    \FormulaRefAuto{SemigroupGreenLDef}[def]{1,3,4,6}}
  \proofstepwidestar[1,2,3,4,5,6]{L_{A,\star}(a)=L_{A,\star}(c)}{%
    \FormulaRefAuto{a=b\dsep b=c\vdash a=c}[thm]{7,8}}
  \proofstepwidestar[1,2,3,4,5,6]{a\mathrel{\mathcal L_{A,\star}}c}{%
    \FormulaRefAuto{SemigroupGreenLDef}[def]{1,2,4,9}}
\end{tabproofwide}

\FormulaThmDeltaK[\(\mathcal L\) ist eine Äquivalenzrelation]{%
  \SemigroupAlg{A}{\star}\vdash\EqRel{A,\mathcal L_{A,\star}}%
}{SemigroupGreenLIsEquivalence}{%
  \DeltaRow{Mengen}{A\dsep a\dsep b\dsep c}
  \DeltaRow{Funktionen}{\star}
  \DeltaRow{Relationsoperatoren}{\mathcal L_{A,\star}}
}
\begin{tabproofwide}
  \proofstepwidestar[1]{\SemigroupAlg{A}{\star}}{\rA}
  \proofstepwidestar[1]{%
    \forall a\in A\,
      \bigl(a\mathrel{\mathcal L_{A,\star}}a\bigr)}{%
    \FormulaRefAuto{SemigroupGreenLReflexive}[thm]{1}}
  \proofstepwidestar[1]{%
    \forall a,b\in A\,
      \bigl(a\mathrel{\mathcal L_{A,\star}}b\rightarrow
        b\mathrel{\mathcal L_{A,\star}}a\bigr)}{%
    \FormulaRefAuto{SemigroupGreenLSymmetric}[thm]{1}}
  \proofstepwidestar[1]{%
    \forall a,b,c\in A\,
      \Bigl(\bigl(a\mathrel{\mathcal L_{A,\star}}b\land
        b\mathrel{\mathcal L_{A,\star}}c\bigr)
        \rightarrow a\mathrel{\mathcal L_{A,\star}}c\Bigr)}{%
    \FormulaRefAuto{SemigroupGreenLTransitive}[thm]{1}}
  \proofstepwidestar[1]{\EqRel{A,\mathcal L_{A,\star}}}{%
    \FormulaRefAuto{EqRelDef}[def]{2,3,4}}
\end{tabproofwide}

\FormulaThmDeltaK[Reflexivität von \(\mathcal R\)]{%
  \SemigroupAlg{A}{\star}\dsep a\in A\vdash
  a\mathrel{\mathcal R_{A,\star}}a%
}{SemigroupGreenRReflexive}{%
  \DeltaRow{Mengen}{A\dsep a}
  \DeltaRow{Funktionen}{\star}
  \DeltaRow{Relationsoperatoren}{\mathcal R_{A,\star}}
}
\begin{tabproofwide}
  \proofstepwidestar[1]{\SemigroupAlg{A}{\star}}{\rA}
  \proofstepwidestar[2]{a\in A}{\rA}
  \proofstepwidestar[]{R_{A,\star}(a)=R_{A,\star}(a)}{\rII}
  \proofstepwidestar[1,2]{a\mathrel{\mathcal R_{A,\star}}a}{%
    \FormulaRefAuto{SemigroupGreenRDef}[def]{1,2,2,3}}
\end{tabproofwide}

\FormulaThmDeltaK[Symmetrie von \(\mathcal R\)]{%
  \SemigroupAlg{A}{\star}\dsep a,b\in A\dsep
  a\mathrel{\mathcal R_{A,\star}}b\vdash
  b\mathrel{\mathcal R_{A,\star}}a%
}{SemigroupGreenRSymmetric}{%
  \DeltaRow{Mengen}{A\dsep a\dsep b}
  \DeltaRow{Funktionen}{\star}
  \DeltaRow{Relationsoperatoren}{\mathcal R_{A,\star}}
}
\begin{tabproofwide}
  \proofstepwidestar[1]{\SemigroupAlg{A}{\star}}{\rA}
  \proofstepwidestar[2]{a\in A}{\rA}
  \proofstepwidestar[3]{b\in A}{\rA}
  \proofstepwidestar[4]{a\mathrel{\mathcal R_{A,\star}}b}{\rA}
  \proofstepwidestar[1,2,3,4]{R_{A,\star}(a)=R_{A,\star}(b)}{%
    \FormulaRefAuto{SemigroupGreenRDef}[def]{1,2,3,4}}
  \proofstepwidestar[1,2,3,4]{R_{A,\star}(b)=R_{A,\star}(a)}{%
    \FormulaRefAuto{a=b\vdash b=a}[thm]{5}}
  \proofstepwidestar[1,2,3,4]{b\mathrel{\mathcal R_{A,\star}}a}{%
    \FormulaRefAuto{SemigroupGreenRDef}[def]{1,3,2,6}}
\end{tabproofwide}

\FormulaThmDeltaK[Transitivität von \(\mathcal R\)]{%
  \begin{aligned}[t]
    &\SemigroupAlg{A}{\star}\dsep a,b,c\in A\dsep
      a\mathrel{\mathcal R_{A,\star}}b\dsep
      b\mathrel{\mathcal R_{A,\star}}c\vdash{}\\[-2pt]
    &a\mathrel{\mathcal R_{A,\star}}c
  \end{aligned}
}{SemigroupGreenRTransitive}{%
  \DeltaRow{Mengen}{A\dsep a\dsep b\dsep c}
  \DeltaRow{Funktionen}{\star}
  \DeltaRow{Relationsoperatoren}{\mathcal R_{A,\star}}
}
\begin{tabproofwide}
  \proofstepwidestar[1]{\SemigroupAlg{A}{\star}}{\rA}
  \proofstepwidestar[2]{a\in A}{\rA}
  \proofstepwidestar[3]{b\in A}{\rA}
  \proofstepwidestar[4]{c\in A}{\rA}
  \proofstepwidestar[5]{a\mathrel{\mathcal R_{A,\star}}b}{\rA}
  \proofstepwidestar[6]{b\mathrel{\mathcal R_{A,\star}}c}{\rA}
  \proofstepwidestar[1,2,3,5]{R_{A,\star}(a)=R_{A,\star}(b)}{%
    \FormulaRefAuto{SemigroupGreenRDef}[def]{1,2,3,5}}
  \proofstepwidestar[1,3,4,6]{R_{A,\star}(b)=R_{A,\star}(c)}{%
    \FormulaRefAuto{SemigroupGreenRDef}[def]{1,3,4,6}}
  \proofstepwidestar[1,2,3,4,5,6]{R_{A,\star}(a)=R_{A,\star}(c)}{%
    \FormulaRefAuto{a=b\dsep b=c\vdash a=c}[thm]{7,8}}
  \proofstepwidestar[1,2,3,4,5,6]{a\mathrel{\mathcal R_{A,\star}}c}{%
    \FormulaRefAuto{SemigroupGreenRDef}[def]{1,2,4,9}}
\end{tabproofwide}

\FormulaThmDeltaK[\(\mathcal R\) ist eine Äquivalenzrelation]{%
  \SemigroupAlg{A}{\star}\vdash\EqRel{A,\mathcal R_{A,\star}}%
}{SemigroupGreenRIsEquivalence}{%
  \DeltaRow{Mengen}{A\dsep a\dsep b\dsep c}
  \DeltaRow{Funktionen}{\star}
  \DeltaRow{Relationsoperatoren}{\mathcal R_{A,\star}}
}
\begin{tabproofwide}
  \proofstepwidestar[1]{\SemigroupAlg{A}{\star}}{\rA}
  \proofstepwidestar[1]{%
    \forall a\in A\,
      \bigl(a\mathrel{\mathcal R_{A,\star}}a\bigr)}{%
    \FormulaRefAuto{SemigroupGreenRReflexive}[thm]{1}}
  \proofstepwidestar[1]{%
    \forall a,b\in A\,
      \bigl(a\mathrel{\mathcal R_{A,\star}}b\rightarrow
        b\mathrel{\mathcal R_{A,\star}}a\bigr)}{%
    \FormulaRefAuto{SemigroupGreenRSymmetric}[thm]{1}}
  \proofstepwidestar[1]{%
    \forall a,b,c\in A\,
      \Bigl(\bigl(a\mathrel{\mathcal R_{A,\star}}b\land
        b\mathrel{\mathcal R_{A,\star}}c\bigr)
        \rightarrow a\mathrel{\mathcal R_{A,\star}}c\Bigr)}{%
    \FormulaRefAuto{SemigroupGreenRTransitive}[thm]{1}}
  \proofstepwidestar[1]{\EqRel{A,\mathcal R_{A,\star}}}{%
    \FormulaRefAuto{EqRelDef}[def]{2,3,4}}
\end{tabproofwide}

\FormulaThmDeltaK[Reflexivität von \(\mathcal J\)]{%
  \SemigroupAlg{A}{\star}\dsep a\in A\vdash
  a\mathrel{\mathcal J_{A,\star}}a%
}{SemigroupGreenJReflexive}{%
  \DeltaRow{Mengen}{A\dsep a}
  \DeltaRow{Funktionen}{\star}
  \DeltaRow{Relationsoperatoren}{\mathcal J_{A,\star}}
}
\begin{tabproofwide}
  \proofstepwidestar[1]{\SemigroupAlg{A}{\star}}{\rA}
  \proofstepwidestar[2]{a\in A}{\rA}
  \proofstepwidestar[]{J_{A,\star}(a)=J_{A,\star}(a)}{\rII}
  \proofstepwidestar[1,2]{a\mathrel{\mathcal J_{A,\star}}a}{%
    \FormulaRefAuto{SemigroupGreenJDef}[def]{1,2,2,3}}
\end{tabproofwide}

\FormulaThmDeltaK[Symmetrie von \(\mathcal J\)]{%
  \SemigroupAlg{A}{\star}\dsep a,b\in A\dsep
  a\mathrel{\mathcal J_{A,\star}}b\vdash
  b\mathrel{\mathcal J_{A,\star}}a%
}{SemigroupGreenJSymmetric}{%
  \DeltaRow{Mengen}{A\dsep a\dsep b}
  \DeltaRow{Funktionen}{\star}
  \DeltaRow{Relationsoperatoren}{\mathcal J_{A,\star}}
}
\begin{tabproofwide}
  \proofstepwidestar[1]{\SemigroupAlg{A}{\star}}{\rA}
  \proofstepwidestar[2]{a\in A}{\rA}
  \proofstepwidestar[3]{b\in A}{\rA}
  \proofstepwidestar[4]{a\mathrel{\mathcal J_{A,\star}}b}{\rA}
  \proofstepwidestar[1,2,3,4]{J_{A,\star}(a)=J_{A,\star}(b)}{%
    \FormulaRefAuto{SemigroupGreenJDef}[def]{1,2,3,4}}
  \proofstepwidestar[1,2,3,4]{J_{A,\star}(b)=J_{A,\star}(a)}{%
    \FormulaRefAuto{a=b\vdash b=a}[thm]{5}}
  \proofstepwidestar[1,2,3,4]{b\mathrel{\mathcal J_{A,\star}}a}{%
    \FormulaRefAuto{SemigroupGreenJDef}[def]{1,3,2,6}}
\end{tabproofwide}

\FormulaThmDeltaK[Transitivität von \(\mathcal J\)]{%
  \begin{aligned}[t]
    &\SemigroupAlg{A}{\star}\dsep a,b,c\in A\dsep
      a\mathrel{\mathcal J_{A,\star}}b\dsep
      b\mathrel{\mathcal J_{A,\star}}c\vdash{}\\[-2pt]
    &a\mathrel{\mathcal J_{A,\star}}c
  \end{aligned}
}{SemigroupGreenJTransitive}{%
  \DeltaRow{Mengen}{A\dsep a\dsep b\dsep c}
  \DeltaRow{Funktionen}{\star}
  \DeltaRow{Relationsoperatoren}{\mathcal J_{A,\star}}
}
\begin{tabproofwide}
  \proofstepwidestar[1]{\SemigroupAlg{A}{\star}}{\rA}
  \proofstepwidestar[2]{a\in A}{\rA}
  \proofstepwidestar[3]{b\in A}{\rA}
  \proofstepwidestar[4]{c\in A}{\rA}
  \proofstepwidestar[5]{a\mathrel{\mathcal J_{A,\star}}b}{\rA}
  \proofstepwidestar[6]{b\mathrel{\mathcal J_{A,\star}}c}{\rA}
  \proofstepwidestar[1,2,3,5]{J_{A,\star}(a)=J_{A,\star}(b)}{%
    \FormulaRefAuto{SemigroupGreenJDef}[def]{1,2,3,5}}
  \proofstepwidestar[1,3,4,6]{J_{A,\star}(b)=J_{A,\star}(c)}{%
    \FormulaRefAuto{SemigroupGreenJDef}[def]{1,3,4,6}}
  \proofstepwidestar[1,2,3,4,5,6]{J_{A,\star}(a)=J_{A,\star}(c)}{%
    \FormulaRefAuto{a=b\dsep b=c\vdash a=c}[thm]{7,8}}
  \proofstepwidestar[1,2,3,4,5,6]{a\mathrel{\mathcal J_{A,\star}}c}{%
    \FormulaRefAuto{SemigroupGreenJDef}[def]{1,2,4,9}}
\end{tabproofwide}

\FormulaThmDeltaK[\(\mathcal J\) ist eine Äquivalenzrelation]{%
  \SemigroupAlg{A}{\star}\vdash\EqRel{A,\mathcal J_{A,\star}}%
}{SemigroupGreenJIsEquivalence}{%
  \DeltaRow{Mengen}{A\dsep a\dsep b\dsep c}
  \DeltaRow{Funktionen}{\star}
  \DeltaRow{Relationsoperatoren}{\mathcal J_{A,\star}}
}
\begin{tabproofwide}
  \proofstepwidestar[1]{\SemigroupAlg{A}{\star}}{\rA}
  \proofstepwidestar[1]{%
    \forall a\in A\,
      \bigl(a\mathrel{\mathcal J_{A,\star}}a\bigr)}{%
    \FormulaRefAuto{SemigroupGreenJReflexive}[thm]{1}}
  \proofstepwidestar[1]{%
    \forall a,b\in A\,
      \bigl(a\mathrel{\mathcal J_{A,\star}}b\rightarrow
        b\mathrel{\mathcal J_{A,\star}}a\bigr)}{%
    \FormulaRefAuto{SemigroupGreenJSymmetric}[thm]{1}}
  \proofstepwidestar[1]{%
    \forall a,b,c\in A\,
      \Bigl(\bigl(a\mathrel{\mathcal J_{A,\star}}b\land
        b\mathrel{\mathcal J_{A,\star}}c\bigr)
        \rightarrow a\mathrel{\mathcal J_{A,\star}}c\Bigr)}{%
    \FormulaRefAuto{SemigroupGreenJTransitive}[thm]{1}}
  \proofstepwidestar[1]{\EqRel{A,\mathcal J_{A,\star}}}{%
    \FormulaRefAuto{EqRelDef}[def]{2,3,4}}
\end{tabproofwide}

\FormulaThmDeltaK[Reflexivität von \(\mathcal H\)]{%
  \SemigroupAlg{A}{\star}\dsep a\in A\vdash
  a\mathrel{\mathcal H_{A,\star}}a%
}{SemigroupGreenHReflexive}{%
  \DeltaRow{Mengen}{A\dsep a}
  \DeltaRow{Funktionen}{\star}
  \DeltaRow{Relationsoperatoren}{\mathcal H_{A,\star}}
}
\begin{tabproofwide}
  \proofstepwidestar[1]{\SemigroupAlg{A}{\star}}{\rA}
  \proofstepwidestar[2]{a\in A}{\rA}
  \proofstepwidestar[1,2]{a\mathrel{\mathcal L_{A,\star}}a}{%
    \FormulaRefAuto{SemigroupGreenLReflexive}[thm]{1,2}}
  \proofstepwidestar[1,2]{a\mathrel{\mathcal R_{A,\star}}a}{%
    \FormulaRefAuto{SemigroupGreenRReflexive}[thm]{1,2}}
  \proofstepwidestar[1,2]{a\mathrel{\mathcal H_{A,\star}}a}{%
    \FormulaRefAuto{SemigroupGreenHDef}[def]{1,2,2,\rAI{3,4}}}
\end{tabproofwide}

\FormulaThmDeltaK[Symmetrie von \(\mathcal H\)]{%
  \SemigroupAlg{A}{\star}\dsep a,b\in A\dsep
  a\mathrel{\mathcal H_{A,\star}}b\vdash
  b\mathrel{\mathcal H_{A,\star}}a%
}{SemigroupGreenHSymmetric}{%
  \DeltaRow{Mengen}{A\dsep a\dsep b}
  \DeltaRow{Funktionen}{\star}
  \DeltaRow{Relationsoperatoren}{\mathcal H_{A,\star}}
}
\begin{tabproofwide}
  \proofstepwidestar[1]{\SemigroupAlg{A}{\star}}{\rA}
  \proofstepwidestar[2]{a\in A}{\rA}
  \proofstepwidestar[3]{b\in A}{\rA}
  \proofstepwidestar[4]{a\mathrel{\mathcal H_{A,\star}}b}{\rA}
  \proofstepwidestar[1,2,3,4]{%
    a\mathrel{\mathcal L_{A,\star}}b\land
    a\mathrel{\mathcal R_{A,\star}}b}{%
    \FormulaRefAuto{SemigroupGreenHDef}[def]{1,2,3,4}}
  \proofstepwidestar[1,2,3,4]{b\mathrel{\mathcal L_{A,\star}}a}{%
    \FormulaRefAuto{SemigroupGreenLSymmetric}[thm]{1,2,3,\rAEa{5}}}
  \proofstepwidestar[1,2,3,4]{b\mathrel{\mathcal R_{A,\star}}a}{%
    \FormulaRefAuto{SemigroupGreenRSymmetric}[thm]{1,2,3,\rAEb{5}}}
  \proofstepwidestar[1,2,3,4]{b\mathrel{\mathcal H_{A,\star}}a}{%
    \FormulaRefAuto{SemigroupGreenHDef}[def]{1,3,2,\rAI{6,7}}}
\end{tabproofwide}

\FormulaThmDeltaK[Transitivität von \(\mathcal H\)]{%
  \begin{aligned}[t]
    &\SemigroupAlg{A}{\star}\dsep a,b,c\in A\dsep
      a\mathrel{\mathcal H_{A,\star}}b\dsep
      b\mathrel{\mathcal H_{A,\star}}c\vdash{}\\[-2pt]
    &a\mathrel{\mathcal H_{A,\star}}c
  \end{aligned}
}{SemigroupGreenHTransitive}{%
  \DeltaRow{Mengen}{A\dsep a\dsep b\dsep c}
  \DeltaRow{Funktionen}{\star}
  \DeltaRow{Relationsoperatoren}{\mathcal H_{A,\star}}
}
\begin{tabproofwide}
  \proofstepwidestar[1]{\SemigroupAlg{A}{\star}}{\rA}
  \proofstepwidestar[2]{a\in A}{\rA}
  \proofstepwidestar[3]{b\in A}{\rA}
  \proofstepwidestar[4]{c\in A}{\rA}
  \proofstepwidestar[5]{a\mathrel{\mathcal H_{A,\star}}b}{\rA}
  \proofstepwidestar[6]{b\mathrel{\mathcal H_{A,\star}}c}{\rA}
  \proofstepwidestar[1,2,3,5]{%
    a\mathrel{\mathcal L_{A,\star}}b\land
    a\mathrel{\mathcal R_{A,\star}}b}{%
    \FormulaRefAuto{SemigroupGreenHDef}[def]{1,2,3,5}}
  \proofstepwidestar[1,3,4,6]{%
    b\mathrel{\mathcal L_{A,\star}}c\land
    b\mathrel{\mathcal R_{A,\star}}c}{%
    \FormulaRefAuto{SemigroupGreenHDef}[def]{1,3,4,6}}
  \proofstepwidestar[1,2,3,4,5,6]{a\mathrel{\mathcal L_{A,\star}}c}{%
    \FormulaRefAuto{SemigroupGreenLTransitive}[thm]{1,2,3,4,\rAEa{7},\rAEa{8}}}
  \proofstepwidestar[1,2,3,4,5,6]{a\mathrel{\mathcal R_{A,\star}}c}{%
    \FormulaRefAuto{SemigroupGreenRTransitive}[thm]{1,2,3,4,\rAEb{7},\rAEb{8}}}
  \proofstepwidestar[1,2,3,4,5,6]{a\mathrel{\mathcal H_{A,\star}}c}{%
    \FormulaRefAuto{SemigroupGreenHDef}[def]{1,2,4,\rAI{9,10}}}
\end{tabproofwide}

\FormulaThmDeltaK[\(\mathcal H\) ist eine Äquivalenzrelation]{%
  \SemigroupAlg{A}{\star}\vdash\EqRel{A,\mathcal H_{A,\star}}%
}{SemigroupGreenHIsEquivalence}{%
  \DeltaRow{Mengen}{A\dsep a\dsep b\dsep c}
  \DeltaRow{Funktionen}{\star}
  \DeltaRow{Relationsoperatoren}{\mathcal H_{A,\star}}
}
\begin{tabproofwide}
  \proofstepwidestar[1]{\SemigroupAlg{A}{\star}}{\rA}
  \proofstepwidestar[1]{%
    \forall a\in A\,
      \bigl(a\mathrel{\mathcal H_{A,\star}}a\bigr)}{%
    \FormulaRefAuto{SemigroupGreenHReflexive}[thm]{1}}
  \proofstepwidestar[1]{%
    \forall a,b\in A\,
      \bigl(a\mathrel{\mathcal H_{A,\star}}b\rightarrow
        b\mathrel{\mathcal H_{A,\star}}a\bigr)}{%
    \FormulaRefAuto{SemigroupGreenHSymmetric}[thm]{1}}
  \proofstepwidestar[1]{%
    \forall a,b,c\in A\,
      \Bigl(\bigl(a\mathrel{\mathcal H_{A,\star}}b\land
        b\mathrel{\mathcal H_{A,\star}}c\bigr)
        \rightarrow a\mathrel{\mathcal H_{A,\star}}c\Bigr)}{%
    \FormulaRefAuto{SemigroupGreenHTransitive}[thm]{1}}
  \proofstepwidestar[1]{\EqRel{A,\mathcal H_{A,\star}}}{%
    \FormulaRefAuto{EqRelDef}[def]{2,3,4}}
\end{tabproofwide}

\section{Morphismen von Halbgruppen}

\FormulaDefDeltaK[Begriff des Halbgruppenhomomorphismus]{%
  \SemigroupHom{f}{A,\star}{B,\diamond}%
}{SemigroupHomDef}{%
  \DeltaRow{Mengen}{A\dsep B\dsep x\dsep y}
  \DeltaRow{Funktionen}{f\dsep\star\dsep\diamond}
  \DeltaRow{\textbf{Axiome}}{}
  \DeltaPrem{Quellhalbgruppe}{\SemigroupAlg{A}{\star}}
  \DeltaPrem{Zielhalbgruppe}{\SemigroupAlg{B}{\diamond}}
  \DeltaPrem{Abbildung}{f\colon A\to B}
  \DeltaRow{Operationserhaltung}{%
    x,y\in A\vdash
    f(x\star y)=f(x)\diamond f(y)}
  \DeltaRow{\textbf{Neues Symbol}}{}
  \DeltaPrem{Halbgruppenhomomorphismus}{%
    \SemigroupHom{f}{A,\star}{B,\diamond}}
}

\FormulaAxiomDeltaK[Quellhalbgruppe]{%
  \SemigroupHom{f}{A,\star}{B,\diamond}
  \vdash\SemigroupAlg{A}{\star}%
}{SemigroupHomSourceAxiom}{%
  \DeltaRow{Mengen}{A\dsep B}
  \DeltaRow{Funktionen}{f\dsep\star\dsep\diamond}
}

\FormulaAxiomDeltaK[Zielhalbgruppe]{%
  \SemigroupHom{f}{A,\star}{B,\diamond}
  \vdash\SemigroupAlg{B}{\diamond}%
}{SemigroupHomTargetAxiom}{%
  \DeltaRow{Mengen}{A\dsep B}
  \DeltaRow{Funktionen}{f\dsep\star\dsep\diamond}
}

\FormulaAxiomDeltaK[Operationserhaltung eines Halbgruppenhomomorphismus]{%
  \SemigroupHom{f}{A,\star}{B,\diamond}\dsep x,y\in A
  \vdash f(x\star y)=f(x)\diamond f(y)%
}{SemigroupHomOperationAxiom}{%
  \DeltaRow{Mengen}{A\dsep B\dsep x\dsep y}
  \DeltaRow{Funktionen}{f\dsep\star\dsep\diamond}
}

\FormulaAxiomDeltaK[Trägerabbildung eines Halbgruppenhomomorphismus]{%
  \SemigroupHom{f}{A,\star}{B,\diamond}
  \vdash f\colon A\to B%
}{SemigroupHomFunctionAxiom}{%
  \DeltaRow{Mengen}{A\dsep B}
  \DeltaRow{Funktionen}{f\dsep\star\dsep\diamond}
}

\FormulaDefDeltaK[Begriff des Halbgruppenisomorphismus]{%
  \SemigroupIso{f}{A,\star}{B,\diamond}%
}{SemigroupIsoDef}{%
  \DeltaRow{Mengen}{A\dsep B}
  \DeltaRow{Funktionen}{f\dsep\star\dsep\diamond}
  \DeltaRow{\textbf{Axiome}}{}
  \DeltaPrem{Halbgruppenhomomorphismus}{%
    \SemigroupHom{f}{A,\star}{B,\diamond}}
  \DeltaPrem{Bijektivität}{f\colon A\bij B}
  \DeltaRow{\textbf{Neues Symbol}}{}
  \DeltaPrem{Halbgruppenisomorphismus}{%
    \SemigroupIso{f}{A,\star}{B,\diamond}}
}

\FormulaAxiomDeltaK[Homomorphismuszugriff]{%
  \SemigroupIso{f}{A,\star}{B,\diamond}
  \vdash\SemigroupHom{f}{A,\star}{B,\diamond}%
}{SemigroupIsoHomAxiom}{%
  \DeltaRow{Mengen}{A\dsep B}
  \DeltaRow{Funktionen}{f\dsep\star\dsep\diamond}
}

\FormulaAxiomDeltaK[Bijektivitätszugriff]{%
  \SemigroupIso{f}{A,\star}{B,\diamond}
  \vdash f\colon A\bij B%
}{SemigroupIsoBijectiveAxiom}{%
  \DeltaRow{Mengen}{A\dsep B}
  \DeltaRow{Funktionen}{f\dsep\star\dsep\diamond}
}

\subsection{Grundlegende Eigenschaften}

\FormulaThmDeltaK[Komposition von Halbgruppenhomomorphismen]{%
  \begin{aligned}[t]
    &\SemigroupHom{f}{A,\star}{B,\diamond}\dsep
      \SemigroupHom{g}{B,\diamond}{C,\bullet}\\
    &\vdash\SemigroupHom{g\circ f}{A,\star}{C,\bullet}
  \end{aligned}
}{SemigroupHomComposition}{%
  \DeltaRow{Mengen}{A\dsep B\dsep C\dsep x\dsep y}
  \DeltaRow{Funktionen}{f\dsep g\dsep\star\dsep\diamond\dsep\bullet}
}
\begin{tabproofsplitwide}
  \proofpartwideindR[Komposition der Trägerabbildungen]{%
    \begin{aligned}[t]
      &\SemigroupHom{f}{A,\star}{B,\diamond}\dsep
        \SemigroupHom{g}{B,\diamond}{C,\bullet}\vdash{}\\[-2pt]
      &g\circ f\colon A\to C
    \end{aligned}
  }{%
    \SemigroupHom{f}{A,\star}{B,\diamond}\dsep
    \SemigroupHom{g}{B,\diamond}{C,\bullet}
    \vdash g\circ f\colon A\to C%
  }[SemigroupHomCompositionFunction]
    \proofstepwidestar[1]{\SemigroupHom{f}{A,\star}{B,\diamond}}{\rA}
    \proofstepwidestar[2]{\SemigroupHom{g}{B,\diamond}{C,\bullet}}{\rA}
    \proofstepwidestar[1]{f\colon A\to B}{%
      \FormulaRefAuto{SemigroupHomFunctionAxiom}[ax]{1}}
    \proofstepwidestar[2]{g\colon B\to C}{%
      \FormulaRefAuto{SemigroupHomFunctionAxiom}[ax]{2}}
    \proofstepwidestar[1,2]{g\circ f\colon A\to C}{%
      \FormulaRefAuto{G\circ F\colon A\to C}[thm]{3,4}}
  \closeproofpartwideindR

  \proofpartwideindR[Einsetzen der ersten Homomorphie]{%
    \begin{aligned}[t]
      &\SemigroupHom{f}{A,\star}{B,\diamond}\dsep
        \SemigroupHom{g}{B,\diamond}{C,\bullet}\dsep x,y\in A\vdash{}\\[-2pt]
      &(g\circ f)(x\star y)=g\bigl(f(x)\diamond f(y)\bigr)
    \end{aligned}
  }{%
    \SemigroupHom{f}{A,\star}{B,\diamond}\dsep
    \SemigroupHom{g}{B,\diamond}{C,\bullet}\dsep x,y\in A
    \vdash (g\circ f)(x\star y)=g\bigl(f(x)\diamond f(y)\bigr)%
  }[SemigroupHomCompositionFirstOperation]
    \proofstepwidestar[1]{\SemigroupHom{f}{A,\star}{B,\diamond}}{\rA}
    \proofstepwidestar[2]{\SemigroupHom{g}{B,\diamond}{C,\bullet}}{\rA}
    \proofstepwidestar[3]{x\in A}{\rA}
    \proofstepwidestar[4]{y\in A}{\rA}
    \proofstepwidestar[1]{f\colon A\to B}{%
      \FormulaRefAuto{SemigroupHomFunctionAxiom}[ax]{1}}
    \proofstepwidestar[2]{g\colon B\to C}{%
      \FormulaRefAuto{SemigroupHomFunctionAxiom}[ax]{2}}
    \proofstepwidestar[1]{\SemigroupAlg{A}{\star}}{%
      \FormulaRefAuto{SemigroupHomSourceAxiom}[ax]{1}}
    \proofstepwidestar[1,3,4]{x\star y\in A}{%
      \FormulaRefAuto{SemigroupClosureAxiom}[ax]{7,3,4}}
    \proofstepwidestar[1,3,4]{f(x\star y)=f(x)\diamond f(y)}{%
      \FormulaRefAuto{SemigroupHomOperationAxiom}[ax]{1,3,4}}
    \proofstepwide[1,2,3,4]{(g\circ f)(x\star y)}{=}{g(f(x\star y))}{%
      \FormulaRefAuto{%
        x\in A \vdash (F\circ G)(x)=F(G(x))%
      }[thm]{5,6,8}}
    \proofstepwide[1,2,3,4]{}{=}{g\bigl(f(x)\diamond f(y)\bigr)}{%
      \rIE{9}}
    \proofstepwidestar[1,2,3,4]{%
      (g\circ f)(x\star y)=g\bigl(f(x)\diamond f(y)\bigr)}{%
      \rChain{10,11}}
  \closeproofpartwideindR

  \proofpartwideindR[Einsetzen der zweiten Homomorphie]{%
    \begin{aligned}[t]
      &\SemigroupHom{f}{A,\star}{B,\diamond}\dsep
        \SemigroupHom{g}{B,\diamond}{C,\bullet}\dsep x,y\in A\vdash{}\\[-2pt]
      &g\bigl(f(x)\diamond f(y)\bigr)
        =(g\circ f)(x)\bullet(g\circ f)(y)
    \end{aligned}
  }{%
    \SemigroupHom{f}{A,\star}{B,\diamond}\dsep
    \SemigroupHom{g}{B,\diamond}{C,\bullet}\dsep x,y\in A
    \vdash g\bigl(f(x)\diamond f(y)\bigr)
      =(g\circ f)(x)\bullet(g\circ f)(y)%
  }[SemigroupHomCompositionSecondOperation]
    \proofstepwidestar[1]{\SemigroupHom{f}{A,\star}{B,\diamond}}{\rA}
    \proofstepwidestar[2]{\SemigroupHom{g}{B,\diamond}{C,\bullet}}{\rA}
    \proofstepwidestar[3]{x\in A}{\rA}
    \proofstepwidestar[4]{y\in A}{\rA}
    \proofstepwidestar[1]{f\colon A\to B}{%
      \FormulaRefAuto{SemigroupHomFunctionAxiom}[ax]{1}}
    \proofstepwidestar[2]{g\colon B\to C}{%
      \FormulaRefAuto{SemigroupHomFunctionAxiom}[ax]{2}}
    \proofstepwidestar[1,3]{f(x)\in B}{%
      \FormulaRefAuto{%
        F\colon A\to B\dsep x\in A\vdash F(x)\in B%
      }[thm]{5,3}}
    \proofstepwidestar[1,4]{f(y)\in B}{%
      \FormulaRefAuto{%
        F\colon A\to B\dsep x\in A\vdash F(x)\in B%
      }[thm]{5,4}}
    \proofstepwidestar[1,2,3]{g(f(x))=(g\circ f)(x)}{%
      \FormulaRefAuto{%
        x\in A \vdash F(G(x))=(F\circ G)(x)%
      }[thm]{5,6,3}}
    \proofstepwidestar[1,2,4]{g(f(y))=(g\circ f)(y)}{%
      \FormulaRefAuto{%
        x\in A \vdash F(G(x))=(F\circ G)(x)%
      }[thm]{5,6,4}}
    \proofstepwide[1,2,3,4]{g\bigl(f(x)\diamond f(y)\bigr)}{=}{%
      g(f(x))\bullet g(f(y))}{%
      \FormulaRefAuto{SemigroupHomOperationAxiom}[ax]{2,7,8}}
    \proofstepwide[1,2,3,4]{}{=}{(g\circ f)(x)\bullet g(f(y))}{\rIE{9}}
    \proofstepwide[1,2,3,4]{}{=}{%
      (g\circ f)(x)\bullet(g\circ f)(y)}{\rIE{10}}
    \proofstepwidestar[1,2,3,4]{%
      g\bigl(f(x)\diamond f(y)\bigr)
      =(g\circ f)(x)\bullet(g\circ f)(y)}{%
      \rChain{11,13}}
  \closeproofpartwideindR

  \proofpartwideindR[Zusammenführung]{%
    \begin{aligned}[t]
      &\SemigroupHom{f}{A,\star}{B,\diamond}\dsep
        \SemigroupHom{g}{B,\diamond}{C,\bullet}\vdash{}\\[-2pt]
      &\SemigroupHom{g\circ f}{A,\star}{C,\bullet}
    \end{aligned}
  }{%
    \SemigroupHom{f}{A,\star}{B,\diamond}\dsep
    \SemigroupHom{g}{B,\diamond}{C,\bullet}
    \vdash \SemigroupHom{g\circ f}{A,\star}{C,\bullet}%
  }[SemigroupHomCompositionConclusion]
    \proofstepwidestar[1]{\SemigroupHom{f}{A,\star}{B,\diamond}}{\rA}
    \proofstepwidestar[2]{\SemigroupHom{g}{B,\diamond}{C,\bullet}}{\rA}
    \proofstepwidestar[1]{\SemigroupAlg{A}{\star}}{%
      \FormulaRefAuto{SemigroupHomSourceAxiom}[ax]{1}}
    \proofstepwidestar[2]{\SemigroupAlg{C}{\bullet}}{%
      \FormulaRefAuto{SemigroupHomTargetAxiom}[ax]{2}}
    \proofstepwidestar[1,2]{g\circ f\colon A\to C}{%
      \FormulaRefAuto{SemigroupHomCompositionFunction}[thm]{1,2}}
    \proofstepwidestar[6]{x\in A}{\rA}
    \proofstepwidestar[7]{y\in A}{\rA}
    \proofstepwidestar[1,2,6,7]{%
      (g\circ f)(x\star y)=g\bigl(f(x)\diamond f(y)\bigr)}{%
      \FormulaRefAuto{SemigroupHomCompositionFirstOperation}[thm]{1,2,6,7}}
    \proofstepwidestar[1,2,6,7]{%
      g\bigl(f(x)\diamond f(y)\bigr)
      =(g\circ f)(x)\bullet(g\circ f)(y)}{%
      \FormulaRefAuto{SemigroupHomCompositionSecondOperation}[thm]{1,2,6,7}}
    \proofstepwidestar[1,2,6,7]{%
      (g\circ f)(x\star y)=(g\circ f)(x)\bullet(g\circ f)(y)}{%
      \rChain{8,9}}
    \proofstepwidestar[1,2]{%
      \SemigroupHom{g\circ f}{A,\star}{C,\bullet}}{%
      \FormulaRefAuto{SemigroupHomDef}[def]{3,4,5,10}}
  \closeproofpartwideindR
\end{tabproofsplitwide}

\FormulaThmDeltaK[Die Identität ist ein Halbgruppenisomorphismus]{%
  \SemigroupAlg{A}{\star}
  \vdash\SemigroupIso{\Id_A}{A,\star}{A,\star}%
}{SemigroupIdentityIsomorphism}{%
  \DeltaRow{Mengen}{A\dsep x\dsep y}
  \DeltaRow{Funktionen}{\star}
}
\begin{tabproofwide}
  \proofstepwidestar[1]{\SemigroupAlg{A}{\star}}{\rA}
  \proofstepwidestar[]{\Id_A\colon A\to A}{%
    \FormulaRefAuto{IdA}[thm]}
  \proofstepwidestar[2]{x\in A}{\rA}
  \proofstepwidestar[3]{y\in A}{\rA}
  \proofstepwidestar[1,2,3]{x\star y\in A}{%
    \FormulaRefAuto{SemigroupClosureAxiom}[ax]{1,2,3}}
  \proofstepwidestar[2]{\Id_A(x)=x}{%
    \FormulaRefAuto{x\in A\vdash \Id_A(x)=x}{2}}
  \proofstepwidestar[2]{x=\Id_A(x)}{%
    \FormulaRefAuto{a=b\vdash b=a}{6}}
  \proofstepwidestar[3]{\Id_A(y)=y}{%
    \FormulaRefAuto{x\in A\vdash \Id_A(x)=x}{3}}
  \proofstepwidestar[3]{y=\Id_A(y)}{%
    \FormulaRefAuto{a=b\vdash b=a}{8}}
  \proofstepwide[1,2,3]{\Id_A(x\star y)}{=}{x\star y}{%
    \FormulaRefAuto{x\in A\vdash \Id_A(x)=x}{5}}
  \proofstepwide[1,2,3]{}{=}{\Id_A(x)\star y}{\rIE{7}}
  \proofstepwide[1,2,3]{}{=}{\Id_A(x)\star\Id_A(y)}{\rIE{9}}
  \proofstepwidestar[1,2,3]{%
    \Id_A(x\star y)=\Id_A(x)\star\Id_A(y)}{\rChain{10,12}}
  \proofstepwidestar[1]{%
    \SemigroupHom{\Id_A}{A,\star}{A,\star}}{%
    \FormulaRefAuto{SemigroupHomDef}[def]{1,1,2,13}}
  \proofstepwidestar[]{\Id_A\colon A\bij A}{%
    \FormulaRefAuto{\Id_A\colon A\bij A}[thm]}
  \proofstepwidestar[1]{%
    \SemigroupIso{\Id_A}{A,\star}{A,\star}}{%
    \FormulaRefAuto{SemigroupIsoDef}[def]{14,15}}
\end{tabproofwide}

\FormulaThmDeltaK[Komposition von Halbgruppenisomorphismen]{%
  \begin{aligned}[t]
    &\SemigroupIso{f}{A,\star}{B,\diamond}\dsep
      \SemigroupIso{g}{B,\diamond}{C,\bullet}\\
    &\vdash\SemigroupIso{g\circ f}{A,\star}{C,\bullet}
  \end{aligned}
}{SemigroupIsoComposition}{%
  \DeltaRow{Mengen}{A\dsep B\dsep C}
  \DeltaRow{Funktionen}{f\dsep g\dsep\star\dsep\diamond\dsep\bullet}
}
\begin{tabproof}
  \proofstep{1}{\SemigroupIso{f}{A,\star}{B,\diamond}}{\rA}
  \proofstep{2}{\SemigroupIso{g}{B,\diamond}{C,\bullet}}{\rA}
  \proofstep{1}{\SemigroupHom{f}{A,\star}{B,\diamond}}{%
    \FormulaRefAuto{SemigroupIsoHomAxiom}[ax]{1}}
  \proofstep{2}{\SemigroupHom{g}{B,\diamond}{C,\bullet}}{%
    \FormulaRefAuto{SemigroupIsoHomAxiom}[ax]{2}}
  \proofstep{1,2}{\SemigroupHom{g\circ f}{A,\star}{C,\bullet}}{%
    \FormulaRefAuto{SemigroupHomComposition}[thm]{3,4}}
  \proofstep{1}{f\colon A\bij B}{%
    \FormulaRefAuto{SemigroupIsoBijectiveAxiom}[ax]{1}}
  \proofstep{2}{g\colon B\bij C}{%
    \FormulaRefAuto{SemigroupIsoBijectiveAxiom}[ax]{2}}
  \proofstep{1,2}{g\circ f\colon A\bij C}{%
    \FormulaRefAuto{F\colon A\bij B \dsep G\colon B\bij C\vdash
      G\circ F\colon A\bij C}[thm]{6,7}}
  \proofstep{1,2}{\SemigroupIso{g\circ f}{A,\star}{C,\bullet}}{%
    \FormulaRefAuto{SemigroupIsoDef}[def]{5,8}}
\end{tabproof}

\section{Koordinatenisomorphismen von Halbgruppen}

Die folgende Konstruktion trennt die beiden Koordinaten eines Produkts:
Die erste Koordinate eines Produkts stammt vom linken, die zweite vom
rechten Faktor.  F\"ur ein festes Element \(e\) werden die daf\"ur ben\"otigten
Koordinaten durch Rechts- beziehungsweise Linksmultiplikation mit \(e\)
gewonnen.  Der Ausdruck \emph{\(e\)-Koordinatenisomorphismus} ist die in
diesem Werk verwendete Bezeichnung f\"ur diese spezielle Darstellung.

\subsection{Axiomatische Definition}

\FormulaDefDeltaK[Koordinatentr\"ager bez\"uglich eines
Halbgruppenelements]{%
  \begin{aligned}[t]
    &\SemigroupAlg{A}{\star}\dsep e\in A\vdash{}\\[-2pt]
    &\SemigroupCoordinateCarrier{A}{\star}{e}\coloneqq
      \bigl\{p\in A\bigm|\exists x\in A\,(p=x\star e)\bigr\}
      \times
      \bigl\{q\in A\bigm|\exists y\in A\,(q=e\star y)\bigr\}
  \end{aligned}%
}{SemigroupCoordinateCarrierDef}{%
  \DeltaRow{Mengen}{A\dsep e\dsep p\dsep q\dsep x\dsep y}
  \DeltaRow{Funktionen}{\star}
  \DeltaRow{\textbf{Neues Symbol}}{}
  \DeltaPrem{Koordinatentr\"ager}{%
    \SemigroupCoordinateCarrier{A}{\star}{e}}
}

\FormulaDefDeltaK[Koordinatenprodukt]{%
  \begin{aligned}[t]
    &\SemigroupAlg{A}{\star}\dsep e\in A\vdash{}\\[-2pt]
    &\SemigroupCoordinateProductOp{A}{\star}{e}\colon
      \SemigroupCoordinateCarrier{A}{\star}{e}
      \times\SemigroupCoordinateCarrier{A}{\star}{e}
      \longrightarrow\SemigroupCoordinateCarrier{A}{\star}{e},\\[-2pt]
    &(p,q),(r,s)\in\SemigroupCoordinateCarrier{A}{\star}{e}\vdash
      (p,q)\SemigroupCoordinateProductOp{A}{\star}{e}(r,s)
      \coloneqq(p,s)
  \end{aligned}%
}{SemigroupCoordinateProductDef}{%
  \DeltaRow{Mengen}{A\dsep e\dsep p\dsep q\dsep r\dsep s}
  \DeltaRow{Funktionen}{\star\dsep
    \SemigroupCoordinateProductOp{A}{\star}{e}}
  \DeltaRow{\textbf{Neues Symbol}}{}
  \DeltaPrem{Koordinatenprodukt}{%
    \SemigroupCoordinateProductOp{A}{\star}{e}}
}

\FormulaDefDeltaK[\(e\)-Koordinatenabbildung]{%
  \begin{aligned}[t]
    &\SemigroupAlg{A}{\star}\dsep e\in A\vdash{}\\[-2pt]
    &\SemigroupCoordinateMap{A}{\star}{e}\colon
      A\longrightarrow\SemigroupCoordinateCarrier{A}{\star}{e},\\[-2pt]
    &\forall x\in A\quad
      \SemigroupCoordinateMap{A}{\star}{e}(x)
      \coloneqq(x\star e,e\star x)
  \end{aligned}%
}{SemigroupCoordinateMapDef}{%
  \DeltaRow{Mengen}{A\dsep e\dsep x}
  \DeltaRow{Funktionen}{\star\dsep
    \SemigroupCoordinateMap{A}{\star}{e}}
  \DeltaRow{\textbf{Neues Symbol}}{}
  \DeltaPrem{\(e\)-Koordinatenabbildung}{%
    \SemigroupCoordinateMap{A}{\star}{e}}
}

\FormulaDefDeltaK[Multiplikationsabbildung der \(e\)-Koordinaten]{%
  \begin{aligned}[t]
    &\SemigroupAlg{A}{\star}\dsep e\in A\vdash{}\\[-2pt]
    &\SemigroupCoordinateMultiplicationMap{A}{\star}{e}\colon
      \SemigroupCoordinateCarrier{A}{\star}{e}\longrightarrow A,\\[-2pt]
    &(p,q)\in\SemigroupCoordinateCarrier{A}{\star}{e}\vdash
      \SemigroupCoordinateMultiplicationMap{A}{\star}{e}(p,q)
      \coloneqq p\star q
  \end{aligned}%
}{SemigroupCoordinateMultiplicationMapDef}{%
  \DeltaRow{Mengen}{A\dsep e\dsep p\dsep q}
  \DeltaRow{Funktionen}{\star\dsep
    \SemigroupCoordinateMultiplicationMap{A}{\star}{e}}
  \DeltaRow{\textbf{Neues Symbol}}{}
  \DeltaPrem{Multiplikationsabbildung}{%
    \SemigroupCoordinateMultiplicationMap{A}{\star}{e}}
}

Der neue Strukturbegriff wird nicht als lange Konjunktion notiert.  Seine
vier Bestandteile werden vielmehr als getrennte Axiome eingef\"uhrt und sind
anschlie\ss end einzeln zug\"anglich.

\FormulaDefDeltaK[Begriff des \(e\)-Koordinatenisomorphismus]{%
  \SemigroupCoordinateIso{A}{\star}{e}%
}{SemigroupCoordinateIsoDef}{%
  \DeltaRow{Mengen}{A\dsep e\dsep x\dsep y}
  \DeltaRow{Funktionen}{\star\dsep
    \SemigroupCoordinateMap{A}{\star}{e}\dsep
    \SemigroupCoordinateProductOp{A}{\star}{e}}
  \DeltaRow{\textbf{Axiome}}{}
  \DeltaPrem{Quellhalbgruppe}{\SemigroupAlg{A}{\star}}
  \DeltaPrem{Basiselement}{e\in A}
  \DeltaPrem{Bijektive Koordinatenabbildung}{%
    \SemigroupCoordinateMap{A}{\star}{e}\colon
    A\bij\SemigroupCoordinateCarrier{A}{\star}{e}}
  \DeltaRow{Operationserhaltung}{%
    \begin{gathered}
      x,y\in A\vdash{}\\[-2pt]
      \SemigroupCoordinateMap{A}{\star}{e}(x\star y)
      =\SemigroupCoordinateMap{A}{\star}{e}(x)
        \SemigroupCoordinateProductOp{A}{\star}{e}
        \SemigroupCoordinateMap{A}{\star}{e}(y)
    \end{gathered}}
  \DeltaRow{\textbf{Neues Symbol}}{}
  \DeltaPrem{\(e\)-Koordinatenisomorphismus}{%
    \SemigroupCoordinateIso{A}{\star}{e}}
}

\FormulaAxiomDeltaK[Zugriff auf die Quellhalbgruppe eines
Koordinatenisomorphismus]{%
  \SemigroupCoordinateIso{A}{\star}{e}
  \vdash\SemigroupAlg{A}{\star}%
}{SemigroupCoordinateIsoSourceAxiom}{%
  \DeltaRow{Mengen}{A\dsep e}
  \DeltaRow{Funktionen}{\star}
}

\FormulaAxiomDeltaK[Zugriff auf das Basiselement eines
Koordinatenisomorphismus]{%
  \SemigroupCoordinateIso{A}{\star}{e}\vdash e\in A%
}{SemigroupCoordinateIsoBaseElementAxiom}{%
  \DeltaRow{Mengen}{A\dsep e}
  \DeltaRow{Funktionen}{\star}
}

\FormulaAxiomDeltaK[Zugriff auf die Bijektivit\"at der
Koordinatenabbildung]{%
  \SemigroupCoordinateIso{A}{\star}{e}\vdash
  \SemigroupCoordinateMap{A}{\star}{e}\colon
  A\bij\SemigroupCoordinateCarrier{A}{\star}{e}%
}{SemigroupCoordinateIsoBijectiveAxiom}{%
  \DeltaRow{Mengen}{A\dsep e}
  \DeltaRow{Funktionen}{\star\dsep
    \SemigroupCoordinateMap{A}{\star}{e}}
}

\FormulaAxiomDeltaK[Zugriff auf die Operationserhaltung der
Koordinatenabbildung]{%
  \begin{aligned}[t]
    &\SemigroupCoordinateIso{A}{\star}{e}\dsep x,y\in A\vdash{}\\[-2pt]
    &\SemigroupCoordinateMap{A}{\star}{e}(x\star y)
      =\SemigroupCoordinateMap{A}{\star}{e}(x)
       \SemigroupCoordinateProductOp{A}{\star}{e}
       \SemigroupCoordinateMap{A}{\star}{e}(y)
  \end{aligned}%
}{SemigroupCoordinateIsoOperationAxiom}{%
  \DeltaRow{Mengen}{A\dsep e\dsep x\dsep y}
  \DeltaRow{Funktionen}{\star\dsep
    \SemigroupCoordinateMap{A}{\star}{e}\dsep
    \SemigroupCoordinateProductOp{A}{\star}{e}}
}

\subsection{Grundlegende Eigenschaften}

\FormulaThmDeltaK[Elementkriterium des Koordinatentr\"agers]{%
  \begin{aligned}[t]
    &\SemigroupAlg{A}{\star}\dsep e\in A\vdash{}\\[-2pt]
    &(p,q)\in\SemigroupCoordinateCarrier{A}{\star}{e}
      \leftrightarrow\\[-2pt]
    &\quad\Bigl(p\in A\land\exists x\in A\,(p=x\star e)\Bigr)
      \land
      \Bigl(q\in A\land\exists y\in A\,(q=e\star y)\Bigr)
  \end{aligned}%
}{SemigroupCoordinateCarrierMembership}{%
  \DeltaRow{Mengen}{A\dsep e\dsep p\dsep q\dsep x\dsep y}
  \DeltaRow{Funktionen}{\star}
}
\begin{proof}
Setze in \FormulaRefAuto{SemigroupCoordinateCarrierDef}[def] die
Definition des Koordinatentr\"agers ein.  Das Elementkriterium des
kartesischen Produkts und zweimal das Elementkriterium einer
ausgesonderten Teilmenge liefern genau die beiden Konjunkte der rechten
Seite.
\end{proof}

\FormulaThmDeltaK[Der Koordinatentr\"ager ist nichtleer]{%
  \SemigroupAlg{A}{\star}\dsep e\in A\vdash
  \SemigroupCoordinateCarrier{A}{\star}{e}\neq\varnothing%
}{SemigroupCoordinateCarrierNonempty}{%
  \DeltaRow{Mengen}{A\dsep e}
  \DeltaRow{Funktionen}{\star}
}
\begin{proof}
Nach der Abgeschlossenheit liegt \(e\star e\) in \(A\).  In beiden
Faktoren des Koordinatentr\"agers dient \(e\) als Zeuge, also liegt
\((e\star e,e\star e)\) nach
\FormulaRefAuto{SemigroupCoordinateCarrierMembership}[thm] im
Koordinatentr\"ager.  Dieser ist daher nichtleer.
\end{proof}

\FormulaThmDeltaK[Auswertung des Koordinatenprodukts]{%
  \begin{aligned}[t]
    &\SemigroupAlg{A}{\star}\dsep e\in A\dsep
      (p,q),(r,s)\in\SemigroupCoordinateCarrier{A}{\star}{e}\vdash{}\\[-2pt]
    &(p,q)\SemigroupCoordinateProductOp{A}{\star}{e}(r,s)=(p,s)
  \end{aligned}%
}{SemigroupCoordinateProductEvaluation}{%
  \DeltaRow{Mengen}{A\dsep e\dsep p\dsep q\dsep r\dsep s}
  \DeltaRow{Funktionen}{\star\dsep
    \SemigroupCoordinateProductOp{A}{\star}{e}}
}
\begin{proof}
Dies ist die Auswertungsregel aus
\FormulaRefAuto{SemigroupCoordinateProductDef}[def].
\end{proof}

\FormulaThmDeltaK[Das Koordinatenprodukt bildet eine Halbgruppe]{%
  \SemigroupAlg{A}{\star}\dsep e\in A\vdash
  \SemigroupAlg{\SemigroupCoordinateCarrier{A}{\star}{e}}
    {\SemigroupCoordinateProductOp{A}{\star}{e}}%
}{SemigroupCoordinateProductSemigroup}{%
  \DeltaRow{Mengen}{A\dsep e\dsep p\dsep q\dsep r\dsep s\dsep
    t\dsep u}
  \DeltaRow{Funktionen}{\star\dsep
    \SemigroupCoordinateProductOp{A}{\star}{e}}
}
\begin{proof}
Die Abgeschlossenheit ist Teil der Typisierung in
\FormulaRefAuto{SemigroupCoordinateProductDef}[def].  Sind
\((p,q),(r,s),(t,u)\) Elemente des Koordinatentr\"agers, so liefert
\FormulaRefAuto{SemigroupCoordinateProductEvaluation}[thm] zweimal
\[
 \bigl((p,q)\SemigroupCoordinateProductOp{A}{\star}{e}(r,s)\bigr)
      \SemigroupCoordinateProductOp{A}{\star}{e}(t,u)
 =(p,u)
 =(p,q)\SemigroupCoordinateProductOp{A}{\star}{e}
      \bigl((r,s)\SemigroupCoordinateProductOp{A}{\star}{e}(t,u)\bigr).
\]
Damit ist das Koordinatenprodukt assoziativ; zusammen mit der
Abgeschlossenheit folgt die Behauptung aus
\FormulaRefAuto{SemigroupDef}[def].
\end{proof}

\FormulaThmDeltaK[Idempotenz des Koordinatenprodukts]{%
  \begin{aligned}[t]
    &\SemigroupAlg{A}{\star}\dsep e\in A\dsep
      u\in\SemigroupCoordinateCarrier{A}{\star}{e}\vdash{}\\[-2pt]
    &u\SemigroupCoordinateProductOp{A}{\star}{e}u=u
  \end{aligned}%
}{SemigroupCoordinateProductIdempotence}{%
  \DeltaRow{Mengen}{A\dsep e\dsep u\dsep p\dsep q}
  \DeltaRow{Funktionen}{\star\dsep
    \SemigroupCoordinateProductOp{A}{\star}{e}}
}
\begin{proof}
Nach dem Elementkriterium des kartesischen Produkts hat \(u\) die Form
\((p,q)\).  Die Auswertungsregel ergibt
\[
 (p,q)\SemigroupCoordinateProductOp{A}{\star}{e}(p,q)=(p,q).
\]
\end{proof}

\FormulaThmDeltaK[Dreitermgesetz des Koordinatenprodukts]{%
  \begin{aligned}[t]
    &\SemigroupAlg{A}{\star}\dsep e\in A\dsep
      u,v,w\in\SemigroupCoordinateCarrier{A}{\star}{e}\vdash{}\\[-2pt]
    &(u\SemigroupCoordinateProductOp{A}{\star}{e}v)
       \SemigroupCoordinateProductOp{A}{\star}{e}w
      =u\SemigroupCoordinateProductOp{A}{\star}{e}w
  \end{aligned}%
}{SemigroupCoordinateProductThreeTermLaw}{%
  \DeltaRow{Mengen}{A\dsep e\dsep u\dsep v\dsep w\dsep
    p\dsep q\dsep r\dsep s\dsep t\dsep z}
  \DeltaRow{Funktionen}{\star\dsep
    \SemigroupCoordinateProductOp{A}{\star}{e}}
}
\begin{proof}
Schreibe \(u=(p,q)\), \(v=(r,s)\) und \(w=(t,z)\).  Zweimalige
Anwendung von
\FormulaRefAuto{SemigroupCoordinateProductEvaluation}[thm] ergibt auf
beiden Seiten \((p,z)\).
\end{proof}

\FormulaThmDeltaK[Rechtecksidentit\"at des Koordinatenprodukts]{%
  \begin{aligned}[t]
    &\SemigroupAlg{A}{\star}\dsep e\in A\dsep
      u,v\in\SemigroupCoordinateCarrier{A}{\star}{e}\vdash{}\\[-2pt]
    &(u\SemigroupCoordinateProductOp{A}{\star}{e}v)
       \SemigroupCoordinateProductOp{A}{\star}{e}u=u
  \end{aligned}%
}{SemigroupCoordinateProductRectangularIdentity}{%
  \DeltaRow{Mengen}{A\dsep e\dsep u\dsep v}
  \DeltaRow{Funktionen}{\star\dsep
    \SemigroupCoordinateProductOp{A}{\star}{e}}
}
\begin{tabproofwide}
  \proofstepwidestar[1]{\SemigroupAlg{A}{\star}}{\rA}
  \proofstepwidestar[2]{e\in A}{\rA}
  \proofstepwidestar[3]{u,v\in
    \SemigroupCoordinateCarrier{A}{\star}{e}}{\rA}
  \proofstepwidestar[1,2,3]{%
    (u\SemigroupCoordinateProductOp{A}{\star}{e}v)
      \SemigroupCoordinateProductOp{A}{\star}{e}u
    =u\SemigroupCoordinateProductOp{A}{\star}{e}u}{%
    \FormulaRefAuto{SemigroupCoordinateProductThreeTermLaw}[thm]{1,2,3}}
  \proofstepwidestar[1,2,3]{%
    u\SemigroupCoordinateProductOp{A}{\star}{e}u=u}{%
    \FormulaRefAuto{SemigroupCoordinateProductIdempotence}[thm]{1,2,3}}
  \proofstepwidestar[1,2,3]{%
    (u\SemigroupCoordinateProductOp{A}{\star}{e}v)
      \SemigroupCoordinateProductOp{A}{\star}{e}u=u}{\rChain{4,5}}
\end{tabproofwide}

\FormulaThmDeltaK[Die \(e\)-Koordinatenabbildung ist eine Abbildung]{%
  \SemigroupAlg{A}{\star}\dsep e\in A\vdash
  \SemigroupCoordinateMap{A}{\star}{e}\colon
  A\to\SemigroupCoordinateCarrier{A}{\star}{e}%
}{SemigroupCoordinateMapFunction}{%
  \DeltaRow{Mengen}{A\dsep e}
  \DeltaRow{Funktionen}{\star\dsep
    \SemigroupCoordinateMap{A}{\star}{e}}
}
\begin{proof}
F\"ur \(x\in A\) liegen \(x\star e\) und \(e\star x\) nach der
Abgeschlossenheit in \(A\).  In den beiden Faktoren des
Koordinatentr\"agers dient jeweils \(x\) als Zeuge.  Somit liegt
\((x\star e,e\star x)\) im Koordinatentr\"ager.  Die Behauptung folgt aus
\FormulaRefAuto{SemigroupCoordinateMapDef}[def].
\end{proof}

\FormulaThmDeltaK[Auswertung der \(e\)-Koordinatenabbildung]{%
  \SemigroupAlg{A}{\star}\dsep e,x\in A\vdash
  \SemigroupCoordinateMap{A}{\star}{e}(x)=(x\star e,e\star x)%
}{SemigroupCoordinateMapEvaluation}{%
  \DeltaRow{Mengen}{A\dsep e\dsep x}
  \DeltaRow{Funktionen}{\star\dsep
    \SemigroupCoordinateMap{A}{\star}{e}}
}
\begin{proof}
Dies ist die Auswertungsregel aus
\FormulaRefAuto{SemigroupCoordinateMapDef}[def].
\end{proof}

\FormulaThmDeltaK[Die Multiplikationsabbildung ist eine Abbildung]{%
  \SemigroupAlg{A}{\star}\dsep e\in A\vdash
  \SemigroupCoordinateMultiplicationMap{A}{\star}{e}\colon
  \SemigroupCoordinateCarrier{A}{\star}{e}\to A%
}{SemigroupCoordinateMultiplicationMapFunction}{%
  \DeltaRow{Mengen}{A\dsep e\dsep p\dsep q}
  \DeltaRow{Funktionen}{\star\dsep
    \SemigroupCoordinateMultiplicationMap{A}{\star}{e}}
}
\begin{proof}
Ist \((p,q)\) im Koordinatentr\"ager, so sind \(p,q\in A\).  Daher liegt
\(p\star q\) nach der Abgeschlossenheit in \(A\).  Dies ist genau die
Wohldefiniertheit der Abbildung aus
\FormulaRefAuto{SemigroupCoordinateMultiplicationMapDef}[def].
\end{proof}

\FormulaThmDeltaK[Auswertung der Multiplikationsabbildung]{%
  \begin{aligned}[t]
    &\SemigroupAlg{A}{\star}\dsep e\in A\dsep
      (p,q)\in\SemigroupCoordinateCarrier{A}{\star}{e}\vdash{}\\[-2pt]
    &\SemigroupCoordinateMultiplicationMap{A}{\star}{e}(p,q)=p\star q
  \end{aligned}%
}{SemigroupCoordinateMultiplicationMapEvaluation}{%
  \DeltaRow{Mengen}{A\dsep e\dsep p\dsep q}
  \DeltaRow{Funktionen}{\star\dsep
    \SemigroupCoordinateMultiplicationMap{A}{\star}{e}}
}
\begin{proof}
Dies ist die Auswertungsregel aus
\FormulaRefAuto{SemigroupCoordinateMultiplicationMapDef}[def].
\end{proof}

\FormulaThmDeltaK[Der Koordinatenisomorphismus ist ein
Halbgruppenisomorphismus]{%
  \SemigroupCoordinateIso{A}{\star}{e}\vdash
  \SemigroupIso{\SemigroupCoordinateMap{A}{\star}{e}}
    {A,\star}
    {\SemigroupCoordinateCarrier{A}{\star}{e},
     \SemigroupCoordinateProductOp{A}{\star}{e}}%
}{SemigroupCoordinateIsoIsSemigroupIso}{%
  \DeltaRow{Mengen}{A\dsep e\dsep x\dsep y}
  \DeltaRow{Funktionen}{\star\dsep
    \SemigroupCoordinateMap{A}{\star}{e}\dsep
    \SemigroupCoordinateProductOp{A}{\star}{e}}
}
\begin{tabproofwide}
  \proofstepwidestar[1]{\SemigroupCoordinateIso{A}{\star}{e}}{\rA}
  \proofstepwidestar[1]{\SemigroupAlg{A}{\star}}{%
    \FormulaRefAuto{SemigroupCoordinateIsoSourceAxiom}[ax]{1}}
  \proofstepwidestar[1]{e\in A}{%
    \FormulaRefAuto{SemigroupCoordinateIsoBaseElementAxiom}[ax]{1}}
  \proofstepwidestar[1]{%
    \SemigroupAlg{\SemigroupCoordinateCarrier{A}{\star}{e}}
      {\SemigroupCoordinateProductOp{A}{\star}{e}}}{%
    \FormulaRefAuto{SemigroupCoordinateProductSemigroup}[thm]{2,3}}
  \proofstepwidestar[1]{%
    \SemigroupCoordinateMap{A}{\star}{e}\colon
    A\bij\SemigroupCoordinateCarrier{A}{\star}{e}}{%
    \FormulaRefAuto{SemigroupCoordinateIsoBijectiveAxiom}[ax]{1}}
  \proofstepwidestar[1]{%
    \SemigroupCoordinateMap{A}{\star}{e}\colon
    A\to\SemigroupCoordinateCarrier{A}{\star}{e}}{%
    \FormulaRefAuto{SemigroupCoordinateMapFunction}[thm]{2,3}}
  \proofstepwidestarbreak[7]{x\in A}{\rA}
  \proofstepwidestarbreak[8]{y\in A}{\rA}
  \proofstepwidestarbreak[1,7,8]{\begin{aligned}[t]
    &\SemigroupCoordinateMap{A}{\star}{e}(x\star y)\\[-2pt]
    &\quad=\SemigroupCoordinateMap{A}{\star}{e}(x)
      \SemigroupCoordinateProductOp{A}{\star}{e}
      \SemigroupCoordinateMap{A}{\star}{e}(y)
    \end{aligned}}{\FormulaRefAuto{SemigroupCoordinateIsoOperationAxiom}[ax]{1,7,8}}
  \proofstepwidestarbreak[1]{%
    \forall x\in A\,\forall y\in A\,
    \Bigl(
      \SemigroupCoordinateMap{A}{\star}{e}(x\star y)
      =\SemigroupCoordinateMap{A}{\star}{e}(x)
       \SemigroupCoordinateProductOp{A}{\star}{e}
       \SemigroupCoordinateMap{A}{\star}{e}(y)
    \Bigr)}{\rUI{\rRI{7,\rUI{\rRI{8,9}}}}}
  \proofstepwidestar[1]{%
    \SemigroupHom{\SemigroupCoordinateMap{A}{\star}{e}}
      {A,\star}
      {\SemigroupCoordinateCarrier{A}{\star}{e},
       \SemigroupCoordinateProductOp{A}{\star}{e}}}{%
    \FormulaRefAuto{SemigroupHomDef}[def]{2,4,6,10}}
  \proofstepwidestar[1]{%
    \SemigroupIso{\SemigroupCoordinateMap{A}{\star}{e}}
      {A,\star}
      {\SemigroupCoordinateCarrier{A}{\star}{e},
       \SemigroupCoordinateProductOp{A}{\star}{e}}}{%
    \FormulaRefAuto{SemigroupIsoDef}[def]{11,5}}
\end{tabproofwide}

\FormulaThmDeltaK[Idempotenz- und Dreitermkriterium f\"ur den
Koordinatenisomorphismus]{%
  \begin{aligned}[t]
    &\SemigroupAlg{A}{\star}\dsep e\in A\dsep
      \forall x\in A\,(x\star x=x)\dsep{}\\[-2pt]
    &\forall x,y,z\in A\,\bigl((x\star y)\star z=x\star z\bigr)
      \vdash\SemigroupCoordinateIso{A}{\star}{e}
  \end{aligned}%
}{SemigroupCoordinateIsoCriterion}{%
  \DeltaRow{Mengen}{A\dsep e\dsep p\dsep q\dsep r\dsep s\dsep
    x\dsep y\dsep z}
  \DeltaRow{Funktionen}{\star\dsep
    \SemigroupCoordinateMap{A}{\star}{e}\dsep
    \SemigroupCoordinateMultiplicationMap{A}{\star}{e}\dsep
    \SemigroupCoordinateProductOp{A}{\star}{e}}
}
\begin{proof}
Schreibe kurz

\[
 K=\SemigroupCoordinateCarrier{A}{\star}{e},\qquad
 \chi=\SemigroupCoordinateMap{A}{\star}{e},\qquad
 \mu=\SemigroupCoordinateMultiplicationMap{A}{\star}{e}.
\]

Nach
\FormulaRefAuto{SemigroupCoordinateMapFunction}[thm] und
\FormulaRefAuto{SemigroupCoordinateMultiplicationMapFunction}[thm]
sind \(\chi\colon A\to K\) und \(\mu\colon K\to A\) Abbildungen.  F\"ur
\(x\in A\) liefern Assoziativit\"at, Dreitermgesetz und Idempotenz die
vollst\"andig geklammerte Kette
\[
\begin{aligned}
 \mu(\chi(x))
 &=(x\star e)\star(e\star x)\\
 &=\bigl((x\star e)\star e\bigr)\star x\\
 &=(x\star e)\star x\\
 &=x\star x=x.
\end{aligned}
\]
Ist \((p,q)\in K\), so gibt es \(a,b\in A\) mit
\(p=a\star e\) und \(q=e\star b\).  Nach dem Dreitermgesetz,
der Assoziativit\"at und der Idempotenz gilt
\[
 p\star e=p,\qquad e\star q=q,\qquad
 (p\star q)\star e=p\star e,
 \qquad e\star(p\star q)=(e\star p)\star q=e\star q.
\]
Mit den Auswertungsregeln von \(\chi\) und \(\mu\) folgt daher
\[
 \chi(\mu(p,q))
 =\bigl((p\star q)\star e,e\star(p\star q)\bigr)
 =(p,q).
\]
Die Allgemeinheit der eingesetzten Elemente und
\FormulaRefAuto{FunctionExtensionality}[thm] liefern aus den beiden
Gleichungen
\[
 \mu\circ\chi=\Id_A,
 \qquad
 \chi\circ\mu=\Id_K.
\]
Der Satz
\FormulaRefAuto{MutuallyInverseCompanionBijection}[thm]
\emph{(zwei wechselseitige Inversengleichungen liefern eine
Bijektion)} ergibt daher \(\chi\colon A\bij K\).

F\"ur \(x,y\in A\) ergibt sich schlie\ss lich
\[
\begin{aligned}
 \chi(x\star y)
 &=\bigl((x\star y)\star e,e\star(x\star y)\bigr)\\
 &=\bigl(x\star e,(e\star x)\star y\bigr)\\
 &=\bigl(x\star e,e\star y\bigr)
  =\chi(x)\SemigroupCoordinateProductOp{A}{\star}{e}\chi(y).
\end{aligned}
\]
In der zweiten Zeile wurden das Dreitermgesetz und die Assoziativit\"at,
in der dritten Zeile erneut das Dreitermgesetz verwendet.  Die
Quellhalbgruppe, das Basiselement, die soeben bewiesene Bijektivit\"at und
die zuletzt gezeigte Operationserhaltung sind genau die vier Axiome aus
\FormulaRefAuto{SemigroupCoordinateIsoDef}[def].
\end{proof}

\FormulaThmDeltaK[Idempotenz aus einem Koordinatenisomorphismus]{%
  \SemigroupCoordinateIso{A}{\star}{e}\vdash
  \forall x\in A\,(x\star x=x)%
}{SemigroupCoordinateIsoIdempotence}{%
  \DeltaRow{Mengen}{A\dsep e\dsep x}
  \DeltaRow{Funktionen}{\star\dsep
    \SemigroupCoordinateMap{A}{\star}{e}\dsep
    \SemigroupCoordinateProductOp{A}{\star}{e}}
}
\begin{proof}
Der abgeleitete Halbgruppenisomorphismus aus
\FormulaRefAuto{SemigroupCoordinateIsoIsSemigroupIso}[thm] ist injektiv
und erh\"alt die Operation.  F\"ur \(x\in A\) ist sein Bild ein Element
des Koordinatentr\"agers.  Nach
\FormulaRefAuto{SemigroupCoordinateProductIdempotence}[thm] gilt daher
\[
 \begin{aligned}
 &\SemigroupCoordinateMap{A}{\star}{e}(x\star x)\\
 &\quad=
 \SemigroupCoordinateMap{A}{\star}{e}(x)
 \SemigroupCoordinateProductOp{A}{\star}{e}
 \SemigroupCoordinateMap{A}{\star}{e}(x)
 =\SemigroupCoordinateMap{A}{\star}{e}(x).
 \end{aligned}
\]
Die Injektivit\"at der Koordinatenabbildung liefert \(x\star x=x\);
die Allgemeinheit
von \(x\) ergibt die Behauptung.
\end{proof}

\FormulaThmDeltaK[Dreitermgesetz aus einem Koordinatenisomorphismus]{%
  \begin{aligned}[t]
    &\SemigroupCoordinateIso{A}{\star}{e}\vdash{}\\[-2pt]
    &\forall x,y,z\in A\,
      \bigl((x\star y)\star z=x\star z\bigr)
  \end{aligned}%
}{SemigroupCoordinateIsoThreeTermLaw}{%
  \DeltaRow{Mengen}{A\dsep e\dsep x\dsep y\dsep z}
  \DeltaRow{Funktionen}{\star\dsep
    \SemigroupCoordinateMap{A}{\star}{e}\dsep
    \SemigroupCoordinateProductOp{A}{\star}{e}}
}
\begin{proof}
Wie im vorigen Satz ist die Koordinatenabbildung ein injektiver
Homomorphismus.  F\"ur \(x,y,z\in A\) liefert das Dreitermgesetz des
Koordinatenprodukts
\[
\begin{aligned}
 &\SemigroupCoordinateMap{A}{\star}{e}
   \bigl((x\star y)\star z\bigr)\\
 &\quad=\Bigl(
   \SemigroupCoordinateMap{A}{\star}{e}(x)
   \SemigroupCoordinateProductOp{A}{\star}{e}
   \SemigroupCoordinateMap{A}{\star}{e}(y)
   \Bigr)
   \SemigroupCoordinateProductOp{A}{\star}{e}
   \SemigroupCoordinateMap{A}{\star}{e}(z)\\
 &\quad=\SemigroupCoordinateMap{A}{\star}{e}(x)
   \SemigroupCoordinateProductOp{A}{\star}{e}
   \SemigroupCoordinateMap{A}{\star}{e}(z)
   =\SemigroupCoordinateMap{A}{\star}{e}(x\star z).
\end{aligned}
\]
Injektivit\"at und Allgemeinheit der drei Elemente ergeben die Aussage.
\end{proof}

\FormulaThmDeltaK[Rechtecksidentit\"at aus einem
Koordinatenisomorphismus]{%
  \begin{aligned}[t]
    &\SemigroupCoordinateIso{A}{\star}{e}\vdash{}\\[-2pt]
    &\forall x,y\in A\,\bigl((x\star y)\star x=x\bigr)
  \end{aligned}%
}{SemigroupCoordinateIsoRectangularIdentity}{%
  \DeltaRow{Mengen}{A\dsep e\dsep x\dsep y}
  \DeltaRow{Funktionen}{\star}
}
\begin{proof}
Nach \FormulaRefAuto{SemigroupCoordinateIsoThreeTermLaw}[thm] und
\FormulaRefAuto{SemigroupCoordinateIsoIdempotence}[thm] gilt f\"ur
\(x,y\in A\)
\[
 (x\star y)\star x=x\star x=x.
\]
Die Allgemeinheit von \(x,y\) liefert die Behauptung.
\end{proof}

\section{Morphismen von Potenzhalbgruppen}

\FormulaThmDeltaK[Die Einermengenkopie ist zur Ausgangshalbgruppe
isomorph]{%
  \SemigroupAlg{A}{\star}\vdash
  \SemigroupIso{\eta_A}
    {A,\star}{\operatorname{Sing}(A),\star_{\mathcal P}}%
}{SemigroupSingletonCopyIsomorphism}{%
  \DeltaRow{Mengen}{A\dsep a\dsep b\dsep X\dsep z}
  \DeltaRow{Funktionen}{\star\dsep\star_{\mathcal P}\dsep\eta_A}
}
\begin{proof}
Seien \(a,b\in A\). Ihre Einermengen liegen in \(\PowNE A\). Ferner ist
\(a\star b\in A\) nach \FormulaRefAuto{SemigroupClosureAxiom}[ax]. Nach
\FormulaRefAuto{PowerSemigroupProductMembership}[thm] gilt für jedes \(z\)
\[
  z\in\{a\}\star_{\mathcal P}\{b\}
  \quad\Longleftrightarrow\quad
  z=a\star b
  \quad\Longleftrightarrow\quad
  z\in\{a\star b\}.
\]
Die Extensionalität liefert somit
\[
  \{a\}\star_{\mathcal P}\{b\}=\{a\star b\}.       \tag{1}
\]
Da \(a\star b\in A\), liegt die rechte Seite von (1) in
\(\operatorname{Sing}(A)\). Damit ist
\(\operatorname{Sing}(A)\) unter \(\star_{\mathcal P}\) abgeschlossen.
Jedes \(X\in\operatorname{Sing}(A)\) ist eine nichtleere Teilmenge von
\(A\), also gilt \(\operatorname{Sing}(A)\subseteq\PowNE A\). Die
Assoziativität aus \FormulaRefAuto{NonemptyPowerSemigroup}[thm] schränkt sich
daher auf \(\operatorname{Sing}(A)\) ein. Nach
\FormulaRefAuto{SemigroupDef}[def] ist folglich
\(\SemigroupAlg{\operatorname{Sing}(A)}{\star_{\mathcal P}}\) erfüllt.

Nach \FormulaRefAuto{SemigroupSingletonCopyDef}[def] gilt
\(\eta_A(a)=\{a\}\). Gleichung (1) ist deshalb genau
\[
  \eta_A(a\star b)
   =\eta_A(a)\star_{\mathcal P}\eta_A(b).
\]
Zusammen mit der Typisierung von \(\eta_A\), der vorausgesetzten
Quellhalbgruppe und der eben bewiesenen Zielhalbgruppe ergibt
\FormulaRefAuto{SemigroupHomDef}[def], dass \(\eta_A\) ein
Halbgruppenhomomorphismus ist. Aus \(\eta_A(a)=\eta_A(b)\), also
\(\{a\}=\{b\}\), folgt \(a=b\); und jedes Element von
\(\operatorname{Sing}(A)\) hat nach seiner Definition die Form \(\eta_A(a)\)
für ein \(a\in A\). Somit ist \(\eta_A\) bijektiv. Mit
\FormulaRefAuto{SemigroupIsoDef}[def] folgt die behauptete Isomorphie.
\end{proof}

\FormulaThmDeltaK[Ein Homomorphismus erhält das Mengenprodukt]{%
  \begin{aligned}[t]
    &\SemigroupHom{f}{A,\star}{B,\diamond}\dsep
      X,Y\in\PowNE{A}\vdash{}\\[-2pt]
    &f[X\star_{\mathcal P}Y]
      =f[X]\diamond_{\mathcal P}f[Y]
  \end{aligned}%
}{SemigroupHomPreservesPowerProduct}{%
  \DeltaRow{Mengen}{A\dsep B\dsep X\dsep Y\dsep
    u\dsep v\dsep w\dsep x\dsep y\dsep z}
  \DeltaRow{Funktionen}{f\dsep\star\dsep\diamond\dsep
    \star_{\mathcal P}\dsep\diamond_{\mathcal P}}
}
\begin{tabproofsplitwide}
  \proofpartwideindR[Bilder fester Faktoren]{%
    \begin{aligned}[t]
      &\SemigroupHom{f}{A,\star}{B,\diamond}\dsep
        X,Y\in\PowNE{A}\dsep x\in X\dsep y\in Y\vdash{}\\[-2pt]
      &f(x)\diamond f(y)\in f[X]\diamond_{\mathcal P}f[Y]
    \end{aligned}
  }{%
    \begin{aligned}[t]
      &\SemigroupHom{f}{A,\star}{B,\diamond}\dsep
        X,Y\in\PowNE{A}\dsep x\in X\dsep y\in Y\vdash{}\\[-2pt]
      &f(x)\diamond f(y)\in f[X]\diamond_{\mathcal P}f[Y]
    \end{aligned}
  }[PowerHomImagesOfFactors]
    \proofstepwidestar[1]{\SemigroupHom{f}{A,\star}{B,\diamond}}{\rA}
    \proofstepwidestar[2]{X\in\PowNE{A}}{\rA}
    \proofstepwidestar[3]{Y\in\PowNE{A}}{\rA}
    \proofstepwidestar[4]{x\in X}{\rA}
    \proofstepwidestar[5]{y\in Y}{\rA}
    \proofstepwidestar[1]{\SemigroupAlg{B}{\diamond}}{%
      \FormulaRefAuto{SemigroupHomTargetAxiom}[ax]{1}}
    \proofstepwidestar[1]{f\colon A\to B}{%
      \FormulaRefAuto{SemigroupHomFunctionAxiom}[ax]{1}}
    \proofstepwidestar[1,2]{f[X]\in\PowNE{B}}{%
      \FormulaRefAuto{NonemptyDirectImage}[thm]{7,2}}
    \proofstepwidestar[1,3]{f[Y]\in\PowNE{B}}{%
      \FormulaRefAuto{NonemptyDirectImage}[thm]{7,3}}
    \proofstepwidestar[2]{X\subseteq A\land X\neq\varnothing}{%
      \FormulaRefAuto{NonemptyPowersetMembership}[thm]{2}}
    \proofstepwidestar[2]{X\subseteq A}{%
      \rAEa{10}}
    \proofstepwidestar[3]{Y\subseteq A\land Y\neq\varnothing}{%
      \FormulaRefAuto{NonemptyPowersetMembership}[thm]{3}}
    \proofstepwidestar[3]{Y\subseteq A}{%
      \rAEa{12}}
    \proofstepwidestar[1,2,4]{f(x)\in f[X]}{%
      \FormulaRefAuto{C\subseteq A\dsep F\colon A\to B\dsep
        x\in C\vdash F(x)\in F[C]}[thm]{11,7,4}}
    \proofstepwidestar[1,3,5]{f(y)\in f[Y]}{%
      \FormulaRefAuto{C\subseteq A\dsep F\colon A\to B\dsep
        x\in C\vdash F(x)\in F[C]}[thm]{13,7,5}}
    \proofstepwidestar[1,2,3,4,5]{%
      f(x)\diamond f(y)\in f[X]\diamond_{\mathcal P}f[Y]}{%
      \FormulaRefAuto{PowerSemigroupProductContains}[thm]{6,8,9,14,15}}
  \closeproofpartwideindR

  \proofpartwideindR[Vorwärtsrichtung mit festen Zeugen]{%
    \begin{aligned}[t]
      &\SemigroupHom{f}{A,\star}{B,\diamond}\dsep X,Y\in\PowNE{A}
        \dsep x\in X\dsep y\in Y\dsep{}\\[-2pt]
      &w=x\star y\dsep z=f(w)
        \vdash z\in f[X]\diamond_{\mathcal P}f[Y]
    \end{aligned}
  }{%
    \begin{aligned}[t]
      &\SemigroupHom{f}{A,\star}{B,\diamond}\dsep X,Y\in\PowNE{A}
        \dsep x\in X\dsep y\in Y\dsep{}\\[-2pt]
      &w=x\star y\dsep z=f(w)
        \vdash z\in f[X]\diamond_{\mathcal P}f[Y]
    \end{aligned}
  }[PowerHomForwardFixedWitnesses]
    \proofstepwidestar[1]{\SemigroupHom{f}{A,\star}{B,\diamond}}{\rA}
    \proofstepwidestar[2]{X\in\PowNE{A}}{\rA}
    \proofstepwidestar[3]{Y\in\PowNE{A}}{\rA}
    \proofstepwidestar[4]{x\in X}{\rA}
    \proofstepwidestar[5]{y\in Y}{\rA}
    \proofstepwidestar[6]{w=x\star y}{\rA}
    \proofstepwidestar[7]{z=f(w)}{\rA}
    \proofstepwidestar[2,4]{x\in A}{%
      \FormulaRefAuto{PowerSemigroupCarrierElement}[thm]{2,4}}
    \proofstepwidestar[3,5]{y\in A}{%
      \FormulaRefAuto{PowerSemigroupCarrierElement}[thm]{3,5}}
    \proofstepwidestar[6,7]{z=f(x\star y)}{\rIE{6,7}}
    \proofstepwidestar[1,2,3,4,5]{%
      f(x\star y)=f(x)\diamond f(y)}{%
      \FormulaRefAuto{SemigroupHomOperationAxiom}[ax]{1,8,9}}
    \proofstepwidestar[1,2,3,4,5,6,7]{z=f(x)\diamond f(y)}{%
      \rChain{10,11}}
    \proofstepwidestar[1,2,3,4,5]{%
      f(x)\diamond f(y)\in f[X]\diamond_{\mathcal P}f[Y]}{%
      \FormulaRefAuto{PowerHomImagesOfFactors}[thm]{1,2,3,4,5}}
    \proofstepwidestar[1,2,3,4,5,6,7]{%
      z\in f[X]\diamond_{\mathcal P}f[Y]}{\rIE{12,13}}
  \closeproofpartwideindR

  \proofpartwideindR[Vorwärtsrichtung mit einem Produktzeugen]{%
    \begin{aligned}[t]
      &\SemigroupHom{f}{A,\star}{B,\diamond}\dsep X,Y\in\PowNE{A}
        \dsep w\in X\star_{\mathcal P}Y\dsep z=f(w)\vdash{}\\[-2pt]
      &z\in f[X]\diamond_{\mathcal P}f[Y]
    \end{aligned}
  }{%
    \begin{aligned}[t]
      &\SemigroupHom{f}{A,\star}{B,\diamond}\dsep X,Y\in\PowNE{A}
        \dsep w\in X\star_{\mathcal P}Y\dsep z=f(w)\vdash{}\\[-2pt]
      &z\in f[X]\diamond_{\mathcal P}f[Y]
    \end{aligned}
  }[PowerHomForwardProductWitness]
    \proofstepwidestar[1]{\SemigroupHom{f}{A,\star}{B,\diamond}}{\rA}
    \proofstepwidestar[2]{X\in\PowNE{A}}{\rA}
    \proofstepwidestar[3]{Y\in\PowNE{A}}{\rA}
    \proofstepwidestar[4]{w\in X\star_{\mathcal P}Y}{\rA}
    \proofstepwidestar[5]{z=f(w)}{\rA}
    \proofstepwidestar[1]{\SemigroupAlg{A}{\star}}{%
      \FormulaRefAuto{SemigroupHomSourceAxiom}[ax]{1}}
    \proofstepwidestar[1,2,3,4]{%
      \exists x\in X\,\exists y\in Y\,(w=x\star y)}{%
      \FormulaRefAuto{PowerSemigroupProductWitnesses}[thm]{6,2,3,4}}
    \proofstepwidestar[8]{x\in X\land\exists y\in Y\,(w=x\star y)}{\rA}
    \proofstepwidestar[8]{x\in X}{\rAEa{8}}
    \proofstepwidestar[8]{\exists y\in Y\,(w=x\star y)}{\rAEb{8}}
    \proofstepwidestar[11]{y\in Y\land w=x\star y}{\rA}
    \proofstepwidestar[11]{y\in Y}{\rAEa{11}}
    \proofstepwidestar[11]{w=x\star y}{\rAEb{11}}
    \proofstepwidestar[1,2,3,5,8,11]{%
      z\in f[X]\diamond_{\mathcal P}f[Y]}{%
      \FormulaRefAuto{PowerHomForwardFixedWitnesses}[thm]{%
        1,2,3,9,12,13,5}}
    \proofstepwidestar[1,2,3,5,8]{%
      z\in f[X]\diamond_{\mathcal P}f[Y]}{\rEE{10,11,14}}
    \proofstepwidestar[1,2,3,4,5]{%
      z\in f[X]\diamond_{\mathcal P}f[Y]}{\rEE{7,8,15}}
  \closeproofpartwideindR

  \proofpartwideindR[Vorwärtsrichtung für Bildelemente]{%
    \begin{aligned}[t]
      &\SemigroupHom{f}{A,\star}{B,\diamond}\dsep X,Y\in\PowNE{A}
        \dsep z\in f[X\star_{\mathcal P}Y]\vdash{}\\[-2pt]
      &z\in f[X]\diamond_{\mathcal P}f[Y]
    \end{aligned}
  }{%
    \begin{aligned}[t]
      &\SemigroupHom{f}{A,\star}{B,\diamond}\dsep X,Y\in\PowNE{A}
        \dsep z\in f[X\star_{\mathcal P}Y]\vdash{}\\[-2pt]
      &z\in f[X]\diamond_{\mathcal P}f[Y]
    \end{aligned}
  }[PowerHomForwardElement]
    \proofstepwidestar[1]{\SemigroupHom{f}{A,\star}{B,\diamond}}{\rA}
    \proofstepwidestar[2]{X\in\PowNE{A}}{\rA}
    \proofstepwidestar[3]{Y\in\PowNE{A}}{\rA}
    \proofstepwidestar[4]{z\in f[X\star_{\mathcal P}Y]}{\rA}
    \proofstepwidestar[1]{\SemigroupAlg{A}{\star}}{%
      \FormulaRefAuto{SemigroupHomSourceAxiom}[ax]{1}}
    \proofstepwidestar[1]{f\colon A\to B}{%
      \FormulaRefAuto{SemigroupHomFunctionAxiom}[ax]{1}}
    \proofstepwidestar[1,2,3]{X\star_{\mathcal P}Y\subseteq A}{%
      \FormulaRefAuto{PowerSemigroupProductSubset}[thm]{5,2,3}}
    \proofstepwidestar[1,2,3,4]{%
      \exists w\in X\star_{\mathcal P}Y\,(z=f(w))}{%
      \FormulaRefAuto{C\subseteq A\dsep F\colon A\to B\dsep
        y\in F[C]\vdash\exists x\in C\,y=F(x)}[thm]{7,6,4}}
    \proofstepwidestar[9]{w\in X\star_{\mathcal P}Y\land z=f(w)}{\rA}
    \proofstepwidestar[9]{w\in X\star_{\mathcal P}Y}{\rAEa{9}}
    \proofstepwidestar[9]{z=f(w)}{\rAEb{9}}
    \proofstepwidestar[1,2,3,9]{z\in f[X]\diamond_{\mathcal P}f[Y]}{%
      \FormulaRefAuto{PowerHomForwardProductWitness}[thm]{1,2,3,10,11}}
    \proofstepwidestar[1,2,3,4]{z\in f[X]\diamond_{\mathcal P}f[Y]}{%
      \rEE{8,9,12}}
  \closeproofpartwideindR

  \proofpartwideindR[Erste Teilmengenrichtung]{%
    \begin{aligned}[t]
      &\SemigroupHom{f}{A,\star}{B,\diamond}\dsep X,Y\in\PowNE{A}
        \vdash{}\\[-2pt]
      &f[X\star_{\mathcal P}Y]
        \subseteq f[X]\diamond_{\mathcal P}f[Y]
    \end{aligned}
  }{%
    \begin{aligned}[t]
      &\SemigroupHom{f}{A,\star}{B,\diamond}\dsep X,Y\in\PowNE{A}
        \vdash{}\\[-2pt]
      &f[X\star_{\mathcal P}Y]
        \subseteq f[X]\diamond_{\mathcal P}f[Y]
    \end{aligned}
  }[PowerHomForwardSubset]
    \proofstepwidestar[1]{\SemigroupHom{f}{A,\star}{B,\diamond}}{\rA}
    \proofstepwidestar[2]{X\in\PowNE{A}}{\rA}
    \proofstepwidestar[3]{Y\in\PowNE{A}}{\rA}
    \proofstepwidestar[4]{z\in f[X\star_{\mathcal P}Y]}{\rA}
    \proofstepwidestar[1,2,3,4]{z\in f[X]\diamond_{\mathcal P}f[Y]}{%
      \FormulaRefAuto{PowerHomForwardElement}[thm]{1,2,3,4}}
    \proofstepwidestar[1,2,3]{%
      f[X\star_{\mathcal P}Y]\subseteq
      f[X]\diamond_{\mathcal P}f[Y]}{%
      \FormulaRefAuto{SubsetIntroductionRule}[thm]{[4]\vdots 5}}
  \closeproofpartwideindR

  \proofpartwideindR[Rückwärtsrichtung der Gleichheit]{%
    \begin{aligned}[t]
      &\SemigroupHom{f}{A,\star}{B,\diamond}\dsep x,y\in A\dsep
        u=f(x)\dsep v=f(y)\dsep z=u\diamond v\vdash{}\\[-2pt]
      &z=f(x\star y)
    \end{aligned}
  }{%
    \begin{aligned}[t]
      &\SemigroupHom{f}{A,\star}{B,\diamond}\dsep x,y\in A\dsep
        u=f(x)\dsep v=f(y)\dsep z=u\diamond v\vdash{}\\[-2pt]
      &z=f(x\star y)
    \end{aligned}
  }[PowerHomBackwardEquality]
    \proofstepwidestar[1]{\SemigroupHom{f}{A,\star}{B,\diamond}}{\rA}
    \proofstepwidestar[2]{x\in A}{\rA}
    \proofstepwidestar[3]{y\in A}{\rA}
    \proofstepwidestar[4]{u=f(x)}{\rA}
    \proofstepwidestar[5]{v=f(y)}{\rA}
    \proofstepwidestar[6]{z=u\diamond v}{\rA}
    \proofstepwidestar[1,2,3]{f(x\star y)=f(x)\diamond f(y)}{%
      \FormulaRefAuto{SemigroupHomOperationAxiom}[ax]{1,2,3}}
    \proofstepwidestar[1,2,3]{f(x)\diamond f(y)=f(x\star y)}{%
      \FormulaRefAuto{a=b\vdash b=a}[thm]{7}}
    \proofstepwidestar[4,6]{z=f(x)\diamond v}{\rIE{4,6}}
    \proofstepwidestar[4,5,6]{z=f(x)\diamond f(y)}{\rIE{5,9}}
    \proofstepwidestar[1,2,3,4,5,6]{z=f(x\star y)}{\rChain{10,8}}
  \closeproofpartwideindR

  \proofpartwideindR[Rückwärtsrichtung mit festen Zeugen]{%
    \begin{aligned}[t]
      &\SemigroupHom{f}{A,\star}{B,\diamond}\dsep X,Y\in\PowNE{A}
        \dsep x\in X\dsep y\in Y\dsep{}\\[-2pt]
      &u=f(x)\dsep v=f(y)\dsep z=u\diamond v
        \vdash z\in f[X\star_{\mathcal P}Y]
    \end{aligned}
  }{%
    \begin{aligned}[t]
      &\SemigroupHom{f}{A,\star}{B,\diamond}\dsep X,Y\in\PowNE{A}
        \dsep x\in X\dsep y\in Y\dsep{}\\[-2pt]
      &u=f(x)\dsep v=f(y)\dsep z=u\diamond v
        \vdash z\in f[X\star_{\mathcal P}Y]
    \end{aligned}
  }[PowerHomBackwardFixedWitnesses]
    \proofstepwidestar[1]{\SemigroupHom{f}{A,\star}{B,\diamond}}{\rA}
    \proofstepwidestar[2]{X\in\PowNE{A}}{\rA}
    \proofstepwidestar[3]{Y\in\PowNE{A}}{\rA}
    \proofstepwidestar[4]{x\in X}{\rA}
    \proofstepwidestar[5]{y\in Y}{\rA}
    \proofstepwidestar[6]{u=f(x)}{\rA}
    \proofstepwidestar[7]{v=f(y)}{\rA}
    \proofstepwidestar[8]{z=u\diamond v}{\rA}
    \proofstepwidestar[1]{\SemigroupAlg{A}{\star}}{%
      \FormulaRefAuto{SemigroupHomSourceAxiom}[ax]{1}}
    \proofstepwidestar[1]{f\colon A\to B}{%
      \FormulaRefAuto{SemigroupHomFunctionAxiom}[ax]{1}}
    \proofstepwidestar[2,4]{x\in A}{%
      \FormulaRefAuto{PowerSemigroupCarrierElement}[thm]{2,4}}
    \proofstepwidestar[3,5]{y\in A}{%
      \FormulaRefAuto{PowerSemigroupCarrierElement}[thm]{3,5}}
    \proofstepwidestar[1,2,3,4,5]{x\star y\in X\star_{\mathcal P}Y}{%
      \FormulaRefAuto{PowerSemigroupProductContains}[thm]{9,2,3,4,5}}
    \proofstepwidestar[1,2,3]{X\star_{\mathcal P}Y\subseteq A}{%
      \FormulaRefAuto{PowerSemigroupProductSubset}[thm]{9,2,3}}
    \proofstepwidestar[1,2,3,4,5]{f(x\star y)\in f[X\star_{\mathcal P}Y]}{%
      \FormulaRefAuto{C\subseteq A\dsep F\colon A\to B\dsep
        x\in C\vdash F(x)\in F[C]}[thm]{14,10,13}}
    \proofstepwidestar[1,2,3,4,5,6,7,8]{z=f(x\star y)}{%
      \FormulaRefAuto{PowerHomBackwardEquality}[thm]{1,11,12,6,7,8}}
    \proofstepwidestar[1,2,3,4,5,6,7,8]{z\in f[X\star_{\mathcal P}Y]}{%
      \rIE{16,15}}
  \closeproofpartwideindR

  \proofpartwideindR[Rückwärtsrichtung mit rechtem Bildzeugen]{%
    \begin{aligned}[t]
      &\SemigroupHom{f}{A,\star}{B,\diamond}\dsep X,Y\in\PowNE{A}
        \dsep x\in X\dsep u=f(x)\dsep v\in f[Y]\dsep{}\\[-2pt]
      &z=u\diamond v\vdash z\in f[X\star_{\mathcal P}Y]
    \end{aligned}
  }{%
    \begin{aligned}[t]
      &\SemigroupHom{f}{A,\star}{B,\diamond}\dsep X,Y\in\PowNE{A}
        \dsep x\in X\dsep u=f(x)\dsep v\in f[Y]\dsep{}\\[-2pt]
      &z=u\diamond v\vdash z\in f[X\star_{\mathcal P}Y]
    \end{aligned}
  }[PowerHomBackwardRightWitness]
    \proofstepwidestar[1]{\SemigroupHom{f}{A,\star}{B,\diamond}}{\rA}
    \proofstepwidestar[2]{X\in\PowNE{A}}{\rA}
    \proofstepwidestar[3]{Y\in\PowNE{A}}{\rA}
    \proofstepwidestar[4]{x\in X}{\rA}
    \proofstepwidestar[5]{u=f(x)}{\rA}
    \proofstepwidestar[6]{v\in f[Y]}{\rA}
    \proofstepwidestar[7]{z=u\diamond v}{\rA}
    \proofstepwidestar[1]{f\colon A\to B}{%
      \FormulaRefAuto{SemigroupHomFunctionAxiom}[ax]{1}}
    \proofstepwidestar[3]{Y\subseteq A\land Y\neq\varnothing}{%
      \FormulaRefAuto{NonemptyPowersetMembership}[thm]{3}}
    \proofstepwidestar[3]{Y\subseteq A}{%
      \rAEa{9}}
    \proofstepwidestar[1,3,6]{\exists y\in Y\,(v=f(y))}{%
      \FormulaRefAuto{C\subseteq A\dsep F\colon A\to B\dsep
        y\in F[C]\vdash\exists x\in C\,y=F(x)}[thm]{10,8,6}}
    \proofstepwidestar[12]{y\in Y\land v=f(y)}{\rA}
    \proofstepwidestar[12]{y\in Y}{\rAEa{12}}
    \proofstepwidestar[12]{v=f(y)}{\rAEb{12}}
    \proofstepwidestar[1,2,3,4,5,7,12]{z\in f[X\star_{\mathcal P}Y]}{%
      \FormulaRefAuto{PowerHomBackwardFixedWitnesses}[thm]{%
        1,2,3,4,13,5,14,7}}
    \proofstepwidestar[1,2,3,4,5,6,7]{z\in f[X\star_{\mathcal P}Y]}{%
      \rEE{11,12,15}}
  \closeproofpartwideindR

  \proofpartwideindR[Rückwärtsrichtung mit zwei Bildzeugen]{%
    \begin{aligned}[t]
      &\SemigroupHom{f}{A,\star}{B,\diamond}\dsep X,Y\in\PowNE{A}
        \dsep u\in f[X]\dsep v\in f[Y]\dsep z=u\diamond v\vdash{}\\[-2pt]
      &z\in f[X\star_{\mathcal P}Y]
    \end{aligned}
  }{%
    \begin{aligned}[t]
      &\SemigroupHom{f}{A,\star}{B,\diamond}\dsep X,Y\in\PowNE{A}
        \dsep u\in f[X]\dsep v\in f[Y]\dsep z=u\diamond v\vdash{}\\[-2pt]
      &z\in f[X\star_{\mathcal P}Y]
    \end{aligned}
  }[PowerHomBackwardImageWitnesses]
    \proofstepwidestar[1]{\SemigroupHom{f}{A,\star}{B,\diamond}}{\rA}
    \proofstepwidestar[2]{X\in\PowNE{A}}{\rA}
    \proofstepwidestar[3]{Y\in\PowNE{A}}{\rA}
    \proofstepwidestar[4]{u\in f[X]}{\rA}
    \proofstepwidestar[5]{v\in f[Y]}{\rA}
    \proofstepwidestar[6]{z=u\diamond v}{\rA}
    \proofstepwidestar[1]{f\colon A\to B}{%
      \FormulaRefAuto{SemigroupHomFunctionAxiom}[ax]{1}}
    \proofstepwidestar[2]{X\subseteq A\land X\neq\varnothing}{%
      \FormulaRefAuto{NonemptyPowersetMembership}[thm]{2}}
    \proofstepwidestar[2]{X\subseteq A}{%
      \rAEa{8}}
    \proofstepwidestar[1,2,4]{\exists x\in X\,(u=f(x))}{%
      \FormulaRefAuto{C\subseteq A\dsep F\colon A\to B\dsep
        y\in F[C]\vdash\exists x\in C\,y=F(x)}[thm]{9,7,4}}
    \proofstepwidestar[11]{x\in X\land u=f(x)}{\rA}
    \proofstepwidestar[11]{x\in X}{\rAEa{11}}
    \proofstepwidestar[11]{u=f(x)}{\rAEb{11}}
    \proofstepwidestar[1,2,3,5,6,11]{z\in f[X\star_{\mathcal P}Y]}{%
      \FormulaRefAuto{PowerHomBackwardRightWitness}[thm]{%
        1,2,3,12,13,5,6}}
    \proofstepwidestar[1,2,3,4,5,6]{z\in f[X\star_{\mathcal P}Y]}{%
      \rEE{10,11,14}}
  \closeproofpartwideindR

  \proofpartwideindR[Auflösung des zweiten Produktzeugen]{%
    \begin{aligned}[t]
      &\SemigroupHom{f}{A,\star}{B,\diamond}\dsep X,Y\in\PowNE{A}
        \dsep u\in f[X]\dsep{}\\[-2pt]
      &\exists v\in f[Y]\,(z=u\diamond v)
        \vdash z\in f[X\star_{\mathcal P}Y]
    \end{aligned}
  }{%
    \begin{aligned}[t]
      &\SemigroupHom{f}{A,\star}{B,\diamond}\dsep X,Y\in\PowNE{A}
        \dsep u\in f[X]\dsep{}\\[-2pt]
      &\exists v\in f[Y]\,(z=u\diamond v)
        \vdash z\in f[X\star_{\mathcal P}Y]
    \end{aligned}
  }[PowerHomBackwardSecondProductWitness]
    \proofstepwidestar[1]{\SemigroupHom{f}{A,\star}{B,\diamond}}{\rA}
    \proofstepwidestar[2]{X\in\PowNE{A}}{\rA}
    \proofstepwidestar[3]{Y\in\PowNE{A}}{\rA}
    \proofstepwidestar[4]{u\in f[X]}{\rA}
    \proofstepwidestar[5]{\exists v\in f[Y]\,(z=u\diamond v)}{\rA}
    \proofstepwidestar[6]{v\in f[Y]\land z=u\diamond v}{\rA}
    \proofstepwidestar[6]{v\in f[Y]}{\rAEa{6}}
    \proofstepwidestar[6]{z=u\diamond v}{\rAEb{6}}
    \proofstepwidestar[1,2,3,4,6]{z\in f[X\star_{\mathcal P}Y]}{%
      \FormulaRefAuto{PowerHomBackwardImageWitnesses}[thm]{1,2,3,4,7,8}}
    \proofstepwidestar[1,2,3,4,5]{z\in f[X\star_{\mathcal P}Y]}{%
      \rEE{5,6,9}}
  \closeproofpartwideindR

  \proofpartwideindR[Rückwärtsrichtung für Produktelemente]{%
    \begin{aligned}[t]
      &\SemigroupHom{f}{A,\star}{B,\diamond}\dsep X,Y\in\PowNE{A}
        \dsep z\in f[X]\diamond_{\mathcal P}f[Y]\vdash{}\\[-2pt]
      &z\in f[X\star_{\mathcal P}Y]
    \end{aligned}
  }{%
    \begin{aligned}[t]
      &\SemigroupHom{f}{A,\star}{B,\diamond}\dsep X,Y\in\PowNE{A}
        \dsep z\in f[X]\diamond_{\mathcal P}f[Y]\vdash{}\\[-2pt]
      &z\in f[X\star_{\mathcal P}Y]
    \end{aligned}
  }[PowerHomBackwardElement]
    \proofstepwidestar[1]{\SemigroupHom{f}{A,\star}{B,\diamond}}{\rA}
    \proofstepwidestar[2]{X\in\PowNE{A}}{\rA}
    \proofstepwidestar[3]{Y\in\PowNE{A}}{\rA}
    \proofstepwidestar[4]{z\in f[X]\diamond_{\mathcal P}f[Y]}{\rA}
    \proofstepwidestar[1]{\SemigroupAlg{B}{\diamond}}{%
      \FormulaRefAuto{SemigroupHomTargetAxiom}[ax]{1}}
    \proofstepwidestar[1]{f\colon A\to B}{%
      \FormulaRefAuto{SemigroupHomFunctionAxiom}[ax]{1}}
    \proofstepwidestar[1,2]{f[X]\in\PowNE{B}}{%
      \FormulaRefAuto{NonemptyDirectImage}[thm]{6,2}}
    \proofstepwidestar[1,3]{f[Y]\in\PowNE{B}}{%
      \FormulaRefAuto{NonemptyDirectImage}[thm]{6,3}}
    \proofstepwidestar[1,2,3,4]{%
      \exists u\in f[X]\,\exists v\in f[Y]\,(z=u\diamond v)}{%
      \FormulaRefAuto{PowerSemigroupProductWitnesses}[thm]{5,7,8,4}}
    \proofstepwidestar[10]{%
      u\in f[X]\land\exists v\in f[Y]\,(z=u\diamond v)}{\rA}
    \proofstepwidestar[10]{u\in f[X]}{\rAEa{10}}
    \proofstepwidestar[10]{\exists v\in f[Y]\,(z=u\diamond v)}{\rAEb{10}}
    \proofstepwidestar[1,2,3,10]{z\in f[X\star_{\mathcal P}Y]}{%
      \FormulaRefAuto{PowerHomBackwardSecondProductWitness}[thm]{%
        1,2,3,11,12}}
    \proofstepwidestar[1,2,3,4]{z\in f[X\star_{\mathcal P}Y]}{%
      \rEE{9,10,13}}
  \closeproofpartwideindR

  \proofpartwideindR[Zweite Teilmengenrichtung]{%
    \begin{aligned}[t]
      &\SemigroupHom{f}{A,\star}{B,\diamond}\dsep X,Y\in\PowNE{A}
        \vdash{}\\[-2pt]
      &f[X]\diamond_{\mathcal P}f[Y]
        \subseteq f[X\star_{\mathcal P}Y]
    \end{aligned}
  }{%
    \begin{aligned}[t]
      &\SemigroupHom{f}{A,\star}{B,\diamond}\dsep X,Y\in\PowNE{A}
        \vdash{}\\[-2pt]
      &f[X]\diamond_{\mathcal P}f[Y]
        \subseteq f[X\star_{\mathcal P}Y]
    \end{aligned}
  }[PowerHomBackwardSubset]
    \proofstepwidestar[1]{\SemigroupHom{f}{A,\star}{B,\diamond}}{\rA}
    \proofstepwidestar[2]{X\in\PowNE{A}}{\rA}
    \proofstepwidestar[3]{Y\in\PowNE{A}}{\rA}
    \proofstepwidestar[4]{z\in f[X]\diamond_{\mathcal P}f[Y]}{\rA}
    \proofstepwidestar[1,2,3,4]{z\in f[X\star_{\mathcal P}Y]}{%
      \FormulaRefAuto{PowerHomBackwardElement}[thm]{1,2,3,4}}
    \proofstepwidestar[1,2,3]{%
      f[X]\diamond_{\mathcal P}f[Y]\subseteq
      f[X\star_{\mathcal P}Y]}{%
      \FormulaRefAuto{SubsetIntroductionRule}[thm]{[4]\vdots 5}}
  \closeproofpartwideindR

  \proofpartwide{Schluss}
    \proofstepwidestar[1]{\SemigroupHom{f}{A,\star}{B,\diamond}}{\rA}
    \proofstepwidestar[2]{X\in\PowNE{A}}{\rA}
    \proofstepwidestar[3]{Y\in\PowNE{A}}{\rA}
    \proofstepwidestar[1,2,3]{%
      f[X\star_{\mathcal P}Y]\subseteq
      f[X]\diamond_{\mathcal P}f[Y]}{%
      \FormulaRefAuto{PowerHomForwardSubset}[thm]{1,2,3}}
    \proofstepwidestar[1,2,3]{%
      f[X]\diamond_{\mathcal P}f[Y]\subseteq
      f[X\star_{\mathcal P}Y]}{%
      \FormulaRefAuto{PowerHomBackwardSubset}[thm]{1,2,3}}
    \proofstepwidestar[1,2,3]{%
      f[X\star_{\mathcal P}Y]=f[X]\diamond_{\mathcal P}f[Y]}{%
      \FormulaRefAuto{A\subseteq B, B\subseteq A\vdash A=B}[thm]{4,5}}
  \closeproofpartwide
\end{tabproofsplitwide}

\FormulaThmDeltaK[Operationserhaltung der induzierten Potenzabbildung]{%
  \begin{aligned}[t]
    &\SemigroupHom{f}{A,\star}{B,\diamond}\dsep
      X,Y\in\PowNE{A}\vdash{}\\[-2pt]
    &\PowMapNE{f}(X\star_{\mathcal P}Y)
      =\PowMapNE{f}(X)\diamond_{\mathcal P}\PowMapNE{f}(Y)
  \end{aligned}%
}{PowerSemigroupInducedMapOperation}{%
  \DeltaRow{Mengen}{A\dsep B\dsep X\dsep Y}
  \DeltaRow{Funktionen}{f\dsep\star\dsep\diamond\dsep
    \star_{\mathcal P}\dsep\diamond_{\mathcal P}\dsep
    \PowMap{f}\dsep\PowMapNE{f}}
}
\begin{tabproofwide}
  \proofstepwidestar[1]{\SemigroupHom{f}{A,\star}{B,\diamond}}{\rA}
  \proofstepwidestar[2]{X\in\PowNE{A}}{\rA}
  \proofstepwidestar[3]{Y\in\PowNE{A}}{\rA}
  \proofstepwidestar[1]{\SemigroupAlg{A}{\star}}{%
    \FormulaRefAuto{SemigroupHomSourceAxiom}[ax]{1}}
  \proofstepwidestar[1]{f\colon A\to B}{%
    \FormulaRefAuto{SemigroupHomFunctionAxiom}[ax]{1}}
  \proofstepwidestar[1,2,3]{X\star_{\mathcal P}Y\in\PowNE{A}}{%
    \FormulaRefAuto{PowerSemigroupProductClosure}[thm]{4,2,3}}
  \proofstepwidestar[1,2]{%
    \PowMapNE{f}(X)=f[X]}{%
    \FormulaRefAuto{NonemptyPowerMapValue}[thm]{5,2}}
  \proofstepwidestar[1,2]{%
    f[X]=\PowMapNE{f}(X)}{%
    \FormulaRefAuto{a=b\vdash b=a}[thm]{7}}
  \proofstepwidestar[1,3]{%
    \PowMapNE{f}(Y)=f[Y]}{%
    \FormulaRefAuto{NonemptyPowerMapValue}[thm]{5,3}}
  \proofstepwidestar[1,3]{%
    f[Y]=\PowMapNE{f}(Y)}{%
    \FormulaRefAuto{a=b\vdash b=a}[thm]{9}}
  \proofstepwide[1,2,3]{%
    \PowMapNE{f}(X\star_{\mathcal P}Y)}{=}{f[X\star_{\mathcal P}Y]}{%
    \FormulaRefAuto{NonemptyPowerMapValue}[thm]{5,6}}
  \proofstepwide[1,2,3]{}{=}{f[X]\diamond_{\mathcal P}f[Y]}{%
    \FormulaRefAuto{SemigroupHomPreservesPowerProduct}[thm]{1,2,3}}
  \proofstepwide[1,2,3]{}{=}{%
    \PowMapNE{f}(X)\diamond_{\mathcal P}f[Y]}{\rLRS{8}}
  \proofstepwide[1,2,3]{}{=}{%
    \PowMapNE{f}(X)\diamond_{\mathcal P}\PowMapNE{f}(Y)}{\rLRS{10}}
  \proofstepwide[1,2,3]{%
    \PowMapNE{f}(X\star_{\mathcal P}Y)}{=}{%
    \PowMapNE{f}(X)\diamond_{\mathcal P}\PowMapNE{f}(Y)}{\rChain{11,14}}
\end{tabproofwide}

\FormulaThmDeltaK[Isomorphismen induzieren Potenzhalbgruppenisomorphismen]{%
  \begin{aligned}[t]
    &\SemigroupIso{f}{A,\star}{B,\diamond}\vdash{}\\[-2pt]
    &\SemigroupIso{\PowMapNE{f}}
      {\PowNE{A},\star_{\mathcal P}}
      {\PowNE{B},\diamond_{\mathcal P}}
  \end{aligned}%
}{SemigroupIsoInducesPowerSemigroupIso}{%
  \DeltaRow{Mengen}{A\dsep B\dsep X\dsep Y}
  \DeltaRow{Funktionen}{f\dsep\star\dsep\diamond\dsep
    \star_{\mathcal P}\dsep\diamond_{\mathcal P}\dsep
    \PowMap{f}\dsep\PowMapNE{f}}
}
\begin{tabproofwide}
  \proofstepwidestarbreak[1]{\SemigroupIso{f}{A,\star}{B,\diamond}}{\rA}
  \proofstepwidestarbreak[1]{\SemigroupHom{f}{A,\star}{B,\diamond}}{%
    \FormulaRefAuto{SemigroupIsoHomAxiom}[ax]{1}}
  \proofstepwidestarbreak[1]{f\colon A\bij B}{%
    \FormulaRefAuto{SemigroupIsoBijectiveAxiom}[ax]{1}}
  \proofstepwidestarbreak[1]{\SemigroupAlg{A}{\star}}{%
    \FormulaRefAuto{SemigroupHomSourceAxiom}[ax]{2}}
  \proofstepwidestarbreak[1]{\SemigroupAlg{B}{\diamond}}{%
    \FormulaRefAuto{SemigroupHomTargetAxiom}[ax]{2}}
  \proofstepwidestarbreak[1]{\SemigroupAlg{\PowNE{A}}{\star_{\mathcal P}}}{%
    \FormulaRefAuto{NonemptyPowerSemigroup}[thm]{4}}
  \proofstepwidestarbreak[1]{\SemigroupAlg{\PowNE{B}}{\diamond_{\mathcal P}}}{%
    \FormulaRefAuto{NonemptyPowerSemigroup}[thm]{5}}
  \proofstepwidestarbreak[1]{%
    \PowMapNE{f}\colon\PowNE{A}\bij\PowNE{B}}{%
    \FormulaRefAuto{NonemptyPowerMapBijective}[thm]{3}}
  \proofstepwidestarbreak[9]{X\in\PowNE{A}}{\rA}
  \proofstepwidestarbreak[10]{Y\in\PowNE{A}}{\rA}
  \proofstepwidestarbreak[1,9,10]{%
    \PowMapNE{f}(X\star_{\mathcal P}Y)
      =\PowMapNE{f}(X)\diamond_{\mathcal P}\PowMapNE{f}(Y)}{%
    \FormulaRefAuto{PowerSemigroupInducedMapOperation}[thm]{2,9,10}}
  \proofstepwidestarbreak[1]{%
    \PowMapNE{f}\colon\PowNE{A}\to\PowNE{B}}{%
    \FormulaRefAuto{Bijektive Funktion}[def]{8}}
  \proofstepwidestarbreak[1]{%
    \begin{aligned}[t]
      &\mathsf{HGrHom}\!\bigl(\PowMapNE{f};\PowNE{A},\star_{\mathcal P};\\[-2pt]
      &\qquad\PowNE{B},\diamond_{\mathcal P}\bigr)
    \end{aligned}}{%
    \FormulaRefAuto{SemigroupHomDef}[def]{6,7,12,11}}
  \proofstepwidestarbreak[1]{%
    \SemigroupIso{\PowMapNE{f}}
      {\PowNE{A},\star_{\mathcal P}}
      {\PowNE{B},\diamond_{\mathcal P}}}{%
    \FormulaRefAuto{SemigroupIsoDef}[def]{13,8}}
\end{tabproofwide}

\subsection{Produktquadrate und ihre Isomorphieinvarianz}

\par\noindent\begin{minipage}{\linewidth}
\FormulaThmDeltaK[Nichtleerheit des Potenzträgers]{%
  A\neq\varnothing\vdash\PowNE{A}\neq\varnothing%
}{NonemptyPowerCarrierNonempty}{%
  \DeltaRow{Mengen}{A}
}
\end{minipage}\par
\begin{tabproofwide}
  \proofstepwidestarbreakkeep[1]{A\neq\varnothing}{\rA}
  \proofstepwidestar[1]{A\subseteq A}{%
      \FormulaRefAuto{A\subseteq A}[thm]}
  \proofstepwidestarbreakkeep[1]{A\in\PowNE{A}}{%
    \FormulaRefAuto{NonemptyPowersetMembership}[thm]{%
      \rAI{2,\allowbreak 1}}}
  \proofstepwidestarbreakkeep[1]{\PowNE{A}\neq\varnothing}{%
    \FormulaRefAuto{a\in A\vdash A\neq\varnothing}[thm]{3}}
\end{tabproofwide}

In dieser Untersektion bezeichnet \(A^2\) das \emph{Produktquadrat}
einer nichtleeren Halbgruppe, also die Menge ihrer Produktwerte.  Es ist
nicht das kartesische Quadrat \(A\times A\).  Die jeweilige Operation
ergibt sich aus der angegebenen Halbgruppe; insbesondere verwendet
\(\PowNE{A}^{\,2}\) die Operation \(\star_{\mathcal P}\).
Auf einem Unterträger \(U\) bezeichnet \(\star\) stets die Einschränkung
der ursprünglichen Operation auf \(U\times U\). Entsprechend ist bei
\(\SemigroupAlg{A^2}{\star}\) die eingeschränkte Operation gemeint.

\par\noindent\begin{minipage}{\linewidth}
\FormulaDefDeltaK[Produktquadrat einer nichtleeren Halbgruppe]{%
  \SemigroupAlg{A}{\star}\dsep A\neq\varnothing\vdash
  A^2\coloneqq A\star_{\mathcal P}A%
}{SemigroupSquareDef}{%
  \DeltaRow{Mengen}{A}
  \DeltaRow{Funktionen}{\star\dsep\star_{\mathcal P}}
  \DeltaRow{Lokale Schreibweise}{A^2}
}
\end{minipage}\par
Die Definition ist zulässig: Aus \(A\subseteq A\) und
\(A\neq\varnothing\) folgt nach
\FormulaRefAuto{NonemptyPowersetMembership}[thm], dass
\(A\in\PowNE{A}\).

\par\noindent\begin{minipage}{\linewidth}
\FormulaThmDeltaK[Elementkriterium des Produktquadrats]{%
  \begin{aligned}[t]
    &\SemigroupAlg{A}{\star}\dsep A\neq\varnothing\vdash{}\\[-2pt]
    &z\in A^2\leftrightarrow
      \bigl(z\in A\land\exists a\in A\,\exists b\in A\,
        (z=a\star b)\bigr)
  \end{aligned}%
}{SemigroupSquareMembership}{%
  \DeltaRow{Mengen}{A\dsep a\dsep b\dsep z}
  \DeltaRow{Funktionen}{\star\dsep\star_{\mathcal P}}
}
\end{minipage}\par
\begin{tabproofwide}
  \proofstepwidestarbreakkeep[1]{\SemigroupAlg{A}{\star}}{\rA}
  \proofstepwidestarbreakkeep[2]{A\neq\varnothing}{\rA}
  \proofstepwidestar[2]{A\subseteq A}{%
      \FormulaRefAuto{A\subseteq A}[thm]}
  \proofstepwidestarbreakkeep[2]{A\in\PowNE{A}}{%
    \FormulaRefAuto{NonemptyPowersetMembership}[thm]{%
      \rAI{3,\allowbreak 2}}}
  \proofstepwidestarbreakkeep[1,2]{A^2=A\star_{\mathcal P}A}{%
    \FormulaRefAuto{SemigroupSquareDef}[def]{1,\allowbreak 2}}
  \proofstepwidestar[1,2]{z\in A\star_{\mathcal P}A\leftrightarrow(z\in A\land\exists a\in A\,\exists b\in A\;(z=a\star b))}{%
      \FormulaRefAuto{PowerSemigroupProductMembership}[thm]{1,\allowbreak 4,\allowbreak 4}}
  \proofstepwidestarbreakkeep[1,2]{%
    z\in A^2\leftrightarrow
    \bigl(z\in A\land\exists a\in A\,\exists b\in A\,
      (z=a\star b)\bigr)}{%
    \rIE{5,\allowbreak 6}}
\end{tabproofwide}

\par\noindent\begin{minipage}{\linewidth}
\FormulaThmDeltaK[Produktzeugen im Quadrat]{%
  \SemigroupAlg{A}{\star}\dsep A\neq\varnothing\dsep z\in A^2
  \vdash\exists a\in A\,\exists b\in A\,(z=a\star b)%
}{SemigroupSquareWitnesses}{%
  \DeltaRow{Mengen}{A\dsep a\dsep b\dsep z}
  \DeltaRow{Funktionen}{\star}
}
\end{minipage}\par
\begin{tabproofwide}
  \proofstepwidestarbreakkeep[1]{\SemigroupAlg{A}{\star}}{\rA}
  \proofstepwidestarbreakkeep[2]{A\neq\varnothing}{\rA}
  \proofstepwidestarbreakkeep[3]{z\in A^2}{\rA}
  \proofstepwidestar[1,2,3]{z\in A\land\exists a\in A\,\exists b\in A\;(z=a\star b)}{%
      \FormulaRefAuto{SemigroupSquareMembership}[thm]{1,\allowbreak 2,\allowbreak 3}}
  \proofstepwidestarbreakkeep[1,2,3]{%
    \exists a\in A\,\exists b\in A\,(z=a\star b)}{%
    \rAEb{4}}
\end{tabproofwide}

\par\noindent\begin{minipage}{\linewidth}
\FormulaThmDeltaK[Produktwerte liegen im Quadrat]{%
  \SemigroupAlg{A}{\star}\dsep A\neq\varnothing\dsep a,b\in A
  \vdash a\star b\in A^2%
}{SemigroupSquareProductMember}{%
  \DeltaRow{Mengen}{A\dsep a\dsep b}
  \DeltaRow{Funktionen}{\star}
}
\end{minipage}\par
\begin{tabproofwide}
  \proofstepwidestarbreakkeep[1]{\SemigroupAlg{A}{\star}}{\rA}
  \proofstepwidestarbreakkeep[2]{A\neq\varnothing}{\rA}
  \proofstepwidestarbreakkeep[3]{a,b\in A}{\rA}
  \proofstepwidestarbreakkeep[1,3]{a\star b\in A}{%
    \FormulaRefAuto{SemigroupClosureAxiom}[ax]{1,\allowbreak 3}}
  \proofstepwidestarbreakkeep[3]{%
    \exists x\in A\,\exists y\in A\,(a\star b=x\star y)}{%
    \rEIStar{3,\allowbreak \rII}}
  \proofstepwidestarbreakkeep[1,2,3]{a\star b\in A^2}{%
    \FormulaRefAuto{SemigroupSquareMembership}[thm]{1,\allowbreak 2,\allowbreak \rAI{4,\allowbreak 5}}}
\end{tabproofwide}

\par\noindent\begin{minipage}{\linewidth}
\FormulaThmDeltaK[Das Produktquadrat liegt im Träger]{%
  \SemigroupAlg{A}{\star}\dsep A\neq\varnothing\vdash A^2\subseteq A%
}{SemigroupSquareSubset}{%
  \DeltaRow{Mengen}{A\dsep z}
  \DeltaRow{Funktionen}{\star}
}
\end{minipage}\par
\begin{tabproofwide}
  \proofstepwidestarbreakkeep[1]{\SemigroupAlg{A}{\star}}{\rA}
  \proofstepwidestarbreakkeep[2]{A\neq\varnothing}{\rA}
  \proofstepwidestarbreakkeep[3]{z\in A^2}{\rA}
  \proofstepwidestar[1,2,3]{z\in A\land\exists a\in A\,\exists b\in A\;(z=a\star b)}{%
      \FormulaRefAuto{SemigroupSquareMembership}[thm]{1,\allowbreak 2,\allowbreak 3}}
  \proofstepwidestarbreakkeep[1,2,3]{z\in A}{%
    \rAEa{4}}
  \proofstepwidestarbreakkeep[1,2]{A^2\subseteq A}{%
    \FormulaRefAuto{SubsetIntroductionRule}[thm]{[3]\vdots 5}}
\end{tabproofwide}

\par\noindent\begin{minipage}{\linewidth}
\FormulaThmDeltaK[Das Produktquadrat ist eine Unterhalbgruppe]{%
  \SemigroupAlg{A}{\star}\dsep A\neq\varnothing
  \vdash\mathsf{UHGr}(A^2;A,\star)%
}{SemigroupSquareIsSubsemigroup}{%
  \DeltaRow{Mengen}{A\dsep x\dsep y}
  \DeltaRow{Funktionen}{\star\dsep\star_{\mathcal P}}
}
\end{minipage}\par
Die Beweisidee ist unmittelbar: Zwei Elemente des Quadrats liegen schon
im ursprünglichen Träger, und ihr Produkt gehört deshalb wieder zum Quadrat.
\begin{tabproofwide}
  \proofstepwidestarbreakkeep[1]{\SemigroupAlg{A}{\star}}{\rA}
  \proofstepwidestarbreakkeep[2]{A\neq\varnothing}{\rA}
  \proofstepwidestarbreakkeep[1,2]{A^2\subseteq A}{%
    \FormulaRefAuto{SemigroupSquareSubset}[thm]{1,\allowbreak 2}}
  \proofstepwidestar[2]{A\subseteq A}{%
      \FormulaRefAuto{A\subseteq A}[thm]}
  \proofstepwidestarbreakkeep[2]{A\in\PowNE{A}}{%
    \FormulaRefAuto{NonemptyPowersetMembership}[thm]{%
      \rAI{4,\allowbreak 2}}}
  \proofstepwidestar[1,2]{A^2=A\star_{\mathcal P}A}{%
      \FormulaRefAuto{SemigroupSquareDef}[def]{1,\allowbreak 2}}
  \proofstepwidestar[1,2]{A\star_{\mathcal P}A\neq\varnothing}{%
      \FormulaRefAuto{PowerSemigroupProductNonempty}[thm]{1,\allowbreak 5,\allowbreak 5}}
  \proofstepwidestarbreakkeep[1,2]{A^2\neq\varnothing}{%
    \rIE{6,\allowbreak %
      7}}
  \proofstepwidestarbreakkeep[9]{x,y\in A^2}{\rA}
  \proofstepwidestarbreakkeep[1,2,9]{x,y\in A}{%
    \FormulaRefAuto{A\subseteq B\dsep x\in A\vdash x\in B}[thm]{3,\allowbreak 9}}
  \proofstepwidestarbreakkeep[1,2,9]{x\star y\in A^2}{%
    \FormulaRefAuto{SemigroupSquareProductMember}[thm]{1,\allowbreak 2,\allowbreak 10}}
  \proofstepwidestarbreakkeep[1,2]{\mathsf{UHGr}(A^2;A,\star)}{%
    \FormulaRefAuto{SemigroupSubsemigroupDef}[def]{1,\allowbreak 3,\allowbreak 8,\allowbreak 11}}
\end{tabproofwide}

\par\noindent\begin{minipage}{\linewidth}
\FormulaThmDeltaK[Das Potenzquadrat liegt über dem Grundquadrat]{%
  \SemigroupAlg{A}{\star}\dsep A\neq\varnothing
  \vdash\PowNE{A}^{\,2}\subseteq\PowNE{A^2}%
}{PowerSemigroupSquareSubset}{%
  \DeltaRow{Mengen}{A\dsep X\dsep Y\dsep Z\dsep z\dsep a\dsep b}
  \DeltaRow{Funktionen}{\star\dsep\star_{\mathcal P}}
}
\end{minipage}\par
Ein Element des Potenzquadrats ist ein Mengenprodukt. Jeder seiner
Einträge ist daher ein Produkt zweier Grundelemente.
\begin{tabproofwide}
  \proofstepwidestarbreakkeep[1]{\SemigroupAlg{A}{\star}}{\rA}
  \proofstepwidestarbreakkeep[2]{A\neq\varnothing}{\rA}
  \proofstepwidestarbreakkeep[1]{\SemigroupAlg{\PowNE{A}}{\star_{\mathcal P}}}{%
    \FormulaRefAuto{NonemptyPowerSemigroup}[thm]{1}}
  \proofstepwidestar[2]{A\subseteq A}{%
      \FormulaRefAuto{A\subseteq A}[thm]}
  \proofstepwidestarbreakkeep[2]{A\in\PowNE{A}}{%
    \FormulaRefAuto{NonemptyPowersetMembership}[thm]{%
      \rAI{4,\allowbreak 2}}}
  \proofstepwidestarbreakkeep[2]{\PowNE{A}\neq\varnothing}{%
    \FormulaRefAuto{a\in A\vdash A\neq\varnothing}[thm]{5}}
  \proofstepwidestarbreakkeep[7]{Z\in\PowNE{A}^{\,2}}{\rA}
  \proofstepwidestarbreakkeep[1,2,7]{%
    \exists X\in\PowNE{A}\,\exists Y\in\PowNE{A}\,
      (Z=X\star_{\mathcal P}Y)}{%
    \FormulaRefAuto{SemigroupSquareWitnesses}[thm]{3,\allowbreak 6,\allowbreak 7}}
  \proofstepwidestarbreakkeep[9]{X,Y\in\PowNE{A}\land Z=X\star_{\mathcal P}Y}{\rA}
  \proofstepwidestarbreakkeep[10]{z\in Z}{\rA}
  \proofstepwidestarbreakkeep[1,9,10]{%
    \exists a\in X\,\exists b\in Y\,(z=a\star b)}{%
    \FormulaRefAuto{PowerSemigroupProductWitnesses}[thm]{1,\allowbreak 9,\allowbreak \rIE{9,\allowbreak 10}}}
  \proofstepwidestarbreakkeep[12]{a\in X\land b\in Y\land z=a\star b}{\rA}
  \proofstepwidestarbreakkeep[9,12]{a,b\in A}{%
    \FormulaRefAuto{PowerSemigroupCarrierElement}[thm]{9,\allowbreak 12}}
  \proofstepwidestar[1,2,9,12]{a\star b\in A^2}{%
      \FormulaRefAuto{SemigroupSquareProductMember}[thm]{1,\allowbreak 2,\allowbreak 13}}
  \proofstepwidestarbreakkeep[1,2,9,12]{z\in A^2}{%
    \rIE{12,\allowbreak 14}}
  \proofstepwidestarbreakkeep[1,2,9,10]{z\in A^2}{\rEEStar{11,\allowbreak 12\text{--}15}}
  \proofstepwidestarbreakkeep[1,2,9]{Z\subseteq A^2}{%
    \FormulaRefAuto{SubsetIntroductionRule}[thm]{[10]\vdots 16}}
  \proofstepwidestar[1,9]{X\star_{\mathcal P}Y\neq\varnothing}{%
      \FormulaRefAuto{PowerSemigroupProductNonempty}[thm]{1,\allowbreak 9}}
  \proofstepwidestarbreakkeep[1,9]{Z\neq\varnothing}{%
    \rIE{9,\allowbreak 18}}
  \proofstepwidestarbreakkeep[1,2,9]{Z\in\PowNE{A^2}}{%
    \FormulaRefAuto{NonemptyPowersetMembership}[thm]{\rAI{17,\allowbreak 19}}}
  \proofstepwidestarbreakkeep[1,2,7]{Z\in\PowNE{A^2}}{\rEEStar{8,\allowbreak 9\text{--}20}}
  \proofstepwidestarbreakkeep[1,2]{\PowNE{A}^{\,2}\subseteq\PowNE{A^2}}{%
    \FormulaRefAuto{SubsetIntroductionRule}[thm]{[7]\vdots 21}}
\end{tabproofwide}

\par\noindent\begin{minipage}{\linewidth}
\FormulaThmDeltaK[Einermengen im Potenzquadrat]{%
  \SemigroupAlg{A}{\star}\dsep A\neq\varnothing\dsep a\in A\vdash
  \bigl(\{a\}\in\PowNE{A}^{\,2}\leftrightarrow a\in A^2\bigr)%
}{PowerSemigroupSquareSingletonCriterion}{%
  \DeltaRow{Mengen}{A\dsep a\dsep x\dsep y}
  \DeltaRow{Funktionen}{\star\dsep\star_{\mathcal P}\dsep\eta_A}
}
\end{minipage}\par
Für die Rückrichtung wird eine Produktdarstellung von \(a\) durch die
Einermengenabbildung in eine Produktdarstellung von \(\{a\}\) überführt.
\begin{tabproofwide}
  \proofstepwidestarbreakkeep[1]{\SemigroupAlg{A}{\star}}{\rA}
  \proofstepwidestarbreakkeep[2]{A\neq\varnothing}{\rA}
  \proofstepwidestarbreakkeep[3]{a\in A}{\rA}
  \proofstepwidestarbreakkeep[1]{\SemigroupAlg{\PowNE{A}}{\star_{\mathcal P}}}{%
    \FormulaRefAuto{NonemptyPowerSemigroup}[thm]{1}}
  \proofstepwidestar[2]{A\subseteq A}{%
      \FormulaRefAuto{A\subseteq A}[thm]}
  \proofstepwidestarbreakkeep[2]{A\in\PowNE{A}}{%
    \FormulaRefAuto{NonemptyPowersetMembership}[thm]{%
      \rAI{5,\allowbreak 2}}}
  \proofstepwidestarbreakkeep[2]{\PowNE{A}\neq\varnothing}{%
    \FormulaRefAuto{a\in A\vdash A\neq\varnothing}[thm]{6}}
  \proofstepwidestarbreakkeep[8]{\{a\}\in\PowNE{A}^{\,2}}{\rA}
  \proofstepwidestar[1,2,8]{\PowNE{A}^{\,2}\subseteq\PowNE{A^2}}{%
      \FormulaRefAuto{PowerSemigroupSquareSubset}[thm]{1,\allowbreak 2}}
  \proofstepwidestarbreakkeep[1,2,8]{\{a\}\in\PowNE{A^2}}{%
    \FormulaRefAuto{A\subseteq B\dsep x\in A\vdash x\in B}[thm]{%
      9,\allowbreak 8}}
  \proofstepwidestar[1,2,8]{a\in\{a\}}{%
      \FormulaRefAuto{a\in\{a\}}[thm]}
  \proofstepwidestarbreakkeep[1,2,8]{a\in A^2}{%
    \FormulaRefAuto{PowerSemigroupCarrierElement}[thm]{10,\allowbreak 11}}
  \proofstepwidestarbreakkeep[1,2]{\{a\}\in\PowNE{A}^{\,2}\rightarrow a\in A^2}{\rRI{8,\allowbreak 12}}
  \proofstepwidestarbreakkeep[14]{a\in A^2}{\rA}
  \proofstepwidestarbreakkeep[1,2,14]{\exists x\in A\,\exists y\in A\,(a=x\star y)}{%
    \FormulaRefAuto{SemigroupSquareWitnesses}[thm]{1,\allowbreak 2,\allowbreak 14}}
  \proofstepwidestarbreakkeep[16]{x,y\in A\land a=x\star y}{\rA}
  \proofstepwidestar[1]{\SemigroupIso{\eta_A}{A,\star}{\operatorname{Sing}(A),\star_{\mathcal P}}}{%
      \FormulaRefAuto{SemigroupSingletonCopyIsomorphism}[thm]{1}}
  \proofstepwidestarbreakkeep[1]{\SemigroupHom{\eta_A}{A,\star}{\operatorname{Sing}(A),\star_{\mathcal P}}}{%
    \FormulaRefAuto{SemigroupIsoHomAxiom}[ax]{%
      17}}
  \proofstepwidestar[1,16]{\forall u\in A\;(\eta_A(u)=\{u\})}{%
      \FormulaRefAuto{SemigroupSingletonCopyDef}[def]{1}}
  \proofstepwidestar[1,16]{\eta_A(x\star y)=\eta_A(x)\star_{\mathcal P}\eta_A(y)}{%
      \FormulaRefAuto{SemigroupHomOperationAxiom}[ax]{18,\allowbreak 16}}
  \proofstepwidestarbreakkeep[1,16]{\{a\}=\{x\}\star_{\mathcal P}\{y\}}{%
    \rIE{16,\allowbreak 19,\allowbreak %
      20}}
  \proofstepwidestar[16]{x\in\{x\}\land y\in\{y\}}{%
      \FormulaRefAuto{a\in\{a\}}[thm]}
  \proofstepwidestar[16]{\{x\}\subseteq A\land\{y\}\subseteq A}{%
      \FormulaRefAuto{a\in A\vdash\{a\}\subseteq A}[thm]{16}}
  \proofstepwidestar[16]{\{x\}\neq\varnothing\land\{y\}\neq\varnothing}{%
      \FormulaRefAuto{a\in A\vdash A\neq\varnothing}[thm]{%
        22}}
  \proofstepwidestarbreakkeep[16]{\{x\},\{y\}\in\PowNE{A}}{%
    \FormulaRefAuto{NonemptyPowersetMembership}[thm]{%
      23,\allowbreak %
      24}}
  \proofstepwidestar[1,2,16]{\{x\}\star_{\mathcal P}\{y\}\in\PowNE{A}^{\,2}}{%
      \FormulaRefAuto{SemigroupSquareProductMember}[thm]{4,\allowbreak 7,\allowbreak 25}}
  \proofstepwidestarbreakkeep[1,2,16]{\{a\}\in\PowNE{A}^{\,2}}{%
    \rIE{21,\allowbreak 26}}
  \proofstepwidestarbreakkeep[1,2,14]{\{a\}\in\PowNE{A}^{\,2}}{\rEEStar{15,\allowbreak 16\text{--}27}}
  \proofstepwidestarbreakkeep[1,2]{a\in A^2\rightarrow\{a\}\in\PowNE{A}^{\,2}}{\rRI{14,\allowbreak 28}}
  \proofstepwidestarbreakkeep[1,2,3]{\{a\}\in\PowNE{A}^{\,2}\leftrightarrow a\in A^2}{\rLRI{13,\allowbreak 29}}
\end{tabproofwide}

\par\noindent\begin{minipage}{\linewidth}
\FormulaThmDeltaK[Isomorphismen erhalten das Produktquadrat]{%
  \SemigroupIso{F}{A,\star}{B,\diamond}\dsep
  A\neq\varnothing\dsep B\neq\varnothing\vdash F[A^2]=B^2%
}{SemigroupIsoSquareImage}{%
  \DeltaRow{Mengen}{A\dsep B}
  \DeltaRow{Funktionen}{F\dsep\star\dsep\diamond}
}
\end{minipage}\par
\begin{tabproofwide}
  \proofstepwidestarbreakkeep[1]{\SemigroupIso{F}{A,\star}{B,\diamond}}{\rA}
  \proofstepwidestarbreakkeep[2]{A\neq\varnothing}{\rA}
  \proofstepwidestarbreakkeep[3]{B\neq\varnothing}{\rA}
  \proofstepwidestarbreakkeep[1]{\SemigroupHom{F}{A,\star}{B,\diamond}}{%
    \FormulaRefAuto{SemigroupIsoHomAxiom}[ax]{1}}
  \proofstepwidestar[1]{\SemigroupAlg{A}{\star}}{%
      \FormulaRefAuto{SemigroupHomSourceAxiom}[ax]{4}}
  \proofstepwidestar[1]{\SemigroupAlg{B}{\diamond}}{%
      \FormulaRefAuto{SemigroupHomTargetAxiom}[ax]{4}}
  \proofstepwidestarbreakkeep[1]{\SemigroupAlg{A}{\star}\land\SemigroupAlg{B}{\diamond}}{%
    \rAI{5,\allowbreak %
      6}}
  \proofstepwidestar[1]{F\colon A\bij B}{%
      \FormulaRefAuto{SemigroupIsoBijectiveAxiom}[ax]{1}}
  \proofstepwidestarbreakkeep[1]{F[A]=B}{%
    \FormulaRefAuto{F\colon A\bij B \vdash F[A]=B}[thm]{%
      8}}
  \proofstepwidestar[2]{A\subseteq A}{%
      \FormulaRefAuto{A\subseteq A}[thm]}
  \proofstepwidestarbreakkeep[2]{A\in\PowNE{A}}{%
    \FormulaRefAuto{NonemptyPowersetMembership}[thm]{%
      \rAI{10,\allowbreak 2}}}
  \proofstepwidestar[1,2,3]{A^2=A\star_{\mathcal P}A}{%
      \FormulaRefAuto{SemigroupSquareDef}[def]{7,\allowbreak 2}}
  \proofstepwidestarbreakkeep[1,2,3]{F[A^2]=F[A\star_{\mathcal P}A]}{%
    \rIE{12}}
  \proofstepwidestarbreakkeep[1,2,3]{\phantom{F[A^2]}=F[A]\diamond_{\mathcal P}F[A]}{%
    \FormulaRefAuto{SemigroupHomPreservesPowerProduct}[thm]{4,\allowbreak 11,\allowbreak 11}}
  \proofstepwidestarbreakkeep[1,2,3]{\phantom{F[A^2]}=B\diamond_{\mathcal P}B}{\rIE{9}}
  \proofstepwidestarbreakkeep[1,2,3]{\phantom{F[A^2]}=B^2}{%
    \FormulaRefAuto{SemigroupSquareDef}[def]{7,\allowbreak 3}}
  \proofstepwidestarbreakkeep[1,2,3]{F[A^2]=B^2}{\rChain{13,\allowbreak 16}}
\end{tabproofwide}

\par\noindent\begin{minipage}{\linewidth}
\FormulaThmDeltaK[Einschränkung eines Isomorphismus auf Halbgruppen]{%
  \begin{aligned}[t]
    &\SemigroupIso{F}{A,\star}{B,\diamond}\dsep
      \SemigroupAlg{U}{\star}\dsep\SemigroupAlg{V}{\diamond}\dsep{}\\[-2pt]
    &U\subseteq A\dsep V\subseteq B\dsep F[U]=V\vdash
      \SemigroupIso{F\restriction_U}{U,\star}{V,\diamond}
  \end{aligned}%
}{SemigroupIsoRestrictionToImage}{%
  \DeltaRow{Mengen}{A\dsep B\dsep U\dsep V\dsep x\dsep y}
  \DeltaRow{Funktionen}{F\dsep\star\dsep\diamond}
}
\end{minipage}\par
Die Einschränkung ist bijektiv auf ihr Bild und übernimmt die
Operationserhaltung. Eine Nichtleerheitsannahme wird hier nicht benötigt.
\begin{tabproofwide}
  \proofstepwidestarbreakkeep[1]{\SemigroupIso{F}{A,\star}{B,\diamond}}{\rA}
  \proofstepwidestarbreakkeep[2]{\SemigroupAlg{U}{\star}}{\rA}
  \proofstepwidestarbreakkeep[3]{\SemigroupAlg{V}{\diamond}}{\rA}
  \proofstepwidestarbreakkeep[4]{U\subseteq A}{\rA}
  \proofstepwidestarbreakkeep[5]{V\subseteq B}{\rA}
  \proofstepwidestarbreakkeep[6]{F[U]=V}{\rA}
  \proofstepwidestar[1]{F\colon A\bij B}{%
      \FormulaRefAuto{SemigroupIsoBijectiveAxiom}[ax]{1}}
  \proofstepwidestarbreakkeep[1]{F\colon A\inj B\land F\colon A\to B}{%
    \FormulaRefAuto{Bijektive Funktion}[def]{%
      7}}
  \proofstepwidestarbreakkeep[1,4]{F\restriction_U\colon U\inj B}{%
    \FormulaRefAuto{F\colon A\inj B\dsep C\subseteq A \vdash F\restriction_C\colon C\inj B}[thm]{8,\allowbreak 4}}
  \proofstepwidestarbreakkeep[1,4]{F\restriction_U\colon U\bij(F\restriction_U)[U]}{%
    \FormulaRefAuto{F\colon A\inj B \vdash F\colon A\bij F[A]}[thm]{9}}
  \proofstepwidestarbreakkeep[1,4]{(F\restriction_U)[U]=F[U]}{%
    \FormulaRefAuto{C\subseteq A \dsep F\colon A\to B \vdash (F\restriction_C)[C]=F[C]}[thm]{4,\allowbreak 8}}
  \proofstepwidestarbreakkeep[1,4,6]{F\restriction_U\colon U\bij V}{\rIE{6,\allowbreak 11,\allowbreak 10}}
  \proofstepwidestarbreakkeep[13]{x,y\in U}{\rA}
  \proofstepwidestarbreakkeep[4,13]{x,y\in A}{%
    \FormulaRefAuto{A\subseteq B\dsep x\in A\vdash x\in B}[thm]{4,\allowbreak 13}}
  \proofstepwidestarbreakkeep[2,13]{x\star y\in U}{%
    \FormulaRefAuto{SemigroupClosureAxiom}[ax]{2,\allowbreak 13}}
  \proofstepwidestarbreakkeep[1,2,4,13]{(F\restriction_U)(x\star y)=F(x\star y)}{%
    \FormulaRefAuto{C\subseteq A \dsep x\in C \vdash F\restriction_C(x)=F(x)}[thm]{8,\allowbreak 4,\allowbreak 15}}
  \proofstepwidestar[1,2,4,13]{\SemigroupHom{F}{A,\star}{B,\diamond}}{%
      \FormulaRefAuto{SemigroupIsoHomAxiom}[ax]{1}}
  \proofstepwidestarbreakkeep[1,2,4,13]{F(x\star y)=F(x)\diamond F(y)}{%
    \FormulaRefAuto{SemigroupHomOperationAxiom}[ax]{%
      17,\allowbreak 14}}
  \proofstepwidestar[1,2,4,13]{(F\restriction_U)(x)=F(x)\land(F\restriction_U)(y)=F(y)}{%
      \FormulaRefAuto{C\subseteq A \dsep x\in C \vdash F\restriction_C(x)=F(x)}[thm]{8,\allowbreak 4,\allowbreak 13}}
  \proofstepwidestarbreakkeep[1,2,4,13]{F(x)\diamond F(y)=(F\restriction_U)(x)\diamond(F\restriction_U)(y)}{%
    \rIE{19}}
  \proofstepwidestarbreakkeep[1,2,4,13]{%
    (F\restriction_U)(x\star y)=(F\restriction_U)(x)\diamond(F\restriction_U)(y)}{\rChain{16,18,20}}
  \proofstepwidestar[1,2,3,4,6]{F\restriction_U\colon U\to V}{%
      \FormulaRefAuto{Bijektive Funktion}[def]{12}}
  \proofstepwidestarbreakkeep[1,2,3,4,6]{\SemigroupHom{F\restriction_U}{U,\star}{V,\diamond}}{%
    \FormulaRefAuto{SemigroupHomDef}[def]{2,\allowbreak 3,\allowbreak %
      22,\allowbreak 21}}
  \proofstepwidestarbreakkeep[1,2,3,4,5,6]{\SemigroupIso{F\restriction_U}{U,\star}{V,\diamond}}{%
    \FormulaRefAuto{SemigroupIsoDef}[def]{23,\allowbreak 12}}
\end{tabproofwide}

\par\noindent\begin{minipage}{\linewidth}
\FormulaThmDeltaK[Isomorphe Einschränkung auf die Produktquadrate]{%
  \begin{aligned}[t]
    &\SemigroupIso{F}{A,\star}{B,\diamond}\dsep
      A\neq\varnothing\dsep B\neq\varnothing\vdash{}\\[-2pt]
    &\SemigroupIso{F\restriction_{A^2}}{A^2,\star}{B^2,\diamond}
  \end{aligned}%
}{SemigroupIsoSquareRestriction}{%
  \DeltaRow{Mengen}{A\dsep B}
  \DeltaRow{Funktionen}{F\dsep\star\dsep\diamond}
}
\end{minipage}\par
\begin{tabproofwide}
  \proofstepwidestarbreakkeep[1]{\SemigroupIso{F}{A,\star}{B,\diamond}}{\rA}
  \proofstepwidestarbreakkeep[2]{A\neq\varnothing}{\rA}
  \proofstepwidestarbreakkeep[3]{B\neq\varnothing}{\rA}
  \proofstepwidestar[1]{\SemigroupHom{F}{A,\star}{B,\diamond}}{%
      \FormulaRefAuto{SemigroupIsoHomAxiom}[ax]{1}}
  \proofstepwidestar[1]{\SemigroupHom{F}{A,\star}{B,\diamond}}{%
      \FormulaRefAuto{SemigroupIsoHomAxiom}[ax]{1}}
  \proofstepwidestar[1]{\SemigroupAlg{A}{\star}}{%
      \FormulaRefAuto{SemigroupHomSourceAxiom}[ax]{%
      4}}
  \proofstepwidestar[1]{\SemigroupAlg{B}{\diamond}}{%
      \FormulaRefAuto{SemigroupHomTargetAxiom}[ax]{%
      5}}
  \proofstepwidestarbreakkeep[1]{\SemigroupAlg{A}{\star}\land\SemigroupAlg{B}{\diamond}}{%
    \rAI{6,\allowbreak %
      7}}
  \proofstepwidestar[1,2]{\mathsf{UHGr}(A^2;A,\star)}{%
      \FormulaRefAuto{SemigroupSquareIsSubsemigroup}[thm]{8,\allowbreak 2}}
  \proofstepwidestarbreakkeep[1,2]{\SemigroupAlg{A^2}{\star}}{%
    \FormulaRefAuto{SubsemigroupFormsSemigroup}[thm]{8,\allowbreak %
      9}}
  \proofstepwidestar[1,3]{\mathsf{UHGr}(B^2;B,\diamond)}{%
      \FormulaRefAuto{SemigroupSquareIsSubsemigroup}[thm]{8,\allowbreak 3}}
  \proofstepwidestarbreakkeep[1,3]{\SemigroupAlg{B^2}{\diamond}}{%
    \FormulaRefAuto{SubsemigroupFormsSemigroup}[thm]{8,\allowbreak %
      11}}
  \proofstepwidestar[1,2,3]{A^2\subseteq A}{%
      \FormulaRefAuto{SemigroupSquareSubset}[thm]{8,\allowbreak 2}}
  \proofstepwidestar[1,2,3]{B^2\subseteq B}{%
      \FormulaRefAuto{SemigroupSquareSubset}[thm]{8,\allowbreak 3}}
  \proofstepwidestar[1,2,3]{F[A^2]=B^2}{%
      \FormulaRefAuto{SemigroupIsoSquareImage}[thm]{1,\allowbreak 2,\allowbreak 3}}
  \proofstepwidestarbreakkeep[1,2,3]{\SemigroupIso{F\restriction_{A^2}}{A^2,\star}{B^2,\diamond}}{%
    \FormulaRefAuto{SemigroupIsoRestrictionToImage}[thm]{1,\allowbreak 10,\allowbreak 12,\allowbreak %
      13,\allowbreak %
      14,\allowbreak %
      15}}
\end{tabproofwide}

\subsection{Transport von Lösungsmengen}

\par\noindent\begin{minipage}{\linewidth}
\FormulaDefDeltaK[Linke Lösungsmengen einer Halbgruppe]{%
  \SemigroupAlg{A}{\star}\dsep a,b\in A\vdash
  L_A(a,b)\coloneqq\{x\in A\mid a\star x=b\}%
}{SemigroupLeftEquationFiberDef}{%
  \DeltaRow{Mengen}{A\dsep a\dsep b\dsep x}
  \DeltaRow{Funktionen}{\star}
  \DeltaRow{Lokale Schreibweise}{L_A(a,b)}
}
\end{minipage}\par
Die Operation in \(L_A(a,b)\) ergibt sich aus der angegebenen
Halbgruppe. Ein Isomorphismus erhält die Gleichung \(a\star x=b\)
in beiden Richtungen. Das punktweise Kriterium aus Band~08 liefert
deshalb eine Bijektion zwischen den Lösungsmengen.

\par\noindent\begin{minipage}{\linewidth}
\FormulaThmDeltaK[Isomorphismen transportieren linke Lösungsmengen]{%
  \begin{aligned}[t]
    &\SemigroupIso{F}{A,\star}{B,\diamond}\dsep a,b\in A\vdash{}\\[-2pt]
    &F\restriction_{L_A(a,b)}\colon L_A(a,b)\bij L_B(F(a),F(b))
  \end{aligned}%
}{SemigroupLeftEquationFiberTransport}{%
  \DeltaRow{Mengen}{A\dsep B\dsep a\dsep b\dsep x}
  \DeltaRow{Funktionen}{F\dsep\star\dsep\diamond}
  \DeltaRow{Lösungsmengen}{L_A(a,b)\dsep L_B(F(a),F(b))}
}
\end{minipage}\par
\begin{tabproofsplitwide}
\proofpartwide{Äquivalenz der beiden Gleichungen}
  \proofstepwidestarbreakkeep[1]{\SemigroupIso{F}{A,\star}{B,\diamond}}{\rA}
  \proofstepwidestarbreakkeep[2]{a,b\in A}{\rA}
  \proofstepwidestarbreakkeep[1]{F\colon A\bij B}{%
    \FormulaRefAuto{SemigroupIsoBijectiveAxiom}[ax]{1}}
  \proofstepwidestar[1]{\SemigroupHom{F}{A,\star}{B,\diamond}}{%
      \FormulaRefAuto{SemigroupIsoHomAxiom}[ax]{1}}
  \proofstepwidestar[1]{\SemigroupHom{F}{A,\star}{B,\diamond}}{%
      \FormulaRefAuto{SemigroupIsoHomAxiom}[ax]{1}}
  \proofstepwidestar[1]{\SemigroupAlg{A}{\star}}{%
      \FormulaRefAuto{SemigroupHomSourceAxiom}[ax]{%
      4}}
  \proofstepwidestar[1]{\SemigroupAlg{B}{\diamond}}{%
      \FormulaRefAuto{SemigroupHomTargetAxiom}[ax]{%
      5}}
  \proofstepwidestarbreakkeep[1]{\SemigroupAlg{A}{\star}\land\SemigroupAlg{B}{\diamond}}{%
    \rAI{6,\allowbreak %
      7}}
  \proofstepwidestarbreakkeep[1,2]{F(a),F(b)\in B}{%
    \FormulaRefAuto{F\colon A\bij B\dsep x\in A\vdash F(x)\in B}[thm]{3,\allowbreak 2}}
  \proofstepwidestarbreakkeep[10]{x\in A}{\rA}
  \proofstepwidestar[1,2,10]{\SemigroupHom{F}{A,\star}{B,\diamond}}{%
      \FormulaRefAuto{SemigroupIsoHomAxiom}[ax]{1}}
  \proofstepwidestarbreakkeep[1,2,10]{F(a\star x)=F(a)\diamond F(x)}{%
    \FormulaRefAuto{SemigroupHomOperationAxiom}[ax]{%
      11,\allowbreak 2,\allowbreak 10}}
  \proofstepwidestarbreakkeep[1,2,10]{a\star x\in A}{%
    \FormulaRefAuto{SemigroupClosureAxiom}[ax]{8,\allowbreak 2,\allowbreak 10}}
  \proofstepwidestarbreakkeep[14]{a\star x=b}{\rA}
  \proofstepwidestarbreakkeep[1,2,10,14]{F(a)\diamond F(x)=F(b)}{\rIE{14,\allowbreak 12}}
  \proofstepwidestarbreakkeep[1,2,10]{a\star x=b\rightarrow F(a)\diamond F(x)=F(b)}{\rRI{14,\allowbreak 15}}
  \proofstepwidestarbreakkeep[17]{F(a)\diamond F(x)=F(b)}{\rA}
  \proofstepwidestarbreakkeep[1,2,10,17]{F(a\star x)=F(b)}{\rChain{12,\allowbreak 17}}
  \proofstepwidestar[1,2,10,17]{F\colon A\inj B}{%
      \FormulaRefAuto{F\colon A\bij B\vdash F\colon A\inj B}[thm]{3}}
  \proofstepwidestarbreakkeep[1,2,10,17]{a\star x=b}{%
    \FormulaRefAuto{F\colon A\inj B\dsep x\in A\dsep y\in A\dsep
      F(x)=F(y)\vdash x=y}[ax]{%
      19,\allowbreak 13,\allowbreak 2,\allowbreak 18}}
  \proofstepwidestarbreakkeep[1,2,10]{F(a)\diamond F(x)=F(b)\rightarrow a\star x=b}{\rRI{17,\allowbreak 20}}
  \proofstepwidestarbreakkeep[1,2,10]{a\star x=b\leftrightarrow F(a)\diamond F(x)=F(b)}{\rLRI{16,\allowbreak 21}}

\closeproofpartwide
\proofpartwide{Lösungsmengenkriterium und eingeschränkte Bijektion}
\setcounter{proofstepnr}{22}
  \proofstepwidestarbreakkeep[1,10]{F(x)\in B}{%
    \FormulaRefAuto{F\colon A\bij B\dsep x\in A\vdash F(x)\in B}[thm]{3,\allowbreak 10}}
  \proofstepwidestar[1,2,10]{L_A(a,b)=\{u\in A\mid a\star u=b\}}{%
      \FormulaRefAuto{SemigroupLeftEquationFiberDef}[def]{8,\allowbreak 2}}
  \proofstepwidestarbreakkeep[1,2,10]{x\in L_A(a,b)\leftrightarrow a\star x=b}{%
    \FormulaRefAuto{x\in\{u\in A\mid P(u)\}\eqvdash x\in A\land P(x)}[thm]{%
      24,\allowbreak 10}}
  \proofstepwidestar[1,2,10]{L_B(F(a),F(b))=\{v\in B\mid F(a)\diamond v=F(b)\}}{%
      \FormulaRefAuto{SemigroupLeftEquationFiberDef}[def]{8,\allowbreak 9}}
  \proofstepwidestarbreakkeep[1,2,10]{F(x)\in L_B(F(a),F(b))
    \leftrightarrow F(a)\diamond F(x)=F(b)}{%
    \FormulaRefAuto{x\in\{u\in A\mid P(u)\}\eqvdash x\in A\land P(x)}[thm]{%
      26,\allowbreak 23}}
  \proofstepwidestarbreakkeep[1,2]{\forall x\in A\,
    (x\in L_A(a,b)\leftrightarrow F(x)\in L_B(F(a),F(b)))}{%
    \rUI{\rRI{10,\allowbreak \rChain{25,\allowbreak 22,\allowbreak 27}}}}
  \proofstepwidestar[1,2]{L_A(a,b)=\{u\in A\mid a\star u=b\}}{%
      \FormulaRefAuto{SemigroupLeftEquationFiberDef}[def]{8,\allowbreak 2}}
  \proofstepwidestar[1,2]{L_B(F(a),F(b))=\{v\in B\mid F(a)\diamond v=F(b)\}}{%
      \FormulaRefAuto{SemigroupLeftEquationFiberDef}[def]{8,\allowbreak 9}}
  \proofstepwidestar[1,2]{L_A(a,b)\subseteq A}{%
      \FormulaRefAuto{\{x\in A\mid P(x)\}\subseteq A}[thm]{%
      29}}
  \proofstepwidestar[1,2]{L_B(F(a),F(b))\subseteq B}{%
      \FormulaRefAuto{\{x\in A\mid P(x)\}\subseteq A}[thm]{%
      30}}
  \proofstepwidestarbreakkeep[1,2]{L_A(a,b)\subseteq A\land
    L_B(F(a),F(b))\subseteq B}{%
    \rAI{31,\allowbreak %
      32}}
  \proofstepwidestarbreakkeep[1,2]{F\restriction_{L_A(a,b)}\colon
    L_A(a,b)\bij L_B(F(a),F(b))}{%
    \FormulaRefAuto{BijectiveRestrictionIffPointwiseCompatibility}[thm]{3,\allowbreak 33,\allowbreak 28}}
\closeproofpartwide
\end{tabproofsplitwide}

\subsection{Rekonstruktion mit festgelegten Einermengenbildern}

\par\noindent\begin{minipage}{\linewidth}
\FormulaThmDeltaK[Die Einermengenkopie liegt in der Potenzhalbgruppe]{%
  \SemigroupAlg{A}{\star}\vdash
  \operatorname{Sing}(A)\subseteq\PowNE{A}%
}{SemigroupSingletonCopySubset}{%
  \DeltaRow{Mengen}{A\dsep X\dsep a}
  \DeltaRow{Funktionen}{\star}
}
\end{minipage}\par
\begin{tabproofwide}
  \proofstepwidestarbreakkeep[1]{\SemigroupAlg{A}{\star}}{\rA}
  \proofstepwidestarbreakkeep[2]{X\in\operatorname{Sing}(A)}{\rA}
  \proofstepwidestarbreakkeep[1,2]{\exists a\in A\,(X=\{a\})}{%
    \FormulaRefAuto{SemigroupSingletonCopyDef}[def]{1,\allowbreak 2}}
  \proofstepwidestarbreakkeep[4]{a\in A\land X=\{a\}}{\rA}
  \proofstepwidestarbreakkeep[4]{\{a\}\subseteq A}{%
    \FormulaRefAuto{a\in A\vdash\{a\}\subseteq A}[thm]{4}}
  \proofstepwidestar[]{a\in\{a\}}{%
      \FormulaRefAuto{a\in\{a\}}[thm]}
  \proofstepwidestarbreakkeep[]{\{a\}\neq\varnothing}{%
    \FormulaRefAuto{a\in A\vdash A\neq\varnothing}[thm]{%
      6}}
  \proofstepwidestar[4]{\{a\}\in\PowNE{A}}{%
      \FormulaRefAuto{NonemptyPowersetMembership}[thm]{\rAI{5,\allowbreak 7}}}
  \proofstepwidestarbreakkeep[4]{X\in\PowNE{A}}{%
    \rIE{4,\allowbreak 8}}
  \proofstepwidestarbreakkeep[1,2]{X\in\PowNE{A}}{\rEE{3,\allowbreak 4,\allowbreak 9}}
  \proofstepwidestarbreakkeep[1]{\operatorname{Sing}(A)\subseteq\PowNE{A}}{%
    \FormulaRefAuto{SubsetIntroductionRule}[thm]{[2]\vdots 10}}
\end{tabproofwide}

\par\noindent\begin{minipage}{\linewidth}
\FormulaThmDeltaK[Rekonstruktion mit vorgeschriebenen Einermengenbildern]{%
  \begin{aligned}[t]
    &\SemigroupAlg{A}{\star}\dsep\SemigroupAlg{B}{\diamond}\dsep{}\\[-2pt]
    &\SemigroupIso{F}{\PowNE{A},\star_{\mathcal P}}
      {\PowNE{B},\diamond_{\mathcal P}}\dsep
      F[\operatorname{Sing}(A)]=\operatorname{Sing}(B)\vdash{}\\[-2pt]
    &\exists f\,\Bigl(\SemigroupIso{f}{A,\star}{B,\diamond}
      \land\forall a\in A\,\bigl(F(\{a\})=\{f(a)\}\bigr)\Bigr)
  \end{aligned}%
}{PowerSemigroupSingletonLayerReconstructionWitness}{%
  \DeltaRow{Mengen}{A\dsep B\dsep a\dsep b}
  \DeltaRow{Funktionen}{\star\dsep\diamond\dsep F\dsep R\dsep f
    \dsep\eta_A\dsep\eta_B}
}
\end{minipage}\par
Die Einschränkung von \(F\) auf die Einermengenkopie wird mit den beiden
Einermengenabbildungen zurück auf die Grundträger transportiert. Setze
für den folgenden Beweis
\[
 R\coloneqq F\restriction_{\operatorname{Sing}(A)},\qquad
 f\coloneqq\eta_B^{-1}\circ(R\circ\eta_A).
\]
Die punktweise Gleichung wird mitgeführt, damit spätere
Rekonstruktionsschritte genau diesen Isomorphismus verwenden können.
\begin{tabproofsplitwide}
\proofpartwide{Konstruktion der bijektiven Trägerabbildung}
  \proofstepwidestarbreakkeep[1]{\SemigroupAlg{A}{\star}}{\rA}
  \proofstepwidestarbreakkeep[2]{\SemigroupAlg{B}{\diamond}}{\rA}
  \proofstepwidestarbreakkeep[3]{\SemigroupIso{F}{\PowNE{A},\star_{\mathcal P}}
    {\PowNE{B},\diamond_{\mathcal P}}}{\rA}
  \proofstepwidestarbreakkeep[4]{F[\operatorname{Sing}(A)]=\operatorname{Sing}(B)}{\rA}
  \proofstepwidestarbreakkeep[1]{\SemigroupIso{\eta_A}{A,\star}
    {\operatorname{Sing}(A),\star_{\mathcal P}}}{%
    \FormulaRefAuto{SemigroupSingletonCopyIsomorphism}[thm]{1}}
  \proofstepwidestarbreakkeep[2]{\SemigroupIso{\eta_B}{B,\diamond}
    {\operatorname{Sing}(B),\diamond_{\mathcal P}}}{%
    \FormulaRefAuto{SemigroupSingletonCopyIsomorphism}[thm]{2}}
  \proofstepwidestar[1,2,3,4]{\SemigroupHom{\eta_A}{A,\star}{\operatorname{Sing}(A),\star_{\mathcal P}}}{%
      \FormulaRefAuto{SemigroupIsoHomAxiom}[ax]{5}}
  \proofstepwidestar[1,2,3,4]{\SemigroupHom{\eta_B}{B,\diamond}{\operatorname{Sing}(B),\diamond_{\mathcal P}}}{%
      \FormulaRefAuto{SemigroupIsoHomAxiom}[ax]{6}}
  \proofstepwidestar[1,2,3,4]{\operatorname{Sing}(A)\subseteq\PowNE{A}}{%
      \FormulaRefAuto{SemigroupSingletonCopySubset}[thm]{1}}
  \proofstepwidestar[1,2,3,4]{\operatorname{Sing}(B)\subseteq\PowNE{B}}{%
      \FormulaRefAuto{SemigroupSingletonCopySubset}[thm]{2}}
  \proofstepwidestar[1,2,3,4]{\SemigroupAlg{\operatorname{Sing}(A)}{\star_{\mathcal P}}}{%
      \FormulaRefAuto{SemigroupHomTargetAxiom}[ax]{%
        7}}
  \proofstepwidestar[1,2,3,4]{\SemigroupAlg{\operatorname{Sing}(B)}{\diamond_{\mathcal P}}}{%
      \FormulaRefAuto{SemigroupHomTargetAxiom}[ax]{%
        8}}
  \proofstepwidestarbreakkeep[1,2,3,4]{\SemigroupIso{R}
    {\operatorname{Sing}(A),\star_{\mathcal P}}
    {\operatorname{Sing}(B),\diamond_{\mathcal P}}}{%
    \FormulaRefAuto{SemigroupIsoRestrictionToImage}[thm]{3,\allowbreak %
      11,\allowbreak %
      12,\allowbreak %
      9,\allowbreak %
      10,\allowbreak 4}}
  \proofstepwidestar[1,2,3,4]{\SemigroupIso{R\circ\eta_A}{A,\star}{\operatorname{Sing}(B),\diamond_{\mathcal P}}}{%
      \FormulaRefAuto{SemigroupIsoComposition}[thm]{5,\allowbreak 13}}
  \proofstepwidestarbreakkeep[1,2,3,4]{R\circ\eta_A\colon A\bij\operatorname{Sing}(B)}{%
    \FormulaRefAuto{SemigroupIsoBijectiveAxiom}[ax]{%
      14}}
  \proofstepwidestarbreakkeep[2]{\eta_B\colon B\bij\operatorname{Sing}(B)}{%
    \FormulaRefAuto{SemigroupIsoBijectiveAxiom}[ax]{6}}
  \proofstepwidestarbreakkeep[2]{\eta_B^{-1}\colon\operatorname{Sing}(B)\bij B}{%
    \FormulaRefAuto{InverseFunctionBijection}[thm]{16}}
  \proofstepwidestarbreakkeep[1,2,3,4]{f\colon A\bij B}{%
    \FormulaRefAuto{F\colon A\bij B\dsep G\colon B\bij C\vdash
      G\circ F\colon A\bij C}[thm]{15,\allowbreak 17}}

\closeproofpartwide
\proofpartwide{Vorgeschriebene Einermengenbilder}
\setcounter{proofstepnr}{18}
  \proofstepwidestarbreakkeep[19]{a\in A}{\rA}
  \proofstepwidestar[1,2,3,4,19]{f\colon A\to B}{%
      \FormulaRefAuto{Bijektive Funktion}[def]{18}}
  \proofstepwidestarbreakkeep[1,2,3,4,19]{f(a)\in B}{%
    \FormulaRefAuto{F\colon A\to B\dsep x\in A\vdash F(x)\in B}[thm]{%
      20,\allowbreak 19}}
  \proofstepwidestarbreakkeep[1,19]{\{a\}\in\operatorname{Sing}(A)}{%
    \FormulaRefAuto{SemigroupSingletonCopyDef}[def]{1,\allowbreak 19}}
  \proofstepwidestar[1,2,3,4,19]{F\colon \PowNE{A}\bij \PowNE{B}}{%
      \FormulaRefAuto{SemigroupIsoBijectiveAxiom}[ax]{3}}
  \proofstepwidestar[1,2,3,4,19]{\operatorname{Sing}(A)\subseteq\PowNE{A}}{%
      \FormulaRefAuto{SemigroupSingletonCopySubset}[thm]{1}}
  \proofstepwidestar[1,2,3,4,19]{F\colon \PowNE{A}\to \PowNE{B}}{%
      \FormulaRefAuto{Bijektive Funktion}[def]{%
        23}}
  \proofstepwidestar[1,2,3,4,19]{F(\{a\})\in F[\operatorname{Sing}(A)]}{%
      \FormulaRefAuto{C\subseteq A\dsep F\colon A\to B\dsep
      x\in C\vdash F(x)\in F[C]}[thm]{%
      24,\allowbreak %
      25,\allowbreak 22}}
  \proofstepwidestarbreakkeep[1,2,3,4,19]{F(\{a\})\in\operatorname{Sing}(B)}{%
    \rIE{4,\allowbreak 26}}
  \proofstepwidestarbreakkeep[1,2,3,4,19]{f(a)=\eta_B^{-1}(F(\{a\}))}{%
    \FormulaRefAuto{x\in A\vdash(F\circ G)(x)=F(G(x))}[thm]{15,\allowbreak 17,\allowbreak 19},\allowbreak %
    \FormulaRefAuto{RestrictionDef}[def]{22},\allowbreak %
    \FormulaRefAuto{SemigroupSingletonCopyDef}[def]{1,\allowbreak 19}}
  \proofstepwidestarbreakkeep[1,2,3,4,19]{\{f(a)\}=\eta_B(f(a))}{%
    \FormulaRefAuto{SemigroupSingletonCopyDef}[def]{2,\allowbreak 21}}
  \proofstepwidestarbreakkeep[1,2,3,4,19]{\phantom{\{f(a)\}}=\eta_B(\eta_B^{-1}(F(\{a\})))}{\rIE{28}}
  \proofstepwidestarbreakkeep[1,2,3,4,19]{\phantom{\{f(a)\}}=F(\{a\})}{%
    \FormulaRefAuto{InverseFunctionRightInverse}[thm]{16,\allowbreak 27}}
  \proofstepwidestar[1,2,3,4,19]{F(\{a\})=\{f(a)\}}{%
      \FormulaRefAuto{a=b\vdash b=a}[thm]{\rChain{29,\allowbreak 31}}}
  \proofstepwidestarbreakkeep[1,2,3,4]{\forall a\in A\,(F(\{a\})=\{f(a)\})}{%
    \rUI{\rRI{19,\allowbreak 32}}}

\closeproofpartwide
\proofpartwide{Operationserhaltung und Existenzschluss}
\setcounter{proofstepnr}{33}
  \proofstepwidestarbreakkeep[34]{a,b\in A}{\rA}
  \proofstepwidestarbreakkeep[1,34]{a\star b\in A}{%
    \FormulaRefAuto{SemigroupClosureAxiom}[ax]{1,\allowbreak 34}}
  \proofstepwidestar[1,2,3,4,34]{f\colon A\to B}{%
      \FormulaRefAuto{Bijektive Funktion}[def]{18}}
  \proofstepwidestarbreakkeep[1,2,3,4,34]{f(a),f(b)\in B}{%
    \FormulaRefAuto{F\colon A\to B\dsep x\in A\vdash F(x)\in B}[thm]{%
      36,\allowbreak 34}}
  \proofstepwidestarbreakkeep[1,2,3,4,34]{\{f(a\star b)\}=F(\{a\star b\})}{%
    \FormulaRefAuto{a=b\vdash b=a}[thm]{\rRE{35,\allowbreak \rUE{33}}}}
  \proofstepwidestar[1,2,3,4,34]{\SemigroupHom{\eta_A}{A,\star}{\operatorname{Sing}(A),\star_{\mathcal P}}}{%
      \FormulaRefAuto{SemigroupIsoHomAxiom}[ax]{5}}
  \proofstepwidestar[1,2,3,4,34]{\forall u\in A\;(\eta_A(u)=\{u\})}{%
      \FormulaRefAuto{SemigroupSingletonCopyDef}[def]{1}}
  \proofstepwidestar[1,2,3,4,34]{\eta_A(a\star b)=\eta_A(a)\star_{\mathcal P}\eta_A(b)}{%
      \FormulaRefAuto{SemigroupHomOperationAxiom}[ax]{%
        39,\allowbreak 34}}
  \proofstepwidestarbreakkeep[1,2,3,4,34]{F(\{a\star b\})=F(\{a\}\star_{\mathcal P}\{b\})}{%
    \rIE{40,\allowbreak %
      41}}
  \proofstepwidestar[1,2,3,4,34]{\SemigroupHom{F}{\PowNE{A},\star_{\mathcal P}}{\PowNE{B},\diamond_{\mathcal P}}}{%
      \FormulaRefAuto{SemigroupIsoHomAxiom}[ax]{3}}
  \proofstepwidestar[1,2,3,4,34]{\operatorname{Sing}(A)\subseteq\PowNE{A}}{%
      \FormulaRefAuto{SemigroupSingletonCopySubset}[thm]{1}}
  \proofstepwidestar[1,2,3,4,34]{\{a\},\{b\}\in\operatorname{Sing}(A)}{%
      \FormulaRefAuto{SemigroupSingletonCopyDef}[def]{1,\allowbreak 34}}
  \proofstepwidestar[1,2,3,4,34]{\{a\},\{b\}\in\PowNE{A}}{%
      \FormulaRefAuto{A\subseteq B\dsep x\in A\vdash x\in B}[thm]{%
        44,\allowbreak %
        45}}
  \proofstepwidestarbreakkeep[1,2,3,4,34]{F(\{a\}\star_{\mathcal P}\{b\})=F(\{a\})\diamond_{\mathcal P}F(\{b\})}{%
    \FormulaRefAuto{SemigroupHomOperationAxiom}[ax]{%
      43,\allowbreak %
      46}}
  \proofstepwidestarbreakkeep[1,2,3,4,34]{F(\{a\})\diamond_{\mathcal P}F(\{b\})=\{f(a)\}\diamond_{\mathcal P}\{f(b)\}}{%
    \rIE{\rRE{34,\allowbreak \rUE{33}}}}
  \proofstepwidestar[1,2,3,4,34]{\SemigroupHom{\eta_B}{B,\diamond}{\operatorname{Sing}(B),\diamond_{\mathcal P}}}{%
      \FormulaRefAuto{SemigroupIsoHomAxiom}[ax]{6}}
  \proofstepwidestar[1,2,3,4,34]{\forall u\in B\;(\eta_B(u)=\{u\})}{%
      \FormulaRefAuto{SemigroupSingletonCopyDef}[def]{2}}
  \proofstepwidestar[1,2,3,4,34]{\eta_B(f(a)\diamond f(b))=\eta_B(f(a))\diamond_{\mathcal P}\eta_B(f(b))}{%
      \FormulaRefAuto{SemigroupHomOperationAxiom}[ax]{%
        49,\allowbreak 37}}
  \proofstepwidestarbreakkeep[1,2,3,4,34]{\{f(a)\}\diamond_{\mathcal P}\{f(b)\}=\{f(a)\diamond f(b)\}}{%
    \rIE{50,\allowbreak %
      51}}
  \proofstepwidestarbreakkeep[1,2,3,4,34]{f(a\star b)=f(a)\diamond f(b)}{%
    \FormulaRefAuto{\{a\}=\{b\}\eqvdash a=b}[thm]{\rChain{38,42,47,48,52}}}
  \proofstepwidestar[1,2,3,4]{f\colon A\to B}{%
      \FormulaRefAuto{Bijektive Funktion}[def]{18}}
  \proofstepwidestar[1,2,3,4]{\SemigroupHom{f}{A,\star}{B,\diamond}}{%
      \FormulaRefAuto{SemigroupHomDef}[def]{1,\allowbreak 2,\allowbreak %
        54,\allowbreak 53}}
  \proofstepwidestarbreakkeep[1,2,3,4]{\SemigroupIso{f}{A,\star}{B,\diamond}}{%
    \FormulaRefAuto{SemigroupIsoDef}[def]{%
      55,\allowbreak 18}}
  \proofstepwidestarbreakkeep[1,2,3,4]{\exists f\,\Bigl(
    \begin{aligned}[t]
      &\SemigroupIso{f}{A,\star}{B,\diamond}\land{}\\[-2pt]
      &\forall a\in A\,(F(\{a\})=\{f(a)\})
    \end{aligned}\Bigr)}{\rEI{\rAI{56,\allowbreak 33}}}
\closeproofpartwide
\end{tabproofsplitwide}

% Allgemeine Isomorphieregeln; eingebunden nach den Potenzmorphismen.
\section{Beispiele von Isomorphismen}

Ein Isomorphismus liefert durch Einschränkung, mengenweise Fortsetzung,
koordinatenweise Anwendung und Transport von Relationen weitere
Isomorphismen. Im Folgenden steht \(f:A\cong B\) abkürzend für
\(\SemigroupIso{f}{A,\star}{B,\diamond}\); die jeweilige Operation
auf einem Teilträger ist stets die Einschränkung der umgebenden Operation.
Die Abkürzung \(h=f^{-1}\) bezeichnet die Umkehrabbildung.
Bei Multiplikationsketten schreiben wir die Operationen auch durch
Nebeneinanderschreiben; auf \(A\) ist damit \(\star\), auf \(B\)
\(\diamond\) gemeint. In den nummerierten Teilbeweisen werden zunächst
Träger, dann Operation und schließlich Bijektivität nachgewiesen;
eine bereits abgeleitete Teilbehauptung wird im nachfolgenden Teil benutzt.
Die Identität und die Komposition von Isomorphismen sind bereits in
\FormulaRefAuto{SemigroupIdentityIsomorphism}[thm] und
\FormulaRefAuto{SemigroupIsoComposition}[thm] bewiesen.

\subsection{Umkehrabbildung und Unterstrukturen}

\FormulaThmDeltaK[Die Umkehrabbildung eines Isomorphismus]{%
  \SemigroupIso{f}{A,\star}{B,\diamond}\vdash
  \SemigroupIso{f^{-1}}{B,\diamond}{A,\star}%
}{SemigroupIsoInverse}{%
  \DeltaRow{Mengen}{A\dsep B\dsep x\dsep y\dsep u\dsep v}
  \DeltaRow{Funktionen}{f\dsep f^{-1}\dsep\star\dsep\diamond}
}
\begin{tabproofsplitwide}
\proofpartwideindR[Bijektive Umkehrabbildung]{%
 f:A\cong B\vdash h:B\bij A%
}{f:A\cong B\vdash h:B\bij A}[SemigroupIsoInverseCarrier]
  \proofstepwidestar[1]{f:A\cong B}{\rA}
  \proofstepwidestar[1]{f:A\bij B}{\FormulaRefAuto{SemigroupIsoBijectiveAxiom}[ax]{1}}
  \proofstepwidestar[1]{h:B\bij A}{\FormulaRefAuto{InverseFunctionBijection}[thm]{2}}
\closeproofpartwideindR
\proofpartwideindR[Operationserhaltung]{%
 f:A\cong B\dsep u,v\in B\vdash h(u\diamond v)=h(u)\star h(v)%
}{f:A\cong B\dsep u,v\in B\vdash h(u\diamond v)=h(u)\star h(v)}[SemigroupIsoInverseOperation]
  \proofstepwidestar[1]{f:A\cong B}{\rA}
  \proofstepwidestar[2]{u\in B}{\rA}
  \proofstepwidestar[3]{v\in B}{\rA}
  \proofstepwidestar[1]{\SemigroupHom{f}{A,\star}{B,\diamond}}{\FormulaRefAuto{SemigroupIsoHomAxiom}[ax]{1}}
  \proofstepwidestar[1]{f:A\bij B}{\FormulaRefAuto{SemigroupIsoBijectiveAxiom}[ax]{1}}
  \proofstepwidestar[1]{h:B\bij A}{\FormulaRefAuto{SemigroupIsoInverseCarrier}[thm]{1}}
  \proofstepwidestar[1]{h:B\to A}{\FormulaRefAuto{Bijektive Funktion}[def]{6}}
  \proofstepwidestar[1,2]{h(u)\in A}{\FormulaRefAuto{F\colon A\to B\dsep x\in A\vdash F(x)\in B}[thm]{7,2}}
  \proofstepwidestar[1,3]{h(v)\in A}{\FormulaRefAuto{F\colon A\to B\dsep x\in A\vdash F(x)\in B}[thm]{7,3}}
  \proofstepwidestar[1]{\SemigroupAlg{A}{\star}}{\FormulaRefAuto{SemigroupHomSourceAxiom}[ax]{4}}
  \proofstepwidestar[1,2,3]{h(u)\star h(v)\in A}{\FormulaRefAuto{SemigroupClosureAxiom}[ax]{10,8,9}}
  \proofstepwidestar[1,2,3]{f(h(u)\star h(v))=f(h(u))\diamond f(h(v))}{\FormulaRefAuto{SemigroupHomOperationAxiom}[ax]{4,8,9}}
  \proofstepwidestar[1,2]{f(h(u))=u}{\FormulaRefAuto{InverseFunctionRightInverse}[thm]{5,2}}
  \proofstepwidestar[1,3]{f(h(v))=v}{\FormulaRefAuto{InverseFunctionRightInverse}[thm]{5,3}}
  \proofstepwidestar[1,2,3]{f(h(u)\star h(v))=u\diamond v}{\rIE{14,\rIE{13,12}}}
  \proofstepwidestar[1,2,3]{h(f(h(u)\star h(v)))=h(u)\star h(v)}{\FormulaRefAuto{InverseFunctionLeftInverse}[thm]{5,11}}
  \proofstepwidestar[1,2,3]{h(u\diamond v)=h(u)\star h(v)}{\rIE{15,16}}
\closeproofpartwideindR
\proofpartwide{Isomorphieschluss}
  \proofstepwidestar[1]{f:A\cong B}{\rA}
  \proofstepwidestar[1]{\SemigroupHom{f}{A,\star}{B,\diamond}}{\FormulaRefAuto{SemigroupIsoHomAxiom}[ax]{1}}
  \proofstepwidestar[1]{\SemigroupAlg{A}{\star}}{\FormulaRefAuto{SemigroupHomSourceAxiom}[ax]{2}}
  \proofstepwidestar[1]{\SemigroupAlg{B}{\diamond}}{\FormulaRefAuto{SemigroupHomTargetAxiom}[ax]{2}}
  \proofstepwidestar[1]{h:B\bij A}{\FormulaRefAuto{SemigroupIsoInverseCarrier}[thm]{1}}
  \proofstepwidestar[1]{h:B\to A}{\FormulaRefAuto{Bijektive Funktion}[def]{5}}
  \proofstepwidestar[7]{u\in B}{\rA}
 \proofstepwidestar[8]{v\in B}{\rA}
 \proofstepwidestar[7,8]{u,v\in B}{\rAIStar{7,8}}
  \proofstepwidestar[1,7,8]{h(u\diamond v)=h(u)\star h(v)}{\FormulaRefAuto{SemigroupIsoInverseOperation}[thm]{1,9}}
  \proofstepwidestar[1]{\forall u,v\in B\;(h(u\diamond v)=h(u)\star h(v))}{\rUI{\rRI{7,\rUI{\rRI{8,10}}}}}
  \proofstepwidestar[1]{\SemigroupHom{h}{B,\diamond}{A,\star}}{\FormulaRefAuto{SemigroupHomDef}[def]{4,3,6,11}}
  \proofstepwidestar[1]{\SemigroupIso{h}{B,\diamond}{A,\star}}{\FormulaRefAuto{SemigroupIsoDef}[def]{12,5}}
\closeproofpartwide
\end{tabproofsplitwide}

\FormulaThmDeltaKR[Elementare Rechenzugriffe auf einen Isomorphismus]{%
 \begin{aligned}[t]
  &f:A\cong B\dsep x,y\in A\dsep U\subseteq A\vdash{}\\[-2pt]
  &\text{(i)}\quad \SemigroupAlg{A}{\star}\\[-2pt]
  &\text{(ii)}\quad \SemigroupAlg{B}{\diamond}\\[-2pt]
  &\text{(iii)}\quad f:A\to B\\[-2pt]
  &\text{(iv)}\quad f(x\star y)=f(x)\diamond f(y)\\[-2pt]
  &\text{(v)}\quad x=y\leftrightarrow f(x)=f(y)\\[-2pt]
  &\text{(vi)}\quad x\in U\leftrightarrow f(x)\in f[U]
 \end{aligned}%
}{%
 \begin{aligned}[t]
 &f:A\cong B\dsep x,y\in A\dsep U\subseteq A\vdash{}\\[-2pt]
 &\SemigroupAlg{A}{\star},\quad\SemigroupAlg{B}{\diamond},\quad f:A\to B,\\[-2pt]
 &f(x\star y)=f(x)\diamond f(y),\\[-2pt]
 &x=y\leftrightarrow f(x)=f(y),\qquad
      x\in U\leftrightarrow f(x)\in f[U].
 \end{aligned}%
}{SemigroupIsoElementaryTransport}{%
 \DeltaRow{Mengen}{A\dsep B\dsep U\dsep x\dsep y\dsep z}
 \DeltaRow{Funktionen}{f\dsep\star\dsep\diamond}
}
\Needspace{8\baselineskip}
\begin{tabproofsplitwide}
\proofpartwideindR[Quellhalbgruppe]{%
 f:A\cong B\vdash\SemigroupAlg{A}{\star}%
}{f:A\cong B\vdash\SemigroupAlg{A}{\star}}[SemigroupIsoSource]
 \proofstepwidestar[1]{f:A\cong B}{\rA}
 \proofstepwidestar[1]{\SemigroupHom{f}{A,\star}{B,\diamond}}{\FormulaRefAuto{SemigroupIsoHomAxiom}[ax]{1}}
 \proofstepwidestar[1]{\SemigroupAlg{A}{\star}}{\FormulaRefAuto{SemigroupHomSourceAxiom}[ax]{2}}
\closeproofpartwideindR
\proofpartwideindR[Zielhalbgruppe]{%
 f:A\cong B\vdash\SemigroupAlg{B}{\diamond}%
}{f:A\cong B\vdash\SemigroupAlg{B}{\diamond}}[SemigroupIsoTarget]
 \proofstepwidestar[1]{f:A\cong B}{\rA}
 \proofstepwidestar[1]{\SemigroupHom{f}{A,\star}{B,\diamond}}{\FormulaRefAuto{SemigroupIsoHomAxiom}[ax]{1}}
 \proofstepwidestar[1]{\SemigroupAlg{B}{\diamond}}{\FormulaRefAuto{SemigroupHomTargetAxiom}[ax]{2}}
\closeproofpartwideindR
\proofpartwideindR[Trägerabbildung]{%
 f:A\cong B\vdash f:A\to B%
}{f:A\cong B\vdash f:A\to B}[SemigroupIsoFunction]
 \proofstepwidestar[1]{f:A\cong B}{\rA}
 \proofstepwidestar[1]{\SemigroupHom{f}{A,\star}{B,\diamond}}{\FormulaRefAuto{SemigroupIsoHomAxiom}[ax]{1}}
 \proofstepwidestar[1]{f:A\to B}{\FormulaRefAuto{SemigroupHomFunctionAxiom}[ax]{2}}
\closeproofpartwideindR
\proofpartwideindR[Vollständiges Bild]{%
 f:A\cong B\vdash f[A]=B%
}{f:A\cong B\vdash f[A]=B}[SemigroupIsoFullImage]
 \proofstepwidestar[1]{f:A\cong B}{\rA}
 \proofstepwidestar[1]{f:A\bij B}{\FormulaRefAuto{SemigroupIsoBijectiveAxiom}[ax]{1}}
 \proofstepwidestar[1]{f:A\sur B}{\FormulaRefAuto{F\colon A\bij B\vdash F\colon A\sur B}[thm]{2}}
 \proofstepwidestar[1]{f:A\to B}{\FormulaRefAuto{SemigroupIsoFunction}[thm]{1}}
 \proofstepwidestar[1]{f[A]=B}{\FormulaRefAuto{SurjectiveImageEqualsCodomain}[thm]{4,3}}
\closeproofpartwideindR
\proofpartwideindR[Produktwerte]{%
 f:A\cong B\dsep x,y\in A\vdash f(x\star y)=f(x)\diamond f(y)%
}{f:A\cong B\dsep x,y\in A\vdash f(x\star y)=f(x)\diamond f(y)}[SemigroupIsoOperation]
 \proofstepwidestar[1]{f:A\cong B}{\rA}
 \proofstepwidestar[2]{x,y\in A}{\rA}
 \proofstepwidestar[1]{\SemigroupHom{f}{A,\star}{B,\diamond}}{\FormulaRefAuto{SemigroupIsoHomAxiom}[ax]{1}}
 \proofstepwidestar[1,2]{f(x\star y)=f(x)\diamond f(y)}{\FormulaRefAuto{SemigroupHomOperationAxiom}[ax]{3,2}}
\closeproofpartwideindR
\proofpartwideindR[Gleichheitsreflexion]{%
 f:A\cong B\dsep x,y\in A\vdash x=y\leftrightarrow f(x)=f(y)%
}{f:A\cong B\dsep x,y\in A\vdash x=y\leftrightarrow f(x)=f(y)}[SemigroupIsoEqualityReflection]
 \proofstepwidestar[1]{f:A\cong B}{\rA}
 \proofstepwidestar[2]{x,y\in A}{\rA}
 \proofstepwidestar[1]{f:A\bij B}{\FormulaRefAuto{SemigroupIsoBijectiveAxiom}[ax]{1}}
 \proofstepwidestar[1]{f:A\inj B}{\FormulaRefAuto{F\colon A\bij B\vdash F\colon A\inj B}[thm]{3}}
 \proofstepwidestar[5]{x=y}{\rA}
 \proofstepwidestar[5]{f(x)=f(y)}{\rIE{5,\rII}}
 \proofstepwidestar[]{x=y\Rightarrow f(x)=f(y)}{\rRI{5,6}}
 \proofstepwidestar[8]{f(x)=f(y)}{\rA}
 \proofstepwidestar[1,2,8]{x=y}{\FormulaRefAuto{F\colon A\inj B\dsep x\in A\dsep y\in A\dsep F(x)=F(y)\vdash x=y}[ax]{4,2,8}}
 \proofstepwidestar[1,2]{f(x)=f(y)\Rightarrow x=y}{\rRI{8,9}}
 \proofstepwidestar[1,2]{x=y\leftrightarrow f(x)=f(y)}{\rLRI{7,10}}
\closeproofpartwideindR
\proofpartwideindR[Bildzugehörigkeit]{%
 f:A\cong B\dsep U\subseteq A\dsep x\in A\vdash
 x\in U\leftrightarrow f(x)\in f[U]%
}{f:A\cong B\dsep U\subseteq A\dsep x\in A\vdash x\in U\leftrightarrow f(x)\in f[U]}[SemigroupIsoImageMembership]
  \proofstepwidestar[1]{f:A\cong B}{\rA}
  \proofstepwidestar[2]{U\subseteq A}{\rA}
  \proofstepwidestar[3]{x\in A}{\rA}
  \proofstepwidestar[1]{f:A\to B}{\FormulaRefAuto{SemigroupIsoFunction}[thm]{1}}
  \proofstepwidestar[5]{x\in U}{\rA}
  \proofstepwidestar[1,2,5]{f(x)\in f[U]}{\FormulaRefAuto{C\subseteq A\dsep F\colon A\to B\dsep x\in C\vdash F(x)\in F[C]}[thm]{2,4,5}}
  \proofstepwidestar[1,2]{x\in U\rightarrow f(x)\in f[U]}{\rRI{5,6}}
  \proofstepwidestar[8]{f(x)\in f[U]}{\rA}
  \proofstepwidestar[1,2,8]{\exists z\in U\;(f(x)=f(z))}{\FormulaRefAuto{C\subseteq A\dsep F\colon A\to B\dsep y\in F[C]\vdash\exists x\in C\,y=F(x)}[thm]{2,4,8}}
  \proofstepwidestar[10]{z\in U\land f(x)=f(z)}{\rA}
  \proofstepwidestar[10]{z\in U}{\rAEa{10}}
  \proofstepwidestar[10]{f(x)=f(z)}{\rAEb{10}}
  \proofstepwidestar[2,10]{z\in A}{\FormulaRefAuto{A\subseteq B,\,x\in A\vdash x\in B}[thm]{2,11}}
  \proofstepwidestar[1,2,3,10]{x=z\leftrightarrow f(x)=f(z)}{\FormulaRefAuto{SemigroupIsoEqualityReflection}[thm]{1,3,13}}
  \proofstepwidestar[1,2,3,10]{x=z}{\rRE{\rLREb{14},12}}
  \proofstepwidestar[1,2,3,10]{z=x}{\FormulaRefAuto{a=b\vdash b=a}[thm]{15}}
  \proofstepwidestar[1,2,3,10]{x\in U}{\rIE{16,11}}
  \proofstepwidestar[1,2,3,8]{x\in U}{\rEE{9,10,17}}
  \proofstepwidestar[1,2,3]{f(x)\in f[U]\rightarrow x\in U}{\rRI{8,18}}
  \proofstepwidestar[1,2,3]{x\in U\leftrightarrow f(x)\in f[U]}{\rLRI{7,19}}
\closeproofpartwideindR
\end{tabproofsplitwide}

\FormulaThmDeltaKR[Transport von Unterhalbgruppen und Idealen]{%
 \begin{aligned}[t]
  &f:A\cong B\dsep U\subseteq A\vdash{}\\[-2pt]
  &\text{(i)}\quad \mathsf{UHGr}(U;A,\star)\leftrightarrow \mathsf{UHGr}(f[U];B,\diamond)\\[-2pt]
  &\text{(ii)}\quad \mathsf{LId}(U;A,\star)\leftrightarrow \mathsf{LId}(f[U];B,\diamond)\\[-2pt]
  &\text{(iii)}\quad \mathsf{RId}(U;A,\star)\leftrightarrow \mathsf{RId}(f[U];B,\diamond)\\[-2pt]
  &\text{(iv)}\quad \mathsf{Id}(U;A,\star)\leftrightarrow \mathsf{Id}(f[U];B,\diamond)
 \end{aligned}%
}{%
 \begin{aligned}[t]
  &f:A\cong B\dsep U\subseteq A\vdash{}\\[-2pt]
  &\mathsf{UHGr}(U;A,\star)\leftrightarrow
    \mathsf{UHGr}(f[U];B,\diamond),\\[-2pt]
  &\mathsf{LId}(U;A,\star)\leftrightarrow
    \mathsf{LId}(f[U];B,\diamond),\\[-2pt]
  &\mathsf{RId}(U;A,\star)\leftrightarrow
    \mathsf{RId}(f[U];B,\diamond),\\[-2pt]
  &\mathsf{Id}(U;A,\star)\leftrightarrow
    \mathsf{Id}(f[U];B,\diamond).
 \end{aligned}%
}{SemigroupIsoSubstructureTransport}{%
 \DeltaRow{Mengen}{A\dsep B\dsep U\dsep x\dsep y\dsep a\dsep b}
 \DeltaRow{Funktionen}{f\dsep h\dsep\star\dsep\diamond}
}
\begin{tabproofsplitwide}
\proofpartwideindR[Bildträger, Rückbild und Nichtleerheit]{%
f:A\cong B\dsep U\subseteq A\vdash f[U]\subseteq B\land h[f[U]]=U\land(U\neq\varnothing\leftrightarrow f[U]\neq\varnothing)%
}{f:A\cong B\dsep U\subseteq A\vdash f[U]\subseteq B\land h[f[U]]=U\land(U\neq\varnothing\leftrightarrow f[U]\neq\varnothing)}[SemigroupIsoImageCarrierData]
 \proofstepwidestar[1]{f:A\cong B}{\rA}
 \proofstepwidestar[2]{U\subseteq A}{\rA}
 \proofstepwidestar[1]{\SemigroupHom{f}{A,\star}{B,\diamond}}{\FormulaRefAuto{SemigroupIsoHomAxiom}[ax]{1}}
 \proofstepwidestar[1]{f:A\to B}{\FormulaRefAuto{SemigroupHomFunctionAxiom}[ax]{3}}
 \proofstepwidestar[1]{f:A\bij B}{\FormulaRefAuto{SemigroupIsoBijectiveAxiom}[ax]{1}}
 \proofstepwidestar[1]{h:B\bij A}{\FormulaRefAuto{SemigroupIsoInverseCarrier}[thm]{1}}
 \proofstepwidestar[1]{h:B\to A}{\FormulaRefAuto{Bijektive Funktion}[def]{6}}
 \proofstepwidestar[1]{\SemigroupAlg{A}{\star}}{\FormulaRefAuto{SemigroupHomSourceAxiom}[ax]{3}}
 \proofstepwidestar[1]{\SemigroupAlg{B}{\diamond}}{\FormulaRefAuto{SemigroupHomTargetAxiom}[ax]{3}}
 \proofstepwidestar[1,2]{f[U]\subseteq B}{\FormulaRefAuto{C\subseteq A\dsep F\colon A\to B\vdash F[C]\subseteq B}[thm]{2,4}}
 \proofstepwidestar[11]{x\in A}{\rA}
 \proofstepwidestar[1,11]{h(f(x))=x}{\FormulaRefAuto{InverseFunctionLeftInverse}[thm]{5,11}}
 \proofstepwidestar[1]{\forall x\in A\;(h(f(x))=x)}{\rUI{\rRI{11,12}}}
 \proofstepwidestar[1,2]{h[f[U]]=U}{\FormulaRefAuto{LeftInverseImages}[thm]{4,7,13,2}}
 \proofstepwidestar[15]{U\neq\varnothing}{\rA}
 \proofstepwidestar[1,2,15]{U\subseteq A\land U\neq\varnothing}{\rAI{2,15}}
  \proofstepwidestar[1,2,15]{U\in\PowNE A}{\FormulaRefAuto{NonemptyPowersetMembership}[thm]{16}}
  \proofstepwidestar[1,2,15]{f[U]\in\PowNE B}{\FormulaRefAuto{NonemptyDirectImage}[thm]{4,17}}
  \proofstepwidestar[1,2,15]{f[U]\subseteq B\land f[U]\neq\varnothing}{\FormulaRefAuto{NonemptyPowersetMembership}[thm]{18}}
  \proofstepwidestar[1,2,15]{f[U]\neq\varnothing}{\rAEb{19}}
 \proofstepwidestar[1,2]{U\neq\varnothing\Rightarrow f[U]\neq\varnothing}{\rRI{15,20}}
 \proofstepwidestar[22]{f[U]\neq\varnothing}{\rA}
 \proofstepwidestar[1,2,22]{f[U]\subseteq B\land f[U]\neq\varnothing}{\rAI{10,22}}
  \proofstepwidestar[1,2,22]{f[U]\in\PowNE B}{\FormulaRefAuto{NonemptyPowersetMembership}[thm]{23}}
  \proofstepwidestar[1,2,22]{h[f[U]]\in\PowNE A}{\FormulaRefAuto{NonemptyDirectImage}[thm]{7,24}}
  \proofstepwidestar[1,2,22]{h[f[U]]\subseteq A\land h[f[U]]\neq\varnothing}{\FormulaRefAuto{NonemptyPowersetMembership}[thm]{25}}
  \proofstepwidestar[1,2,22]{U\neq\varnothing}{\rIE{14,\rAEb{26}}}
 \proofstepwidestar[1,2]{f[U]\neq\varnothing\Rightarrow U\neq\varnothing}{\rRI{22,27}}
 \proofstepwidestar[1,2]{U\neq\varnothing\leftrightarrow f[U]\neq\varnothing}{\rLRI{21,28}}
 \proofstepwidestar[1,2]{f[U]\subseteq B\land h[f[U]]=U}{\rAI{10,14}}
 \proofstepwidestar[1,2]{f[U]\subseteq B\land h[f[U]]=U\land(U\neq\varnothing\leftrightarrow f[U]\neq\varnothing)}{\rAI{10,\rAI{14,29}}}
\closeproofpartwideindR
\proofpartwideindR[Urbild eines Bildzeugen]{%
f:A\cong B\dsep U\subseteq A\dsep u\in f[U]\vdash h(u)\in U%
}{f:A\cong B\dsep U\subseteq A\dsep u\in f[U]\vdash h(u)\in U}[SemigroupIsoInverseImageMember]
 \proofstepwidestar[1]{f:A\cong B}{\rA}
 \proofstepwidestar[2]{U\subseteq A}{\rA}
 \proofstepwidestar[3]{u\in f[U]}{\rA}
 \proofstepwidestar[1]{\SemigroupHom{f}{A,\star}{B,\diamond}}{\FormulaRefAuto{SemigroupIsoHomAxiom}[ax]{1}}
 \proofstepwidestar[1]{f:A\to B}{\FormulaRefAuto{SemigroupHomFunctionAxiom}[ax]{4}}
 \proofstepwidestar[1]{f:A\bij B}{\FormulaRefAuto{SemigroupIsoBijectiveAxiom}[ax]{1}}
 \proofstepwidestar[1]{h:B\bij A}{\FormulaRefAuto{SemigroupIsoInverseCarrier}[thm]{1}}
 \proofstepwidestar[1]{h:B\to A}{\FormulaRefAuto{Bijektive Funktion}[def]{7}}
 \proofstepwidestar[1]{\SemigroupAlg{A}{\star}}{\FormulaRefAuto{SemigroupHomSourceAxiom}[ax]{4}}
 \proofstepwidestar[1]{\SemigroupAlg{B}{\diamond}}{\FormulaRefAuto{SemigroupHomTargetAxiom}[ax]{4}}
 \proofstepwidestar[1,2]{f[U]\subseteq B}{\FormulaRefAuto{C\subseteq A\dsep F\colon A\to B\vdash F[C]\subseteq B}[thm]{2,5}}
 \proofstepwidestar[1,2,3]{u\in B}{\FormulaRefAuto{A\subseteq B,\,x\in A\vdash x\in B}[thm]{11,3}}
 \proofstepwidestar[1,2,3]{h(u)\in A}{\FormulaRefAuto{F\colon A\to B\dsep x\in A\vdash F(x)\in B}[thm]{8,12}}
 \proofstepwidestar[1,2,3]{f(h(u))=u}{\FormulaRefAuto{InverseFunctionRightInverse}[thm]{6,12}}
 \proofstepwidestar[1,2,3]{h(u)\in U\leftrightarrow f(h(u))\in f[U]}{\FormulaRefAuto{SemigroupIsoImageMembership}[thm]{1,2,13}}
 \proofstepwidestar[1,2,3]{u=f(h(u))}{\FormulaRefAuto{a=b\vdash b=a}[thm]{14}}
  \proofstepwidestar[1,2,3]{f(h(u))\in f[U]}{\rIE{16,3}}
 \proofstepwidestar[1,2,3]{h(u)\in U}{\rRE{\rLREb{15},17}}
\closeproofpartwideindR
\proofpartwideindR[Produktabschluss]{%
f:A\cong B\dsep \mathsf{UHGr}(U;A,\star)\vdash \mathsf{UHGr}(f[U];B,\diamond)%
}{f:A\cong B\dsep \mathsf{UHGr}(U;A,\star)\vdash \mathsf{UHGr}(f[U];B,\diamond)}[SemigroupIsoUHGrImage]
 \proofstepwidestar[1]{f:A\cong B}{\rA}
 \proofstepwidestar[2]{\mathsf{UHGr}(U;A,\star)}{\rA}
 \proofstepwidestar[1]{\SemigroupHom{f}{A,\star}{B,\diamond}}{\FormulaRefAuto{SemigroupIsoHomAxiom}[ax]{1}}
 \proofstepwidestar[1]{f:A\to B}{\FormulaRefAuto{SemigroupHomFunctionAxiom}[ax]{3}}
 \proofstepwidestar[1]{f:A\bij B}{\FormulaRefAuto{SemigroupIsoBijectiveAxiom}[ax]{1}}
 \proofstepwidestar[1]{h:B\bij A}{\FormulaRefAuto{SemigroupIsoInverseCarrier}[thm]{1}}
 \proofstepwidestar[1]{h:B\to A}{\FormulaRefAuto{Bijektive Funktion}[def]{6}}
 \proofstepwidestar[1]{\SemigroupAlg{A}{\star}}{\FormulaRefAuto{SemigroupHomSourceAxiom}[ax]{3}}
 \proofstepwidestar[1]{\SemigroupAlg{B}{\diamond}}{\FormulaRefAuto{SemigroupHomTargetAxiom}[ax]{3}}
 \proofstepwidestar[1,2]{U\subseteq A}{\FormulaRefAuto{SemigroupSubsemigroupSubsetAxiom}[ax]{8,2}}
 \proofstepwidestar[1,2]{U\neq\varnothing}{\FormulaRefAuto{SemigroupSubsemigroupNonemptyAxiom}[ax]{8,2}}
 \proofstepwidestar[1,2]{f[U]\subseteq B\land h[f[U]]=U\land(U\neq\varnothing\leftrightarrow f[U]\neq\varnothing)}{\FormulaRefAuto{SemigroupIsoImageCarrierData}[thm]{1,10}}
 \proofstepwidestar[1,2]{f[U]\subseteq B}{\rAEa{12}}
 \proofstepwidestar[1,2]{f[U]\neq\varnothing}{\rRE{\rLREa{\rAEb{\rAEb{12}}},11}}
 \proofstepwidestar[15]{u\in f[U]}{\rA}
 \proofstepwidestar[16]{v\in f[U]}{\rA}
 \proofstepwidestar[1,2,15]{h(u)\in U}{\FormulaRefAuto{SemigroupIsoInverseImageMember}[thm]{1,10,15}}
 \proofstepwidestar[1,2,16]{h(v)\in U}{\FormulaRefAuto{SemigroupIsoInverseImageMember}[thm]{1,10,16}}
 \proofstepwidestar[1,2,15,16]{h(u)\star h(v)\in U}{\FormulaRefAuto{SemigroupSubsemigroupClosureAxiom}[ax]{8,2,17,18}}
 \proofstepwidestar[1,2,15,16]{h(u)\in A}{\FormulaRefAuto{A\subseteq B,\,x\in A\vdash x\in B}[thm]{10,17}}
  \proofstepwidestar[1,2,15,16]{h(v)\in A}{\FormulaRefAuto{A\subseteq B,\,x\in A\vdash x\in B}[thm]{10,18}}
  \proofstepwidestar[1,2,15,16]{h(u),h(v)\in A}{\rAI{20,21}}
 \proofstepwidestar[1,2,15,16]{f(h(u)\star h(v))=f(h(u))\diamond f(h(v))}{\FormulaRefAuto{SemigroupIsoOperation}[thm]{1,22}}
 \proofstepwidestar[1,2,15,16]{u\in B}{\FormulaRefAuto{A\subseteq B,\,x\in A\vdash x\in B}[thm]{13,15}}
  \proofstepwidestar[1,2,15,16]{v\in B}{\FormulaRefAuto{A\subseteq B,\,x\in A\vdash x\in B}[thm]{13,16}}
  \proofstepwidestar[1,2,15,16]{u,v\in B}{\rAI{24,25}}
 \proofstepwidestar[1,2,15,16]{f(h(u))=u}{\FormulaRefAuto{InverseFunctionRightInverse}[thm]{5,\rAEa{26}}}
 \proofstepwidestar[1,2,15,16]{f(h(v))=v}{\FormulaRefAuto{InverseFunctionRightInverse}[thm]{5,\rAEb{26}}}
 \proofstepwidestar[1,2,15,16]{f(h(u)\star h(v))\in f[U]}{\FormulaRefAuto{C\subseteq A\dsep F\colon A\to B\dsep x\in C\vdash F(x)\in F[C]}[thm]{10,4,19}}
 \proofstepwidestar[1,2,15,16]{u\diamond v\in f[U]}{\rIE{28,\rIE{27,\rIE{23,29}}}}
 \proofstepwidestar[1,2]{\forall u,v\in f[U]\;(u\diamond v\in f[U])}{\rUI{\rRI{15,\rUI{\rRI{16,30}}}}}
 \proofstepwidestar[1,2]{f[U]\subseteq B\land f[U]\neq\varnothing\land \forall u,v\in f[U]\;(u\diamond v\in f[U])}{\rAI{13,\rAI{14,31}}}
 \proofstepwidestar[1]{\mathsf{UHGr}(f[U];B,\diamond)\leftrightarrow(f[U]\subseteq B\land f[U]\neq\varnothing\land \forall u,v\in f[U]\;(u\diamond v\in f[U]))}{\FormulaRefAuto{SemigroupSubsemigroupCriterion}[thm]{9}}
 \proofstepwidestar[1,2]{\mathsf{UHGr}(f[U];B,\diamond)}{\rRE{\rLREb{33},32}}
\closeproofpartwideindR
\proofpartwideindR[Linksabsorption]{%
f:A\cong B\dsep \mathsf{LId}(U;A,\star)\vdash \mathsf{LId}(f[U];B,\diamond)%
}{f:A\cong B\dsep \mathsf{LId}(U;A,\star)\vdash \mathsf{LId}(f[U];B,\diamond)}[SemigroupIsoLIdImage]
 \proofstepwidestar[1]{f:A\cong B}{\rA}
 \proofstepwidestar[2]{\mathsf{LId}(U;A,\star)}{\rA}
 \proofstepwidestar[1]{\SemigroupHom{f}{A,\star}{B,\diamond}}{\FormulaRefAuto{SemigroupIsoHomAxiom}[ax]{1}}
 \proofstepwidestar[1]{f:A\to B}{\FormulaRefAuto{SemigroupHomFunctionAxiom}[ax]{3}}
 \proofstepwidestar[1]{f:A\bij B}{\FormulaRefAuto{SemigroupIsoBijectiveAxiom}[ax]{1}}
 \proofstepwidestar[1]{h:B\bij A}{\FormulaRefAuto{SemigroupIsoInverseCarrier}[thm]{1}}
 \proofstepwidestar[1]{h:B\to A}{\FormulaRefAuto{Bijektive Funktion}[def]{6}}
 \proofstepwidestar[1]{\SemigroupAlg{A}{\star}}{\FormulaRefAuto{SemigroupHomSourceAxiom}[ax]{3}}
 \proofstepwidestar[1]{\SemigroupAlg{B}{\diamond}}{\FormulaRefAuto{SemigroupHomTargetAxiom}[ax]{3}}
 \proofstepwidestar[1,2]{U\subseteq A}{\FormulaRefAuto{SemigroupLeftIdealSubsetAxiom}[ax]{8,2}}
 \proofstepwidestar[1,2]{U\neq\varnothing}{\FormulaRefAuto{SemigroupLeftIdealNonemptyAxiom}[ax]{8,2}}
 \proofstepwidestar[1,2]{f[U]\subseteq B\land h[f[U]]=U\land(U\neq\varnothing\leftrightarrow f[U]\neq\varnothing)}{\FormulaRefAuto{SemigroupIsoImageCarrierData}[thm]{1,10}}
 \proofstepwidestar[1,2]{f[U]\subseteq B}{\rAEa{12}}
 \proofstepwidestar[1,2]{f[U]\neq\varnothing}{\rRE{\rLREa{\rAEb{\rAEb{12}}},11}}
 \proofstepwidestar[15]{b\in B}{\rA}
 \proofstepwidestar[16]{u\in f[U]}{\rA}
 \proofstepwidestar[1,15]{h(b)\in A}{\FormulaRefAuto{F\colon A\to B\dsep x\in A\vdash F(x)\in B}[thm]{7,15}}
 \proofstepwidestar[1,2,16]{h(u)\in U}{\FormulaRefAuto{SemigroupIsoInverseImageMember}[thm]{1,10,16}}
 \proofstepwidestar[1,2,15,16]{h(b)\star h(u)\in U}{\FormulaRefAuto{SemigroupLeftIdealClosureAxiom}[ax]{8,2,17,18}}
 \proofstepwidestar[1,2,15,16]{h(b)\in A}{\rAEa{\rAI{17,17}}}
  \proofstepwidestar[1,2,15,16]{h(u)\in A}{\FormulaRefAuto{A\subseteq B,\,x\in A\vdash x\in B}[thm]{10,18}}
  \proofstepwidestar[1,2,15,16]{h(b),h(u)\in A}{\rAI{20,21}}
 \proofstepwidestar[1,2,15,16]{f(h(b)\star h(u))=f(h(b))\diamond f(h(u))}{\FormulaRefAuto{SemigroupIsoOperation}[thm]{1,22}}
 \proofstepwidestar[1,2,15,16]{u\in B}{\FormulaRefAuto{A\subseteq B,\,x\in A\vdash x\in B}[thm]{13,16}}
  \proofstepwidestar[1,2,15,16]{b,u\in B}{\rAI{15,24}}
 \proofstepwidestar[1,2,15,16]{f(h(b))=b}{\FormulaRefAuto{InverseFunctionRightInverse}[thm]{5,\rAEa{25}}}
 \proofstepwidestar[1,2,15,16]{f(h(u))=u}{\FormulaRefAuto{InverseFunctionRightInverse}[thm]{5,\rAEb{25}}}
 \proofstepwidestar[1,2,15,16]{f(h(b)\star h(u))\in f[U]}{\FormulaRefAuto{C\subseteq A\dsep F\colon A\to B\dsep x\in C\vdash F(x)\in F[C]}[thm]{10,4,19}}
 \proofstepwidestar[1,2,15,16]{b\diamond u\in f[U]}{\rIE{27,\rIE{26,\rIE{23,28}}}}
 \proofstepwidestar[1,2]{\forall b\in B\;\forall u\in f[U]\;(b\diamond u\in f[U])}{\rUI{\rRI{15,\rUI{\rRI{16,29}}}}}
 \proofstepwidestar[1,2]{f[U]\subseteq B\land f[U]\neq\varnothing\land \forall b\in B\;\forall u\in f[U]\;(b\diamond u\in f[U])}{\rAI{13,\rAI{14,30}}}
 \proofstepwidestar[1]{\mathsf{LId}(f[U];B,\diamond)\leftrightarrow(f[U]\subseteq B\land f[U]\neq\varnothing\land \forall b\in B\;\forall u\in f[U]\;(b\diamond u\in f[U]))}{\FormulaRefAuto{SemigroupLeftIdealCriterion}[thm]{9}}
 \proofstepwidestar[1,2]{\mathsf{LId}(f[U];B,\diamond)}{\rRE{\rLREb{32},31}}
\closeproofpartwideindR
\proofpartwideindR[Rechtsabsorption]{%
f:A\cong B\dsep \mathsf{RId}(U;A,\star)\vdash \mathsf{RId}(f[U];B,\diamond)%
}{f:A\cong B\dsep \mathsf{RId}(U;A,\star)\vdash \mathsf{RId}(f[U];B,\diamond)}[SemigroupIsoRIdImage]
 \proofstepwidestar[1]{f:A\cong B}{\rA}
 \proofstepwidestar[2]{\mathsf{RId}(U;A,\star)}{\rA}
 \proofstepwidestar[1]{\SemigroupHom{f}{A,\star}{B,\diamond}}{\FormulaRefAuto{SemigroupIsoHomAxiom}[ax]{1}}
 \proofstepwidestar[1]{f:A\to B}{\FormulaRefAuto{SemigroupHomFunctionAxiom}[ax]{3}}
 \proofstepwidestar[1]{f:A\bij B}{\FormulaRefAuto{SemigroupIsoBijectiveAxiom}[ax]{1}}
 \proofstepwidestar[1]{h:B\bij A}{\FormulaRefAuto{SemigroupIsoInverseCarrier}[thm]{1}}
 \proofstepwidestar[1]{h:B\to A}{\FormulaRefAuto{Bijektive Funktion}[def]{6}}
 \proofstepwidestar[1]{\SemigroupAlg{A}{\star}}{\FormulaRefAuto{SemigroupHomSourceAxiom}[ax]{3}}
 \proofstepwidestar[1]{\SemigroupAlg{B}{\diamond}}{\FormulaRefAuto{SemigroupHomTargetAxiom}[ax]{3}}
 \proofstepwidestar[1,2]{U\subseteq A}{\FormulaRefAuto{SemigroupRightIdealSubsetAxiom}[ax]{8,2}}
 \proofstepwidestar[1,2]{U\neq\varnothing}{\FormulaRefAuto{SemigroupRightIdealNonemptyAxiom}[ax]{8,2}}
 \proofstepwidestar[1,2]{f[U]\subseteq B\land h[f[U]]=U\land(U\neq\varnothing\leftrightarrow f[U]\neq\varnothing)}{\FormulaRefAuto{SemigroupIsoImageCarrierData}[thm]{1,10}}
 \proofstepwidestar[1,2]{f[U]\subseteq B}{\rAEa{12}}
 \proofstepwidestar[1,2]{f[U]\neq\varnothing}{\rRE{\rLREa{\rAEb{\rAEb{12}}},11}}
 \proofstepwidestar[15]{b\in B}{\rA}
 \proofstepwidestar[16]{u\in f[U]}{\rA}
 \proofstepwidestar[1,2,16]{h(u)\in U}{\FormulaRefAuto{SemigroupIsoInverseImageMember}[thm]{1,10,16}}
 \proofstepwidestar[1,15]{h(b)\in A}{\FormulaRefAuto{F\colon A\to B\dsep x\in A\vdash F(x)\in B}[thm]{7,15}}
 \proofstepwidestar[1,2,15,16]{h(u)\star h(b)\in U}{\FormulaRefAuto{SemigroupRightIdealClosureAxiom}[ax]{8,2,17,18}}
 \proofstepwidestar[1,2,15,16]{h(u)\in A}{\FormulaRefAuto{A\subseteq B,\,x\in A\vdash x\in B}[thm]{10,17}}
  \proofstepwidestar[1,2,15,16]{h(b)\in A}{\rAEa{\rAI{18,18}}}
  \proofstepwidestar[1,2,15,16]{h(u),h(b)\in A}{\rAI{20,21}}
 \proofstepwidestar[1,2,15,16]{f(h(u)\star h(b))=f(h(u))\diamond f(h(b))}{\FormulaRefAuto{SemigroupIsoOperation}[thm]{1,22}}
 \proofstepwidestar[1,2,15,16]{u\in B}{\FormulaRefAuto{A\subseteq B,\,x\in A\vdash x\in B}[thm]{13,16}}
  \proofstepwidestar[1,2,15,16]{u,b\in B}{\rAI{24,15}}
 \proofstepwidestar[1,2,15,16]{f(h(u))=u}{\FormulaRefAuto{InverseFunctionRightInverse}[thm]{5,\rAEa{25}}}
 \proofstepwidestar[1,2,15,16]{f(h(b))=b}{\FormulaRefAuto{InverseFunctionRightInverse}[thm]{5,\rAEb{25}}}
 \proofstepwidestar[1,2,15,16]{f(h(u)\star h(b))\in f[U]}{\FormulaRefAuto{C\subseteq A\dsep F\colon A\to B\dsep x\in C\vdash F(x)\in F[C]}[thm]{10,4,19}}
 \proofstepwidestar[1,2,15,16]{u\diamond b\in f[U]}{\rIE{27,\rIE{26,\rIE{23,28}}}}
 \proofstepwidestar[1,2]{\forall b\in B\;\forall u\in f[U]\;(u\diamond b\in f[U])}{\rUI{\rRI{15,\rUI{\rRI{16,29}}}}}
 \proofstepwidestar[1,2]{f[U]\subseteq B\land f[U]\neq\varnothing\land \forall b\in B\;\forall u\in f[U]\;(u\diamond b\in f[U])}{\rAI{13,\rAI{14,30}}}
 \proofstepwidestar[1]{\mathsf{RId}(f[U];B,\diamond)\leftrightarrow(f[U]\subseteq B\land f[U]\neq\varnothing\land \forall b\in B\;\forall u\in f[U]\;(u\diamond b\in f[U]))}{\FormulaRefAuto{SemigroupRightIdealCriterion}[thm]{9}}
 \proofstepwidestar[1,2]{\mathsf{RId}(f[U];B,\diamond)}{\rRE{\rLREb{32},31}}
\closeproofpartwideindR
\proofpartwideindR[Zweiseitige Ideale]{%
f:A\cong B\dsep\mathsf{Id}(U;A,\star)\vdash\mathsf{Id}(f[U];B,\diamond)%
}{f:A\cong B\dsep\mathsf{Id}(U;A,\star)\vdash\mathsf{Id}(f[U];B,\diamond)}[SemigroupIsoIdImage]
 \proofstepwidestar[1]{f:A\cong B}{\rA}
 \proofstepwidestar[2]{\mathsf{Id}(U;A,\star)}{\rA}
 \proofstepwidestar[1]{\SemigroupHom{f}{A,\star}{B,\diamond}}{\FormulaRefAuto{SemigroupIsoHomAxiom}[ax]{1}}
 \proofstepwidestar[1]{\SemigroupAlg{A}{\star}}{\FormulaRefAuto{SemigroupHomSourceAxiom}[ax]{3}}
 \proofstepwidestar[1]{\SemigroupAlg{B}{\diamond}}{\FormulaRefAuto{SemigroupHomTargetAxiom}[ax]{3}}
 \proofstepwidestar[1,2]{\mathsf{LId}(U;A,\star)}{\FormulaRefAuto{SemigroupIdealLeftAxiom}[ax]{4,2}}
 \proofstepwidestar[1,2]{\mathsf{RId}(U;A,\star)}{\FormulaRefAuto{SemigroupIdealRightAxiom}[ax]{4,2}}
 \proofstepwidestar[1,2]{\mathsf{LId}(f[U];B,\diamond)}{\FormulaRefAuto{SemigroupIsoLIdImage}[thm]{1,6}}
 \proofstepwidestar[1,2]{\mathsf{RId}(f[U];B,\diamond)}{\FormulaRefAuto{SemigroupIsoRIdImage}[thm]{1,7}}
 \proofstepwidestar[1,2]{\mathsf{LId}(f[U];B,\diamond)\land\mathsf{RId}(f[U];B,\diamond)}{\rAI{8,9}}
 \proofstepwidestar[1,2]{\mathsf{Id}(f[U];B,\diamond)}{\FormulaRefAuto{SemigroupIdealDef}[def]{5,10}}
\closeproofpartwideindR
\proofpartwideindR[Rückrichtung für UHGr]{%
f:A\cong B\dsep U\subseteq A\vdash \mathsf{UHGr}(U;A,\star)\leftrightarrow \mathsf{UHGr}(f[U];B,\diamond)%
}{f:A\cong B\dsep U\subseteq A\vdash \mathsf{UHGr}(U;A,\star)\leftrightarrow \mathsf{UHGr}(f[U];B,\diamond)}[SemigroupIsoUHGrTransport]
 \proofstepwidestar[1]{f:A\cong B}{\rA}
 \proofstepwidestar[2]{U\subseteq A}{\rA}
 \proofstepwidestar[1]{h:B\cong A}{\FormulaRefAuto{SemigroupIsoInverse}[thm]{1}}
 \proofstepwidestar[1,2]{f[U]\subseteq B\land h[f[U]]=U\land(U\neq\varnothing\leftrightarrow f[U]\neq\varnothing)}{\FormulaRefAuto{SemigroupIsoImageCarrierData}[thm]{1,2}}
 \proofstepwidestar[1,2]{h[f[U]]=U}{\rAEa{\rAEb{4}}}
 \proofstepwidestar[6]{\mathsf{UHGr}(U;A,\star)}{\rA}
 \proofstepwidestar[1,6]{\mathsf{UHGr}(f[U];B,\diamond)}{\FormulaRefAuto{SemigroupIsoUHGrImage}[thm]{1,6}}
 \proofstepwidestar[1]{\mathsf{UHGr}(U;A,\star)\Rightarrow \mathsf{UHGr}(f[U];B,\diamond)}{\rRI{6,7}}
 \proofstepwidestar[9]{\mathsf{UHGr}(f[U];B,\diamond)}{\rA}
 \proofstepwidestar[1,9]{\mathsf{UHGr}(h[f[U]];A,\star)}{\FormulaRefAuto{SemigroupIsoUHGrImage}[thm]{3,9}}
 \proofstepwidestar[1,2,9]{\mathsf{UHGr}(U;A,\star)}{\rIE{5,10}}
 \proofstepwidestar[1,2]{\mathsf{UHGr}(f[U];B,\diamond)\Rightarrow \mathsf{UHGr}(U;A,\star)}{\rRI{9,11}}
 \proofstepwidestar[1,2]{\mathsf{UHGr}(U;A,\star)\leftrightarrow \mathsf{UHGr}(f[U];B,\diamond)}{\rLRI{8,12}}
\closeproofpartwideindR
\proofpartwideindR[Rückrichtung für LId]{%
f:A\cong B\dsep U\subseteq A\vdash \mathsf{LId}(U;A,\star)\leftrightarrow \mathsf{LId}(f[U];B,\diamond)%
}{f:A\cong B\dsep U\subseteq A\vdash \mathsf{LId}(U;A,\star)\leftrightarrow \mathsf{LId}(f[U];B,\diamond)}[SemigroupIsoLIdTransport]
 \proofstepwidestar[1]{f:A\cong B}{\rA}
 \proofstepwidestar[2]{U\subseteq A}{\rA}
 \proofstepwidestar[1]{h:B\cong A}{\FormulaRefAuto{SemigroupIsoInverse}[thm]{1}}
 \proofstepwidestar[1,2]{f[U]\subseteq B\land h[f[U]]=U\land(U\neq\varnothing\leftrightarrow f[U]\neq\varnothing)}{\FormulaRefAuto{SemigroupIsoImageCarrierData}[thm]{1,2}}
 \proofstepwidestar[1,2]{h[f[U]]=U}{\rAEa{\rAEb{4}}}
 \proofstepwidestar[6]{\mathsf{LId}(U;A,\star)}{\rA}
 \proofstepwidestar[1,6]{\mathsf{LId}(f[U];B,\diamond)}{\FormulaRefAuto{SemigroupIsoLIdImage}[thm]{1,6}}
 \proofstepwidestar[1]{\mathsf{LId}(U;A,\star)\Rightarrow \mathsf{LId}(f[U];B,\diamond)}{\rRI{6,7}}
 \proofstepwidestar[9]{\mathsf{LId}(f[U];B,\diamond)}{\rA}
 \proofstepwidestar[1,9]{\mathsf{LId}(h[f[U]];A,\star)}{\FormulaRefAuto{SemigroupIsoLIdImage}[thm]{3,9}}
 \proofstepwidestar[1,2,9]{\mathsf{LId}(U;A,\star)}{\rIE{5,10}}
 \proofstepwidestar[1,2]{\mathsf{LId}(f[U];B,\diamond)\Rightarrow \mathsf{LId}(U;A,\star)}{\rRI{9,11}}
 \proofstepwidestar[1,2]{\mathsf{LId}(U;A,\star)\leftrightarrow \mathsf{LId}(f[U];B,\diamond)}{\rLRI{8,12}}
\closeproofpartwideindR
\proofpartwideindR[Rückrichtung für RId]{%
f:A\cong B\dsep U\subseteq A\vdash \mathsf{RId}(U;A,\star)\leftrightarrow \mathsf{RId}(f[U];B,\diamond)%
}{f:A\cong B\dsep U\subseteq A\vdash \mathsf{RId}(U;A,\star)\leftrightarrow \mathsf{RId}(f[U];B,\diamond)}[SemigroupIsoRIdTransport]
 \proofstepwidestar[1]{f:A\cong B}{\rA}
 \proofstepwidestar[2]{U\subseteq A}{\rA}
 \proofstepwidestar[1]{h:B\cong A}{\FormulaRefAuto{SemigroupIsoInverse}[thm]{1}}
 \proofstepwidestar[1,2]{f[U]\subseteq B\land h[f[U]]=U\land(U\neq\varnothing\leftrightarrow f[U]\neq\varnothing)}{\FormulaRefAuto{SemigroupIsoImageCarrierData}[thm]{1,2}}
 \proofstepwidestar[1,2]{h[f[U]]=U}{\rAEa{\rAEb{4}}}
 \proofstepwidestar[6]{\mathsf{RId}(U;A,\star)}{\rA}
 \proofstepwidestar[1,6]{\mathsf{RId}(f[U];B,\diamond)}{\FormulaRefAuto{SemigroupIsoRIdImage}[thm]{1,6}}
 \proofstepwidestar[1]{\mathsf{RId}(U;A,\star)\Rightarrow \mathsf{RId}(f[U];B,\diamond)}{\rRI{6,7}}
 \proofstepwidestar[9]{\mathsf{RId}(f[U];B,\diamond)}{\rA}
 \proofstepwidestar[1,9]{\mathsf{RId}(h[f[U]];A,\star)}{\FormulaRefAuto{SemigroupIsoRIdImage}[thm]{3,9}}
 \proofstepwidestar[1,2,9]{\mathsf{RId}(U;A,\star)}{\rIE{5,10}}
 \proofstepwidestar[1,2]{\mathsf{RId}(f[U];B,\diamond)\Rightarrow \mathsf{RId}(U;A,\star)}{\rRI{9,11}}
 \proofstepwidestar[1,2]{\mathsf{RId}(U;A,\star)\leftrightarrow \mathsf{RId}(f[U];B,\diamond)}{\rLRI{8,12}}
\closeproofpartwideindR
\proofpartwideindR[Rückrichtung für Id]{%
f:A\cong B\dsep U\subseteq A\vdash \mathsf{Id}(U;A,\star)\leftrightarrow \mathsf{Id}(f[U];B,\diamond)%
}{f:A\cong B\dsep U\subseteq A\vdash \mathsf{Id}(U;A,\star)\leftrightarrow \mathsf{Id}(f[U];B,\diamond)}[SemigroupIsoIdTransport]
 \proofstepwidestar[1]{f:A\cong B}{\rA}
 \proofstepwidestar[2]{U\subseteq A}{\rA}
 \proofstepwidestar[1]{h:B\cong A}{\FormulaRefAuto{SemigroupIsoInverse}[thm]{1}}
 \proofstepwidestar[1,2]{f[U]\subseteq B\land h[f[U]]=U\land(U\neq\varnothing\leftrightarrow f[U]\neq\varnothing)}{\FormulaRefAuto{SemigroupIsoImageCarrierData}[thm]{1,2}}
 \proofstepwidestar[1,2]{h[f[U]]=U}{\rAEa{\rAEb{4}}}
 \proofstepwidestar[6]{\mathsf{Id}(U;A,\star)}{\rA}
 \proofstepwidestar[1,6]{\mathsf{Id}(f[U];B,\diamond)}{\FormulaRefAuto{SemigroupIsoIdImage}[thm]{1,6}}
 \proofstepwidestar[1]{\mathsf{Id}(U;A,\star)\Rightarrow \mathsf{Id}(f[U];B,\diamond)}{\rRI{6,7}}
 \proofstepwidestar[9]{\mathsf{Id}(f[U];B,\diamond)}{\rA}
 \proofstepwidestar[1,9]{\mathsf{Id}(h[f[U]];A,\star)}{\FormulaRefAuto{SemigroupIsoIdImage}[thm]{3,9}}
 \proofstepwidestar[1,2,9]{\mathsf{Id}(U;A,\star)}{\rIE{5,10}}
 \proofstepwidestar[1,2]{\mathsf{Id}(f[U];B,\diamond)\Rightarrow \mathsf{Id}(U;A,\star)}{\rRI{9,11}}
 \proofstepwidestar[1,2]{\mathsf{Id}(U;A,\star)\leftrightarrow \mathsf{Id}(f[U];B,\diamond)}{\rLRI{8,12}}
\closeproofpartwideindR
\end{tabproofsplitwide}

\FormulaThmDeltaK[Einschränkung auf einen abgeschlossenen Teilträger]{%
 f:A\cong B\dsep U\subseteq A\dsep\SemigroupAlg{U}{\star}\vdash
 f|_U:U\cong f[U]%
}{SemigroupIsoClosedCarrierRestriction}{%
 \DeltaRow{Mengen}{A\dsep B\dsep U\dsep x\dsep y\dsep z}
 \DeltaRow{Funktionen}{f\dsep\star\dsep\diamond}
}
\begin{tabproofsplitwide}
\proofpartwideindR[Halbgruppenstruktur auf dem Bild]{%
 f:A\cong B\dsep U\subseteq A\dsep\SemigroupAlg{U}{\star}\vdash
 \SemigroupAlg{f[U]}{\diamond}%
}{f:A\cong B\dsep U\subseteq A\dsep\SemigroupAlg{U}{\star}\vdash\SemigroupAlg{f[U]}{\diamond}}[SemigroupIsoClosedCarrierImage]
 \proofstepwidestar[1]{f:A\cong B}{\rA}
 \proofstepwidestar[2]{U\subseteq A}{\rA}
 \proofstepwidestar[3]{\SemigroupAlg{U}{\star}}{\rA}
 \proofstepwidestar[1]{\SemigroupAlg{B}{\diamond}}{\FormulaRefAuto{SemigroupIsoTarget}[thm]{1}}
 \proofstepwidestar[1]{f:A\to B}{\FormulaRefAuto{SemigroupIsoFunction}[thm]{1}}
 \proofstepwidestar[1,2]{f[U]\subseteq B}{\FormulaRefAuto{C\subseteq A\dsep F\colon A\to B\vdash F[C]\subseteq B}[thm]{2,5}}
 \proofstepwidestar[7]{x\in f[U]}{\rA}
 \proofstepwidestar[8]{y\in f[U]}{\rA}
 \proofstepwidestar[7,8]{x,y\in f[U]}{\rAIStar{7,8}}
 \proofstepwidestar[1,2,7,8]{h(x),h(y)\in U}{\FormulaRefAuto{SemigroupIsoInverseImageMember}[thm]{1,2,9}}
 \proofstepwidestar[1,2,3,7,8]{h(x)\star h(y)\in U}{\FormulaRefAuto{SemigroupClosureAxiom}[ax]{3,10}}
 \proofstepwidestar[1,2,7,8]{x,y\in B}{\FormulaRefAuto{A\subseteq B,\,x\in A\vdash x\in B}[thm]{6,9}}
 \proofstepwidestar[1,2,7,8]{h(x\diamond y)=h(x)\star h(y)}{\FormulaRefAuto{SemigroupIsoInverseOperation}[thm]{1,12}}
 \proofstepwidestar[1,2,7,8]{h(x)\star h(y)=h(x\diamond y)}{\FormulaRefAuto{a=b\vdash b=a}[thm]{13}}
  \proofstepwidestar[1,2,3,7,8]{h(x\diamond y)\in U}{\rIE{14,11}}
 \proofstepwidestar[1]{h:B\cong A}{\FormulaRefAuto{SemigroupIsoInverse}[thm]{1}}
 \proofstepwidestar[1,2,7,8]{x\diamond y\in B}{\FormulaRefAuto{SemigroupClosureAxiom}[ax]{4,12}}
 \proofstepwidestar[1,2]{h[f[U]]=U}{\rAEa{\rAEb{\FormulaRefAuto{SemigroupIsoImageCarrierData}[thm]{1,2}}}}
 \proofstepwidestar[1,2,7,8]{x\diamond y\in f[U]\leftrightarrow h(x\diamond y)\in h[f[U]]}{\FormulaRefAuto{SemigroupIsoImageMembership}[thm]{16,6,17}}
  \proofstepwidestar[1,2,7,8]{x\diamond y\in f[U]\leftrightarrow h(x\diamond y)\in U}{\rIE{18,19}}
 \proofstepwidestar[1,2,3,7,8]{x\diamond y\in f[U]}{\rRE{\rLREb{20},15}}
 \proofstepwidestar[1,2,3]{\forall x,y\in f[U]\;(x\diamond y\in f[U])}{\rUI{\rRI{7,\rUI{\rRI{8,21}}}}}
 \proofstepwidestar[23]{x\in f[U]}{\rA}
 \proofstepwidestar[24]{y\in f[U]}{\rA}
 \proofstepwidestar[25]{z\in f[U]}{\rA}
 \proofstepwidestar[23,24,25]{x,y,z\in f[U]}{\rAIStar{23,24,25}}
 \proofstepwidestar[1,2,23,24,25]{x,y,z\in B}{\FormulaRefAuto{A\subseteq B,\,x\in A\vdash x\in B}[thm]{6,26}}
 \proofstepwidestar[1,2,23,24,25]{(x\diamond y)\diamond z=x\diamond(y\diamond z)}{\FormulaRefAuto{SemigroupAssociativityAxiom}[ax]{4,27}}
 \proofstepwidestar[1,2]{\forall x,y,z\in f[U]\;((x\diamond y)\diamond z=x\diamond(y\diamond z))}{\rUI{\rRI{23,\rUI{\rRI{24,\rUI{\rRI{25,28}}}}}}}
 \proofstepwidestar[1,2,3]{\SemigroupAlg{f[U]}{\diamond}}{\FormulaRefAuto{SemigroupDef}[def]{22,29}}
\closeproofpartwideindR
\proofpartwide{Einschränkungsisomorphismus}
 \proofstepwidestar[1]{f:A\cong B}{\rA}
 \proofstepwidestar[2]{U\subseteq A}{\rA}
 \proofstepwidestar[3]{\SemigroupAlg{U}{\star}}{\rA}
 \proofstepwidestar[1]{f:A\to B}{\FormulaRefAuto{SemigroupIsoFunction}[thm]{1}}
 \proofstepwidestar[1,2]{f[U]\subseteq B}{\FormulaRefAuto{C\subseteq A\dsep F\colon A\to B\vdash F[C]\subseteq B}[thm]{2,4}}
 \proofstepwidestar[1,2,3]{\SemigroupAlg{f[U]}{\diamond}}{\FormulaRefAuto{SemigroupIsoClosedCarrierImage}[thm]{1,2,3}}
 \proofstepwidestar[]{f[U]=f[U]}{\rII}
 \proofstepwidestar[1,2,3]{f|_U:U\cong f[U]}{\FormulaRefAuto{SemigroupIsoRestrictionToImage}[thm]{1,3,6,2,5,7}}
\closeproofpartwide
\end{tabproofsplitwide}

\FormulaThmDeltaK[Abgeschlossene Teilträger erben die Halbgruppenstruktur]{%
 \SemigroupAlg{A}{\star}\dsep U\subseteq A\dsep
 \forall x,y\in U\;(x\star y\in U)\vdash\SemigroupAlg{U}{\star}%
}{SemigroupClosedSubsetIsSemigroup}{%
 \DeltaRow{Mengen}{A\dsep U\dsep x\dsep y\dsep z}
 \DeltaRow{Funktionen}{\star}
}
\begin{tabproofwide}
 \proofstepwidestar[1]{\SemigroupAlg{A}{\star}}{\rA}
 \proofstepwidestar[2]{U\subseteq A}{\rA}
 \proofstepwidestar[3]{\forall x,y\in U\;(x\star y\in U)}{\rA}
 \proofstepwidestar[4]{x\in U}{\rA}
 \proofstepwidestar[5]{y\in U}{\rA}
 \proofstepwidestar[6]{z\in U}{\rA}
 \proofstepwidestar[4,5,6]{x,y,z\in U}{\rAIStar{4,5,6}}
 \proofstepwidestar[2,4,5,6]{x,y,z\in A}{\FormulaRefAuto{A\subseteq B,\,x\in A\vdash x\in B}[thm]{2,7}}
 \proofstepwidestar[1,2,4,5,6]{(x\star y)\star z=x\star(y\star z)}{\FormulaRefAuto{SemigroupAssociativityAxiom}[ax]{1,8}}
 \proofstepwidestar[1,2]{\forall x,y,z\in U\;((x\star y)\star z=x\star(y\star z))}{\rUI{\rRI{4,\rUI{\rRI{5,\rUI{\rRI{6,9}}}}}}}
 \proofstepwidestar[1,2,3]{\SemigroupAlg{U}{\star}}{\FormulaRefAuto{SemigroupDef}[def]{3,10}}
\end{tabproofwide}

\subsection{Produktpotenzen und Erzeugnisse}

\FormulaDefDeltaK[Positive Produktpotenzen eines Halbgruppenträgers]{%
 \begin{aligned}[t]
 &\SemigroupAlg{A}{\star}\dsep A\neq\varnothing\vdash{}\\[-2pt]
 &A^1\coloneqq A,\qquad A^{n+1}\coloneqq A^n\star_{\mathcal P}A
       \quad(n\in\mathbb N,\ n\neq0).
 \end{aligned}%
}{SemigroupPositiveCarrierPowersDef}{%
 \DeltaRow{Mengen}{A\dsep n}
 \DeltaRow{Funktionen}{\star\dsep\star_{\mathcal P}}
}
Hier bezeichnet \(A^n\) die Menge der Produkte aus \(n\) Faktoren.
Das bereits definierte \(A^2\) bleibt unverändert. Ein kartesisches
Potenzobjekt wird unten durch \(A^I\) mit ausdrücklich genannter
Indexmenge \(I\) bezeichnet.

\FormulaThmDeltaKR[Abschluss und Rechenregeln der Produktpotenzen]{%
 \begin{aligned}[t]
  &\SemigroupAlg{A}{\star}\dsep A\neq\varnothing\dsep m,n\geq1\vdash{}\\[-2pt]
  &\text{(i)}\quad A^m A^n=A^{m+n}\\[-2pt]
  &\text{(ii)}\quad A^{n+1}\subseteq A^n\\[-2pt]
  &\text{(iii)}\quad \mathsf{UHGr}(A^n;A,\star)
 \end{aligned}%
}{%
 \begin{aligned}[t]
 &\SemigroupAlg{A}{\star}\dsep A\neq\varnothing\dsep m,n\geq1\vdash{}\\[-2pt]
 &A^m A^n=A^{m+n},\qquad A^{n+1}\subseteq A^n,\qquad
      \mathsf{UHGr}(A^n;A,\star).
 \end{aligned}%
}{SemigroupCarrierPowersCalculus}{%
 \DeltaRow{Mengen}{A\dsep m,n\in\mathbb N\dsep x\dsep a}
 \DeltaRow{Funktionen}{\star\dsep\star_{\mathcal P}}
}
\begin{tabproofsplitwide}
\proofpartwideindR[Nichtleere Teilträger]{%
\SemigroupAlg{A}{\star}\dsep A\neq\varnothing\dsep n\geq1\vdash\varnothing\neq A^n\subseteq A%
}{\SemigroupAlg{A}{\star}\dsep A\neq\varnothing\dsep n\geq1\vdash\varnothing\neq A^n\subseteq A}[SemigroupCarrierPowersCarriers]
 \proofstepwidestar[1]{\SemigroupAlg{A}{\star}}{\rA}
 \proofstepwidestar[2]{A\neq\varnothing}{\rA}
 \proofstepwidestar[1,2]{\varnothing\neq A^1\subseteq A}{\FormulaRefAuto{SemigroupPositiveCarrierPowersDef}[def]{1,2}}
 \proofstepwidestar[4]{n\geq1}{\rA}
 \proofstepwidestar[5]{\varnothing\neq A^n\subseteq A}{\rA}
 \proofstepwidestar[4,5]{n\geq1\,,\quad\varnothing\neq A^n\subseteq A}{\rAIStar{4,5}}
 \proofstepwidestar[1,2,4,5]{A^n,A\in\PowNE A}{\text{\parbox[t]{\linewidth}{\raggedright Nichtleere Teilmengen aus 6,2; Definition von \(\PowNE A\).}}}
 \proofstepwidestar[1,2,4,5]{A^n A\in\PowNE A}{\FormulaRefAuto{PowerSemigroupProductClosure}[thm]{1,7}}
 \proofstepwidestar[1,2,4,5]{\varnothing\neq A^{n+1}\subseteq A}{\FormulaRefAuto{SemigroupPositiveCarrierPowersDef}[def]{8}}
 \proofstepwidestar[1,2]{\forall n\geq1\;(\varnothing\neq A^n\subseteq A\Rightarrow\varnothing\neq A^{n+1}\subseteq A)}{\rUI{\rRI{4,\rRI{5,9}}}}
 \proofstepwidestar[1,2]{\forall n\geq1\;(\varnothing\neq A^n\subseteq A)}{\FormulaRefAuto{PeanoInductionFromOne}[thm]{3,10}}
\closeproofpartwideindR
\proofpartwideindR[Addition der Produktlängen]{%
\SemigroupAlg{A}{\star}\dsep A\neq\varnothing\dsep m,n\geq1\vdash A^mA^n=A^{m+n}%
}{\SemigroupAlg{A}{\star}\dsep A\neq\varnothing\dsep m,n\geq1\vdash A^mA^n=A^{m+n}}[SemigroupCarrierPowersAddition]
 \proofstepwidestar[1]{\SemigroupAlg{A}{\star}}{\rA}
 \proofstepwidestar[2]{A\neq\varnothing}{\rA}
 \proofstepwidestar[3]{m\geq1}{\rA}
 \proofstepwidestar[1,2,3]{A^m A^1=A^m A=A^{m+1}}{\FormulaRefAuto{SemigroupPositiveCarrierPowersDef}[def]{3}}
 \proofstepwidestar[5]{n\geq1}{\rA}
 \proofstepwidestar[6]{A^m A^n=A^{m+n}}{\rA}
 \proofstepwidestar[5,6]{n\geq1\,,\quad A^m A^n=A^{m+n}}{\rAIStar{5,6}}
 \proofstepwidestar[1,2,3]{\varnothing\neq A^m\subseteq A}{\FormulaRefAuto{SemigroupCarrierPowersCarriers}[thm]{1,2,3}}
 \proofstepwidestar[1,2,5,6]{\varnothing\neq A^n\subseteq A}{\FormulaRefAuto{SemigroupCarrierPowersCarriers}[thm]{1,2,7}}
 \proofstepwidestar[1,2,3,5,6]{A^m,A^n,A\in\PowNE A}{\text{\parbox[t]{\linewidth}{\raggedright Definition von \(\PowNE A\) aus 8,9,2}}}
 \proofstepwidestar[1,2,3,5,6]{A^m(A^n A)=(A^m A^n)A}{\FormulaRefAuto{PowerSemigroupAssociativity}[thm]{1,10}}
 \proofstepwidestar[1,2,3,5,6]{A^m A^{n+1}=A^{m+n+1}}{\FormulaRefAuto{SemigroupPositiveCarrierPowersDef}[def]{11,7}}
 \proofstepwidestar[1,2,3]{\forall n\geq1\;(A^m A^n=A^{m+n}\Rightarrow A^m A^{n+1}=A^{m+n+1})}{\rUI{\rRI{5,\rRI{6,12}}}}
 \proofstepwidestar[1,2,3]{\forall n\geq1\;(A^m A^n=A^{m+n})}{\FormulaRefAuto{PeanoInductionFromOne}[thm]{4,13}}
\closeproofpartwideindR
\proofpartwideindR[Absteigende Folge und Unterhalbgruppen]{%
\SemigroupAlg{A}{\star}\dsep A\neq\varnothing\dsep n\geq1\vdash\mathsf{UHGr}(A^n;A,\star)%
}{\SemigroupAlg{A}{\star}\dsep A\neq\varnothing\dsep n\geq1\vdash\mathsf{UHGr}(A^n;A,\star)}[SemigroupCarrierPowersSubsemigroups]
 \proofstepwidestar[1]{\SemigroupAlg{A}{\star}}{\rA}
 \proofstepwidestar[2]{A\neq\varnothing}{\rA}
 \proofstepwidestar[1,2]{A^1=A}{\FormulaRefAuto{SemigroupPositiveCarrierPowersDef}[def]{1,2}}
 \proofstepwidestar[1,2]{A^2\subseteq A}{\FormulaRefAuto{SemigroupSquareSubset}[thm]{1,2}}
 \proofstepwidestar[1,2]{A=A^1}{\FormulaRefAuto{a=b\vdash b=a}[thm]{3}}
  \proofstepwidestar[1,2]{A^2\subseteq A^1}{\rIE{5,4}}
 \proofstepwidestar[7]{n\geq1}{\rA}
 \proofstepwidestar[8]{A^{n+1}\subseteq A^n}{\rA}
 \proofstepwidestar[7,8]{n\geq1\,,\quad A^{n+1}\subseteq A^n}{\rAIStar{7,8}}
 \proofstepwidestar[1,2,7]{A^{n+2}=A^{n+1}A}{\FormulaRefAuto{SemigroupPositiveCarrierPowersDef}[def]{1,2,7}}
  \proofstepwidestar[1,2,7]{A^{n+1}A=A^{n+2}}{\FormulaRefAuto{a=b\vdash b=a}[thm]{10}}
  \proofstepwidestar[1,2,7,8]{A^{n+1}A\subseteq A^n A}{\text{\parbox[t]{\linewidth}{\raggedright Derselbe Produktzeuge liegt wegen 9 im rechten Produkt.}}}
  \proofstepwidestar[1,2,7]{A^{n+1}=A^n A}{\FormulaRefAuto{SemigroupPositiveCarrierPowersDef}[def]{1,2,7}}
  \proofstepwidestar[1,2,7]{A^n A=A^{n+1}}{\FormulaRefAuto{a=b\vdash b=a}[thm]{13}}
  \proofstepwidestar[1,2,7,8]{A^{n+2}\subseteq A^{n+1}}{\rIE{11,14,12}}
 \proofstepwidestar[1,2]{\forall n\geq1\;(A^{n+1}\subseteq A^n\Rightarrow A^{n+2}\subseteq A^{n+1})}{\rUI{\rRI{7,\rRI{8,15}}}}
 \proofstepwidestar[1,2]{\forall n\geq1\;(A^{n+1}\subseteq A^n)}{\FormulaRefAuto{PeanoInductionFromOne}[thm]{6,16}}
 \proofstepwidestar[18]{n\geq1}{\rA}
 \proofstepwidestar[1,2,18]{A^{n+1}\subseteq A^n}{\rRE{\rUE{17},18}}
 \proofstepwidestar[20]{k\geq1}{\rA}
 \proofstepwidestar[21]{A^{n+k}\subseteq A^n}{\rA}
 \proofstepwidestar[20,21]{k\geq1\,,\quad A^{n+k}\subseteq A^n}{\rAIStar{20,21}}
 \proofstepwidestar[1,2,18,20]{A^{n+k+1}\subseteq A^{n+k}}{\text{\parbox[t]{\linewidth}{\raggedright 17 bei \(n+k\)}}}
  \proofstepwidestar[1,2,18,20,21]{A^{n+k+1}\subseteq A^n}{\FormulaRefAuto{A\subseteq B\dsep B\subseteq C\vdash A\subseteq C}[thm]{23,21}}
 \proofstepwidestar[1,2,18]{\forall k\geq1\;(A^{n+k}\subseteq A^n\Rightarrow A^{n+k+1}\subseteq A^n)}{\rUI{\rRI{20,\rRI{21,24}}}}
 \proofstepwidestar[1,2,18]{\forall k\geq1\;(A^{n+k}\subseteq A^n)}{\FormulaRefAuto{PeanoInductionFromOne}[thm]{19,25}}
 \proofstepwidestar[1,2,18]{A^{2n}\subseteq A^n}{\text{\parbox[t]{\linewidth}{\raggedright 26 bei \(k=n\)}}}
 \proofstepwidestar[1,2,18]{A^n A^n=A^{2n}}{\FormulaRefAuto{SemigroupCarrierPowersAddition}[thm]{1,2,18}}
 \proofstepwidestar[1,2,18]{\varnothing\neq A^n\land A^n\subseteq A}{\FormulaRefAuto{SemigroupCarrierPowersCarriers}[thm]{1,2,18}}
  \proofstepwidestar[1,2,18]{A^n\subseteq A}{\rAEb{29}}
  \proofstepwidestar[1,2,18]{\varnothing\neq A^n}{\rAEa{29}}
  \proofstepwidestar[1,2,18]{A^n\neq\varnothing}{\FormulaRefAuto{a \neq b \vdash b \neq a}[thm]{31}}
  \proofstepwidestar[1,2,18]{A^n\subseteq A\land A^n\neq\varnothing}{\rAI{30,32}}
  \proofstepwidestar[1,2,18]{A^n\in\PowNE A}{\FormulaRefAuto{NonemptyPowersetMembership}[thm]{33}}
 \proofstepwidestar[35]{x\in A^n}{\rA}
  \proofstepwidestar[36]{y\in A^n}{\rA}
  \proofstepwidestar[1,2,18,35,36]{xy\in A^n A^n}{\FormulaRefAuto{PowerSemigroupProductContains}[thm]{1,34,34,35,36}}
  \proofstepwidestar[1,2,18,35,36]{xy\in A^{2n}}{\rIE{28,37}}
  \proofstepwidestar[1,2,18,35,36]{xy\in A^n}{\FormulaRefAuto{A\subseteq B,\,x\in A\vdash x\in B}[thm]{27,38}}
  \proofstepwidestar[1,2,18]{\forall x,y\in A^n\;(xy\in A^n)}{\rUI{\rRI{35,\rUI{\rRI{36,39}}}}}
 \proofstepwidestar[1,2,18]{A^n\subseteq A\land A^n\neq\varnothing\land\forall x,y\in A^n\;(xy\in A^n)}{\rAI{30,\rAI{32,40}}}
 \proofstepwidestar[1,2,18]{\mathsf{UHGr}(A^n;A,\star)}{\FormulaRefAuto{SemigroupSubsemigroupCriterion}[thm]{1,41}}
\closeproofpartwideindR
\end{tabproofsplitwide}

\FormulaThmDeltaKR[Isomorphe Einschränkungen auf alle Produktpotenzen]{%
 \begin{aligned}[t]
  &f:A\cong B\dsep A\neq\varnothing\dsep n\in\mathbb N\dsep n\neq0 \vdash{}\\[-2pt]
  &\text{(i)}\quad f[A^n]=B^n\\[-2pt]
  &\text{(ii)}\quad f|_{A^n}:A^n\cong B^n
 \end{aligned}%
}{%
 \begin{aligned}[t]
 &f:A\cong B\dsep A\neq\varnothing\dsep n\in\mathbb N\dsep n\neq0
       \vdash{}\\[-2pt]
 &f[A^n]=B^n\,,\qquad f|_{A^n}:A^n\cong B^n.
 \end{aligned}%
}{SemigroupIsoPositiveCarrierPowers}{%
 \DeltaRow{Mengen}{A\dsep B\dsep n\dsep x\dsep y}
 \DeltaRow{Funktionen}{f\dsep\star\dsep\diamond}
}
\begin{tabproofsplitwide}
\proofpartwideindR[Abschluss der Produktpotenzen]{%
f:A\cong B\dsep A\neq\varnothing\dsep n\geq1\vdash\mathsf{UHGr}(A^n;A,\star)\land\mathsf{UHGr}(B^n;B,\diamond)%
}{f:A\cong B\dsep A\neq\varnothing\dsep n\geq1\vdash\mathsf{UHGr}(A^n;A,\star)\land\mathsf{UHGr}(B^n;B,\diamond)}[SemigroupIsoPositivePowersCarriers]
 \proofstepwidestar[1]{f:A\cong B}{\rA}
 \proofstepwidestar[1]{\SemigroupAlg{A}{\star}}{\FormulaRefAuto{SemigroupIsoSource}[thm]{1}}
 \proofstepwidestar[1]{\SemigroupAlg{B}{\diamond}}{\FormulaRefAuto{SemigroupIsoTarget}[thm]{1}}
 \proofstepwidestar[1]{f:A\to B}{\FormulaRefAuto{SemigroupIsoFunction}[thm]{1}}
 \proofstepwidestar[1]{f:A\bij B}{\FormulaRefAuto{SemigroupIsoBijectiveAxiom}[ax]{1}}
 \proofstepwidestar[1]{\SemigroupHom{f}{A,\star}{B,\diamond}}{\FormulaRefAuto{SemigroupIsoHomAxiom}[ax]{1}}
 \proofstepwidestar[1]{f[A]=B}{\FormulaRefAuto{SemigroupIsoFullImage}[thm]{1}}
 \proofstepwidestar[8]{A\neq\varnothing}{\rA}
 \proofstepwidestar[9]{n\geq1}{\rA}
 \proofstepwidestar[1,8]{B\neq\varnothing}{\text{\parbox[t]{\linewidth}{\raggedright 7; ein Element von \(A\) liefert einen Bildzeugen.}}}
 \proofstepwidestar[1,8,9]{\mathsf{UHGr}(A^n;A,\star)}{\FormulaRefAuto{SemigroupCarrierPowersSubsemigroups}[thm]{2,8,9}}
 \proofstepwidestar[1,8,9]{\mathsf{UHGr}(B^n;B,\diamond)}{\FormulaRefAuto{SemigroupCarrierPowersSubsemigroups}[thm]{3,10,9}}
 \proofstepwidestar[1,8,9]{\mathsf{UHGr}(A^n;A,\star)\land\mathsf{UHGr}(B^n;B,\diamond)}{\rAI{11,12}}
\closeproofpartwideindR
\proofpartwideindR[Induktion für die Bildgleichheit]{%
f:A\cong B\dsep A\neq\varnothing\dsep n\geq1\vdash f[A^n]=B^n%
}{f:A\cong B\dsep A\neq\varnothing\dsep n\geq1\vdash f[A^n]=B^n}[SemigroupIsoPositivePowersImage]
 \proofstepwidestar[1]{f:A\cong B}{\rA}
 \proofstepwidestar[1]{\SemigroupAlg{A}{\star}}{\FormulaRefAuto{SemigroupIsoSource}[thm]{1}}
 \proofstepwidestar[1]{\SemigroupAlg{B}{\diamond}}{\FormulaRefAuto{SemigroupIsoTarget}[thm]{1}}
 \proofstepwidestar[1]{f:A\to B}{\FormulaRefAuto{SemigroupIsoFunction}[thm]{1}}
 \proofstepwidestar[1]{f:A\bij B}{\FormulaRefAuto{SemigroupIsoBijectiveAxiom}[ax]{1}}
 \proofstepwidestar[1]{\SemigroupHom{f}{A,\star}{B,\diamond}}{\FormulaRefAuto{SemigroupIsoHomAxiom}[ax]{1}}
 \proofstepwidestar[1]{f[A]=B}{\FormulaRefAuto{SemigroupIsoFullImage}[thm]{1}}
 \proofstepwidestar[8]{A\neq\varnothing}{\rA}
 \proofstepwidestar[1,8]{f[A^1]=B^1}{\FormulaRefAuto{SemigroupPositiveCarrierPowersDef}[def]{2,3,7}}
 \proofstepwidestar[10]{n\geq1}{\rA}
 \proofstepwidestar[11]{ f[A^n]=B^n}{\rA}
 \proofstepwidestar[10,11]{n\geq1\land f[A^n]=B^n}{\rAIStar{10,11}}
 \proofstepwidestar[1,8,10,11]{\varnothing\neq A^n\subseteq A}{\FormulaRefAuto{SemigroupCarrierPowersCarriers}[thm]{2,8,12}}
 \proofstepwidestar[1,8,10,11]{A^n,A\in\PowNE A}{\text{\parbox[t]{\linewidth}{\raggedright 13,8; Definition der nichtleeren Potenzmenge}}}
 \proofstepwidestar[1,8,10,11]{f[A^nA]=f[A^n]f[A]}{\FormulaRefAuto{SemigroupHomPreservesPowerProduct}[thm]{6,14}}
 \proofstepwidestar[1,8,10,11]{f[A^{n+1}]=B^{n+1}}{\FormulaRefAuto{SemigroupPositiveCarrierPowersDef}[def]{15,12,7}}
 \proofstepwidestar[1,8]{\forall n\geq1\;(f[A^n]=B^n\Rightarrow f[A^{n+1}]=B^{n+1})}{\rUI{\rRI{10,\rRI{11,16}}}}
 \proofstepwidestar[1,8]{\forall n\geq1\;(f[A^n]=B^n)}{\FormulaRefAuto{PeanoInductionFromOne}[thm]{9,17}}
\closeproofpartwideindR
\proofpartwide{Isomorphe Einschränkung}
 \proofstepwidestar[1]{f:A\cong B}{\rA}
 \proofstepwidestar[1]{\SemigroupAlg{A}{\star}}{\FormulaRefAuto{SemigroupIsoSource}[thm]{1}}
 \proofstepwidestar[1]{\SemigroupAlg{B}{\diamond}}{\FormulaRefAuto{SemigroupIsoTarget}[thm]{1}}
 \proofstepwidestar[1]{f:A\to B}{\FormulaRefAuto{SemigroupIsoFunction}[thm]{1}}
 \proofstepwidestar[1]{f:A\bij B}{\FormulaRefAuto{SemigroupIsoBijectiveAxiom}[ax]{1}}
 \proofstepwidestar[1]{\SemigroupHom{f}{A,\star}{B,\diamond}}{\FormulaRefAuto{SemigroupIsoHomAxiom}[ax]{1}}
 \proofstepwidestar[1]{f[A]=B}{\FormulaRefAuto{SemigroupIsoFullImage}[thm]{1}}
 \proofstepwidestar[8]{A\neq\varnothing}{\rA}
 \proofstepwidestar[9]{n\geq1}{\rA}
 \proofstepwidestar[1,8,9]{\mathsf{UHGr}(A^n;A,\star)}{\FormulaRefAuto{SemigroupCarrierPowersSubsemigroups}[thm]{2,8,9}}
 \proofstepwidestar[1,8,9]{A^n\subseteq A}{\FormulaRefAuto{SemigroupSubsemigroupSubsetAxiom}[ax]{2,10}}
 \proofstepwidestar[1,8,9]{\SemigroupAlg{A^n}{\star}}{\FormulaRefAuto{SubsemigroupFormsSemigroup}[thm]{2,10}}
 \proofstepwidestar[1,8,9]{f|_{A^n}:A^n\cong f[A^n]}{\FormulaRefAuto{SemigroupIsoClosedCarrierRestriction}[thm]{1,11,12}}
 \proofstepwidestar[1,8,9]{f[A^n]=B^n}{\FormulaRefAuto{SemigroupIsoPositivePowersImage}[thm]{1,8,9}}
 \proofstepwidestar[1,8,9]{f|_{A^n}:A^n\cong B^n}{\rIE{14,13}}
\closeproofpartwide
\end{tabproofsplitwide}

\FormulaThmDeltaKR[Isomorphismen transportieren Erzeugnisse und Monogene]{%
 \begin{aligned}[t]
  &f:A\cong B\dsep\varnothing\neq X\subseteq A\dsep a\in A\vdash{}\\[-2pt]
  &\text{(i)}\quad f[\langle X\rangle_{A,\star}]=\langle f[X]\rangle_{B,\diamond}\\[-2pt]
  &\text{(ii)}\quad f|_{\langle X\rangle}:\langle X\rangle_{A,\star} \cong\langle f[X]\rangle_{B,\diamond}\\[-2pt]
  &\text{(iii)}\quad f(a^n)=f(a)^n\quad(n\geq1)\\[-2pt]
  &\text{(iv)}\quad f|_{\langle a\rangle}:\langle a\rangle_{A,\star} \cong\langle f(a)\rangle_{B,\diamond}
 \end{aligned}%
}{%
 \begin{aligned}[t]
 &f:A\cong B\dsep\varnothing\neq X\subseteq A\dsep a\in A\vdash{}\\[-2pt]
 &f[\langle X\rangle_{A,\star}]=\langle f[X]\rangle_{B,\diamond},\\[-2pt]
 &f|_{\langle X\rangle}:\langle X\rangle_{A,\star}
      \cong\langle f[X]\rangle_{B,\diamond},\\[-2pt]
 &f(a^n)=f(a)^n\quad(n\geq1),\qquad
   f|_{\langle a\rangle}:\langle a\rangle_{A,\star}
      \cong\langle f(a)\rangle_{B,\diamond}.
 \end{aligned}%
}{SemigroupIsoGeneratedSubsemigroups}{%
 \DeltaRow{Mengen}{A\dsep B\dsep X\dsep a\dsep n\dsep U\dsep V}
 \DeltaRow{Funktionen}{f\dsep h\dsep\star\dsep\diamond}
}
\begin{tabproofsplitwide}
\proofpartwideindR[Die beiden minimalen Unterhalbgruppen]{%
f:A\cong B\dsep\varnothing\neq X\subseteq A\vdash f|_{\langle X\rangle}:\langle X\rangle\cong\langle f[X]\rangle%
}{f:A\cong B\dsep\varnothing\neq X\subseteq A\vdash f|_{\langle X\rangle}:\langle X\rangle\cong\langle f[X]\rangle}[SemigroupIsoGeneratedImage]
 \proofstepwidestar[1]{f:A\cong B}{\rA}
 \proofstepwidestar[1]{\SemigroupAlg{A}{\star}}{\FormulaRefAuto{SemigroupIsoSource}[thm]{1}}
 \proofstepwidestar[1]{\SemigroupAlg{B}{\diamond}}{\FormulaRefAuto{SemigroupIsoTarget}[thm]{1}}
 \proofstepwidestar[1]{f:A\to B}{\FormulaRefAuto{SemigroupIsoFunction}[thm]{1}}
 \proofstepwidestar[1]{f:A\bij B}{\FormulaRefAuto{SemigroupIsoBijectiveAxiom}[ax]{1}}
 \proofstepwidestar[1]{\SemigroupHom{f}{A,\star}{B,\diamond}}{\FormulaRefAuto{SemigroupIsoHomAxiom}[ax]{1}}
 \proofstepwidestar[1]{f[A]=B}{\FormulaRefAuto{SemigroupIsoFullImage}[thm]{1}}
 \proofstepwidestar[8]{X\subseteq A}{\rA}
 \proofstepwidestar[9]{X\neq\varnothing}{\rA}
 \proofstepwidestar[1,8,9]{\mathsf{UHGr}(\langle X\rangle;A,\star)}{\FormulaRefAuto{GeneratedSubsemigroupIsSubsemigroup}[thm]{2,8,9}}
 \proofstepwidestar[1,8,9]{X\subseteq\langle X\rangle}{\FormulaRefAuto{GeneratedSubsemigroupContainsGenerators}[thm]{2,8,9}}
 \proofstepwidestar[1,8,9]{\langle X\rangle\subseteq A}{\FormulaRefAuto{SemigroupSubsemigroupSubsetAxiom}[ax]{2,10}}
 \proofstepwidestar[1,8,9]{\mathsf{UHGr}(f[\langle X\rangle];B,\diamond)}{\FormulaRefAuto{SemigroupIsoUHGrImage}[thm]{1,10}}
 \proofstepwidestar[1,8,9]{f[X]\subseteq f[\langle X\rangle]}{\FormulaRefAuto{F\colon U\to V\dsep A\subseteq B\dsep B\subseteq U\vdash F[A]\subseteq F[B]}[thm]{4,11,12}}
 \proofstepwidestar[8,9]{X\in\PowNE A}{\text{\parbox[t]{\linewidth}{\raggedright Definition der nichtleeren Potenzmenge}}}
 \proofstepwidestar[1,8,9]{f[X]\in\PowNE B}{\FormulaRefAuto{NonemptyDirectImage}[thm]{4,15}}
 \proofstepwidestar[1,8,9]{f[X]\subseteq B\land f[X]\neq\varnothing}{\FormulaRefAuto{NonemptyPowersetMembership}[thm]{16}}
  \proofstepwidestar[1,8,9]{f[X]\subseteq B}{\rAEa{17}}
 \proofstepwidestar[1,8,9]{f[X]\neq\varnothing}{\rAEb{17}}
 \proofstepwidestar[1,8,9]{\langle f[X]\rangle\subseteq f[\langle X\rangle]}{\FormulaRefAuto{GeneratedSubsemigroupMinimal}[thm]{3,18,19,13,14}}
 \proofstepwidestar[1]{h:B\cong A}{\FormulaRefAuto{SemigroupIsoInverse}[thm]{1}}
 \proofstepwidestar[1,8,9]{\mathsf{UHGr}(\langle f[X]\rangle;B,\diamond)}{\FormulaRefAuto{GeneratedSubsemigroupIsSubsemigroup}[thm]{3,18,19}}
 \proofstepwidestar[1,8,9]{f[X]\subseteq\langle f[X]\rangle}{\FormulaRefAuto{GeneratedSubsemigroupContainsGenerators}[thm]{3,18,19}}
 \proofstepwidestar[1,8,9]{\langle f[X]\rangle\subseteq B}{\FormulaRefAuto{SemigroupSubsemigroupSubsetAxiom}[ax]{3,22}}
 \proofstepwidestar[1,8,9]{\mathsf{UHGr}(h[\langle f[X]\rangle];A,\star)}{\FormulaRefAuto{SemigroupIsoUHGrImage}[thm]{21,22}}
 \proofstepwidestar[1]{h:B\to A}{\FormulaRefAuto{SemigroupIsoFunction}[thm]{21}}
 \proofstepwidestar[1,8,9]{h[f[X]]\subseteq h[\langle f[X]\rangle]}{\FormulaRefAuto{F\colon U\to V\dsep A\subseteq B\dsep B\subseteq U\vdash F[A]\subseteq F[B]}[thm]{26,23,24}}
 \proofstepwidestar[1,8]{f[X]\subseteq B\land h[f[X]]=X\land(X\neq\varnothing\leftrightarrow f[X]\neq\varnothing)}{\FormulaRefAuto{SemigroupIsoImageCarrierData}[thm]{1,8}}
 \proofstepwidestar[1,8]{h[f[X]]=X}{\rAEa{\rAEb{28}}}
 \proofstepwidestar[1,8,9]{X\subseteq h[\langle f[X]\rangle]}{\rIE{29,27}}
 \proofstepwidestar[1,8,9]{\langle X\rangle\subseteq h[\langle f[X]\rangle]}{\FormulaRefAuto{GeneratedSubsemigroupMinimal}[thm]{2,8,9,25,30}}
 \proofstepwidestar[1,8,9]{f[\langle X\rangle]\subseteq\langle f[X]\rangle}{\text{\parbox[t]{\linewidth}{\raggedright Bildmonotonie aus 31; Umkehrgleichung auf \(B\)}}}
 \proofstepwidestar[1,8,9]{f[\langle X\rangle]=\langle f[X]\rangle}{\FormulaRefAuto{A\subseteq B,\,B\subseteq A\vdash A=B}[thm]{32,20}}
 \proofstepwidestar[1,8,9]{\SemigroupAlg{\langle X\rangle}{\star}}{\FormulaRefAuto{SubsemigroupFormsSemigroup}[thm]{2,10}}
 \proofstepwidestar[1,8,9]{f|_{\langle X\rangle}:\langle X\rangle\cong f[\langle X\rangle]}{\FormulaRefAuto{SemigroupIsoClosedCarrierRestriction}[thm]{1,12,34}}
 \proofstepwidestar[1,8,9]{f|_{\langle X\rangle}:\langle X\rangle\cong\langle f[X]\rangle}{\rIE{33,35}}
\closeproofpartwideindR
\proofpartwide{Elementpotenzen und Monogene}
 \proofstepwidestar[1]{f:A\cong B}{\rA}
 \proofstepwidestar[1]{\SemigroupAlg{A}{\star}}{\FormulaRefAuto{SemigroupIsoSource}[thm]{1}}
 \proofstepwidestar[1]{\SemigroupAlg{B}{\diamond}}{\FormulaRefAuto{SemigroupIsoTarget}[thm]{1}}
 \proofstepwidestar[1]{f:A\to B}{\FormulaRefAuto{SemigroupIsoFunction}[thm]{1}}
 \proofstepwidestar[1]{f:A\bij B}{\FormulaRefAuto{SemigroupIsoBijectiveAxiom}[ax]{1}}
 \proofstepwidestar[1]{\SemigroupHom{f}{A,\star}{B,\diamond}}{\FormulaRefAuto{SemigroupIsoHomAxiom}[ax]{1}}
 \proofstepwidestar[1]{f[A]=B}{\FormulaRefAuto{SemigroupIsoFullImage}[thm]{1}}
 \proofstepwidestar[8]{a\in A}{\rA}
 \proofstepwidestar[1,8]{f(a)\in B}{\FormulaRefAuto{F\colon A\to B\dsep x\in A\vdash F(x)\in B}[thm]{4,8}}
 \proofstepwidestar[1,8]{a^1=a}{\FormulaRefAuto{SemigroupPowerOne}[thm]{2,8}}
 \proofstepwidestar[1,8]{f(a)^1=f(a)}{\FormulaRefAuto{SemigroupPowerOne}[thm]{3,9}}
 \proofstepwidestar[1,8]{f(a^1)=f(a)}{\rIE{10,\rII}}
  \proofstepwidestar[1,8]{f(a)=f(a)^1}{\FormulaRefAuto{a=b\vdash b=a}[thm]{11}}
  \proofstepwidestar[1,8]{f(a^1)=f(a)^1}{\rIE{13,12}}
 \proofstepwidestar[15]{n\geq1}{\rA}
 \proofstepwidestar[16]{ f(a^n)=f(a)^n}{\rA}
 \proofstepwidestar[15,16]{n\geq1\land f(a^n)=f(a)^n}{\rAIStar{15,16}}
 \proofstepwidestar[1,8,15]{a^n\in A}{\FormulaRefAuto{SemigroupPowerClosure}[thm]{2,8,15}}
 \proofstepwidestar[1,8,15]{a^{n+1}=a^n\star a}{\FormulaRefAuto{SemigroupPowerSuccessor}[thm]{2,8,15}}
 \proofstepwidestar[1,8,15]{f(a^n\star a)=f(a^n)\diamond f(a)}{\FormulaRefAuto{SemigroupIsoOperation}[thm]{1,18,8}}
 \proofstepwidestar[1,8,15]{f(a)^{n+1}=f(a)^n\diamond f(a)}{\FormulaRefAuto{SemigroupPowerSuccessor}[thm]{3,9,15}}
 \proofstepwidestar[1,8,15]{a^n\star a=a^{n+1}}{\FormulaRefAuto{a=b\vdash b=a}[thm]{19}}
  \proofstepwidestar[1,8,15]{f(a)^n\diamond f(a)=f(a)^{n+1}}{\FormulaRefAuto{a=b\vdash b=a}[thm]{21}}
  \proofstepwidestar[1,8,15,16]{f(a^{n+1})=f(a)^{n+1}}{\rIE{23,\rIE{16,\rIE{22,20}}}}
 \proofstepwidestar[1,8]{\forall n\geq1\;(f(a^n)=f(a)^n\Rightarrow f(a^{n+1})=f(a)^{n+1})}{\rUI{\rRI{15,\rRI{16,24}}}}
 \proofstepwidestar[1,8]{\forall n\geq1\;(f(a^n)=f(a)^n)}{\FormulaRefAuto{PeanoInductionFromOne}[thm]{14,25}}
 \proofstepwidestar[8]{\{a\}\subseteq A}{\FormulaRefAuto{a\in A\vdash\{a\}\subseteq A}[thm]{8}}
 \proofstepwidestar[]{a\in\{a\}}{\FormulaRefAuto{a\in\{a\}}[thm]}
  \proofstepwidestar[]{\{a\}\neq\varnothing}{\FormulaRefAuto{a\in A\vdash A\neq\varnothing}[thm]{28}}
 \proofstepwidestar[1,8]{f|_{\langle\{a\}\rangle}:\langle\{a\}\rangle\cong\langle f[\{a\}]\rangle}{\FormulaRefAuto{SemigroupIsoGeneratedImage}[thm]{1,27,29}}
 \proofstepwidestar[1,8]{f[\{a\}]=\{f(a)\}}{\FormulaRefAuto{x\in A\dsep F\colon A\to B\vdash F[\{x\}]=\{F(x)\}}[thm]{8,4}}
 \proofstepwidestar[1,8]{f|_{\langle a\rangle}:\langle a\rangle\cong\langle f(a)\rangle}{\FormulaRefAuto{SemigroupMonogenicSubsemigroupDef}[def]{2,3,30,31}}
\closeproofpartwide
\end{tabproofsplitwide}

\subsection{Hauptideale, Teilbarkeit und Green-Relationen}

\FormulaThmDeltaK[Einseitige Ideale als Unterhalbgruppen]{%
 \SemigroupAlg{A}{\star}\dsep
 (\mathsf{LId}(I;A,\star)\lor\mathsf{RId}(I;A,\star)\lor\mathsf{Id}(I;A,\star))
 \vdash\mathsf{UHGr}(I;A,\star)%
}{SemigroupIdealsAreSubsemigroups}{%
 \DeltaRow{Mengen}{A\dsep I\dsep x\dsep y}
 \DeltaRow{Funktionen}{\star}
}
\begin{tabproofsplitwide}
\proofpartwideindR[Linksideale]{%
\SemigroupAlg{A}{\star}\dsep\mathsf{LId}(I;A,\star)\vdash\mathsf{UHGr}(I;A,\star)%
}{\SemigroupAlg{A}{\star}\dsep\mathsf{LId}(I;A,\star)\vdash\mathsf{UHGr}(I;A,\star)}[SemigroupLeftIdealIsSubsemigroup]
 \proofstepwidestar[1]{\SemigroupAlg{A}{\star}}{\rA}
 \proofstepwidestar[2]{\mathsf{LId}(I;A,\star)}{\rA}
 \proofstepwidestar[1,2]{I\subseteq A}{\FormulaRefAuto{SemigroupLeftIdealSubsetAxiom}[ax]{1,2}}
 \proofstepwidestar[1,2]{I\neq\varnothing}{\FormulaRefAuto{SemigroupLeftIdealNonemptyAxiom}[ax]{1,2}}
 \proofstepwidestar[5]{x\in I}{\rA}
 \proofstepwidestar[6]{y\in I}{\rA}
 \proofstepwidestar[5,6]{x,y\in I}{\rAIStar{5,6}}
 \proofstepwidestar[1,2,5,6]{x,y\in A}{\FormulaRefAuto{A\subseteq B,\,x\in A\vdash x\in B}[thm]{3,7}}
 \proofstepwidestar[1,2,5,6]{x\star y\in I}{\FormulaRefAuto{SemigroupLeftIdealClosureAxiom}[ax]{1,2,8,7}}
 \proofstepwidestar[1,2]{\forall x,y\in I\;(x\star y\in I)}{\rUI{\rRI{5,\rUI{\rRI{6,9}}}}}
 \proofstepwidestar[1,2]{I\subseteq A\land I\neq\varnothing\land\forall x,y\in I\;(x\star y\in I)}{\rAI{3,\rAI{4,10}}}
 \proofstepwidestar[1,2]{\mathsf{UHGr}(I;A,\star)}{\FormulaRefAuto{SemigroupSubsemigroupCriterion}[thm]{1,11}}
\closeproofpartwideindR
\proofpartwideindR[Rechtsideale]{%
\SemigroupAlg{A}{\star}\dsep\mathsf{RId}(I;A,\star)\vdash\mathsf{UHGr}(I;A,\star)%
}{\SemigroupAlg{A}{\star}\dsep\mathsf{RId}(I;A,\star)\vdash\mathsf{UHGr}(I;A,\star)}[SemigroupRightIdealIsSubsemigroup]
 \proofstepwidestar[1]{\SemigroupAlg{A}{\star}}{\rA}
 \proofstepwidestar[2]{\mathsf{RId}(I;A,\star)}{\rA}
 \proofstepwidestar[1,2]{I\subseteq A}{\FormulaRefAuto{SemigroupRightIdealSubsetAxiom}[ax]{1,2}}
 \proofstepwidestar[1,2]{I\neq\varnothing}{\FormulaRefAuto{SemigroupRightIdealNonemptyAxiom}[ax]{1,2}}
 \proofstepwidestar[5]{x\in I}{\rA}
 \proofstepwidestar[6]{y\in I}{\rA}
 \proofstepwidestar[5,6]{x,y\in I}{\rAIStar{5,6}}
 \proofstepwidestar[1,2,5,6]{x,y\in A}{\FormulaRefAuto{A\subseteq B,\,x\in A\vdash x\in B}[thm]{3,7}}
 \proofstepwidestar[1,2,5,6]{x\star y\in I}{\FormulaRefAuto{SemigroupRightIdealClosureAxiom}[ax]{1,2,8,7}}
 \proofstepwidestar[1,2]{\forall x,y\in I\;(x\star y\in I)}{\rUI{\rRI{5,\rUI{\rRI{6,9}}}}}
 \proofstepwidestar[1,2]{I\subseteq A\land I\neq\varnothing\land\forall x,y\in I\;(x\star y\in I)}{\rAI{3,\rAI{4,10}}}
 \proofstepwidestar[1,2]{\mathsf{UHGr}(I;A,\star)}{\FormulaRefAuto{SemigroupSubsemigroupCriterion}[thm]{1,11}}
\closeproofpartwideindR
\proofpartwideindR[Zweiseitige Ideale]{%
\SemigroupAlg{A}{\star}\dsep\mathsf{Id}(I;A,\star)\vdash\mathsf{UHGr}(I;A,\star)%
}{\SemigroupAlg{A}{\star}\dsep\mathsf{Id}(I;A,\star)\vdash\mathsf{UHGr}(I;A,\star)}[SemigroupIdealIsSubsemigroup]
 \proofstepwidestar[1]{\SemigroupAlg{A}{\star}}{\rA}
 \proofstepwidestar[2]{\mathsf{Id}(I;A,\star)}{\rA}
 \proofstepwidestar[1,2]{\mathsf{LId}(I;A,\star)}{\FormulaRefAuto{SemigroupIdealLeftAxiom}[ax]{1,2}}
 \proofstepwidestar[1,2]{\mathsf{UHGr}(I;A,\star)}{\FormulaRefAuto{SemigroupLeftIdealIsSubsemigroup}[thm]{1,3}}
\closeproofpartwideindR
\end{tabproofsplitwide}



\FormulaThmDeltaKR[Hauptideale und Green-Relationen]{%
 \begin{aligned}[t]
  &f:A\cong B\dsep a,b\in A\vdash{}\\[-2pt]
  &\text{(i)}\quad f[K_{A,\star}(a)]=K_{B,\diamond}(f(a))\quad(K\in\{L,R,J\})\\[-2pt]
  &\text{(ii)}\quad f|_{K_{A,\star}(a)}:K_{A,\star}(a)\cong K_{B,\diamond}(f(a))\\[-2pt]
  &\text{(iii)}\quad a\mathrel{\mathcal K_{A,\star}}b\leftrightarrow f(a)\mathrel{\mathcal K_{B,\diamond}}f(b)\\[-2pt]
  &\qquad(\mathcal K\in\{\mathcal L,\mathcal R,\mathcal J,\mathcal H\})
 \end{aligned}%
}{%
 \begin{aligned}[t]
 &f:A\cong B\dsep a,b\in A\vdash{}\\[-2pt]
 &f[K_{A,\star}(a)]=K_{B,\diamond}(f(a))\quad(K\in\{L,R,J\}),\\[-2pt]
 &f|_{K_{A,\star}(a)}:K_{A,\star}(a)\cong K_{B,\diamond}(f(a)),\\[-2pt]
 &a\mathrel{\mathcal K_{A,\star}}b\leftrightarrow
      f(a)\mathrel{\mathcal K_{B,\diamond}}f(b)
       \quad(\mathcal K\in\{\mathcal L,\mathcal R,\mathcal J,\mathcal H\}).
 \end{aligned}%
}{SemigroupIsoPrincipalIdealsAndGreen}{%
 \DeltaRow{Mengen}{A\dsep B\dsep a\dsep b\dsep x\dsep y}
 \DeltaRow{Funktionen}{f\dsep h\dsep\star\dsep\diamond}
}
\begin{tabproofsplitwide}
\proofpartwideindR[Transport der vier möglichen Produktformen]{%
f:A\cong B\dsep a\in A\vdash\forall K\in\{L,R,J\}\;(f[K_A(a)]=K_B(f(a)))%
}{f:A\cong B\dsep a\in A\vdash\forall K\in\{L,R,J\}\;(f[K_A(a)]=K_B(f(a)))}[SemigroupIsoPrincipalIdealImages]
 \proofstepwidestar[1]{f:A\cong B}{\rA}
 \proofstepwidestar[1]{\SemigroupAlg{A}{\star}}{\FormulaRefAuto{SemigroupIsoSource}[thm]{1}}
 \proofstepwidestar[1]{\SemigroupAlg{B}{\diamond}}{\FormulaRefAuto{SemigroupIsoTarget}[thm]{1}}
 \proofstepwidestar[1]{f:A\to B}{\FormulaRefAuto{SemigroupIsoFunction}[thm]{1}}
 \proofstepwidestar[1]{f:A\bij B}{\FormulaRefAuto{SemigroupIsoBijectiveAxiom}[ax]{1}}
 \proofstepwidestar[1]{\SemigroupHom{f}{A,\star}{B,\diamond}}{\FormulaRefAuto{SemigroupIsoHomAxiom}[ax]{1}}
 \proofstepwidestar[1]{f[A]=B}{\FormulaRefAuto{SemigroupIsoFullImage}[thm]{1}}
 \proofstepwidestar[8]{a\in A}{\rA}
 \proofstepwidestar[1,8]{f(a)\in B}{\FormulaRefAuto{F\colon A\to B\dsep x\in A\vdash F(x)\in B}[thm]{4,8}}
 \proofstepwidestar[1,8]{f[\{a\}]=\{f(a)\}}{\FormulaRefAuto{x\in A\dsep F\colon A\to B\vdash F[\{x\}]=\{F(x)\}}[thm]{8,4}}
 \proofstepwidestar[8]{A,\{a\}\in\PowNE A}{\text{\parbox[t]{\linewidth}{\raggedright Die Teilmengen enthalten \(a\); Definition von \(\PowNE A\).}}}
 \proofstepwidestar[1,8]{f[A\{a\}]=f[A]f[\{a\}]}{\FormulaRefAuto{SemigroupHomPreservesPowerProduct}[thm]{6,11}}
 \proofstepwidestar[1,8]{f[\{a\}A]=f[\{a\}]f[A]}{\FormulaRefAuto{SemigroupHomPreservesPowerProduct}[thm]{6,11}}
 \proofstepwidestar[1,8]{A\{a\}\in\PowNE A}{\FormulaRefAuto{PowerSemigroupProductClosure}[thm]{2,11}}
 \proofstepwidestar[1,8]{f[(A\{a\})A]=f[A\{a\}]f[A]}{\FormulaRefAuto{SemigroupHomPreservesPowerProduct}[thm]{6,14,11}}
 \proofstepwidestar[1,8]{L_A(a)=\{a\}\cup A\{a\}}{\text{\parbox[t]{\linewidth}{\raggedright \FormulaRefAuto{SemigroupPrincipalLeftIdealDef}[def]{2,8}; Mengenproduktdefinition und Abschluss.}}}
 \proofstepwidestar[1,8]{L_B(f(a))=\{a\}\cup B\{a\}}{\text{\parbox[t]{\linewidth}{\raggedright \FormulaRefAuto{SemigroupPrincipalLeftIdealDef}[def]{3,9}; Mengenproduktdefinition und Abschluss.}}}
 \proofstepwidestar[1,8]{f[L_A(a)]=L_B(f(a))}{\text{\parbox[t]{\linewidth}{\raggedright 10,7,12,13,15,16,17; dieselben Existenzzeugen in jeder Vereinigung.}}}
 \proofstepwidestar[1,8]{R_A(a)=\{a\}\cup\{a\}A}{\text{\parbox[t]{\linewidth}{\raggedright \FormulaRefAuto{SemigroupPrincipalRightIdealDef}[def]{2,8}; Mengenproduktdefinition und Abschluss.}}}
 \proofstepwidestar[1,8]{R_B(f(a))=\{a\}\cup\{a\}B}{\text{\parbox[t]{\linewidth}{\raggedright \FormulaRefAuto{SemigroupPrincipalRightIdealDef}[def]{3,9}; Mengenproduktdefinition und Abschluss.}}}
 \proofstepwidestar[1,8]{f[R_A(a)]=R_B(f(a))}{\text{\parbox[t]{\linewidth}{\raggedright 10,7,12,13,15,19,20; dieselben Existenzzeugen in jeder Vereinigung.}}}
 \proofstepwidestar[1,8]{J_A(a)=\{a\}\cup A\{a\}\cup\{a\}A\cup(A\{a\})A}{\text{\parbox[t]{\linewidth}{\raggedright \FormulaRefAuto{SemigroupPrincipalTwoSidedIdealDef}[def]{2,8}; Mengenproduktdefinition und Abschluss.}}}
 \proofstepwidestar[1,8]{J_B(f(a))=\{a\}\cup B\{a\}\cup\{a\}B\cup(B\{a\})B}{\text{\parbox[t]{\linewidth}{\raggedright \FormulaRefAuto{SemigroupPrincipalTwoSidedIdealDef}[def]{3,9}; Mengenproduktdefinition und Abschluss.}}}
 \proofstepwidestar[1,8]{f[J_A(a)]=J_B(f(a))}{\text{\parbox[t]{\linewidth}{\raggedright 10,7,12,13,15,22,23; dieselben Existenzzeugen in jeder Vereinigung.}}}
 \proofstepwidestar[1,8]{\forall K\in\{L,R,J\}\;(f[K_A(a)]=K_B(f(a)))}{\text{\parbox[t]{\linewidth}{\raggedright Die drei Fälle 18,21,24}}}
\closeproofpartwideindR
\proofpartwide{Einschränkungen und Relationstransport}
 \proofstepwidestar[1]{f:A\cong B}{\rA}
 \proofstepwidestar[1]{\SemigroupAlg{A}{\star}}{\FormulaRefAuto{SemigroupIsoSource}[thm]{1}}
 \proofstepwidestar[1]{\SemigroupAlg{B}{\diamond}}{\FormulaRefAuto{SemigroupIsoTarget}[thm]{1}}
 \proofstepwidestar[1]{f:A\to B}{\FormulaRefAuto{SemigroupIsoFunction}[thm]{1}}
 \proofstepwidestar[1]{f:A\bij B}{\FormulaRefAuto{SemigroupIsoBijectiveAxiom}[ax]{1}}
 \proofstepwidestar[1]{\SemigroupHom{f}{A,\star}{B,\diamond}}{\FormulaRefAuto{SemigroupIsoHomAxiom}[ax]{1}}
 \proofstepwidestar[1]{f[A]=B}{\FormulaRefAuto{SemigroupIsoFullImage}[thm]{1}}
 \proofstepwidestar[8]{a,b\in A}{\rA}
 \proofstepwidestar[1,8]{\forall K\in\{L,R,J\}\;(f[K_A(a)]=K_B(f(a)))}{\FormulaRefAuto{SemigroupIsoPrincipalIdealImages}[thm]{1,8}}
 \proofstepwidestar[1,8]{\forall K\in\{L,R,J\}\;(f[K_A(b)]=K_B(f(b)))}{\FormulaRefAuto{SemigroupIsoPrincipalIdealImages}[thm]{1,8}}
 \proofstepwidestar[1,8]{\mathsf{LId}(L_A(a);A,\star)}{\FormulaRefAuto{SemigroupPrincipalLeftIdealIsLeftIdeal}[thm]{2,8}}
 \proofstepwidestar[1,8]{\mathsf{UHGr}(L_A(a);A,\star)}{\FormulaRefAuto{SemigroupLeftIdealIsSubsemigroup}[thm]{2,11}}
 \proofstepwidestar[1,8]{L_A(a)\subseteq A}{\FormulaRefAuto{SemigroupSubsemigroupSubsetAxiom}[ax]{2,12}}
 \proofstepwidestar[1,8]{\SemigroupAlg{L_A(a)}{\star}}{\FormulaRefAuto{SubsemigroupFormsSemigroup}[thm]{2,12}}
 \proofstepwidestar[1,8]{f|_{L_A(a)}:L_A(a)\cong f[L_A(a)]}{\FormulaRefAuto{SemigroupIsoClosedCarrierRestriction}[thm]{1,13,14}}
 \proofstepwidestar[1,8]{f|_{L_A(a)}:L_A(a)\cong L_B(f(a))}{\rIE{9,15}}
 \proofstepwidestar[1,8]{\mathsf{RId}(R_A(a);A,\star)}{\FormulaRefAuto{SemigroupPrincipalRightIdealIsRightIdeal}[thm]{2,8}}
 \proofstepwidestar[1,8]{\mathsf{UHGr}(R_A(a);A,\star)}{\FormulaRefAuto{SemigroupRightIdealIsSubsemigroup}[thm]{2,17}}
 \proofstepwidestar[1,8]{R_A(a)\subseteq A}{\FormulaRefAuto{SemigroupSubsemigroupSubsetAxiom}[ax]{2,18}}
 \proofstepwidestar[1,8]{\SemigroupAlg{R_A(a)}{\star}}{\FormulaRefAuto{SubsemigroupFormsSemigroup}[thm]{2,18}}
 \proofstepwidestar[1,8]{f|_{R_A(a)}:R_A(a)\cong f[R_A(a)]}{\FormulaRefAuto{SemigroupIsoClosedCarrierRestriction}[thm]{1,19,20}}
 \proofstepwidestar[1,8]{f|_{R_A(a)}:R_A(a)\cong R_B(f(a))}{\rIE{9,21}}
 \proofstepwidestar[1,8]{\mathsf{Id}(J_A(a);A,\star)}{\FormulaRefAuto{SemigroupPrincipalTwoSidedIdealIsIdeal}[thm]{2,8}}
 \proofstepwidestar[1,8]{\mathsf{UHGr}(J_A(a);A,\star)}{\FormulaRefAuto{SemigroupIdealIsSubsemigroup}[thm]{2,23}}
 \proofstepwidestar[1,8]{J_A(a)\subseteq A}{\FormulaRefAuto{SemigroupSubsemigroupSubsetAxiom}[ax]{2,24}}
 \proofstepwidestar[1,8]{\SemigroupAlg{J_A(a)}{\star}}{\FormulaRefAuto{SubsemigroupFormsSemigroup}[thm]{2,24}}
 \proofstepwidestar[1,8]{f|_{J_A(a)}:J_A(a)\cong f[J_A(a)]}{\FormulaRefAuto{SemigroupIsoClosedCarrierRestriction}[thm]{1,25,26}}
 \proofstepwidestar[1,8]{f|_{J_A(a)}:J_A(a)\cong J_B(f(a))}{\rIE{9,27}}
 \proofstepwidestar[1]{h:B\to A}{\text{\parbox[t]{\linewidth}{\raggedright \FormulaRefAuto{SemigroupIsoInverseCarrier}[thm]{1}; Abbildungseigenschaft der Bijektion.}}}
 \proofstepwidestar[1,8]{K_A(a)=K_A(b)\leftrightarrow f[K_A(a)]=f[K_A(b)]}{\text{\parbox[t]{\linewidth}{\raggedright Vorwärts: Gleichheitseinsetzung. Rückwärts: auf beide Bilder \(h\) anwenden; \(h(f(x))=x\) für jeden Zeugen.}}}
 \proofstepwidestar[1,8]{K_A(a)\subseteq K_A(b)\leftrightarrow f[K_A(a)]\subseteq f[K_A(b)]}{\text{\parbox[t]{\linewidth}{\raggedright Vorwärts: derselbe Bildzeuge. Rückwärts: \(f(x)=f(y)\) impliziert \(x=y\) wegen 5.}}}
 \proofstepwidestar[1,8]{a\mathrel{\mathcal L_A}b\leftrightarrow f(a)\mathrel{\mathcal L_B}f(b)}{\FormulaRefAuto{SemigroupGreenLDef}[def]{9,10,30}}
 \proofstepwidestar[1,8]{a\mathrel{\mathcal R_A}b\leftrightarrow f(a)\mathrel{\mathcal R_B}f(b)}{\FormulaRefAuto{SemigroupGreenRDef}[def]{9,10,30}}
 \proofstepwidestar[1,8]{a\mathrel{\mathcal J_A}b\leftrightarrow f(a)\mathrel{\mathcal J_B}f(b)}{\FormulaRefAuto{SemigroupGreenJDef}[def]{9,10,30}}
 \proofstepwidestar[1,8]{a\mathrel{\mathcal H_A}b\leftrightarrow f(a)\mathrel{\mathcal H_B}f(b)}{\FormulaRefAuto{SemigroupGreenHDef}[def]{32,33}}
\closeproofpartwide
\end{tabproofsplitwide}
Die Inklusionsaussage im zweiten Beweisteil liefert zugleich den Transport
der zugehörigen Teilbarkeitsvorordnungen. Die Einschränkung von \(f\) auf
eine Green-Klasse ist eine Bijektion auf die entsprechende Klasse in
\(B\). Eine Green-Klasse ist im Allgemeinen keine Unterhalbgruppe.

\subsection{Potenzhalbgruppen und endliche Teilmengen}

Für sämtliche nichtleeren Teilmengen ist die Abbildung \(X\mapsto f[X]\)
bereits durch \FormulaRefAuto{SemigroupIsoInducesPowerSemigroupIso}[thm]
als Isomorphismus bewiesen. Dies lässt sich beliebig oft anwenden.
Die Einermengenkopie wird durch
\FormulaRefAuto{SemigroupSingletonCopyIsomorphism}[thm] erfasst.

\FormulaDefDeltaK[Endliche nichtleere Potenzhalbgruppe]{%
 \mathcal P_{\mathrm{fin},+}(A)\coloneqq
 \{X\in\PowNE A\mid\Finite X\}%
}{SemigroupFinitePowerCarrierDef}{%
 \DeltaRow{Mengen}{A\dsep X}
 \DeltaRow{Operation}{\star_{\mathcal P}\text{ ist das Mengenprodukt.}}
}

\FormulaThmDeltaK[Mengenprodukte endlicher Teilmengen sind endlich]{%
 \SemigroupAlg{A}{\star}\dsep X,Y\subseteq A\dsep\Finite X\dsep\Finite Y\vdash
 \Finite{\{x\star y\mid x\in X,\ y\in Y\}}%
}{SemigroupFiniteSetProduct}{%
 \DeltaRow{Mengen}{A\dsep X\dsep Y\dsep U\dsep a\dsep x\dsep y}
 \DeltaRow{Funktionen}{\star\dsep\lambda_a}
}
Für diesen Beweis setzen wir
\(M(U,Y)=\{u\star y\mid u\in U,\ y\in Y\}\) und
\(P(U)\equiv(U\subseteq A\Rightarrow\Finite{M(U,Y)})\).
\begin{tabproofsplitwide}
\proofpartwideindR[Leerer erster Faktor]{%
\SemigroupAlg{A}{\star}\dsep Y\subseteq A\vdash\Finite{M(\varnothing,Y)}%
}{\SemigroupAlg{A}{\star}\dsep Y\subseteq A\vdash\Finite{M(\varnothing,Y)}}[SemigroupFiniteSetProductBase]
 \proofstepwidestar[1]{\SemigroupAlg{A}{\star}}{\rA}
 \proofstepwidestar[2]{Y\subseteq A}{\rA}
 \proofstepwidestar[]{\Finite\varnothing}{\FormulaRefAuto{FiniteEmptyAxiom}[ax]}
 \proofstepwidestar[]{M(\varnothing,Y)=\varnothing}{\text{\parbox[t]{\linewidth}{\raggedright Elementkriterium: es gibt keinen ersten Faktor.}}}
 \proofstepwidestar[]{\varnothing=M(\varnothing,Y)}{\FormulaRefAuto{a=b\vdash b=a}[thm]{4}}
  \proofstepwidestar[]{\Finite{M(\varnothing,Y)}}{\rIE{5,3}}

\closeproofpartwideindR
\proofpartwideindR[Adjunktion eines ersten Faktors]{%
\begin{aligned}[t]
 &\SemigroupAlg{A}{\star}\dsep Y\subseteq A\dsep\Finite Y\dsep U\subseteq A\dsep{}\\[-2pt]
 &\Finite{M(U,Y)}\dsep a\in A\vdash\Finite{M(U\cup\{a\},Y)}
\end{aligned}%
}{\SemigroupAlg{A}{\star}\dsep Y\subseteq A\dsep\Finite Y\dsep U\subseteq A\dsep\Finite{M(U,Y)}\dsep a\in A\vdash\Finite{M(U\cup\{a\},Y)}}[SemigroupFiniteSetProductStep]
 \proofstepwidestar[1]{\SemigroupAlg{A}{\star}}{\rA}
 \proofstepwidestar[2]{Y\subseteq A}{\rA}
 \proofstepwidestar[3]{\Finite Y}{\rA}
 \proofstepwidestar[4]{U\subseteq A}{\rA}
 \proofstepwidestar[5]{\Finite{M(U,Y)}}{\rA}
 \proofstepwidestar[6]{a\in A}{\rA}
 \proofstepwidestar[1,2,6]{\lambda_a:Y\to A,\quad\lambda_a(y)=a\star y}{\text{\parbox[t]{\linewidth}{\raggedright Abgeschlossenheit; Funktionsdefinition durch die eindeutigen Produktwerte}}}
 \proofstepwidestar[]{Y\subseteq Y}{\FormulaRefAuto{A\subseteq A}[thm]}
 \proofstepwidestar[1,2,3,6]{\Finite{\lambda_a[Y]}}{\FormulaRefAuto{FiniteImage}[thm]{7,8,3}}
 \proofstepwidestar[1,2,3,5,6]{\Finite{M(U,Y)\cup\lambda_a[Y]}}{\FormulaRefAuto{FiniteUnion}[thm]{5,9}}
 \proofstepwidestar[1,2,4,6]{M(U\cup\{a\},Y)=M(U,Y)\cup\lambda_a[Y]}{\text{\parbox[t]{\linewidth}{\raggedright Elementkriterium: der erste Faktor liegt in \(U\) oder ist \(a\).}}}
 \proofstepwidestar[1,2,4,6]{M(U,Y)\cup\lambda_a[Y]=M(U\cup\{a\},Y)}{\FormulaRefAuto{a=b\vdash b=a}[thm]{11}}
  \proofstepwidestar[1,2,3,4,5,6]{\Finite{M(U\cup\{a\},Y)}}{\rIE{12,10}}

\closeproofpartwideindR
\proofpartwide{Endlichkeitsinduktion}
 \proofstepwidestar[1]{\SemigroupAlg{A}{\star}}{\rA}
 \proofstepwidestar[2]{X\subseteq A}{\rA}
 \proofstepwidestar[3]{Y\subseteq A}{\rA}
 \proofstepwidestar[4]{\Finite X}{\rA}
 \proofstepwidestar[5]{\Finite Y}{\rA}
 \proofstepwidestar[6]{\varnothing\subseteq A}{\rA}
  \proofstepwidestar[1,3]{\Finite{M(\varnothing,Y)}}{\FormulaRefAuto{SemigroupFiniteSetProductBase}[thm]{1,3}}
  \proofstepwidestar[1,3]{P(\varnothing)}{\rRI{6,7}}
 \proofstepwidestar[9]{\Finite U\land P(U)\land a\notin U}{\rA}
 \proofstepwidestar[10]{U\cup\{a\}\subseteq A}{\rA}
 \proofstepwidestar[]{U\subseteq U\cup\{a\}}{\FormulaRefAuto{A\subseteq A\cup B}[thm]}
  \proofstepwidestar[10]{U\subseteq A}{\FormulaRefAuto{A\subseteq B\dsep B\subseteq C\vdash A\subseteq C}[thm]{11,10}}
 \proofstepwidestar[]{a\in U\cup\{a\}}{\FormulaRefAuto{a \in A \cup \{a\}}[thm]}
  \proofstepwidestar[10]{a\in A}{\FormulaRefAuto{A\subseteq B,\,x\in A\vdash x\in B}[thm]{10,13}}
 \proofstepwidestar[9,10]{\Finite{M(U,Y)}}{\rRE{\rAEa{\rAEb{9}},12}}
 \proofstepwidestar[1,3,5,9,10]{\Finite{M(U\cup\{a\},Y)}}{\FormulaRefAuto{SemigroupFiniteSetProductStep}[thm]{1,3,5,12,15,14}}
 \proofstepwidestar[1,3,5,9]{P(U\cup\{a\})}{\rRI{10,16}}
 \proofstepwidestar[1,3,4,5]{P(X)}{\FormulaRefAuto{FiniteInductionRule}[ax]{8,9,17,4}}
 \proofstepwidestar[1,2,3,4,5]{\Finite{M(X,Y)}}{\rRE{18,2}}
\closeproofpartwide
\end{tabproofsplitwide}

\FormulaThmDeltaK[Endliche Potenzhalbgruppen sind isomorphieinvariant]{%
 \begin{aligned}[t]
 &f:A\cong B\vdash{}\\[-2pt]
 &\SemigroupIso{\PowMapNE f|_{\mathcal P_{\mathrm{fin},+}(A)}}
 {\mathcal P_{\mathrm{fin},+}(A),\star_{\mathcal P}}
 {\mathcal P_{\mathrm{fin},+}(B),\diamond_{\mathcal P}}.
 \end{aligned}%
}{SemigroupIsoFinitePowerSemigroups}{%
 \DeltaRow{Mengen}{A\dsep B\dsep X\dsep Y}
 \DeltaRow{Funktionen}{f\dsep h\dsep\star_{\mathcal P}\dsep\diamond_{\mathcal P}}
}
\begin{tabproofsplitwide}
\proofpartwideindR[Abschluss der endlichen Teilmengen]{%
\SemigroupAlg{A}{\star}\vdash\SemigroupAlg{\mathcal P_{\mathrm{fin},+}(A)}{\star_{\mathcal P}}%
}{\SemigroupAlg{A}{\star}\vdash\SemigroupAlg{\mathcal P_{\mathrm{fin},+}(A)}{\star_{\mathcal P}}}[SemigroupFinitePowerSemigroup]
 \proofstepwidestar[1]{\SemigroupAlg{A}{\star}}{\rA}
 \proofstepwidestar[2]{X\in\mathcal P_{\mathrm{fin},+}(A)}{\rA}
 \proofstepwidestar[3]{Y\in\mathcal P_{\mathrm{fin},+}(A)}{\rA}
 \proofstepwidestar[2]{X\in\PowNE A}{\FormulaRefAuto{SemigroupFinitePowerCarrierDef}[def]{2}}
 \proofstepwidestar[3]{Y\in\PowNE A}{\FormulaRefAuto{SemigroupFinitePowerCarrierDef}[def]{3}}
 \proofstepwidestar[2]{\Finite X}{\FormulaRefAuto{SemigroupFinitePowerCarrierDef}[def]{2}}
 \proofstepwidestar[3]{\Finite Y}{\FormulaRefAuto{SemigroupFinitePowerCarrierDef}[def]{3}}
 \proofstepwidestar[2,3]{X,Y\subseteq A}{\text{\parbox[t]{\linewidth}{\raggedright Definition der nichtleeren Potenzmenge aus 4,5}}}
 \proofstepwidestar[1,2,3]{\Finite{XY}}{\FormulaRefAuto{SemigroupFiniteSetProduct}[thm]{1,8,6,7}}
 \proofstepwidestar[1,2,3]{XY\in\PowNE A}{\FormulaRefAuto{PowerSemigroupProductClosure}[thm]{1,4,5}}
 \proofstepwidestar[1,2,3]{XY\in\mathcal P_{\mathrm{fin},+}(A)}{\FormulaRefAuto{SemigroupFinitePowerCarrierDef}[def]{10,9}}
 \proofstepwidestar[1]{\forall X,Y\in\mathcal P_{\mathrm{fin},+}(A)\;(XY\in\mathcal P_{\mathrm{fin},+}(A))}{\rUI{\rRI{2,\rUI{\rRI{3,11}}}}}
 \proofstepwidestar[]{\mathcal P_{\mathrm{fin},+}(A)\subseteq\PowNE A}{\FormulaRefAuto{SemigroupFinitePowerCarrierDef}[def]}
 \proofstepwidestar[1]{\SemigroupAlg{\PowNE A}{\star_{\mathcal P}}}{\FormulaRefAuto{NonemptyPowerSemigroup}[thm]{1}}
 \proofstepwidestar[1]{\SemigroupAlg{\mathcal P_{\mathrm{fin},+}(A)}{\star_{\mathcal P}}}{\FormulaRefAuto{SemigroupClosedSubsetIsSemigroup}[thm]{14,13,12}}
\closeproofpartwideindR
\proofpartwide{Bild und Einschränkung}
 \proofstepwidestar[1]{f:A\cong B}{\rA}
 \proofstepwidestar[1]{\SemigroupAlg{A}{\star}}{\FormulaRefAuto{SemigroupIsoSource}[thm]{1}}
 \proofstepwidestar[1]{\SemigroupAlg{B}{\diamond}}{\FormulaRefAuto{SemigroupIsoTarget}[thm]{1}}
 \proofstepwidestar[1]{f:A\to B}{\FormulaRefAuto{SemigroupIsoFunction}[thm]{1}}
 \proofstepwidestar[1]{f:A\bij B}{\FormulaRefAuto{SemigroupIsoBijectiveAxiom}[ax]{1}}
 \proofstepwidestar[1]{\SemigroupHom{f}{A,\star}{B,\diamond}}{\FormulaRefAuto{SemigroupIsoHomAxiom}[ax]{1}}
 \proofstepwidestar[1]{f[A]=B}{\FormulaRefAuto{SemigroupIsoFullImage}[thm]{1}}
 \proofstepwidestar[1]{\PowMapNE f:\PowNE A\cong\PowNE B}{\FormulaRefAuto{SemigroupIsoInducesPowerSemigroupIso}[thm]{1}}
 \proofstepwidestar[9]{X\in\mathcal P_{\mathrm{fin},+}(A)}{\rA}
 \proofstepwidestar[9]{X\in\PowNE A}{\FormulaRefAuto{SemigroupFinitePowerCarrierDef}[def]{9}}
 \proofstepwidestar[9]{X\subseteq A}{\text{\parbox[t]{\linewidth}{\raggedright Definition der nichtleeren Potenzmenge aus 10}}}
 \proofstepwidestar[9]{\Finite X}{\FormulaRefAuto{SemigroupFinitePowerCarrierDef}[def]{9}}
 \proofstepwidestar[1,9]{\Finite{f[X]}}{\FormulaRefAuto{FiniteImage}[thm]{4,11,12}}
 \proofstepwidestar[1,9]{f[X]\in\PowNE B}{\FormulaRefAuto{NonemptyDirectImage}[thm]{4,10}}
 \proofstepwidestar[1,9]{f[X]\in\mathcal P_{\mathrm{fin},+}(B)}{\FormulaRefAuto{SemigroupFinitePowerCarrierDef}[def]{14,13}}
 \proofstepwidestar[1]{\forall X\in\mathcal P_{\mathrm{fin},+}(A)\;(f[X]\in\mathcal P_{\mathrm{fin},+}(B))}{\rUI{\rRI{9,15}}}
 \proofstepwidestar[1]{h:B\cong A}{\FormulaRefAuto{SemigroupIsoInverse}[thm]{1}}
 \proofstepwidestar[1]{h:B\to A}{\FormulaRefAuto{SemigroupIsoFunction}[thm]{17}}
 \proofstepwidestar[19]{Y\in\mathcal P_{\mathrm{fin},+}(B)}{\rA}
 \proofstepwidestar[19]{Y\in\PowNE B}{\FormulaRefAuto{SemigroupFinitePowerCarrierDef}[def]{19}}
 \proofstepwidestar[19]{Y\subseteq B}{\text{\parbox[t]{\linewidth}{\raggedright Definition der nichtleeren Potenzmenge aus 20}}}
 \proofstepwidestar[19]{\Finite Y}{\FormulaRefAuto{SemigroupFinitePowerCarrierDef}[def]{19}}
 \proofstepwidestar[1,19]{\Finite{h[Y]}}{\FormulaRefAuto{FiniteImage}[thm]{18,21,22}}
 \proofstepwidestar[1,19]{h[Y]\in\PowNE A}{\FormulaRefAuto{NonemptyDirectImage}[thm]{18,20}}
 \proofstepwidestar[1,19]{h[Y]\in\mathcal P_{\mathrm{fin},+}(A)}{\FormulaRefAuto{SemigroupFinitePowerCarrierDef}[def]{24,23}}
 \proofstepwidestar[1,19]{f[h[Y]]=Y}{\FormulaRefAuto{Y\subseteq B\dsep F\colon A\bij B\vdash F[F^{-1}[Y]]=Y}[thm]{21,5}}
 \proofstepwidestar[1]{\PowMapNE f[\mathcal P_{\mathrm{fin},+}(A)]=\mathcal P_{\mathrm{fin},+}(B)}{\text{\parbox[t]{\linewidth}{\raggedright Bilddefinition aus 16; für jedes 19 liefern 25,26 einen Urbildzeugen.}}}
 \proofstepwidestar[1]{\SemigroupAlg{\mathcal P_{\mathrm{fin},+}(A)}{\star_{\mathcal P}}}{\FormulaRefAuto{SemigroupFinitePowerSemigroup}[thm]{2}}
 \proofstepwidestar[1]{\SemigroupAlg{\mathcal P_{\mathrm{fin},+}(B)}{\diamond_{\mathcal P}}}{\FormulaRefAuto{SemigroupFinitePowerSemigroup}[thm]{3}}
 \proofstepwidestar[]{\mathcal P_{\mathrm{fin},+}(A)\subseteq\PowNE A}{\FormulaRefAuto{SemigroupFinitePowerCarrierDef}[def]}
 \proofstepwidestar[]{\mathcal P_{\mathrm{fin},+}(B)\subseteq\PowNE B}{\FormulaRefAuto{SemigroupFinitePowerCarrierDef}[def]}
 \proofstepwidestar[1]{\PowMapNE f|_{\mathcal P_{\mathrm{fin},+}(A)}:\mathcal P_{\mathrm{fin},+}(A)\cong\mathcal P_{\mathrm{fin},+}(B)}{\FormulaRefAuto{SemigroupIsoRestrictionToImage}[thm]{8,28,29,30,31,27}}
\closeproofpartwide
\end{tabproofsplitwide}

\subsection{Direkte Produkte, kartesische Potenzen und Gegenhalbgruppen}

\FormulaDefDeltaK[Direktes Produkt einer Halbgruppenfamilie]{%
 \begin{aligned}[t]
 &(\forall i\in I\;\SemigroupAlg{A_i}{\star_i})\vdash{}\\[-2pt]
 &\prod_{i\in I}A_i\coloneqq\{x:I\to\bigcup_{i\in I}A_i\mid
       \forall i\in I\;(x(i)\in A_i)\},\\[-2pt]
 &(x\odot y)(i)\coloneqq x(i)\star_i y(i).
 \end{aligned}%
}{SemigroupDirectProductDef}{%
 \DeltaRow{Mengen}{I\dsep A_i\dsep x\dsep y\dsep i}
 \DeltaRow{Funktionen}{\star_i\dsep\odot}
}
Bei konstanter Familie ist dies die kartesische Potenz \(A^I\).
Das leere Produkt besteht aus der leeren Funktion und ist eine
einelementige Halbgruppe. Eine nichtleere Indexmenge mit leerem
Produktträger verursacht für den hier verwendeten Halbgruppenbegriff
keinen Sonderfall: Abgeschlossenheit und Assoziativität gelten dann leer.


\FormulaDefDeltaK[Zugehörigkeit zur indizierten Vereinigung]{%
 u\in\bigcup_{i\in I}A_i\coloneqq
 \exists i\in I\;(u\in A_i)%
}{SemigroupIndexedUnionNotationDef}{%
 \DeltaRow{Mengenfamilie}{(A_i)_{i\in I}}
 \DeltaRow{Mengen}{I\dsep u\dsep i}
 \DeltaRow{Schreibweise}{\bigcup_{i\in I}A_i=\bigcup\{A_i\mid i\in I\}}
}

\FormulaDefDeltaK[Graphterm der punktweisen Produktoperation]{%
 \begin{aligned}[t]
 &(\forall i\in I\;\SemigroupAlg{A_i}{\star_i})\dsep
     x,y\in\prod_{i\in I}A_i\vdash{}\\[-2pt]
 &x\odot y\coloneqq
   \bigl\{(i,a)\in I\times\bigcup_{j\in I}A_j
     \mid a=x(i)\star_i y(i)\bigr\}.
 \end{aligned}%
}{SemigroupDirectProductGraphDef}{%
 \DeltaRow{Mengenfamilie}{(A_i)_{i\in I}}
 \DeltaRow{Mengen}{I\dsep x\dsep y\dsep i\dsep j\dsep a}
 \DeltaRow{Operationen}{\star_i\dsep\odot}
}

\FormulaThmDeltaK[Direkte Produkte tragen eine Halbgruppenoperation]{%
 (\forall i\in I\;\SemigroupAlg{A_i}{\star_i})\vdash
 \SemigroupAlg{\prod_{i\in I}A_i}{\odot}%
}{SemigroupDirectProductIsSemigroup}{%
 \DeltaRow{Mengen}{I\dsep A_i\dsep x\dsep y\dsep z\dsep i}
 \DeltaRow{Funktionen}{\star_i\dsep\odot}
}
\begin{tabproofwide}
 \proofstepwidestar[1]{%
  \forall i\in I\;\SemigroupAlg{A_i}{\star_i}}{%
  \rA}
 \proofstepwidestar[2]{%
  x\in\prod_{i\in I}A_i}{%
  \rA}
 \proofstepwidestar[3]{%
  y\in\prod_{i\in I}A_i}{%
  \rA}
 \proofstepwidestar[4]{%
  i\in I}{%
  \rA}
 \proofstepwidestar[1,2,4]{%
  x(i)\in A_i}{%
  \FormulaRefAuto{SemigroupDirectProductDef}[def]{1,2,4}}
 \proofstepwidestar[1,3,4]{%
  y(i)\in A_i}{%
  \FormulaRefAuto{SemigroupDirectProductDef}[def]{1,3,4}}
 \proofstepwidestar[1,4]{%
  \SemigroupAlg{A_i}{\star_i}}{%
  \rRE{\rUE{1},4}}
 \proofstepwidestar[1,2,3,4]{%
  x(i)\star_i y(i)\in A_i}{%
  \FormulaRefAuto{SemigroupClosureAxiom}[ax]{7,5,6}}
 \proofstepwidestar[1,2,3,4]{%
  \exists j\in I\;(x(i)\star_i y(i)\in A_j)}{%
  \rEI{\rAI{4,8}}}
 \proofstepwidestar[1,2,3,4]{%
  x(i)\star_i y(i)\in \bigcup_{i\in I}A_i}{%
  \FormulaRefAuto{SemigroupIndexedUnionNotationDef}[def]{9}}
 \proofstepwidestar[1,2,3]{%
  \forall i\in I\;(x(i)\star_i y(i)\in A_i)}{%
  \rUI{\rRI{4,8}}}
 \proofstepwidestar[1,2,3]{%
  \forall i\in I\;(x(i)\star_i y(i)\in \bigcup_{i\in I}A_i)}{%
  \rUI{\rRI{4,10}}}
 \proofstepwidestar[1,2,3]{%
  x\odot y\colon I\to\bigcup_{i\in I}A_i}{%
  \FormulaRefAuto{TypedTermGraphFunction}[thm]{12},\ \FormulaRefAuto{SemigroupDirectProductGraphDef}[def]}
 \proofstepwidestar[1,2,3]{%
  \forall i\in I\;((x\odot y)(i)=x(i)\star_i y(i))}{%
  \FormulaRefAuto{TypedTermGraphValues}[thm]{12},\ \FormulaRefAuto{SemigroupDirectProductGraphDef}[def]}
 \proofstepwidestar[1,2,3,4]{%
  (x\odot y)(i)=x(i)\star_i y(i)}{%
  \rRE{\rUE{14},4}}
 \proofstepwidestar[1,2,3,4]{%
  x(i)\star_i y(i)=(x\odot y)(i)}{%
  \FormulaRefAuto{a=b\vdash b=a}[thm]{15}}
 \proofstepwidestar[1,2,3,4]{%
  (x\odot y)(i)\in A_i}{%
  \rIE{16,8}}
 \proofstepwidestar[1,2,3]{%
  \forall i\in I\;((x\odot y)(i)\in A_i)}{%
  \rUI{\rRI{4,17}}}
 \proofstepwidestar[1,2,3]{%
  x\odot y\in\prod_{i\in I}A_i}{%
  \FormulaRefAuto{SemigroupDirectProductDef}[def]{1,13,18}}
 \proofstepwidestar[1]{%
  \forall x\in\prod_{i\in I}A_i\,\forall y\in\prod_{i\in I}A_i\;(x\odot y\in\prod_{i\in I}A_i)}{%
  \rUI{\rRI{2,\rUI{\rRI{3,19}}}}}
 \proofstepwidestar[21]{%
  x\in\prod_{i\in I}A_i}{%
  \rA}
 \proofstepwidestar[22]{%
  y\in\prod_{i\in I}A_i}{%
  \rA}
 \proofstepwidestar[23]{%
  z\in\prod_{i\in I}A_i}{%
  \rA}
 \proofstepwidestar[1,21,22]{%
  x\odot y\in\prod_{i\in I}A_i}{%
  \rRE{\rUE{\rRE{\rUE{20},21}},22}}
 \proofstepwidestar[1,22,23]{%
  y\odot z\in\prod_{i\in I}A_i}{%
  \rRE{\rUE{\rRE{\rUE{20},22}},23}}
 \proofstepwidestar[1,21,22,23]{%
  (x\odot y)\odot z\in\prod_{i\in I}A_i}{%
  \rRE{\rUE{\rRE{\rUE{20},24}},23}}
 \proofstepwidestar[1,21,22,23]{%
  x\odot(y\odot z)\in\prod_{i\in I}A_i}{%
  \rRE{\rUE{\rRE{\rUE{20},21}},25}}
 \proofstepwidestar[1,21,22,23]{%
  (x\odot y)\odot z\colon I\to\bigcup_{i\in I}A_i}{%
  \FormulaRefAuto{SemigroupDirectProductDef}[def]{1,26}}
 \proofstepwidestar[1,21,22,23]{%
  x\odot(y\odot z)\colon I\to\bigcup_{i\in I}A_i}{%
  \FormulaRefAuto{SemigroupDirectProductDef}[def]{1,27}}
 \proofstepwidestar[30]{%
  i\in I}{%
  \rA}
 \proofstepwidestar[1,21,30]{%
  x(i)\in A_i}{%
  \FormulaRefAuto{SemigroupDirectProductDef}[def]{1,21,30}}
 \proofstepwidestar[1,22,30]{%
  y(i)\in A_i}{%
  \FormulaRefAuto{SemigroupDirectProductDef}[def]{1,22,30}}
 \proofstepwidestar[1,23,30]{%
  z(i)\in A_i}{%
  \FormulaRefAuto{SemigroupDirectProductDef}[def]{1,23,30}}
 \proofstepwidestar[1,30]{%
  \SemigroupAlg{A_i}{\star_i}}{%
  \rRE{\rUE{1},30}}
 \proofstepwidestar[1,21,22,23,30]{%
  (x(i)\star_i y(i))\star_i z(i)=x(i)\star_i(y(i)\star_i z(i))}{%
  \FormulaRefAuto{SemigroupAssociativityAxiom}[ax]{34,31,32,33}}
 \proofstepwidestar[1,21,22,23,30]{%
  ((x\odot y)\odot z)(i)=(x(i)\star_i y(i))\star_i z(i)}{%
  \FormulaRefAuto{SemigroupDirectProductDef}[def]{1,21,22,23,24,30}}
 \proofstepwidestar[1,21,22,23,30]{%
  (x\odot(y\odot z))(i)=x(i)\star_i(y(i)\star_i z(i))}{%
  \FormulaRefAuto{SemigroupDirectProductDef}[def]{1,21,22,23,25,30}}
 \proofstepwidestar[1,21,22,23,30]{%
  (x(i)\star_i y(i))\star_i z(i)=((x\odot y)\odot z)(i)}{%
  \FormulaRefAuto{a=b\vdash b=a}[thm]{36}}
 \proofstepwidestar[1,21,22,23,30]{%
  x(i)\star_i(y(i)\star_i z(i))=(x\odot(y\odot z))(i)}{%
  \FormulaRefAuto{a=b\vdash b=a}[thm]{37}}
 \proofstepwidestar[1,21,22,23,30]{%
  ((x\odot y)\odot z)(i)=(x\odot(y\odot z))(i)}{%
  \rIE{38,39,35}}
 \proofstepwidestar[1,21,22,23]{%
  \forall i\in I\;(((x\odot y)\odot z)(i)=(x\odot(y\odot z))(i))}{%
  \rUI{\rRI{30,40}}}
 \proofstepwidestar[1,21,22,23]{%
  (x\odot y)\odot z=x\odot(y\odot z)}{%
  \FormulaRefAuto{FunctionExtensionality}[thm]{28,29,41}}
 \proofstepwidestar[1]{%
  \forall x\in\prod_{i\in I}A_i\,\forall y\in\prod_{i\in I}A_i\,\forall z\in\prod_{i\in I}A_i\;((x\odot y)\odot z=x\odot(y\odot z))}{%
  \rUI{\rRI{21,\rUI{\rRI{22,\rUI{\rRI{23,42}}}}}}}
 \proofstepwidestar[1]{%
  \SemigroupAlg{\prod_{i\in I}A_i}{\odot}}{%
  \FormulaRefAuto{SemigroupDef}[def]{20,43}}
\end{tabproofwide}

\FormulaThmDeltaKR[Produkte und Umindizierungen von Isomorphismen]{%
 \begin{aligned}[t]
  &(\forall i\in I\;f_i:A_i\cong B_i)\vdash{}\\[-2pt]
  &\text{(i)}\quad \prod_i f_i:\prod_i A_i\cong\prod_i B_i\\[-2pt]
  &\text{(ii)}\quad (\prod_i f_i)(x)(i)=f_i(x(i))\\[-2pt]
  &(\forall i\in I\;\SemigroupAlg{A_i}{\star_i})\dsep\sigma:J\bij I\vdash{}\\[-2pt]
  &\text{(iii)}\quad \hspace{12pt} (x\mapsto x\circ\sigma):\prod_{i\in I}A_i \cong\prod_{j\in J}A_{\sigma(j)}
 \end{aligned}%
}{%
 \begin{aligned}[t]
 &(\forall i\in I\;f_i:A_i\cong B_i)\vdash{}\\[-2pt]
 &\prod_i f_i:\prod_i A_i\cong\prod_i B_i,
       \qquad (\prod_i f_i)(x)(i)=f_i(x(i));\\[-2pt]
 &(\forall i\in I\;\SemigroupAlg{A_i}{\star_i})\dsep\sigma:J\bij I\vdash{}\\[-2pt]
 &\hspace{12pt}
       (x\mapsto x\circ\sigma):\prod_{i\in I}A_i
            \cong\prod_{j\in J}A_{\sigma(j)}.
 \end{aligned}%
}{SemigroupIsoDirectProducts}{%
 \DeltaRow{Mengen}{I\dsep J\dsep A_i\dsep B_i\dsep x\dsep y\dsep i\dsep j}
 \DeltaRow{Funktionen}{f_i\dsep\sigma\dsep\odot}
}
\Needspace{20\baselineskip}
\FormulaDefDeltaK[Graphterme der koordinatenweisen Produktabbildungen]{%
 \begin{aligned}[t]
 &\Phi_f(x)\coloneqq
   \bigl\{(i,b)\in I\times\bigcup_{j\in I}B_j
     \mid b=f_i(x(i))\bigr\},\\[-2pt]
 &\Psi_f(y)\coloneqq
   \bigl\{(i,a)\in I\times\bigcup_{j\in I}A_j
     \mid a=f_i^{-1}(y(i))\bigr\},\\[-2pt]
 &F\coloneqq
   \bigl\{(x,z)\in\prod_{i\in I}A_i\times\prod_{i\in I}B_i
     \mid z=\Phi_f(x)\bigr\},\\[-2pt]
 &H\coloneqq
   \bigl\{(y,z)\in\prod_{i\in I}B_i\times\prod_{i\in I}A_i
     \mid z=\Psi_f(y)\bigr\},\\[-2pt]
 &\prod_{i\in I}f_i\coloneqq F.
 \end{aligned}%
}{SemigroupCoordinateProductMapsDef}{%
 \DeltaRow{Mengenfamilien}{(A_i)_{i\in I}\dsep(B_i)_{i\in I}}
 \DeltaRow{Mengen}{I\dsep x\dsep y\dsep z\dsep i\dsep j\dsep a\dsep b}
 \DeltaRow{Funktionsterme}{f_i\dsep f_i^{-1}\dsep\Phi_f(x)\dsep\Psi_f(y)\dsep F\dsep H}
}

\Needspace{12\baselineskip}
\begin{tabproofsplitwide}
\proofpartwideindR[Koordinatenabbildung]{%
\begin{aligned}[t]&\forall i\in I\;(f_i:A_i\cong B_i)\vdash{}\\[-2pt]&F\colon\prod_{i\in I}A_i\to\prod_{i\in I}B_i\land{}\\[-2pt]&\forall x\in\prod_{i\in I}A_i\;\forall i\in I\;(F(x)(i)=f_i(x(i)))\end{aligned}%
}{\forall i\in I\;(f_i:A_i\cong B_i)\vdash F\colon\prod_{i\in I}A_i\to\prod_{i\in I}B_i\land \forall x\in\prod_{i\in I}A_i\;\forall i\in I\;(F(x)(i)=f_i(x(i)))}[SemigroupCoordinateProductForwardFacts]
 \proofstepwidestar[1]{%
  \forall i\in I\;(f_i:A_i\cong B_i)}{%
  \rA}
 \proofstepwidestar[2]{%
  i\in I}{%
  \rA}
 \proofstepwidestar[1,2]{%
  f_i:A_i\cong B_i}{%
  \rRE{\rUE{1},2}}
 \proofstepwidestar[1,2]{%
  \SemigroupAlg{A_i}{\star_i}}{%
  \FormulaRefAuto{SemigroupIsoSource}[thm]{3}}
 \proofstepwidestar[1,2]{%
  \SemigroupAlg{B_i}{\diamond_i}}{%
  \FormulaRefAuto{SemigroupIsoTarget}[thm]{3}}
 \proofstepwidestar[1]{%
  \forall i\in I\;\SemigroupAlg{A_i}{\star_i}}{%
  \rUI{\rRI{2,4}}}
 \proofstepwidestar[1]{%
  \forall i\in I\;\SemigroupAlg{B_i}{\diamond_i}}{%
  \rUI{\rRI{2,5}}}
 \proofstepwidestar[8]{%
  x\in\prod_{i\in I}A_i}{%
  \rA}
 \proofstepwidestar[9]{%
  i\in I}{%
  \rA}
 \proofstepwidestar[1,9]{%
  f_i:A_i\cong B_i}{%
  \rRE{\rUE{1},9}}
 \proofstepwidestar[1,9]{%
  f_i:A_i\bij B_i}{%
  \FormulaRefAuto{SemigroupIsoBijectiveAxiom}[ax]{10}}
 \proofstepwidestar[1,9]{%
  f_i:A_i\to B_i}{%
  \FormulaRefAuto{Bijektive Funktion}[def]{11}}
 \proofstepwidestar[1,8,9]{%
  x(i)\in A_i}{%
  \FormulaRefAuto{SemigroupDirectProductDef}[def]{6,8,9}}
 \proofstepwidestar[1,8,9]{%
  f_i(x(i))\in B_i}{%
  \FormulaRefAuto{F\colon A\to B\dsep x\in A\vdash F(x)\in B}[thm]{12,13}}
 \proofstepwidestar[1,8,9]{%
  \exists j\in I\;(f_i(x(i))\in B_j)}{%
  \rEI{\rAI{9,14}}}
 \proofstepwidestar[1,8,9]{%
  f_i(x(i))\in \bigcup_{i\in I}B_i}{%
  \FormulaRefAuto{SemigroupIndexedUnionNotationDef}[def]{15}}
 \proofstepwidestar[1,8]{%
  \forall i\in I\;(f_i(x(i))\in \bigcup_{i\in I}B_i)}{%
  \rUI{\rRI{9,16}}}
 \proofstepwidestar[1,8]{%
  \Phi_f(x)\colon I\to\bigcup_{i\in I}B_i}{%
  \FormulaRefAuto{TypedTermGraphFunction}[thm]{17},\ \FormulaRefAuto{SemigroupCoordinateProductMapsDef}[def]}
 \proofstepwidestar[1,8]{%
  \forall i\in I\;(\Phi_f(x)(i)=f_i(x(i)))}{%
  \FormulaRefAuto{TypedTermGraphValues}[thm]{17},\ \FormulaRefAuto{SemigroupCoordinateProductMapsDef}[def]}
 \proofstepwidestar[1,8,9]{%
  \Phi_f(x)(i)=f_i(x(i))}{%
  \rRE{\rUE{19},9}}
 \proofstepwidestar[1,8,9]{%
  f_i(x(i))=\Phi_f(x)(i)}{%
  \FormulaRefAuto{a=b\vdash b=a}[thm]{20}}
 \proofstepwidestar[1,8,9]{%
  \Phi_f(x)(i)\in B_i}{%
  \rIE{21,14}}
 \proofstepwidestar[1,8]{%
  \forall i\in I\;(\Phi_f(x)(i)\in B_i)}{%
  \rUI{\rRI{9,22}}}
 \proofstepwidestar[1,8]{%
  \Phi_f(x)\in\prod_{i\in I}B_i}{%
  \FormulaRefAuto{SemigroupDirectProductDef}[def]{7,18,23}}
 \proofstepwidestar[1]{%
  \forall x\in\prod_{i\in I}A_i\;(\Phi_f(x)\in\prod_{i\in I}B_i)}{%
  \rUI{\rRI{8,24}}}
 \proofstepwidestar[1]{%
  F\colon\prod_{i\in I}A_i\to\prod_{i\in I}B_i}{%
  \FormulaRefAuto{TypedTermGraphFunction}[thm]{25},\ \FormulaRefAuto{SemigroupCoordinateProductMapsDef}[def]}
 \proofstepwidestar[1]{%
  \forall x\in\prod_{i\in I}A_i\;(F(x)=\Phi_f(x))}{%
  \FormulaRefAuto{TypedTermGraphValues}[thm]{25},\ \FormulaRefAuto{SemigroupCoordinateProductMapsDef}[def]}
 \proofstepwidestar[1,8]{%
  F(x)=\Phi_f(x)}{%
  \rRE{\rUE{27},8}}
 \proofstepwidestar[1,8]{%
  \Phi_f(x)=F(x)}{%
  \FormulaRefAuto{a=b\vdash b=a}[thm]{28}}
 \proofstepwidestar[1,8,9]{%
  F(x)(i)=f_i(x(i))}{%
  \rIE{29,20}}
 \proofstepwidestar[1]{%
  \forall x\in\prod_{i\in I}A_i\;\forall i\in I\;(F(x)(i)=f_i(x(i)))}{%
  \rUI{\rRI{8,\rUI{\rRI{9,30}}}}}
 \proofstepwidestar[1]{%
  \begin{aligned}[t]&F\colon\prod_{i\in I}A_i\to\prod_{i\in I}B_i\land{}\\[-2pt]&\forall x\in\prod_{i\in I}A_i\;\forall i\in I\;(F(x)(i)=f_i(x(i)))\end{aligned}}{%
  \rAI{26,31}}
\closeproofpartwideindR

\proofpartwideindR[Inverse Koordinatenabbildung]{%
\begin{aligned}[t]&\forall i\in I\;(f_i:A_i\cong B_i)\vdash{}\\[-2pt]&H\colon\prod_{i\in I}B_i\to\prod_{i\in I}A_i\land{}\\[-2pt]&\forall y\in\prod_{i\in I}B_i\;\forall i\in I\;(H(y)(i)=f_i^{-1}(y(i)))\end{aligned}%
}{\forall i\in I\;(f_i:A_i\cong B_i)\vdash H\colon\prod_{i\in I}B_i\to\prod_{i\in I}A_i\land \forall y\in\prod_{i\in I}B_i\;\forall i\in I\;(H(y)(i)=f_i^{-1}(y(i)))}[SemigroupCoordinateProductInverseFacts]
 \proofstepwidestar[1]{%
  \forall i\in I\;(f_i:A_i\cong B_i)}{%
  \rA}
 \proofstepwidestar[2]{%
  i\in I}{%
  \rA}
 \proofstepwidestar[1,2]{%
  f_i:A_i\cong B_i}{%
  \rRE{\rUE{1},2}}
 \proofstepwidestar[1,2]{%
  \SemigroupAlg{B_i}{\diamond_i}}{%
  \FormulaRefAuto{SemigroupIsoTarget}[thm]{3}}
 \proofstepwidestar[1,2]{%
  \SemigroupAlg{A_i}{\star_i}}{%
  \FormulaRefAuto{SemigroupIsoSource}[thm]{3}}
 \proofstepwidestar[1]{%
  \forall i\in I\;\SemigroupAlg{B_i}{\diamond_i}}{%
  \rUI{\rRI{2,4}}}
 \proofstepwidestar[1]{%
  \forall i\in I\;\SemigroupAlg{A_i}{\star_i}}{%
  \rUI{\rRI{2,5}}}
 \proofstepwidestar[8]{%
  y\in\prod_{i\in I}B_i}{%
  \rA}
 \proofstepwidestar[9]{%
  i\in I}{%
  \rA}
 \proofstepwidestar[1,9]{%
  f_i:A_i\cong B_i}{%
  \rRE{\rUE{1},9}}
 \proofstepwidestar[1,9]{%
  f_i^{-1}:B_i\bij A_i}{%
  \FormulaRefAuto{SemigroupIsoInverseCarrier}[thm]{10}}
 \proofstepwidestar[1,9]{%
  f_i^{-1}:B_i\to A_i}{%
  \FormulaRefAuto{Bijektive Funktion}[def]{11}}
 \proofstepwidestar[1,8,9]{%
  y(i)\in B_i}{%
  \FormulaRefAuto{SemigroupDirectProductDef}[def]{6,8,9}}
 \proofstepwidestar[1,8,9]{%
  f_i^{-1}(y(i))\in A_i}{%
  \FormulaRefAuto{F\colon A\to B\dsep x\in A\vdash F(x)\in B}[thm]{12,13}}
 \proofstepwidestar[1,8,9]{%
  \exists j\in I\;(f_i^{-1}(y(i))\in A_j)}{%
  \rEI{\rAI{9,14}}}
 \proofstepwidestar[1,8,9]{%
  f_i^{-1}(y(i))\in \bigcup_{i\in I}A_i}{%
  \FormulaRefAuto{SemigroupIndexedUnionNotationDef}[def]{15}}
 \proofstepwidestar[1,8]{%
  \forall i\in I\;(f_i^{-1}(y(i))\in \bigcup_{i\in I}A_i)}{%
  \rUI{\rRI{9,16}}}
 \proofstepwidestar[1,8]{%
  \Psi_f(y)\colon I\to\bigcup_{i\in I}A_i}{%
  \FormulaRefAuto{TypedTermGraphFunction}[thm]{17},\ \FormulaRefAuto{SemigroupCoordinateProductMapsDef}[def]}
 \proofstepwidestar[1,8]{%
  \forall i\in I\;(\Psi_f(y)(i)=f_i^{-1}(y(i)))}{%
  \FormulaRefAuto{TypedTermGraphValues}[thm]{17},\ \FormulaRefAuto{SemigroupCoordinateProductMapsDef}[def]}
 \proofstepwidestar[1,8,9]{%
  \Psi_f(y)(i)=f_i^{-1}(y(i))}{%
  \rRE{\rUE{19},9}}
 \proofstepwidestar[1,8,9]{%
  f_i^{-1}(y(i))=\Psi_f(y)(i)}{%
  \FormulaRefAuto{a=b\vdash b=a}[thm]{20}}
 \proofstepwidestar[1,8,9]{%
  \Psi_f(y)(i)\in A_i}{%
  \rIE{21,14}}
 \proofstepwidestar[1,8]{%
  \forall i\in I\;(\Psi_f(y)(i)\in A_i)}{%
  \rUI{\rRI{9,22}}}
 \proofstepwidestar[1,8]{%
  \Psi_f(y)\in\prod_{i\in I}A_i}{%
  \FormulaRefAuto{SemigroupDirectProductDef}[def]{7,18,23}}
 \proofstepwidestar[1]{%
  \forall y\in\prod_{i\in I}B_i\;(\Psi_f(y)\in\prod_{i\in I}A_i)}{%
  \rUI{\rRI{8,24}}}
 \proofstepwidestar[1]{%
  H\colon\prod_{i\in I}B_i\to\prod_{i\in I}A_i}{%
  \FormulaRefAuto{TypedTermGraphFunction}[thm]{25},\ \FormulaRefAuto{SemigroupCoordinateProductMapsDef}[def]}
 \proofstepwidestar[1]{%
  \forall y\in\prod_{i\in I}B_i\;(H(y)=\Psi_f(y))}{%
  \FormulaRefAuto{TypedTermGraphValues}[thm]{25},\ \FormulaRefAuto{SemigroupCoordinateProductMapsDef}[def]}
 \proofstepwidestar[1,8]{%
  H(y)=\Psi_f(y)}{%
  \rRE{\rUE{27},8}}
 \proofstepwidestar[1,8]{%
  \Psi_f(y)=H(y)}{%
  \FormulaRefAuto{a=b\vdash b=a}[thm]{28}}
 \proofstepwidestar[1,8,9]{%
  H(y)(i)=f_i^{-1}(y(i))}{%
  \rIE{29,20}}
 \proofstepwidestar[1]{%
  \forall y\in\prod_{i\in I}B_i\;\forall i\in I\;(H(y)(i)=f_i^{-1}(y(i)))}{%
  \rUI{\rRI{8,\rUI{\rRI{9,30}}}}}
 \proofstepwidestar[1]{%
  \begin{aligned}[t]&H\colon\prod_{i\in I}B_i\to\prod_{i\in I}A_i\land{}\\[-2pt]&\forall y\in\prod_{i\in I}B_i\;\forall i\in I\;(H(y)(i)=f_i^{-1}(y(i)))\end{aligned}}{%
  \rAI{26,31}}
\closeproofpartwideindR

\proofpartwideindR[Produktoperation und inverse Produktabbildung]{%
(\forall i\in I\;f_i:A_i\cong B_i)\vdash F:\prod_iA_i\bij\prod_iB_i%
}{(\forall i\in I\;f_i:A_i\cong B_i)\vdash F:\prod_iA_i\bij\prod_iB_i}[SemigroupDirectProductCarrierBijection]
 \proofstepwidestar[1]{%
  \forall i\in I\;(f_i:A_i\cong B_i)}{%
  \rA}
 \proofstepwidestar[2]{%
  i\in I}{%
  \rA}
 \proofstepwidestar[1,2]{%
  f_i:A_i\cong B_i}{%
  \rRE{\rUE{1},2}}
 \proofstepwidestar[1,2]{%
  \SemigroupAlg{A_i}{\star_i}}{%
  \FormulaRefAuto{SemigroupIsoSource}[thm]{3}}
 \proofstepwidestar[1,2]{%
  \SemigroupAlg{B_i}{\diamond_i}}{%
  \FormulaRefAuto{SemigroupIsoTarget}[thm]{3}}
 \proofstepwidestar[1]{%
  \forall i\in I\;\SemigroupAlg{A_i}{\star_i}}{%
  \rUI{\rRI{2,4}}}
 \proofstepwidestar[1]{%
  \forall i\in I\;\SemigroupAlg{B_i}{\diamond_i}}{%
  \rUI{\rRI{2,5}}}
 \proofstepwidestar[1,2]{%
  f_i:A_i\bij B_i}{%
  \FormulaRefAuto{SemigroupIsoBijectiveAxiom}[ax]{3}}
 \proofstepwidestar[1]{%
  \forall i\in I\;(f_i:A_i\bij B_i)}{%
  \rUI{\rRI{2,8}}}
 \proofstepwidestar[1]{%
  \begin{aligned}[t]&F\colon\prod_{i\in I}A_i\to\prod_{i\in I}B_i\land{}\\[-2pt]&\forall x\in\prod_{i\in I}A_i\;\forall i\in I\;(F(x)(i)=f_i(x(i)))\end{aligned}}{%
  \FormulaRefAuto{SemigroupCoordinateProductForwardFacts}[thm]{1}}
 \proofstepwidestar[1]{%
  F\colon\prod_{i\in I}A_i\to\prod_{i\in I}B_i}{%
  \rAEa{10}}
 \proofstepwidestar[1]{%
  \forall x\in\prod_{i\in I}A_i\;\forall i\in I\;(F(x)(i)=f_i(x(i)))}{%
  \rAEb{10}}
 \proofstepwidestar[1]{%
  \begin{aligned}[t]&H\colon\prod_{i\in I}B_i\to\prod_{i\in I}A_i\land{}\\[-2pt]&\forall y\in\prod_{i\in I}B_i\;\forall i\in I\;(H(y)(i)=f_i^{-1}(y(i)))\end{aligned}}{%
  \FormulaRefAuto{SemigroupCoordinateProductInverseFacts}[thm]{1}}
 \proofstepwidestar[1]{%
  H\colon\prod_{i\in I}B_i\to\prod_{i\in I}A_i}{%
  \rAEa{13}}
 \proofstepwidestar[1]{%
  \forall y\in\prod_{i\in I}B_i\;\forall i\in I\;(H(y)(i)=f_i^{-1}(y(i)))}{%
  \rAEb{13}}
 \proofstepwidestar[16]{%
  x\in\prod_{i\in I}A_i}{%
  \rA}
 \proofstepwidestar[1,16]{%
  F(x)\in\prod_{i\in I}B_i}{%
  \FormulaRefAuto{F\colon A\to B\dsep x\in A\vdash F(x)\in B}[thm]{11,16}}
 \proofstepwidestar[1,16]{%
  H(F(x))\in\prod_{i\in I}A_i}{%
  \FormulaRefAuto{F\colon A\to B\dsep x\in A\vdash F(x)\in B}[thm]{14,17}}
 \proofstepwidestar[1,16]{%
  H(F(x))\colon I\to\bigcup_{i\in I}A_i}{%
  \FormulaRefAuto{SemigroupDirectProductDef}[def]{6,18}}
 \proofstepwidestar[1,16]{%
  x\colon I\to\bigcup_{i\in I}A_i}{%
  \FormulaRefAuto{SemigroupDirectProductDef}[def]{6,16}}
 \proofstepwidestar[21]{%
  i\in I}{%
  \rA}
 \proofstepwidestar[1,21]{%
  f_i:A_i\bij B_i}{%
  \rRE{\rUE{9},21}}
 \proofstepwidestar[1,16,21]{%
  x(i)\in A_i}{%
  \FormulaRefAuto{SemigroupDirectProductDef}[def]{6,16,21}}
 \proofstepwidestar[1,16,21]{%
  f_i^{-1}(f_i(x(i)))=x(i)}{%
  \FormulaRefAuto{InverseFunctionLeftInverse}[thm]{22,23}}
 \proofstepwidestar[1,16,21]{%
  F(x)(i)=f_i(x(i))}{%
  \rRE{\rUE{\rRE{\rUE{12},16}},21}}
 \proofstepwidestar[1,16,21]{%
  H(F(x))(i)=f_i^{-1}(F(x)(i))}{%
  \rRE{\rUE{\rRE{\rUE{15},17}},21}}
 \proofstepwidestar[1,16,21]{%
  f_i(x(i))=F(x)(i)}{%
  \FormulaRefAuto{a=b\vdash b=a}[thm]{25}}
 \proofstepwidestar[1,16,21]{%
  f_i^{-1}(F(x)(i))=H(F(x))(i)}{%
  \FormulaRefAuto{a=b\vdash b=a}[thm]{26}}
 \proofstepwidestar[1,16,21]{%
  H(F(x))(i)=x(i)}{%
  \rIE{27,28,24}}
 \proofstepwidestar[1,16]{%
  \forall i\in I\;(H(F(x))(i)=x(i))}{%
  \rUI{\rRI{21,29}}}
 \proofstepwidestar[1,16]{%
  H(F(x))=x}{%
  \FormulaRefAuto{FunctionExtensionality}[thm]{19,20,30}}
 \proofstepwidestar[1]{%
  H\circ F\colon\prod_{i\in I}A_i\to\prod_{i\in I}A_i}{%
  \FormulaRefAuto{CompositionDef}[def]{11,14}}
 \proofstepwidestar[]{%
  \Id_{\prod_{i\in I}A_i}\colon\prod_{i\in I}A_i\to\prod_{i\in I}A_i}{%
  \FormulaRefAuto{IdA}[thm]}
 \proofstepwidestar[1,16]{%
  (H\circ F)(x)=H(F(x))}{%
  \FormulaRefAuto{CompositionDef}[def]{11,14,16}}
 \proofstepwidestar[16]{%
  \Id_{\prod_{i\in I}A_i}(x)=x}{%
  \FormulaRefAuto{x\in A\vdash\Id_A(x)=x}[thm]{16}}
 \proofstepwidestar[1,16]{%
  H(F(x))=(H\circ F)(x)}{%
  \FormulaRefAuto{a=b\vdash b=a}[thm]{34}}
 \proofstepwidestar[16]{%
  x=\Id_{\prod_{i\in I}A_i}(x)}{%
  \FormulaRefAuto{a=b\vdash b=a}[thm]{35}}
 \proofstepwidestar[1,16]{%
  (H\circ F)(x)=\Id_{\prod_{i\in I}A_i}(x)}{%
  \rIE{36,37,31}}
 \proofstepwidestar[1]{%
  \forall x\in\prod_{i\in I}A_i\;((H\circ F)(x)=\Id_{\prod_{i\in I}A_i}(x))}{%
  \rUI{\rRI{16,38}}}
 \proofstepwidestar[1]{%
  H\circ F=\Id_{\prod_{i\in I}A_i}}{%
  \FormulaRefAuto{FunctionExtensionality}[thm]{32,33,39}}
 \proofstepwidestar[41]{%
  y\in\prod_{i\in I}B_i}{%
  \rA}
 \proofstepwidestar[1,41]{%
  H(y)\in\prod_{i\in I}A_i}{%
  \FormulaRefAuto{F\colon A\to B\dsep x\in A\vdash F(x)\in B}[thm]{14,41}}
 \proofstepwidestar[1,41]{%
  F(H(y))\in\prod_{i\in I}B_i}{%
  \FormulaRefAuto{F\colon A\to B\dsep x\in A\vdash F(x)\in B}[thm]{11,42}}
 \proofstepwidestar[1,41]{%
  F(H(y))\colon I\to\bigcup_{i\in I}B_i}{%
  \FormulaRefAuto{SemigroupDirectProductDef}[def]{7,43}}
 \proofstepwidestar[1,41]{%
  y\colon I\to\bigcup_{i\in I}B_i}{%
  \FormulaRefAuto{SemigroupDirectProductDef}[def]{7,41}}
 \proofstepwidestar[46]{%
  i\in I}{%
  \rA}
 \proofstepwidestar[1,46]{%
  f_i:A_i\bij B_i}{%
  \rRE{\rUE{9},46}}
 \proofstepwidestar[1,41,46]{%
  y(i)\in B_i}{%
  \FormulaRefAuto{SemigroupDirectProductDef}[def]{7,41,46}}
 \proofstepwidestar[1,41,46]{%
  f_i(f_i^{-1}(y(i)))=y(i)}{%
  \FormulaRefAuto{InverseFunctionRightInverse}[thm]{47,48}}
 \proofstepwidestar[1,41,46]{%
  H(y)(i)=f_i^{-1}(y(i))}{%
  \rRE{\rUE{\rRE{\rUE{15},41}},46}}
 \proofstepwidestar[1,41,46]{%
  F(H(y))(i)=f_i(H(y)(i))}{%
  \rRE{\rUE{\rRE{\rUE{12},42}},46}}
 \proofstepwidestar[1,41,46]{%
  f_i^{-1}(y(i))=H(y)(i)}{%
  \FormulaRefAuto{a=b\vdash b=a}[thm]{50}}
 \proofstepwidestar[1,41,46]{%
  f_i(H(y)(i))=F(H(y))(i)}{%
  \FormulaRefAuto{a=b\vdash b=a}[thm]{51}}
 \proofstepwidestar[1,41,46]{%
  F(H(y))(i)=y(i)}{%
  \rIE{52,53,49}}
 \proofstepwidestar[1,41]{%
  \forall i\in I\;(F(H(y))(i)=y(i))}{%
  \rUI{\rRI{46,54}}}
 \proofstepwidestar[1,41]{%
  F(H(y))=y}{%
  \FormulaRefAuto{FunctionExtensionality}[thm]{44,45,55}}
 \proofstepwidestar[1]{%
  F\circ H\colon\prod_{i\in I}B_i\to\prod_{i\in I}B_i}{%
  \FormulaRefAuto{CompositionDef}[def]{14,11}}
 \proofstepwidestar[]{%
  \Id_{\prod_{i\in I}B_i}\colon\prod_{i\in I}B_i\to\prod_{i\in I}B_i}{%
  \FormulaRefAuto{IdA}[thm]}
 \proofstepwidestar[1,41]{%
  (F\circ H)(y)=F(H(y))}{%
  \FormulaRefAuto{CompositionDef}[def]{14,11,41}}
 \proofstepwidestar[41]{%
  \Id_{\prod_{i\in I}B_i}(y)=y}{%
  \FormulaRefAuto{x\in A\vdash\Id_A(x)=x}[thm]{41}}
 \proofstepwidestar[1,41]{%
  F(H(y))=(F\circ H)(y)}{%
  \FormulaRefAuto{a=b\vdash b=a}[thm]{59}}
 \proofstepwidestar[41]{%
  y=\Id_{\prod_{i\in I}B_i}(y)}{%
  \FormulaRefAuto{a=b\vdash b=a}[thm]{60}}
 \proofstepwidestar[1,41]{%
  (F\circ H)(y)=\Id_{\prod_{i\in I}B_i}(y)}{%
  \rIE{61,62,56}}
 \proofstepwidestar[1]{%
  \forall y\in\prod_{i\in I}B_i\;((F\circ H)(y)=\Id_{\prod_{i\in I}B_i}(y))}{%
  \rUI{\rRI{41,63}}}
 \proofstepwidestar[1]{%
  F\circ H=\Id_{\prod_{i\in I}B_i}}{%
  \FormulaRefAuto{FunctionExtensionality}[thm]{57,58,64}}
 \proofstepwidestar[1]{%
  F:\prod_{i\in I}A_i\bij\prod_{i\in I}B_i}{%
  \FormulaRefAuto{MutuallyInverseFunctionBijection}[thm]{14,11,65,40}}
\closeproofpartwideindR
\proofpartwide{Operationserhaltung}
 \proofstepwidestar[1]{%
  \forall i\in I\;(f_i:A_i\cong B_i)}{%
  \rA}
 \proofstepwidestar[2]{%
  i\in I}{%
  \rA}
 \proofstepwidestar[1,2]{%
  f_i:A_i\cong B_i}{%
  \rRE{\rUE{1},2}}
 \proofstepwidestar[1,2]{%
  \SemigroupAlg{A_i}{\star_i}}{%
  \FormulaRefAuto{SemigroupIsoSource}[thm]{3}}
 \proofstepwidestar[1,2]{%
  \SemigroupAlg{B_i}{\diamond_i}}{%
  \FormulaRefAuto{SemigroupIsoTarget}[thm]{3}}
 \proofstepwidestar[1]{%
  \forall i\in I\;\SemigroupAlg{A_i}{\star_i}}{%
  \rUI{\rRI{2,4}}}
 \proofstepwidestar[1]{%
  \forall i\in I\;\SemigroupAlg{B_i}{\diamond_i}}{%
  \rUI{\rRI{2,5}}}
 \proofstepwidestar[1]{%
  \SemigroupAlg{\prod_{i\in I}A_i}{\odot_A}}{%
  \FormulaRefAuto{SemigroupDirectProductIsSemigroup}[thm]{6}}
 \proofstepwidestar[1]{%
  \SemigroupAlg{\prod_{i\in I}B_i}{\odot_B}}{%
  \FormulaRefAuto{SemigroupDirectProductIsSemigroup}[thm]{7}}
 \proofstepwidestar[1]{%
  F:\prod_{i\in I}A_i\bij\prod_{i\in I}B_i}{%
  \FormulaRefAuto{SemigroupDirectProductCarrierBijection}[thm]{1}}
 \proofstepwidestar[1]{%
  F:\prod_{i\in I}A_i\to\prod_{i\in I}B_i}{%
  \FormulaRefAuto{Bijektive Funktion}[def]{10}}
 \proofstepwidestar[1]{%
  \begin{aligned}[t]&F\colon\prod_{i\in I}A_i\to\prod_{i\in I}B_i\land{}\\[-2pt]&\forall x\in\prod_{i\in I}A_i\;\forall i\in I\;(F(x)(i)=f_i(x(i)))\end{aligned}}{%
  \FormulaRefAuto{SemigroupCoordinateProductForwardFacts}[thm]{1}}
 \proofstepwidestar[1]{%
  \forall x\in\prod_{i\in I}A_i\;\forall i\in I\;(F(x)(i)=f_i(x(i)))}{%
  \rAEb{12}}
 \proofstepwidestar[14]{%
  x\in\prod_{i\in I}A_i}{%
  \rA}
 \proofstepwidestar[15]{%
  y\in\prod_{i\in I}A_i}{%
  \rA}
 \proofstepwidestar[1,14,15]{%
  x\odot_A y\in\prod_{i\in I}A_i}{%
  \FormulaRefAuto{SemigroupClosureAxiom}[ax]{8,14,15}}
 \proofstepwidestar[1,14]{%
  F(x)\in\prod_{i\in I}B_i}{%
  \FormulaRefAuto{F\colon A\to B\dsep x\in A\vdash F(x)\in B}[thm]{11,14}}
 \proofstepwidestar[1,15]{%
  F(y)\in\prod_{i\in I}B_i}{%
  \FormulaRefAuto{F\colon A\to B\dsep x\in A\vdash F(x)\in B}[thm]{11,15}}
 \proofstepwidestar[1,14,15]{%
  F(x)\odot_B F(y)\in\prod_{i\in I}B_i}{%
  \FormulaRefAuto{SemigroupClosureAxiom}[ax]{9,17,18}}
 \proofstepwidestar[1,14,15]{%
  F(x\odot_A y)\in\prod_{i\in I}B_i}{%
  \FormulaRefAuto{F\colon A\to B\dsep x\in A\vdash F(x)\in B}[thm]{11,16}}
 \proofstepwidestar[1,14,15]{%
  F(x\odot_A y)\colon I\to\bigcup_{i\in I}B_i}{%
  \FormulaRefAuto{SemigroupDirectProductDef}[def]{7,20}}
 \proofstepwidestar[1,14,15]{%
  F(x)\odot_B F(y)\colon I\to\bigcup_{i\in I}B_i}{%
  \FormulaRefAuto{SemigroupDirectProductDef}[def]{7,19}}
 \proofstepwidestar[23]{%
  i\in I}{%
  \rA}
 \proofstepwidestar[1,23]{%
  f_i:A_i\cong B_i}{%
  \rRE{\rUE{1},23}}
 \proofstepwidestar[1,14,23]{%
  x(i)\in A_i}{%
  \FormulaRefAuto{SemigroupDirectProductDef}[def]{6,14,23}}
 \proofstepwidestar[1,15,23]{%
  y(i)\in A_i}{%
  \FormulaRefAuto{SemigroupDirectProductDef}[def]{6,15,23}}
 \proofstepwidestar[1,14,15,23]{%
  f_i(x(i)\star_i y(i))=f_i(x(i))\diamond_i f_i(y(i))}{%
  \FormulaRefAuto{SemigroupIsoOperation}[thm]{24,25,26}}
 \proofstepwidestar[1,14,15,23]{%
  (x\odot_A y)(i)=x(i)\star_i y(i)}{%
  \FormulaRefAuto{SemigroupDirectProductDef}[def]{6,14,15,23}}
 \proofstepwidestar[1,14,23]{%
  F(x)(i)=f_i(x(i))}{%
  \rRE{\rUE{\rRE{\rUE{13},14}},23}}
 \proofstepwidestar[1,15,23]{%
  F(y)(i)=f_i(y(i))}{%
  \rRE{\rUE{\rRE{\rUE{13},15}},23}}
 \proofstepwidestar[1,14,15,23]{%
  F(x\odot_A y)(i)=f_i((x\odot_A y)(i))}{%
  \rRE{\rUE{\rRE{\rUE{13},16}},23}}
 \proofstepwidestar[1,14,15,23]{%
  (F(x)\odot_B F(y))(i)=F(x)(i)\diamond_i F(y)(i)}{%
  \FormulaRefAuto{SemigroupDirectProductDef}[def]{7,17,18,23}}
 \proofstepwidestar[1,14,15,23]{%
  F(x\odot_A y)(i)=f_i(x(i)\star_i y(i))}{%
  \rIE{28,31}}
 \proofstepwidestar[1,14,15,23]{%
  (F(x)\odot_B F(y))(i)=f_i(x(i))\diamond_i f_i(y(i))}{%
  \rIE{29,30,32}}
 \proofstepwidestar[1,14,15,23]{%
  f_i(x(i)\star_i y(i))=F(x\odot_A y)(i)}{%
  \FormulaRefAuto{a=b\vdash b=a}[thm]{33}}
 \proofstepwidestar[1,14,15,23]{%
  f_i(x(i))\diamond_i f_i(y(i))=(F(x)\odot_B F(y))(i)}{%
  \FormulaRefAuto{a=b\vdash b=a}[thm]{34}}
 \proofstepwidestar[1,14,15,23]{%
  F(x\odot_A y)(i)=(F(x)\odot_B F(y))(i)}{%
  \rIE{35,36,27}}
 \proofstepwidestar[1,14,15]{%
  \forall i\in I\;(F(x\odot_A y)(i)=(F(x)\odot_B F(y))(i))}{%
  \rUI{\rRI{23,37}}}
 \proofstepwidestar[1,14,15]{%
  F(x\odot_A y)=F(x)\odot_B F(y)}{%
  \FormulaRefAuto{FunctionExtensionality}[thm]{21,22,38}}
 \proofstepwidestar[1]{%
  \forall x\in\prod_{i\in I}A_i\;\forall y\in\prod_{i\in I}A_i\;(F(x\odot_A y)=F(x)\odot_B F(y))}{%
  \rUI{\rRI{14,\rUI{\rRI{15,39}}}}}
 \proofstepwidestar[1]{%
  \SemigroupHom{F}{\prod_{i\in I}A_i,\odot_A}{\prod_{i\in I}B_i,\odot_B}}{%
  \FormulaRefAuto{SemigroupHomDef}[def]{8,9,11,40}}
 \proofstepwidestar[1]{%
  \forall x\in\prod_{i\in I}A_i\;\forall i\in I\;(F(x)(i)=f_i(x(i)))}{%
  \rAEb{12}}
 \proofstepwidestar[1]{%
  F:\prod_{i\in I}A_i\cong\prod_{i\in I}B_i}{%
  \FormulaRefAuto{SemigroupIsoDef}[def]{41,10}}
\closeproofpartwide
\proofpartwide{Umindizierung}
 \proofstepwidestar[1]{\forall i\in I\;\SemigroupAlg{A_i}{\star_i}}{\rA}
 \proofstepwidestar[2]{\sigma:J\bij I}{\rA}
 \proofstepwidestar[1]{\SemigroupAlg{\prod_iA_i}{\odot_A}}{\FormulaRefAuto{SemigroupDirectProductIsSemigroup}[thm]{1}}
 \proofstepwidestar[1,2]{\forall j\in J\;\SemigroupAlg{A_{\sigma(j)}}{\star_{\sigma(j)}}}{\text{\parbox[t]{\linewidth}{\raggedright 1,2; All-Elimination in 1: Aus der Abbildung in 2 folgt für jedes \(j\in J\), dass \(\sigma(j)\in I\).}}}
 \proofstepwidestar[1,2]{\SemigroupAlg{\prod_jA_{\sigma(j)}}{\odot_B}}{\FormulaRefAuto{SemigroupDirectProductIsSemigroup}[thm]{4}}
 \proofstepwidestar[2]{P_\sigma(x)=x\circ\sigma,\quad Q_\sigma(y)=y\circ\sigma^{-1}}{\text{\parbox[t]{\linewidth}{\raggedright Funktionsdefinitionen auf den jeweils angegebenen Produktträgern}}}
 \proofstepwidestar[2]{Q_\sigma(P_\sigma(x))(i)=x(\sigma(\sigma^{-1}(i)))=x(i)\quad(i\in I)}{\text{\parbox[t]{\linewidth}{\raggedright 6; rechte Umkehrgleichung von \(\sigma\)}}}
 \proofstepwidestar[2]{P_\sigma(Q_\sigma(y))(j)=y(\sigma^{-1}(\sigma(j)))=y(j)\quad(j\in J)}{\text{\parbox[t]{\linewidth}{\raggedright 6; linke Umkehrgleichung von \(\sigma\)}}}
 \proofstepwidestar[2]{P_\sigma:\prod_iA_i\bij\prod_jA_{\sigma(j)}}{\text{\parbox[t]{\linewidth}{\raggedright 7,8; Funktionsextensionalität und wechselseitige Inversen}}}
 \proofstepwidestar[10]{x\in \prod_iA_i}{\rA}
 \proofstepwidestar[11]{y\in \prod_iA_i}{\rA}
 \proofstepwidestar[10,11]{x,y\in\prod_iA_i}{\rAIStar{10,11}}
 \proofstepwidestar[2,10,11]{P_\sigma(x\odot_A y)(j)=x(\sigma(j))\star_{\sigma(j)}y(\sigma(j))\quad(j\in J)}{\FormulaRefAuto{SemigroupDirectProductDef}[def]{6,12}}
 \proofstepwidestar[2,10,11]{P_\sigma(x\odot_A y)=P_\sigma(x)\odot_B P_\sigma(y)}{\text{\parbox[t]{\linewidth}{\raggedright 13; Produktdefinition und Funktionsextensionalität}}}
 \proofstepwidestar[2]{\forall x,y\in\prod_iA_i\;(P_\sigma(x\odot_A y)=P_\sigma(x)\odot_B P_\sigma(y))}{\rUI{\rRI{10,\rUI{\rRI{11,14}}}}}
 \proofstepwidestar[2]{P_\sigma:\prod_iA_i\to\prod_jA_{\sigma(j)}}{\FormulaRefAuto{Bijektive Funktion}[def]{9}}
 \proofstepwidestar[1,2]{\SemigroupHom{P_\sigma}{\prod_iA_i,\odot_A}{\prod_jA_{\sigma(j)},\odot_B}}{\FormulaRefAuto{SemigroupHomDef}[def]{3,5,16,15}}
 \proofstepwidestar[1,2]{P_\sigma:\prod_iA_i\cong\prod_jA_{\sigma(j)}}{\FormulaRefAuto{SemigroupIsoDef}[def]{17,9}}
\closeproofpartwide
\end{tabproofsplitwide}
Damit ergeben sich insbesondere \(A\times C\cong B\times D\) aus
\(A\cong B\) und \(C\cong D\), sowie \(A^I\cong B^I\).

\FormulaDefDeltaK[Binäres direktes Produkt]{%
 (a,b)\odot(a',b')\coloneqq(a\star_A a',b\star_B b')%
}{SemigroupBinaryProductDef}{%
 \DeltaRow{Voraussetzung}{\SemigroupAlg{A}{\star_A}\dsep\SemigroupAlg{B}{\star_B}}
 \DeltaRow{Mengen}{a,a'\in A\dsep b,b'\in B}
 \DeltaRow{Träger}{A\times B}
}

\FormulaThmDeltaK[Halbgruppenstruktur des binären Produkts]{%
\SemigroupAlg{A}{\star_A}\dsep\SemigroupAlg{B}{\star_B}\vdash\SemigroupAlg{A\times B}{\odot}%
}{SemigroupBinaryProductIsSemigroup}{%
 \DeltaRow{Mengen}{A\dsep B\dsep a\dsep a'\dsep a''\dsep b\dsep b'\dsep b''}\DeltaRow{Funktionen}{\star_A\dsep\star_B\dsep\odot}
}
\begin{tabproofwide}
 \proofstepwidestar[1]{%
  \SemigroupAlg{A}{\star_A}}{%
  \rA}
 \proofstepwidestar[2]{%
  \SemigroupAlg{B}{\star_B}}{%
  \rA}
 \proofstepwidestar[3]{%
  x\in A\times B}{%
  \rA}
 \proofstepwidestar[3]{%
  \exists a\in A\,\exists b\in B\;x=(a,b)}{%
  \FormulaRefAuto{x\in A\times B\eqvdash \exists a\in A\,\exists b\in B\,x=(a,b)}[thm]{3}}
 \proofstepwidestar[5]{%
  a\in A\land\exists b\in B\;x=(a,b)}{%
  \rA}
 \proofstepwidestar[5]{%
  a\in A}{%
  \rAEa{5}}
 \proofstepwidestar[5]{%
  \exists b\in B\;x=(a,b)}{%
  \rAEb{5}}
 \proofstepwidestar[8]{%
  b\in B\land x=(a,b)}{%
  \rA}
 \proofstepwidestar[8]{%
  b\in B}{%
  \rAEa{8}}
 \proofstepwidestar[8]{%
  x=(a,b)}{%
  \rAEb{8}}
 \proofstepwidestar[8]{%
  (a,b)=x}{%
  \FormulaRefAuto{a=b\vdash b=a}[thm]{10}}
 \proofstepwidestar[12]{%
  y\in A\times B}{%
  \rA}
 \proofstepwidestar[12]{%
  \exists a'\in A\,\exists b'\in B\;y=(a',b')}{%
  \FormulaRefAuto{x\in A\times B\eqvdash \exists a\in A\,\exists b\in B\,x=(a,b)}[thm]{12}}
 \proofstepwidestar[14]{%
  a'\in A\land\exists b'\in B\;y=(a',b')}{%
  \rA}
 \proofstepwidestar[14]{%
  a'\in A}{%
  \rAEa{14}}
 \proofstepwidestar[14]{%
  \exists b'\in B\;y=(a',b')}{%
  \rAEb{14}}
 \proofstepwidestar[17]{%
  b'\in B\land y=(a',b')}{%
  \rA}
 \proofstepwidestar[17]{%
  b'\in B}{%
  \rAEa{17}}
 \proofstepwidestar[17]{%
  y=(a',b')}{%
  \rAEb{17}}
 \proofstepwidestar[17]{%
  (a',b')=y}{%
  \FormulaRefAuto{a=b\vdash b=a}[thm]{19}}
 \proofstepwidestar[1,5,14]{%
  a\star_A a'\in A}{%
  \FormulaRefAuto{SemigroupClosureAxiom}[ax]{1,6,15}}
 \proofstepwidestar[2,8,17]{%
  b\star_B b'\in B}{%
  \FormulaRefAuto{SemigroupClosureAxiom}[ax]{2,9,18}}
 \proofstepwidestar[1,2,5,8,14,17]{%
  (a\star_A a',b\star_B b')\in A\times B}{%
  \FormulaRefAuto{a\in A\dsep b\in B\vdash(a,b)\in A\times B}[thm]{21,22}}
 \proofstepwidestar[1,2,5,8,14,17]{%
  (a,b)\odot(a',b')=(a\star_A a',b\star_B b')}{%
  \FormulaRefAuto{SemigroupBinaryProductDef}[def]{1,2,6,9,15,18}}
 \proofstepwidestar[1,2,5,8,14,17]{%
  (a\star_A a',b\star_B b')=(a,b)\odot(a',b')}{%
  \FormulaRefAuto{a=b\vdash b=a}[thm]{24}}
 \proofstepwidestar[1,2,5,8,14,17]{%
  (a,b)\odot(a',b')\in A\times B}{%
  \rIE{25,23}}
 \proofstepwidestar[1,2,5,8,14,17]{%
  x\odot y\in A\times B}{%
  \rIE{11,20,26}}
 \proofstepwidestar[1,2,5,8,14]{%
  x\odot y\in A\times B}{%
  \rEE{16,17,27}}
 \proofstepwidestar[1,2,5,8,12]{%
  x\odot y\in A\times B}{%
  \rEE{13,14,28}}
 \proofstepwidestar[1,2,5,12]{%
  x\odot y\in A\times B}{%
  \rEE{7,8,29}}
 \proofstepwidestar[1,2,12,3]{%
  x\odot y\in A\times B}{%
  \rEE{4,5,30}}
 \proofstepwidestar[1,2]{%
  \forall x\in A\times B\,\forall y\in A\times B\;(x\odot y\in A\times B)}{%
  \rUI{\rRI{3,\rUI{\rRI{12,31}}}}}
 \proofstepwidestar[33]{%
  x\in A\times B}{%
  \rA}
 \proofstepwidestar[33]{%
  \exists a\in A\,\exists b\in B\;x=(a,b)}{%
  \FormulaRefAuto{x\in A\times B\eqvdash \exists a\in A\,\exists b\in B\,x=(a,b)}[thm]{33}}
 \proofstepwidestar[35]{%
  a\in A\land\exists b\in B\;x=(a,b)}{%
  \rA}
 \proofstepwidestar[35]{%
  a\in A}{%
  \rAEa{35}}
 \proofstepwidestar[35]{%
  \exists b\in B\;x=(a,b)}{%
  \rAEb{35}}
 \proofstepwidestar[38]{%
  b\in B\land x=(a,b)}{%
  \rA}
 \proofstepwidestar[38]{%
  b\in B}{%
  \rAEa{38}}
 \proofstepwidestar[38]{%
  x=(a,b)}{%
  \rAEb{38}}
 \proofstepwidestar[38]{%
  (a,b)=x}{%
  \FormulaRefAuto{a=b\vdash b=a}[thm]{40}}
 \proofstepwidestar[42]{%
  y\in A\times B}{%
  \rA}
 \proofstepwidestar[42]{%
  \exists a'\in A\,\exists b'\in B\;y=(a',b')}{%
  \FormulaRefAuto{x\in A\times B\eqvdash \exists a\in A\,\exists b\in B\,x=(a,b)}[thm]{42}}
 \proofstepwidestar[44]{%
  a'\in A\land\exists b'\in B\;y=(a',b')}{%
  \rA}
 \proofstepwidestar[44]{%
  a'\in A}{%
  \rAEa{44}}
 \proofstepwidestar[44]{%
  \exists b'\in B\;y=(a',b')}{%
  \rAEb{44}}
 \proofstepwidestar[47]{%
  b'\in B\land y=(a',b')}{%
  \rA}
 \proofstepwidestar[47]{%
  b'\in B}{%
  \rAEa{47}}
 \proofstepwidestar[47]{%
  y=(a',b')}{%
  \rAEb{47}}
 \proofstepwidestar[47]{%
  (a',b')=y}{%
  \FormulaRefAuto{a=b\vdash b=a}[thm]{49}}
 \proofstepwidestar[51]{%
  z\in A\times B}{%
  \rA}
 \proofstepwidestar[51]{%
  \exists a''\in A\,\exists b''\in B\;z=(a'',b'')}{%
  \FormulaRefAuto{x\in A\times B\eqvdash \exists a\in A\,\exists b\in B\,x=(a,b)}[thm]{51}}
 \proofstepwidestar[53]{%
  a''\in A\land\exists b''\in B\;z=(a'',b'')}{%
  \rA}
 \proofstepwidestar[53]{%
  a''\in A}{%
  \rAEa{53}}
 \proofstepwidestar[53]{%
  \exists b''\in B\;z=(a'',b'')}{%
  \rAEb{53}}
 \proofstepwidestar[56]{%
  b''\in B\land z=(a'',b'')}{%
  \rA}
 \proofstepwidestar[56]{%
  b''\in B}{%
  \rAEa{56}}
 \proofstepwidestar[56]{%
  z=(a'',b'')}{%
  \rAEb{56}}
 \proofstepwidestar[56]{%
  (a'',b'')=z}{%
  \FormulaRefAuto{a=b\vdash b=a}[thm]{58}}
 \proofstepwidestar[1,35,44]{%
  a\star_A a'\in A}{%
  \FormulaRefAuto{SemigroupClosureAxiom}[ax]{1,36,45}}
 \proofstepwidestar[1,44,53]{%
  a'\star_A a''\in A}{%
  \FormulaRefAuto{SemigroupClosureAxiom}[ax]{1,45,54}}
 \proofstepwidestar[2,38,47]{%
  b\star_B b'\in B}{%
  \FormulaRefAuto{SemigroupClosureAxiom}[ax]{2,39,48}}
 \proofstepwidestar[2,47,56]{%
  b'\star_B b''\in B}{%
  \FormulaRefAuto{SemigroupClosureAxiom}[ax]{2,48,57}}
 \proofstepwidestar[1,35,44,53]{%
  (a\star_A a')\star_A a''=a\star_A(a'\star_A a'')}{%
  \FormulaRefAuto{SemigroupAssociativityAxiom}[ax]{1,36,45,54}}
 \proofstepwidestar[2,38,47,56]{%
  (b\star_B b')\star_B b''=b\star_B(b'\star_B b'')}{%
  \FormulaRefAuto{SemigroupAssociativityAxiom}[ax]{2,39,48,57}}
 \proofstepwidestar[1,2,35,38,44,47,53,56]{%
  \begin{aligned}[t]&((a,b)\odot(a',b'))\odot(a'',b'')\\[-2pt]&\quad=((a\star_A a')\star_A a'',(b\star_B b')\star_B b'')\end{aligned}}{%
  \FormulaRefAuto{SemigroupBinaryProductDef}[def]{1,2,36,39,45,48,54,57,60,62}}
 \proofstepwidestar[1,2,35,38,44,47,53,56]{%
  \begin{aligned}[t]&(a,b)\odot((a',b')\odot(a'',b''))\\[-2pt]&\quad=(a\star_A(a'\star_A a''),b\star_B(b'\star_B b''))\end{aligned}}{%
  \FormulaRefAuto{SemigroupBinaryProductDef}[def]{1,2,36,39,45,48,54,57,61,63}}
 \proofstepwidestar[1,2,35,38,44,47,53,56]{%
  \begin{aligned}[t]&((a,b)\odot(a',b'))\odot(a'',b'')\\[-2pt]&\quad=(a\star_A(a'\star_A a''),b\star_B(b'\star_B b''))\end{aligned}}{%
  \rIE{64,65,66}}
 \proofstepwidestar[1,2,35,38,44,47,53,56]{%
  \begin{aligned}[t]&(a\star_A(a'\star_A a''),b\star_B(b'\star_B b''))\\[-2pt]&\quad=(a,b)\odot((a',b')\odot(a'',b''))\end{aligned}}{%
  \FormulaRefAuto{a=b\vdash b=a}[thm]{67}}
 \proofstepwidestar[1,2,35,38,44,47,53,56]{%
  \begin{aligned}[t]&((a,b)\odot(a',b'))\odot(a'',b'')\\[-2pt]&\quad=(a,b)\odot((a',b')\odot(a'',b''))\end{aligned}}{%
  \rIE{69,68}}
 \proofstepwidestar[1,2,35,38,44,47,53,56]{%
  (x\odot y)\odot z=x\odot(y\odot z)}{%
  \rIE{41,50,59,70}}
 \proofstepwidestar[1,2,35,38,44,47,53]{%
  (x\odot y)\odot z=x\odot(y\odot z)}{%
  \rEE{55,56,71}}
 \proofstepwidestar[1,2,35,38,44,47,51]{%
  (x\odot y)\odot z=x\odot(y\odot z)}{%
  \rEE{52,53,72}}
 \proofstepwidestar[1,2,35,38,44,51]{%
  (x\odot y)\odot z=x\odot(y\odot z)}{%
  \rEE{46,47,73}}
 \proofstepwidestar[1,2,35,38,51,42]{%
  (x\odot y)\odot z=x\odot(y\odot z)}{%
  \rEE{43,44,74}}
 \proofstepwidestar[1,2,35,51,42]{%
  (x\odot y)\odot z=x\odot(y\odot z)}{%
  \rEE{37,38,75}}
 \proofstepwidestar[1,2,51,42,33]{%
  (x\odot y)\odot z=x\odot(y\odot z)}{%
  \rEE{34,35,76}}
 \proofstepwidestar[1,2]{%
  \begin{aligned}[t]&\forall x\in A\times B\,\forall y\in A\times B\,\forall z\in A\times B\;\\[-2pt]&\quad((x\odot y)\odot z=x\odot(y\odot z))\end{aligned}}{%
  \rUI{\rRI{33,\rUI{\rRI{42,\rUI{\rRI{51,77}}}}}}}
 \proofstepwidestar[1,2]{%
  \SemigroupAlg{A\times B}{\odot}}{%
  \FormulaRefAuto{SemigroupDef}[def]{32,78}}
\end{tabproofwide}

\FormulaThmDeltaK[Halbgruppenstruktur kartesischer Potenzen]{%
\SemigroupAlg{A}{\star}\vdash\SemigroupAlg{A^I}{\odot}%
}{SemigroupCartesianPowerIsSemigroup}{%
 \DeltaRow{Mengen}{A\dsep I}\DeltaRow{Funktionen}{\star\dsep\odot}
}
\begin{tabproofwide}
 \proofstepwidestar[1]{\SemigroupAlg{A}{\star}}{\rA}
 \proofstepwidestar[1]{\forall i\in I\;\SemigroupAlg{A}{\star}}{\rUI{1}}
 \proofstepwidestar[1]{\SemigroupAlg{\prod_{i\in I}A}{\odot}}{\FormulaRefAuto{SemigroupDirectProductIsSemigroup}[thm]{2}}
 \proofstepwidestar[1]{\SemigroupAlg{A^I}{\odot}}{\text{\parbox[t]{\linewidth}{\raggedright \FormulaRefAuto{SemigroupDirectProductDef}[def]{3}; konstante Familie.}}}
\end{tabproofwide}

\FormulaThmDeltaK[Isomorphiekriterium mit angegebener Umkehrabbildung]{%
\begin{aligned}[t]&\SemigroupAlg{S}{\cdot_S}\dsep\SemigroupAlg{T}{\cdot_T}\dsep F:S\to T\dsep G:T\to S,\\&G\circ F=\Id_S\dsep F\circ G=\Id_T,\\&\forall x,y\in S\;(F(x\cdot_S y)=F(x)\cdot_T F(y))\vdash\SemigroupIso{F}{S,\cdot_S}{T,\cdot_T}\end{aligned}%
}{SemigroupExplicitInverseIsomorphismCriterion}{%
 \DeltaRow{Mengen}{S\dsep T\dsep x\dsep y}\DeltaRow{Funktionen}{F\dsep G\dsep\cdot_S\dsep\cdot_T}
}
\begin{tabproofwide}
 \proofstepwidestar[1]{\SemigroupAlg{S}{\cdot_S}}{\rA}
 \proofstepwidestar[2]{\SemigroupAlg{T}{\cdot_T}}{\rA}
 \proofstepwidestar[3]{F:S\to T}{\rA}
 \proofstepwidestar[4]{G:T\to S}{\rA}
 \proofstepwidestar[5]{G\circ F=\Id_S}{\rA}
 \proofstepwidestar[6]{F\circ G=\Id_T}{\rA}
 \proofstepwidestar[7]{\forall x,y\in S\;(F(x\cdot_S y)=F(x)\cdot_T F(y))}{\rA}
 \proofstepwidestar[3,4,5,6]{F:S\bij T}{\FormulaRefAuto{MutuallyInverseFunctionBijection}[thm]{4,3,6,5}}
 \proofstepwidestar[1,2,3,7]{\SemigroupHom{F}{S,\cdot_S}{T,\cdot_T}}{\FormulaRefAuto{SemigroupHomDef}[def]{1,2,3,7}}
 \proofstepwidestar[1,2,3,4,5,6,7]{\SemigroupIso{F}{S,\cdot_S}{T,\cdot_T}}{\FormulaRefAuto{SemigroupIsoDef}[def]{9,8}}
\end{tabproofwide}

\FormulaThmDeltaKR[Kanonische Isomorphismen direkter Produkte]{%
 \begin{aligned}[t]
  &\SemigroupAlg{A}{\star_A}\dsep\SemigroupAlg{B}{\star_B} \dsep\SemigroupAlg{C}{\star_C}\vdash{}\\[-2pt]
  &\text{(i)}\quad A\times B\cong B\times A\\[-2pt]
  &\text{(ii)}\quad (A\times B)\times C\cong A\times(B\times C)\\[-2pt]
  &\text{(iii)}\quad A\times\{*\}\cong A\\[-2pt]
  &\text{(iv)}\quad A^{I\times J}\cong(A^J)^I\\[-2pt]
  &\text{(v)}\quad (A\times B)^I\cong A^I\times B^I
 \end{aligned}%
}{%
 \begin{aligned}[t]
 &\SemigroupAlg{A}{\star_A}\dsep\SemigroupAlg{B}{\star_B}
       \dsep\SemigroupAlg{C}{\star_C}\vdash{}\\[-2pt]
 &A\times B\cong B\times A,\qquad
       (A\times B)\times C\cong A\times(B\times C),\\[-2pt]
 &A\times\{*\}\cong A,\qquad
       A^{I\times J}\cong(A^J)^I,\qquad
       (A\times B)^I\cong A^I\times B^I.
 \end{aligned}%
}{SemigroupCanonicalProductIsomorphisms}{%
 \DeltaRow{Mengen}{A\dsep B\dsep C\dsep I\dsep J\dsep a\dsep b\dsep c}
 \DeltaRow{Operationen}{\text{\parbox[t]{.60\linewidth}{Alle Produkte und Potenzen tragen die koordinatenweisen Operationen.}}}
 \DeltaRow{Einpunktträger}{*\cdot *\coloneqq *}
}
\begin{tabproofsplitwide}
\proofpartwide{Vertauschung}
 \proofstepwidestar[1]{\SemigroupAlg{A}{\star_A}}{\rA}
 \proofstepwidestar[2]{\SemigroupAlg{B}{\star_B}}{\rA}
 \proofstepwidestar[1,2]{\SemigroupAlg{A\times B}{\odot}}{\FormulaRefAuto{SemigroupBinaryProductIsSemigroup}[thm]{1,2}}
 \proofstepwidestar[1,2]{\SemigroupAlg{B\times A}{\odot}}{\FormulaRefAuto{SemigroupBinaryProductIsSemigroup}[thm]{2,1}}
 \proofstepwidestar[1,2]{S:A\times B\to B\times A\,,\quad T:B\times A\to A\times B}{\text{\parbox[t]{\linewidth}{\raggedright Definition durch \(S(a,b)=(b,a)\), \(T(b,a)=(a,b)\).}}}
 \proofstepwidestar[1,2]{\forall a,a'\in A\;\forall b,b'\in B\;(S((a,b)(a',b'))=(bb',aa')=S(a,b)S(a',b'))}{\text{\parbox[t]{\linewidth}{\raggedright \FormulaRefAuto{SemigroupBinaryProductDef}[def]; Definition von \(S\).}}}
 \proofstepwidestar[1,2]{T\circ S=\Id_{A\times B}}{\text{\parbox[t]{\linewidth}{\raggedright Auswertung der definierten Abbildungen auf jedem Element; Funktionsextensionalität}}}
 \proofstepwidestar[1,2]{S\circ T=\Id_{B\times A}}{\text{\parbox[t]{\linewidth}{\raggedright Auswertung der definierten Abbildungen auf jedem Element; Funktionsextensionalität}}}
 \proofstepwidestar[1,2]{\forall x,y\in A\times B\;(S(xy)=S(x)S(y))}{\text{\parbox[t]{\linewidth}{\raggedright Aus 6 durch Komponentenzerlegung und Funktionsextensionalität}}}
 \proofstepwidestar[1,2]{S:A\times B\cong B\times A}{\FormulaRefAuto{SemigroupExplicitInverseIsomorphismCriterion}[thm]{3,4,5,7,8,9}}
\closeproofpartwide
\Needspace{5\baselineskip}
\proofpartwide{Umklammerung}
 \proofstepwidestar[1]{\SemigroupAlg{A}{\star_A}}{\rA}
 \proofstepwidestar[2]{\SemigroupAlg{B}{\star_B}}{\rA}
 \proofstepwidestar[3]{\SemigroupAlg{C}{\star_C}}{\rA}
 \proofstepwidestar[1,2]{\SemigroupAlg{A\times B}{\odot}}{\FormulaRefAuto{SemigroupBinaryProductIsSemigroup}[thm]{1,2}}
 \proofstepwidestar[2,3]{\SemigroupAlg{B\times C}{\odot}}{\FormulaRefAuto{SemigroupBinaryProductIsSemigroup}[thm]{2,3}}
 \proofstepwidestar[1,2,3]{\SemigroupAlg{(A\times B)\times C}{\odot}}{\FormulaRefAuto{SemigroupBinaryProductIsSemigroup}[thm]{4,3}}
 \proofstepwidestar[1,2,3]{\SemigroupAlg{A\times(B\times C)}{\odot}}{\FormulaRefAuto{SemigroupBinaryProductIsSemigroup}[thm]{1,5}}
 \proofstepwidestar[1,2,3]{R:(A\times B)\times C\to A\times(B\times C)\,,\quad Q:A\times(B\times C)\to(A\times B)\times C}{\text{\parbox[t]{\linewidth}{\raggedright Definition durch \(R((a,b),c)=(a,(b,c))\), \(Q(a,(b,c))=((a,b),c)\).}}}
 \proofstepwidestar[1,2,3]{\begin{aligned}[t]&\forall a,a'\in A\;\forall b,b'\in B\;\forall c,c'\in C\\&\quad(R(((a,b),c)((a',b'),c'))\\&\qquad=(aa',(bb',cc'))\\&\qquad=R((a,b),c)R((a',b'),c'))\end{aligned}}{\text{\parbox[t]{\linewidth}{\raggedright \FormulaRefAuto{SemigroupBinaryProductDef}[def]; Definition von \(R\).}}}
 \proofstepwidestar[1,2,3]{Q\circ R=\Id_{(A\times B)\times C}}{\text{\parbox[t]{\linewidth}{\raggedright Auswertung der definierten Abbildungen auf jedem Element; Funktionsextensionalität}}}
 \proofstepwidestar[1,2,3]{R\circ Q=\Id_{A\times(B\times C)}}{\text{\parbox[t]{\linewidth}{\raggedright Auswertung der definierten Abbildungen auf jedem Element; Funktionsextensionalität}}}
 \proofstepwidestar[1,2,3]{\forall x,y\in (A\times B)\times C\;(R(xy)=R(x)R(y))}{\text{\parbox[t]{\linewidth}{\raggedright Aus 9 durch Komponentenzerlegung und Funktionsextensionalität}}}
 \proofstepwidestar[1,2,3]{R:(A\times B)\times C\cong A\times(B\times C)}{\FormulaRefAuto{SemigroupExplicitInverseIsomorphismCriterion}[thm]{6,7,8,10,11,12}}
\closeproofpartwide
\proofpartwide{Einpunktfaktor}
 \proofstepwidestar[1]{\SemigroupAlg{A}{\star_A}}{\rA}
 \proofstepwidestar[1]{\SemigroupAlg{\{*\}}{\cdot}}{\text{\parbox[t]{\linewidth}{\raggedright \FormulaRefAuto{SemigroupDef}[def]; einziges Produkt \(*\cdot *=*\), beide Dreierprodukte sind \(*\).}}}
 \proofstepwidestar[1]{\SemigroupAlg{A\times\{*\}}{\odot}}{\FormulaRefAuto{SemigroupBinaryProductIsSemigroup}[thm]{1,2}}
 \proofstepwidestar[1]{\SemigroupAlg{A}{\star_A}}{\rAEa{\rAI{1,1}}}
 \proofstepwidestar[1]{P:A\times\{*\}\to A\,,\quad E:A\to A\times\{*\}}{\text{\parbox[t]{\linewidth}{\raggedright Definition durch \(P(a,*)=a\), \(E(a)=(a,*)\).}}}
 \proofstepwidestar[1]{\forall a,b\in A\;(P((a,*)(b,*))=ab=P(a,*)P(b,*))}{\text{\parbox[t]{\linewidth}{\raggedright \FormulaRefAuto{SemigroupBinaryProductDef}[def]; Definition von \(P\).}}}
 \proofstepwidestar[1]{E\circ P=\Id_{A\times\{*\}}}{\text{\parbox[t]{\linewidth}{\raggedright Auswertung der definierten Abbildungen auf jedem Element; Funktionsextensionalität}}}
 \proofstepwidestar[1]{P\circ E=\Id_{A}}{\text{\parbox[t]{\linewidth}{\raggedright Auswertung der definierten Abbildungen auf jedem Element; Funktionsextensionalität}}}
 \proofstepwidestar[1]{\forall x,y\in A\times\{*\}\;(P(xy)=P(x)P(y))}{\text{\parbox[t]{\linewidth}{\raggedright Aus 6 durch Komponentenzerlegung und Funktionsextensionalität}}}
 \proofstepwidestar[1]{P:A\times\{*\}\cong A}{\FormulaRefAuto{SemigroupExplicitInverseIsomorphismCriterion}[thm]{3,4,5,7,8,9}}
\closeproofpartwide
\proofpartwide{Verschachtelte kartesische Potenzen}
 \proofstepwidestar[1]{\SemigroupAlg{A}{\star_A}}{\rA}
 \proofstepwidestar[1]{\SemigroupAlg{A^{I\times J}}{\odot}}{\FormulaRefAuto{SemigroupCartesianPowerIsSemigroup}[thm]{1}}
 \proofstepwidestar[1]{\SemigroupAlg{A^J}{\odot}}{\FormulaRefAuto{SemigroupCartesianPowerIsSemigroup}[thm]{1}}
 \proofstepwidestar[1]{\SemigroupAlg{(A^J)^I}{\odot}}{\FormulaRefAuto{SemigroupCartesianPowerIsSemigroup}[thm]{3}}
 \proofstepwidestar[1]{K:A^{I\times J}\to(A^J)^I\,,\quad L:(A^J)^I\to A^{I\times J}}{\text{\parbox[t]{\linewidth}{\raggedright Definition durch \(K(u)(i)(j)=u(i,j)\), \(L(v)(i,j)=v(i)(j)\); jeder Wert liegt in \(A\).}}}
 \proofstepwidestar[1]{\forall u,w\in A^{I\times J}\;\forall i\in I\;\forall j\in J\;(K(uw)(i)(j)=u(i,j)w(i,j)=(K(u)K(w))(i)(j))}{\text{\parbox[t]{\linewidth}{\raggedright \FormulaRefAuto{SemigroupDirectProductDef}[def]; Definition von \(K\).}}}
 \proofstepwidestar[1]{L\circ K=\Id_{A^{I\times J}}}{\text{\parbox[t]{\linewidth}{\raggedright Auswertung der definierten Abbildungen auf jedem Element; Funktionsextensionalität}}}
 \proofstepwidestar[1]{K\circ L=\Id_{(A^J)^I}}{\text{\parbox[t]{\linewidth}{\raggedright Auswertung der definierten Abbildungen auf jedem Element; Funktionsextensionalität}}}
 \proofstepwidestar[1]{\forall x,y\in A^{I\times J}\;(K(xy)=K(x)K(y))}{\text{\parbox[t]{\linewidth}{\raggedright Aus 6 durch Komponentenzerlegung und Funktionsextensionalität}}}
 \proofstepwidestar[1]{K:A^{I\times J}\cong (A^J)^I}{\FormulaRefAuto{SemigroupExplicitInverseIsomorphismCriterion}[thm]{2,4,5,7,8,9}}
\closeproofpartwide
\proofpartwide{Aufteilung eines Paars von Koordinaten}
 \proofstepwidestar[1]{\SemigroupAlg{A}{\star_A}}{\rA}
 \proofstepwidestar[2]{\SemigroupAlg{B}{\star_B}}{\rA}
 \proofstepwidestar[1,2]{\SemigroupAlg{A\times B}{\odot}}{\FormulaRefAuto{SemigroupBinaryProductIsSemigroup}[thm]{1,2}}
 \proofstepwidestar[1,2]{\SemigroupAlg{(A\times B)^I}{\odot}}{\FormulaRefAuto{SemigroupCartesianPowerIsSemigroup}[thm]{3}}
 \proofstepwidestar[1]{\SemigroupAlg{A^I}{\odot}}{\FormulaRefAuto{SemigroupCartesianPowerIsSemigroup}[thm]{1}}
 \proofstepwidestar[2]{\SemigroupAlg{B^I}{\odot}}{\FormulaRefAuto{SemigroupCartesianPowerIsSemigroup}[thm]{2}}
 \proofstepwidestar[1,2]{\SemigroupAlg{A^I\times B^I}{\odot}}{\FormulaRefAuto{SemigroupBinaryProductIsSemigroup}[thm]{5,6}}
 \proofstepwidestar[1,2]{D:(A\times B)^I\to A^I\times B^I\,,\quad E:A^I\times B^I\to(A\times B)^I}{\text{\parbox[t]{\linewidth}{\raggedright Definition durch \(D(u)=(\pi_1\circ u,\pi_2\circ u)\), \(E(v,w)(i)=(v(i),w(i))\).}}}
 \proofstepwidestar[1,2]{\forall u,u'\in(A\times B)^I\;\forall i\in I\;\forall k\in\{1,2\}\;(\pi_k((uu')(i))=\pi_k(u(i))\pi_k(u'(i)))}{\text{\parbox[t]{\linewidth}{\raggedright \FormulaRefAuto{SemigroupBinaryProductDef}[def]; koordinatenweise Operation.}}}
 \proofstepwidestar[1,2]{E\circ D=\Id_{(A\times B)^I}}{\text{\parbox[t]{\linewidth}{\raggedright Auswertung der definierten Abbildungen auf jedem Element; Funktionsextensionalität}}}
 \proofstepwidestar[1,2]{D\circ E=\Id_{A^I\times B^I}}{\text{\parbox[t]{\linewidth}{\raggedright Auswertung der definierten Abbildungen auf jedem Element; Funktionsextensionalität}}}
 \proofstepwidestar[1,2]{\forall x,y\in (A\times B)^I\;(D(xy)=D(x)D(y))}{\text{\parbox[t]{\linewidth}{\raggedright Aus 9 durch Komponentenzerlegung und Funktionsextensionalität}}}
 \proofstepwidestar[1,2]{D:(A\times B)^I\cong A^I\times B^I}{\FormulaRefAuto{SemigroupExplicitInverseIsomorphismCriterion}[thm]{4,7,8,10,11,12}}
\closeproofpartwide
\end{tabproofsplitwide}

\FormulaDefDeltaK[Gegenhalbgruppe]{%
 (A,\star)^{\mathrm{op}}\coloneqq(A,\star^{\mathrm{op}}),
 \qquad x\star^{\mathrm{op}}y\coloneqq y\star x%
}{SemigroupOppositeDef}{%
 \DeltaRow{Mengen}{A\dsep x\dsep y}
 \DeltaRow{Voraussetzung}{\SemigroupAlg{A}{\star}}
}
\FormulaThmDeltaK[Isomorphismen der Gegenhalbgruppen]{%
 f:A\cong B\vdash
 \SemigroupIso{f}{A,\star^{\mathrm{op}}}{B,\diamond^{\mathrm{op}}}%
}{SemigroupIsoOpposite}{%
 \DeltaRow{Mengen}{A\dsep B\dsep x\dsep y\dsep z}
 \DeltaRow{Funktionen}{f\dsep\star\dsep\diamond}
}
\begin{tabproofsplitwide}
\proofpartwideindR[Die umgekehrte Operation]{%
\SemigroupAlg{A}{\star}\vdash\SemigroupAlg{A}{\star^{\mathrm{op}}}%
}{\SemigroupAlg{A}{\star}\vdash\SemigroupAlg{A}{\star^{\mathrm{op}}}}[SemigroupOppositeIsSemigroup]
 \proofstepwidestar[1]{\SemigroupAlg{A}{\star}}{\rA}
 \proofstepwidestar[2]{x\in A}{\rA}
 \proofstepwidestar[3]{y\in A}{\rA}
 \proofstepwidestar[2,3]{x,y\in A}{\rAIStar{2,3}}
 \proofstepwidestar[1,2,3]{y\star x\in A}{\FormulaRefAuto{SemigroupClosureAxiom}[ax]{1,4}}
 \proofstepwidestar[1,2,3]{x\star^{\mathrm{op}}y\in A}{\FormulaRefAuto{SemigroupOppositeDef}[def]{5}}
 \proofstepwidestar[1]{\forall x,y\in A\;(x\star^{\mathrm{op}}y\in A)}{\rUI{\rRI{2,\rUI{\rRI{3,6}}}}}
 \proofstepwidestar[8]{x\in A}{\rA}
 \proofstepwidestar[9]{y\in A}{\rA}
 \proofstepwidestar[10]{z\in A}{\rA}
 \proofstepwidestar[8,9,10]{x,y,z\in A}{\rAIStar{8,9,10}}
 \proofstepwidestar[1,8,9,10]{(x\star^{\mathrm{op}}y)\star^{\mathrm{op}}z=z\star(y\star x)}{\FormulaRefAuto{SemigroupOppositeDef}[def]{1,11}}
 \proofstepwidestar[1,8,9,10]{(z\star y)\star x=z\star(y\star x)}{\FormulaRefAuto{SemigroupAssociativityAxiom}[ax]{1,11}}
 \proofstepwidestar[1,8,9,10]{z\star(y\star x)=(z\star y)\star x}{\FormulaRefAuto{a=b\vdash b=a}[thm]{13}}
 \proofstepwidestar[1,8,9,10]{(z\star y)\star x=x\star^{\mathrm{op}}(y\star^{\mathrm{op}}z)}{\FormulaRefAuto{SemigroupOppositeDef}[def]{1,11}}
 \proofstepwidestar[1,8,9,10]{(x\star^{\mathrm{op}}y)\star^{\mathrm{op}}z=x\star^{\mathrm{op}}(y\star^{\mathrm{op}}z)}{\rChain{12,14,15}}
 \proofstepwidestar[1]{\forall x,y,z\in A\;((x\star^{\mathrm{op}}y)\star^{\mathrm{op}}z=x\star^{\mathrm{op}}(y\star^{\mathrm{op}}z))}{\rUI{\rRI{8,\rUI{\rRI{9,\rUI{\rRI{10,16}}}}}}}
 \proofstepwidestar[1]{\SemigroupAlg{A}{\star^{\mathrm{op}}}}{\FormulaRefAuto{SemigroupDef}[def]{7,17}}
\closeproofpartwideindR
\proofpartwide{Die unveränderte Trägerabbildung}
 \proofstepwidestar[1]{f:A\cong B}{\rA}
 \proofstepwidestar[1]{\SemigroupAlg{A}{\star}}{\FormulaRefAuto{SemigroupIsoSource}[thm]{1}}
 \proofstepwidestar[1]{\SemigroupAlg{B}{\diamond}}{\FormulaRefAuto{SemigroupIsoTarget}[thm]{1}}
 \proofstepwidestar[1]{f:A\to B}{\FormulaRefAuto{SemigroupIsoFunction}[thm]{1}}
 \proofstepwidestar[1]{f:A\bij B}{\FormulaRefAuto{SemigroupIsoBijectiveAxiom}[ax]{1}}
 \proofstepwidestar[1]{\SemigroupHom{f}{A,\star}{B,\diamond}}{\FormulaRefAuto{SemigroupIsoHomAxiom}[ax]{1}}
 \proofstepwidestar[1]{f[A]=B}{\FormulaRefAuto{SemigroupIsoFullImage}[thm]{1}}
 \proofstepwidestar[1]{\SemigroupAlg{A}{\star^{\mathrm{op}}}}{\FormulaRefAuto{SemigroupOppositeIsSemigroup}[thm]{2}}
 \proofstepwidestar[1]{\SemigroupAlg{B}{\diamond^{\mathrm{op}}}}{\FormulaRefAuto{SemigroupOppositeIsSemigroup}[thm]{3}}
 \proofstepwidestar[10]{x\in A}{\rA}
 \proofstepwidestar[11]{y\in A}{\rA}
 \proofstepwidestar[10,11]{x,y\in A}{\rAIStar{10,11}}
 \proofstepwidestar[1,10,11]{f(x\star^{\mathrm{op}}y)=f(y\star x)}{\FormulaRefAuto{SemigroupOppositeDef}[def]{2,12}}
 \proofstepwidestar[1,10,11]{f(y\star x)=f(y)\diamond f(x)}{\FormulaRefAuto{SemigroupIsoOperation}[thm]{1,12}}
 \proofstepwidestar[1,10,11]{f(y)\diamond f(x)=f(x)\diamond^{\mathrm{op}}f(y)}{\FormulaRefAuto{SemigroupOppositeDef}[def]{3}}
 \proofstepwidestar[1,10,11]{f(x\star^{\mathrm{op}}y)=f(x)\diamond^{\mathrm{op}}f(y)}{\rChain{13,14,15}}
 \proofstepwidestar[1]{\forall x,y\in A\;(f(x\star^{\mathrm{op}}y)=f(x)\diamond^{\mathrm{op}}f(y))}{\rUI{\rRI{10,\rUI{\rRI{11,16}}}}}
 \proofstepwidestar[1]{\SemigroupHom{f}{A,\star^{\mathrm{op}}}{B,\diamond^{\mathrm{op}}}}{\FormulaRefAuto{SemigroupHomDef}[def]{8,9,4,17}}
 \proofstepwidestar[1]{\SemigroupIso{f}{A,\star^{\mathrm{op}}}{B,\diamond^{\mathrm{op}}}}{\FormulaRefAuto{SemigroupIsoDef}[def]{18,5}}
\closeproofpartwide
\end{tabproofsplitwide}

\subsection{Zentralisatoren, Zentrum und lokale Teilhalbgruppen}

\FormulaDefDeltaK[Zentralisator und Zentrum einer Halbgruppe]{%
 \begin{aligned}[t]
 &\SemigroupAlg{A}{\star}\dsep X\subseteq A\vdash{}\\[-2pt]
 &C_A(X)\coloneqq\{a\in A\mid\forall x\in X\;(a\star x=x\star a)\},
       \qquad Z(A)\coloneqq C_A(A).
 \end{aligned}%
}{SemigroupCentralizerDef}{%
 \DeltaRow{Mengen}{A\dsep X\dsep a\dsep x}
 \DeltaRow{Funktionen}{\star}
}
Zentralisatoren können leer sein. Die folgenden Aussagen verwenden deshalb
die auf den Teilträgern abgeschlossenen assoziativen Operationen;
Nichtleerheit wird nur bei der Bezeichnung Unterhalbgruppe verlangt.

\FormulaThmDeltaKR[Isomorphie von Zentralisatoren und Zentren]{%
 \begin{aligned}[t]
  &f:A\cong B\dsep X\subseteq A\vdash{}\\[-2pt]
  &\text{(i)}\quad f[C_A(X)]=C_B(f[X])\\[-2pt]
  &\text{(ii)}\quad f|_{C_A(X)}:C_A(X)\cong C_B(f[X])\\[-2pt]
  &\text{(iii)}\quad f[Z(A)]=Z(B)\\[-2pt]
  &\text{(iv)}\quad f|_{Z(A)}:Z(A)\cong Z(B)
 \end{aligned}%
}{%
 \begin{aligned}[t]
 &f:A\cong B\dsep X\subseteq A\vdash{}\\[-2pt]
 &f[C_A(X)]=C_B(f[X]),\qquad
       f|_{C_A(X)}:C_A(X)\cong C_B(f[X]),\\[-2pt]
 &f[Z(A)]=Z(B),\qquad f|_{Z(A)}:Z(A)\cong Z(B).
 \end{aligned}%
}{SemigroupIsoCentralizers}{%
 \DeltaRow{Mengen}{A\dsep B\dsep X\dsep a\dsep b\dsep x}
 \DeltaRow{Funktionen}{f\dsep\star\dsep\diamond}
}
\begin{tabproofsplitwide}
\proofpartwideindR[Produktabschluss eines Zentralisators]{%
\SemigroupAlg{A}{\star}\dsep X\subseteq A\vdash\SemigroupAlg{C_A(X)}{\star}%
}{\SemigroupAlg{A}{\star}\dsep X\subseteq A\vdash\SemigroupAlg{C_A(X)}{\star}}[SemigroupCentralizerIsSemigroup]
 \proofstepwidestar[1]{\SemigroupAlg{A}{\star}}{\rA}
 \proofstepwidestar[2]{X\subseteq A}{\rA}
 \proofstepwidestar[3]{a\in C_A(X)}{\rA}
 \proofstepwidestar[4]{b\in C_A(X)}{\rA}
 \proofstepwidestar[3,4]{a,b\in C_A(X)}{\rAIStar{3,4}}
 \proofstepwidestar[3,4]{a,b\in A}{\FormulaRefAuto{SemigroupCentralizerDef}[def]{5}}
 \proofstepwidestar[1,3,4]{a\star b\in A}{\FormulaRefAuto{SemigroupClosureAxiom}[ax]{1,6}}
 \proofstepwidestar[8]{x\in X}{\rA}
 \proofstepwidestar[2,8]{x\in A}{\FormulaRefAuto{A\subseteq B,\,x\in A\vdash x\in B}[thm]{2,8}}
 \proofstepwidestar[3,4,8]{b\star x=x\star b}{\FormulaRefAuto{SemigroupCentralizerDef}[def]{5,8}}
 \proofstepwidestar[3,4,8]{a\star x=x\star a}{\FormulaRefAuto{SemigroupCentralizerDef}[def]{5,8}}
 \proofstepwidestar[1,2,3,4,8]{(a\star b)\star x=a\star(b\star x)}{\FormulaRefAuto{SemigroupAssociativityAxiom}[ax]{1,6,9}}
 \proofstepwidestar[1,2,3,4,8]{a\star(b\star x)=a\star(x\star b)}{\rIE{10}}
 \proofstepwidestar[1,2,3,4,8]{a\star(x\star b)=(a\star x)\star b}{\text{\parbox[t]{\linewidth}{\raggedright \FormulaRefAuto{SemigroupAssociativityAxiom}[ax]{1,6,9}; Symmetrie}}}
 \proofstepwidestar[1,2,3,4,8]{(a\star x)\star b=(x\star a)\star b}{\rIE{11}}
 \proofstepwidestar[1,2,3,4,8]{(x\star a)\star b=x\star(a\star b)}{\FormulaRefAuto{SemigroupAssociativityAxiom}[ax]{1,9,6}}
 \proofstepwidestar[1,2,3,4,8]{(a\star b)\star x=x\star(a\star b)}{\rChain{12,13,14,15,16}}
 \proofstepwidestar[1,2,3,4]{\forall x\in X\;((a\star b)\star x=x\star(a\star b))}{\rUI{\rRI{8,17}}}
 \proofstepwidestar[1,2,3,4]{a\star b\in C_A(X)}{\FormulaRefAuto{SemigroupCentralizerDef}[def]{7,18}}
 \proofstepwidestar[1,2]{\forall a,b\in C_A(X)\;(a\star b\in C_A(X))}{\rUI{\rRI{3,\rUI{\rRI{4,19}}}}}
 \proofstepwidestar[]{C_A(X)\subseteq A}{\FormulaRefAuto{SemigroupCentralizerDef}[def]}
 \proofstepwidestar[1,2]{\SemigroupAlg{C_A(X)}{\star}}{\FormulaRefAuto{SemigroupClosedSubsetIsSemigroup}[thm]{1,21,20}}
\closeproofpartwideindR
\proofpartwideindR[Bildkriterium und Einschränkung]{%
f:A\cong B\dsep X\subseteq A\vdash f|_{C_A(X)}:C_A(X)\cong C_B(f[X])%
}{f:A\cong B\dsep X\subseteq A\vdash f|_{C_A(X)}:C_A(X)\cong C_B(f[X])}[SemigroupIsoCentralizerRestriction]
 \proofstepwidestar[1]{f:A\cong B}{\rA}
 \proofstepwidestar[1]{\SemigroupAlg{A}{\star}}{\FormulaRefAuto{SemigroupIsoSource}[thm]{1}}
 \proofstepwidestar[1]{\SemigroupAlg{B}{\diamond}}{\FormulaRefAuto{SemigroupIsoTarget}[thm]{1}}
 \proofstepwidestar[1]{f:A\to B}{\FormulaRefAuto{SemigroupIsoFunction}[thm]{1}}
 \proofstepwidestar[1]{f:A\bij B}{\FormulaRefAuto{SemigroupIsoBijectiveAxiom}[ax]{1}}
 \proofstepwidestar[1]{\SemigroupHom{f}{A,\star}{B,\diamond}}{\FormulaRefAuto{SemigroupIsoHomAxiom}[ax]{1}}
 \proofstepwidestar[1]{f[A]=B}{\FormulaRefAuto{SemigroupIsoFullImage}[thm]{1}}
 \proofstepwidestar[8]{X\subseteq A}{\rA}
 \proofstepwidestar[1,8]{f[X]\subseteq B}{\FormulaRefAuto{C\subseteq A\dsep F\colon A\to B\vdash F[C]\subseteq B}[thm]{8,4}}
 \proofstepwidestar[10]{a\in A}{\rA}
 \proofstepwidestar[11]{x\in X}{\rA}
 \proofstepwidestar[8,11]{x\in A}{\FormulaRefAuto{A\subseteq B,\,x\in A\vdash x\in B}[thm]{8,11}}
 \proofstepwidestar[1,8,10,11]{a\star x\in A}{\FormulaRefAuto{SemigroupClosureAxiom}[ax]{2,10,12}}
 \proofstepwidestar[1,8,10,11]{x\star a\in A}{\FormulaRefAuto{SemigroupClosureAxiom}[ax]{2,12,10}}
 \proofstepwidestar[1,8,10,11]{a\star x=x\star a\leftrightarrow f(a\star x)=f(x\star a)}{\FormulaRefAuto{SemigroupIsoEqualityReflection}[thm]{1,13,14}}
 \proofstepwidestar[1,8,10,11]{f(a\star x)=f(a)\diamond f(x)}{\FormulaRefAuto{SemigroupIsoOperation}[thm]{1,10,12}}
 \proofstepwidestar[1,8,10,11]{f(x\star a)=f(x)\diamond f(a)}{\FormulaRefAuto{SemigroupIsoOperation}[thm]{1,12,10}}
 \proofstepwidestar[1,8,10,11]{a\star x=x\star a\leftrightarrow f(a)\diamond f(x)=f(x)\diamond f(a)}{\rIE{16,17,15}}
 \proofstepwidestar[1,8,10]{(\forall x\in X\;(a\star x=x\star a))\leftrightarrow(\forall y\in f[X]\;(f(a)\diamond y=y\diamond f(a)))}{\text{\parbox[t]{\linewidth}{\raggedright 18 für alle 11; jedes \(y\in f[X]\) besitzt einen Bildzeugen.}}}
 \proofstepwidestar[1,10]{f(a)\in B}{\FormulaRefAuto{F\colon A\to B\dsep x\in A\vdash F(x)\in B}[thm]{4,10}}
 \proofstepwidestar[1,8,10]{a\in C_A(X)\leftrightarrow f(a)\in C_B(f[X])}{\FormulaRefAuto{SemigroupCentralizerDef}[def]{2,3,8,9,10,20,19}}
 \proofstepwidestar[1,8]{f[C_A(X)]=C_B(f[X])}{\text{\parbox[t]{\linewidth}{\raggedright 21 für alle 10; Bilddefinition und 7}}}
 \proofstepwidestar[]{C_A(X)\subseteq A}{\FormulaRefAuto{SemigroupCentralizerDef}[def]}
 \proofstepwidestar[1,8]{\SemigroupAlg{C_A(X)}{\star}}{\FormulaRefAuto{SemigroupCentralizerIsSemigroup}[thm]{2,8}}
 \proofstepwidestar[1,8]{f|_{C_A(X)}:C_A(X)\cong f[C_A(X)]}{\FormulaRefAuto{SemigroupIsoClosedCarrierRestriction}[thm]{1,23,24}}
 \proofstepwidestar[1,8]{f|_{C_A(X)}:C_A(X)\cong C_B(f[X])}{\rIE{22,25}}
\closeproofpartwideindR
\proofpartwide{Spezialisierung auf das Zentrum}
 \proofstepwidestar[1]{f:A\cong B}{\rA}
 \proofstepwidestar[1]{\SemigroupAlg{A}{\star}}{\FormulaRefAuto{SemigroupIsoSource}[thm]{1}}
 \proofstepwidestar[1]{\SemigroupAlg{B}{\diamond}}{\FormulaRefAuto{SemigroupIsoTarget}[thm]{1}}
 \proofstepwidestar[1]{f:A\to B}{\FormulaRefAuto{SemigroupIsoFunction}[thm]{1}}
 \proofstepwidestar[1]{f:A\bij B}{\FormulaRefAuto{SemigroupIsoBijectiveAxiom}[ax]{1}}
 \proofstepwidestar[1]{\SemigroupHom{f}{A,\star}{B,\diamond}}{\FormulaRefAuto{SemigroupIsoHomAxiom}[ax]{1}}
 \proofstepwidestar[1]{f[A]=B}{\FormulaRefAuto{SemigroupIsoFullImage}[thm]{1}}
 \proofstepwidestar[]{A\subseteq A}{\FormulaRefAuto{A\subseteq A}[thm]}
 \proofstepwidestar[1]{f|_{C_A(A)}:C_A(A)\cong C_B(f[A])}{\FormulaRefAuto{SemigroupIsoCentralizerRestriction}[thm]{1,8}}
 \proofstepwidestar[1]{f|_{Z(A)}:Z(A)\cong Z(B)}{\FormulaRefAuto{SemigroupCentralizerDef}[def]{9,7}}
\closeproofpartwide
\end{tabproofsplitwide}

\FormulaDefDeltaK[Lokaler Träger eines idempotenten Elements]{%
 \begin{aligned}[t]
 &\SemigroupAlg{A}{\star}\dsep e\in A\dsep e\star e=e\vdash{}\\[-2pt]
 &eAe\coloneqq\{(e\star a)\star e\mid a\in A\}.
 \end{aligned}%
}{SemigroupLocalCarrierDef}{%
 \DeltaRow{Mengen}{A\dsep e\dsep a}
 \DeltaRow{Funktionen}{\star}
}

\FormulaThmDeltaKR[Isomorphie der lokalen Träger]{%
 \begin{aligned}[t]
  &f:A\cong B\dsep e\in A\dsep e\star e=e\vdash{}\\[-2pt]
  &\text{(i)}\quad f[eAe]=f(e)Bf(e)\\[-2pt]
  &\text{(ii)}\quad f|_{eAe}:eAe\cong f(e)Bf(e)\\[-2pt]
  &\text{(iii)}\quad \forall u\in eAe\;(e\star u=u\land u\star e=u)
 \end{aligned}%
}{%
 \begin{aligned}[t]
 &f:A\cong B\dsep e\in A\dsep e\star e=e\vdash{}\\[-2pt]
 &f[eAe]=f(e)Bf(e),\qquad
   f|_{eAe}:eAe\cong f(e)Bf(e),\\[-2pt]
 &\forall u\in eAe\;(e\star u=u\land u\star e=u).
 \end{aligned}%
}{SemigroupIsoLocalCarriers}{%
 \DeltaRow{Mengen}{A\dsep B\dsep e\dsep a\dsep b\dsep u\dsep v}
 \DeltaRow{Funktionen}{f\dsep\star\dsep\diamond}
}
\begin{tabproofsplitwide}
\proofpartwideindR[Lokaler Abschluss und neutrales Element]{%
\SemigroupAlg{A}{\star}\dsep e\in A\dsep ee=e\vdash\mathsf{UHGr}(eAe;A,\star)%
}{\SemigroupAlg{A}{\star}\dsep e\in A\dsep ee=e\vdash\mathsf{UHGr}(eAe;A,\star)}[SemigroupLocalCarrierStructure]
  \proofstepwidestar[1]{\SemigroupAlg{A}{\star}}{\rA}
  \proofstepwidestar[2]{e\in A}{\rA}
  \proofstepwidestar[3]{ee=e}{\rA}
  \proofstepwidestar[4]{u\in eAe}{\rA}
  \proofstepwidestar[1,2,3,4]{\exists a\in A\;(u=(ea)e)}{\FormulaRefAuto{SemigroupLocalCarrierDef}[def]{1,2,3,4}}
  \proofstepwidestar[6]{a\in A\land u=(ea)e}{\rA}
  \proofstepwidestar[6]{a\in A}{\rAEa{6}}
  \proofstepwidestar[6]{u=(ea)e}{\rAEb{6}}
  \proofstepwidestar[1,2,6]{ea\in A}{\FormulaRefAuto{SemigroupClosureAxiom}[ax]{1,2,7}}
  \proofstepwidestar[1,2,6]{(ea)e\in A}{\FormulaRefAuto{SemigroupClosureAxiom}[ax]{1,9,2}}
  \proofstepwidestar[6]{(ea)e=u}{\FormulaRefAuto{a=b\vdash b=a}[thm]{8}}
  \proofstepwidestar[1,2,6]{u\in A}{\rIE{11,10}}
  \proofstepwidestar[1,2,6]{(ee)a=e(ea)}{\FormulaRefAuto{SemigroupAssociativityAxiom}[ax]{1,2,2,7}}
  \proofstepwidestar[1,2,6]{e(ea)=(ee)a}{\FormulaRefAuto{a=b\vdash b=a}[thm]{13}}
  \proofstepwidestar[1,2,3,6]{e(ea)=ea}{\rIE{3,14}}
  \proofstepwidestar[1,2,6]{(e(ea))e=e((ea)e)}{\FormulaRefAuto{SemigroupAssociativityAxiom}[ax]{1,2,9,2}}
  \proofstepwidestar[1,2,6]{e((ea)e)=(e(ea))e}{\FormulaRefAuto{a=b\vdash b=a}[thm]{16}}
  \proofstepwidestar[1,2,3,6]{eu=u}{\rIE{11,15,11,17}}
  \proofstepwidestar[1,2,6]{((ea)e)e=(ea)(ee)}{\FormulaRefAuto{SemigroupAssociativityAxiom}[ax]{1,9,2,2}}
  \proofstepwidestar[1,2,3,6]{ue=u}{\rIE{11,3,11,19}}
  \proofstepwidestar[1,2,3,6]{u\in A\land eu=u\land ue=u}{\rAI{12,\rAI{18,20}}}
  \proofstepwidestar[1,2,3,4]{u\in A\land eu=u\land ue=u}{\rEE{5,6,21}}
  \proofstepwidestar[1,2,3]{\forall u\in eAe\;(u\in A\land eu=u\land ue=u)}{\rUI{\rRI{4,22}}}
  \proofstepwidestar[1,2,3,4]{u\in A}{\rAEa{22}}
  \proofstepwidestar[1,2,3]{eAe\subseteq A}{\FormulaRefAuto{SubsetIntroductionRule}[thm]{[4]\vdots 24}}
  \proofstepwidestar[1,2,3,4]{eu=u\land ue=u}{\rAEb{22}}
  \proofstepwidestar[1,2,3]{\forall u\in eAe\;(eu=u\land ue=u)}{\rUI{\rRI{4,26}}}
  \proofstepwidestar[]{(ee)e=(ee)e}{\rII}
  \proofstepwidestar[3]{e=(ee)e}{\rIE{3,3,28}}
  \proofstepwidestar[2,3]{\exists a\in A\;(e=(ea)e)}{\rEIStar{2,29}}
  \proofstepwidestar[1,2,3]{e\in eAe}{\FormulaRefAuto{SemigroupLocalCarrierDef}[def]{1,2,3,30}}
  \proofstepwidestar[1,2,3]{eAe\neq\varnothing}{\FormulaRefAuto{a\in A\vdash A\neq\varnothing}[thm]{31}}
  \proofstepwidestar[33]{u\in eAe}{\rA}
  \proofstepwidestar[34]{v\in eAe}{\rA}
  \proofstepwidestar[1,2,3,33]{u\in A\land eu=u\land ue=u}{\rRE{\rUE{23},33}}
  \proofstepwidestar[1,2,3,34]{v\in A\land ev=v\land ve=v}{\rRE{\rUE{23},34}}
  \proofstepwidestar[1,2,3,33]{u\in A}{\rAEa{35}}
  \proofstepwidestar[1,2,3,34]{v\in A}{\rAEa{36}}
  \proofstepwidestar[1,2,3,33]{eu=u}{\rAEa{\rAEb{35}}}
  \proofstepwidestar[1,2,3,34]{ve=v}{\rAEb{\rAEb{36}}}
  \proofstepwidestar[1,2,3,33,34]{uv\in A}{\FormulaRefAuto{SemigroupClosureAxiom}[ax]{1,37,38}}
  \proofstepwidestar[1,2,3,33,34]{(eu)v=e(uv)}{\FormulaRefAuto{SemigroupAssociativityAxiom}[ax]{1,2,37,38}}
  \proofstepwidestar[1,2,3,33,34]{e(uv)=(eu)v}{\FormulaRefAuto{a=b\vdash b=a}[thm]{42}}
  \proofstepwidestar[1,2,3,33,34]{e(uv)=uv}{\rIE{39,43}}
  \proofstepwidestar[1,2,3,33,34]{uv=e(uv)}{\FormulaRefAuto{a=b\vdash b=a}[thm]{44}}
  \proofstepwidestar[1,2,3,33,34]{(uv)e=u(ve)}{\FormulaRefAuto{SemigroupAssociativityAxiom}[ax]{1,37,38,2}}
  \proofstepwidestar[1,2,3,33,34]{(uv)e=uv}{\rIE{40,46}}
  \proofstepwidestar[1,2,3,33,34]{(e(uv))e=uv}{\rIE{45,47}}
  \proofstepwidestar[1,2,3,33,34]{uv=(e(uv))e}{\FormulaRefAuto{a=b\vdash b=a}[thm]{48}}
  \proofstepwidestar[1,2,3,33,34]{\exists a\in A\;(uv=(ea)e)}{\rEIStar{41,49}}
  \proofstepwidestar[1,2,3,33,34]{uv\in eAe}{\FormulaRefAuto{SemigroupLocalCarrierDef}[def]{1,2,3,50}}
  \proofstepwidestar[1,2,3]{\forall u,v\in eAe\;(uv\in eAe)}{\rUI{\rRI{33,\rUI{\rRI{34,51}}}}}
  \proofstepwidestar[1,2,3]{eAe\subseteq A\land eAe\neq\varnothing\land\forall u,v\in eAe\;(uv\in eAe)}{\rAI{25,\rAI{32,52}}}
  \proofstepwidestar[1]{%
    \begin{aligned}[t]
      &\mathsf{UHGr}(eAe;A,\star)\leftrightarrow{}\\[-2pt]
      &\quad\bigl(eAe\subseteq A\land eAe\neq\varnothing\land{}\\[-2pt]
      &\qquad\forall u\in eAe\,\forall v\in eAe\;(uv\in eAe)\bigr)
    \end{aligned}}{\FormulaRefAuto{SemigroupSubsemigroupCriterion}[thm]{1}}
  \proofstepwidestar[1,2,3]{\mathsf{UHGr}(eAe;A,\star)}{\rRE{\rLREb{54},53}}
\closeproofpartwideindR
\proofpartwide{Bildgleichheit und Isomorphie}
 \proofstepwidestar[1]{f:A\cong B}{\rA}
 \proofstepwidestar[1]{\SemigroupAlg{A}{\star}}{\FormulaRefAuto{SemigroupIsoSource}[thm]{1}}
 \proofstepwidestar[1]{\SemigroupAlg{B}{\diamond}}{\FormulaRefAuto{SemigroupIsoTarget}[thm]{1}}
 \proofstepwidestar[1]{f:A\to B}{\FormulaRefAuto{SemigroupIsoFunction}[thm]{1}}
 \proofstepwidestar[1]{f:A\bij B}{\FormulaRefAuto{SemigroupIsoBijectiveAxiom}[ax]{1}}
 \proofstepwidestar[1]{\SemigroupHom{f}{A,\star}{B,\diamond}}{\FormulaRefAuto{SemigroupIsoHomAxiom}[ax]{1}}
 \proofstepwidestar[1]{f[A]=B}{\FormulaRefAuto{SemigroupIsoFullImage}[thm]{1}}
 \proofstepwidestar[8]{e\in A}{\rA}
 \proofstepwidestar[9]{e\star e=e}{\rA}
 \proofstepwidestar[1,8]{f(e)\in B}{\FormulaRefAuto{F\colon A\to B\dsep x\in A\vdash F(x)\in B}[thm]{4,8}}
 \proofstepwidestar[1,8]{f(e\star e)=f(e)\diamond f(e)}{\FormulaRefAuto{SemigroupIsoOperation}[thm]{1,8}}
 \proofstepwidestar[1,8,9]{f(e)\diamond f(e)=f(e)}{\rIE{9,11}}
 \proofstepwidestar[13]{a\in A}{\rA}
 \proofstepwidestar[1,8,13]{e\star a\in A}{\FormulaRefAuto{SemigroupClosureAxiom}[ax]{2,8,13}}
 \proofstepwidestar[1,8,13]{f((e\star a)\star e)=f(e\star a)\diamond f(e)}{\FormulaRefAuto{SemigroupIsoOperation}[thm]{1,14,8}}
 \proofstepwidestar[1,8,13]{f(e\star a)=f(e)\diamond f(a)}{\FormulaRefAuto{SemigroupIsoOperation}[thm]{1,8,13}}
 \proofstepwidestar[1,8,13]{f((e\star a)\star e)=(f(e)\diamond f(a))\diamond f(e)}{\rIE{16,15}}
 \proofstepwidestar[1,8,9]{f[eAe]=f(e)Bf(e)}{\text{\parbox[t]{\linewidth}{\raggedright 17 für alle 13; 7; lokale Trägerdefinition}}}
 \proofstepwidestar[1,8,9]{\mathsf{UHGr}(eAe;A,\star)}{\FormulaRefAuto{SemigroupLocalCarrierStructure}[thm]{2,8,9}}
 \proofstepwidestar[1,8,9]{eAe\subseteq A}{\FormulaRefAuto{SemigroupSubsemigroupSubsetAxiom}[ax]{2,19}}
 \proofstepwidestar[1,8,9]{\SemigroupAlg{eAe}{\star}}{\FormulaRefAuto{SubsemigroupFormsSemigroup}[thm]{2,19}}
 \proofstepwidestar[1,8,9]{f|_{eAe}:eAe\cong f[eAe]}{\FormulaRefAuto{SemigroupIsoClosedCarrierRestriction}[thm]{1,20,21}}
 \proofstepwidestar[1,8,9]{f|_{eAe}:eAe\cong f(e)Bf(e)}{\rIE{18,22}}
\closeproofpartwide
\end{tabproofsplitwide}

\subsection{Neutrale Elemente, Einheiten und Adjunktionen}

\FormulaDefDeltaK[Einheiten bezüglich eines neutralen Elements]{%
 \begin{aligned}[t]
 &\SemigroupAlg{A}{\star}\dsep e\in A\dsep
  \forall x\in A\;(e\star x=x\land x\star e=x)\vdash{}\\[-2pt]
 &U(A,e)\coloneqq\{u\in A\mid\exists v\in A\;
                  (u\star v=e\land v\star u=e)\}.
 \end{aligned}%
}{SemigroupUnitsAtIdentityDef}{%
 \DeltaRow{Mengen}{A\dsep e\dsep u\dsep v\dsep x}
 \DeltaRow{Funktionen}{\star}
}
Hier werden ausschließlich Multiplikation und Bijektivität benutzt.
Der spätere Gruppenbegriff ist für die folgende Aussage nicht erforderlich.

\FormulaThmDeltaKR[Transport neutraler Elemente, Nullen und Einheiten]{%
 \begin{aligned}[t]
  &f:A\cong B\dsep e,z\in A\vdash{}\\[-2pt]
  &\text{(i)}\quad (\forall x\in A\;(ex=x\land xe=x))\leftrightarrow (\forall y\in B\;(f(e)y=y\land yf(e)=y))\\[-2pt]
  &\text{(ii)}\quad (\forall x\in A\;(zx=z\land xz=z))\leftrightarrow (\forall y\in B\;(f(z)y=f(z)\land yf(z)=f(z)))\\[-2pt]
  &\text{(iii)}\quad
    \begin{aligned}[t]
      &(\forall x\in A\;(ex=x\land xe=x))\\[-2pt]
      &\quad\rightarrow\bigl(f|_{U(A,e)}:U(A,e)\cong U(B,f(e))\bigr)
    \end{aligned}
 \end{aligned}%
}{%
 \begin{aligned}[t]
 &f:A\cong B\dsep e,z\in A\vdash{}\\[-2pt]
 &(\forall x\in A\;(ex=x=xe))\leftrightarrow
    (\forall y\in B\;(f(e)y=y=yf(e))),\\[-2pt]
 &(\forall x\in A\;(zx=z=xz))\leftrightarrow
    (\forall y\in B\;(f(z)y=f(z)=yf(z)));\\[-2pt]
 &e\text{ neutral in }A\Rightarrow
    \bigl(f|_{U(A,e)}:U(A,e)\cong U(B,f(e))\bigr).
 \end{aligned}%
}{SemigroupIsoIdentitiesZerosUnits}{%
 \DeltaRow{Mengen}{A\dsep B\dsep e\dsep z\dsep u\dsep v\dsep x\dsep y}
 \DeltaRow{Funktionen}{f\dsep h\dsep\star\dsep\diamond}
}
\begin{tabproofsplitwide}
\proofpartwideindR[Globale Multiplikationsbedingungen]{%
\begin{aligned}[t]
 &f:A\cong B\dsep e,z\in A\vdash{}\\[-2pt]
 &\text{(i)}\quad (\forall x\in A\;(ex=x\land xe=x))\\[-2pt]
 &\qquad\leftrightarrow(\forall y\in B\;(f(e)y=y\land yf(e)=y))\\[-2pt]
 &\text{(ii)}\quad (\forall x\in A\;(zx=z\land xz=z))\\[-2pt]
 &\qquad\leftrightarrow(\forall y\in B\;(f(z)y=f(z)\land yf(z)=f(z)))
\end{aligned}%
}{f:A\cong B\dsep e,z\in A\vdash\begin{aligned}[t]&(\forall x\in A\;(ex=x\land xe=x))\leftrightarrow(\forall y\in B\;(f(e)y=y\land yf(e)=y)),\\&(\forall x\in A\;(zx=z\land xz=z))\leftrightarrow(\forall y\in B\;(f(z)y=f(z)\land yf(z)=f(z)))\end{aligned}}[SemigroupIsoIdentityZeroConditions]
  \proofstepwidestar[1]{f:A\cong B}{\rA}
  \proofstepwidestar[2]{e\in A}{\rA}
  \proofstepwidestar[3]{z\in A}{\rA}
  \proofstepwidestar[1]{\SemigroupAlg{A}{\star}}{\FormulaRefAuto{SemigroupIsoSource}[thm]{1}}
  \proofstepwidestar[1]{f:A\to B}{\FormulaRefAuto{SemigroupIsoFunction}[thm]{1}}
  \proofstepwidestar[1]{f:A\bij B}{\FormulaRefAuto{SemigroupIsoBijectiveAxiom}[ax]{1}}
  \proofstepwidestar[1]{h:B\bij A}{\FormulaRefAuto{SemigroupIsoInverseCarrier}[thm]{1}}
  \proofstepwidestar[1]{h:B\to A}{\FormulaRefAuto{Bijektive Funktion}[def]{7}}
  \proofstepwidestar[9]{\forall x\in A\;(ex=x\land xe=x)}{\rA}
  \proofstepwidestar[10]{y\in B}{\rA}
  \proofstepwidestar[1,10]{h(y)\in A}{\FormulaRefAuto{F\colon A\to B\dsep x\in A\vdash F(x)\in B}[thm]{8,10}}
  \proofstepwidestar[1,10]{f(h(y))=y}{\FormulaRefAuto{InverseFunctionRightInverse}[thm]{6,10}}
  \proofstepwidestar[1,9,10]{eh(y)=h(y)\land h(y)e=h(y)}{\rRE{\rUE{9},11}}
  \proofstepwidestar[1,9,10]{eh(y)=h(y)}{\rAEa{13}}
  \proofstepwidestar[1,9,10]{h(y)e=h(y)}{\rAEb{13}}
  \proofstepwidestar[1,2,10]{f(eh(y))=f(e)f(h(y))}{\FormulaRefAuto{SemigroupIsoOperation}[thm]{1,2,11}}
  \proofstepwidestar[1,2,10]{f(h(y)e)=f(h(y))f(e)}{\FormulaRefAuto{SemigroupIsoOperation}[thm]{1,11,2}}
  \proofstepwidestar[1,2,9,10]{y=f(e)y}{\rIE{14,12,16}}
  \proofstepwidestar[1,2,9,10]{y=yf(e)}{\rIE{15,12,17}}
  \proofstepwidestar[1,2,9,10]{f(e)y=y}{\FormulaRefAuto{a=b\vdash b=a}[thm]{18}}
  \proofstepwidestar[1,2,9,10]{yf(e)=y}{\FormulaRefAuto{a=b\vdash b=a}[thm]{19}}
  \proofstepwidestar[1,2,9,10]{f(e)y=y\land yf(e)=y}{\rAI{20,21}}
  \proofstepwidestar[1,2,9]{\forall y\in B\;(f(e)y=y\land yf(e)=y)}{\rUI{\rRI{10,22}}}
  \proofstepwidestar[1,2]{\begin{aligned}[t]&(\forall x\in A\;(ex=x\land xe=x))\\[-2pt]&\quad\rightarrow(\forall y\in B\;(f(e)y=y\land yf(e)=y))\end{aligned}}{\rRI{9,23}}
  \proofstepwidestar[25]{\forall y\in B\;(f(e)y=y\land yf(e)=y)}{\rA}
  \proofstepwidestar[26]{x\in A}{\rA}
  \proofstepwidestar[1,26]{f(x)\in B}{\FormulaRefAuto{F\colon A\to B\dsep x\in A\vdash F(x)\in B}[thm]{5,26}}
  \proofstepwidestar[1,25,26]{f(e)f(x)=f(x)\land f(x)f(e)=f(x)}{\rRE{\rUE{25},27}}
  \proofstepwidestar[1,25,26]{f(e)f(x)=f(x)}{\rAEa{28}}
  \proofstepwidestar[1,25,26]{f(x)f(e)=f(x)}{\rAEb{28}}
  \proofstepwidestar[1,2,26]{ex\in A}{\FormulaRefAuto{SemigroupClosureAxiom}[ax]{4,2,26}}
  \proofstepwidestar[1,2,26]{xe\in A}{\FormulaRefAuto{SemigroupClosureAxiom}[ax]{4,26,2}}
  \proofstepwidestar[1,2,26]{f(ex)=f(e)f(x)}{\FormulaRefAuto{SemigroupIsoOperation}[thm]{1,2,26}}
  \proofstepwidestar[1,2,26]{f(xe)=f(x)f(e)}{\FormulaRefAuto{SemigroupIsoOperation}[thm]{1,26,2}}
  \proofstepwidestar[1,2,25,26]{f(ex)=f(x)}{\rIE{29,33}}
  \proofstepwidestar[1,2,25,26]{f(xe)=f(x)}{\rIE{30,34}}
  \proofstepwidestar[1,2,26]{ex=x\leftrightarrow f(ex)=f(x)}{\FormulaRefAuto{SemigroupIsoEqualityReflection}[thm]{1,31,26}}
  \proofstepwidestar[1,2,26]{xe=x\leftrightarrow f(xe)=f(x)}{\FormulaRefAuto{SemigroupIsoEqualityReflection}[thm]{1,32,26}}
  \proofstepwidestar[1,2,25,26]{ex=x}{\rRE{\rLREb{37},35}}
  \proofstepwidestar[1,2,25,26]{xe=x}{\rRE{\rLREb{38},36}}
  \proofstepwidestar[1,2,25,26]{ex=x\land xe=x}{\rAI{39,40}}
  \proofstepwidestar[1,2,25]{\forall x\in A\;(ex=x\land xe=x)}{\rUI{\rRI{26,41}}}
  \proofstepwidestar[1,2]{\begin{aligned}[t]&(\forall y\in B\;(f(e)y=y\land yf(e)=y))\\[-2pt]&\quad\rightarrow(\forall x\in A\;(ex=x\land xe=x))\end{aligned}}{\rRI{25,42}}
  \proofstepwidestar[1,2]{\begin{aligned}[t]&(\forall x\in A\;(ex=x\land xe=x))\\[-2pt]&\quad\leftrightarrow(\forall y\in B\;(f(e)y=y\land yf(e)=y))\end{aligned}}{\rLRI{24,43}}
  \proofstepwidestar[45]{\forall x\in A\;(zx=z\land xz=z)}{\rA}
  \proofstepwidestar[46]{y\in B}{\rA}
  \proofstepwidestar[1,46]{h(y)\in A}{\FormulaRefAuto{F\colon A\to B\dsep x\in A\vdash F(x)\in B}[thm]{8,46}}
  \proofstepwidestar[1,46]{f(h(y))=y}{\FormulaRefAuto{InverseFunctionRightInverse}[thm]{6,46}}
  \proofstepwidestar[1,45,46]{zh(y)=z\land h(y)z=z}{\rRE{\rUE{45},47}}
  \proofstepwidestar[1,45,46]{zh(y)=z}{\rAEa{49}}
  \proofstepwidestar[1,45,46]{h(y)z=z}{\rAEb{49}}
  \proofstepwidestar[1,3,46]{f(zh(y))=f(z)f(h(y))}{\FormulaRefAuto{SemigroupIsoOperation}[thm]{1,3,47}}
  \proofstepwidestar[1,3,46]{f(h(y)z)=f(h(y))f(z)}{\FormulaRefAuto{SemigroupIsoOperation}[thm]{1,47,3}}
  \proofstepwidestar[1,3,45,46]{f(z)=f(z)y}{\rIE{50,48,52}}
  \proofstepwidestar[1,3,45,46]{f(z)=yf(z)}{\rIE{51,48,53}}
  \proofstepwidestar[1,3,45,46]{f(z)y=f(z)}{\FormulaRefAuto{a=b\vdash b=a}[thm]{54}}
  \proofstepwidestar[1,3,45,46]{yf(z)=f(z)}{\FormulaRefAuto{a=b\vdash b=a}[thm]{55}}
  \proofstepwidestar[1,3,45,46]{f(z)y=f(z)\land yf(z)=f(z)}{\rAI{56,57}}
  \proofstepwidestar[1,3,45]{\forall y\in B\;(f(z)y=f(z)\land yf(z)=f(z))}{\rUI{\rRI{46,58}}}
  \proofstepwidestar[1,3]{\begin{aligned}[t]&(\forall x\in A\;(zx=z\land xz=z))\\[-2pt]&\quad\rightarrow(\forall y\in B\;(f(z)y=f(z)\land yf(z)=f(z)))\end{aligned}}{\rRI{45,59}}
  \proofstepwidestar[61]{\forall y\in B\;(f(z)y=f(z)\land yf(z)=f(z))}{\rA}
  \proofstepwidestar[62]{x\in A}{\rA}
  \proofstepwidestar[1,62]{f(x)\in B}{\FormulaRefAuto{F\colon A\to B\dsep x\in A\vdash F(x)\in B}[thm]{5,62}}
  \proofstepwidestar[1,61,62]{f(z)f(x)=f(z)\land f(x)f(z)=f(z)}{\rRE{\rUE{61},63}}
  \proofstepwidestar[1,61,62]{f(z)f(x)=f(z)}{\rAEa{64}}
  \proofstepwidestar[1,61,62]{f(x)f(z)=f(z)}{\rAEb{64}}
  \proofstepwidestar[1,3,62]{zx\in A}{\FormulaRefAuto{SemigroupClosureAxiom}[ax]{4,3,62}}
  \proofstepwidestar[1,3,62]{xz\in A}{\FormulaRefAuto{SemigroupClosureAxiom}[ax]{4,62,3}}
  \proofstepwidestar[1,3,62]{f(zx)=f(z)f(x)}{\FormulaRefAuto{SemigroupIsoOperation}[thm]{1,3,62}}
  \proofstepwidestar[1,3,62]{f(xz)=f(x)f(z)}{\FormulaRefAuto{SemigroupIsoOperation}[thm]{1,62,3}}
  \proofstepwidestar[1,3,61,62]{f(zx)=f(z)}{\rIE{65,69}}
  \proofstepwidestar[1,3,61,62]{f(xz)=f(z)}{\rIE{66,70}}
  \proofstepwidestar[1,3,62]{zx=z\leftrightarrow f(zx)=f(z)}{\FormulaRefAuto{SemigroupIsoEqualityReflection}[thm]{1,67,3}}
  \proofstepwidestar[1,3,62]{xz=z\leftrightarrow f(xz)=f(z)}{\FormulaRefAuto{SemigroupIsoEqualityReflection}[thm]{1,68,3}}
  \proofstepwidestar[1,3,61,62]{zx=z}{\rRE{\rLREb{73},71}}
  \proofstepwidestar[1,3,61,62]{xz=z}{\rRE{\rLREb{74},72}}
  \proofstepwidestar[1,3,61,62]{zx=z\land xz=z}{\rAI{75,76}}
  \proofstepwidestar[1,3,61]{\forall x\in A\;(zx=z\land xz=z)}{\rUI{\rRI{62,77}}}
  \proofstepwidestar[1,3]{\begin{aligned}[t]&(\forall y\in B\;(f(z)y=f(z)\land yf(z)=f(z)))\\[-2pt]&\quad\rightarrow(\forall x\in A\;(zx=z\land xz=z))\end{aligned}}{\rRI{61,78}}
  \proofstepwidestar[1,3]{\begin{aligned}[t]&(\forall x\in A\;(zx=z\land xz=z))\\[-2pt]&\quad\leftrightarrow(\forall y\in B\;(f(z)y=f(z)\land yf(z)=f(z)))\end{aligned}}{\rLRI{60,79}}
\closeproofpartwideindR
\proofpartwideindR[Der Träger der Einheiten]{%
\SemigroupAlg{A}{\star}\dsep e\in A\dsep \forall x\in A\;(ex=x\land xe=x)\vdash U(A,e)\subseteq A\land\SemigroupAlg{U(A,e)}{\star}%
}{\SemigroupAlg{A}{\star}\dsep e\in A\dsep \forall x\in A\;(ex=x\land xe=x)\vdash U(A,e)\subseteq A\land\SemigroupAlg{U(A,e)}{\star}}[SemigroupUnitCarrierStructure]
  \proofstepwidestar[1]{\SemigroupAlg{A}{\star}}{\rA}
  \proofstepwidestar[2]{e\in A}{\rA}
  \proofstepwidestar[3]{\forall x\in A\;(ex=x\land xe=x)}{\rA}
  \proofstepwidestar[1,2,3]{U(A,e)\subseteq A}{\FormulaRefAuto{SemigroupUnitsAtIdentityDef}[def]{1,2,3}}
  \proofstepwidestar[2,3]{ee=e}{\rAEa{\rRE{\rUE{3},2}}}
  \proofstepwidestar[2,3]{\exists v\in A\;(ev=e\land ve=e)}{\rEIStar{2,\rAI{5,5}}}
  \proofstepwidestar[1,2,3]{e\in U(A,e)}{\FormulaRefAuto{SemigroupUnitsAtIdentityDef}[def]{1,2,3,6}}
  \proofstepwidestar[1,2,3]{U(A,e)\neq\varnothing}{\FormulaRefAuto{a\in A\vdash A\neq\varnothing}[thm]{7}}
  \proofstepwidestar[9]{u\in U(A,e)}{\rA}
  \proofstepwidestar[10]{v\in U(A,e)}{\rA}
  \proofstepwidestar[1,2,3,9]{u\in A}{\FormulaRefAuto{A\subseteq B,\,x\in A\vdash x\in B}[thm]{4,9}}
  \proofstepwidestar[1,2,3,10]{v\in A}{\FormulaRefAuto{A\subseteq B,\,x\in A\vdash x\in B}[thm]{4,10}}
  \proofstepwidestar[1,2,3,9]{\exists u^{\prime}\in A\;(uu^{\prime}=e\land u^{\prime}u=e)}{\FormulaRefAuto{SemigroupUnitsAtIdentityDef}[def]{1,2,3,9}}
  \proofstepwidestar[14]{u^{\prime}\in A\land uu^{\prime}=e\land u^{\prime}u=e}{\rA}
  \proofstepwidestar[14]{u^{\prime}\in A}{\rAEa{14}}
  \proofstepwidestar[14]{uu^{\prime}=e}{\rAEa{\rAEb{14}}}
  \proofstepwidestar[14]{u^{\prime}u=e}{\rAEb{\rAEb{14}}}
  \proofstepwidestar[1,2,3,10]{\exists v^{\prime}\in A\;(vv^{\prime}=e\land v^{\prime}v=e)}{\FormulaRefAuto{SemigroupUnitsAtIdentityDef}[def]{1,2,3,10}}
  \proofstepwidestar[19]{v^{\prime}\in A\land vv^{\prime}=e\land v^{\prime}v=e}{\rA}
  \proofstepwidestar[19]{v^{\prime}\in A}{\rAEa{19}}
  \proofstepwidestar[19]{vv^{\prime}=e}{\rAEa{\rAEb{19}}}
  \proofstepwidestar[19]{v^{\prime}v=e}{\rAEb{\rAEb{19}}}
  \proofstepwidestar[1,2,3,9,10]{uv\in A}{\FormulaRefAuto{SemigroupClosureAxiom}[ax]{1,11,12}}
  \proofstepwidestar[1,14,19]{v^{\prime}u^{\prime}\in A}{\FormulaRefAuto{SemigroupClosureAxiom}[ax]{1,20,15}}
  \proofstepwidestar[1,2,3,9,10,14,19]{(uv)(v^{\prime}u^{\prime})=u(v(v^{\prime}u^{\prime}))}{\FormulaRefAuto{SemigroupAssociativityAxiom}[ax]{1,11,12,24}}
  \proofstepwidestar[1,2,3,10,14,19]{(vv^{\prime})u^{\prime}=v(v^{\prime}u^{\prime})}{\FormulaRefAuto{SemigroupAssociativityAxiom}[ax]{1,12,20,15}}
  \proofstepwidestar[1,2,3,10,14,19]{v(v^{\prime}u^{\prime})=(vv^{\prime})u^{\prime}}{\FormulaRefAuto{a=b\vdash b=a}[thm]{26}}
  \proofstepwidestar[3,14]{eu^{\prime}=u^{\prime}}{\rAEa{\rRE{\rUE{3},15}}}
  \proofstepwidestar[1,2,3,10,14,19]{v(v^{\prime}u^{\prime})=u^{\prime}}{\rIE{21,28,27}}
  \proofstepwidestar[1,2,3,9,10,14,19]{(uv)(v^{\prime}u^{\prime})=e}{\rIE{29,16,25}}
  \proofstepwidestar[1,2,3,9,10,14,19]{(v^{\prime}u^{\prime})(uv)=v^{\prime}(u^{\prime}(uv))}{\FormulaRefAuto{SemigroupAssociativityAxiom}[ax]{1,20,15,23}}
  \proofstepwidestar[1,2,3,9,10,14]{(u^{\prime}u)v=u^{\prime}(uv)}{\FormulaRefAuto{SemigroupAssociativityAxiom}[ax]{1,15,11,12}}
  \proofstepwidestar[1,2,3,9,10,14]{u^{\prime}(uv)=(u^{\prime}u)v}{\FormulaRefAuto{a=b\vdash b=a}[thm]{32}}
  \proofstepwidestar[1,2,3,10]{ev=v}{\rAEa{\rRE{\rUE{3},12}}}
  \proofstepwidestar[1,2,3,9,10,14]{u^{\prime}(uv)=v}{\rIE{17,34,33}}
  \proofstepwidestar[1,2,3,9,10,14,19]{(v^{\prime}u^{\prime})(uv)=e}{\rIE{35,22,31}}
  \proofstepwidestar[1,2,3,9,10,14,19]{\exists w\in A\;((uv)w=e\land w(uv)=e)}{\rEIStar{24,\rAI{30,36}}}
  \proofstepwidestar[1,2,3,9,10,14,19]{uv\in U(A,e)}{\FormulaRefAuto{SemigroupUnitsAtIdentityDef}[def]{1,2,3,23,37}}
  \proofstepwidestar[1,2,3,9,10,14]{uv\in U(A,e)}{\rEE{18,19,38}}
  \proofstepwidestar[1,2,3,9,10]{uv\in U(A,e)}{\rEE{13,14,39}}
  \proofstepwidestar[1,2,3]{\forall u,v\in U(A,e)\;(uv\in U(A,e))}{\rUI{\rRI{9,\rUI{\rRI{10,40}}}}}
  \proofstepwidestar[1,2,3]{\SemigroupAlg{U(A,e)}{\star}}{\FormulaRefAuto{SemigroupClosedSubsetIsSemigroup}[thm]{1,4,41}}
  \proofstepwidestar[1,2,3]{U(A,e)\subseteq A\land\SemigroupAlg{U(A,e)}{\star}}{\rAI{4,42}}
\closeproofpartwideindR
\proofpartwide{Abschluss und Transport der Einheiten}
  \proofstepwidestar[1]{f:A\cong B}{\rA}
  \proofstepwidestar[2]{e\in A}{\rA}
  \proofstepwidestar[3]{\forall x\in A\;(ex=x\land xe=x)}{\rA}
  \proofstepwidestar[1]{\SemigroupAlg{A}{\star}}{\FormulaRefAuto{SemigroupIsoSource}[thm]{1}}
  \proofstepwidestar[1]{\SemigroupAlg{B}{\diamond}}{\FormulaRefAuto{SemigroupIsoTarget}[thm]{1}}
  \proofstepwidestar[1]{f:A\to B}{\FormulaRefAuto{SemigroupIsoFunction}[thm]{1}}
  \proofstepwidestar[1]{f:A\bij B}{\FormulaRefAuto{SemigroupIsoBijectiveAxiom}[ax]{1}}
  \proofstepwidestar[1]{h:B\bij A}{\FormulaRefAuto{SemigroupIsoInverseCarrier}[thm]{1}}
  \proofstepwidestar[1]{h:B\to A}{\FormulaRefAuto{Bijektive Funktion}[def]{8}}
  \proofstepwidestar[1,2]{f(e)\in B}{\FormulaRefAuto{F\colon A\to B\dsep x\in A\vdash F(x)\in B}[thm]{6,2}}
  \proofstepwidestar[1,2]{\begin{aligned}[t]&(\forall x\in A\;(ex=x\land xe=x))\\[-2pt]&\quad\leftrightarrow(\forall y\in B\;(f(e)y=y\land yf(e)=y))\end{aligned}}{\FormulaRefAuto{SemigroupIsoIdentityZeroConditions}[thm]{1,\rAI{2,2}}}
  \proofstepwidestar[1,2,3]{\forall y\in B\;(f(e)y=y\land yf(e)=y)}{\rRE{\rLREa{11},3}}
  \proofstepwidestar[1,2,3]{U(A,e)\subseteq A\land\SemigroupAlg{U(A,e)}{\star}}{\FormulaRefAuto{SemigroupUnitCarrierStructure}[thm]{4,2,3}}
  \proofstepwidestar[1,2,3]{U(A,e)\subseteq A}{\rAEa{13}}
  \proofstepwidestar[1,2,3]{\SemigroupAlg{U(A,e)}{\star}}{\rAEb{13}}
  \proofstepwidestar[1,2,3]{U(B,f(e))\subseteq B\land\SemigroupAlg{U(B,f(e))}{\diamond}}{\FormulaRefAuto{SemigroupUnitCarrierStructure}[thm]{5,10,12}}
  \proofstepwidestar[1,2,3]{U(B,f(e))\subseteq B}{\rAEa{16}}
  \proofstepwidestar[18]{u\in A}{\rA}
  \proofstepwidestar[19]{v\in A}{\rA}
  \proofstepwidestar[1,18,19]{uv\in A}{\FormulaRefAuto{SemigroupClosureAxiom}[ax]{4,18,19}}
  \proofstepwidestar[1,18,19]{vu\in A}{\FormulaRefAuto{SemigroupClosureAxiom}[ax]{4,19,18}}
  \proofstepwidestar[1,18,19]{f(uv)=f(u)f(v)}{\FormulaRefAuto{SemigroupIsoOperation}[thm]{1,18,19}}
  \proofstepwidestar[1,18,19]{f(vu)=f(v)f(u)}{\FormulaRefAuto{SemigroupIsoOperation}[thm]{1,19,18}}
  \proofstepwidestar[1,2,18,19]{uv=e\leftrightarrow f(uv)=f(e)}{\FormulaRefAuto{SemigroupIsoEqualityReflection}[thm]{1,20,2}}
  \proofstepwidestar[1,2,18,19]{vu=e\leftrightarrow f(vu)=f(e)}{\FormulaRefAuto{SemigroupIsoEqualityReflection}[thm]{1,21,2}}
  \proofstepwidestar[1,2,18,19]{uv=e\leftrightarrow f(u)f(v)=f(e)}{\rIE{22,24}}
  \proofstepwidestar[1,2,18,19]{vu=e\leftrightarrow f(v)f(u)=f(e)}{\rIE{23,25}}
  \proofstepwidestar[]{(uv=e\land vu=e)\leftrightarrow(uv=e\land vu=e)}{\FormulaRefAuto{P\leftrightarrow P}[thm]}
  \proofstepwidestar[1,2,18,19]{(uv=e\land vu=e)\leftrightarrow(f(u)f(v)=f(e)\land f(v)f(u)=f(e))}{\rLRS{26,\rLRS{27,28}}}
  \proofstepwidestar[1,2]{\begin{aligned}[t]&\forall u,v\in A\;\bigl((uv=e\land vu=e)\\[-2pt]&\quad\leftrightarrow(f(u)f(v)=f(e)\land f(v)f(u)=f(e))\bigr)\end{aligned}}{\rUI{\rRI{18,\rUI{\rRI{19,29}}}}}
  \proofstepwidestar[31]{u\in A}{\rA}
  \proofstepwidestar[1,31]{f(u)\in B}{\FormulaRefAuto{F\colon A\to B\dsep x\in A\vdash F(x)\in B}[thm]{6,31}}
  \proofstepwidestar[33]{u\in U(A,e)}{\rA}
  \proofstepwidestar[1,2,3,33]{\exists v\in A\;(uv=e\land vu=e)}{\FormulaRefAuto{SemigroupUnitsAtIdentityDef}[def]{4,2,3,33}}
  \proofstepwidestar[35]{v\in A\land uv=e\land vu=e}{\rA}
  \proofstepwidestar[35]{v\in A}{\rAEa{35}}
  \proofstepwidestar[35]{uv=e\land vu=e}{\rAEb{35}}
  \proofstepwidestar[1,35]{f(v)\in B}{\FormulaRefAuto{F\colon A\to B\dsep x\in A\vdash F(x)\in B}[thm]{6,36}}
  \proofstepwidestar[1,2,31,35]{(uv=e\land vu=e)\leftrightarrow(f(u)f(v)=f(e)\land f(v)f(u)=f(e))}{\rRE{\rUE{\rRE{\rUE{30},31}},36}}
  \proofstepwidestar[1,2,31,35]{f(u)f(v)=f(e)\land f(v)f(u)=f(e)}{\rRE{\rLREa{39},37}}
  \proofstepwidestar[1,2,31,35]{\exists w\in B\;(f(u)w=f(e)\land wf(u)=f(e))}{\rEIStar{38,40}}
  \proofstepwidestar[1,2,3,31,35]{f(u)\in U(B,f(e))}{\FormulaRefAuto{SemigroupUnitsAtIdentityDef}[def]{5,10,12,32,41}}
  \proofstepwidestar[1,2,3,31,33]{f(u)\in U(B,f(e))}{\rEE{34,35,42}}
  \proofstepwidestar[1,2,3,31]{u\in U(A,e)\rightarrow f(u)\in U(B,f(e))}{\rRI{33,43}}
  \proofstepwidestar[45]{f(u)\in U(B,f(e))}{\rA}
  \proofstepwidestar[1,2,3,45]{\exists w\in B\;(f(u)w=f(e)\land wf(u)=f(e))}{\FormulaRefAuto{SemigroupUnitsAtIdentityDef}[def]{5,10,12,45}}
  \proofstepwidestar[47]{w\in B\land f(u)w=f(e)\land wf(u)=f(e)}{\rA}
  \proofstepwidestar[47]{w\in B}{\rAEa{47}}
  \proofstepwidestar[47]{f(u)w=f(e)\land wf(u)=f(e)}{\rAEb{47}}
  \proofstepwidestar[1,47]{h(w)\in A}{\FormulaRefAuto{F\colon A\to B\dsep x\in A\vdash F(x)\in B}[thm]{9,48}}
  \proofstepwidestar[1,47]{f(h(w))=w}{\FormulaRefAuto{InverseFunctionRightInverse}[thm]{7,48}}
  \proofstepwidestar[1,47]{w=f(h(w))}{\FormulaRefAuto{a=b\vdash b=a}[thm]{51}}
  \proofstepwidestar[1,47]{f(u)f(h(w))=f(e)\land f(h(w))f(u)=f(e)}{\rIE{52,49}}
  \proofstepwidestar[1,2,31,47]{\begin{aligned}[t]&(uh(w)=e\land h(w)u=e)\\[-2pt]&\quad\leftrightarrow(f(u)f(h(w))=f(e)\\[-2pt]&\qquad\land f(h(w))f(u)=f(e))\end{aligned}}{\rRE{\rUE{\rRE{\rUE{30},31}},50}}
  \proofstepwidestar[1,2,31,47]{uh(w)=e\land h(w)u=e}{\rRE{\rLREb{54},53}}
  \proofstepwidestar[1,2,31,47]{\exists v\in A\;(uv=e\land vu=e)}{\rEIStar{50,55}}
  \proofstepwidestar[1,2,3,31,47]{u\in U(A,e)}{\FormulaRefAuto{SemigroupUnitsAtIdentityDef}[def]{4,2,3,31,56}}
  \proofstepwidestar[1,2,3,31,45]{u\in U(A,e)}{\rEE{46,47,57}}
  \proofstepwidestar[1,2,3,31]{f(u)\in U(B,f(e))\rightarrow u\in U(A,e)}{\rRI{45,58}}
  \proofstepwidestar[1,2,3,31]{u\in U(A,e)\leftrightarrow f(u)\in U(B,f(e))}{\rLRI{44,59}}
  \proofstepwidestar[1,2,3]{\forall u\in A\;(u\in U(A,e)\leftrightarrow f(u)\in U(B,f(e)))}{\rUI{\rRI{31,60}}}
  \proofstepwidestar[62]{y\in f[U(A,e)]}{\rA}
  \proofstepwidestar[1,2,3,62]{\exists u\in U(A,e)\;(y=f(u))}{\FormulaRefAuto{C\subseteq A\dsep F\colon A\to B\dsep y\in F[C]\vdash\exists x\in C\,y=F(x)}[thm]{14,6,62}}
  \proofstepwidestar[64]{u\in U(A,e)\land y=f(u)}{\rA}
  \proofstepwidestar[64]{u\in U(A,e)}{\rAEa{64}}
  \proofstepwidestar[64]{y=f(u)}{\rAEb{64}}
  \proofstepwidestar[64]{f(u)=y}{\FormulaRefAuto{a=b\vdash b=a}[thm]{66}}
  \proofstepwidestar[1,2,3,64]{u\in A}{\FormulaRefAuto{A\subseteq B,\,x\in A\vdash x\in B}[thm]{14,65}}
  \proofstepwidestar[1,2,3,64]{u\in U(A,e)\leftrightarrow f(u)\in U(B,f(e))}{\rRE{\rUE{61},68}}
  \proofstepwidestar[1,2,3,64]{f(u)\in U(B,f(e))}{\rRE{\rLREa{69},65}}
  \proofstepwidestar[1,2,3,64]{y\in U(B,f(e))}{\rIE{67,70}}
  \proofstepwidestar[1,2,3,62]{y\in U(B,f(e))}{\rEE{63,64,71}}
  \proofstepwidestar[1,2,3]{y\in f[U(A,e)]\rightarrow y\in U(B,f(e))}{\rRI{62,72}}
  \proofstepwidestar[74]{y\in U(B,f(e))}{\rA}
  \proofstepwidestar[1,2,3,74]{y\in B}{\FormulaRefAuto{A\subseteq B,\,x\in A\vdash x\in B}[thm]{17,74}}
  \proofstepwidestar[1,2,3,74]{h(y)\in A}{\FormulaRefAuto{F\colon A\to B\dsep x\in A\vdash F(x)\in B}[thm]{9,75}}
  \proofstepwidestar[1,2,3,74]{f(h(y))=y}{\FormulaRefAuto{InverseFunctionRightInverse}[thm]{7,75}}
  \proofstepwidestar[1,2,3,74]{y=f(h(y))}{\FormulaRefAuto{a=b\vdash b=a}[thm]{77}}
  \proofstepwidestar[1,2,3,74]{f(h(y))\in U(B,f(e))}{\rIE{78,74}}
  \proofstepwidestar[1,2,3,74]{h(y)\in U(A,e)\leftrightarrow f(h(y))\in U(B,f(e))}{\rRE{\rUE{61},76}}
  \proofstepwidestar[1,2,3,74]{h(y)\in U(A,e)}{\rRE{\rLREb{80},79}}
  \proofstepwidestar[1,2,3,74]{f(h(y))\in f[U(A,e)]}{\FormulaRefAuto{C\subseteq A\dsep F\colon A\to B\dsep x\in C\vdash F(x)\in F[C]}[thm]{14,6,81}}
  \proofstepwidestar[1,2,3,74]{y\in f[U(A,e)]}{\rIE{77,82}}
  \proofstepwidestar[1,2,3]{y\in U(B,f(e))\rightarrow y\in f[U(A,e)]}{\rRI{74,83}}
  \proofstepwidestar[1,2,3]{y\in f[U(A,e)]\leftrightarrow y\in U(B,f(e))}{\rLRI{73,84}}
  \proofstepwidestar[1,2,3]{\forall y\;(y\in f[U(A,e)]\leftrightarrow y\in U(B,f(e)))}{\rUI{85}}
  \proofstepwidestar[1,2,3]{f[U(A,e)]=U(B,f(e))}{\FormulaRefAuto{\forall x\,(x\in A\leftrightarrow x\in B)\vdash A=B}[ax]{86}}
  \proofstepwidestar[1,2,3]{f|_{U(A,e)}:U(A,e)\cong f[U(A,e)]}{\FormulaRefAuto{SemigroupIsoClosedCarrierRestriction}[thm]{1,14,15}}
  \proofstepwidestar[1,2,3]{f|_{U(A,e)}:U(A,e)\cong U(B,f(e))}{\rIE{87,88}}
  \proofstepwidestar[1,2]{%
    \begin{aligned}[t]
      &(\forall x\in A\;(ex=x\land xe=x))\\[-2pt]
      &\quad\rightarrow\bigl(f|_{U(A,e)}:U(A,e)\cong U(B,f(e))\bigr)
    \end{aligned}}{\rRI{3,89}}
\closeproofpartwide
\end{tabproofsplitwide}

\FormulaDefDeltaK[Adjunktion eines frischen neutralen oder Nullelements]{%
 \begin{aligned}[t]
 &A^{+1}\coloneqq(A\times\{0\})\cup\{(\varnothing,1)\},
       \quad\iota_A(a)\coloneqq(a,0),\quad 1_A\coloneqq(\varnothing,1),\\[-2pt]
 &\iota_A(a)\iota_A(b)\coloneqq\iota_A(a\star b),
       \qquad 1_Ax=x1_A\coloneqq x;\\[-2pt]
 &A^{+0}\coloneqq(A\times\{0\})\cup\{(\varnothing,1)\},
       \quad 0_A\coloneqq(\varnothing,1),\\[-2pt]
 &\iota_A(a)\iota_A(b)\coloneqq\iota_A(a\star b),
       \qquad 0_Ax=x0_A\coloneqq0_A.
 \end{aligned}%
}{SemigroupFreshAdjunctionDef}{%
 \DeltaRow{Voraussetzung}{\SemigroupAlg{A}{\star}}
 \DeltaRow{Mengen}{A\dsep a\dsep b\dsep x}
 \DeltaRow{Funktionen}{\star\dsep\iota_A}
}
Die beiden Träger sind als Mengen gleich. Für \(A\neq\varnothing\)
sind ihre Operationen verschieden; für \(A=\varnothing\) ergeben beide
Adjunktionen die einelementige Halbgruppe.
Die Tags \(0\) und \(1\) halten das neue Element von der Kopie von \(A\)
getrennt. Es wird auch dann ein frisches Element adjungiert, wenn \(A\)
bereits ein neutrales Element beziehungsweise eine Null besitzt.
Damit kollidiert \(A^{+1}\) nicht mit der Produktpotenz \(A^1=A\).

\FormulaThmDeltaKR[Fortsetzung auf adjungierte Elemente]{%
 \begin{aligned}[t]
  &f:A\cong B\vdash{}\\[-2pt]
  &\text{(i)}\quad f^{+1}:A^{+1}\cong B^{+1}\\[-2pt]
  &\text{(ii)}\quad f^{+0}:A^{+0}\cong B^{+0}\\[-2pt]
  &\text{(iii)}\quad f^{+\varepsilon}(\iota_A(a))=\iota_B(f(a))\\[-2pt]
  &\text{(iv)}\quad f^{+1}(1_A)=1_B\\[-2pt]
  &\text{(v)}\quad f^{+0}(0_A)=0_B
 \end{aligned}%
}{%
 \begin{aligned}[t]
 &f:A\cong B\vdash f^{+1}:A^{+1}\cong B^{+1},
       \qquad f^{+0}:A^{+0}\cong B^{+0},\\[-2pt]
 &f^{+\varepsilon}(\iota_A(a))=\iota_B(f(a)),
       \quad f^{+1}(1_A)=1_B,\quad f^{+0}(0_A)=0_B.
 \end{aligned}%
}{SemigroupIsoFreshAdjunctions}{%
 \DeltaRow{Mengen}{A\dsep B\dsep a\dsep b\dsep x\dsep y\dsep z}
 \DeltaRow{Funktionen}{f\dsep h\dsep f^{+1}\dsep f^{+0}\dsep\iota_A\dsep\iota_B}
}
\begin{tabproofsplitwide}
\proofpartwideindR[Abschluss und Assoziativität der Adjunktionen]{%
\SemigroupAlg{A}{\star}\vdash\SemigroupAlg{A^{+1}}{\cdot}\land\SemigroupAlg{A^{+0}}{\cdot}%
}{\SemigroupAlg{A}{\star}\vdash\SemigroupAlg{A^{+1}}{\cdot}\land\SemigroupAlg{A^{+0}}{\cdot}}[SemigroupFreshAdjunctionsAreSemigroups]
 \proofstepwidestar[1]{\SemigroupAlg{A}{\star}}{\rA}
 \proofstepwidestar[2]{a\in A}{\rA}
  \proofstepwidestar[3]{b\in A}{\rA}
  \proofstepwidestar[1,2,3]{ab\in A}{\FormulaRefAuto{SemigroupClosureAxiom}[ax]{1,2,3}}
  \proofstepwidestar[1]{\forall a,b\in A\;(ab\in A)}{\rUI{\rRI{2,\rUI{\rRI{3,4}}}}}
 \proofstepwidestar[6]{a\in A}{\rA}
  \proofstepwidestar[7]{b\in A}{\rA}
  \proofstepwidestar[8]{c\in A}{\rA}
  \proofstepwidestar[1,6,7,8]{(ab)c=a(bc)}{\FormulaRefAuto{SemigroupAssociativityAxiom}[ax]{1,6,7,8}}
  \proofstepwidestar[1]{\forall a,b,c\in A\;((ab)c=a(bc))}{\rUI{\rRI{6,\rUI{\rRI{7,\rUI{\rRI{8,9}}}}}}}
 \proofstepwidestar[1]{\forall a,b\in A\;(\iota_A(a)\iota_A(b)=\iota_A(ab)\in\iota_A[A])}{\FormulaRefAuto{SemigroupFreshAdjunctionDef}[def]{5}}
 \proofstepwidestar[1]{\forall x\in A^{+1}\;(1_Ax=x=x1_A)\,,\quad\forall x\in A^{+0}\;(0_Ax=0_A=x0_A)}{\FormulaRefAuto{SemigroupFreshAdjunctionDef}[def]}
 \proofstepwidestar[1]{\forall\varepsilon\in\{0,1\}\;\forall x,y\in A^{+\varepsilon}\;(xy\in A^{+\varepsilon})}{\text{\parbox[t]{\linewidth}{\raggedright 11,12; jeder Faktor hat einen alten oder den neuen Tag.}}}
 \proofstepwidestar[1]{\forall a,b,c\in A\;((\iota_A(a)\iota_A(b))\iota_A(c)=\iota_A(a)(\iota_A(b)\iota_A(c)))}{\FormulaRefAuto{SemigroupFreshAdjunctionDef}[def]{10}}
 \proofstepwidestar[1]{\forall x,y\in A^{+1}\;\begin{aligned}[t]&(1_Ax)y=xy=1_A(xy),\\&(x1_A)y=xy=x(1_Ay),\\&(xy)1_A=xy=x(y1_A)\end{aligned}}{\FormulaRefAuto{SemigroupFreshAdjunctionDef}[def]}
 \proofstepwidestar[1]{\forall x,y,z\in A^{+0}\;((x=0_A\lor y=0_A\lor z=0_A)\Rightarrow(xy)z=0_A=x(yz))}{\text{\parbox[t]{\linewidth}{\raggedright \FormulaRefAuto{SemigroupFreshAdjunctionDef}[def]; beide Klammerungen absorbieren den Nullfaktor.}}}
 \proofstepwidestar[1]{\forall\varepsilon\in\{0,1\}\;\forall x,y,z\in A^{+\varepsilon}\;((xy)z=x(yz))}{\text{\parbox[t]{\linewidth}{\raggedright 14,15,16; Fallunterscheidung nach den Tags}}}
 \proofstepwidestar[1]{\SemigroupAlg{A^{+1}}{\cdot}\land\SemigroupAlg{A^{+0}}{\cdot}}{\FormulaRefAuto{SemigroupDef}[def]{13,17}}
\closeproofpartwideindR
\proofpartwide{Bijektion und Produktfälle}
 \proofstepwidestar[1]{f:A\cong B}{\rA}
 \proofstepwidestar[1]{\SemigroupAlg{A}{\star}}{\FormulaRefAuto{SemigroupIsoSource}[thm]{1}}
 \proofstepwidestar[1]{\SemigroupAlg{B}{\diamond}}{\FormulaRefAuto{SemigroupIsoTarget}[thm]{1}}
 \proofstepwidestar[1]{f:A\to B}{\FormulaRefAuto{SemigroupIsoFunction}[thm]{1}}
 \proofstepwidestar[1]{f:A\bij B}{\FormulaRefAuto{SemigroupIsoBijectiveAxiom}[ax]{1}}
 \proofstepwidestar[1]{\SemigroupHom{f}{A,\star}{B,\diamond}}{\FormulaRefAuto{SemigroupIsoHomAxiom}[ax]{1}}
 \proofstepwidestar[1]{f[A]=B}{\FormulaRefAuto{SemigroupIsoFullImage}[thm]{1}}
 \proofstepwidestar[8]{\varepsilon\in\{0,1\}}{\rA}
 \proofstepwidestar[1]{\SemigroupAlg{A^{+1}}{\cdot}\land\SemigroupAlg{A^{+0}}{\cdot}}{\FormulaRefAuto{SemigroupFreshAdjunctionsAreSemigroups}[thm]{2}}
 \proofstepwidestar[1]{\SemigroupAlg{B^{+1}}{\cdot}\land\SemigroupAlg{B^{+0}}{\cdot}}{\FormulaRefAuto{SemigroupFreshAdjunctionsAreSemigroups}[thm]{3}}
 \proofstepwidestar[1,8]{\SemigroupAlg{A^{+\varepsilon}}{\cdot}}{\text{\parbox[t]{\linewidth}{\raggedright 9,8; Konjunktionselimination und Fallunterscheidung}}}
 \proofstepwidestar[1,8]{\SemigroupAlg{B^{+\varepsilon}}{\cdot}}{\text{\parbox[t]{\linewidth}{\raggedright 10,8; Konjunktionselimination und Fallunterscheidung}}}
 \proofstepwidestar[1]{h:B\bij A}{\FormulaRefAuto{SemigroupIsoInverseCarrier}[thm]{1}}
 \proofstepwidestar[1]{h:B\to A}{\FormulaRefAuto{Bijektive Funktion}[def]{13}}
 \proofstepwidestar[1,8]{f^{+\varepsilon}:A^{+\varepsilon}\to B^{+\varepsilon}\,,\quad h^{+\varepsilon}:B^{+\varepsilon}\to A^{+\varepsilon}}{\text{\parbox[t]{\linewidth}{\raggedright \FormulaRefAuto{SemigroupFreshAdjunctionDef}[def]{4,14}; die beiden Tagfälle sind disjunkt und erschöpfend.}}}
 \proofstepwidestar[16]{a\in A}{\rA}
  \proofstepwidestar[1,16]{h(f(a))=a}{\FormulaRefAuto{InverseFunctionLeftInverse}[thm]{5,16}}
  \proofstepwidestar[1]{\forall a\in A\;(h(f(a))=a)}{\rUI{\rRI{16,17}}}
 \proofstepwidestar[19]{b\in B}{\rA}
  \proofstepwidestar[1,19]{f(h(b))=b}{\FormulaRefAuto{InverseFunctionRightInverse}[thm]{5,19}}
  \proofstepwidestar[1]{\forall b\in B\;(f(h(b))=b)}{\rUI{\rRI{19,20}}}
 \proofstepwidestar[1,8]{h^{+\varepsilon}\circ f^{+\varepsilon}=\Id_{A^{+\varepsilon}}\,,\quad f^{+\varepsilon}\circ h^{+\varepsilon}=\Id_{B^{+\varepsilon}}}{\text{\parbox[t]{\linewidth}{\raggedright Auf alten Tags 18,21; auf dem neuen Tag die Definition; Funktionsextensionalität.}}}
 \proofstepwidestar[1,8]{f^{+\varepsilon}:A^{+\varepsilon}\bij B^{+\varepsilon}}{\FormulaRefAuto{MutuallyInverseFunctionBijection}[thm]{\rAEb{15},\rAEa{15},\rAEb{22},\rAEa{22}}}
 \proofstepwidestar[24]{a\in A}{\rA}
  \proofstepwidestar[25]{b\in A}{\rA}
  \proofstepwidestar[1,24,25]{f(ab)=f(a)f(b)}{\FormulaRefAuto{SemigroupIsoOperation}[thm]{1,24,25}}
  \proofstepwidestar[1]{\forall a,b\in A\;(f(ab)=f(a)f(b))}{\rUI{\rRI{24,\rUI{\rRI{25,26}}}}}
 \proofstepwidestar[1,8]{\forall a,b\in A\;(f^{+\varepsilon}(\iota_A(a)\iota_A(b))=\iota_B(f(ab))=f^{+\varepsilon}(\iota_A(a))f^{+\varepsilon}(\iota_A(b)))}{\FormulaRefAuto{SemigroupFreshAdjunctionDef}[def]{27}}
 \proofstepwidestar[1]{\forall x\in A^{+1}\;\begin{aligned}[t]&f^{+1}(1_Ax)=f^{+1}(x)=1_Bf^{+1}(x),\\&f^{+1}(x1_A)=f^{+1}(x)=f^{+1}(x)1_B\end{aligned}}{\FormulaRefAuto{SemigroupFreshAdjunctionDef}[def]}
 \proofstepwidestar[1]{\forall x\in A^{+0}\;\begin{aligned}[t]&f^{+0}(0_Ax)=0_B=0_Bf^{+0}(x),\\&f^{+0}(x0_A)=0_B=f^{+0}(x)0_B\end{aligned}}{\FormulaRefAuto{SemigroupFreshAdjunctionDef}[def]}
 \proofstepwidestar[1,8]{\forall x,y\in A^{+\varepsilon}\;(f^{+\varepsilon}(xy)=f^{+\varepsilon}(x)f^{+\varepsilon}(y))}{\text{\parbox[t]{\linewidth}{\raggedright 28,29,30; erschöpfende Tagfälle}}}
 \proofstepwidestar[1,8]{\SemigroupHom{f^{+\varepsilon}}{A^{+\varepsilon},\cdot}{B^{+\varepsilon},\cdot}}{\FormulaRefAuto{SemigroupHomDef}[def]{11,12,15,31}}
 \proofstepwidestar[1,8]{f^{+\varepsilon}:A^{+\varepsilon}\cong B^{+\varepsilon}}{\FormulaRefAuto{SemigroupIsoDef}[def]{32,23}}
\closeproofpartwide
\end{tabproofsplitwide}

\subsection{Kongruenzen und Quotienten}

\FormulaDefDeltaK[Kongruenz und Quotient einer Halbgruppe]{%
 \begin{aligned}[t]
 &\rho\in\operatorname{Con}(A)\coloneqq
    \EqRel{A,\rho}\land{}\\[-2pt]
 &\hspace{22pt}\forall a,a',b,b'\in A\;
    (a\mathrel\rho a'\land b\mathrel\rho b'\Rightarrow
                    ab\mathrel\rho a'b');\\[-2pt]
 &[a]_\rho\coloneqq\EqClass{a}{\rho}{A},\qquad A/\rho\coloneqq\{[a]_\rho\mid a\in A\},
       \qquad [a]_\rho[b]_\rho\coloneqq[ab]_\rho.
 \end{aligned}%
}{SemigroupCongruenceQuotientDef}{%
 \DeltaRow{Voraussetzung}{\begin{aligned}[t]&\SemigroupAlg{A}{\star};\\[-2pt]
     &\text{für den Quotienten zusätzlich }\rho\in\operatorname{Con}(A)\end{aligned}}
 \DeltaRow{Mengen}{A\dsep a\dsep a'\dsep b\dsep b'}
 \DeltaRow{Relationen}{\rho\subseteq A\times A}
}

\FormulaDefDeltaK[Transport einer beliebigen Relation]{%
 \begin{aligned}[t]
 &f:A\to B\dsep\rho\subseteq A\times A\vdash{}\\[-2pt]
 &f_*\rho\coloneqq\{(f(a),f(b))\mid a,b\in A,\ a\mathrel\rho b\}.
 \end{aligned}%
}{SemigroupTransportedRelationDef}{%
 \DeltaRow{Mengen}{A\dsep B\dsep a\dsep b}
 \DeltaRow{Funktionen}{f}
 \DeltaRow{Relationen}{\rho}
}

\FormulaThmDeltaKR[Transport von Kongruenzen und Quotienten]{%
 \begin{aligned}[t]
  &f:A\cong B\dsep\rho\in\operatorname{Con}(A)\vdash{}\\[-2pt]
  &\text{(i)}\quad f_*\rho\in\operatorname{Con}(B)\\[-2pt]
  &\text{(ii)}\quad \bar f:A/\rho\cong B/(f_*\rho),\quad \bar f([a]_\rho)\coloneqq[f(a)]_{f_*\rho}
 \end{aligned}%
}{%
 \begin{aligned}[t]
 &f:A\cong B\dsep\rho\in\operatorname{Con}(A)\vdash{}\\[-2pt]
 &f_*\rho\in\operatorname{Con}(B),\\[-2pt]
 &\bar f:A/\rho\cong B/(f_*\rho),
       \qquad\bar f([a]_\rho)\coloneqq[f(a)]_{f_*\rho}.
 \end{aligned}%
}{SemigroupIsoCongruenceQuotients}{%
 \DeltaRow{Mengen}{A\dsep B\dsep a\dsep b\dsep c\dsep d}
 \DeltaRow{Funktionen}{f\dsep h\dsep\bar f\dsep\star\dsep\diamond}
 \DeltaRow{Relationen}{\rho\dsep f_*\rho}
}
\begin{tabproofsplitwide}
\proofpartwideindR[Kriterium für das Relationsbild]{%
\begin{aligned}[t]
 &f:A\cong B\dsep\rho\subseteq A\times A\vdash{}\\[-2pt]
 &\text{(i)}\quad \forall u,v\in B\;(u\mathrel{f_*\rho}v\leftrightarrow h(u)\mathrel\rho h(v))\\[-2pt]
 &\text{(ii)}\quad \forall a,b\in A\;(a\mathrel\rho b\leftrightarrow f(a)\mathrel{f_*\rho}f(b))
\end{aligned}%
}{f:A\cong B\dsep\rho\subseteq A\times A\vdash\begin{aligned}[t]&\forall u,v\in B\;(u\mathrel{f_*\rho}v\leftrightarrow h(u)\mathrel\rho h(v)),\\&\forall a,b\in A\;(a\mathrel\rho b\leftrightarrow f(a)\mathrel{f_*\rho}f(b))\end{aligned}}[SemigroupTransportedRelationCriterion]
 \proofstepwidestar[1]{f:A\cong B}{\rA}
 \proofstepwidestar[1]{\SemigroupAlg{A}{\star}}{\FormulaRefAuto{SemigroupIsoSource}[thm]{1}}
 \proofstepwidestar[1]{\SemigroupAlg{B}{\diamond}}{\FormulaRefAuto{SemigroupIsoTarget}[thm]{1}}
 \proofstepwidestar[1]{f:A\to B}{\FormulaRefAuto{SemigroupIsoFunction}[thm]{1}}
 \proofstepwidestar[1]{f:A\bij B}{\FormulaRefAuto{SemigroupIsoBijectiveAxiom}[ax]{1}}
 \proofstepwidestar[1]{\SemigroupHom{f}{A,\star}{B,\diamond}}{\FormulaRefAuto{SemigroupIsoHomAxiom}[ax]{1}}
 \proofstepwidestar[1]{f[A]=B}{\FormulaRefAuto{SemigroupIsoFullImage}[thm]{1}}
 \proofstepwidestar[8]{\rho\subseteq A\times A}{\rA}
 \proofstepwidestar[1]{h:B\bij A}{\FormulaRefAuto{SemigroupIsoInverseCarrier}[thm]{1}}
 \proofstepwidestar[1]{h:B\to A}{\FormulaRefAuto{Bijektive Funktion}[def]{9}}
 \proofstepwidestar[11]{u\in B}{\rA}
 \proofstepwidestar[12]{v\in B}{\rA}
 \proofstepwidestar[11,12]{u,v\in B}{\rAIStar{11,12}}
 \proofstepwidestar[1,11,12]{h(u),h(v)\in A}{\FormulaRefAuto{F\colon A\to B\dsep x\in A\vdash F(x)\in B}[thm]{10,13}}
 \proofstepwidestar[1,11,12]{f(h(u))=u\,,\quad f(h(v))=v}{\rAI{\FormulaRefAuto{InverseFunctionRightInverse}[thm]{5,\rAEa{13}},\FormulaRefAuto{InverseFunctionRightInverse}[thm]{5,\rAEb{13}}}}
 \proofstepwidestar[1,8,11,12]{u\mathrel{f_*\rho}v\leftrightarrow h(u)\mathrel\rho h(v)}{\text{\parbox[t]{\linewidth}{\raggedright Definition des Relationsbildes: vorwärts eindeutige Urbilder, rückwärts die Zeugen aus 14,15.}}}
 \proofstepwidestar[1,8]{\forall u,v\in B\;(u\mathrel{f_*\rho}v\leftrightarrow h(u)\mathrel\rho h(v))}{\rUI{\rRI{11,\rUI{\rRI{12,16}}}}}
 \proofstepwidestar[18]{a\in A}{\rA}
 \proofstepwidestar[19]{b\in A}{\rA}
 \proofstepwidestar[18,19]{a,b\in A}{\rAIStar{18,19}}
 \proofstepwidestar[1,18,19]{f(a),f(b)\in B}{\FormulaRefAuto{F\colon A\to B\dsep x\in A\vdash F(x)\in B}[thm]{4,20}}
 \proofstepwidestar[1,18,19]{h(f(a))=a\,,\quad h(f(b))=b}{\rAI{\FormulaRefAuto{InverseFunctionLeftInverse}[thm]{5,\rAEa{20}},\FormulaRefAuto{InverseFunctionLeftInverse}[thm]{5,\rAEb{20}}}}
 \proofstepwidestar[1,8,18,19]{a\mathrel\rho b\leftrightarrow f(a)\mathrel{f_*\rho}f(b)}{\text{\parbox[t]{\linewidth}{\raggedright All-Elimination in 17 mit 21; Gleichheitseinsetzung 22}}}
 \proofstepwidestar[1,8]{\forall a,b\in A\;(a\mathrel\rho b\leftrightarrow f(a)\mathrel{f_*\rho}f(b))}{\rUI{\rRI{18,\rUI{\rRI{19,23}}}}}
\closeproofpartwideindR
\proofpartwideindR[Die transportierte Kongruenz]{%
f:A\cong B\dsep\rho\in\operatorname{Con}(A)\vdash f_*\rho\in\operatorname{Con}(B)%
}{f:A\cong B\dsep\rho\in\operatorname{Con}(A)\vdash f_*\rho\in\operatorname{Con}(B)}[SemigroupTransportedCongruence]
 \proofstepwidestar[1]{f:A\cong B}{\rA}
 \proofstepwidestar[1]{\SemigroupAlg{A}{\star}}{\FormulaRefAuto{SemigroupIsoSource}[thm]{1}}
 \proofstepwidestar[1]{\SemigroupAlg{B}{\diamond}}{\FormulaRefAuto{SemigroupIsoTarget}[thm]{1}}
 \proofstepwidestar[1]{f:A\to B}{\FormulaRefAuto{SemigroupIsoFunction}[thm]{1}}
 \proofstepwidestar[1]{f:A\bij B}{\FormulaRefAuto{SemigroupIsoBijectiveAxiom}[ax]{1}}
 \proofstepwidestar[1]{\SemigroupHom{f}{A,\star}{B,\diamond}}{\FormulaRefAuto{SemigroupIsoHomAxiom}[ax]{1}}
 \proofstepwidestar[1]{f[A]=B}{\FormulaRefAuto{SemigroupIsoFullImage}[thm]{1}}
 \proofstepwidestar[8]{\rho\in\operatorname{Con}(A)}{\rA}
 \proofstepwidestar[1,8]{\EqRel{A,\rho}}{\FormulaRefAuto{SemigroupCongruenceQuotientDef}[def]{2,8}}
 \proofstepwidestar[1,8]{\rho\subseteq A\times A}{\FormulaRefAuto{EqRelDef}[def]{9}}
 \proofstepwidestar[1]{h:B\cong A}{\FormulaRefAuto{SemigroupIsoInverse}[thm]{1}}
 \proofstepwidestar[1]{h:B\to A}{\FormulaRefAuto{SemigroupIsoFunction}[thm]{11}}
 \proofstepwidestar[1,8]{\begin{aligned}[t]&\forall u,v\in B\;(u\mathrel{f_*\rho}v\leftrightarrow h(u)\mathrel\rho h(v)),\\&\forall a,b\in A\;(a\mathrel\rho b\leftrightarrow f(a)\mathrel{f_*\rho}f(b))\end{aligned}}{\FormulaRefAuto{SemigroupTransportedRelationCriterion}[thm]{1,10}}
 \proofstepwidestar[1,8]{f_*\rho\subseteq B\times B}{\text{\parbox[t]{\linewidth}{\raggedright Definition des Relationsbildes und 4,10}}}
 \proofstepwidestar[15]{u\in B}{\rA}
 \proofstepwidestar[1,15]{h(u)\in A}{\FormulaRefAuto{F\colon A\to B\dsep x\in A\vdash F(x)\in B}[thm]{12,15}}
 \proofstepwidestar[1,8,15]{h(u)\mathrel\rho h(u)}{\FormulaRefAuto{EqRelReflexivity}[ax]{9,16}}
 \proofstepwidestar[1,8,15]{u\mathrel{f_*\rho}u}{13,17}
 \proofstepwidestar[1,8]{\forall u\in B\;(u\mathrel{f_*\rho}u)}{\rUI{\rRI{15,18}}}
 \proofstepwidestar[20]{u\in B}{\rA}
 \proofstepwidestar[21]{v\in B}{\rA}
 \proofstepwidestar[22]{u\mathrel{f_*\rho}v}{\rA}
 \proofstepwidestar[20,21,22]{u,v\in B\,,\quad u\mathrel{f_*\rho}v}{\rAIStar{20,21,22}}
 \proofstepwidestar[1,8,20,21,22]{h(u),h(v)\in A\,,\quad h(u)\mathrel\rho h(v)}{12,13,23}
 \proofstepwidestar[1,8,20,21,22]{h(v)\mathrel\rho h(u)}{\FormulaRefAuto{EqRelSymmetry}[ax]{9,24}}
 \proofstepwidestar[1,8,20,21,22]{v\mathrel{f_*\rho}u}{13,25}
 \proofstepwidestar[1,8]{\forall u,v\in B\;(u\mathrel{f_*\rho}v\Rightarrow v\mathrel{f_*\rho}u)}{\rUI{\rRI{20,\rUI{\rRI{21,\rRI{22,26}}}}}}
 \proofstepwidestar[28]{u\in B}{\rA}
 \proofstepwidestar[29]{v\in B}{\rA}
 \proofstepwidestar[30]{w\in B}{\rA}
 \proofstepwidestar[31]{u\mathrel{f_*\rho}v\land v\mathrel{f_*\rho}w}{\rA}
 \proofstepwidestar[28,29,30,31]{u,v,w\in B\,,\quad u\mathrel{f_*\rho}v\,,\quad v\mathrel{f_*\rho}w}{\rAIStar{28,29,30,31}}
 \proofstepwidestar[1,8,28,29,30,31]{h(u),h(v),h(w)\in A\,,\quad h(u)\mathrel\rho h(v)\,,\quad h(v)\mathrel\rho h(w)}{12,13,32}
 \proofstepwidestar[1,8,28,29,30,31]{h(u)\mathrel\rho h(w)}{\FormulaRefAuto{EqRelTransitivity}[ax]{9,33}}
 \proofstepwidestar[1,8,28,29,30,31]{u\mathrel{f_*\rho}w}{13,34}
 \proofstepwidestar[1,8]{\forall u,v,w\in B\;(u\mathrel{f_*\rho}v\land v\mathrel{f_*\rho}w\Rightarrow u\mathrel{f_*\rho}w)}{\rUI{\rRI{28,\rUI{\rRI{29,\rUI{\rRI{30,\rRI{31,35}}}}}}}}
 \proofstepwidestar[1,8]{\EqRel{B,f_*\rho}}{\FormulaRefAuto{EqRelDef}[def]{14,19,27,36}}
 \proofstepwidestar[38]{u\in B}{\rA}
 \proofstepwidestar[39]{u'\in B}{\rA}
 \proofstepwidestar[40]{v\in B}{\rA}
 \proofstepwidestar[41]{v'\in B}{\rA}
 \proofstepwidestar[42]{u\mathrel{f_*\rho}u'\land v\mathrel{f_*\rho}v'}{\rA}
 \proofstepwidestar[38,39,40,41,42]{u,u',v,v'\in B\,,\quad u\mathrel{f_*\rho}u'\,,\quad v\mathrel{f_*\rho}v'}{\rAIStar{38,39,40,41,42}}
 \proofstepwidestar[1,8,38,39,40,41,42]{h(u),h(u'),h(v),h(v')\in A\,,\quad h(u)\mathrel\rho h(u')\,,\quad h(v)\mathrel\rho h(v')}{12,13,43}
 \proofstepwidestar[1,8,38,39,40,41,42]{h(u)h(v)\mathrel\rho h(u')h(v')}{\FormulaRefAuto{SemigroupCongruenceQuotientDef}[def]{8,44}}
 \proofstepwidestar[1,38,39,40,41,42]{h(uv)=h(u)h(v)\,,\quad h(u' v')=h(u')h(v')}{\FormulaRefAuto{SemigroupIsoOperation}[thm]{11,43}}
 \proofstepwidestar[1,8,38,39,40,41,42]{h(uv)\mathrel\rho h(u' v')}{\rIE{46,45}}
 \proofstepwidestar[1,38,39,40,41,42]{uv,u' v'\in B}{\FormulaRefAuto{SemigroupClosureAxiom}[ax]{3,43}}
 \proofstepwidestar[1,8,38,39,40,41,42]{uv\mathrel{f_*\rho}u' v'}{13,48,47}
 \proofstepwidestar[1,8]{\forall u,u',v,v'\in B\;(u\mathrel{f_*\rho}u'\land v\mathrel{f_*\rho}v'\Rightarrow uv\mathrel{f_*\rho}u' v')}{\rUI{\rRI{38,\rUI{\rRI{39,\rUI{\rRI{40,\rUI{\rRI{41,\rRI{42,49}}}}}}}}}}
 \proofstepwidestar[1,8]{f_*\rho\in\operatorname{Con}(B)}{\FormulaRefAuto{SemigroupCongruenceQuotientDef}[def]{37,50}}
\closeproofpartwideindR
\proofpartwideindR[Wohldefiniertheit der Quotientenoperationen]{%
\SemigroupAlg{A}{\star}\dsep\rho\in\operatorname{Con}(A)\vdash\SemigroupAlg{A/\rho}{\cdot}%
}{\SemigroupAlg{A}{\star}\dsep\rho\in\operatorname{Con}(A)\vdash\SemigroupAlg{A/\rho}{\cdot}}[SemigroupQuotientIsSemigroup]
  \proofstepwidestar[1]{\SemigroupAlg{A}{\star}}{\rA}
  \proofstepwidestar[2]{\rho\in\operatorname{Con}(A)}{\rA}
  \proofstepwidestar[1,2]{\EqRel{A,\rho}}{\FormulaRefAuto{SemigroupCongruenceQuotientDef}[def]{1,2}}
  \proofstepwidestar[4]{a\in A}{\rA}
  \proofstepwidestar[5]{a'\in A}{\rA}
  \proofstepwidestar[6]{b\in A}{\rA}
  \proofstepwidestar[7]{b'\in A}{\rA}
  \proofstepwidestar[8]{[a]_\rho=[a']_\rho\land[b]_\rho=[b']_\rho}{\rA}
  \proofstepwidestar[8]{[a]_\rho=[a']_\rho}{\rAEa{8}}
  \proofstepwidestar[8]{[b]_\rho=[b']_\rho}{\rAEb{8}}
  \proofstepwidestar[1,2,4,5,8]{a\mathrel\rho a'}{\FormulaRefAuto{EqClassRelFromEqual}[thm]{3,4,5,9}}
  \proofstepwidestar[1,2,6,7,8]{b\mathrel\rho b'}{\FormulaRefAuto{EqClassRelFromEqual}[thm]{3,6,7,10}}
  \proofstepwidestar[1,2,4,5,6,7,8]{ab\mathrel\rho a'b'}{\FormulaRefAuto{SemigroupCongruenceQuotientDef}[def]{1,2,4,5,6,7,11,12}}
  \proofstepwidestar[1,4,6]{ab\in A}{\FormulaRefAuto{SemigroupClosureAxiom}[ax]{1,4,6}}
  \proofstepwidestar[1,5,7]{a'b'\in A}{\FormulaRefAuto{SemigroupClosureAxiom}[ax]{1,5,7}}
  \proofstepwidestar[1,2,4,5,6,7,8]{[ab]_\rho=[a'b']_\rho}{\FormulaRefAuto{EqClassEqualFromRel}[thm]{3,14,15,13}}
  \proofstepwidestar[1,2]{\forall a,a',b,b'\in A\;([a]_\rho=[a']_\rho\land[b]_\rho=[b']_\rho\Rightarrow[ab]_\rho=[a'b']_\rho)}{\rUI{\rRI{4,\rUI{\rRI{5,\rUI{\rRI{6,\rUI{\rRI{7,\rRI{8,16}}}}}}}}}}
  \proofstepwidestar[18]{X\in A/\rho}{\rA}
  \proofstepwidestar[1,2,18]{\exists a\in A\;(X=[a]_\rho)}{\FormulaRefAuto{QuotientRepresentative}[thm]{3,18}}
  \proofstepwidestar[20]{Y\in A/\rho}{\rA}
  \proofstepwidestar[1,2,20]{\exists b\in A\;(Y=[b]_\rho)}{\FormulaRefAuto{QuotientRepresentative}[thm]{3,20}}
  \proofstepwidestar[22]{a\in A\land X=[a]_\rho}{\rA}
  \proofstepwidestar[22]{a\in A}{\rAEa{22}}
  \proofstepwidestar[22]{X=[a]_\rho}{\rAEb{22}}
  \proofstepwidestar[22]{[a]_\rho=X}{\FormulaRefAuto{a=b\vdash b=a}[thm]{24}}
  \proofstepwidestar[26]{b\in A\land Y=[b]_\rho}{\rA}
  \proofstepwidestar[26]{b\in A}{\rAEa{26}}
  \proofstepwidestar[26]{Y=[b]_\rho}{\rAEb{26}}
  \proofstepwidestar[26]{[b]_\rho=Y}{\FormulaRefAuto{a=b\vdash b=a}[thm]{28}}
  \proofstepwidestar[1,22,26]{ab\in A}{\FormulaRefAuto{SemigroupClosureAxiom}[ax]{1,23,27}}
  \proofstepwidestar[1,2,22,26]{[ab]_\rho\in A/\rho}{\FormulaRefAuto{EqClassInQuotient}[thm]{3,30}}
  \proofstepwidestar[1,2,22,26]{[a]_\rho[b]_\rho=[ab]_\rho}{\FormulaRefAuto{SemigroupCongruenceQuotientDef}[def]{1,2,17,23,27}}
  \proofstepwidestar[1,2,22,26]{[ab]_\rho=[a]_\rho[b]_\rho}{\FormulaRefAuto{a=b\vdash b=a}[thm]{32}}
  \proofstepwidestar[1,2,22,26]{XY\in A/\rho}{\rIE{29,\rIE{25,\rIE{33,31}}}}
  \proofstepwidestar[1,2,20,22]{XY\in A/\rho}{\rEE{21,26,34}}
  \proofstepwidestar[1,2,18,20]{XY\in A/\rho}{\rEE{19,22,35}}
  \proofstepwidestar[1,2]{\forall X,Y\in A/\rho\;(XY\in A/\rho)}{\rUI{\rRI{18,\rUI{\rRI{20,36}}}}}
  \proofstepwidestar[38]{X\in A/\rho}{\rA}
  \proofstepwidestar[1,2,38]{\exists a\in A\;(X=[a]_\rho)}{\FormulaRefAuto{QuotientRepresentative}[thm]{3,38}}
  \proofstepwidestar[40]{Y\in A/\rho}{\rA}
  \proofstepwidestar[1,2,40]{\exists b\in A\;(Y=[b]_\rho)}{\FormulaRefAuto{QuotientRepresentative}[thm]{3,40}}
  \proofstepwidestar[42]{Z\in A/\rho}{\rA}
  \proofstepwidestar[1,2,42]{\exists c\in A\;(Z=[c]_\rho)}{\FormulaRefAuto{QuotientRepresentative}[thm]{3,42}}
  \proofstepwidestar[44]{a\in A\land X=[a]_\rho}{\rA}
  \proofstepwidestar[44]{a\in A}{\rAEa{44}}
  \proofstepwidestar[44]{X=[a]_\rho}{\rAEb{44}}
  \proofstepwidestar[44]{[a]_\rho=X}{\FormulaRefAuto{a=b\vdash b=a}[thm]{46}}
  \proofstepwidestar[48]{b\in A\land Y=[b]_\rho}{\rA}
  \proofstepwidestar[48]{b\in A}{\rAEa{48}}
  \proofstepwidestar[48]{Y=[b]_\rho}{\rAEb{48}}
  \proofstepwidestar[48]{[b]_\rho=Y}{\FormulaRefAuto{a=b\vdash b=a}[thm]{50}}
  \proofstepwidestar[52]{c\in A\land Z=[c]_\rho}{\rA}
  \proofstepwidestar[52]{c\in A}{\rAEa{52}}
  \proofstepwidestar[52]{Z=[c]_\rho}{\rAEb{52}}
  \proofstepwidestar[52]{[c]_\rho=Z}{\FormulaRefAuto{a=b\vdash b=a}[thm]{54}}
  \proofstepwidestar[1,44,48]{ab\in A}{\FormulaRefAuto{SemigroupClosureAxiom}[ax]{1,45,49}}
  \proofstepwidestar[1,48,52]{bc\in A}{\FormulaRefAuto{SemigroupClosureAxiom}[ax]{1,49,53}}
  \proofstepwidestar[1,44,48,52]{(ab)c=a(bc)}{\FormulaRefAuto{SemigroupAssociativityAxiom}[ax]{1,45,49,53}}
  \proofstepwidestar[1,2,44,48]{[a]_\rho[b]_\rho=[ab]_\rho}{\FormulaRefAuto{SemigroupCongruenceQuotientDef}[def]{1,2,17,45,49}}
  \proofstepwidestar[1,2,44,48,52]{[ab]_\rho[c]_\rho=[(ab)c]_\rho}{\FormulaRefAuto{SemigroupCongruenceQuotientDef}[def]{1,2,17,56,53}}
  \proofstepwidestar[1,2,48,52]{[b]_\rho[c]_\rho=[bc]_\rho}{\FormulaRefAuto{SemigroupCongruenceQuotientDef}[def]{1,2,17,49,53}}
  \proofstepwidestar[1,2,44,48,52]{[a]_\rho[bc]_\rho=[a(bc)]_\rho}{\FormulaRefAuto{SemigroupCongruenceQuotientDef}[def]{1,2,17,45,57}}
  \proofstepwidestar[1,2,44,48,52]{[a(bc)]_\rho=[a]_\rho[bc]_\rho}{\FormulaRefAuto{a=b\vdash b=a}[thm]{62}}
  \proofstepwidestar[1,2,48,52]{[bc]_\rho=[b]_\rho[c]_\rho}{\FormulaRefAuto{a=b\vdash b=a}[thm]{61}}
  \proofstepwidestar[1,2,44,48]{([a]_\rho[b]_\rho)[c]_\rho=[ab]_\rho[c]_\rho}{\rIE{59,\rII}}
  \proofstepwidestar[1,2,44,48,52]{([a]_\rho[b]_\rho)[c]_\rho=[(ab)c]_\rho}{\rIE{60,65}}
  \proofstepwidestar[1,2,44,48,52]{([a]_\rho[b]_\rho)[c]_\rho=[a(bc)]_\rho}{\rIE{58,66}}
  \proofstepwidestar[1,2,44,48,52]{([a]_\rho[b]_\rho)[c]_\rho=[a]_\rho([b]_\rho[c]_\rho)}{\rIE{64,\rIE{63,67}}}
  \proofstepwidestar[1,2,44,48,52]{(XY)Z=X(YZ)}{\rIE{55,\rIE{51,\rIE{47,68}}}}
  \proofstepwidestar[1,2,42,44,48]{(XY)Z=X(YZ)}{\rEE{43,52,69}}
  \proofstepwidestar[1,2,40,42,44]{(XY)Z=X(YZ)}{\rEE{41,48,70}}
  \proofstepwidestar[1,2,38,40,42]{(XY)Z=X(YZ)}{\rEE{39,44,71}}
  \proofstepwidestar[1,2]{\forall X,Y,Z\in A/\rho\;((XY)Z=X(YZ))}{\rUI{\rRI{38,\rUI{\rRI{40,\rUI{\rRI{42,72}}}}}}}
  \proofstepwidestar[1,2]{\SemigroupAlg{A/\rho}{\cdot}}{\FormulaRefAuto{SemigroupDef}[def]{17,37,73}}
\closeproofpartwideindR
\proofpartwideindR[Bijektive Klassenabbildung]{%
f:A\cong B\dsep\rho\in\operatorname{Con}(A)\vdash\bar f:A/\rho\bij B/(f_*\rho)%
}{f:A\cong B\dsep\rho\in\operatorname{Con}(A)\vdash\bar f:A/\rho\bij B/(f_*\rho)}[SemigroupIsoQuotientBijection]
 \proofstepwidestar[1]{f:A\cong B}{\rA}
 \proofstepwidestar[1]{\SemigroupAlg{A}{\star}}{\FormulaRefAuto{SemigroupIsoSource}[thm]{1}}
 \proofstepwidestar[1]{\SemigroupAlg{B}{\diamond}}{\FormulaRefAuto{SemigroupIsoTarget}[thm]{1}}
 \proofstepwidestar[1]{f:A\to B}{\FormulaRefAuto{SemigroupIsoFunction}[thm]{1}}
 \proofstepwidestar[1]{f:A\bij B}{\FormulaRefAuto{SemigroupIsoBijectiveAxiom}[ax]{1}}
 \proofstepwidestar[1]{\SemigroupHom{f}{A,\star}{B,\diamond}}{\FormulaRefAuto{SemigroupIsoHomAxiom}[ax]{1}}
 \proofstepwidestar[1]{f[A]=B}{\FormulaRefAuto{SemigroupIsoFullImage}[thm]{1}}
 \proofstepwidestar[8]{\rho\in\operatorname{Con}(A)}{\rA}
 \proofstepwidestar[1,8]{\EqRel{A,\rho}}{\FormulaRefAuto{SemigroupCongruenceQuotientDef}[def]{2,8}}
 \proofstepwidestar[1,8]{\rho\subseteq A\times A}{\FormulaRefAuto{EqRelDef}[def]{9}}
 \proofstepwidestar[1,8]{f_*\rho\in\operatorname{Con}(B)}{\FormulaRefAuto{SemigroupTransportedCongruence}[thm]{1,8}}
 \proofstepwidestar[1,8]{\EqRel{B,f_*\rho}}{\FormulaRefAuto{SemigroupCongruenceQuotientDef}[def]{3,11}}
 \proofstepwidestar[1,8]{\begin{aligned}[t]&\forall u,v\in B\;(u\mathrel{f_*\rho}v\leftrightarrow h(u)\mathrel\rho h(v)),\\&\forall a,b\in A\;(a\mathrel\rho b\leftrightarrow f(a)\mathrel{f_*\rho}f(b))\end{aligned}}{\FormulaRefAuto{SemigroupTransportedRelationCriterion}[thm]{1,10}}
 \proofstepwidestar[14]{a\in A}{\rA}
 \proofstepwidestar[15]{b\in A}{\rA}
 \proofstepwidestar[14,15]{a,b\in A}{\rAIStar{14,15}}
 \proofstepwidestar[1,14,15]{f(a),f(b)\in B}{\FormulaRefAuto{F\colon A\to B\dsep x\in A\vdash F(x)\in B}[thm]{4,16}}
 \proofstepwidestar[1,8,14,15]{[a]_\rho=[b]_\rho\leftrightarrow a\mathrel\rho b}{\FormulaRefAuto{EqClassEqualIffRel}[thm]{9,16}}
 \proofstepwidestar[1,8,14,15]{a\mathrel\rho b\leftrightarrow f(a)\mathrel{f_*\rho}f(b)}{\text{\parbox[t]{\linewidth}{\raggedright All-Elimination in 13 mit 16}}}
 \proofstepwidestar[1,8,14,15]{f(a)\mathrel{f_*\rho}f(b)\leftrightarrow[f(a)]_{f_*\rho}=[f(b)]_{f_*\rho}}{\FormulaRefAuto{EqClassEqualIffRel}[thm]{12,17}}
 \proofstepwidestar[1,8,14,15]{[a]_\rho=[b]_\rho\leftrightarrow[f(a)]_{f_*\rho}=[f(b)]_{f_*\rho}}{\text{\parbox[t]{\linewidth}{\raggedright Äquivalenztransitivität aus 18,19,20}}}
 \proofstepwidestar[1,8]{\forall a,b\in A\;([a]_\rho=[b]_\rho\leftrightarrow[f(a)]_{f_*\rho}=[f(b)]_{f_*\rho})}{\rUI{\rRI{14,\rUI{\rRI{15,21}}}}}
 \proofstepwidestar[1,8]{\bar f:A/\rho\to B/(f_*\rho)}{\text{\parbox[t]{\linewidth}{\raggedright Definition der Klassenabbildung: Vertreterexistenz; Wohldefiniertheit aus der Vorwärtsrichtung von 22.}}}
 \proofstepwidestar[1,8]{\bar f:A/\rho\inj B/(f_*\rho)}{\text{\parbox[t]{\linewidth}{\raggedright Injektivitätsdefinition und Rückwärtsrichtung von 22}}}
 \proofstepwidestar[1]{h:B\bij A}{\FormulaRefAuto{SemigroupIsoInverseCarrier}[thm]{1}}
 \proofstepwidestar[1]{h:B\to A}{\FormulaRefAuto{Bijektive Funktion}[def]{25}}
 \proofstepwidestar[27]{u\in B}{\rA}
 \proofstepwidestar[1,27]{h(u)\in A}{\FormulaRefAuto{F\colon A\to B\dsep x\in A\vdash F(x)\in B}[thm]{26,27}}
 \proofstepwidestar[1,27]{f(h(u))=u}{\FormulaRefAuto{InverseFunctionRightInverse}[thm]{5,27}}
 \proofstepwidestar[1,8,27]{[h(u)]_\rho\in A/\rho\,,\quad\bar f([h(u)]_\rho)=[u]_{f_*\rho}}{\text{\parbox[t]{\linewidth}{\raggedright \FormulaRefAuto{SemigroupCongruenceQuotientDef}[def]{28,29}; Definition von \(\bar f\).}}}
 \proofstepwidestar[1,8]{\bar f:A/\rho\sur B/(f_*\rho)}{\text{\parbox[t]{\linewidth}{\raggedright Jede Zielklasse hat einen Vertreter; 27--30 liefern ein Urbild.}}}
 \proofstepwidestar[1,8]{\bar f:A/\rho\bij B/(f_*\rho)}{\FormulaRefAuto{Bijektive Funktion}[def]{23,24,31}}
\closeproofpartwideindR
\proofpartwide{Operationserhaltung}
  \proofstepwidestar[1]{f:A\cong B}{\rA}
  \proofstepwidestar[1]{\SemigroupAlg{A}{\star}}{\FormulaRefAuto{SemigroupIsoSource}[thm]{1}}
  \proofstepwidestar[1]{\SemigroupAlg{B}{\diamond}}{\FormulaRefAuto{SemigroupIsoTarget}[thm]{1}}
  \proofstepwidestar[1]{f:A\to B}{\FormulaRefAuto{SemigroupIsoFunction}[thm]{1}}
  \proofstepwidestar[5]{\rho\in\operatorname{Con}(A)}{\rA}
  \proofstepwidestar[1,5]{f_*\rho\in\operatorname{Con}(B)}{\FormulaRefAuto{SemigroupTransportedCongruence}[thm]{1,5}}
  \proofstepwidestar[1,5]{\EqRel{A,\rho}}{\FormulaRefAuto{SemigroupCongruenceQuotientDef}[def]{2,5}}
  \proofstepwidestar[1,5]{\SemigroupAlg{A/\rho}{\cdot}}{\FormulaRefAuto{SemigroupQuotientIsSemigroup}[thm]{2,5}}
  \proofstepwidestar[1,5]{\SemigroupAlg{B/(f_*\rho)}{\cdot}}{\FormulaRefAuto{SemigroupQuotientIsSemigroup}[thm]{3,6}}
  \proofstepwidestar[1,5]{\bar f:A/\rho\bij B/(f_*\rho)}{\FormulaRefAuto{SemigroupIsoQuotientBijection}[thm]{1,5}}
  \proofstepwidestar[1,5]{\bar f:A/\rho\to B/(f_*\rho)}{\FormulaRefAuto{Bijektive Funktion}[def]{10}}
  \proofstepwidestar[12]{X\in A/\rho}{\rA}
  \proofstepwidestar[1,5,12]{\exists a\in A\;(X=[a]_\rho)}{\FormulaRefAuto{QuotientRepresentative}[thm]{7,12}}
  \proofstepwidestar[14]{Y\in A/\rho}{\rA}
  \proofstepwidestar[1,5,14]{\exists b\in A\;(Y=[b]_\rho)}{\FormulaRefAuto{QuotientRepresentative}[thm]{7,14}}
  \proofstepwidestar[16]{a\in A\land X=[a]_\rho}{\rA}
  \proofstepwidestar[16]{a\in A}{\rAEa{16}}
  \proofstepwidestar[16]{X=[a]_\rho}{\rAEb{16}}
  \proofstepwidestar[16]{[a]_\rho=X}{\FormulaRefAuto{a=b\vdash b=a}[thm]{18}}
  \proofstepwidestar[20]{b\in A\land Y=[b]_\rho}{\rA}
  \proofstepwidestar[20]{b\in A}{\rAEa{20}}
  \proofstepwidestar[20]{Y=[b]_\rho}{\rAEb{20}}
  \proofstepwidestar[20]{[b]_\rho=Y}{\FormulaRefAuto{a=b\vdash b=a}[thm]{22}}
  \proofstepwidestar[1,16,20]{ab\in A}{\FormulaRefAuto{SemigroupClosureAxiom}[ax]{2,17,21}}
  \proofstepwidestar[1,16]{f(a)\in B}{\FormulaRefAuto{F\colon A\to B\dsep x\in A\vdash F(x)\in B}[thm]{4,17}}
  \proofstepwidestar[1,20]{f(b)\in B}{\FormulaRefAuto{F\colon A\to B\dsep x\in A\vdash F(x)\in B}[thm]{4,21}}
  \proofstepwidestar[1,16,20]{f(ab)=f(a)f(b)}{\FormulaRefAuto{SemigroupIsoOperation}[thm]{1,17,21}}
  \proofstepwidestar[1,5,16,20]{[a]_\rho[b]_\rho=[ab]_\rho}{\FormulaRefAuto{SemigroupCongruenceQuotientDef}[def]{2,5,17,21}}
  \proofstepwidestar[1,5,16]{\bar f([a]_\rho)=[f(a)]_{f_*\rho}}{\text{\parbox[t]{\linewidth}{\raggedright Definitionsgleichung von \(\bar f\)}}}
  \proofstepwidestar[1,5,20]{\bar f([b]_\rho)=[f(b)]_{f_*\rho}}{\text{\parbox[t]{\linewidth}{\raggedright Definitionsgleichung von \(\bar f\)}}}
  \proofstepwidestar[1,5,16,20]{\bar f([ab]_\rho)=[f(ab)]_{f_*\rho}}{\text{\parbox[t]{\linewidth}{\raggedright Definitionsgleichung von \(\bar f\)}}}
  \proofstepwidestar[1,5,16,20]{[f(a)]_{f_*\rho}[f(b)]_{f_*\rho}=[f(a)f(b)]_{f_*\rho}}{\FormulaRefAuto{SemigroupCongruenceQuotientDef}[def]{3,6,25,26}}
  \proofstepwidestar[1,5,16,20]{[f(a)f(b)]_{f_*\rho}=[f(a)]_{f_*\rho}[f(b)]_{f_*\rho}}{\FormulaRefAuto{a=b\vdash b=a}[thm]{32}}
  \proofstepwidestar[1,5,16]{[f(a)]_{f_*\rho}=\bar f([a]_\rho)}{\FormulaRefAuto{a=b\vdash b=a}[thm]{29}}
  \proofstepwidestar[1,5,20]{[f(b)]_{f_*\rho}=\bar f([b]_\rho)}{\FormulaRefAuto{a=b\vdash b=a}[thm]{30}}
  \proofstepwidestar[1,5,16,20]{\bar f([a]_\rho[b]_\rho)=\bar f([ab]_\rho)}{\rIE{28,\rII}}
  \proofstepwidestar[1,5,16,20]{\bar f([a]_\rho[b]_\rho)=[f(ab)]_{f_*\rho}}{\rIE{31,36}}
  \proofstepwidestar[1,5,16,20]{\bar f([a]_\rho[b]_\rho)=[f(a)f(b)]_{f_*\rho}}{\rIE{27,37}}
  \proofstepwidestar[1,5,16,20]{\bar f([a]_\rho[b]_\rho)=\bar f([a]_\rho)\bar f([b]_\rho)}{\rIE{35,\rIE{34,\rIE{33,38}}}}
  \proofstepwidestar[1,5,16,20]{\bar f(XY)=\bar f(X)\bar f(Y)}{\rIE{23,\rIE{19,39}}}
  \proofstepwidestar[1,5,14,16]{\bar f(XY)=\bar f(X)\bar f(Y)}{\rEE{15,20,40}}
  \proofstepwidestar[1,5,12,14]{\bar f(XY)=\bar f(X)\bar f(Y)}{\rEE{13,16,41}}
  \proofstepwidestar[1,5]{\forall X,Y\in A/\rho\;(\bar f(XY)=\bar f(X)\bar f(Y))}{\rUI{\rRI{12,\rUI{\rRI{14,42}}}}}
  \proofstepwidestar[1,5]{\SemigroupHom{\bar f}{A/\rho,\cdot}{B/(f_*\rho),\cdot}}{\FormulaRefAuto{SemigroupHomDef}[def]{8,9,11,43}}
  \proofstepwidestar[1,5]{\bar f:A/\rho\cong B/(f_*\rho)}{\FormulaRefAuto{SemigroupIsoDef}[def]{44,10}}
\closeproofpartwide
\end{tabproofsplitwide}
Ist zusätzlich eine Kongruenz \(\sigma\) auf \(B\) vorgegeben, so ist
die Voraussetzung \(f_*\rho=\sigma\) genau
\(a\mathrel\rho b\leftrightarrow f(a)\mathrel\sigma f(b)\).
Unter dieser Voraussetzung lautet der Zielquotient \(B/\sigma\).

\FormulaDefDeltaK[Rees-Kongruenz und Rees-Quotient]{%
 \begin{aligned}[t]
 &\mathsf{Id}(I;A,\star)\vdash{}\\[-2pt]
 &a\mathrel{\rho_I}b\coloneqq
    (a,b\in A\land(a=b\lor(a\in I\land b\in I))),
       \qquad A/I\coloneqq A/\rho_I.
 \end{aligned}%
}{SemigroupReesQuotientDef}{%
 \DeltaRow{Mengen}{A\dsep I\dsep a\dsep b}
 \DeltaRow{Voraussetzung}{\SemigroupAlg{A}{\star}}
}

\FormulaThmDeltaK[Isomorphie von Rees-Quotienten]{%
 f:A\cong B\dsep\mathsf{Id}(I;A,\star)\vdash
 ([a]_{\rho_I}\mapsto[f(a)]_{\rho_{f[I]}}):A/I\cong B/f[I]%
}{SemigroupIsoReesQuotients}{%
 \DeltaRow{Mengen}{A\dsep B\dsep I\dsep a\dsep b\dsep c\dsep d}
 \DeltaRow{Funktionen}{f\dsep\star\dsep\diamond}
}
\begin{tabproofsplitwide}
\proofpartwideindR[Rees-Kongruenz]{%
\SemigroupAlg{A}{\star}\dsep\mathsf{Id}(I;A,\star)\vdash\rho_I\in\operatorname{Con}(A)%
}{\SemigroupAlg{A}{\star}\dsep\mathsf{Id}(I;A,\star)\vdash\rho_I\in\operatorname{Con}(A)}[SemigroupReesCongruence]
 \proofstepwidestar[1]{\SemigroupAlg{A}{\star}}{\rA}
 \proofstepwidestar[2]{\mathsf{Id}(I;A,\star)}{\rA}
 \proofstepwidestar[1,2]{I\subseteq A}{\FormulaRefAuto{SemigroupIdealSubset}[thm]{1,2}}
 \proofstepwidestar[1,2]{\rho_I\subseteq A\times A}{\FormulaRefAuto{SemigroupReesQuotientDef}[def]{1,2}}
 \proofstepwidestar[1,2]{\forall a\in A\;(a\mathrel{\rho_I}a)}{\text{\parbox[t]{\linewidth}{\raggedright \FormulaRefAuto{SemigroupReesQuotientDef}[def]; Gleichheitsfall.}}}
 \proofstepwidestar[1,2]{\forall a,b\in A\;(a\mathrel{\rho_I}b\Rightarrow b\mathrel{\rho_I}a)}{\text{\parbox[t]{\linewidth}{\raggedright \FormulaRefAuto{SemigroupReesQuotientDef}[def]; Symmetrie der Gleichheit und der beiden Mitgliedschaften.}}}
 \proofstepwidestar[1,2]{\forall a,b,c\in A\;(a\mathrel{\rho_I}b\land b\mathrel{\rho_I}c\Rightarrow a\mathrel{\rho_I}c)}{\text{\parbox[t]{\linewidth}{\raggedright \FormulaRefAuto{SemigroupReesQuotientDef}[def]; Gleichheitsfälle einsetzen, sonst liegen \(a,c\) in \(I\).}}}
 \proofstepwidestar[1,2]{\EqRel{A,\rho_I}}{\FormulaRefAuto{EqRelDef}[def]{4,5,6,7}}
 \proofstepwidestar[9]{a\in A}{\rA}
 \proofstepwidestar[10]{a'\in A}{\rA}
 \proofstepwidestar[11]{b\in A}{\rA}
 \proofstepwidestar[12]{b'\in A}{\rA}
 \proofstepwidestar[13]{a\mathrel{\rho_I}a'\land b\mathrel{\rho_I}b'}{\rA}
 \proofstepwidestar[9,10,11,12,13]{a,a',b,b'\in A\,,\quad a\mathrel{\rho_I}a'\,,\quad b\mathrel{\rho_I}b'}{\rAIStar{9,10,11,12,13}}
 \proofstepwidestar[1,9,10,11,12,13]{ab,a'b'\in A}{\FormulaRefAuto{SemigroupClosureAxiom}[ax]{1,14}}
 \proofstepwidestar[1,2,9,10,11,12,13]{(a=a'\land b=b')\lor(a,a'\in I)\lor(b,b'\in I)}{\text{\parbox[t]{\linewidth}{\raggedright \FormulaRefAuto{SemigroupReesQuotientDef}[def]{14}; Distributivität der Disjunktionen.}}}
 \proofstepwidestar[1,9,10,11,12,13]{a=a'\land b=b'\Rightarrow ab=a'b'}{\text{\parbox[t]{\linewidth}{\raggedright Gleichheitseinsetzung}}}
 \proofstepwidestar[1,2,9,10,11,12,13]{a,a'\in I\Rightarrow ab,a'b'\in I}{\FormulaRefAuto{SemigroupIdealRightAbsorption}[thm]{1,2,14}}
 \proofstepwidestar[1,2,9,10,11,12,13]{b,b'\in I\Rightarrow ab,a'b'\in I}{\FormulaRefAuto{SemigroupIdealLeftAbsorption}[thm]{1,2,14}}
 \proofstepwidestar[1,2,9,10,11,12,13]{ab\mathrel{\rho_I}a'b'}{\FormulaRefAuto{SemigroupReesQuotientDef}[def]{15,16,17,18,19}}
 \proofstepwidestar[1,2]{\forall a,a',b,b'\in A\;(a\mathrel{\rho_I}a'\land b\mathrel{\rho_I}b'\Rightarrow ab\mathrel{\rho_I}a'b')}{\rUI{\rRI{9,\rUI{\rRI{10,\rUI{\rRI{11,\rUI{\rRI{12,\rRI{13,20}}}}}}}}}}
 \proofstepwidestar[1,2]{\rho_I\in\operatorname{Con}(A)}{\FormulaRefAuto{SemigroupCongruenceQuotientDef}[def]{8,21}}
\closeproofpartwideindR
\proofpartwide{Transport der zusammengezogenen Klasse}
 \proofstepwidestar[1]{f:A\cong B}{\rA}
 \proofstepwidestar[1]{\SemigroupAlg{A}{\star}}{\FormulaRefAuto{SemigroupIsoSource}[thm]{1}}
 \proofstepwidestar[1]{\SemigroupAlg{B}{\diamond}}{\FormulaRefAuto{SemigroupIsoTarget}[thm]{1}}
 \proofstepwidestar[1]{f:A\to B}{\FormulaRefAuto{SemigroupIsoFunction}[thm]{1}}
 \proofstepwidestar[1]{f:A\bij B}{\FormulaRefAuto{SemigroupIsoBijectiveAxiom}[ax]{1}}
 \proofstepwidestar[1]{\SemigroupHom{f}{A,\star}{B,\diamond}}{\FormulaRefAuto{SemigroupIsoHomAxiom}[ax]{1}}
 \proofstepwidestar[1]{f[A]=B}{\FormulaRefAuto{SemigroupIsoFullImage}[thm]{1}}
 \proofstepwidestar[8]{\mathsf{Id}(I;A,\star)}{\rA}
 \proofstepwidestar[1,8]{I\subseteq A}{\FormulaRefAuto{SemigroupIdealSubset}[thm]{2,8}}
 \proofstepwidestar[1,8]{\mathsf{Id}(f[I];B,\diamond)}{\FormulaRefAuto{SemigroupIsoIdImage}[thm]{1,8}}
 \proofstepwidestar[1,8]{\rho_I\in\operatorname{Con}(A)}{\FormulaRefAuto{SemigroupReesCongruence}[thm]{2,8}}
 \proofstepwidestar[1,8]{\rho_{f[I]}\in\operatorname{Con}(B)}{\FormulaRefAuto{SemigroupReesCongruence}[thm]{3,10}}
 \proofstepwidestar[13]{a\in A}{\rA}
 \proofstepwidestar[14]{b\in A}{\rA}
 \proofstepwidestar[13,14]{a,b\in A}{\rAIStar{13,14}}
 \proofstepwidestar[1,13,14]{a=b\leftrightarrow f(a)=f(b)}{\FormulaRefAuto{SemigroupIsoEqualityReflection}[thm]{1,15}}
 \proofstepwidestar[1,8,13,14]{a\in I\leftrightarrow f(a)\in f[I]}{\FormulaRefAuto{SemigroupIsoImageMembership}[thm]{1,9,15}}
 \proofstepwidestar[1,8,13,14]{b\in I\leftrightarrow f(b)\in f[I]}{\FormulaRefAuto{SemigroupIsoImageMembership}[thm]{1,9,15}}
 \proofstepwidestar[1,8,13,14]{a\mathrel{\rho_I}b\leftrightarrow f(a)\mathrel{\rho_{f[I]}}f(b)}{\FormulaRefAuto{SemigroupReesQuotientDef}[def]{16,17,18}}
 \proofstepwidestar[1,8]{\forall a,b\in A\;(a\mathrel{\rho_I}b\leftrightarrow f(a)\mathrel{\rho_{f[I]}}f(b))}{\rUI{\rRI{13,\rUI{\rRI{14,19}}}}}
 \proofstepwidestar[1,8]{f_*\rho_I=\rho_{f[I]}}{\text{\parbox[t]{\linewidth}{\raggedright Relationsextensionalität: 20 auf jedem Bildpaar; 7 liefert alle Paare in \(B\).}}}
 \proofstepwidestar[1,8]{\bar f:A/\rho_I\cong B/(f_*\rho_I)}{\FormulaRefAuto{SemigroupIsoCongruenceQuotients}[thm]{1,11}}
 \proofstepwidestar[1,8]{([a]_{\rho_I}\mapsto[f(a)]_{\rho_{f[I]}}):A/I\cong B/f[I]}{\FormulaRefAuto{SemigroupReesQuotientDef}[def]{22,21}}
\closeproofpartwide
\end{tabproofsplitwide}

\subsection{Endomorphismen und Automorphismen}

\FormulaDefDeltaK[Endomorphismen und Automorphismen einer Halbgruppe]{%
 \begin{aligned}[t]
 &\operatorname{End}(A)\coloneqq
       \{u:A\to A\mid\SemigroupHom{u}{A,\star}{A,\star}\},\\[-2pt]
 &\operatorname{Aut}(A)\coloneqq
       \{u:A\to A\mid\SemigroupIso{u}{A,\star}{A,\star}\}.
 \end{aligned}%
}{SemigroupEndomorphismCarriersDef}{%
 \DeltaRow{Voraussetzung}{\SemigroupAlg{A}{\star}}
 \DeltaRow{Mengen}{A\dsep u}
 \DeltaRow{Operation}{\circ\text{ ist die Abbildungskomposition.}}
}

\FormulaThmDeltaKR[Konjugation von Endomorphismen und Automorphismen]{%
 \begin{aligned}[t]
  &f:A\cong B\vdash{}\\[-2pt]
  &\text{(i)}\quad C_f:(\operatorname{End}(A),\circ)\cong(\operatorname{End}(B),\circ),\\[-2pt]
  &\qquad C_f(u)\coloneqq f\circ u\circ f^{-1}\\[-2pt]
  &\text{(ii)}\quad C_f|_{\operatorname{Aut}(A)}: (\operatorname{Aut}(A),\circ)\cong(\operatorname{Aut}(B),\circ)
 \end{aligned}%
}{%
 \begin{aligned}[t]
 &f:A\cong B\vdash{}\\[-2pt]
 &C_f:(\operatorname{End}(A),\circ)\cong(\operatorname{End}(B),\circ),
       \quad C_f(u)\coloneqq f\circ u\circ f^{-1},\\[-2pt]
 &C_f|_{\operatorname{Aut}(A)}:
       (\operatorname{Aut}(A),\circ)\cong(\operatorname{Aut}(B),\circ).
 \end{aligned}%
}{SemigroupIsoEndomorphismConjugation}{%
 \DeltaRow{Mengen}{A\dsep B\dsep u\dsep v}
 \DeltaRow{Funktionen}{f\dsep h\dsep C_f\dsep C_h\dsep\circ}
}
\begin{tabproofsplitwide}
\proofpartwideindR[Die beiden Kompositionshalbgruppen]{%
\SemigroupAlg{A}{\star}\vdash\SemigroupAlg{\operatorname{End}(A)}{\circ}\land\operatorname{Aut}(A)\subseteq\operatorname{End}(A)\land\SemigroupAlg{\operatorname{Aut}(A)}{\circ}%
}{\SemigroupAlg{A}{\star}\vdash\SemigroupAlg{\operatorname{End}(A)}{\circ}\land\operatorname{Aut}(A)\subseteq\operatorname{End}(A)\land\SemigroupAlg{\operatorname{Aut}(A)}{\circ}}[SemigroupEndomorphismSemigroups]
 \proofstepwidestar[1]{\SemigroupAlg{A}{\star}}{\rA}
 \proofstepwidestar[2]{u\in \operatorname{End}(A)}{\rA}
 \proofstepwidestar[3]{v\in \operatorname{End}(A)}{\rA}
 \proofstepwidestar[2,3]{u,v\in\operatorname{End}(A)}{\rAIStar{2,3}}
 \proofstepwidestar[1,2,3]{\SemigroupHom{u}{A,\star}{A,\star}}{\FormulaRefAuto{SemigroupEndomorphismCarriersDef}[def]{4}}
 \proofstepwidestar[1,2,3]{\SemigroupHom{v}{A,\star}{A,\star}}{\FormulaRefAuto{SemigroupEndomorphismCarriersDef}[def]{4}}
 \proofstepwidestar[1,2,3]{\SemigroupHom{u\circ v}{A,\star}{A,\star}}{\FormulaRefAuto{SemigroupHomComposition}[thm]{6,5}}
 \proofstepwidestar[1,2,3]{u\circ v\in\operatorname{End}(A)}{\FormulaRefAuto{SemigroupEndomorphismCarriersDef}[def]{7}}
 \proofstepwidestar[1]{\forall u,v\in\operatorname{End}(A)\;(u\circ v\in\operatorname{End}(A))}{\rUI{\rRI{2,\rUI{\rRI{3,8}}}}}
 \proofstepwidestar[10]{u\in \operatorname{End}(A)}{\rA}
 \proofstepwidestar[11]{v\in \operatorname{End}(A)}{\rA}
 \proofstepwidestar[12]{w\in \operatorname{End}(A)}{\rA}
 \proofstepwidestar[10,11,12]{u,v,w\in\operatorname{End}(A)}{\rAIStar{10,11,12}}
 \proofstepwidestar[1,10,11,12]{u:A\to A\land(v:A\to A\land w:A\to A)}{\FormulaRefAuto{SemigroupEndomorphismCarriersDef}[def]{13}}
 \proofstepwidestar[1,10,11,12]{u\circ(v\circ w)=(u\circ v)\circ w}{\FormulaRefAuto{F\colon A\to B\dsep G\colon B\to C\dsep H\colon C\to D\vdash H\circ(G\circ F)=(H\circ G)\circ F}[thm]{\rAEb{\rAEb{14}},\rAEa{\rAEb{14}},\rAEa{14}}}
  \proofstepwidestar[1,10,11,12]{(u\circ v)\circ w=u\circ(v\circ w)}{\FormulaRefAuto{a=b\vdash b=a}[thm]{15}}
 \proofstepwidestar[1]{\forall u,v,w\in\operatorname{End}(A)\;((u\circ v)\circ w=u\circ(v\circ w))}{\rUI{\rRI{10,\rUI{\rRI{11,\rUI{\rRI{12,16}}}}}}}
 \proofstepwidestar[1]{\SemigroupAlg{\operatorname{End}(A)}{\circ}}{\FormulaRefAuto{SemigroupDef}[def]{9,17}}
 \proofstepwidestar[19]{u\in\operatorname{Aut}(A)}{\rA}
  \proofstepwidestar[1,19]{\SemigroupIso{u}{A,\star}{A,\star}}{\FormulaRefAuto{SemigroupEndomorphismCarriersDef}[def]{19}}
  \proofstepwidestar[1,19]{\SemigroupHom{u}{A,\star}{A,\star}}{\FormulaRefAuto{SemigroupIsoHomAxiom}[ax]{20}}
  \proofstepwidestar[1,19]{u\in\operatorname{End}(A)}{\FormulaRefAuto{SemigroupEndomorphismCarriersDef}[def]{21}}
  \proofstepwidestar[1]{\operatorname{Aut}(A)\subseteq\operatorname{End}(A)}{\FormulaRefAuto{SubsetIntroductionRule}[thm]{[19]\vdots 22}}
 \proofstepwidestar[24]{u\in \operatorname{Aut}(A)}{\rA}
 \proofstepwidestar[25]{v\in \operatorname{Aut}(A)}{\rA}
 \proofstepwidestar[24,25]{u,v\in\operatorname{Aut}(A)}{\rAIStar{24,25}}
 \proofstepwidestar[1,24,25]{\SemigroupIso{u}{A,\star}{A,\star}}{\FormulaRefAuto{SemigroupEndomorphismCarriersDef}[def]{26}}
 \proofstepwidestar[1,24,25]{\SemigroupIso{v}{A,\star}{A,\star}}{\FormulaRefAuto{SemigroupEndomorphismCarriersDef}[def]{26}}
 \proofstepwidestar[1,24,25]{\SemigroupIso{u\circ v}{A,\star}{A,\star}}{\FormulaRefAuto{SemigroupIsoComposition}[thm]{28,27}}
 \proofstepwidestar[1,24,25]{u\circ v\in\operatorname{Aut}(A)}{\FormulaRefAuto{SemigroupEndomorphismCarriersDef}[def]{29}}
 \proofstepwidestar[1]{\forall u,v\in\operatorname{Aut}(A)\;(u\circ v\in\operatorname{Aut}(A))}{\rUI{\rRI{24,\rUI{\rRI{25,30}}}}}
 \proofstepwidestar[1]{\SemigroupAlg{\operatorname{Aut}(A)}{\circ}}{\FormulaRefAuto{SemigroupClosedSubsetIsSemigroup}[thm]{18,23,31}}
  \proofstepwidestar[1]{\SemigroupAlg{\operatorname{End}(A)}{\circ}\land\operatorname{Aut}(A)\subseteq\operatorname{End}(A)\land\SemigroupAlg{\operatorname{Aut}(A)}{\circ}}{\rAI{18,\rAI{23,32}}}
\closeproofpartwideindR
\proofpartwideindR[Abschluss und Träger der Konjugation]{%
f:A\cong B\vdash C_f[\operatorname{End}(A)]\subseteq\operatorname{End}(B)\land C_f[\operatorname{Aut}(A)]\subseteq\operatorname{Aut}(B)%
}{f:A\cong B\vdash C_f[\operatorname{End}(A)]\subseteq\operatorname{End}(B)\land C_f[\operatorname{Aut}(A)]\subseteq\operatorname{Aut}(B)}[SemigroupConjugationCarriers]
 \proofstepwidestar[1]{f:A\cong B}{\rA}
 \proofstepwidestar[1]{\SemigroupAlg{A}{\star}}{\FormulaRefAuto{SemigroupIsoSource}[thm]{1}}
 \proofstepwidestar[1]{\SemigroupAlg{B}{\diamond}}{\FormulaRefAuto{SemigroupIsoTarget}[thm]{1}}
 \proofstepwidestar[1]{f:A\to B}{\FormulaRefAuto{SemigroupIsoFunction}[thm]{1}}
 \proofstepwidestar[1]{f:A\bij B}{\FormulaRefAuto{SemigroupIsoBijectiveAxiom}[ax]{1}}
 \proofstepwidestar[1]{\SemigroupHom{f}{A,\star}{B,\diamond}}{\FormulaRefAuto{SemigroupIsoHomAxiom}[ax]{1}}
 \proofstepwidestar[1]{f[A]=B}{\FormulaRefAuto{SemigroupIsoFullImage}[thm]{1}}
 \proofstepwidestar[1]{h:B\cong A}{\FormulaRefAuto{SemigroupIsoInverse}[thm]{1}}
 \proofstepwidestar[1]{\SemigroupHom{h}{B,\diamond}{A,\star}}{\FormulaRefAuto{SemigroupIsoHomAxiom}[ax]{8}}
 \proofstepwidestar[10]{u\in\operatorname{End}(A)}{\rA}
 \proofstepwidestar[1,10]{\SemigroupHom{u}{A,\star}{A,\star}}{\FormulaRefAuto{SemigroupEndomorphismCarriersDef}[def]{10}}
 \proofstepwidestar[1,10]{\SemigroupHom{u\circ h}{B,\diamond}{A,\star}}{\FormulaRefAuto{SemigroupHomComposition}[thm]{9,11}}
 \proofstepwidestar[1,10]{\SemigroupHom{f\circ(u\circ h)}{B,\diamond}{B,\diamond}}{\FormulaRefAuto{SemigroupHomComposition}[thm]{12,6}}
 \proofstepwidestar[1,10]{C_f(u)\in\operatorname{End}(B)}{\text{\parbox[t]{\linewidth}{\raggedright \FormulaRefAuto{SemigroupEndomorphismCarriersDef}[def]{13}; Definition von \(C_f\).}}}
 \proofstepwidestar[1]{\forall u\in\operatorname{End}(A)\;(C_f(u)\in\operatorname{End}(B))}{\rUI{\rRI{10,14}}}
  \proofstepwidestar[1]{C_f:\operatorname{End}(A)\to\operatorname{End}(B)}{\text{\parbox[t]{\linewidth}{\raggedright \FormulaRefAuto{\forall u\in A\,F(u)\in B\vdash\Graph(F)\colon A\to B}[thm]{15}; Definition von \(C_f\) als Wertzuordnung}}}
  \proofstepwidestar[1]{C_f[\operatorname{End}(A)]\subseteq\operatorname{End}(B)}{\FormulaRefAuto{F\colon A\to B\vdash F[A]\subseteq B}[thm]{16}}
 \proofstepwidestar[18]{u\in\operatorname{Aut}(A)}{\rA}
 \proofstepwidestar[1,18]{\SemigroupIso{u}{A,\star}{A,\star}}{\FormulaRefAuto{SemigroupEndomorphismCarriersDef}[def]{18}}
 \proofstepwidestar[1,18]{\SemigroupIso{u\circ h}{B,\diamond}{A,\star}}{\FormulaRefAuto{SemigroupIsoComposition}[thm]{8,19}}
 \proofstepwidestar[1,18]{\SemigroupIso{f\circ(u\circ h)}{B,\diamond}{B,\diamond}}{\FormulaRefAuto{SemigroupIsoComposition}[thm]{20,1}}
 \proofstepwidestar[1,18]{C_f(u)\in\operatorname{Aut}(B)}{\text{\parbox[t]{\linewidth}{\raggedright \FormulaRefAuto{SemigroupEndomorphismCarriersDef}[def]{21}; Definition von \(C_f\).}}}
 \proofstepwidestar[1]{\forall u\in\operatorname{Aut}(A)\;(C_f(u)\in\operatorname{Aut}(B))}{\rUI{\rRI{18,22}}}
  \proofstepwidestar[1]{\operatorname{Aut}(A)\subseteq\operatorname{End}(A)}{\rAEa{\rAEb{\FormulaRefAuto{SemigroupEndomorphismSemigroups}[thm]{2}}}}
  \proofstepwidestar[25]{w\in C_f[\operatorname{Aut}(A)]}{\rA}
  \proofstepwidestar[1,25]{\exists u\in\operatorname{Aut}(A)\;(w=C_f(u))}{\FormulaRefAuto{C\subseteq A\dsep F\colon A\to B\dsep y\in F[C]\vdash\exists x\in C\,y=F(x)}[thm]{24,16,25}}
  \proofstepwidestar[27]{u\in\operatorname{Aut}(A)\land w=C_f(u)}{\rA}
  \proofstepwidestar[1,27]{C_f(u)\in\operatorname{Aut}(B)}{\rRE{\rUE{23},\rAEa{27}}}
  \proofstepwidestar[1,27]{w\in\operatorname{Aut}(B)}{\rIE{\FormulaRefAuto{a=b\vdash b=a}[thm]{\rAEb{27}},28}}
  \proofstepwidestar[1,25]{w\in\operatorname{Aut}(B)}{\rEE{26,27,29}}
  \proofstepwidestar[1]{C_f[\operatorname{Aut}(A)]\subseteq\operatorname{Aut}(B)}{\FormulaRefAuto{SubsetIntroductionRule}[thm]{[25]\vdots 30}}
  \proofstepwidestar[1]{C_f[\operatorname{End}(A)]\subseteq\operatorname{End}(B)\land C_f[\operatorname{Aut}(A)]\subseteq\operatorname{Aut}(B)}{\rAI{17,31}}
\closeproofpartwideindR
\proofpartwide{Inverse Konjugation und Produkt}
 \proofstepwidestar[1]{f:A\cong B}{\rA}
 \proofstepwidestar[1]{\SemigroupAlg{A}{\star}}{\FormulaRefAuto{SemigroupIsoSource}[thm]{1}}
 \proofstepwidestar[1]{\SemigroupAlg{B}{\diamond}}{\FormulaRefAuto{SemigroupIsoTarget}[thm]{1}}
 \proofstepwidestar[1]{f:A\to B}{\FormulaRefAuto{SemigroupIsoFunction}[thm]{1}}
 \proofstepwidestar[1]{f:A\bij B}{\FormulaRefAuto{SemigroupIsoBijectiveAxiom}[ax]{1}}
 \proofstepwidestar[1]{\SemigroupHom{f}{A,\star}{B,\diamond}}{\FormulaRefAuto{SemigroupIsoHomAxiom}[ax]{1}}
 \proofstepwidestar[1]{f[A]=B}{\FormulaRefAuto{SemigroupIsoFullImage}[thm]{1}}
 \proofstepwidestar[1]{h:B\cong A}{\FormulaRefAuto{SemigroupIsoInverse}[thm]{1}}
 \proofstepwidestar[1]{h:B\to A}{\FormulaRefAuto{SemigroupIsoFunction}[thm]{8}}
 \proofstepwidestar[10]{x\in A}{\rA}
  \proofstepwidestar[1,10]{h(f(x))=x}{\FormulaRefAuto{InverseFunctionLeftInverse}[thm]{5,10}}
  \proofstepwidestar[1]{\forall x\in A\;(h(f(x))=x)}{\rUI{\rRI{10,11}}}
 \proofstepwidestar[13]{y\in B}{\rA}
  \proofstepwidestar[1,13]{f(h(y))=y}{\FormulaRefAuto{InverseFunctionRightInverse}[thm]{5,13}}
  \proofstepwidestar[1]{\forall y\in B\;(f(h(y))=y)}{\rUI{\rRI{13,14}}}
 \proofstepwidestar[1]{h:B\bij A}{\FormulaRefAuto{SemigroupIsoBijectiveAxiom}[ax]{8}}
  \proofstepwidestar[1]{f=h^{-1}}{\FormulaRefAuto{InverseFunctionRightInverseUniqueness}[thm]{16,4,12}}
  \proofstepwidestar[1]{h^{-1}=f}{\FormulaRefAuto{a=b\vdash b=a}[thm]{17}}
 \proofstepwidestar[1]{\SemigroupAlg{\operatorname{End}(A)}{\circ}\land\operatorname{Aut}(A)\subseteq\operatorname{End}(A)\land\SemigroupAlg{\operatorname{Aut}(A)}{\circ}}{\FormulaRefAuto{SemigroupEndomorphismSemigroups}[thm]{2}}
 \proofstepwidestar[1]{\SemigroupAlg{\operatorname{End}(B)}{\circ}\land\operatorname{Aut}(B)\subseteq\operatorname{End}(B)\land\SemigroupAlg{\operatorname{Aut}(B)}{\circ}}{\FormulaRefAuto{SemigroupEndomorphismSemigroups}[thm]{3}}
 \proofstepwidestar[1]{\SemigroupAlg{\operatorname{End}(A)}{\circ}}{\rAEa{19}}
 \proofstepwidestar[1]{\SemigroupAlg{\operatorname{End}(B)}{\circ}}{\rAEa{20}}
 \proofstepwidestar[1]{\operatorname{Aut}(A)\subseteq\operatorname{End}(A)}{\rAEa{\rAEb{19}}}
 \proofstepwidestar[1]{\SemigroupAlg{\operatorname{Aut}(A)}{\circ}}{\rAEb{\rAEb{19}}}
 \proofstepwidestar[1]{C_f[\operatorname{End}(A)]\subseteq\operatorname{End}(B)\land C_f[\operatorname{Aut}(A)]\subseteq\operatorname{Aut}(B)}{\FormulaRefAuto{SemigroupConjugationCarriers}[thm]{1}}
 \proofstepwidestar[1]{C_h[\operatorname{End}(B)]\subseteq\operatorname{End}(A)\land C_h[\operatorname{Aut}(B)]\subseteq\operatorname{Aut}(A)}{\FormulaRefAuto{SemigroupConjugationCarriers}[thm]{8}}
 \proofstepwidestar[1]{C_f:\operatorname{End}(A)\to\operatorname{End}(B)\land C_h:\operatorname{End}(B)\to\operatorname{End}(A)}{\text{\parbox[t]{\linewidth}{\raggedright Definition der Konjugationsabbildungen und 25,26}}}
 \proofstepwidestar[28]{u\in\operatorname{End}(A)}{\rA}
 \proofstepwidestar[1,28]{u:A\to A}{\FormulaRefAuto{SemigroupEndomorphismCarriersDef}[def]{28}}
 \proofstepwidestar[1,28]{C_f(u)\in\operatorname{End}(B)}{\FormulaRefAuto{F\colon A\to B\dsep x\in A\vdash F(x)\in B}[thm]{\rAEa{27},28}}
  \proofstepwidestar[1,28]{C_f(u):B\to B}{\FormulaRefAuto{SemigroupEndomorphismCarriersDef}[def]{30}}
  \proofstepwidestar[1,28]{C_h(C_f(u))\in\operatorname{End}(A)}{\FormulaRefAuto{F\colon A\to B\dsep x\in A\vdash F(x)\in B}[thm]{\rAEb{27},30}}
  \proofstepwidestar[1,28]{C_h(C_f(u)):A\to A}{\FormulaRefAuto{SemigroupEndomorphismCarriersDef}[def]{32}}
  \proofstepwidestar[34]{x\in A}{\rA}
  \proofstepwidestar[1,28,34]{f(x)\in B}{\FormulaRefAuto{F\colon A\to B\dsep x\in A\vdash F(x)\in B}[thm]{4,34}}
  \proofstepwidestar[1,28,34]{h(f(x))=x}{\rRE{\rUE{12},34}}
  \proofstepwidestar[1,28,34]{u(x)\in A}{\FormulaRefAuto{F\colon A\to B\dsep x\in A\vdash F(x)\in B}[thm]{29,34}}
  \proofstepwidestar[1,28,34]{h(f(u(x)))=u(x)}{\rRE{\rUE{12},37}}
  \proofstepwidestar[1,28,34]{C_f(u)=(f\circ u)\circ h}{\text{Definition von \(C_f(u)\)}}
  \proofstepwidestar[1,28,34]{((f\circ u)\circ h)(f(x))=f(u(h(f(x))))}{\FormulaRefAuto{x\in A\vdash ((F\circ G)\circ H)(x)=F(G(H(x)))}[thm]{9,29,4,35}}
  \proofstepwidestar[1,28,34]{C_f(u)(f(x))=f(u(h(f(x))))}{\rIE{\FormulaRefAuto{a=b\vdash b=a}[thm]{39},40}}
  \proofstepwidestar[1,28,34]{C_h(C_f(u))=(h\circ C_f(u))\circ h^{-1}}{\text{\parbox[t]{\linewidth}{\raggedright Definition von \(C_h(C_f(u))\)}}}
  \proofstepwidestar[1,28,34]{C_h(C_f(u))=(h\circ C_f(u))\circ f}{\rIE{18,42}}
  \proofstepwidestar[1,28,34]{((h\circ C_f(u))\circ f)(x)=h(C_f(u)(f(x)))}{\FormulaRefAuto{x\in A\vdash ((F\circ G)\circ H)(x)=F(G(H(x)))}[thm]{4,31,9,34}}
  \proofstepwidestar[1,28,34]{C_h(C_f(u))(x)=h(C_f(u)(f(x)))}{\rIE{\FormulaRefAuto{a=b\vdash b=a}[thm]{43},44}}
  \proofstepwidestar[1,28,34]{C_h(C_f(u))(x)=u(x)}{\rIE{41,36,38,45}}
  \proofstepwidestar[1,28]{\forall x\in A\;(C_h(C_f(u))(x)=u(x))}{\rUI{\rRI{34,46}}}
 \proofstepwidestar[1,28]{C_h(C_f(u))=u}{\FormulaRefAuto{FunctionExtensionality}[thm]{33,29,47}}
 \proofstepwidestar[1]{C_h\circ C_f:\operatorname{End}(A)\to \operatorname{End}(A)}{\FormulaRefAuto{CompositionDef}[def]{\rAEa{27},\rAEb{27}}}
  \proofstepwidestar[]{\Id_{\operatorname{End}(A)}:\operatorname{End}(A)\to \operatorname{End}(A)}{\FormulaRefAuto{IdA}[thm]}
  \proofstepwidestar[1,28]{(C_h\circ C_f)(u)=C_h(C_f(u))}{\FormulaRefAuto{CompositionDef}[def]{\rAEa{27},\rAEb{27},28}}
  \proofstepwidestar[28]{u=\Id_{\operatorname{End}(A)}(u)}{\FormulaRefAuto{x\in A\vdash x=\Id_A(x)}[thm]{28}}
  \proofstepwidestar[1,28]{(C_h\circ C_f)(u)=\Id_{\operatorname{End}(A)}(u)}{\rIE{48,52,51}}
  \proofstepwidestar[1]{\forall u\in \operatorname{End}(A)\;((C_h\circ C_f)(u)=\Id_{\operatorname{End}(A)}(u))}{\rUI{\rRI{28,53}}}
  \proofstepwidestar[1]{C_h\circ C_f=\Id_{\operatorname{End}(A)}}{\FormulaRefAuto{FunctionExtensionality}[thm]{49,50,54}}
 \proofstepwidestar[56]{w\in\operatorname{End}(B)}{\rA}
 \proofstepwidestar[1,56]{w:B\to B}{\FormulaRefAuto{SemigroupEndomorphismCarriersDef}[def]{56}}
 \proofstepwidestar[1,56]{C_h(w)\in\operatorname{End}(A)}{\FormulaRefAuto{F\colon A\to B\dsep x\in A\vdash F(x)\in B}[thm]{\rAEb{27},56}}
  \proofstepwidestar[1,56]{C_h(w):A\to A}{\FormulaRefAuto{SemigroupEndomorphismCarriersDef}[def]{58}}
  \proofstepwidestar[1,56]{C_f(C_h(w))\in\operatorname{End}(B)}{\FormulaRefAuto{F\colon A\to B\dsep x\in A\vdash F(x)\in B}[thm]{\rAEa{27},58}}
  \proofstepwidestar[1,56]{C_f(C_h(w)):B\to B}{\FormulaRefAuto{SemigroupEndomorphismCarriersDef}[def]{60}}
  \proofstepwidestar[62]{y\in B}{\rA}
  \proofstepwidestar[1,56,62]{h(y)\in A}{\FormulaRefAuto{F\colon A\to B\dsep x\in A\vdash F(x)\in B}[thm]{9,62}}
  \proofstepwidestar[1,56,62]{f(h(y))=y}{\rRE{\rUE{15},62}}
  \proofstepwidestar[1,56,62]{w(y)\in B}{\FormulaRefAuto{F\colon A\to B\dsep x\in A\vdash F(x)\in B}[thm]{57,62}}
  \proofstepwidestar[1,56,62]{f(h(w(y)))=w(y)}{\rRE{\rUE{15},65}}
  \proofstepwidestar[1,56,62]{C_h(w)=(h\circ w)\circ h^{-1}}{\text{Definition von \(C_h(w)\)}}
  \proofstepwidestar[1,56,62]{C_h(w)=(h\circ w)\circ f}{\rIE{18,67}}
  \proofstepwidestar[1,56,62]{((h\circ w)\circ f)(h(y))=h(w(f(h(y))))}{\FormulaRefAuto{x\in A\vdash ((F\circ G)\circ H)(x)=F(G(H(x)))}[thm]{4,57,9,63}}
  \proofstepwidestar[1,56,62]{C_h(w)(h(y))=h(w(f(h(y))))}{\rIE{\FormulaRefAuto{a=b\vdash b=a}[thm]{68},69}}
  \proofstepwidestar[1,56,62]{C_f(C_h(w))=(f\circ C_h(w))\circ h}{\text{\parbox[t]{\linewidth}{\raggedright Definition von \(C_f(C_h(w))\)}}}
  \proofstepwidestar[1,56,62]{((f\circ C_h(w))\circ h)(y)=f(C_h(w)(h(y)))}{\FormulaRefAuto{x\in A\vdash ((F\circ G)\circ H)(x)=F(G(H(x)))}[thm]{9,59,4,62}}
  \proofstepwidestar[1,56,62]{C_f(C_h(w))(y)=f(C_h(w)(h(y)))}{\rIE{\FormulaRefAuto{a=b\vdash b=a}[thm]{71},72}}
  \proofstepwidestar[1,56,62]{C_f(C_h(w))(y)=w(y)}{\rIE{70,64,66,73}}
  \proofstepwidestar[1,56]{\forall y\in B\;(C_f(C_h(w))(y)=w(y))}{\rUI{\rRI{62,74}}}
 \proofstepwidestar[1,56]{C_f(C_h(w))=w}{\FormulaRefAuto{FunctionExtensionality}[thm]{61,57,75}}
 \proofstepwidestar[1]{C_f\circ C_h:\operatorname{End}(B)\to \operatorname{End}(B)}{\FormulaRefAuto{CompositionDef}[def]{\rAEb{27},\rAEa{27}}}
  \proofstepwidestar[]{\Id_{\operatorname{End}(B)}:\operatorname{End}(B)\to \operatorname{End}(B)}{\FormulaRefAuto{IdA}[thm]}
  \proofstepwidestar[1,56]{(C_f\circ C_h)(w)=C_f(C_h(w))}{\FormulaRefAuto{CompositionDef}[def]{\rAEb{27},\rAEa{27},56}}
  \proofstepwidestar[56]{w=\Id_{\operatorname{End}(B)}(w)}{\FormulaRefAuto{x\in A\vdash x=\Id_A(x)}[thm]{56}}
  \proofstepwidestar[1,56]{(C_f\circ C_h)(w)=\Id_{\operatorname{End}(B)}(w)}{\rIE{76,80,79}}
  \proofstepwidestar[1]{\forall w\in \operatorname{End}(B)\;((C_f\circ C_h)(w)=\Id_{\operatorname{End}(B)}(w))}{\rUI{\rRI{56,81}}}
  \proofstepwidestar[1]{C_f\circ C_h=\Id_{\operatorname{End}(B)}}{\FormulaRefAuto{FunctionExtensionality}[thm]{77,78,82}}
 \proofstepwidestar[1]{C_f:\operatorname{End}(A)\bij\operatorname{End}(B)}{\FormulaRefAuto{MutuallyInverseFunctionBijection}[thm]{\rAEb{27},\rAEa{27},83,55}}
 \proofstepwidestar[85]{u\in \operatorname{End}(A)}{\rA}
 \proofstepwidestar[86]{v\in \operatorname{End}(A)}{\rA}
 \proofstepwidestar[85,86]{u,v\in\operatorname{End}(A)}{\rAIStar{85,86}}
 \proofstepwidestar[1,85,86]{u:A\to A\land v:A\to A}{\FormulaRefAuto{SemigroupEndomorphismCarriersDef}[def]{87}}
 \proofstepwidestar[1,85,86]{u\circ v\in\operatorname{End}(A)}{\FormulaRefAuto{SemigroupClosureAxiom}[ax]{21,87}}
  \proofstepwidestar[1,85,86]{C_f(u)\in\operatorname{End}(B)}{\FormulaRefAuto{F\colon A\to B\dsep x\in A\vdash F(x)\in B}[thm]{\rAEa{27},85}}
  \proofstepwidestar[1,85,86]{C_f(u):B\to B}{\FormulaRefAuto{SemigroupEndomorphismCarriersDef}[def]{90}}
  \proofstepwidestar[1,85,86]{C_f(v)\in\operatorname{End}(B)}{\FormulaRefAuto{F\colon A\to B\dsep x\in A\vdash F(x)\in B}[thm]{\rAEa{27},86}}
  \proofstepwidestar[1,85,86]{C_f(v):B\to B}{\FormulaRefAuto{SemigroupEndomorphismCarriersDef}[def]{92}}
  \proofstepwidestar[1,85,86]{C_f(u\circ v)\in\operatorname{End}(B)}{\FormulaRefAuto{F\colon A\to B\dsep x\in A\vdash F(x)\in B}[thm]{\rAEa{27},89}}
  \proofstepwidestar[1,85,86]{C_f(u\circ v):B\to B}{\FormulaRefAuto{SemigroupEndomorphismCarriersDef}[def]{94}}
  \proofstepwidestar[1,85,86]{u\circ v:A\to A}{\FormulaRefAuto{CompositionDef}[def]{\rAEb{88},\rAEa{88}}}
  \proofstepwidestar[1,85,86]{C_f(u)\circ C_f(v):B\to B}{\FormulaRefAuto{CompositionDef}[def]{93,91}}
  \proofstepwidestar[98]{y\in B}{\rA}
  \proofstepwidestar[1,85,86,98]{h(y)\in A}{\FormulaRefAuto{F\colon A\to B\dsep x\in A\vdash F(x)\in B}[thm]{9,98}}
  \proofstepwidestar[1,85,86,98]{v(h(y))\in A}{\FormulaRefAuto{F\colon A\to B\dsep x\in A\vdash F(x)\in B}[thm]{\rAEb{88},99}}
  \proofstepwidestar[1,85,86,98]{f(v(h(y)))\in B}{\FormulaRefAuto{F\colon A\to B\dsep x\in A\vdash F(x)\in B}[thm]{4,100}}
  \proofstepwidestar[1,85,86,98]{h(f(v(h(y))))=v(h(y))}{\rRE{\rUE{12},100}}
  \proofstepwidestar[1,85,86,98]{C_f(v)=(f\circ v)\circ h}{\text{Definition von \(C_f(v)\)}}
  \proofstepwidestar[1,85,86,98]{((f\circ v)\circ h)(y)=f(v(h(y)))}{\FormulaRefAuto{x\in A\vdash ((F\circ G)\circ H)(x)=F(G(H(x)))}[thm]{9,\rAEb{88},4,98}}
  \proofstepwidestar[1,85,86,98]{C_f(v)(y)=f(v(h(y)))}{\rIE{\FormulaRefAuto{a=b\vdash b=a}[thm]{103},104}}
  \proofstepwidestar[1,85,86,98]{C_f(u)=(f\circ u)\circ h}{\text{Definition von \(C_f(u)\)}}
  \proofstepwidestar[1,85,86,98]{((f\circ u)\circ h)(f(v(h(y))))=f(u(h(f(v(h(y))))))}{\FormulaRefAuto{x\in A\vdash ((F\circ G)\circ H)(x)=F(G(H(x)))}[thm]{9,\rAEa{88},4,101}}
  \proofstepwidestar[1,85,86,98]{C_f(u)(f(v(h(y))))=f(u(h(f(v(h(y))))))}{\rIE{\FormulaRefAuto{a=b\vdash b=a}[thm]{106},107}}
  \proofstepwidestar[1,85,86,98]{(C_f(u)\circ C_f(v))(y)=C_f(u)(C_f(v)(y))}{\FormulaRefAuto{CompositionDef}[def]{93,91,98}}
  \proofstepwidestar[1,85,86,98]{(C_f(u)\circ C_f(v))(y)=f(u(v(h(y))))}{\rIE{105,108,102,109}}
  \proofstepwidestar[1,85,86,98]{C_f((u\circ v))=(f\circ (u\circ v))\circ h}{\text{\parbox[t]{\linewidth}{\raggedright Definition von \(C_f((u\circ v))\)}}}
  \proofstepwidestar[1,85,86,98]{((f\circ (u\circ v))\circ h)(y)=f((u\circ v)(h(y)))}{\FormulaRefAuto{x\in A\vdash ((F\circ G)\circ H)(x)=F(G(H(x)))}[thm]{9,96,4,98}}
  \proofstepwidestar[1,85,86,98]{C_f((u\circ v))(y)=f((u\circ v)(h(y)))}{\rIE{\FormulaRefAuto{a=b\vdash b=a}[thm]{111},112}}
  \proofstepwidestar[1,85,86,98]{(u\circ v)(h(y))=u(v(h(y)))}{\FormulaRefAuto{CompositionDef}[def]{\rAEb{88},\rAEa{88},99}}
  \proofstepwidestar[1,85,86,98]{C_f(u\circ v)(y)=f(u(v(h(y))))}{\rIE{114,113}}
  \proofstepwidestar[1,85,86,98]{C_f(u\circ v)(y)=(C_f(u)\circ C_f(v))(y)}{\FormulaRefAuto{a=b, c=b\vdash a=c}[thm]{115,110}}
  \proofstepwidestar[1,85,86]{\forall y\in B\;(C_f(u\circ v)(y)=(C_f(u)\circ C_f(v))(y))}{\rUI{\rRI{98,116}}}
 \proofstepwidestar[1,85,86]{C_f(u\circ v)=C_f(u)\circ C_f(v)}{\FormulaRefAuto{FunctionExtensionality}[thm]{95,97,117}}
 \proofstepwidestar[1]{\forall u,v\in\operatorname{End}(A)\;(C_f(u\circ v)=C_f(u)\circ C_f(v))}{\rUI{\rRI{85,\rUI{\rRI{86,118}}}}}
 \proofstepwidestar[1]{\SemigroupHom{C_f}{\operatorname{End}(A),\circ}{\operatorname{End}(B),\circ}}{\FormulaRefAuto{SemigroupHomDef}[def]{21,22,\rAEa{27},119}}
 \proofstepwidestar[1]{C_f:\operatorname{End}(A)\cong\operatorname{End}(B)}{\FormulaRefAuto{SemigroupIsoDef}[def]{120,84}}
 \proofstepwidestar[1]{C_f[\operatorname{Aut}(A)]\subseteq\operatorname{Aut}(B)}{\rAEb{25}}
  \proofstepwidestar[1]{\operatorname{Aut}(B)\subseteq\operatorname{End}(B)}{\rAEa{\rAEb{20}}}
  \proofstepwidestar[124]{w\in\operatorname{Aut}(B)}{\rA}
  \proofstepwidestar[1,124]{w\in\operatorname{End}(B)}{\FormulaRefAuto{A\subseteq B,\,x\in A\vdash x\in B}[thm]{123,124}}
  \proofstepwidestar[1,124]{C_h(w)\in C_h[\operatorname{Aut}(B)]}{\FormulaRefAuto{C\subseteq A\dsep F\colon A\to B\dsep x\in C\vdash F(x)\in F[C]}[thm]{123,\rAEb{27},124}}
  \proofstepwidestar[1,124]{C_h(w)\in\operatorname{Aut}(A)}{\FormulaRefAuto{A\subseteq B,\,x\in A\vdash x\in B}[thm]{\rAEb{26},126}}
  \proofstepwidestar[1,124]{(C_f\circ C_h)(w)=C_f(C_h(w))}{\FormulaRefAuto{CompositionDef}[def]{\rAEb{27},\rAEa{27},125}}
  \proofstepwidestar[1,124]{\Id_{\operatorname{End}(B)}(w)=w}{\FormulaRefAuto{x\in A\vdash\Id_A(x)=x}[thm]{125}}
  \proofstepwidestar[1,124]{w=C_f(C_h(w))}{\rIE{83,129,128}}
  \proofstepwidestar[1,124]{C_f(C_h(w))\in C_f[\operatorname{Aut}(A)]}{\FormulaRefAuto{C\subseteq A\dsep F\colon A\to B\dsep x\in C\vdash F(x)\in F[C]}[thm]{23,\rAEa{27},127}}
  \proofstepwidestar[1,124]{w\in C_f[\operatorname{Aut}(A)]}{\rIE{\FormulaRefAuto{a=b\vdash b=a}[thm]{130},131}}
  \proofstepwidestar[1]{\operatorname{Aut}(B)\subseteq C_f[\operatorname{Aut}(A)]}{\FormulaRefAuto{SubsetIntroductionRule}[thm]{[124]\vdots 132}}
  \proofstepwidestar[1]{C_f[\operatorname{Aut}(A)]=\operatorname{Aut}(B)}{\FormulaRefAuto{A\subseteq B, B\subseteq A\vdash A=B}[thm]{122,133}}
 \proofstepwidestar[1]{C_f|_{\operatorname{Aut}(A)}:\operatorname{Aut}(A)\cong C_f[\operatorname{Aut}(A)]}{\FormulaRefAuto{SemigroupIsoClosedCarrierRestriction}[thm]{121,23,24}}
 \proofstepwidestar[1]{C_f|_{\operatorname{Aut}(A)}:\operatorname{Aut}(A)\cong\operatorname{Aut}(B)}{\rIE{134,135}}
\closeproofpartwide
\end{tabproofsplitwide}

\subsection{Ordnungen von Unterstrukturen}

\FormulaDefDeltaK[Geordnete Mengen von Unterstrukturen]{%
 \begin{aligned}[t]
 &\operatorname{Sub}(A)\coloneqq\{U\subseteq A\mid\mathsf{UHGr}(U;A,\star)\},\\[-2pt]
 &\operatorname{LId}(A)\coloneqq\{I\subseteq A\mid\mathsf{LId}(I;A,\star)\},\\[-2pt]
 &\operatorname{RId}(A)\coloneqq\{I\subseteq A\mid\mathsf{RId}(I;A,\star)\},\\[-2pt]
 &\operatorname{Id}(A)\coloneqq\{I\subseteq A\mid\mathsf{Id}(I;A,\star)\}.
 \end{aligned}%
}{SemigroupSubstructurePosetsDef}{%
 \DeltaRow{Voraussetzung}{\SemigroupAlg{A}{\star}}
 \DeltaRow{Ordnung}{\subseteq\text{ auf allen vier Mengen sowie auf }\operatorname{Con}(A)}
}
Ein Ordnungsisomorphismus bedeutet hier ausdrücklich eine Bijektion,
die die Inklusionsrelation in beiden Richtungen erhält; dies ist kein
Halbgruppenisomorphismus ohne zusätzliche Operation auf der Strukturmenge.

\FormulaThmDeltaKR[Isomorphe Inklusionsordnungen der Unterstrukturen]{%
 \begin{aligned}[t]
  &f:A\cong B\dsep U,V\subseteq A\dsep \rho,\sigma\in\operatorname{Con}(A)\vdash{}\\[-2pt]
  &\text{(i)}\quad (U\mapsto f[U]):\mathcal S(A)\bij\mathcal S(B)\\[-2pt]
  &\qquad(\mathcal S=\operatorname{Sub},\operatorname{LId},
    \operatorname{RId},\operatorname{Id})\\[-2pt]
  &\text{(ii)}\quad U\subseteq V\leftrightarrow f[U]\subseteq f[V]\\[-2pt]
  &\text{(iii)}\quad (\rho\mapsto f_*\rho):\operatorname{Con}(A)\bij\operatorname{Con}(B)\\[-2pt]
  &\text{(iv)}\quad \rho\subseteq\sigma\leftrightarrow f_*\rho\subseteq f_*\sigma
 \end{aligned}%
}{%
 \begin{aligned}[t]
 &f:A\cong B\dsep U,V\subseteq A\dsep
       \rho,\sigma\in\operatorname{Con}(A)\vdash{}\\[-2pt]
 &(U\mapsto f[U]):\mathcal S(A)\bij\mathcal S(B)
       \quad(\mathcal S=\operatorname{Sub},\operatorname{LId},
          \operatorname{RId},\operatorname{Id}),\\[-2pt]
 &U\subseteq V\leftrightarrow f[U]\subseteq f[V],\\[-2pt]
 &(\rho\mapsto f_*\rho):\operatorname{Con}(A)\bij\operatorname{Con}(B),
       \qquad\rho\subseteq\sigma\leftrightarrow f_*\rho\subseteq f_*\sigma.
 \end{aligned}%
}{SemigroupIsoSubstructureOrders}{%
 \DeltaRow{Mengen}{A\dsep B\dsep U\dsep V}
 \DeltaRow{Funktionen}{f\dsep h}
 \DeltaRow{Relationen}{\rho\dsep\sigma}
}
\begin{tabproofsplitwide}
\proofpartwide{Unterhalbgruppen und Ideale}
 \proofstepwidestar[1]{f:A\cong B}{\rA}
 \proofstepwidestar[1]{\SemigroupAlg{A}{\star}}{\FormulaRefAuto{SemigroupIsoSource}[thm]{1}}
 \proofstepwidestar[1]{\SemigroupAlg{B}{\diamond}}{\FormulaRefAuto{SemigroupIsoTarget}[thm]{1}}
 \proofstepwidestar[1]{f:A\to B}{\FormulaRefAuto{SemigroupIsoFunction}[thm]{1}}
 \proofstepwidestar[1]{f:A\bij B}{\FormulaRefAuto{SemigroupIsoBijectiveAxiom}[ax]{1}}
 \proofstepwidestar[1]{\SemigroupHom{f}{A,\star}{B,\diamond}}{\FormulaRefAuto{SemigroupIsoHomAxiom}[ax]{1}}
 \proofstepwidestar[1]{f[A]=B}{\FormulaRefAuto{SemigroupIsoFullImage}[thm]{1}}
 \proofstepwidestar[1]{h:B\cong A}{\FormulaRefAuto{SemigroupIsoInverse}[thm]{1}}
 \proofstepwidestar[1]{h:B\to A}{\FormulaRefAuto{SemigroupIsoFunction}[thm]{8}}
 \proofstepwidestar[10]{\mathsf{UHGr}(W;A,\star)}{\rA}
 \proofstepwidestar[1,10]{\mathsf{UHGr}(f[W];B,\diamond)}{\FormulaRefAuto{SemigroupIsoUHGrImage}[thm]{1,10}}
 \proofstepwidestar[1]{\forall W\;(\mathsf{UHGr}(W;A,\star)\Rightarrow\mathsf{UHGr}(f[W];B,\diamond))}{\rUI{\rRI{10,11}}}
 \proofstepwidestar[13]{\mathsf{UHGr}(W;B,\diamond)}{\rA}
 \proofstepwidestar[1,13]{\mathsf{UHGr}(h[W];A,\star)}{\FormulaRefAuto{SemigroupIsoUHGrImage}[thm]{8,13}}
 \proofstepwidestar[1]{\forall W\;(\mathsf{UHGr}(W;B,\diamond)\Rightarrow\mathsf{UHGr}(h[W];A,\star))}{\rUI{\rRI{13,14}}}
 \proofstepwidestar[16]{\mathsf{LId}(W;A,\star)}{\rA}
 \proofstepwidestar[1,16]{\mathsf{LId}(f[W];B,\diamond)}{\FormulaRefAuto{SemigroupIsoLIdImage}[thm]{1,16}}
 \proofstepwidestar[1]{\forall W\;(\mathsf{LId}(W;A,\star)\Rightarrow\mathsf{LId}(f[W];B,\diamond))}{\rUI{\rRI{16,17}}}
 \proofstepwidestar[19]{\mathsf{LId}(W;B,\diamond)}{\rA}
 \proofstepwidestar[1,19]{\mathsf{LId}(h[W];A,\star)}{\FormulaRefAuto{SemigroupIsoLIdImage}[thm]{8,19}}
 \proofstepwidestar[1]{\forall W\;(\mathsf{LId}(W;B,\diamond)\Rightarrow\mathsf{LId}(h[W];A,\star))}{\rUI{\rRI{19,20}}}
 \proofstepwidestar[22]{\mathsf{RId}(W;A,\star)}{\rA}
 \proofstepwidestar[1,22]{\mathsf{RId}(f[W];B,\diamond)}{\FormulaRefAuto{SemigroupIsoRIdImage}[thm]{1,22}}
 \proofstepwidestar[1]{\forall W\;(\mathsf{RId}(W;A,\star)\Rightarrow\mathsf{RId}(f[W];B,\diamond))}{\rUI{\rRI{22,23}}}
 \proofstepwidestar[25]{\mathsf{RId}(W;B,\diamond)}{\rA}
 \proofstepwidestar[1,25]{\mathsf{RId}(h[W];A,\star)}{\FormulaRefAuto{SemigroupIsoRIdImage}[thm]{8,25}}
 \proofstepwidestar[1]{\forall W\;(\mathsf{RId}(W;B,\diamond)\Rightarrow\mathsf{RId}(h[W];A,\star))}{\rUI{\rRI{25,26}}}
 \proofstepwidestar[28]{\mathsf{Id}(W;A,\star)}{\rA}
 \proofstepwidestar[1,28]{\mathsf{Id}(f[W];B,\diamond)}{\FormulaRefAuto{SemigroupIsoIdImage}[thm]{1,28}}
 \proofstepwidestar[1]{\forall W\;(\mathsf{Id}(W;A,\star)\Rightarrow\mathsf{Id}(f[W];B,\diamond))}{\rUI{\rRI{28,29}}}
 \proofstepwidestar[31]{\mathsf{Id}(W;B,\diamond)}{\rA}
 \proofstepwidestar[1,31]{\mathsf{Id}(h[W];A,\star)}{\FormulaRefAuto{SemigroupIsoIdImage}[thm]{8,31}}
 \proofstepwidestar[1]{\forall W\;(\mathsf{Id}(W;B,\diamond)\Rightarrow\mathsf{Id}(h[W];A,\star))}{\rUI{\rRI{31,32}}}
 \proofstepwidestar[1]{\Phi:\mathcal S(A)\to\mathcal S(B)\land\Psi:\mathcal S(B)\to\mathcal S(A)}{\text{\parbox[t]{\linewidth}{\raggedright \FormulaRefAuto{SemigroupSubstructurePosetsDef}[def]{12,18,24,30,15,21,27,33}; \(\Phi(W)=f[W]\), \(\Psi(W)=h[W]\).}}}
 \proofstepwidestar[35]{W\in\mathcal S(A)}{\rA}
 \proofstepwidestar[1,35]{W\subseteq A}{\FormulaRefAuto{SemigroupSubstructurePosetsDef}[def]{35}}
 \proofstepwidestar[1,35]{f[W]\subseteq B\land h[f[W]]=W\land(W\neq\varnothing\leftrightarrow f[W]\neq\varnothing)}{\FormulaRefAuto{SemigroupIsoImageCarrierData}[thm]{1,36}}
 \proofstepwidestar[1,35]{h[f[W]]=W}{\rAEa{\rAEb{37}}}
  \proofstepwidestar[1,35]{\Psi(\Phi(W))=h[f[W]]}{\text{Definitionen von \(\Phi,\Psi\)}}
  \proofstepwidestar[1,35]{\Psi(\Phi(W))=W}{\rIE{38,39}}
 \proofstepwidestar[1]{\Psi\circ \Phi:\mathcal S(A)\to \mathcal S(A)}{\FormulaRefAuto{CompositionDef}[def]{\rAEa{34},\rAEb{34}}}
  \proofstepwidestar[]{\Id_{\mathcal S(A)}:\mathcal S(A)\to \mathcal S(A)}{\FormulaRefAuto{IdA}[thm]}
  \proofstepwidestar[1,35]{(\Psi\circ \Phi)(W)=\Psi(\Phi(W))}{\FormulaRefAuto{CompositionDef}[def]{\rAEa{34},\rAEb{34},35}}
  \proofstepwidestar[35]{W=\Id_{\mathcal S(A)}(W)}{\FormulaRefAuto{x\in A\vdash x=\Id_A(x)}[thm]{35}}
  \proofstepwidestar[1,35]{(\Psi\circ \Phi)(W)=\Id_{\mathcal S(A)}(W)}{\rIE{40,44,43}}
  \proofstepwidestar[1]{\forall W\in \mathcal S(A)\;((\Psi\circ \Phi)(W)=\Id_{\mathcal S(A)}(W))}{\rUI{\rRI{35,45}}}
  \proofstepwidestar[1]{\Psi\circ \Phi=\Id_{\mathcal S(A)}}{\FormulaRefAuto{FunctionExtensionality}[thm]{41,42,46}}
 \proofstepwidestar[48]{W\in\mathcal S(B)}{\rA}
 \proofstepwidestar[1,48]{W\subseteq B}{\FormulaRefAuto{SemigroupSubstructurePosetsDef}[def]{48}}
 \proofstepwidestar[1,48]{f[h[W]]=W}{\FormulaRefAuto{Y\subseteq B\dsep F\colon A\bij B\vdash F[F^{-1}[Y]]=Y}[thm]{49,5}}
 \proofstepwidestar[1,48]{\Phi(\Psi(W))=f[h[W]]}{\text{Definitionen von \(\Phi,\Psi\)}}
  \proofstepwidestar[1,48]{\Phi(\Psi(W))=W}{\rIE{50,51}}
  \proofstepwidestar[1]{\Phi\circ \Psi:\mathcal S(B)\to \mathcal S(B)}{\FormulaRefAuto{CompositionDef}[def]{\rAEb{34},\rAEa{34}}}
  \proofstepwidestar[]{\Id_{\mathcal S(B)}:\mathcal S(B)\to \mathcal S(B)}{\FormulaRefAuto{IdA}[thm]}
  \proofstepwidestar[1,48]{(\Phi\circ \Psi)(W)=\Phi(\Psi(W))}{\FormulaRefAuto{CompositionDef}[def]{\rAEb{34},\rAEa{34},48}}
  \proofstepwidestar[48]{W=\Id_{\mathcal S(B)}(W)}{\FormulaRefAuto{x\in A\vdash x=\Id_A(x)}[thm]{48}}
  \proofstepwidestar[1,48]{(\Phi\circ \Psi)(W)=\Id_{\mathcal S(B)}(W)}{\rIE{52,56,55}}
  \proofstepwidestar[1]{\forall W\in \mathcal S(B)\;((\Phi\circ \Psi)(W)=\Id_{\mathcal S(B)}(W))}{\rUI{\rRI{48,57}}}
  \proofstepwidestar[1]{\Phi\circ \Psi=\Id_{\mathcal S(B)}}{\FormulaRefAuto{FunctionExtensionality}[thm]{53,54,58}}
 \proofstepwidestar[1]{\Phi:\mathcal S(A)\bij\mathcal S(B)}{\FormulaRefAuto{MutuallyInverseFunctionBijection}[thm]{\rAEb{34},\rAEa{34},59,47}}
 \proofstepwidestar[61]{U,V\subseteq A}{\rA}
 \proofstepwidestar[61]{U\subseteq A}{\rAEa{61}}
  \proofstepwidestar[61]{V\subseteq A}{\rAEb{61}}
  \proofstepwidestar[64]{U\subseteq V}{\rA}
  \proofstepwidestar[1,61,64]{f[U]\subseteq f[V]}{\FormulaRefAuto{F\colon M\to N\dsep A\subseteq B\dsep B\subseteq M\vdash F[A]\subseteq F[B]}[thm]{4,64,63}}
  \proofstepwidestar[1,61]{U\subseteq V\Rightarrow f[U]\subseteq f[V]}{\rRI{64,65}}
 \proofstepwidestar[1,61]{f[U]\subseteq B\land h[f[U]]=U\land(U\neq\varnothing\leftrightarrow f[U]\neq\varnothing)}{\FormulaRefAuto{SemigroupIsoImageCarrierData}[thm]{1,62}}
  \proofstepwidestar[1,61]{f[V]\subseteq B\land h[f[V]]=V\land(V\neq\varnothing\leftrightarrow f[V]\neq\varnothing)}{\FormulaRefAuto{SemigroupIsoImageCarrierData}[thm]{1,63}}
  \proofstepwidestar[69]{f[U]\subseteq f[V]}{\rA}
  \proofstepwidestar[1,61,69]{h[f[U]]\subseteq h[f[V]]}{\FormulaRefAuto{F\colon M\to N\dsep A\subseteq B\dsep B\subseteq M\vdash F[A]\subseteq F[B]}[thm]{9,69,\rAEa{68}}}
  \proofstepwidestar[1,61,69]{U\subseteq V}{\rIE{\rAEa{\rAEb{67}},\rAEa{\rAEb{68}},70}}
  \proofstepwidestar[1,61]{f[U]\subseteq f[V]\Rightarrow U\subseteq V}{\rRI{69,71}}
 \proofstepwidestar[1,61]{U\subseteq V\leftrightarrow f[U]\subseteq f[V]}{\rLRI{66,72}}
\closeproofpartwide
\proofpartwideindR[Monotonie des Relationsbildes]{%
\begin{aligned}[t]&f:A\cong B\dsep\rho,\sigma\in\operatorname{Con}(A)\dsep\rho\subseteq\sigma\\[-2pt]&\vdash f_*\rho\subseteq f_*\sigma\end{aligned}%
}{\begin{aligned}[t]&f:A\cong B\dsep\rho,\sigma\in\operatorname{Con}(A)\dsep\rho\subseteq\sigma\\[-2pt]&\vdash f_*\rho\subseteq f_*\sigma\end{aligned}}[SemigroupTransportedRelationMonotonicity]
  \proofstepwidestar[1]{f:A\cong B}{\rA}
  \proofstepwidestar[2]{\rho\in\operatorname{Con}(A)}{\rA}
  \proofstepwidestar[3]{\sigma\in\operatorname{Con}(A)}{\rA}
  \proofstepwidestar[1]{\SemigroupAlg{A}{\star}}{\FormulaRefAuto{SemigroupIsoSource}[thm]{1}}
  \proofstepwidestar[1]{\SemigroupAlg{B}{\diamond}}{\FormulaRefAuto{SemigroupIsoTarget}[thm]{1}}
  \proofstepwidestar[1]{f:A\to B}{\FormulaRefAuto{SemigroupIsoFunction}[thm]{1}}
  \proofstepwidestar[1]{f:A\bij B}{\FormulaRefAuto{SemigroupIsoBijectiveAxiom}[ax]{1}}
  \proofstepwidestar[1,2]{\EqRel{A,\rho}}{\FormulaRefAuto{SemigroupCongruenceQuotientDef}[def]{4,2}}
  \proofstepwidestar[1,2]{\rho\subseteq A\times A}{\FormulaRefAuto{EqRelDef}[def]{8}}
  \proofstepwidestar[1,3]{\EqRel{A,\sigma}}{\FormulaRefAuto{SemigroupCongruenceQuotientDef}[def]{4,3}}
  \proofstepwidestar[1,3]{\sigma\subseteq A\times A}{\FormulaRefAuto{EqRelDef}[def]{10}}
  \proofstepwidestar[12]{\rho\subseteq\sigma}{\rA}
  \proofstepwidestar[13]{p\in f_*\rho}{\rA}
  \proofstepwidestar[1,2,13]{\exists a\in A\,\exists b\in A\;(p=(f(a),f(b))\land(a,b)\in\rho)}{\FormulaRefAuto{SemigroupTransportedRelationDef}[def]{6,9,13}}
  \proofstepwidestar[15]{a\in A\land(b\in A\land(p=(f(a),f(b))\land(a,b)\in\rho))}{\rA}
  \proofstepwidestar[15]{a\in A}{\rAEa{15}}
  \proofstepwidestar[15]{b\in A}{\rAEa{\rAEb{15}}}
  \proofstepwidestar[15]{p=(f(a),f(b))}{\rAEa{\rAEb{\rAEb{15}}}}
  \proofstepwidestar[15]{(a,b)\in\rho}{\rAEb{\rAEb{\rAEb{15}}}}
  \proofstepwidestar[12,15]{(a,b)\in\sigma}{\FormulaRefAuto{A\subseteq B,\,x\in A\vdash x\in B}[thm]{12,19}}
  \proofstepwidestar[12,15]{a\in A\land(b\in A\land(p=(f(a),f(b))\land(a,b)\in\sigma))}{\rAI{16,\rAI{17,\rAI{18,20}}}}
  \proofstepwidestar[12,15]{\exists a\in A\,\exists b\in A\;(p=(f(a),f(b))\land(a,b)\in\sigma)}{\rEIStar{21}}
  \proofstepwidestar[1,3,12,15]{p\in f_*\sigma}{\FormulaRefAuto{SemigroupTransportedRelationDef}[def]{6,11,22}}
  \proofstepwidestar[1,2,3,12,13]{p\in f_*\sigma}{\rEEStar{14,15,23}}
  \proofstepwidestar[1,2,3,12]{f_*\rho\subseteq f_*\sigma}{\FormulaRefAuto{SubsetIntroductionRule}[thm]{[13]\vdots 24}}
\closeproofpartwideindR
\proofpartwideindR[Rücktransport einer Kongruenz]{%
f:A\cong B\dsep\rho\in\operatorname{Con}(A)\vdash h_*(f_*\rho)=\rho%
}{f:A\cong B\dsep\rho\in\operatorname{Con}(A)\vdash h_*(f_*\rho)=\rho}[SemigroupTransportedCongruenceRoundTrip]
  \proofstepwidestar[1]{f:A\cong B}{\rA}
  \proofstepwidestar[2]{\rho\in\operatorname{Con}(A)}{\rA}
  \proofstepwidestar[1]{\SemigroupAlg{A}{\star}}{\FormulaRefAuto{SemigroupIsoSource}[thm]{1}}
  \proofstepwidestar[1]{\SemigroupAlg{B}{\diamond}}{\FormulaRefAuto{SemigroupIsoTarget}[thm]{1}}
  \proofstepwidestar[1]{f:A\to B}{\FormulaRefAuto{SemigroupIsoFunction}[thm]{1}}
  \proofstepwidestar[1]{f:A\bij B}{\FormulaRefAuto{SemigroupIsoBijectiveAxiom}[ax]{1}}
  \proofstepwidestar[1,2]{\EqRel{A,\rho}}{\FormulaRefAuto{SemigroupCongruenceQuotientDef}[def]{3,2}}
  \proofstepwidestar[1,2]{\rho\subseteq A\times A}{\FormulaRefAuto{EqRelDef}[def]{7}}
  \proofstepwidestar[1]{h:B\cong A}{\FormulaRefAuto{SemigroupIsoInverse}[thm]{1}}
  \proofstepwidestar[1]{h:B\to A}{\FormulaRefAuto{SemigroupIsoFunction}[thm]{9}}
  \proofstepwidestar[1,2]{f_*\rho\in\operatorname{Con}(B)}{\FormulaRefAuto{SemigroupTransportedCongruence}[thm]{1,2}}
  \proofstepwidestar[1,2]{\EqRel{B,f_*\rho}}{\FormulaRefAuto{SemigroupCongruenceQuotientDef}[def]{4,11}}
  \proofstepwidestar[1,2]{f_*\rho\subseteq B\times B}{\FormulaRefAuto{EqRelDef}[def]{12}}
  \proofstepwidestar[1,2]{\forall u,v\in B\;(u\mathrel{f_*\rho}v\leftrightarrow h(u)\mathrel\rho h(v))}{\rAEn{\FormulaRefAuto{SemigroupTransportedRelationCriterion}[thm]{1,8}}}
  \proofstepwidestar[1,2]{\forall a,b\in A\;(a\mathrel\rho b\leftrightarrow f(a)\mathrel{f_*\rho}f(b))}{\rAEn{\FormulaRefAuto{SemigroupTransportedRelationCriterion}[thm]{1,8}}}
  \proofstepwidestar[16]{p\in h_*(f_*\rho)}{\rA}
  \proofstepwidestar[1,2,16]{\exists u\in B\,\exists v\in B\;(p=(h(u),h(v))\land u\mathrel{f_*\rho}v)}{\FormulaRefAuto{SemigroupTransportedRelationDef}[def]{10,13,16}}
  \proofstepwidestar[18]{u\in B\land(v\in B\land(p=(h(u),h(v))\land u\mathrel{f_*\rho}v))}{\rA}
  \proofstepwidestar[18]{u\in B}{\rAEa{18}}
  \proofstepwidestar[18]{v\in B}{\rAEa{\rAEb{18}}}
  \proofstepwidestar[18]{p=(h(u),h(v))}{\rAEa{\rAEb{\rAEb{18}}}}
  \proofstepwidestar[18]{u\mathrel{f_*\rho}v}{\rAEb{\rAEb{\rAEb{18}}}}
  \proofstepwidestar[1,2,18]{u\mathrel{f_*\rho}v\leftrightarrow h(u)\mathrel\rho h(v)}{\rRE{\rUE{\rRE{\rUE{14},19}},20}}
  \proofstepwidestar[1,2,18]{(h(u),h(v))\in\rho}{\rRE{\rLREa{23},22}}
  \proofstepwidestar[1,2,18]{p\in\rho}{\rIE{\FormulaRefAuto{a=b\vdash b=a}[thm]{21},24}}
  \proofstepwidestar[1,2,16]{p\in\rho}{\rEEStar{17,18,25}}
  \proofstepwidestar[1,2]{h_*(f_*\rho)\subseteq\rho}{\FormulaRefAuto{SubsetIntroductionRule}[thm]{[16]\vdots 26}}
  \proofstepwidestar[28]{q\in\rho}{\rA}
  \proofstepwidestar[1,2,28]{q\in A\times A}{\FormulaRefAuto{A\subseteq B,\,x\in A\vdash x\in B}[thm]{8,28}}
  \proofstepwidestar[1,2,28]{\exists a\in A\,\exists b\in A\;(q=(a,b))}{\FormulaRefAuto{x\in A\times B\eqvdash\exists a\in A\,\exists b\in B\,x=(a,b)}[thm]{29}}
  \proofstepwidestar[31]{a\in A\land(b\in A\land q=(a,b))}{\rA}
  \proofstepwidestar[31]{a\in A}{\rAEa{31}}
  \proofstepwidestar[31]{b\in A}{\rAEa{\rAEb{31}}}
  \proofstepwidestar[31]{q=(a,b)}{\rAEb{\rAEb{31}}}
  \proofstepwidestar[28,31]{a\mathrel\rho b}{\rIE{34,28}}
  \proofstepwidestar[1,2,31]{a\mathrel\rho b\leftrightarrow f(a)\mathrel{f_*\rho}f(b)}{\rRE{\rUE{\rRE{\rUE{15},32}},33}}
  \proofstepwidestar[1,2,28,31]{f(a)\mathrel{f_*\rho}f(b)}{\rRE{\rLREa{36},35}}
  \proofstepwidestar[1,31]{f(a)\in B}{\FormulaRefAuto{F\colon A\to B\dsep x\in A\vdash F(x)\in B}[thm]{5,32}}
  \proofstepwidestar[1,31]{f(b)\in B}{\FormulaRefAuto{F\colon A\to B\dsep x\in A\vdash F(x)\in B}[thm]{5,33}}
  \proofstepwidestar[1,31]{h(f(a))=a}{\FormulaRefAuto{InverseFunctionLeftInverse}[thm]{6,32}}
  \proofstepwidestar[1,31]{h(f(b))=b}{\FormulaRefAuto{InverseFunctionLeftInverse}[thm]{6,33}}
  \proofstepwidestar[1,31]{q=(h(f(a)),h(f(b)))}{\rIE{\FormulaRefAuto{a=b\vdash b=a}[thm]{40},\FormulaRefAuto{a=b\vdash b=a}[thm]{41},34}}
  \proofstepwidestar[1,2,28,31]{f(a)\in B\land(f(b)\in B\land(q=(h(f(a)),h(f(b)))\land f(a)\mathrel{f_*\rho}f(b)))}{\rAI{38,\rAI{39,\rAI{42,37}}}}
  \proofstepwidestar[1,2,28,31]{\exists u\in B\,\exists v\in B\;(q=(h(u),h(v))\land u\mathrel{f_*\rho}v)}{\rEIStar{43}}
  \proofstepwidestar[1,2,28,31]{q\in h_*(f_*\rho)}{\FormulaRefAuto{SemigroupTransportedRelationDef}[def]{10,13,44}}
  \proofstepwidestar[1,2,28]{q\in h_*(f_*\rho)}{\rEEStar{30,31,45}}
  \proofstepwidestar[1,2]{\rho\subseteq h_*(f_*\rho)}{\FormulaRefAuto{SubsetIntroductionRule}[thm]{[28]\vdots 46}}
  \proofstepwidestar[1,2]{h_*(f_*\rho)=\rho}{\FormulaRefAuto{A\subseteq B, B\subseteq A\vdash A=B}[thm]{27,47}}
\closeproofpartwideindR
\proofpartwide{Kongruenzen}
 \proofstepwidestar[1]{f:A\cong B}{\rA}
 \proofstepwidestar[1]{\SemigroupAlg{A}{\star}}{\FormulaRefAuto{SemigroupIsoSource}[thm]{1}}
 \proofstepwidestar[1]{\SemigroupAlg{B}{\diamond}}{\FormulaRefAuto{SemigroupIsoTarget}[thm]{1}}
 \proofstepwidestar[1]{f:A\to B}{\FormulaRefAuto{SemigroupIsoFunction}[thm]{1}}
 \proofstepwidestar[1]{f:A\bij B}{\FormulaRefAuto{SemigroupIsoBijectiveAxiom}[ax]{1}}
 \proofstepwidestar[1]{\SemigroupHom{f}{A,\star}{B,\diamond}}{\FormulaRefAuto{SemigroupIsoHomAxiom}[ax]{1}}
 \proofstepwidestar[1]{f[A]=B}{\FormulaRefAuto{SemigroupIsoFullImage}[thm]{1}}
 \proofstepwidestar[1]{h:B\cong A}{\FormulaRefAuto{SemigroupIsoInverse}[thm]{1}}
 \proofstepwidestar[9]{\rho\in\operatorname{Con}(A)}{\rA}
 \proofstepwidestar[1,9]{f_*\rho\in\operatorname{Con}(B)}{\FormulaRefAuto{SemigroupTransportedCongruence}[thm]{1,9}}
 \proofstepwidestar[1]{\forall\rho\in\operatorname{Con}(A)\;(f_*\rho\in\operatorname{Con}(B))}{\rUI{\rRI{9,10}}}
 \proofstepwidestar[12]{\tau\in\operatorname{Con}(B)}{\rA}
 \proofstepwidestar[1,12]{h_*\tau\in\operatorname{Con}(A)}{\FormulaRefAuto{SemigroupTransportedCongruence}[thm]{8,12}}
 \proofstepwidestar[1]{\forall\tau\in\operatorname{Con}(B)\;(h_*\tau\in\operatorname{Con}(A))}{\rUI{\rRI{12,13}}}
 \proofstepwidestar[1]{\Theta:\operatorname{Con}(A)\to\operatorname{Con}(B)\land\Lambda:\operatorname{Con}(B)\to\operatorname{Con}(A)}{\text{\parbox[t]{\linewidth}{\raggedright 11,14; \(\Theta(\rho)=f_*\rho\), \(\Lambda(\tau)=h_*\tau\).}}}
 \proofstepwidestar[16]{a\in A}{\rA}
  \proofstepwidestar[1,16]{h(f(a))=a}{\FormulaRefAuto{InverseFunctionLeftInverse}[thm]{5,16}}
  \proofstepwidestar[1]{\forall a\in A\;(h(f(a))=a)}{\rUI{\rRI{16,17}}}
 \proofstepwidestar[19]{b\in B}{\rA}
  \proofstepwidestar[1,19]{f(h(b))=b}{\FormulaRefAuto{InverseFunctionRightInverse}[thm]{5,19}}
  \proofstepwidestar[1]{\forall b\in B\;(f(h(b))=b)}{\rUI{\rRI{19,20}}}
 \proofstepwidestar[22]{\rho\in\operatorname{Con}(A)}{\rA}
  \proofstepwidestar[1,22]{h_*(f_*\rho)=\rho}{\FormulaRefAuto{SemigroupTransportedCongruenceRoundTrip}[thm]{1,22}}
  \proofstepwidestar[1,22]{\Lambda(\Theta(\rho))=h_*(f_*\rho)}{\text{Definitionen von \(\Theta,\Lambda\)}}
  \proofstepwidestar[1,22]{\Lambda(\Theta(\rho))=\rho}{\rIE{23,24}}
  \proofstepwidestar[1]{\Lambda\circ \Theta:\operatorname{Con}(A)\to \operatorname{Con}(A)}{\FormulaRefAuto{CompositionDef}[def]{\rAEa{15},\rAEb{15}}}
  \proofstepwidestar[]{\Id_{\operatorname{Con}(A)}:\operatorname{Con}(A)\to \operatorname{Con}(A)}{\FormulaRefAuto{IdA}[thm]}
  \proofstepwidestar[1,22]{(\Lambda\circ \Theta)(\rho)=\Lambda(\Theta(\rho))}{\FormulaRefAuto{CompositionDef}[def]{\rAEa{15},\rAEb{15},22}}
  \proofstepwidestar[22]{\rho=\Id_{\operatorname{Con}(A)}(\rho)}{\FormulaRefAuto{x\in A\vdash x=\Id_A(x)}[thm]{22}}
  \proofstepwidestar[1,22]{(\Lambda\circ \Theta)(\rho)=\Id_{\operatorname{Con}(A)}(\rho)}{\rIE{25,29,28}}
  \proofstepwidestar[1]{\forall \rho\in \operatorname{Con}(A)\;((\Lambda\circ \Theta)(\rho)=\Id_{\operatorname{Con}(A)}(\rho))}{\rUI{\rRI{22,30}}}
  \proofstepwidestar[1]{\Lambda\circ \Theta=\Id_{\operatorname{Con}(A)}}{\FormulaRefAuto{FunctionExtensionality}[thm]{26,27,31}}
 \proofstepwidestar[1]{h:B\bij A}{\FormulaRefAuto{SemigroupIsoBijectiveAxiom}[ax]{8}}
  \proofstepwidestar[1]{f=h^{-1}}{\FormulaRefAuto{InverseFunctionRightInverseUniqueness}[thm]{33,4,18}}
  \proofstepwidestar[1]{h^{-1}=f}{\FormulaRefAuto{a=b\vdash b=a}[thm]{34}}
  \proofstepwidestar[36]{\tau\in\operatorname{Con}(B)}{\rA}
  \proofstepwidestar[1,36]{(h^{-1})_*(h_*\tau)=\tau}{\FormulaRefAuto{SemigroupTransportedCongruenceRoundTrip}[thm]{8,36}}
  \proofstepwidestar[1,36]{f_*(h_*\tau)=\tau}{\rIE{35,37}}
  \proofstepwidestar[1,36]{\Theta(\Lambda(\tau))=f_*(h_*\tau)}{\text{Definitionen von \(\Theta,\Lambda\)}}
  \proofstepwidestar[1,36]{\Theta(\Lambda(\tau))=\tau}{\rIE{38,39}}
  \proofstepwidestar[1]{\Theta\circ \Lambda:\operatorname{Con}(B)\to \operatorname{Con}(B)}{\FormulaRefAuto{CompositionDef}[def]{\rAEb{15},\rAEa{15}}}
  \proofstepwidestar[]{\Id_{\operatorname{Con}(B)}:\operatorname{Con}(B)\to \operatorname{Con}(B)}{\FormulaRefAuto{IdA}[thm]}
  \proofstepwidestar[1,36]{(\Theta\circ \Lambda)(\tau)=\Theta(\Lambda(\tau))}{\FormulaRefAuto{CompositionDef}[def]{\rAEb{15},\rAEa{15},36}}
  \proofstepwidestar[36]{\tau=\Id_{\operatorname{Con}(B)}(\tau)}{\FormulaRefAuto{x\in A\vdash x=\Id_A(x)}[thm]{36}}
  \proofstepwidestar[1,36]{(\Theta\circ \Lambda)(\tau)=\Id_{\operatorname{Con}(B)}(\tau)}{\rIE{40,44,43}}
  \proofstepwidestar[1]{\forall \tau\in \operatorname{Con}(B)\;((\Theta\circ \Lambda)(\tau)=\Id_{\operatorname{Con}(B)}(\tau))}{\rUI{\rRI{36,45}}}
  \proofstepwidestar[1]{\Theta\circ \Lambda=\Id_{\operatorname{Con}(B)}}{\FormulaRefAuto{FunctionExtensionality}[thm]{41,42,46}}
 \proofstepwidestar[1]{\Theta:\operatorname{Con}(A)\bij\operatorname{Con}(B)}{\FormulaRefAuto{MutuallyInverseFunctionBijection}[thm]{\rAEb{15},\rAEa{15},47,32}}
 \proofstepwidestar[49]{\rho,\sigma\in\operatorname{Con}(A)}{\rA}
 \proofstepwidestar[50]{\rho\subseteq\sigma}{\rA}
  \proofstepwidestar[1,49,50]{f_*\rho\subseteq f_*\sigma}{\FormulaRefAuto{SemigroupTransportedRelationMonotonicity}[thm]{1,\rAEa{49},\rAEb{49},50}}
  \proofstepwidestar[1,49]{\rho\subseteq\sigma\Rightarrow f_*\rho\subseteq f_*\sigma}{\rRI{50,51}}
 \proofstepwidestar[1,49]{f_*\rho\in\operatorname{Con}(B)}{\FormulaRefAuto{SemigroupTransportedCongruence}[thm]{1,\rAEa{49}}}
  \proofstepwidestar[1,49]{f_*\sigma\in\operatorname{Con}(B)}{\FormulaRefAuto{SemigroupTransportedCongruence}[thm]{1,\rAEb{49}}}
  \proofstepwidestar[1,49]{h_*(f_*\rho)=\rho}{\FormulaRefAuto{SemigroupTransportedCongruenceRoundTrip}[thm]{1,\rAEa{49}}}
  \proofstepwidestar[1,49]{h_*(f_*\sigma)=\sigma}{\FormulaRefAuto{SemigroupTransportedCongruenceRoundTrip}[thm]{1,\rAEb{49}}}
  \proofstepwidestar[57]{f_*\rho\subseteq f_*\sigma}{\rA}
  \proofstepwidestar[1,49,57]{h_*(f_*\rho)\subseteq h_*(f_*\sigma)}{\FormulaRefAuto{SemigroupTransportedRelationMonotonicity}[thm]{8,53,54,57}}
  \proofstepwidestar[1,49,57]{\rho\subseteq\sigma}{\rIE{55,56,58}}
  \proofstepwidestar[1,49]{f_*\rho\subseteq f_*\sigma\Rightarrow\rho\subseteq\sigma}{\rRI{57,59}}
 \proofstepwidestar[1,49]{\rho\subseteq\sigma\leftrightarrow f_*\rho\subseteq f_*\sigma}{\rLRI{52,60}}
\closeproofpartwide
\end{tabproofsplitwide}
Folglich werden vorhandene kleinste und größte Elemente, minimale und
maximale Unterstrukturen sowie vorhandene Infima und Suprema transportiert:
Die definierenden Inklusionsbedingungen sind auf beiden Seiten äquivalent.
Bei den hier ausdrücklich nichtleeren Unterhalbgruppen und Idealen ist
ein leerer Durchschnitt nicht automatisch wieder eine Unterstruktur.

\subsection{Gleichungen und ihre Lösungsmengen}

\FormulaDefDeltaK[Halbgruppenterme mit Parametern]{%
 \begin{aligned}[t]
 &t::=x_i\mid c_a\mid(t\cdot t),\qquad i\in I,\ a\in A;\\[-2pt]
 &x_i^A(\alpha)\coloneqq\alpha(i),\quad c_a^A(\alpha)\coloneqq a,
       \quad(st)^A(\alpha)\coloneqq s^A(\alpha)\star t^A(\alpha).
 \end{aligned}%
}{SemigroupTermsWithParametersDef}{%
 \DeltaRow{Voraussetzung}{\SemigroupAlg{A}{\star}\,,\quad\alpha:I\to A}
 \DeltaRow{Konstruktion}{\text{\parbox[t]{.60\linewidth}{Terme sind endliche, nach diesen drei Regeln gebildete Ausdrücke.}}}
 \DeltaRow{Transport}{t^f\text{ ersetzt jeden Parameter }c_a\text{ durch }c_{f(a)}.}
}
Ein Gleichungssystem \(E\) ist eine Menge von Paaren solcher Terme.
Wir setzen
\[
 \begin{aligned}[t]
 \operatorname{Sol}_A(E)&\coloneqq
 \{\alpha:I\to A\mid\forall(s,t)\in E\;(s^A(\alpha)=t^A(\alpha))\},
 \\[-2pt]
 E^f&\coloneqq\{(s^f,t^f)\mid(s,t)\in E\}.
 \end{aligned}
\]

\FormulaThmDeltaKR[Isomorphismen transportieren alle Gleichungslösungen]{%
 \begin{aligned}[t]
  &f:A\cong B\vdash{}\\[-2pt]
  &\text{(i)}\quad f(t^A(\alpha))=(t^f)^B(f\circ\alpha)\\[-2pt]
  &\text{(ii)}\quad \alpha\in\operatorname{Sol}_A(E)\leftrightarrow f\circ\alpha\in\operatorname{Sol}_B(E^f)\\[-2pt]
  &\text{(iii)}\quad (\alpha\mapsto f\circ\alpha): \operatorname{Sol}_A(E)\bij\operatorname{Sol}_B(E^f)
 \end{aligned}%
}{%
 \begin{aligned}[t]
 &f:A\cong B\vdash{}\\[-2pt]
 &f(t^A(\alpha))=(t^f)^B(f\circ\alpha),\\[-2pt]
 &\alpha\in\operatorname{Sol}_A(E)\leftrightarrow
       f\circ\alpha\in\operatorname{Sol}_B(E^f),\\[-2pt]
 &(\alpha\mapsto f\circ\alpha):
       \operatorname{Sol}_A(E)\bij\operatorname{Sol}_B(E^f).
 \end{aligned}%
}{SemigroupIsoEquationSolutionTransport}{%
 \DeltaRow{Mengen}{A\dsep B\dsep I\dsep E}
 \DeltaRow{Funktionen}{f\dsep h\dsep\alpha:I\to A}
 \DeltaRow{Terme}{s\dsep t\text{ nach der vorstehenden Definition}}
}
\Needspace{8\baselineskip}
\begin{tabproofsplitwide}
\proofpartwideindR[Auswertungen liegen im Träger]{%
\SemigroupAlg{A}{\star}\dsep\alpha:I\to A\vdash t^A(\alpha)\in A%
}{\SemigroupAlg{A}{\star}\dsep\alpha:I\to A\vdash t^A(\alpha)\in A}[SemigroupTermEvaluationCarrier]
 \proofstepwidestar[1]{\SemigroupAlg{A}{\star}}{\rA}
 \proofstepwidestar[2]{\alpha:I\to A}{\rA}
 \proofstepwidestar[3]{i\in I}{\rA}
  \proofstepwidestar[2,3]{\alpha(i)\in A}{\FormulaRefAuto{F\colon A\to B\dsep x\in A\vdash F(x)\in B}[thm]{2,3}}
  \proofstepwidestar[2,3]{x_i^A(\alpha)=\alpha(i)}{\FormulaRefAuto{SemigroupTermsWithParametersDef}[def]{2,3}}
  \proofstepwidestar[2,3]{x_i^A(\alpha)\in A}{\rIE{\FormulaRefAuto{a=b\vdash b=a}[thm]{5},4}}
  \proofstepwidestar[2]{\forall i\in I\;(x_i^A(\alpha)\in A)}{\rUI{\rRI{3,6}}}
 \proofstepwidestar[8]{a\in A}{\rA}
  \proofstepwidestar[2,8]{c_a^A(\alpha)=a}{\FormulaRefAuto{SemigroupTermsWithParametersDef}[def]{2,8}}
  \proofstepwidestar[2,8]{c_a^A(\alpha)\in A}{\rIE{\FormulaRefAuto{a=b\vdash b=a}[thm]{9},8}}
  \proofstepwidestar[2]{\forall a\in A\;(c_a^A(\alpha)\in A)}{\rUI{\rRI{8,10}}}
 \proofstepwidestar[12]{s^A(\alpha),t^A(\alpha)\in A}{\rA}
 \proofstepwidestar[1,12]{s^A(\alpha)\star t^A(\alpha)\in A}{\FormulaRefAuto{SemigroupClosureAxiom}[ax]{1,12}}
 \proofstepwidestar[1,2,12]{(st)^A(\alpha)\in A}{\FormulaRefAuto{SemigroupTermsWithParametersDef}[def]{13}}
 \proofstepwidestar[1,2]{t^A(\alpha)\in A\quad\text{für jeden Term }t}{\text{\parbox[t]{\linewidth}{\raggedright Induktion über die endliche Konstruktion der Terme: 7,11 sind die Anfangsfälle, 12--14 der einzige Konstruktionsschritt.}}}
\closeproofpartwideindR
\proofpartwideindR[Strukturelle Induktion über einen Term]{%
f:A\cong B\dsep\alpha:I\to A\vdash f(t^A(\alpha))=(t^f)^B(f\circ\alpha)%
}{f:A\cong B\dsep\alpha:I\to A\vdash f(t^A(\alpha))=(t^f)^B(f\circ\alpha)}[SemigroupIsoTermEvaluation]
 \proofstepwidestar[1]{f:A\cong B}{\rA}
 \proofstepwidestar[1]{\SemigroupAlg{A}{\star}}{\FormulaRefAuto{SemigroupIsoSource}[thm]{1}}
 \proofstepwidestar[1]{\SemigroupAlg{B}{\diamond}}{\FormulaRefAuto{SemigroupIsoTarget}[thm]{1}}
 \proofstepwidestar[1]{f:A\to B}{\FormulaRefAuto{SemigroupIsoFunction}[thm]{1}}
 \proofstepwidestar[1]{f:A\bij B}{\FormulaRefAuto{SemigroupIsoBijectiveAxiom}[ax]{1}}
 \proofstepwidestar[1]{\SemigroupHom{f}{A,\star}{B,\diamond}}{\FormulaRefAuto{SemigroupIsoHomAxiom}[ax]{1}}
 \proofstepwidestar[1]{f[A]=B}{\FormulaRefAuto{SemigroupIsoFullImage}[thm]{1}}
 \proofstepwidestar[8]{\alpha:I\to A}{\rA}
 \proofstepwidestar[1,8]{f\circ\alpha:I\to B}{\FormulaRefAuto{CompositionDef}[def]{8,4}}
 \proofstepwidestar[10]{i\in I}{\rA}
  \proofstepwidestar[8,10]{x_i^A(\alpha)=\alpha(i)}{\FormulaRefAuto{SemigroupTermsWithParametersDef}[def]{8,10}}
  \proofstepwidestar[8,10]{f(x_i^A(\alpha))=f(\alpha(i))}{\rIE{11,\rII}}
  \proofstepwidestar[1,8,10]{(x_i^f)^B(f\circ\alpha)=(f\circ\alpha)(i)}{\FormulaRefAuto{SemigroupTermsWithParametersDef}[def]{9,10}}
  \proofstepwidestar[1,8,10]{(f\circ\alpha)(i)=f(\alpha(i))}{\FormulaRefAuto{CompositionDef}[def]{8,4,10}}
  \proofstepwidestar[1,8,10]{(x_i^f)^B(f\circ\alpha)=f(\alpha(i))}{\rIE{14,13}}
  \proofstepwidestar[1,8,10]{f(x_i^A(\alpha))=(x_i^f)^B(f\circ\alpha)}{\FormulaRefAuto{a=b, c=b\vdash a=c}[thm]{12,15}}
  \proofstepwidestar[1,8]{\forall i\in I\;(f(x_i^A(\alpha))=(x_i^f)^B(f\circ\alpha))}{\rUI{\rRI{10,16}}}
 \proofstepwidestar[18]{a\in A}{\rA}
  \proofstepwidestar[8,18]{c_a^A(\alpha)=a}{\FormulaRefAuto{SemigroupTermsWithParametersDef}[def]{8,18}}
  \proofstepwidestar[8,18]{f(c_a^A(\alpha))=f(a)}{\rIE{19,\rII}}
  \proofstepwidestar[1,18]{f(a)\in B}{\FormulaRefAuto{F\colon A\to B\dsep x\in A\vdash F(x)\in B}[thm]{4,18}}
  \proofstepwidestar[1,8,18]{(c_a^f)^B(f\circ\alpha)=f(a)}{\FormulaRefAuto{SemigroupTermsWithParametersDef}[def]{9,21}}
  \proofstepwidestar[1,8,18]{f(c_a^A(\alpha))=(c_a^f)^B(f\circ\alpha)}{\FormulaRefAuto{a=b, c=b\vdash a=c}[thm]{20,22}}
  \proofstepwidestar[1,8]{\forall a\in A\;(f(c_a^A(\alpha))=(c_a^f)^B(f\circ\alpha))}{\rUI{\rRI{18,23}}}
 \proofstepwidestar[25]{f(s^A(\alpha))=(s^f)^B(f\circ\alpha)}{\rA}
 \proofstepwidestar[26]{f(t^A(\alpha))=(t^f)^B(f\circ\alpha)}{\rA}
 \proofstepwidestar[1,8]{s^A(\alpha),t^A(\alpha)\in A}{\FormulaRefAuto{SemigroupTermEvaluationCarrier}[thm]{2,8}}
 \proofstepwidestar[1,8]{f(s^A(\alpha)\star t^A(\alpha))=f(s^A(\alpha))\diamond f(t^A(\alpha))}{\FormulaRefAuto{SemigroupIsoOperation}[thm]{1,27}}
 \proofstepwidestar[1,8,25,26]{f((st)^A(\alpha))=((st)^f)^B(f\circ\alpha)}{\FormulaRefAuto{SemigroupTermsWithParametersDef}[def]{28,25,26}}
 \proofstepwidestar[1,8]{f(t^A(\alpha))=(t^f)^B(f\circ\alpha)\quad\text{für jeden Term }t}{\text{\parbox[t]{\linewidth}{\raggedright Induktion über die endliche Termkonstruktion: 17,24 und 25--29}}}
\closeproofpartwideindR
\proofpartwide{Gleichheit, Systeme und Bijektivität}
 \proofstepwidestar[1]{f:A\cong B}{\rA}
 \proofstepwidestar[1]{\SemigroupAlg{A}{\star}}{\FormulaRefAuto{SemigroupIsoSource}[thm]{1}}
 \proofstepwidestar[1]{\SemigroupAlg{B}{\diamond}}{\FormulaRefAuto{SemigroupIsoTarget}[thm]{1}}
 \proofstepwidestar[1]{f:A\to B}{\FormulaRefAuto{SemigroupIsoFunction}[thm]{1}}
 \proofstepwidestar[1]{f:A\bij B}{\FormulaRefAuto{SemigroupIsoBijectiveAxiom}[ax]{1}}
 \proofstepwidestar[1]{\SemigroupHom{f}{A,\star}{B,\diamond}}{\FormulaRefAuto{SemigroupIsoHomAxiom}[ax]{1}}
 \proofstepwidestar[1]{f[A]=B}{\FormulaRefAuto{SemigroupIsoFullImage}[thm]{1}}
 \proofstepwidestar[8]{\alpha:I\to A}{\rA}
 \proofstepwidestar[1]{h:B\cong A}{\FormulaRefAuto{SemigroupIsoInverse}[thm]{1}}
 \proofstepwidestar[1]{h:B\to A}{\FormulaRefAuto{SemigroupIsoFunction}[thm]{9}}
 \proofstepwidestar[1,8]{s^A(\alpha),t^A(\alpha)\in A}{\FormulaRefAuto{SemigroupTermEvaluationCarrier}[thm]{2,8}}
 \proofstepwidestar[1,8]{s^A(\alpha)=t^A(\alpha)\leftrightarrow f(s^A(\alpha))=f(t^A(\alpha))}{\FormulaRefAuto{SemigroupIsoEqualityReflection}[thm]{1,11}}
 \proofstepwidestar[1,8]{f(s^A(\alpha))=(s^f)^B(f\circ\alpha)}{\FormulaRefAuto{SemigroupIsoTermEvaluation}[thm]{1,8}}
 \proofstepwidestar[1,8]{f(t^A(\alpha))=(t^f)^B(f\circ\alpha)}{\FormulaRefAuto{SemigroupIsoTermEvaluation}[thm]{1,8}}
 \proofstepwidestar[1,8]{s^A(\alpha)=t^A(\alpha)\leftrightarrow(s^f)^B(f\circ\alpha)=(t^f)^B(f\circ\alpha)}{\rIE{13,14,12}}
 \proofstepwidestar[1,8]{\alpha\in\operatorname{Sol}_A(E)\leftrightarrow f\circ\alpha\in\operatorname{Sol}_B(E^f)}{\text{\parbox[t]{\linewidth}{\raggedright Definition der Lösungsmengen; 15 für jedes Paar in \(E\).}}}
 \proofstepwidestar[1]{\forall\alpha:I\to A\;(\alpha\in\operatorname{Sol}_A(E)\leftrightarrow f\circ\alpha\in\operatorname{Sol}_B(E^f))}{\rUI{\rRI{8,16}}}
 \proofstepwidestar[18]{x\in A}{\rA}
  \proofstepwidestar[1,18]{h(f(x))=x}{\FormulaRefAuto{InverseFunctionLeftInverse}[thm]{5,18}}
  \proofstepwidestar[1]{\forall x\in A\;(h(f(x))=x)}{\rUI{\rRI{18,19}}}
 \proofstepwidestar[21]{y\in B}{\rA}
  \proofstepwidestar[1,21]{f(h(y))=y}{\FormulaRefAuto{InverseFunctionRightInverse}[thm]{5,21}}
  \proofstepwidestar[1]{\forall y\in B\;(f(h(y))=y)}{\rUI{\rRI{21,22}}}
 \proofstepwidestar[1,8]{f\circ \alpha:I\to B}{\FormulaRefAuto{CompositionDef}[def]{8,4}}
  \proofstepwidestar[1,8]{h\circ(f\circ \alpha):I\to A}{\FormulaRefAuto{CompositionDef}[def]{24,10}}
  \proofstepwidestar[26]{i\in I}{\rA}
  \proofstepwidestar[1,8,26]{\alpha(i)\in A}{\FormulaRefAuto{F\colon A\to B\dsep x\in A\vdash F(x)\in B}[thm]{8,26}}
  \proofstepwidestar[1,8,26]{h(f(\alpha(i)))=\alpha(i)}{\rRE{\rUE{20},27}}
  \proofstepwidestar[1,8,26]{(f\circ \alpha)(i)=f(\alpha(i))}{\FormulaRefAuto{CompositionDef}[def]{8,4,26}}
  \proofstepwidestar[1,8,26]{(h\circ(f\circ \alpha))(i)=h((f\circ \alpha)(i))}{\FormulaRefAuto{CompositionDef}[def]{24,10,26}}
  \proofstepwidestar[1,8,26]{(h\circ(f\circ \alpha))(i)=\alpha(i)}{\rIE{29,28,30}}
  \proofstepwidestar[1,8]{\forall i\in I\;((h\circ(f\circ \alpha))(i)=\alpha(i))}{\rUI{\rRI{26,31}}}
  \proofstepwidestar[1,8]{h\circ(f\circ \alpha)=\alpha}{\FormulaRefAuto{FunctionExtensionality}[thm]{25,8,32}}
 \proofstepwidestar[34]{\beta:I\to B}{\rA}
 \proofstepwidestar[1,34]{h\circ\beta:I\to A}{\FormulaRefAuto{CompositionDef}[def]{34,10}}
 \proofstepwidestar[1,34]{h\circ \beta:I\to A}{\FormulaRefAuto{CompositionDef}[def]{34,10}}
  \proofstepwidestar[1,34]{f\circ(h\circ \beta):I\to B}{\FormulaRefAuto{CompositionDef}[def]{36,4}}
  \proofstepwidestar[38]{i\in I}{\rA}
  \proofstepwidestar[1,34,38]{\beta(i)\in B}{\FormulaRefAuto{F\colon A\to B\dsep x\in A\vdash F(x)\in B}[thm]{34,38}}
  \proofstepwidestar[1,34,38]{f(h(\beta(i)))=\beta(i)}{\rRE{\rUE{23},39}}
  \proofstepwidestar[1,34,38]{(h\circ \beta)(i)=h(\beta(i))}{\FormulaRefAuto{CompositionDef}[def]{34,10,38}}
  \proofstepwidestar[1,34,38]{(f\circ(h\circ \beta))(i)=f((h\circ \beta)(i))}{\FormulaRefAuto{CompositionDef}[def]{36,4,38}}
  \proofstepwidestar[1,34,38]{(f\circ(h\circ \beta))(i)=\beta(i)}{\rIE{41,40,42}}
  \proofstepwidestar[1,34]{\forall i\in I\;((f\circ(h\circ \beta))(i)=\beta(i))}{\rUI{\rRI{38,43}}}
  \proofstepwidestar[1,34]{f\circ(h\circ \beta)=\beta}{\FormulaRefAuto{FunctionExtensionality}[thm]{37,34,44}}
 \proofstepwidestar[1,34]{h\circ\beta\in\operatorname{Sol}_A(E)\leftrightarrow f\circ(h\circ\beta)\in\operatorname{Sol}_B(E^f)}{\rRE{\rUE{17},35}}
  \proofstepwidestar[1,34]{h\circ\beta\in\operatorname{Sol}_A(E)\leftrightarrow\beta\in\operatorname{Sol}_B(E^f)}{\rIE{45,46}}
  \proofstepwidestar[1,34]{\beta\in\operatorname{Sol}_B(E^f)\Rightarrow h\circ\beta\in\operatorname{Sol}_A(E)}{\rLREb{47}}
 \proofstepwidestar[1]{\forall\beta:I\to B\;(\beta\in\operatorname{Sol}_B(E^f)\Rightarrow h\circ\beta\in\operatorname{Sol}_A(E))}{\rUI{\rRI{34,48}}}
  \proofstepwidestar[1]{\Phi:\operatorname{Sol}_A(E)\to\operatorname{Sol}_B(E^f)\land\Psi:\operatorname{Sol}_B(E^f)\to\operatorname{Sol}_A(E)}{\text{\parbox[t]{\linewidth}{\raggedright Definitionen von \(\Phi,\Psi\) aus 17,49}}}
 \proofstepwidestar[1]{\forall\alpha:I\to A\;(h\circ(f\circ\alpha)=\alpha)}{\rUI{\rRI{8,33}}}
  \proofstepwidestar[1]{\forall\beta:I\to B\;(f\circ(h\circ\beta)=\beta)}{\rUI{\rRI{34,45}}}
  \proofstepwidestar[53]{\alpha\in\operatorname{Sol}_A(E)}{\rA}
  \proofstepwidestar[53]{\alpha:I\to A}{\text{\parbox[t]{\linewidth}{\raggedright Definition von \(\operatorname{Sol}_A(E)\) aus 53}}}
  \proofstepwidestar[1,53]{h\circ(f\circ\alpha)=\alpha}{\rRE{\rUE{51},54}}
  \proofstepwidestar[1,53]{\Psi(\Phi(\alpha))=h\circ(f\circ\alpha)}{\text{Definitionen von \(\Phi,\Psi\)}}
  \proofstepwidestar[1,53]{\Psi(\Phi(\alpha))=\alpha}{\rIE{55,56}}
  \proofstepwidestar[1]{\Psi\circ \Phi:\operatorname{Sol}_A(E)\to \operatorname{Sol}_A(E)}{\FormulaRefAuto{CompositionDef}[def]{\rAEa{50},\rAEb{50}}}
  \proofstepwidestar[]{\Id_{\operatorname{Sol}_A(E)}:\operatorname{Sol}_A(E)\to \operatorname{Sol}_A(E)}{\FormulaRefAuto{IdA}[thm]}
  \proofstepwidestar[1,53]{(\Psi\circ \Phi)(\alpha)=\Psi(\Phi(\alpha))}{\FormulaRefAuto{CompositionDef}[def]{\rAEa{50},\rAEb{50},53}}
  \proofstepwidestar[53]{\alpha=\Id_{\operatorname{Sol}_A(E)}(\alpha)}{\FormulaRefAuto{x\in A\vdash x=\Id_A(x)}[thm]{53}}
  \proofstepwidestar[1,53]{(\Psi\circ \Phi)(\alpha)=\Id_{\operatorname{Sol}_A(E)}(\alpha)}{\rIE{57,61,60}}
  \proofstepwidestar[1]{\forall \alpha\in \operatorname{Sol}_A(E)\;((\Psi\circ \Phi)(\alpha)=\Id_{\operatorname{Sol}_A(E)}(\alpha))}{\rUI{\rRI{53,62}}}
  \proofstepwidestar[1]{\Psi\circ \Phi=\Id_{\operatorname{Sol}_A(E)}}{\FormulaRefAuto{FunctionExtensionality}[thm]{58,59,63}}
  \proofstepwidestar[65]{\beta\in\operatorname{Sol}_B(E^f)}{\rA}
  \proofstepwidestar[65]{\beta:I\to B}{\text{\parbox[t]{\linewidth}{\raggedright Definition von \(\operatorname{Sol}_B(E^f)\) aus 65}}}
  \proofstepwidestar[1,65]{f\circ(h\circ\beta)=\beta}{\rRE{\rUE{52},66}}
  \proofstepwidestar[1,65]{\Phi(\Psi(\beta))=f\circ(h\circ\beta)}{\text{Definitionen von \(\Phi,\Psi\)}}
  \proofstepwidestar[1,65]{\Phi(\Psi(\beta))=\beta}{\rIE{67,68}}
  \proofstepwidestar[1]{\Phi\circ \Psi:\operatorname{Sol}_B(E^f)\to \operatorname{Sol}_B(E^f)}{\FormulaRefAuto{CompositionDef}[def]{\rAEb{50},\rAEa{50}}}
  \proofstepwidestar[]{\Id_{\operatorname{Sol}_B(E^f)}:\operatorname{Sol}_B(E^f)\to \operatorname{Sol}_B(E^f)}{\FormulaRefAuto{IdA}[thm]}
  \proofstepwidestar[1,65]{(\Phi\circ \Psi)(\beta)=\Phi(\Psi(\beta))}{\FormulaRefAuto{CompositionDef}[def]{\rAEb{50},\rAEa{50},65}}
  \proofstepwidestar[65]{\beta=\Id_{\operatorname{Sol}_B(E^f)}(\beta)}{\FormulaRefAuto{x\in A\vdash x=\Id_A(x)}[thm]{65}}
  \proofstepwidestar[1,65]{(\Phi\circ \Psi)(\beta)=\Id_{\operatorname{Sol}_B(E^f)}(\beta)}{\rIE{69,73,72}}
  \proofstepwidestar[1]{\forall \beta\in \operatorname{Sol}_B(E^f)\;((\Phi\circ \Psi)(\beta)=\Id_{\operatorname{Sol}_B(E^f)}(\beta))}{\rUI{\rRI{65,74}}}
  \proofstepwidestar[1]{\Phi\circ \Psi=\Id_{\operatorname{Sol}_B(E^f)}}{\FormulaRefAuto{FunctionExtensionality}[thm]{70,71,75}}
  \proofstepwidestar[1]{\Psi\circ\Phi=\Id_{\operatorname{Sol}_A(E)}\land\Phi\circ\Psi=\Id_{\operatorname{Sol}_B(E^f)}}{\rAI{64,76}}
 \proofstepwidestar[1]{(\alpha\mapsto f\circ\alpha):\operatorname{Sol}_A(E)\bij\operatorname{Sol}_B(E^f)}{\FormulaRefAuto{MutuallyInverseFunctionBijection}[thm]{\rAEb{50},\rAEa{50},\rAEb{77},\rAEa{77}}}
\closeproofpartwide
\end{tabproofsplitwide}
Dies erfasst insbesondere \(x^2=x\), \(xy=yx\), \(ax=b\), \(xa=b\)
und \(xy=e=yx\). Die Lösungsmengen sind im Allgemeinen nur Mengen.
Wenn eine solche Lösungsmenge unter der koordinatenweisen Operation
abgeschlossen ist, liefert zusätzlich der Einschränkungssatz einen
Halbgruppenisomorphismus. Die Äquivalenz für Gleichungen bleibt bei
Negation, Konjunktion, Disjunktion und Quantifizierung erhalten:
Für Existenzquantoren transportiert man den Zeugen mit \(f\), für
Allquantoren verwendet man die Surjektivität. Damit werden auch
Kürzbarkeit und andere durch solche Formeln beschriebene Eigenschaften
in beiden Richtungen erhalten.


\section{Das Isomorphieproblem für Potenzhalbgruppen}

Zwei Halbgruppen heißen \emph{global isomorph}, wenn ihre
Potenzhalbgruppen isomorph sind.  Die folgende Richtung gilt ohne jede
Endlichkeitsvoraussetzung.

\FormulaThmDeltaK[Die isomorphe Richtung des Potenzhalbgruppenproblems]{%
  \begin{aligned}[t]
    &\exists f\,\SemigroupIso{f}{A,\star}{B,\diamond}\vdash{}\\[-2pt]
    &\exists F\,\SemigroupIso{F}
      {\PowNE{A},\star_{\mathcal P}}
      {\PowNE{B},\diamond_{\mathcal P}}
  \end{aligned}%
}{SemigroupIsoImpliesPowerSemigroupIso}{%
  \DeltaRow{Mengen}{A\dsep B}
  \DeltaRow{Funktionen}{F\dsep f\dsep\star\dsep\diamond\dsep
    \star_{\mathcal P}\dsep\diamond_{\mathcal P}\dsep
    \PowMap{f}\dsep\PowMapNE{f}}
}
\begin{tabproofwide}
  \proofstepwidestar[1]{\exists f\,\SemigroupIso{f}{A,\star}{B,\diamond}}{\rA}
  \proofstepwidestar[2]{\SemigroupIso{f}{A,\star}{B,\diamond}}{\rA}
  \proofstepwidestar[2]{%
    \SemigroupIso{\PowMapNE{f}}
      {\PowNE{A},\star_{\mathcal P}}
      {\PowNE{B},\diamond_{\mathcal P}}}{%
    \FormulaRefAuto{SemigroupIsoInducesPowerSemigroupIso}[thm]{2}}
  \proofstepwidestar[2]{%
    \exists F\,\SemigroupIso{F}
      {\PowNE{A},\star_{\mathcal P}}
      {\PowNE{B},\diamond_{\mathcal P}}}{\rEI{3}}
  \proofstepwidestar[1]{%
    \exists F\,\SemigroupIso{F}
      {\PowNE{A},\star_{\mathcal P}}
      {\PowNE{B},\diamond_{\mathcal P}}}{\rEE{1,2,4}}
\end{tabproofwide}

\subsection{Inflationen und translationelle Ununterscheidbarkeit}

Eine Erweiterung kann neue Elemente enthalten, ohne neue Multiplikationswerte
zu erzeugen.  Der folgende Begriff fasst diese Situation zusammen.  Die
Abbildung \(r\) zieht jedes neue Element auf ein altes Element zurück; die
Multiplikation der Erweiterung hängt nur von diesen Rückzugswerten ab.

\FormulaDefDeltaK[Inflation einer Halbgruppe]{%
  \mathsf{Infl}_{r}\bigl((A,\star),(B,\diamond)\bigr)%
}{SemigroupInflationDef}{%
  \DeltaRow{Mengen}{A\dsep B\dsep x\dsep y\dsep a}
  \DeltaRow{Funktionen}{r\dsep\star\dsep\diamond}
  \DeltaRow{\textbf{Axiome}}{}
  \DeltaPrem{Quellhalbgruppe}{\SemigroupAlg{A}{\star}}
  \DeltaPrem{Zielhalbgruppe}{\SemigroupAlg{B}{\diamond}}
  \DeltaPrem{Erweiterung}{A\subseteq B}
  \DeltaPrem{Rückzug}{r\colon B\to A}
  \DeltaRow{Retraktion}{a\in A\vdash r(a)=a}
  \DeltaRow{Produkt über dem Rückzug}{%
    x,y\in B\vdash x\diamond y=r(x)\star r(y)}
  \DeltaRow{\textbf{Neues Symbol}}{}
  \DeltaPrem{Inflation}{%
    \mathsf{Infl}_{r}\bigl((A,\star),(B,\diamond)\bigr)}
}

\paragraph{Isomorphismen aus passenden Inflationsfasern.}
Die Multiplikation einer Inflation hängt nur von den Rückzugswerten ab.
Deshalb genügt es, passende Faserbijektionen zu verkleben und dabei den
gegebenen Isomorphismus auf der Grundhalbgruppe beizubehalten.
Die mengentheoretische Verklebung ist in Band~09 bereits bewiesen.

\par\noindent\begin{minipage}{\linewidth}
\FormulaThmDeltaK[Rekonstruktion aus passenden Inflationsfasern]{%
  \begin{aligned}[t]
    &\mathsf{Infl}_{r}((G,\star),(S,\star))\dsep
      \mathsf{Infl}_{s}((H,\diamond),(T,\diamond))\dsep{}\\[-2pt]
    &\SemigroupIso{\varphi}{G,\star}{H,\diamond}\dsep{}\\[-2pt]
    &\forall g\in G\,\exists u\,
      \bigl(u\colon r^{-1}[\{g\}]\bij s^{-1}[\{\varphi(g)\}]\bigr)\vdash{}\\[-2pt]
    &\exists f\,\SemigroupIso{f}{S,\star}{T,\diamond}
  \end{aligned}%
}{InflationFiberIsomorphism}{%
  \DeltaRow{Mengen}{S\dsep T\dsep G\dsep H\dsep x\dsep y\dsep g}
  \DeltaRow{Funktionen}{\star\dsep\diamond\dsep r\dsep s\dsep\varphi\dsep f\dsep u}
}
\end{minipage}\par
\begin{tabproofsplitwide}
\proofpartwide{Konstruktion einer kompatiblen Bijektion}
  \proofstepwidestarbreakkeep[1]{\mathsf{Infl}_{r}((G,\star),(S,\star))}{\rA}
  \proofstepwidestarbreakkeep[2]{\mathsf{Infl}_{s}((H,\diamond),(T,\diamond))}{\rA}
  \proofstepwidestarbreakkeep[3]{\SemigroupIso{\varphi}{G,\star}{H,\diamond}}{\rA}
  \proofstepwidestarbreakkeep[4]{\forall g\in G\,\exists u\,
    \bigl(u\colon r^{-1}[\{g\}]\bij s^{-1}[\{\varphi(g)\}]\bigr)}{\rA}
  \proofstepwidestarbreakkeep[1]{G\subseteq S\land r\colon S\to G\land
    \forall g\in G\,(r(g)=g)}{\FormulaRefAuto{SemigroupInflationDef}[def]{1}}
  \proofstepwidestarbreakkeep[2]{H\subseteq T\land s\colon T\to H\land
    \forall h\in H\,(s(h)=h)}{\FormulaRefAuto{SemigroupInflationDef}[def]{2}}
  \proofstepwidestarbreakkeep[3]{\varphi\colon G\bij H}{%
    \FormulaRefAuto{SemigroupIsoBijectiveAxiom}[ax]{3}}
  \proofstepwidestarbreakkeep[1,2,3,4]{\exists f\,\Bigl(
    \begin{aligned}[t]
      &f\colon S\bij T\land{}\\[-2pt]
      &\forall x\in S\,(s(f(x))=\varphi(r(x)))\land{}\\[-2pt]
      &\forall g\in G\,(f(g)=\varphi(g))
    \end{aligned}\Bigr)}{%
    \FormulaRefAuto{RetractionFiberBijection}[thm]{5,\allowbreak 6,\allowbreak 7,\allowbreak 4}}
  \proofstepwidestarbreakkeep[9]{\begin{aligned}[t]
    &f\colon S\bij T\land{}\\[-2pt]
    &\forall x\in S\,(s(f(x))=\varphi(r(x)))\land{}\\[-2pt]
    &\forall g\in G\,(f(g)=\varphi(g))
  \end{aligned}}{\rA}

\closeproofpartwide
\proofpartwide{Trägerzugriffe und Produktwerte}
\setcounter{proofstepnr}{9}
  \proofstepwidestarbreakkeep[1]{\SemigroupAlg{G}{\star}\land\SemigroupAlg{S}{\star}}{%
    \FormulaRefAuto{SemigroupInflationDef}[def]{1}}
  \proofstepwidestarbreakkeep[2]{\SemigroupAlg{H}{\diamond}\land\SemigroupAlg{T}{\diamond}}{%
    \FormulaRefAuto{SemigroupInflationDef}[def]{2}}
  \proofstepwidestarbreakkeep[12]{x,y\in S}{\rA}
  \proofstepwidestarbreakkeep[1,12]{r(x),r(y)\in G}{%
    \FormulaRefAuto{F\colon A\to B\dsep x\in A\vdash F(x)\in B}[thm]{5,\allowbreak 12}}
  \proofstepwidestarbreakkeep[9,12]{f(x),f(y)\in T}{%
    \FormulaRefAuto{F\colon A\bij B\dsep x\in A\vdash F(x)\in B}[thm]{\rAEa{9},\allowbreak 12}}
  \proofstepwidestarbreakkeep[1,12]{x\star y=r(x)\star r(y)}{%
    \FormulaRefAuto{SemigroupInflationDef}[def]{1,\allowbreak 12}}
  \proofstepwidestarbreakkeep[1,12]{r(x)\star r(y)\in G}{%
    \FormulaRefAuto{SemigroupClosureAxiom}[ax]{10,\allowbreak 13}}
  \proofstepwidestarbreakkeep[9]{\forall g\in G\,(f(g)=\varphi(g))}{\rAEb{\rAEb{9}}}
  \proofstepwidestarbreakkeep[9]{\forall x\in S\,(s(f(x))=\varphi(r(x)))}{\rAEa{\rAEb{9}}}

\closeproofpartwide
\proofpartwide{Operationserhaltung und Existenzschluss}
\setcounter{proofstepnr}{18}
  \proofstepwidestarbreakkeep[1,9,12]{f(x\star y)=f(r(x)\star r(y))}{\rIE{15}}
  \proofstepwidestarbreakkeep[1,9,12]{\phantom{f(x\star y)}=\varphi(r(x)\star r(y))}{\rRE{16,\allowbreak \rUE{17}}}
  \proofstepwidestar[1,3,9,12]{\SemigroupHom{\varphi}{G,\star}{H,\diamond}}{%
      \FormulaRefAuto{SemigroupIsoHomAxiom}[ax]{3}}
  \proofstepwidestarbreakkeep[1,3,9,12]{\varphi(r(x)\star r(y))=\varphi(r(x))\diamond\varphi(r(y))}{%
    \FormulaRefAuto{SemigroupHomOperationAxiom}[ax]{%
      21,\allowbreak 13}}
  \proofstepwidestarbreakkeep[1,3,9,12]{\varphi(r(x))\diamond\varphi(r(y))=s(f(x))\diamond s(f(y))}{\rIE{\rRE{12,\allowbreak \rUE{18}}}}
  \proofstepwidestarbreakkeep[1,2,3,9,12]{s(f(x))\diamond s(f(y))=f(x)\diamond f(y)}{%
    \FormulaRefAuto{SemigroupInflationDef}[def]{2,\allowbreak 14}}
  \proofstepwidestarbreakkeep[1,2,3,9,12]{f(x\star y)=f(x)\diamond f(y)}{\rChain{19,20,22,23,24}}
  \proofstepwidestar[1,2,3,9]{f\colon S\to T}{%
      \FormulaRefAuto{Bijektive Funktion}[def]{\rAEa{9}}}
  \proofstepwidestarbreakkeep[1,2,3,9]{\SemigroupHom{f}{S,\star}{T,\diamond}}{%
    \FormulaRefAuto{SemigroupHomDef}[def]{10,\allowbreak 11,\allowbreak %
      26,\allowbreak 25}}
  \proofstepwidestarbreakkeep[1,2,3,9]{\SemigroupIso{f}{S,\star}{T,\diamond}}{%
    \FormulaRefAuto{SemigroupIsoDef}[def]{27,\allowbreak \rAEa{9}}}
  \proofstepwidestarbreakkeep[1,2,3,9]{\exists f\,\SemigroupIso{f}{S,\star}{T,\diamond}}{\rEI{28}}
  \proofstepwidestarbreakkeep[1,2,3,4]{\exists f\,\SemigroupIso{f}{S,\star}{T,\diamond}}{\rEE{8,\allowbreak 9,\allowbreak 29}}
\closeproofpartwide
\end{tabproofsplitwide}

Für die nachfolgende Anwendung ist nicht die Gleichheit zweier Elemente
entscheidend, sondern die Gleichheit ihrer Links- und Rechtstranslationen.

\FormulaDefDeltaK[Translationelle Ununterscheidbarkeit]{%
  \begin{aligned}[t]
    &\SemigroupAlg{B}{\diamond}\dsep x,y\in B\vdash{}\\[-2pt]
    &x\mathrel{\equiv^{\mathrm{tr}}_{B,\diamond}}y
      \coloneqq
      \forall z\in B\,\bigl(
        x\diamond z=y\diamond z\land
        z\diamond x=z\diamond y\bigr)
  \end{aligned}%
}{SemigroupTranslationIndistinguishableDef}{%
  \DeltaRow{Mengen}{B\dsep x\dsep y\dsep z}
  \DeltaRow{Funktionen}{\diamond}
  \DeltaRow{Relationsoperatoren}{%
    \equiv^{\mathrm{tr}}_{B,\diamond}}
}

\FormulaThmDeltaK[Zugriff auf die gemeinsamen Translationen]{%
  \begin{aligned}[t]
    &\SemigroupAlg{B}{\diamond}\dsep
      x,y\in B\dsep
      x\mathrel{\equiv^{\mathrm{tr}}_{B,\diamond}}y\dsep z\in B
      \vdash{}\\[-2pt]
    &x\diamond z=y\diamond z\land z\diamond x=z\diamond y
  \end{aligned}%
}{SemigroupTranslationIndistinguishableActions}{%
  \DeltaRow{Mengen}{B\dsep x\dsep y\dsep z}
  \DeltaRow{Funktionen}{\diamond}
}
\begin{tabproofwide}
  \proofstepwidestar[1]{\SemigroupAlg{B}{\diamond}}{\rA}
  \proofstepwidestar[2]{x\in B}{\rA}
  \proofstepwidestar[3]{y\in B}{\rA}
  \proofstepwidestar[4]{%
    x\mathrel{\equiv^{\mathrm{tr}}_{B,\diamond}}y}{\rA}
  \proofstepwidestar[5]{z\in B}{\rA}
  \proofstepwidestar[1,2,3,4]{%
    \forall u\in B\,\bigl(
      x\diamond u=y\diamond u\land
      u\diamond x=u\diamond y\bigr)}{%
    \FormulaRefAuto{SemigroupTranslationIndistinguishableDef}[def]{1,2,3,4}}
  \proofstepwidestar[1,2,3,4,5]{%
    x\diamond z=y\diamond z\land z\diamond x=z\diamond y}{%
    \rRE{5,\rUE{6}}}
\end{tabproofwide}

Zunächst halten wir die für die Anwendung benötigte Reflexivität des
Translationsprofils fest.

\FormulaThmDeltaK[Reflexivität der translationellen Ununterscheidbarkeit]{%
  \SemigroupAlg{B}{\diamond}\dsep x\in B\vdash
  x\mathrel{\equiv^{\mathrm{tr}}_{B,\diamond}}x%
}{SemigroupTranslationIndistinguishableReflexive}{%
  \DeltaRow{Mengen}{B\dsep x\dsep z}
  \DeltaRow{Funktionen}{\diamond}
  \DeltaRow{Relationsoperatoren}{\equiv^{\mathrm{tr}}_{B,\diamond}}
}
\begin{tabproofwide}
  \proofstepwidestar[1]{\SemigroupAlg{B}{\diamond}}{\rA}
  \proofstepwidestar[2]{x\in B}{\rA}
  \proofstepwidestar[3]{z\in B}{\rA}
  \proofstepwidestar[3]{x\diamond z=x\diamond z}{\rII}
  \proofstepwidestar[3]{z\diamond x=z\diamond x}{\rII}
  \proofstepwidestar[3]{%
    x\diamond z=x\diamond z\land z\diamond x=z\diamond x}{\rAI{4,5}}
  \proofstepwidestar[]{%
    \forall z\in B\,\bigl(
      x\diamond z=x\diamond z\land z\diamond x=z\diamond x\bigr)}{%
    \rUI{\rRI{3,6}}}
  \proofstepwidestar[1,2]{%
    x\mathrel{\equiv^{\mathrm{tr}}_{B,\diamond}}x}{%
    \FormulaRefAuto{SemigroupTranslationIndistinguishableDef}[def]{1,2,2,7}}
\end{tabproofwide}

Der folgende Satz ist das algebraische Bindeglied zur späteren
Reservefamilie.  Er verlangt nicht, dass die Bijektion die einzelnen
Elemente festhält.  Es genügt, sie nur innerhalb ihrer Translationsprofile
zu verschieben und sämtliche wirklichen Produktwerte festzulassen.

\FormulaThmDeltaK[Isomorphiekriterium durch Fixierung der Produktwerte]{%
  \begin{aligned}[t]
    &\SemigroupAlg{A}{\star}\dsep\SemigroupAlg{B}{\diamond}\dsep
      A\subseteq B\dsep F\colon A\bij B\dsep{}\\[-2pt]
    &\forall x,y\in A\,(x\star y=x\diamond y)\dsep{}\\[-2pt]
    &\forall x\in A\,
      \bigl(x\mathrel{\equiv^{\mathrm{tr}}_{B,\diamond}}F(x)\bigr)
      \dsep{}\\[-2pt]
    &\forall x,y\in A\,F(x\star y)=x\star y
      \vdash \SemigroupIso{F}{A,\star}{B,\diamond}
  \end{aligned}%
}{SemigroupFixedProductsIsomorphismCriterion}{%
  \DeltaRow{Mengen}{A\dsep B\dsep x\dsep y}
  \DeltaRow{Funktionen}{F\dsep\star\dsep\diamond}
}
\begin{tabproofsplitwide}
\proofpartwide{Trägerzugriffe und Translationsprofile}
  \proofstepwidestar[1]{\SemigroupAlg{A}{\star}}{\rA}
  \proofstepwidestar[2]{\SemigroupAlg{B}{\diamond}}{\rA}
  \proofstepwidestar[3]{A\subseteq B}{\rA}
  \proofstepwidestar[4]{F\colon A\bij B}{\rA}
  \proofstepwidestar[5]{\forall u,v\in A\,(u\star v=u\diamond v)}{\rA}
  \proofstepwidestar[6]{\forall u\in A\,
    \bigl(u\mathrel{\equiv^{\mathrm{tr}}_{B,\diamond}}F(u)\bigr)}{\rA}
  \proofstepwidestar[7]{\forall u,v\in A\,F(u\star v)=u\star v}{\rA}
  \proofstepwidestar[4]{F\colon A\to B}{%
    \FormulaRefAuto{Bijektive Funktion}[def]{4}}
  \proofstepwidestar[9]{x\in A}{\rA}
  \proofstepwidestar[10]{y\in A}{\rA}
  \proofstepwidestar[3,9]{x\in B}{%
    \FormulaRefAuto{A\subseteq B\dsep x\in A\vdash x\in B}[thm]{3,9}}
  \proofstepwidestar[3,10]{y\in B}{%
    \FormulaRefAuto{A\subseteq B\dsep x\in A\vdash x\in B}[thm]{3,10}}
  \proofstepwidestar[4,9]{F(x)\in B}{%
    \FormulaRefAuto{F\colon A\to B\dsep x\in A\vdash F(x)\in B}[thm]{8,9}}
  \proofstepwidestar[4,10]{F(y)\in B}{%
    \FormulaRefAuto{F\colon A\to B\dsep x\in A\vdash F(x)\in B}[thm]{8,10}}
  \proofstepwidestar[6,9]{%
    x\mathrel{\equiv^{\mathrm{tr}}_{B,\diamond}}F(x)}{\rRE{9,\rUE{6}}}
  \proofstepwidestar[6,10]{%
    y\mathrel{\equiv^{\mathrm{tr}}_{B,\diamond}}F(y)}{\rRE{10,\rUE{6}}}

\closeproofpartwide
\proofpartwide{Produktvergleich und Isomorphieschluss}
\setcounter{proofstepnr}{16}
  \proofstepwidestar[2,3,4,6,9,10]{x\diamond F(y)=F(x)\diamond F(y)\land F(y)\diamond x=F(y)\diamond F(x)}{%
      \FormulaRefAuto{SemigroupTranslationIndistinguishableActions}[thm]{%
      2,11,13,15,14}}
  \proofstepwidestar[2,3,4,6,9,10]{%
    x\diamond F(y)=F(x)\diamond F(y)}{%
    \rAEa{17}}
  \proofstepwidestar[2,3,4,6,9,10]{%
    F(x)\diamond F(y)=x\diamond F(y)}{%
    \FormulaRefAuto{a=b\vdash b=a}[thm]{18}}
  \proofstepwidestar[2,3,4,6,9,10]{y\diamond x=F(y)\diamond x\land x\diamond y=x\diamond F(y)}{%
      \FormulaRefAuto{SemigroupTranslationIndistinguishableActions}[thm]{%
      2,12,14,16,11}}
  \proofstepwidestar[2,3,4,6,9,10]{x\diamond y=x\diamond F(y)}{%
    \rAEb{20}}
  \proofstepwidestar[2,3,4,6,9,10]{x\diamond F(y)=x\diamond y}{%
    \FormulaRefAuto{a=b\vdash b=a}[thm]{21}}
  \proofstepwidestar[5,9,10]{x\star y=x\diamond y}{%
    \rRE{10,\rUE{\rRE{9,\rUE{5}}}}}
  \proofstepwidestar[5,9,10]{x\diamond y=x\star y}{%
    \FormulaRefAuto{a=b\vdash b=a}[thm]{23}}
  \proofstepwidestar[7,9,10]{F(x\star y)=x\star y}{%
    \rRE{10,\rUE{\rRE{9,\rUE{7}}}}}
  \proofstepwidestar[7,9,10]{x\star y=F(x\star y)}{%
    \FormulaRefAuto{a=b\vdash b=a}[thm]{25}}
  \proofstepwidestar[2,3,4,5,6,7,9,10]{%
    F(x\star y)=F(x)\diamond F(y)}{%
    \rChain{25,23,21,18}}
  \proofstepwidestar[1,2,3,4,5,6,7]{%
    \SemigroupHom{F}{A,\star}{B,\diamond}}{%
    \FormulaRefAuto{SemigroupHomDef}[def]{1,2,8,27}}
  \proofstepwidestar[1,2,3,4,5,6,7]{%
    \SemigroupIso{F}{A,\star}{B,\diamond}}{%
    \FormulaRefAuto{SemigroupIsoDef}[def]{28,4}}
\closeproofpartwide
\end{tabproofsplitwide}

\FormulaThmDeltaK[Gleiche Translationsprofile in der Potenzhalbgruppe]{%
  \begin{aligned}[t]
    &\SemigroupAlg{B}{\diamond}\dsep X,Y\in\PowNE{B}\dsep{}\\[-2pt]
    &\forall x\in X\,\exists y\in Y\,
      \bigl(x\mathrel{\equiv^{\mathrm{tr}}_{B,\diamond}}y\bigr)
      \dsep{}\\[-2pt]
    &\forall y\in Y\,\exists x\in X\,
      \bigl(y\mathrel{\equiv^{\mathrm{tr}}_{B,\diamond}}x\bigr)
      \vdash
      X\mathrel{\equiv^{\mathrm{tr}}_{\PowNE{B},
        \diamond_{\mathcal P}}}Y
  \end{aligned}%
}{PowerSemigroupTranslationSupportCriterion}{%
  \DeltaRow{Mengen}{B\dsep X\dsep Y\dsep Z\dsep
    x\dsep y\dsep z\dsep w}
  \DeltaRow{Funktionen}{\diamond\dsep\diamond_{\mathcal P}}
}
\begin{tabproofsplitwide}
  \proofpartwideindR[Rechte Mengenprodukte]{%
    \begin{aligned}[t]
      &\SemigroupAlg{B}{\diamond}\dsep X,Y,Z\in\PowNE{B}\dsep{}\\[-2pt]
      &\forall x\in X\,\exists y\in Y\,
        \bigl(x\mathrel{\equiv^{\mathrm{tr}}_{B,\diamond}}y\bigr)
        \dsep{}\\[-2pt]
      &\forall y\in Y\,\exists x\in X\,
        \bigl(y\mathrel{\equiv^{\mathrm{tr}}_{B,\diamond}}x\bigr)
        \vdash X\diamond_{\mathcal P}Z=Y\diamond_{\mathcal P}Z
    \end{aligned}%
  }{%
    \begin{aligned}[t]
      &\SemigroupAlg{B}{\diamond}\dsep X,Y,Z\in\PowNE{B}\dsep{}\\[-2pt]
      &\forall x\in X\,\exists y\in Y\,
        \bigl(x\mathrel{\equiv^{\mathrm{tr}}_{B,\diamond}}y\bigr)
        \dsep{}\\[-2pt]
      &\forall y\in Y\,\exists x\in X\,
        \bigl(y\mathrel{\equiv^{\mathrm{tr}}_{B,\diamond}}x\bigr)
        \vdash X\diamond_{\mathcal P}Z=Y\diamond_{\mathcal P}Z
    \end{aligned}%
  }[PowerSemigroupTranslationSupportRightProduct]
    \proofstepwidestarbreak[1]{\SemigroupAlg{B}{\diamond}}{\rA}
    \proofstepwidestarbreak[2]{X\in\PowNE{B}}{\rA}
    \proofstepwidestarbreak[3]{Y\in\PowNE{B}}{\rA}
    \proofstepwidestarbreak[4]{Z\in\PowNE{B}}{\rA}
    \proofstepwidestarbreak[5]{\forall u\in X\,\exists v\in Y\,
      \bigl(u\mathrel{\equiv^{\mathrm{tr}}_{B,\diamond}}v\bigr)}{\rA}
    \proofstepwidestarbreak[6]{\forall v\in Y\,\exists u\in X\,
      \bigl(v\mathrel{\equiv^{\mathrm{tr}}_{B,\diamond}}u\bigr)}{\rA}

    \proofstepwidestarbreak[7]{w\in X\diamond_{\mathcal P}Z}{\rA}
    \proofstepwidestarbreak[1,2,4,7]{%
      \exists x\in X\,\exists z\in Z\,(w=x\diamond z)}{%
      \FormulaRefAuto{PowerSemigroupProductWitnesses}[thm]{1,2,4,7}}
    \proofstepwidestarbreak[8]{x\in X}{\rA}
    \proofstepwidestarbreak[8]{z\in Z}{\rA}
    \proofstepwidestarbreak[8]{w=x\diamond z}{\rA}
    \proofstepwidestarbreak[5,8]{\exists y\in Y\,
      \bigl(x\mathrel{\equiv^{\mathrm{tr}}_{B,\diamond}}y\bigr)}{%
      \rRE{9,\rUE{5}}}
    \proofstepwidestarbreak[12]{y\in Y}{\rA}
    \proofstepwidestarbreak[12]{%
      x\mathrel{\equiv^{\mathrm{tr}}_{B,\diamond}}y}{\rA}
    \proofstepwidestarbreak[2,8]{x\in B}{%
      \FormulaRefAuto{PowerSemigroupCarrierElement}[thm]{2,9}}
    \proofstepwidestarbreak[3,12]{y\in B}{%
      \FormulaRefAuto{PowerSemigroupCarrierElement}[thm]{3,13}}
    \proofstepwidestarbreak[4,8]{z\in B}{%
      \FormulaRefAuto{PowerSemigroupCarrierElement}[thm]{4,10}}
    \proofstepwidestar[1,2,3,4,8,12]{x\diamond z=y\diamond z\land z\diamond x=z\diamond y}{%
      \FormulaRefAuto{SemigroupTranslationIndistinguishableActions}[thm]{%
        1,15,16,14,17}}
    \proofstepwidestarbreak[1,2,3,4,8,12]{x\diamond z=y\diamond z}{%
      \rAEa{18}}
    \proofstepwidestarbreak[1,2,3,4,8,12]{w=y\diamond z}{\rChain{11,19}}
    \proofstepwidestarbreak[1,3,4,8,12]{%
      y\diamond z\in Y\diamond_{\mathcal P}Z}{%
      \FormulaRefAuto{PowerSemigroupProductContains}[thm]{1,3,4,13,10}}
    \proofstepwidestarbreak[1,2,3,4,8,12]{%
      w\in Y\diamond_{\mathcal P}Z}{\rIE{20,21}}
    \proofstepwidestarbreak[1,2,3,4,5,8]{%
      w\in Y\diamond_{\mathcal P}Z}{\rEEStar{12\text{--}22}}
    \proofstepwidestarbreak[1,2,3,4,5,7]{%
      w\in Y\diamond_{\mathcal P}Z}{\rEEStar{8\text{--}23}}
    \proofstepwidestarbreak[1,2,3,4,5]{%
      X\diamond_{\mathcal P}Z\subseteq Y\diamond_{\mathcal P}Z}{%
      \FormulaRefAuto{SubsetIntroductionRule}[thm]{[7]\vdots 24}}

    \proofstepwidestarbreak[26]{w\in Y\diamond_{\mathcal P}Z}{\rA}
    \proofstepwidestarbreak[1,3,4,26]{%
      \exists y\in Y\,\exists z\in Z\,(w=y\diamond z)}{%
      \FormulaRefAuto{PowerSemigroupProductWitnesses}[thm]{1,3,4,26}}
    \proofstepwidestarbreak[27]{y\in Y}{\rA}
    \proofstepwidestarbreak[27]{z\in Z}{\rA}
    \proofstepwidestarbreak[27]{w=y\diamond z}{\rA}
    \proofstepwidestarbreak[6,27]{\exists x\in X\,
      \bigl(y\mathrel{\equiv^{\mathrm{tr}}_{B,\diamond}}x\bigr)}{%
      \rRE{28,\rUE{6}}}
    \proofstepwidestarbreak[31]{x\in X}{\rA}
    \proofstepwidestarbreak[31]{%
      y\mathrel{\equiv^{\mathrm{tr}}_{B,\diamond}}x}{\rA}
    \proofstepwidestarbreak[3,27]{y\in B}{%
      \FormulaRefAuto{PowerSemigroupCarrierElement}[thm]{3,28}}
    \proofstepwidestarbreak[2,31]{x\in B}{%
      \FormulaRefAuto{PowerSemigroupCarrierElement}[thm]{2,32}}
    \proofstepwidestarbreak[4,27]{z\in B}{%
      \FormulaRefAuto{PowerSemigroupCarrierElement}[thm]{4,29}}
    \proofstepwidestar[1,2,3,4,27,31]{y\diamond z=x\diamond z\land z\diamond y=z\diamond x}{%
      \FormulaRefAuto{SemigroupTranslationIndistinguishableActions}[thm]{%
        1,34,35,33,36}}
    \proofstepwidestarbreak[1,2,3,4,27,31]{y\diamond z=x\diamond z}{%
      \rAEa{37}}
    \proofstepwidestarbreak[1,2,3,4,27,31]{w=x\diamond z}{\rChain{30,38}}
    \proofstepwidestarbreak[1,2,4,27,31]{%
      x\diamond z\in X\diamond_{\mathcal P}Z}{%
      \FormulaRefAuto{PowerSemigroupProductContains}[thm]{1,2,4,32,29}}
    \proofstepwidestarbreak[1,2,3,4,27,31]{%
      w\in X\diamond_{\mathcal P}Z}{\rIE{39,40}}
    \proofstepwidestarbreak[1,2,3,4,6,27]{%
      w\in X\diamond_{\mathcal P}Z}{\rEEStar{31\text{--}41}}
    \proofstepwidestarbreak[1,2,3,4,6,26]{%
      w\in X\diamond_{\mathcal P}Z}{\rEEStar{27\text{--}42}}
    \proofstepwidestarbreak[1,2,3,4,6]{%
      Y\diamond_{\mathcal P}Z\subseteq X\diamond_{\mathcal P}Z}{%
      \FormulaRefAuto{SubsetIntroductionRule}[thm]{[26]\vdots 43}}
    \proofstepwidestarbreak[1,2,3,4,5,6]{%
      X\diamond_{\mathcal P}Z=Y\diamond_{\mathcal P}Z}{%
      \FormulaRefAuto{A\subseteq B\dsep B\subseteq A\vdash A=B}[thm]{25,44}}
  \closeproofpartwideindR

  \Needspace{16\baselineskip}
  \proofpartwideindR[Linke Mengenprodukte]{%
    \begin{aligned}[t]
      &\SemigroupAlg{B}{\diamond}\dsep X,Y,Z\in\PowNE{B}\dsep{}\\[-2pt]
      &\forall x\in X\,\exists y\in Y\,
        \bigl(x\mathrel{\equiv^{\mathrm{tr}}_{B,\diamond}}y\bigr)
        \dsep{}\\[-2pt]
      &\forall y\in Y\,\exists x\in X\,
        \bigl(y\mathrel{\equiv^{\mathrm{tr}}_{B,\diamond}}x\bigr)
        \vdash Z\diamond_{\mathcal P}X=Z\diamond_{\mathcal P}Y
    \end{aligned}%
  }{%
    \begin{aligned}[t]
      &\SemigroupAlg{B}{\diamond}\dsep X,Y,Z\in\PowNE{B}\dsep{}\\[-2pt]
      &\forall x\in X\,\exists y\in Y\,
        \bigl(x\mathrel{\equiv^{\mathrm{tr}}_{B,\diamond}}y\bigr)
        \dsep{}\\[-2pt]
      &\forall y\in Y\,\exists x\in X\,
        \bigl(y\mathrel{\equiv^{\mathrm{tr}}_{B,\diamond}}x\bigr)
        \vdash Z\diamond_{\mathcal P}X=Z\diamond_{\mathcal P}Y
    \end{aligned}%
  }[PowerSemigroupTranslationSupportLeftProduct]
    \proofstepwidestarbreak[1]{\SemigroupAlg{B}{\diamond}}{\rA}
    \proofstepwidestarbreak[2]{X\in\PowNE{B}}{\rA}
    \proofstepwidestarbreak[3]{Y\in\PowNE{B}}{\rA}
    \proofstepwidestarbreak[4]{Z\in\PowNE{B}}{\rA}
    \proofstepwidestarbreak[5]{\forall u\in X\,\exists v\in Y\,
      \bigl(u\mathrel{\equiv^{\mathrm{tr}}_{B,\diamond}}v\bigr)}{\rA}
    \proofstepwidestarbreak[6]{\forall v\in Y\,\exists u\in X\,
      \bigl(v\mathrel{\equiv^{\mathrm{tr}}_{B,\diamond}}u\bigr)}{\rA}

    \proofstepwidestarbreak[7]{w\in Z\diamond_{\mathcal P}X}{\rA}
    \proofstepwidestarbreak[1,2,4,7]{%
      \exists z\in Z\,\exists x\in X\,(w=z\diamond x)}{%
      \FormulaRefAuto{PowerSemigroupProductWitnesses}[thm]{1,4,2,7}}
    \proofstepwidestarbreak[8]{z\in Z}{\rA}
    \proofstepwidestarbreak[8]{x\in X}{\rA}
    \proofstepwidestarbreak[8]{w=z\diamond x}{\rA}
    \proofstepwidestarbreak[5,8]{\exists y\in Y\,
      \bigl(x\mathrel{\equiv^{\mathrm{tr}}_{B,\diamond}}y\bigr)}{%
      \rRE{10,\rUE{5}}}
    \proofstepwidestarbreak[12]{y\in Y}{\rA}
    \proofstepwidestarbreak[12]{%
      x\mathrel{\equiv^{\mathrm{tr}}_{B,\diamond}}y}{\rA}
    \proofstepwidestarbreak[2,8]{x\in B}{%
      \FormulaRefAuto{PowerSemigroupCarrierElement}[thm]{2,10}}
    \proofstepwidestarbreak[3,12]{y\in B}{%
      \FormulaRefAuto{PowerSemigroupCarrierElement}[thm]{3,13}}
    \proofstepwidestarbreak[4,8]{z\in B}{%
      \FormulaRefAuto{PowerSemigroupCarrierElement}[thm]{4,9}}
    \proofstepwidestar[1,2,3,4,8,12]{x\diamond z=y\diamond z\land z\diamond x=z\diamond y}{%
      \FormulaRefAuto{SemigroupTranslationIndistinguishableActions}[thm]{%
        1,15,16,14,17}}
    \proofstepwidestarbreak[1,2,3,4,8,12]{z\diamond x=z\diamond y}{%
      \rAEb{18}}
    \proofstepwidestarbreak[1,2,3,4,8,12]{w=z\diamond y}{\rChain{11,19}}
    \proofstepwidestarbreak[1,3,4,8,12]{%
      z\diamond y\in Z\diamond_{\mathcal P}Y}{%
      \FormulaRefAuto{PowerSemigroupProductContains}[thm]{1,4,3,9,13}}
    \proofstepwidestarbreak[1,2,3,4,8,12]{%
      w\in Z\diamond_{\mathcal P}Y}{\rIE{20,21}}
    \proofstepwidestarbreak[1,2,3,4,5,8]{%
      w\in Z\diamond_{\mathcal P}Y}{\rEEStar{12\text{--}22}}
    \proofstepwidestarbreak[1,2,3,4,5,7]{%
      w\in Z\diamond_{\mathcal P}Y}{\rEEStar{8\text{--}23}}
    \proofstepwidestarbreak[1,2,3,4,5]{%
      Z\diamond_{\mathcal P}X\subseteq Z\diamond_{\mathcal P}Y}{%
      \FormulaRefAuto{SubsetIntroductionRule}[thm]{[7]\vdots 24}}

    \proofstepwidestarbreak[26]{w\in Z\diamond_{\mathcal P}Y}{\rA}
    \proofstepwidestarbreak[1,3,4,26]{%
      \exists z\in Z\,\exists y\in Y\,(w=z\diamond y)}{%
      \FormulaRefAuto{PowerSemigroupProductWitnesses}[thm]{1,4,3,26}}
    \proofstepwidestarbreak[27]{z\in Z}{\rA}
    \proofstepwidestarbreak[27]{y\in Y}{\rA}
    \proofstepwidestarbreak[27]{w=z\diamond y}{\rA}
    \proofstepwidestarbreak[6,27]{\exists x\in X\,
      \bigl(y\mathrel{\equiv^{\mathrm{tr}}_{B,\diamond}}x\bigr)}{%
      \rRE{29,\rUE{6}}}
    \proofstepwidestarbreak[31]{x\in X}{\rA}
    \proofstepwidestarbreak[31]{%
      y\mathrel{\equiv^{\mathrm{tr}}_{B,\diamond}}x}{\rA}
    \proofstepwidestarbreak[3,27]{y\in B}{%
      \FormulaRefAuto{PowerSemigroupCarrierElement}[thm]{3,29}}
    \proofstepwidestarbreak[2,31]{x\in B}{%
      \FormulaRefAuto{PowerSemigroupCarrierElement}[thm]{2,32}}
    \proofstepwidestarbreak[4,27]{z\in B}{%
      \FormulaRefAuto{PowerSemigroupCarrierElement}[thm]{4,28}}
    \proofstepwidestar[1,2,3,4,27,31]{y\diamond z=x\diamond z\land z\diamond y=z\diamond x}{%
      \FormulaRefAuto{SemigroupTranslationIndistinguishableActions}[thm]{%
        1,34,35,33,36}}
    \proofstepwidestarbreak[1,2,3,4,27,31]{z\diamond y=z\diamond x}{%
      \rAEb{37}}
    \proofstepwidestarbreak[1,2,3,4,27,31]{w=z\diamond x}{\rChain{30,38}}
    \proofstepwidestarbreak[1,2,4,27,31]{%
      z\diamond x\in Z\diamond_{\mathcal P}X}{%
      \FormulaRefAuto{PowerSemigroupProductContains}[thm]{1,4,2,28,32}}
    \proofstepwidestarbreak[1,2,3,4,27,31]{%
      w\in Z\diamond_{\mathcal P}X}{\rIE{39,40}}
    \proofstepwidestarbreak[1,2,3,4,6,27]{%
      w\in Z\diamond_{\mathcal P}X}{\rEEStar{31\text{--}41}}
    \proofstepwidestarbreak[1,2,3,4,6,26]{%
      w\in Z\diamond_{\mathcal P}X}{\rEEStar{27\text{--}42}}
    \proofstepwidestarbreak[1,2,3,4,6]{%
      Z\diamond_{\mathcal P}Y\subseteq Z\diamond_{\mathcal P}X}{%
      \FormulaRefAuto{SubsetIntroductionRule}[thm]{[26]\vdots 43}}
    \proofstepwidestarbreak[1,2,3,4,5,6]{%
      Z\diamond_{\mathcal P}X=Z\diamond_{\mathcal P}Y}{%
      \FormulaRefAuto{A\subseteq B\dsep B\subseteq A\vdash A=B}[thm]{25,44}}
  \closeproofpartwideindR

  \proofpartwide{Schluss}
    \proofstepwidestarbreak[1]{\SemigroupAlg{B}{\diamond}}{\rA}
    \proofstepwidestarbreak[2]{X\in\PowNE{B}}{\rA}
    \proofstepwidestarbreak[3]{Y\in\PowNE{B}}{\rA}
    \proofstepwidestarbreak[4]{\forall x\in X\,\exists y\in Y\,
      \bigl(x\mathrel{\equiv^{\mathrm{tr}}_{B,\diamond}}y\bigr)}{\rA}
    \proofstepwidestarbreak[5]{\forall y\in Y\,\exists x\in X\,
      \bigl(y\mathrel{\equiv^{\mathrm{tr}}_{B,\diamond}}x\bigr)}{\rA}
    \proofstepwidestarbreak[1]{%
      \SemigroupAlg{\PowNE{B}}{\diamond_{\mathcal P}}}{%
      \FormulaRefAuto{NonemptyPowerSemigroup}[thm]{1}}
    \proofstepwidestarbreak[7]{Z\in\PowNE{B}}{\rA}
    \proofstepwidestarbreak[1,2,3,4,5,7]{%
      X\diamond_{\mathcal P}Z=Y\diamond_{\mathcal P}Z}{%
      \FormulaRefAuto{PowerSemigroupTranslationSupportRightProduct}[thm]{%
        1,2,3,7,4,5}}
    \proofstepwidestarbreak[1,2,3,4,5,7]{%
      Z\diamond_{\mathcal P}X=Z\diamond_{\mathcal P}Y}{%
      \FormulaRefAuto{PowerSemigroupTranslationSupportLeftProduct}[thm]{%
        1,2,3,7,4,5}}
    \proofstepwidestarbreak[1,2,3,4,5,7]{%
      \begin{aligned}[t]
        &X\diamond_{\mathcal P}Z=Y\diamond_{\mathcal P}Z\land{}\\[-2pt]
        &Z\diamond_{\mathcal P}X=Z\diamond_{\mathcal P}Y
      \end{aligned}}{\rAI{8,9}}
    \proofstepwidestarbreak[1,2,3,4,5]{%
      \forall Z\in\PowNE{B}\,\left(
        \begin{aligned}[t]
          &X\diamond_{\mathcal P}Z=Y\diamond_{\mathcal P}Z\land{}\\[-2pt]
          &Z\diamond_{\mathcal P}X=Z\diamond_{\mathcal P}Y
        \end{aligned}\right)}{%
      \rUI{\rRI{7,10}}}
    \proofstepwidestarbreak[1,2,3,4,5]{%
      X\mathrel{\equiv^{\mathrm{tr}}_{\PowNE{B},
        \diamond_{\mathcal P}}}Y}{%
      \FormulaRefAuto{SemigroupTranslationIndistinguishableDef}[def]{%
        6,2,3,11}}
  \closeproofpartwide
\end{tabproofsplitwide}

\subsection{Ein unendliches Gegenbeispiel zur Umkehrung}

Ein Isomorphismus zweier Halbgruppen induziert einen Isomorphismus ihrer
Potenzhalbgruppen. Die Umkehrung gilt im Allgemeinen nicht, wie das
Gegenbeispiel von Mogiljanskaja zeigt.

Die Operationen \(\star_{\mathcal P}\) und \(\diamond_{\mathcal P}\)
bezeichnen dabei die aus \(\star\) und \(\diamond\) abgeleiteten
Mengenprodukte gemäß \FormulaRefAuto{PowerSemigroupProductDef}[def].

\Needspace{20\baselineskip}
\FormulaThmDeltaK[Das Mogiljanskaja-Gegenbeispiel]{%
  \begin{aligned}[t]
    &\exists S,T\,\exists\star,\diamond\;\Bigl(\\[-2pt]
    &\quad\SemigroupAlg{S}{\star}\land
      \SemigroupAlg{T}{\diamond}\land{}\\[-2pt]
    &\quad\neg\Finite{S}\land\neg\Finite{T}\land{}\\[-2pt]
    &\quad\neg\exists f\,
      \SemigroupIso{f}{S,\star}{T,\diamond}\land{}\\[-2pt]
    &\quad\exists F\,\SemigroupIso{F}
      {\PowNE{S},\star_{\mathcal P}}
      {\PowNE{T},\diamond_{\mathcal P}}\Bigr)
  \end{aligned}
}{MogiljanskajaPairCounterexample}{%
  \DeltaRow{Mengen}{S\dsep T}
  \DeltaRow{Funktionen}{\star\dsep\diamond\dsep f\dsep F}
}

Es gibt also zwei unendliche, nicht isomorphe Halbgruppen, deren
Potenzhalbgruppen der nichtleeren Teilmengen isomorph sind. Die
Potenzhalbgruppe bestimmt ihre Grundhalbgruppe somit im Allgemeinen
nicht bis auf Isomorphie. Über den entsprechenden endlichen Fall
entscheidet dieses Gegenbeispiel nicht.

\noindent\emph{Beweis:} Die \MogReadingLink{} enthält die Konstruktion
und den vollständigen erklärenden Beweis. Die \MogProofsLink{} führen
die Definitionen, Hilfsaussagen und Ableitungen aus; ihr
\href{\MogEditionDirectory_B28-mog-proofs.pdf\#mog.statement.MogiljanskajaConstructedPairCounterexample}%
  {konkretes Gegenbeispiel}
liefert die Zeugen des Existenzsatzes.

\section{Produktquadrate mit höchstens zwei Werten}
Die Potenzmenge des Produktquadrats darf im Allgemeinen nicht mit dem
Produktquadrat der Potenzhalbgruppe gleichgesetzt werden. Bei einer
Zweiermenge reichen jedoch die Einermengen und der ganze Träger aus.

\Needspace{14\baselineskip}
\FormulaThmDeltaK[Einermengen und Gesamtmenge im Potenzquadrat]{%
\begin{aligned}[t]
&\SemigroupAlg{A}{\star}\dsep A\neq\varnothing\vdash{}\\[-2pt]
&A^2\neq\varnothing\land\PowNE{A}^{\,2}\subseteq\PowNE{A^2}\\[-2pt]
&\qquad{}\land A^2\in\PowNE{A}^{\,2}
 \land\forall a\in A^2\,(\{a\}\in\PowNE{A}^{\,2})
\end{aligned}
}{SemigroupSquarePowerFamily}{%
  \DeltaRow{Mengen}{A\dsep a}
  \DeltaRow{Funktionen}{\star\dsep\star_{\mathcal P}}
}
\begin{tabproofwide}
  \proofstepwidestar[1]{\SemigroupAlg{A}{\star}}{\rA}
  \proofstepwidestar[2]{A\neq\varnothing}{\rA}
  \proofstepwidestar[2]{A\in\PowNE{A}}{\FormulaRefAuto{NonemptyPowersetMembership}[thm]{\rAI{\FormulaRefAuto{A\subseteq A}[thm],2}}}
  \proofstepwidestar[1,2]{A^2=A\star_{\mathcal P}A}{\FormulaRefAuto{SemigroupSquareDef}[def]{1,2}}
  \proofstepwidestar[1,2]{A\star_{\mathcal P}A\neq\varnothing}{\FormulaRefAuto{PowerSemigroupProductNonempty}[thm]{1,3,3}}
  \proofstepwidestar[1,2]{A^2\neq\varnothing}{\rIE{\FormulaRefAuto{a=b\vdash b=a}[thm]{4},5}}
  \proofstepwidestar[1,2]{\PowNE{A}^{\,2}\subseteq\PowNE{A^2}}{\FormulaRefAuto{PowerSemigroupSquareSubset}[thm]{1,2}}
  \proofstepwidestar[1]{\SemigroupAlg{\PowNE{A}}{\star_{\mathcal P}}}{\FormulaRefAuto{NonemptyPowerSemigroup}[thm]{1}}
  \proofstepwidestar[2]{\PowNE{A}\neq\varnothing}{\FormulaRefAuto{NonemptyPowerCarrierNonempty}[thm]{2}}
  \proofstepwidestar[1,2]{A\star_{\mathcal P}A\in\PowNE{A}^{\,2}}{\FormulaRefAuto{SemigroupSquareProductMember}[thm]{8,9,\rAI{3,3}}}
  \proofstepwidestar[1,2]{A^2\in\PowNE{A}^{\,2}}{\rIE{\FormulaRefAuto{a=b\vdash b=a}[thm]{4},10}}
  \proofstepwidestar[12]{a\in A^2}{\rA}
  \proofstepwidestar[1,2]{A^2\subseteq A}{\FormulaRefAuto{SemigroupSquareSubset}[thm]{1,2}}
  \proofstepwidestar[1,2,12]{a\in A}{\FormulaRefAuto{A\subseteq B\dsep x\in A\vdash x\in B}[thm]{13,12}}
  \proofstepwidestar[1,2,12]{\{a\}\in\PowNE{A}^{\,2}}{\rRE{\rLREb{\FormulaRefAuto{PowerSemigroupSquareSingletonCriterion}[thm]{1,2,14}},12}}
  \proofstepwidestar[1,2]{\forall a\in A^2\,(\{a\}\in\PowNE{A}^{\,2})}{\rUI{\rRI{12,15}}}
  \proofstepwidestar[1,2]{\begin{gathered}A^2\neq\varnothing\land\PowNE{A}^{\,2}\subseteq\PowNE{A^2}\\{}\land A^2\in\PowNE{A}^{\,2}\\
 {}\land\forall a\in A^2\,(\{a\}\in\PowNE{A}^{\,2})\end{gathered}}{\rAI{6,\rAI{7,\rAI{11,16}}}}
\end{tabproofwide}
\Needspace{14\baselineskip}
\FormulaThmDeltaK[Bei einem zweielementigen Quadrat gibt es keine fehlenden Teilmengen]{%
\begin{aligned}[t]
&\SemigroupAlg{A}{\star}\dsep A\neq\varnothing\dsep A^2=\{a,b\}\\[-2pt]
&\vdash\PowNE{A}^{\,2}=\PowNE{A^2}
\end{aligned}
}{PairProductSquarePowerEquality}{%
  \DeltaRow{Mengen}{A\dsep a\dsep b\dsep X}
  \DeltaRow{Funktionen}{\star\dsep\star_{\mathcal P}}
}
Der Beweis gilt auch für \(a=b\). In der Tabelle steht
\(D=A^2\) und \(K=\PowNE{A}^{\,2}\); dies sind ausschließlich Abkürzungen.

\begin{tabproofwide}
  \proofstepwidestar[1]{\SemigroupAlg{A}{\star}}{\rA}
  \proofstepwidestar[2]{A\neq\varnothing}{\rA}
  \proofstepwidestar[3]{D=\{a,b\}}{\rA}
  \proofstepwidestar[1,2]{\begin{gathered}D\neq\varnothing\land K\subseteq\PowNE{D}\\{}\land D\in K\land\forall d\in D\,(\{d\}\in K)\end{gathered}}{\FormulaRefAuto{SemigroupSquarePowerFamily}[thm]{1,2}}
  \proofstepwidestar[3]{a\in D}{\rIE{\FormulaRefAuto{a=b\vdash b=a}[thm]{3},\FormulaRefAuto{a\in\{a,b\}}[thm]}}
  \proofstepwidestar[3]{b\in D}{\rIE{\FormulaRefAuto{a=b\vdash b=a}[thm]{3},\FormulaRefAuto{b\in\{a,b\}}[thm]}}
  \proofstepwidestar[1,2,3]{\{a\}\in K}{\rRE{5,\rUE{\rAEn{4}}}}
  \proofstepwidestar[1,2,3]{\{b\}\in K}{\rRE{6,\rUE{\rAEn{4}}}}
  \proofstepwidestar[1,2,3]{\{a,b\}\in K}{\rIE{3,\rAEn{4}}}
  \proofstepwidestar[10]{X\in\PowNE{D}}{\rA}
  \proofstepwidestar[3,10]{X\in\PowNE{\{a,b\}}}{\rIE{3,10}}
  \proofstepwidestar[3,10]{X\in\bigl\{\{a\},\{b\},\{a,b\}\bigr\}}{\rIE{\FormulaRefAuto{NonemptyPowersetPair}[thm],11}}
  \proofstepwidestar[3,10]{X=\{a\}\lor(X=\{b\}\lor X=\{a,b\})}{\FormulaRefAuto{TripleSetMembership}[thm]{12}}
  \proofstepwidestar[14]{X=\{a\}}{\rA}
  \proofstepwidestar[1,2,3,14]{X\in K}{\rIE{\FormulaRefAuto{a=b\vdash b=a}[thm]{14},7}}
  \proofstepwidestar[16]{X=\{b\}\lor X=\{a,b\}}{\rA}
  \proofstepwidestar[17]{X=\{b\}}{\rA}
  \proofstepwidestar[1,2,3,17]{X\in K}{\rIE{\FormulaRefAuto{a=b\vdash b=a}[thm]{17},8}}
  \proofstepwidestar[19]{X=\{a,b\}}{\rA}
  \proofstepwidestar[1,2,3,19]{X\in K}{\rIE{\FormulaRefAuto{a=b\vdash b=a}[thm]{19},9}}
  \proofstepwidestar[1,2,3,16]{X\in K}{\rOE{16,17,18,19,20}}
  \proofstepwidestar[1,2,3,10]{X\in K}{\rOE{13,14,15,16,21}}
  \nopagebreak
  \proofstepwidestar[1,2,3]{\PowNE{D}\subseteq K}{\FormulaRefAuto{SubsetIntroductionRule}[thm]{[10]\vdots 22}}
  \nopagebreak
  \proofstepwidestar[1,2,3]{K=\PowNE{D}}{\FormulaRefAuto{A\subseteq B,\, B\subseteq A\vdash A=B}[thm]{\rAEn{4},23}}
\end{tabproofwide}


\ProofWideReasonsAtRight

\end{document}
